% This LaTeX document needs to be compiled with XeLaTeX.
\documentclass[10pt]{article}
\usepackage[utf8]{inputenc}
\usepackage{graphicx}
\usepackage[export]{adjustbox}
\graphicspath{ {./images/} }
\usepackage{amsmath}
\usepackage{amsfonts}
\usepackage{amssymb}
\usepackage[version=4]{mhchem}
\usepackage{extpfeil}
\usepackage{stmaryrd}
\usepackage{bbold}
\usepackage{fvextra, csquotes}
\usepackage[fallback]{xeCJK}
\usepackage{polyglossia}
\usepackage{fontspec}
\IfFontExistsTF{Noto Serif CJK JP}
{\setCJKmainfont{Noto Serif CJK JP}}
{\IfFontExistsTF{STSong}
  {\setCJKmainfont{STSong}}
  {\IfFontExistsTF{Droid Sans Fallback}
    {\setCJKmainfont{Droid Sans Fallback}}
    {\setCJKmainfont{SimSun}}
}}

\setmainlanguage{english}
\IfFontExistsTF{CMU Serif}
{\setmainfont{CMU Serif}}
{\IfFontExistsTF{DejaVu Sans}
  {\setmainfont{DejaVu Sans}}
  {\setmainfont{Georgia}}
}

%New command to display footnote whose markers will always be hidden
\let\svthefootnote\thefootnote
\newcommand\blfootnotetext[1]{%
  \let\thefootnote\relax\footnote{#1}%
  \addtocounter{footnote}{-1}%
  \let\thefootnote\svthefootnote%
}

%Overriding the \footnotetext command to hide the marker if its value is `0`
\let\svfootnotetext\footnotetext
\renewcommand\footnotetext[2][?]{%
  \if\relax#1\relax%
    \ifnum\value{footnote}=0\blfootnotetext{#2}\else\svfootnotetext{#2}\fi%
  \else%
    \if?#1\ifnum\value{footnote}=0\blfootnotetext{#2}\else\svfootnotetext{#2}\fi%
    \else\svfootnotetext[#1]{#2}\fi%
  \fi
}

\begin{document}
\begin{center}
\includegraphics[max width=\textwidth]{2025_02_16_7c9aaf92bc99dc07f048g-001(1)}
\end{center}

\section*{数学思维－应试方法}
作者：Jonjo\\
时间：October 19， 2024\\
原作者：阳哥／阿不\\
\includegraphics[max width=\textwidth, center]{2025_02_16_7c9aaf92bc99dc07f048g-001}\\
第1章 高中数学经验分享 ..... 2\\
1.1 动静视角 ..... 2\\
1.2 充要替换 ..... 12\\
1.3 即时检查 ..... 16\\
1.4 结构意识 ..... 19\\
1.5 考场习惯 ..... 27\\
1.6 封装命题 ..... 36\\
第2章 高中解题思想与杂谈 ..... 47\\
2.1 辅助思想 ..... 47\\
2.2 具体方法 ..... 47\\
2.3 阶段性总结 ..... 47\\
2.4 好题分享 ..... 47\\
2.5 漫游难题 ..... 47\\
2.6 数学联赛 ..... 47\\
2.7 吹水／杂谈 ..... 47\\
第3章 计算基本法—杜杆 ..... 48\\
3.1 主元思想 ..... 48\\
3.2 对称式 ..... 51\\
3.3 齐次式 ..... 54\\
3.4 共轭根式 ..... 57\\
3.5 自由度，消元 ..... 60\\
3.6 因式分解 ..... 64\\
3.7 二次式 ..... 70\\
第4章 计算基本法—实战 ..... 74\\
4.1 函数的性质 ..... 74\\
4.2 不等式 ..... 81\\
4.3 三角函数，解三角形 ..... 84\\
4.4 立体几何 ..... 90\\
4.5 数列 ..... 94\\
4.6 概率与统计 ..... 99\\
4.7 导数 ..... 103\\
4.8 圆锥曲线 ..... 109\\
第5章 选填破题基本法 ..... 127\\
5.1 思想试验 ..... 127\\
5.2 遍历与临界 ..... 136\\
5.3 再论自由度 ..... 148\\
5.4 动态图像 ..... 154\\
5.5 数形结合 ..... 160\\
5.6 基底思想 ..... 167\\
5.7 分类讨论 ..... 170\\
5.8 充要与化归 ..... 175\\
5.9 证明方法 ..... 191\\
第6章 科学的思维方式解题 ..... 197\\
6.1 排列组合与计数问题 ..... 197\\
6.2 数列难题 ..... 202\\
6.3 三角函数难题 ..... 219\\
6.4 新定义问题 ..... 235\\
6.5 圆锥曲线 ..... 245\\
6.6 函数与导数 ..... 258\\
6.7 立体几何与概率 ..... 275\\
6.8 不等式问题 ..... 278\\
6.9 综合题目的思维练习 ..... 282\\
第7章 套卷相关的综合讲解 ..... 296\\
7.1 2023年全国甲卷 ..... 296\\
7.2 2023年新高考 I 卷 ..... 323\\
7.32024 年新高考 I 卷 ..... 324\\
7.4 2023 年四省联考 ..... 325\\
7.5 阳哥高质量原创题 ..... 326\\
7.62022 年新高考 I 卷 ..... 328\\
第8章 大学数学科普系列 ..... 330\\
8.1 数学的基础 ..... 330\\
8.2 重要的概念 ..... 337\\
8.3 联系万物 ..... 337\\
8.4 优美的仿射几何 ..... 338\\
8.5 实数理论 ..... 338

\section*{第1章 高中数学经验分享}
高中数学想要考到高分，有三个重要的因素。前两个是思维方式与逻辑能力，他们都是为了帮你找出解题方法，正确解出题目的。本讲义的内容重点强调稳定性，也就是尽量快，尽量准地找出，能在考场上稳定解出题目，拿到分数的办法。这是数学硬实力的重要方面，降低失误率，确保获取尽量多的，自我能力范围内的分数，这是应试能力的重要方面。

本章内容是围绕 up 主 nichijou 的涂山苏苏总结出的一套思维方法的有关思维方式与逻辑能力的讲解。其中，前五节是最重要的关于高考数学思维与逻辑的内容，包含自由度，要素关系网，充要替换等。建议所有来学习高中数学的读者都认真阅读一下。其中，第 \(1,2,4\) 节讲解的是他思考高中数学问题的思维方式与逻辑判断；第 3 节与第 5 节则是他对临场应试的一些基本的理解与经验，包含正确率，答题速度，时间规划等方面。尤其是第五节还涉及到如何学习掌握 up 主讲的思维方式的学习理念。因此，这五节内容是未来所有内容的思想基础，而章节＜科学的思维方式解题＞就是对本章内容的直接应用，可以同时学习，有助于理解。

\section*{1.1 动静视角}
读者可以先思考一下我们说的这个视角是指什么，你自己如何看待高中数学中的视角，而在我们正式开始学习视角之前，请先回答以下这个问题。

\section*{1．1．1 这个世界是死的吗？}
这个问题虽然听上去非常中二，但是问题中所说的世界，并不是一般意义上的世界，而是在一个数学题目中的世界。

\section*{1．1．1．1 搭积木}
为了更好解释这个问题，我们来玩一玩这个玩具。\\
例1．1以下是两种已知条件，请描述出看到这些条件时的感觉。\\
（1）已知 \(\triangle A B C, A B=3, B C=4, A C=5\) ．\\
（2）已知 \(\triangle A B C, A B=3, B C=4\) ．\\
先来看条件（1），我相信很多初中生都能立马反应过来：这个三角形是一个直角三角形。所以这道题无论条件后面的问题是什么，只需根据条件本身，就能在大脑中迅速的反应出来一个图景。在你脑子里可能有这样一个世界，这个世界非常单调，中间孤零零的屹立着这样的一个三角形，而我所说的世界就是这样的一个世界，他是指数学上的系统描述出的一个世界。

在一道数学题中，他会给你各种各样的东西，比如一个三角形，函数，点等等，我们把这些东西称为要素，各种各样的要素构成了这样的一个世界，要素之间还有彼此的一些关系，而他们本身也具有一些性质，比如这个要素就是一个三角形，而理论上这个三角形的三边和三个角我们都能知道，那么我们如何把它构建出来呢？比较自然的想法就是拿三个棍子形状的积木，去拼出一个三角形，那么我们要怎样去拼呢？

首先，找一个接口将长度为 3 和 4 的积木连接在一起，充当两个积木的关节，再把长度为 4 的积木固定在\\
\includegraphics[max width=\textwidth, center]{2025_02_16_7c9aaf92bc99dc07f048g-004(1)}\\
（1）卡死的积木

地面上，用长度为 3 的积木绕着关节在一个平面内旋转，最后再把长度为 5 的积木搭在长度为 3 和 4 的积木上。我们知道三角形是具有稳定性的，那这样一搭的话，这个图形是不是就被确定下来了？没错，因此我们可以得出结论：（1）SSS 判定三角形全等；（2）三角形不能活动；（3）这个系统是死的！\\
\includegraphics[max width=\textwidth, center]{2025_02_16_7c9aaf92bc99dc07f048g-004(4)}

再来看条件（2），这不就是我们搭积木游戏的前面几步：在没有长度为 5 的积木的情况下，长度为 3 的积木可以绕着关节在一个平面内随意旋转。也就是说，这个 \(A B\) 是有活动余地的，它可以绕着点 \(B\) 转动，此时这个世界就比上边卡死的世界多了一份生机，所以说这个世界是活的，这就形成了一个对应。

此时大家应该能稍微感觉到，这些系统都是分死活的。我们再举出一个复杂点的例子，用来加深这种感觉。\\
例1．2 四棱雉 \(P-A B C D\) ，（1）四边形 \(A B C D\) 是菱形，边长为1．（2）\(\angle A B C=60^{\circ}\) ．（3）\(P B \perp\) 面 \(A B C D, P B=1\) ．\\
（1）若：\(P B\) 长度不知？\\
（2）若：\(\angle A B C\) 大小不知？\\
\includegraphics[max width=\textwidth, center]{2025_02_16_7c9aaf92bc99dc07f048g-004(2)}\\
（1）可挤压\\
\includegraphics[max width=\textwidth, center]{2025_02_16_7c9aaf92bc99dc07f048g-004}\\
（1）＋（2）不动了！\\
\includegraphics[max width=\textwidth, center]{2025_02_16_7c9aaf92bc99dc07f048g-004(3)}\\
（1）+ （2）+ （3）高度确定，底面不动

这是一个立体几何的例子，当我们读完条件（1）时，我们的脑子里同样去构建一个世界：首先有一个四棱雉的底座，他是边长为 1 的菱形，同时他是有变化余地的，因为我们并不能确定它的形状。我们依旧把图形的边想象成棍子形状的积木，那这个菱形就是由四个边长为 1 的积木搭成的，相邻的两个积木间都有一个接口，同样充当它们的关节，然后把做好的这个小玩具放在桌面上，往下挤一挤再往外拉一拉。在挤压的过程中，大家就可以看到它就是一个动态的物体。但现在我又加了一个条件（2）\(\angle A B C=60^{\circ}\) 。这个时候你就不能随便挤了，只能保持两角为 \(60^{\circ}\) 的状态，所以说这个底座就不动了。然后呢？我们现在再给这个底座加个高怎么样？果然，条件（3）就是如此，我们找一根长度为 1 的旗杆，插到点 \(B\) 上，同时保证跟底座垂直，此时发现这是一根位置和长度都固定的旗杆，这个世界也是死的，这个系统没有活动的余地了。我们把 \(P A, ~ P C, ~ P D\) 分别连接之后构建成的四棱雉依旧没有活动的余地，如果把这个玩具给一个小孩子玩，只要质量够好，无论那小孩怎么摔，玩具都不会形变，也就是它具有稳定性。

当（1）成立时，这个旗杆就是个可伸缩的旗杆，虽然底座不动，但是这个锥点可以上下移动。这就类似在宿舍安装床帘的时候，需要拿一个铁杆撑一下，此时铁杆的高度就可以根据天花板的高度来调节。当（2）成立时，\\
\includegraphics[max width=\textwidth, center]{2025_02_16_7c9aaf92bc99dc07f048g-005(1)}\\
（1）不动的底座，跑动的锥点\\
\includegraphics[max width=\textwidth, center]{2025_02_16_7c9aaf92bc99dc07f048g-005(3)}\\
（2）不动的旗杆，可压可拉的底座

这个旗杆就是个固定的旗杆，但这个底座不老实，可以像之前那样挤一挤拉一拉，它就会发生形变。以上两种情况就可以说这个世界（系统）是活的！

\section*{1．1．1．2 可视化的参数}
刚才举了一些几何上的例子，现在我们再举一个函数上的例子。

例1．3已知函数（1）\(y=\ln x\) ；（2）\(y=a \ln x \quad(a>0)\)\\
\includegraphics[max width=\textwidth, center]{2025_02_16_7c9aaf92bc99dc07f048g-005(2)}\\
（1）这是一个死寂的世界！\\
\includegraphics[max width=\textwidth, center]{2025_02_16_7c9aaf92bc99dc07f048g-005}\\
（2）含有参数的函数，本质是活的！

注意函数（1）\(y=\ln x\) ，它就死死的躺在这个地方，一动不动，你不可能翟它一下让它变个形，因为那就不是 \(y=\ln x\) 了。当我们用软件 desmos 绘制（2）\(y=a \ln x \quad(a>0)\) 的图像的时，通过调节 \(a\) 的大小可以发现其对应的图像整个都运动起来了。这种情况的本质就是含参数的函数是活的！这个函数只是个单参数的函数，而在导数题里可能还会遇到双参数之类的问题，其实是一样的，都是这种运动的过程。

注 参数在背后操纵着运动的过程。有了运动，便有了生机！例如：令例1．1中的 \(\angle A B C=\alpha, \alpha\) 变化让系统发生了变化。又或：令 \(A C=b, b\) 也可以操纵系统的变化（积木的变化）。

总结\\
一道数学题总有一些背景，由此构建了一个系统（世界），系统受到一系列参数（边长，角度，系数……）的控制，一些参数可以影响到另一些参数。比如面积，它就和边长有：如果给你一个正方形，边长和面积各是一个参数，而此时的面积等于边长的平方，这就是参数之间的一些影响。

那么题目的已知，本质来讲是什么东西呢？它是阐述各个参数的基本信息（位置，数量关系。基本性质……）的。当然，还有它们的基本性质，要素之间的关系等等，这就是题目的已知的含义。

我们刚才讲了一些死活的问题，如果题干信息足以将系统固化（积木卡死，毫无变化余地），则称该系统是死的（静止的）。就像例1．2那样，用条件（1）＋（2）＋（3）构建的图像（四棱雉），完全没有变化的余地，无论如何

都不会形变，还是那个样子。\\
那么反过来呢？如果加上题干信息仍让系统有活动余地，则称此系统是活动的（运动的）。这些就是我们对题目的一个新的分类方式：按照这种分类标准，题目就只有死的和活的两种情况。

\section*{1．1．2 自由度初步}
\section*{1．1．2．1 如何归类高中数学题目？}
本节我们会放出一套高考试卷用于归纳高中数学题目。而在此之前，我们先思考一个问题：如何判断一大堆条件是动是静？我们有以下三个方法：

方法一：脑力搭积木，判断几参数。如之前的几何问题。\\
方法二：直接计算法，通常读完题就能直接完成。此处我们引入一个概念——自由度（大学物理学和数学中经常涉及）：一般的，有 \(n\) 个参数（自由度）的时候，用 \(n\) 个方程（等量关系）就可以解出这 \(n\) 个未知数的值；因此，每有一个方程存在，本质上可以减少一个参数的存在。如：题目中有 \(a, b\) 两个参数（两个自由度），又有已知 \(a+b=1\) ，于是有 \(b=1-a\)（ \(b\) 被 \(a\) 表示了），实质上的参数只剩下 \(a(b\) 被 \(1-a\) 代替了）一个实际的参数。所以未知数（参数）个数减去方程（等式）个数就等于实际参数个数，即实际自由度个数。

在一个题目里，如果抛开已知的等式，如何判断自由度个数？举一个例子，已知植圆 \(C: \frac{x^{2}}{a^{2}}+\frac{y^{2}}{b^{2}}=1\) ，我们先不看后续条件里的等式，在椭圆里有两个变化的参数，这两个参数可以分别在两个维度进行运动，所以本质上就是有两个自由度；但如果题干中给了一些信息，比如离心率是 \(\frac{1}{2}\) ，此时这个椭圆就要减少一个自由度，即只剩下大小这个自由度了。没错，大小可以看作一个自由度，这里可以理解为只有形状这个自由度确定，而大小仍在变化。如果再来一个等量关系，比如长轴长度，短轴长度，两个焦点之间的距离或者上顶点和两个焦点之间围成三角形面积等等，总之就是随便给你一个等式，它唯一的自由度就被减掉了，也就是说它就静止了。

所以我们可以通过自由度来判断某一个题目是运动的还是静止的，其中运动的题目我们也可以判断出来它有多少个自由度来。上述过程，读完题就可以直接完成，毕竟这些都是题目中的显性的已知条件。

方法三：结合问号。依旧是已知椭圆 \(C: \frac{x^{2}}{a^{2}}+\frac{y^{2}}{b^{2}}=1\) ，这个时候题目中有两个自由度，然后给你一些条件，最后给出问号：＂栯圆的标准方程是什么？＂。这个问号能说明什么呢？这个椭圆到最后肯定会静止下来，即自由度减二，可以根据减去的自由度个数反向理解后续条件。上述就是结合问号的方法，虽然听上去挺巧妙，但不能算是那种正统的逻辑上的方法，所以这是一个邪法。

现在我们正式开始高考数学之旅：\\
例1．4「2020年全国卷理科数学」\\
1．若 \(z=1+i\) ，则 \(\left|z^{2}-2 z\right|=\)\\
A， 0\\
B， 1\\
C，\(\sqrt{2}\)\\
D， 2

很明显 z 是静止的，没有任何变化余地。\\
2．设集合 \(A=\left\{x \mid x^{2}-4 \leq 0\right\}, B=\{x \mid 2 x+a \leq 0\}, A \cap B=\{x \mid-2 \leq x \leq 1\}, a=\)\\
A，-4\\
B，-2\\
C， 2\\
D， 4

集合 \(A\) 是静止的，同样没有变化余地；但集合 \(B\) 不同，\(y=2 x+a\) 的图像会随着 \(a\) 的变化而变化，它是一个单自由度，而 \(A \cap B\) 中 \(x \leq 1\) 就把 \(B\) 给兒住了，这个 1 就相当于一个等量关系，\(B\) 只能运动到 1 这个位置，从而确定了 \(a\) 的值，本质就是把运动的参数静止化。

3．埃及胡夫金字塔是古代世界建筑奇迹之一，它的形状可视为一个正四棱雉，以该四棱雉的高为边长的正方形面积等于该四校锥一个侧面三角形的面积，则其侧面三角形底边上的高与底面正方形的边长的比值为\\
A，\(\frac{\sqrt{5}-1}{4}\)\\
B，\(\frac{\sqrt{5}-1}{2}\)\\
C，\(\frac{\sqrt{5}+1}{4}\)\\
D，\(\frac{\sqrt{5}+1}{2}\)

之前提到过的＂大小也是一个自由度＂，比如三角形的三边长，本身有三个自由度，因为 SSS，ASA，SAS…这些全都是三个条件，也就是说三个参数可以决定一个三角形。如果三角形三边长是 \(3: 4: 5\) ，本质上就是两个等式，按照之前的法二的结论就是减少了两个自由度，只剩一个自由度了，那就是大小。那么这个时候我任给一个长度，即这个三角形的任一边的长度或者面积，这个形状理论上就可以确定下来了。\\
言归正传，这个胡夫金字塔是一个正四棱锥，它是几个自由度的东西呢？明显是两个对吧：底座的大小和四棱锥的高度；而题目中的条件＂一个面积等于另一个面积＂，本质就是一个比例；根据已知条件减去一个自由度后，这个四棱锥只剩下唯一自由度，那就是大小，此时你应该能反应出来这个四棱锥的形状被确定了！最后它的问题不出所料，就是边长比值，和大小无关。

4．已知 \(A\) 为抛物线 \(C: y^{2}=2 p x(p>0)\) 上一点，点 \(A\) 到 \(C\) 的焦点距离为 12 ，到 \(y\) 轴的距离为 9 ，则 \(p=\)\\
A， 2\\
B， 3\\
C， 6\\
D， 9

很明显 \(p\) 是一个自由度，它唯一控制着曲线的运动，而在这个自由度的基础上，点 \(A\) 也包含一个自由度，可以是它的横坐标，因为横坐标确定之后点也被确定了。但由于曲线的横坐标一般有范围限制，所以它最多可以确定有限个点。因此，条件＂\(A\) 为抛物线 \(C: y^{2}=2 p x(p>0)\) 上一点＂包含了两个自由度：曲线的自由度和点的自由度。后面的条件本质就是两个方程，两个自由度就减没了，最终就是让你求静止后 \(p\) 是多少。虽然这类题目是很简单的，画个图就能看出来，但其实我是想让你知道这类题目也是我们动静视角体系下的东西，所以也是值得一提的。

5．某校一个课外学习小组为研究某作物种子的发芽率 \(y\) 和温度 \(x\)（单位：\({ }^{\circ} C\) ）的关系，在 20 个不同的温度条件下进行种子发芽实验，由实验数据 \(\left(x_{i}, y_{i}\right)(i=1,2, \cdots, 20)\) 得到下面的散点图：（略）\\
由此散点图，在 \(10^{\circ} \mathrm{C}\) 至 \(40^{\circ} \mathrm{C}\) 之间，下面四个回归方程类型中最适宜作为发芽率 \(y\) 和温度 \(x\) 的回归方程类型的是\\
A，\(y=a+b x\)\\
B ，\(y=a+b x^{2}\)\\
C,\(y=a+b e^{x}\)\\
D，\(y=a+b \ln x\)

这道题严格来讲不能归到我们的动静视角体系里面，因为它主要就只考你基本的概念，不过后续章节也会介绍关于这类问题的同样本质的方法。

6．函数 \(f(x)=x^{4}-2 x^{3}\) 的图像在点 \((1, f(1))\) 处的切线方程为\\
A，\(y=-2 x-1\)\\
B ，\(y=-2 x+1\)\\
C，\(y=2 x-3\)\\
D，\(y=2 x+1\)

函数 \(f(x)\) 的表达式告诉你了，这是一个死的，死的东西它问什么问题都不用奇怪，因为它的任何信息理论上都是可以被我们知道的。虽然现在有些世界未解之谜可能还是跟一些静止的东西有关，其实就只是涉及的东西

比较多，如果开个＂上帝视角＂肯定都可以知道。而像我们这种高考题，严格来讲，它问的都是可以直接用从它里面挖掘出的信息来解决的。

7．设函数 \(f(x)=\cos \left(\omega x+\frac{\pi}{6}\right)\) 在 \([-\pi, \pi]\) 的图像大致如图所示则 \(f(x)\) 的最小正周期为\\
A，\(\frac{10 \pi}{9}\)\\
B ，\(\frac{7 \pi}{6}\)\\
C，\(\frac{4 \pi}{3}\)\\
D，\(\frac{3 \pi}{2}\)\\
\includegraphics[max width=\textwidth, center]{2025_02_16_7c9aaf92bc99dc07f048g-008}

第 7 题图\\
\includegraphics[max width=\textwidth, center]{2025_02_16_7c9aaf92bc99dc07f048g-008(1)}

第 16 题图

图像问题。 \(f(x)\) 只包含一个自由度 \(\omega\) ，而这个图像中比较有用的信息只有它的零点 \(-\frac{4 \pi}{9}\) ，也是题目中唯一的等量关系。结合给的图像上单调性的一些基本信息和这个等量关系足够确定 \(f(x)\) 。

8．\(\left(x+\frac{y^{2}}{x}\right)(x+y)^{5}\) 的展开式中 \(x^{3} y^{3}\) 的系数为\\
A， 5\\
B ， 10\\
C， 15\\
D， 20

静止的东西，同样不用奇怪后面的问题。

9．已知 \(\alpha \in(0, \pi)\) ，且 \(3 \cos 2 \alpha-8 \cos \alpha=5\) ，则 \(\sin \alpha=\)\\
A，\(\frac{\sqrt{5}}{3}\)\\
B，\(\frac{2}{3}\)\\
C，\(\frac{1}{3}\)\\
D ，\(\frac{\sqrt{5}}{9}\)

显然题目中 \(\alpha\) 是唯一自由度，条件里恰好有一个方程，自由度减没了。因此，只需解这个方程，从而确定 \(\alpha\) 的大概位置，使得 \(\alpha\) 静止下来即可，这时候就算让你求 \(\tan \alpha\) 也是正常的。

10．已知 \(A, B, C\) 为球 \(O\) 的球面上的三个点，\(\odot O_{1}\) 为 \(\triangle A B C\) 的外接圆．若 \(\odot O_{1}\) 的面积为 \(4 \pi, A B=B C=A C=O O_{1}\) ，则球 \(O\) 的表面积为\\
A， \(64 \pi\)\\
B ， \(48 \pi\)\\
C， \(36 \pi\)\\
D， \(32 \pi\)

这道题比较适合在大脑中建模。先在脑海中绘制出一个球面，并在球面中截取一张平面，平面与球面的交线 \(\odot O_{1}\) 上有三个点 \(A, B, C\)（暂时不用管＂\(\odot O_{1}\) 的面积为 \(4 \pi\)＂这个等量关系）。再根据条件 \(A B=B C=A C=O O_{1}\) ，得知这三个点构成了一个等边三角形，而这个三角形显然跟形状相关；球心到 \(A B C\) 所在平面的距离 \(O O_{1}\) 的大小与三角形三边大小相等。如果刨去大小这一个自由度，这个题目就只剩一个自由度了，即 \(\triangle A B C\) 所在的这个平面是可以上下移动的，在移动过程中带动它的截线圆周以及上面的正三角形的运动。这时候可以意识到，题目主要考察的是三角形这个平面与球心的距离，如果不看 \(4 \pi\) 的话，整个的大小也是已经被确定了的。综上就是这样一个图景：球面被平面 \(A B C\) 所截，截了一个圆出来，在这个圆上有一个正三角形；这个三角形还一个旋转不变性（以过球心并垂直于平面 \(A B C\) 的直线为轴中心旋转）。所以平面 \(A B C\) 可以沿着轴中心上下平移，

这个正三角形也会跟着伸缩，\(\left|O O_{1}\right|\) 也会跟着变化，这就是在抛开大小的情况下自由度的关键所在。平移到什么程度为止呢？平移到三角形的边长等于 \(\left|O O_{1}\right|\) 为止（加入大小后平移结束）。\\
我们现在再补一个面积 \(4 \pi\) ，把大小这个自由度给刨掉，就静止下来了。最后题目让你求表面积，既然都静止下来了，那么就只需要像往常一样找关键方程即可。总之这种题目就在大脑里面好好建一下模，然后想象一下关键自由度，而这道题目的几何性和直观性算是比较强的。

11．已知 \(\odot M: x^{2}+y^{2}-2 x-2 y-2=0\) ，直线 \(l: 2 x+y+2=0, P\) 为 \(l\) 上的动点．过点 \(P\) 作 \(\odot M\) 的切线 \(P A, P B\) ，切点为 \(A, B\) ，当 \(|P M| \cdot|A B|\) 最小时，直线 \(A B\) 的方程为\\
A， \(2 x-y-1=0\)\\
B， \(2 x+y-1=0\)\\
C， \(2 x-y+1=0\)\\
D， \(2 x+y+1=0\)

已知 \(\odot M\) 和直线 \(l \Leftrightarrow\) 圆，直线都是静止的；\(P\) 为 \(l\) 上的动点 \(\Leftrightarrow\) 点 \(P\) 是运动的，一个自由度；接着是一些附属条件，最后给出问题。动点是一个自由度的话，理论上来讲它就可以表示为点 \(p\) 所在自由度的函数 \(f\) ，而 \(|P M| \cdot|A B|_{\min } \Leftrightarrow f_{\min }\) 。这类单一自由度的问题，整个图形里的各个要素，理论上都可以由这一个自由度表示，因为一旦这个自由度静止下来，整个图形就静止下来了，所以这个图形就可以把它看作是以这个自由度为自变量的函数，这个 \(|P M| \cdot|A B|\) 也不例外。这道题就是一个动态的问题，一个动态的过程，或者说单自由度的一个动态关系。关键是分析好最小值这个自由度，并进行一些转化，从而更快更稳的解决相关问题。

12．若 \(2^{a}+\log _{2} a=4^{b}+2 \log _{4} b\) ，则\\
A，\(a>2 b\)\\
B，\(a<2 b\)\\
C，\(a>2 b^{2}\)\\
D，\(a<b^{2}\)

两个未知数，一个变量，所以它只有一个自由度，但这道题并不太需要分析自由度，直接按照对应方法做题即可。

13．若 \(x, y\) 满足约束条件 \(\left\{\begin{array}{l}2 x+y-2 \leqslant 0, \\ x-y-1 \geqslant 0, \\ y+1 \geqslant 0 .\end{array}\right.\) 则 \(\mathrm{z}=x+7 y\) 的最大值为 \(\quad —\) ．没有参数，\(x, y\) 的运动范围是死的（约束条件），而 \(\mathrm{z}=x+7 y\) 是个变化的方程，所以 z 是一个单参数的一个自由度，所以这道题就是是问的取等条件时，这个方程是多少。因为求最值本质上就是最值的取等条件成立的时候，整个系统就被确定下来了，最值也就被确定下来了。

14．设 \(a, b\) 为单位向量，且 \(|a+b|=1\) ，则 \(|a-b|=\) \(\qquad\)

如果不看单位向量 \(a, b\) 的位置的话，那么题目本中只有一个自由度，即这两个向量彼此之间的位置关系，而 \(|a+b|=1 \Leftrightarrow\) 一个方程，直接静止化了。

15．已知 \(F\) 为双曲线 \(C: \frac{x^{2}}{a^{2}}-\frac{y^{2}}{b^{2}}=1(a>0, b>0)\) 的右焦点，\(A\) 为 \(C\) 的右顶点，\(B\) 为 \(C\) 上的点且 \(B F\) 垂直于 \(x\) 轴．若 \(A B\) 的斜率为 3 ，则 \(C\) 的离心率为 \(\qquad\)

双曲线明显有两个自由度，而在＂\(B\) 为 \(C\) 上的点＂之前的条件都是一些附属物品，或者说一些要素，之后便多了一个自由度。而平行关系和垂直关系都是可以视为一个等量关系的东西，本质都是减一个自由度，因为如果

用向量语言写的话都表示一个等式，一个是这个对应的比例，另一个是点乘，他们都可以代表方程，即等量关系。到这已经是三个自由度减一个自由度了，后面的斜率再减一个自由度，最后只剩一个自由度了。\\
这个自由度是啥呢？看问号：本题中问你离心率，而离心率就牵扯到大小，所以用这个邪法可以猜出只剩下大小这一个自由度了。换个角度也能理解，前面这些已知条件和大小都无关，其中双曲线和垂直这俩条件放大与缩小都无所谓，这个斜率肯定是不变的，所以最后剩下来的自由度就是大小。

16．如图，在三棱雉 \(P-A B C\) 的平面展开图中，\(A C=1, A B=A D=\sqrt{3}, A B \perp A C, A B \perp A D, \angle C A E=30^{\circ}\) ，则 \(\cos \angle F C B=\) \(\qquad\)

三棱锥本身是 6 个自由度，因为底面是 3 个自由度，点 \(p\) 是在三维空间上随便找的一个点，也是 3 个自由度，所以整体正好 6 个自由度。然后根据后面的条件在脑海中建一个模，因为 \(A B D C\) 这一圈都是固定的，已经死了，如果要还原三棱锥，即把如图所示的这个展开图立起来的话，只需要绕着 \(A B\) 旋转即可，想象两个纸壳子，铺在这个平面 \(A B D\) 上，然后沿着 \(A B\) 折叠，这折叠的过程中，本质上是二面角来操控的。\\
所以说前面的一大堆东西，本质上最后只剩下一个自由度了，三棱锥底座有了，面也有了，无非就是在这个面上转，都是垂直的，二面角就是操控这个运动的主要参数。而 \(C A E=30^{\circ}\) 不就叫二面角吗？把这个纸壳子往下压，压到 \(30^{\circ}\) 的位置为止，此时这个图形就死了，他问你啥都行，比如三棱锥的面积，体积，外接球等等。

17．（12 分）\\
设 \(\left\{a_{n}\right\}\) 是公比不为 1 的等比数列，\(a_{1}\) 为 \(a_{2}, a_{3}\) 的等差中项．\\
（1）求 \(\left\{a_{n}\right\}\) 的公比；\\
（2）若 \(a_{1}=1\) ，求数列 \(\left\{n a_{n}\right\}\) 的前 \(n\) 项和．

等比数列是两个自由度，首项与公比或者任意项与公比都可以，然后一个方程减去了一个自由度。结合第一问求公比，可知公比这个自由度在前面已经静止掉了，所以前面的等量关系应该是把公比确定出来了的。第二问求前 \(n\) 项和，\(a_{1}=1\) 再减一个自由度，直接静止掉了。最后利用相关知识＂错位相减＂解答即可。

18．（12 分）\\
如图，\(D\) 为圆锥的顶点，\(O\) 是圆锥底面的圆心，\(A E\) 为底面直径，\(A E=A D . \triangle A B C\) 是底面的内接正三角形，\(P\)为 \(D O\) 上一点，\(P O=\frac{\sqrt{6}}{6} D O\) ．\\
（1）证明：\(P A \perp\) 平面 \(P B C\) ；\\
（2）求二面角 \(B-P C-E\) 的余弦值．，\\
\includegraphics[max width=\textwidth, center]{2025_02_16_7c9aaf92bc99dc07f048g-010(1)}

第 18 题图\\
\includegraphics[max width=\textwidth, center]{2025_02_16_7c9aaf92bc99dc07f048g-010}

第 23 题图

圆锥是两个自由度，底面半径和高两个自由度，\(O\) 是底面的圆心，\(A E\) 是底面半径也没啥，\(A E=A D\) ，可以看到这是一个边长的比例，所以减少一个自由度，减哪个自由度呢？最后只剩下这个大小自由度了。本身圆锥两个自由度，这里有一个等量关系，说明这个圆锥的形状是确定的，所以只剩一个大小了，因为这个地方是边长相等本质上就相当于一个比例。 \(P O=\frac{\sqrt{6}}{6} D O\) ，这部相当于告诉你点 \(p\) 的相对位置吗，点 \(p\) 是一个自由度，告诉你一个比例，就是完全不考虑大小是多少，就看这些比例关系，然后第一问 \(P A\) 垂直与平面 \(P B C\) ，让你证明，第二问让你求二面角的余弦值。因为形状都固定了，那么让你干啥就千啥，只要他问的问题和大小，长度，面积无关的话，那么这些角度，实际上这些和垂直，二面角类似的都是与位置角度相关的东西，而和大小无关。准确的来说，前面的自由度都由齐次条件决定的，就是形状条件决定的，剩下的一个自由度可以通过大小来确定，就是考比值，形状。这道题的余弦值本质上也是比值。

19．（12 分）\\
甲，乙，丙三位同学进行羽毛球比赛，约定赛制如下：累计负两场者被淘汰；比赛前抽签决定首先比赛的两人，另一人轮空；每场比赛的胜者与轮空者进行下一场比赛，负者下一轮轮空，直至有一人被淘汰；当一人被淘汰后，剩余的两人继续比赛，直至其中一人被淘汰，另一人最终获胜，比赛结束。\\
经抽签，甲，乙首先比赛，丙轮空．设每场比赛双方获胜的概率都为 \(\frac{1}{2}\) 。\\
（1）求甲连胜四场的概率；\\
（2）求需要进行第五场比赛的概率；\\
（3）求丙最终获胜的概率．

这是一个概率问题，一般来讲不太适合纳入到这种运动和静止的体系中。不过概率问题中涉及的模型是一个确定的东西，理论上也是可以类比成静止的问题，。

20．（12 分）\\
已知 \(A, B\) 分别为椭圆 \(E: \frac{x^{2}}{a^{2}}+y^{2}=1(a>1)\) 的左，右顶点，\(G\) 为 \(E\) 的上顶点， \(\overrightarrow{A G} \cdot \overrightarrow{G B}=8 . P\) 为直线 \(x=6\) 上的动点，\(P A\) 与 \(E\) 的另一交点为 \(C, P B\) 与 \(E\) 的另一交点为 \(D\) ．\\
（1）求 \(E\) 的方程；\\
（2）证明：直线 \(C D\) 过定点．

本题中的椭圆只有一个自由度，\(G\) 为 \(E\) 的上顶点一个自由度，后面接着一个方程，此时这个 \(a\) 已经可以被我确定了，然后求 \(e\) 的方程后，第一问就解决了。然后点 \(p\) 是是一个自由度的动点，前面已经没有自由度了。但是它让你证明过定点，这在难题中经常出现的就是动态问题中的不变量，这种问题在运动和静止这种视角下，本质上就是一个动态过程之中的不变量，后续我们会在题目中讲解。

21．（12 分）已知函数 \(f(x)=\mathrm{e}^{x}+a x^{2}-x\) ．\\
（1）当 \(a=1\) 时，讨论 \(f(x)\) 的单调性；\\
（2）当 \(x \geqslant 0\) 时，\(f(x) \geqslant \frac{1}{2} x^{3}+1\) ，求 \(a\) 的取值范围．\\
\(f(x)\) 只有一个自由度 \(a\) ，而第一问 \(a=1\) 使得 \(f(x)\) 静止了；第二问，恒成立问题，它是一个不等式的关系，不是类似于等式的限制，它是不减少自由度的。但是它会把条件 \(a\) 给加强，就是说 \(a\) 的范围受到了限制，使得这

个不等式恒成立的这样一个范围。也就是说，和不等式相关的条件，再有一种就是类似于有两个零点，有几个极值点这种几个的，就不是严格意义上的等量关系，都是不减少自由度的。

22．（选修 4－4：坐标系与参数方程）（10 分）\\
在直角坐标系 \(x O y\) 中，曲线 \(C_{1}\) 的参数方程为 \(\left\{\begin{array}{l}x=\cos ^{k} t \\ y=\sin ^{k} t\end{array}\right.\)（ \(t\) 为参数）．以坐标原点为极点，\(x\) 轴正半轴为极轴建立极坐标系，曲线 \(C_{2}\) 的极坐标方程为 \(4 \rho \cos \theta-16 \rho \sin \theta+3=0\)\\
（1）当 \(k=1\) 时，\(C_{1}\) 是什么曲线？\\
（2）当 \(k=4\) 时，求 \(C_{1}\) 与 \(C_{2}\) 的公共点的直角坐标．

参数方程中这种的 \(t\) 是一个假参数，类比函数里的 \(x\) ，是参数方程里固有参数，并不算做一个自由度，但 \(k\) 是一个真参数，即一个自由度。然后极坐标没有参数，这是一个死的，第一问第二问给了条件，都成静态了，接下来该怎么做就怎么做就可以了。

23．（选修 4－5：不等式选讲）（10 分）\\
已知函数 \(f(x)=|3 x+1|-2|x-1|\) ．\\
（1）面出 \(y=f(x)\) 的图像；\\
（2）求不等式 \(f(x)>f(x+1)\) 的解集．

这两道题都是静态的，第二问也是一个不等式，用固定的套路做题即可。

读完理数题后，思考题目分布：\\
（1）静态（死题）或可以被方程消除参数的题占大多数。而我们可以观察到，死题的解法基本上都是比较死板的，一般通过我们的推理把积木（要素）搭起来即可。\\
（2）压轴型（难题）常见类型：动态问题分析，即动态过程中的分类讨论与分段讨论等的分析手段。研究动态问题的静态不变量问题，这类问题比较特殊，和处理静态问题时的用到的思想不太一样，即时是在运动过程中也有一些不变的东西。而类似于 2019 年新一卷的导数压轴题的题目，就是典型的复杂静态问题，考验的是学生的基本功，而并非花里胡哨的技巧。\\
（3）特殊题型（难以单纯归类）：应用题（常见于概率统计），以读清题干，建模为关键。

\section*{1．1．2．2 视角有用吗？}
没用！但是每道题都会用。具体作用如下：\\
（1）动一静视角帮助我们更好把握题目中的关系，帮助我们尽快找准简单题与中档题的着手方向。\\
（2）动一静视角提前一览题目全貌，增强我们做题信心。\\
使用方法：读完全题，快速建立基本关系，搞清：死题？几个参数（自由度）？

\section*{1.2 充要替换 \({ }^{1}\)}
本节将带领大家从逻辑的角度探讨什么是易错点，以及我们如何来规避易错点，还有另一个很重要的一点，就是按照我的习惯，如何去写高中数学大题的解题步骤。

\section*{讨论 \(A \subseteq B\) 时忽略空集}
\[
\left\{\begin{array}{l}
\text { 均值不等式 } \frac{a+b}{2} \geq \sqrt{a b}, \text { 要求 } a, b>0 \\
\boldsymbol{a} \cdot \boldsymbol{b}<0 \text { 并不能推出 } \boldsymbol{a} \text { 与 } \boldsymbol{b} \text { 的夹角为钝角 } \\
\ldots \ldots
\end{array}\right.
\]

高中数学中有大量易错点，而易错点到底是什么？如何规避？

在讨论这个之前，我们先复习如下这个知识点：

\section*{充分必要性}
有 2 个命题：\(p, q\) 。简单来讲，就是看 \(p\) 与 \(q\) 之间的逻辑关系：

\begin{itemize}
  \item 若 \(p \Rightarrow q\) ．则称 \(p\) 是 \(q\) 的充分条件（ \(p\) 可以充分说明 \(q\) ）
  \item 若 \(q \Rightarrow p\) ．则称 \(p\) 是 \(q\) 的必要条件（欲有 \(q\) ，则必须有 \(p\) ）
  \item 若 \(p \Leftrightarrow q\) ．则称 \(p\) 是 \(q\) 的充分必要求件，简称充要条件（二者等价）
  \item 若 \(p \nRightarrow q\) 且 \(q \nRightarrow p\) ．则称 \(p\) 是 \(q\) 的既不充分也不必要条件（逻辑无关联）
\end{itemize}

充分必要性是逻辑推理的根本。事实上，我们由上边的定义，可以很容易的有有以下两个结论：\\
（i）若在一个推理过程中某一步用到了 \(p, p\) 与 \(q\) 互为充要条件，则可以把推理中的 \(p\) 改成 \(q\) ，并不影响推理本身，反之亦然。\\
（ii）一个连续推理过程中，前一步总要是后一步的充分条件，否则推理不成立

上述两个结论看似简单，实则用处无穷。接着我们回答以下问题：

\section*{1．2．1 数学的易错点本质何在？}
分为概念型易错点和推理型易错点两大类：

\footnotetext{\({ }^{1}\) 保证逻辑严密性与下笔严谨性。常用的手段是遍历法，遍历整个自由度所在区域。
}\section*{1．2．1．1 概念性易错点}
例如：零点是数，如果让你求一个函数 \(f(x)\) 的零点，你就不能写它的零点对应的坐标，比如你算出来它的零点是一，你就不能说它的零点是 \((1,0)\) ，而是 \(x=1\) ；再比如截距是数之类的东西。这些东西实际上就是与概念相关的，在物理，化学，生物上就有很多概念型易错点，熟练掌握就可以完全规避掉。

但是，还有一种易错点是数学，物理等等演绎推理学科中独有的，那就是接下来要讲的推理型易错点。

\section*{1．2．1．2 推理型易错点（重点）}
凡是涉及演绎推理的可能都有，化学和生物里面也会有一点，但是没有数学中这么多。那什么叫推理型易错点呢？我们来看一个非常非常简单的例子。

例1．5集合 \(A=\{x \mid 2 a-1 \leq x \leq a+1\}\) ，集合 \(B=\cdots \cdots\) ，若 \(A \subseteq B\) ，求 \(a\) 取值范围．

易漏 \(A=\emptyset\) ，即 \(2 a-1>a+1\) 。实际上，\(A \subseteq B \Leftrightarrow " A=\emptyset "\) 或＂\(A \neq \emptyset\) 且 \(A \subseteq B\)＂，但是，常见错误是我们默认了 \(A \subseteq B \Leftrightarrow " A \neq \emptyset\) 且 \(A \subseteq B "\) ，并且在接下来的推理中，又忽略了 \(A \neq \emptyset\) 的条件。我们取 \(B=[2,4]\) ．\\
（i）错解：

\[
\begin{aligned}
& \because A \subseteq B \\
& \therefore\left\{\begin{array}{l}
2 a-1 \geq 2 \\
a+1 \leq 4
\end{array} \quad \Rightarrow \frac{3}{2} \leq a \leq 3\right.
\end{aligned}
\]

这里的 \(\because \Rightarrow \therefore\) 先漏掉了 \(\emptyset\) 的情形，又在推理中漏了 \(A=\emptyset\) 的情形．\\
（ii）正解：

\[
A \subseteq B \Leftrightarrow\left\{\begin{array}{l}
A \neq \emptyset \text { 且 } A \subseteq B \Leftrightarrow " 2 a-1 \leq a+1, ~ \text { 且 }\left\{\begin{array}{l}
2 a-1 \geq 2 \\
a+1 \leq 4
\end{array} \quad " \Leftrightarrow \quad " \frac{3}{2} \leq a \leq 2 "\right. \\
\text { 或 } \\
A=\emptyset \Leftrightarrow 2 a-1>a+1 \Leftrightarrow a>2
\end{array} \quad \Leftrightarrow a \geq \frac{3}{2}\right.
\]

严格来讲，讨论完 \(A=\emptyset\) 后，必须再把 \(A=\emptyset\) 部分加上，否则要扣分。因为不这样是不严谨的。

从上面的两个解来看，错解逻辑极为混乱，各步无法构成充分条件。再看正解，行文逻辑清晰，各步均保持严格的等价关系，由繁人简，保证不漏，不增解，成功规避掉空集易错点。

结论 推理型易错点是在解题的逻辑过程中，因人的思维局限性，容易导致前后构成不严格的＂\(\Rightarrow\)＂关系而出现的错误。进一步：其实绝大部分数学解题中的错误都是导致前后＂\(\Rightarrow\)＂关系破碎的问题，甚至包括计算错误： \(x+1<0 \nLeftarrow x<0\)（不充要！）

\section*{1．2．2 如何规避易错点？}
说完易错点本质后，我们来看看如何规避易错点。

\section*{1．2．2．1 非推理型易错点（读题不细，抄错数，基础概念不清）}
基础概念需自己的努力与认真的学习（本讲义第三部分将带你把高中数学里的基础概念全面梳理一遍）。审题，计算错误需要纠错与检查能力（见本章后几节的内容）。

\section*{1．2．2．2 推理型易错点规避方法——充要替换思想}
回顾先前集合小题的求解，我们保证了每一步都是充要的！之前提到，\(p\) 与 \(q\) 充要时，在推理中便占据同等地位，于是，在很多问题中，保持充分必要性是十分重要的！如

\[
x^{2}-1 \leq 0 \Leftrightarrow-1 \leq x \leq 1
\]

这当然是充要的。但是，后者明显比前者更简单，于是后续的过程，我们便可以抛弃原来的 \(x^{2}-1 \leq 0\) ，转而用 \(-1 \leq x \leq 1\) 代替。而上述内容的前提是真的是推理过程真的是充要的，像第一种的错解，就完全不充要，只是逻辑混乱地嗐写。

例1．6等比数列 \(\left\{a_{n}\right\}, a_{1}\) 是 \(a_{2}, ~ a_{3}\) 的等差中项，求 \(q\) 。

解

\[
\begin{aligned}
a_{1} \text { 是 } a_{2}, ~ a_{3} \text { 的等差中项 } & \Leftrightarrow 2 a_{1}=a_{2}+a_{3} \\
& \Leftrightarrow 2 a_{1}=a_{1} q+a_{2} q^{2} \\
& \left.\Leftrightarrow 2=q+q^{2} \quad \text { (注: 由于默认首项非 } 0, \text { 此处才充要 }\right) \\
& \Leftrightarrow q=1 \text { 或 } q=-2 \quad \text { (题干若有大前提 } q \neq 1, \text { 则左式 } \Leftrightarrow q=-2 \text { ) }
\end{aligned}
\]

例1．7 P－ABC 平面展开图，之前已知 \(\Leftrightarrow\)（1）．\\
\includegraphics[max width=\textwidth, center]{2025_02_16_7c9aaf92bc99dc07f048g-015(1)}\\
（1）\(A B P\) 绕 \(A B\) 转动，单参数自由度\\
\includegraphics[max width=\textwidth, center]{2025_02_16_7c9aaf92bc99dc07f048g-015}\\
\(\angle C A E=30^{\circ} \Leftrightarrow \angle C A P=30^{\circ}\)（卡死过程也是充要的！）\\
卡死后静止，随便做了（用 \(\overrightarrow{A C}, ~ \overrightarrow{A B}, ~ \overrightarrow{A P}\) 进行向量计算即可）。为什么能卡死？实际上脑内搭积木也是充要替换！

论充要替换：

\begin{itemize}
  \item 充要替换法本身是解题时帮助我们更有逻辑地利用条件，更清晰地简化条件，更准确地得到答案。
  \item 坚持随时检验，及时抛弃的原则。
  \item 实际应用：
\end{itemize}

例1．8（2020 III 卷理数）椭圆 \(C: \frac{x^{2}}{25}+\frac{y^{2}}{m^{2}}=1(0<m<5)\) ，离心率 \(\frac{\sqrt{15}}{4}, A, B\) 分别为 \(C\) 的左右顶点。\\
（1）求 \(C\) 方程。\\
（2）若 \(P\) 在 \(C\) 上，\(Q\) 在直线 \(x=6\) 上，\(|B P|=|B Q|, B P \perp B Q\) ，求 \(\triangle A P Q\) 面积。

1．读题：\\
（1）\(C\) ：单参数，已知：一方程 \(\Rightarrow\) 可静止\\
（2）\(P\) ：单自由度，\(Q\) ：单自由度，双方程 \(\Rightarrow\) 可静止\\
结论：本题是彻底的可静止化问题\\
2．下笔：\\
（1）经典：

\[
\begin{aligned}
a=\frac{c}{a}=\frac{\sqrt{15}}{4} & \Leftrightarrow \frac{a^{2}-b^{2}}{a^{2}}=\frac{15}{16} \\
& \Leftrightarrow \frac{25-m^{2}}{25}=\frac{15}{16}
\end{aligned}
\]

解得 \(m^{2}=\frac{25}{16}, \quad \therefore C: \frac{x^{2}}{25}+\frac{16 y^{2}}{25}=1\)\\
（2）先作图，看关系\\
\includegraphics[max width=\textwidth, center]{2025_02_16_7c9aaf92bc99dc07f048g-016(1)}\\
（1）\\
\includegraphics[max width=\textwidth, center]{2025_02_16_7c9aaf92bc99dc07f048g-016}\\
（2）\\
（1）虽然静止，不知何处，遍历法大致观察 \(|B Q|=|P Q|, B P \perp B Q\) ，更易控制相等？想象 \(P\) 走一圈，发现上下可对称！但 \(y_{p}>0, y_{p}<0\) 不太影响，于是由对称性，不妨设 \(y_{p}>0\) ，作出大致图。\\
（2）等腰直角？随便切。经典全等三角形，过 P 作 PM 于，如图，恒有 \(\triangle P M B \sim \Delta B H Q\)

\[
B P=B Q \Leftrightarrow \triangle P M B \cong \triangle B H Q \Leftrightarrow P M=B H
\]

『

\[
\left.\begin{array}{l}
P(-3,1) \text { 时, }|M B|=8 \\
Q(6,8) \\
P(3,1) \text { 时, } Q(6,2)
\end{array}\right\} x_{p= \pm 3 \longleftarrow y_{p}=1}
\]

思考两处 \(\Leftrightarrow\) 是为什么；解出 \(y_{p}=1\) 时，定量化基本完成，不难解出 \(x_{p}= \pm 3\) ，虽然是双解，但由于一直等价，所以双解都要！结束时别忘了 \(|M B|=|Q H|\) ．\\
注 推理型易错点仍需要积累，错过之后就把这一理性漏洞补上，而我们接下来要讲的就是规避推理性易错点的方法—遍历法。

\section*{1．2．3 遍历法初步}
所谓遍历，是指遇到求范围的题目时，在大脑中想象参数 \(a\) 从 \(-\infty\) 跑到 \(+\infty\) 导致的系统行为变化．

例1．9 \(f(x) \geq a\) 恒成立，求 \(a\) 范围．

该条件的转化相信不少同学十分熟悉，但其实根本无需背，考场时现场反映即可：\\
\includegraphics[max width=\textwidth, center]{2025_02_16_7c9aaf92bc99dc07f048g-017}\\
\(a\) 从 \(+\infty\) 开始跑，很大，必然 \(a\) 越大越难让 \(f(x) \geq a\) 恒成立，于是 \(a\) 一路下降。若 \(f(x)\) 如上图，\(a\) 一路下降，肯定到了 \(f(x)\) 最低点就行了，\(a\) 仅需往下跑，更对了！\(a \rightarrow-\infty\) 也没问题！\(\therefore a \leq f(x)_{\text {min }}\) 就是充要条件。对于刚才的 \(A=\{2 a-1 \leq x \leq a+1\}\) ，让 \(a\) 跑得很大就马上发现 \(A \rightarrow \emptyset\) 了，自然便会想起 \(A=\emptyset\) 这一易错点．（遍历一遍 \(R\) 后，充要条件便能正确写出了！）

\section*{1．2．4 了解化归思想}
所谓化归思想 \({ }^{2}\) ，简单说，就是把不熟悉的问题，转化为熟悉的问题。这是数学学习中最重要的思想之一。如果说，新高考模式下有哪些能力是最需要培养起来的，那么化归的能力一定是其中之一。为什么？在新高考背景下，套路的题目少，而新题目多（如 19 题等），这要求大家在面对新问题时，也能够找到合适的解决方案。怎么找？就是将＂新问题＂，一个你不熟悉的问题，转化为＂旧问题＂，一个你熟悉的问题。只要是你熟悉的问题，就可以把你熟悉的工具和手段，运用上去了。\\
注 详细讲解请见本讲义第三章第八节。

总结\\
－充要替换的理念是十分重要的思想。它保证可以不增解，不漏解的完成题目，而且有助于我们条理清晰实成题目，尽力规避易错点。

\begin{itemize}
  \item 所谓数学易错点，有不少是直接与推理的严密性挂钧的，需要在积累的同时，提高自己的充要辩别能力。
  \item 本节内容主要针对简单题与中档题，帮助大家提高该部分题目的准确率。
\end{itemize}

\section*{1.3 即时检查}
前两节内容主要围绕做题方法展开，包括动静视角（迅速帮我们宏观上把握整道题目的脉络）和充要替换 （从逻辑的角度把我们更加清晰，更加严密以及高准确率的完成一到题目的论证或计算）。本节将讲解使用这些方

\footnotetext{\({ }^{2}\) 比充要替换更本质的思维其实就是化归思想。
}法做完题后 up 主是如何提高自己的正确率的，主要从＂一遍过＂与＂检查＂ 2 个角度展开。

\section*{1．3．1 如何提高做对的几率？}
在我们做一道题目的时候，我们要尽可能的一次性做对。而提升做对题目的几率是非常重要的事情，因为你不能把得高分的希望寄托在检查上。主要思想：

\section*{1．3．1．1 动一静视角俯瞰全题，把握全题脉络}
稍微读题后就马上下笔的做法可取吗？答案是否定的，我们必须首先利用这种视角搞清楚这道题目是运动的还是静止的，它是一个有几个自由度，参数控制的问题；然后大致在脑子里建一个模型，想象一下搭一个积木，把握全体脉络。此时，我们应该列几个方程写几个函数之类的东西把握后就要开始做这个题了。

做题的时候要注意随时检查充要替换是否到位（见前面讲过的例 1.8 ）\(\Rightarrow\) 尽量减少读题失误与方法失误！！！这个东西如何避免？其实就是要提高我们读题的仔细程度。一方面你要把题目中所有的已知条件都给标记出来，因为高考题目中已知条件，就是在动静视角帮我们建构全体的脉络的过程中，一般没有一个条件是多余的。那么对于简单题和中档题，方法上的失误是不被允许的，比如像二项分布，超几何分布等这些常见的分布或者数列求和要用到的错位相减等等，都是比较固定的方法。

还有一类比较特殊的题目——应用题，比如纯粹的数学理论性的问题。它要加一些情景，所以此类题大家要多读几遍，保证题意理解准确 \(\Rightarrow\) 建好模后回归充要替换。

\section*{1．3．1．2 即时性检查一计算错误？（区别于下面的二轮检查）}
（1）快速重读全题：是否有条件没用上，或在充要替换里没有出现过？（结合自由度判断）\\
（2）检查得数：（1） \(17 \times 23=\) ？\\
（2）\(a_{n}=n+1, \quad a_{1}=2\) ？\(a_{2}=3\) ？\\
（2）\(\frac{1+i}{1-i}=i, \quad(1-i) i=i-i^{2}=i+1\) ．\\
让计算过程发生本质性变化！\\
（3）判断得数有无原则性错误？如：水池体积是 \(1 \mathrm{~cm}^{3}\) ．

\section*{1．3．2 一轮做题，二轮检查}
一道题最好不要马上二轮检查，因为可能被做题思路带偏！至少全部选择题结束再检查选择题（让题目先忘一忘），让做题过程与检查过程发生本质性变化！现在我们来看一看常见的检查方法。

\begin{itemize}
  \item 法一：见1．b－（2）
  \item 法二：见自由度替换：
\end{itemize}

例 \(1.10 \triangle A B C, A B=3, B C=5, \angle B=120^{\circ}\) ，则 \(A C=\)

初做：

\[
\begin{aligned}
A C & =A B^{2}+B C^{2}-2 A B \cdot B C \cdot \cos B \\
& =9+25-30 \times\left(-\frac{1}{2}\right)
\end{aligned}
\]

于是填了 49 ．\\
检查时： 49 怎么可能？ \(49>3+5\) ．改 49 为 7\\
检查：

\[
\begin{gathered}
\text { 去掉已知: } \underbrace{\angle B=120^{\circ}}_{\text {已知少一个自由度 }}, \quad \text { 增加已知: } \underbrace{A C=7}_{\text {已知加一个自由度 }} \\
\cos B=\frac{9+25-49}{2 \times 3 \times 5}=\frac{-15}{2 \times 15}=-\frac{1}{2} \Rightarrow B=120^{\circ} . \checkmark \text { 吻合! }
\end{gathered}
\]

注 自由度替换后若检查不吻合，则先看检查过程是否有错！！！\\
－法三：改变解法：

例1．11 \(\triangle A B C, A B=3, B C=5, \angle B=120^{\circ}\) ，则 \(A C=\) \(\qquad\) ．\\
\includegraphics[max width=\textwidth, center]{2025_02_16_7c9aaf92bc99dc07f048g-019}

做题：余弦定理\\
检查：过 \(A\) 作 \(A D \perp B C\) 于 \(D, A C=\frac{3}{2} \sqrt{3}, B D=\frac{3}{2}, \therefore C D=\frac{13}{2}\)

\[
A C=\sqrt{A D^{2}+C D^{2}}=\sqrt{\frac{27}{4}+\frac{169}{4}=7}
\]

（可以借助二级结论方便检查，如抛物线上点到焦点距离公式极坐标形式）\\
常见解法改变：几何法 \(\longleftrightarrow\) 计算，代数，公式……\\
注 一轮做题时选自己最擅长的方法！！！\\
－法四（选择题专用）：选项代入法／自信排除法\\
针对部分填人流程的选择题，只需验证所选选项是否符合题意即可，正常做出选择题结果后再选择 \(\underbrace{\text { v．s }}_{\text {也可用于检查 }}\) 排除法，代人法．

例 \(1.12 \cdots \cdots\) ，求 \(x\) 最大值（）\\
A， 1\\
B，2\\
C，3\\
D， 4

做题选了 D ：则只要验证 D 可以取到即可；做题选了 C ：（1）验证 C 可以取到 \(\Rightarrow A, B\) 不对．（2）验证 D 取不到 \((x=4\) 时无解 \() \Rightarrow D\) 不对。\\
－法五：退而求其次法：\\
在圆锥曲线题中，当你画出一个函数 \(S_{\triangle A P Q}=f(t)\) 时，可以对其验证两个特殊值。如：\(t=0\) 时，\(t \rightarrow \infty\) 时，看面积是否吻合。若吻合，那基本上 \(f(t)\) 就是对的。

\section*{1.3 .3 注意事项}
－二轮检查应量力而行。如果做题时间比较紧张，则把更多精力放在提高一轮做题精度上。熟练掌握动—静视角观察与充要替换能有效提高一轮做题效率，实时检查能有效提高一轮做题正确率。

\begin{itemize}
  \item 很多方法在即时性检查与二轮检查中是通用的，要随机应变，灵活处理。
  \item 检查的动力来源于对分数的渴望，直接受考场状态影响：自信心，激情，精气神都会影响做题与检查质量。
\end{itemize}

\section*{1.4 结构意识}
事前声明：\\
在第一，二节讲解中，我们为大家提供了一种新的看问题视角——动静视角，并为大家深度讲解了解题中逻辑推理的形式——充要替换（高级：充分替换／充要转化）。本期我们将系统带领大家感受这两大工具是如何在解题中大放光彩的，并进一步加深对这两大工具的认识。

\section*{1．4．1 要素关系网}
\section*{1．4．1．1 观察题目本身}
为了更深人理解＂做题＂这一行为，我们有必要对＂题目＂本身进行一个观察。

\begin{itemize}
  \item 一道题：已知 \(\cdots \cdots\) ，求 \(\cdots \cdots\)／证明 \(\cdots \cdots\) ．
  \item 回忆一下：我们的动静视角在干嘛？
\end{itemize}

例 \(1.13 \triangle A B C, B C=3, A B=4, \angle B=90^{\circ}\) ，求 \(A\) ．

已知条件实际是描述了两个东西：

\[
B C=3 \quad \underbrace{\text { 谁? } B C}_{\text {要素 }(\text { 对象 })} \underbrace{\text { 怎么样? 等于 } 3}_{\text {谓词 }(\text { 谓语 })}
\]

例1．14 若 \(f(x)=x^{2}-a x+1 \geq 0\) 恒成立，求 \(a\) 的范围。

已知描述了 2 个东西：

\[
\underbrace{\text { 谁? } a(\text { 参数 })}_{\text {要素 }(\text { 对象 })} \underbrace{\text { 怎么样? 使 } x^{2}-a x+1 \geq 0 \text { 恒成立 }}_{\text {谓词 }(\text { 谓语 })}
\]

例1．15 \(\triangle A B C\) ，\(D\) 为 \(B C\) 上一点。

\[
\underbrace{\text { 谁? } D}_{\text {要素 (对象) }} \underbrace{\text { 怎么样? 是 } B C \text { 上一点 }}_{\text {谓词 }(\text { 谓语 })}
\]

但这个例子与之前两个不同，似乎并没有讲什么实质性的东西。第一个例子的谓词描述了一个等量关系；第二个例子中的谓词描述了一个不等关系；第三个例子的谓词描述了一个对象关系。

一个已知条件描述了一个对象的某些特征，但是在此之有更加一般的事情：对象 解题是利用要素（对象）的性质（由谓词给出）进行推演的。但在此之前，对象本身是更加一般的情况，只有题目中设定了一系列对象，我们才能给对象附加性质（由谓词给出的已知条件），才能利用已知条件去推演。

\section*{1．4．1．2 观察对象之间的关系——要素关系网}
举一个例子：在 \(\triangle A B C\) 中，\(D\) 是 \(B C\) 中点，\(E\) 在 \(A D\) 上，\(A E=2 E D\) ，连接 \(B E\) 交 \(A C\) 于 \(F\) ．\\
\includegraphics[max width=\textwidth, center]{2025_02_16_7c9aaf92bc99dc07f048g-021(1)}

已知的条件：

\[
\begin{array}{r}
\left.\quad \begin{array}{r}
D \text { 中点 } \\
E=A E=2 D E
\end{array}\right\} \\
F=B E \cap A C
\end{array}
\]

题中全部对象：\(\triangle A B C, D, E, F\)

对于一个问题中的全部要素（对象），它们并不是＂平行＂——平起平坐的，而是有一个上下级的关系。比如 \(\mathrm{D}: ~ \mathrm{D}\) 是基于 \(\mathrm{B}, ~ \mathrm{C}\) ，寻找中点而的得，因此 D 应当是比 BC 靠后出现的， F 更是如此，最后才出现。因此，这个问题的对象之间有一个相互制约的关系网：\\
\includegraphics[max width=\textwidth, center]{2025_02_16_7c9aaf92bc99dc07f048g-021}\\
（要素关系网）

D 从 ABC 里生出； E 从 \(\mathrm{A}, ~ \mathrm{D}\) 里生出； F 从 E 中生出。\\
要素关系网反映了一个问题中对象之间的上下级关系。 \(\rightarrow\) 后的东西完全被 \(\rightarrow\) 前的东西所决定，\(\rightarrow\) 包含了要素之前的控制关系。（这种控制是与一只条件相结合的）\\
D 被 \(\triangle A B C\) 控制：等量关系：中点； E 被 D 控制：等分点； F 被 E 控制：相对位置（交点）。好比盖大楼（理解方式），先盖地基，再盖一层，再盖第二层……\\
注 要素关系网可能不唯一：不像现实生活中那样，必须＂父＋母 \(\rightarrow\) 子＂，数学中的要素关系未必如此。

例 1.16 已知线段 \(A B, C\) 是 \(A B\) 中点。

我们自然可以 \(A B \rightarrow C\) ，作为要素关系。但是，已知条件我们可以改一下：已知线段 \(A C\) ，沿长 \(A C\) 至 \(B\) ，使得 \(\overrightarrow{A C}=\overrightarrow{C B}\) ，于是要素关系网成了：\(A C \rightarrow B\) 。\\
所以，有时候对已知条件的不同理解可能改变要素关系。就好比 \(y\) 是 \(x\) 对的函数 \(y=x+1(x \rightarrow y)\) ，反过来，\(x\)是 \(y\) 的函数 \(x=y-1(y \rightarrow x)\) 。有时候，要素关系网的改写将会简化问题。\\
另外，要素关系网只是一种理清对象之间关系的手段，并不是僵化的，只要能搞清题中对象关系即可。

\section*{1．4．1．3 要素关系与动静关系——控制与牵制}
我们借助要素关系，来进一步阐明动静视角的核心。要素关系重点描述题目中所有对象的构建顺序（先有 \(\cdots \cdots\) ，后有 \(\cdots \cdots\) ，再有），而动静视角侧重于题中各要素的活动情况。真正在问题中，要素关系的活动于控制便是动静之间的反应。

例1．17椭圆 \(C: \frac{x^{2}}{4}+\frac{y^{2}}{3}=1, P, ~ Q\) 分别在 \(C\) 和直线 \(x=4\) 上，\(O Q / / A P, M\) 为 \(A P\) 中点，求 \(O M\) 于 \(O F\) 交点 \(R\)的轨迹。（ \(A\) 是左顶点，\(F\) 是右焦点）\\
\includegraphics[max width=\textwidth, center]{2025_02_16_7c9aaf92bc99dc07f048g-022}

我们首先看一下要素关系。最早的地基是椭圆，\(A \cdot F\) 这些静止物下一层似乎是 \(P, ~ Q\) 两个动点。但发现 \(P, ~ Q\)之间有一个牵制关系

\[
O Q / / A P
\]

如果设定了这个牵制后，\(P, ~ Q\) 也被放进这个体系内，\(M\) 直接从 生出最后一个位置关系

\[
R=M O \cap Q F
\]

大致看下来，可以画出一个要素关系网来：\\
\includegraphics[max width=\textwidth, center]{2025_02_16_7c9aaf92bc99dc07f048g-022(1)}

直接画要素关系网\\
\includegraphics[max width=\textwidth, center]{2025_02_16_7c9aaf92bc99dc07f048g-022(2)}

或者整理一下

但是，这和我们看到的要素关系图不同。\\
回顾（2）中的 \(\triangle A B C\) 例子，每一个箭头都是＂十分清晰＂的决定，感觉很舒服，但这个题似乎没这么友好：（1）里面有＂运气＂成分— \(P, ~ Q\) 无法定下来，而是以自由度 \(1=2-1\) 运动。于是，这个要素关系网需要添加上运动成分，也就是自由度。\\
\includegraphics[max width=\textwidth, center]{2025_02_16_7c9aaf92bc99dc07f048g-023}

为了更清晰，全面表达要素之间的关系，我们现在给要素关系网添加运动与静止的关系。（这部分大家自由一点，按自己喜好标注即可，无所谓）或者能记住的，不标注也行，关键是认识到这个图里的要素的动静关系是组成的一部分。因此，大家可以看到，我们原本运动与静止的观点，现在可以与要素关系网进行结合。其中 \(P\) ， \(Q\) 两点通过平行关系得到牵制，原本是两个独立的动点，现在之间有一个关系制约，就好比原来两个独立的参数 \(a, ~ b\) ，被一个已知的关系 \(a=b\) 牵制，本质上一样。当然，如果去掉已知的平行条件，那么 \(P, ~ Q\) 就是独立的动态对象了，彼此之间没有牵制关系（单纯是个动点）。

\section*{小总结}
要素关系网（看谁生成谁），盖大楼的理解方式，要素关系网中的静，动态对象。已知条件可以赋予动态对象之间的牵制；已知条件是赋给要素关系网的＂灵魂＂。

\section*{1．4．1．4 动静关系如何控制系统？}
当我们对一个题建立要素关系网后，我们不难看出：\(\rightarrow\) 后的已知条件不影响 \(\rightarrow\) 前的事物。如刚才的椭圆 （静）题：\(P \rightarrow M: M\) 是 \(A P\) 中点，与 \(Q\) 无关，也就是后面不影响前面（二楼怎么盖关一楼啥事）。我们在动静视角下看：\(p\) 与 \(a\) 一共是 1 个自由度（实际被 1 个参数控制），而这一个自由度同时控制着后方所有内容。因此，当我们用要素关系网来重新审视动静视角时：\(P\) 与 \(Q\) 可以运动，它们之间的相互牵制，导致运动是单参数控制的，这个运动同步随着要素关系控制后方内容。

如果抛开动静视角，那么这种控制仍然保留！\(P \rightarrow M\) ，即一个 \(P\) 确定一个 ；类似于函数 \(y=f(x)\) ，一个 \(x\)确定一个 \(y, \therefore M\) 是 \(P\) 的＂函数＂。现在 \(P\) 被一个单参数控制（ \(P\) 在椭圆上，\(P Q\) 手拉手平行运动），\(P\) 用数字表示，可以是 \(x\) 坐标，\(A P\) 斜率等等。是 \(P\) 的函数 \(\Rightarrow\) 的所有特征也都是这个参数，它就等同于 \(P\) 的函数，参数 \(\rightarrow P\) 的所有特征 \(\rightarrow\) 的所有特征。

因此，当我们放在要素关系网中时，后一级的要素作为前一级要素们的＂函数＂。前者的所有性质按照某些规则（数学规则）确定后一级要素的所有性质，而我们关心的解题，就是研究这些性质的。也必须按照要素关系网的顺序，由前件决定后件，前级要素性质决定后级要素性质。

当要素关系网中的对象被已知赋予了动态，静态关系后，静态的前一级要素将唯一决定后一级要素（好比已知 \(f(x)\) ，求 \(f(1)\) ）。而若是动态的（例如本题的 \(P\) ），我们需要找一些参数来刻画 \(P\) 的所有信息（本题中，\(P\) 是椭圆上的动点，可以用 \(x\) 坐标，\(A P\) 斜率等代表）。这些用来刻画 P 的信息（参数信息），让 \(P\) 的一切可计算。有了 \(P\) 的一切，就可以沿要素关系去推算更靠后的要素。\\
＊那么，我们究竟需要多少参数呢？取决于这个问题的动态关系有多复杂。平面上一动点 \(P\) ，需要两个参量 （横纵坐标，极坐标…‥）；曲线上一动点 \(P\) ，仅一个参量；这些参量的总数，称为自由度（与我们第一节讲的自由度相同）。

\[
\left\{\begin{array}{l}
\text { 所有这些自由度, 对应到这么多参量; } \\
\text { 分散到整个要素关系网中, 控制了整个系统。 }
\end{array}\right.
\]

注：牵制关系作为变量之间的等量关系，会减少系统的自由度，一个牵制等式会减少一个自由度。例如：AP／／OQ，\\
\(P\) 与 \(Q\) 之间存在一个等式的牵制（平行即向量坐标成比例，算作一个等式），于是 \(P, ~ Q\) 总共是一个自由度。

\section*{总结}
要素关系网中，后一级要素完全由它前面的要素唯一决定（类比函数）。当要素被一系列参数所决定时，后一级要素的所有性质可以被参数们决定。

对于一个具体题目，设定已知后要素们可能还有运动余地。我们从最前面的要素开始，分析其运动需几个参量控制，则顺着要素关系网可以定下全要素的全性质，这便是自由度控制要素关系。

\section*{1．4．2 简单题与中档题的解题技巧}
讲过要素关系网后，我们利用这个工具来阐述简单题与中档题的求解方法。\\
回顾：要素关系网中，靠前的要素完全决定靠后的要素（类比函数）；利用参数（找一些具体的量），如长度，角度，题设参数 \(a\) 等，可以控制靠前的要素进而控制靠后的要素（靠后的要素一律作为他们的＂函数＂），不同理解也许网会不同。因此，做题应当：先研究靠前的要素，再按 \(\rightarrow\) 去研究靠后的要素。但是，这样实操难度较大，正向思考不利于寻找方向。因此，可以采用以下的倒推方法：

\section*{1．4．2．1 静态问题（沿要素关系逆推）}
静态问题指没有任何运动成分的问题。这意味着从最初给定的要素出发，可以决定整个要素关系网内的事物，对于复杂静态问题，我们常常可以这样逆向考虑：

例1．18 三棱雉 \(A B C-A_{1} B_{1} C_{1}, \triangle A B C\) 是边长为 2 的等边三角形，\(O\) 是 \(A C\) 中点，\(O B_{1} \perp\) 面 \(A B C, O B_{1}=\sqrt{2}\) ， \(T, ~ M\) 分别是 \(B_{1} C, ~ A A_{1}\) 中点，求 \(\triangle O M T\) 面积。（填空）\\
\includegraphics[max width=\textwidth, center]{2025_02_16_7c9aaf92bc99dc07f048g-024(1)}

我们首先建立要素关系网，注意要得当建立，使问题简化：首先，底座 \(\triangle A B C\) 是完全确定（静止）的，先画它一定没错；然后，再 \(A C\) 上取一个中点（有了中点可建立 \(B_{1}\) ，进而确定上层 \(A_{1} B_{1} C_{1}\) ）。这样，整个问题被确定，剩下几个中点：\\
\includegraphics[max width=\textwidth, center]{2025_02_16_7c9aaf92bc99dc07f048g-024(2)}\\
\(A B C \rightarrow O \rightarrow B_{1} \rightarrow A_{1}, C_{1}\)\\
\includegraphics[max width=\textwidth, center]{2025_02_16_7c9aaf92bc99dc07f048g-024}\\
\(A B C \rightarrow O \rightarrow B_{1}\)

从这里看，\(A B C\) 静止，后面的都按题述规则唯一确定，\(\therefore\) 这个题是静态问题（我们没有分析自由度，便静止）。这实质上是从动静视角的本源去考察，直接利用靠前要素的动静关系，以及 \(\rightarrow\) 之间的控制，新要素是否添加新的动态关系来判断。入手要素关系网后，我们开始＂倒推＂：求 \(\triangle O M T\) 面积，我们需要知道什么？可以有很多选择，常见的：三边，\(S A S\)（两边夹一角），\(A S A\)（两角一边）\(\cdots \cdots\) 这里我选 \(|O M|, ~|O T|, ~ \sin <M O T\) 。注意一下： \(|M T|\) 不好，因为如果把 \(M T\) 放到要素关系网中，仍是最后一级要素，但是 \(O M=\frac{1}{2} A_{1} C\) 少一级，\(O T=\frac{1}{2} A B_{1}\)少两级，级越少越好算。夹角 \(\angle M O T\) 不好确定，但夹角 \(\leftrightarrow\) 向量数量积，所以可以考虑向量求解。倒不如说，这个题向量求解是比较好的，因为向量的加减可以往低级要素转化

\[
\begin{aligned}
& \overrightarrow{O M}=\frac{1}{2} \overrightarrow{C A_{1}}(\text { 少一级 })=\frac{1}{2}\left(\overrightarrow{C A}+\overrightarrow{A A_{1}}\right)=\frac{1}{2}\left(\overrightarrow{C A}+\overrightarrow{B B_{1}}\right) \\
& \overrightarrow{O T}=\frac{1}{2} \overrightarrow{A B_{1}}(\text { 少了 })=\frac{1}{2}\left(\overrightarrow{A B}+\overrightarrow{B B_{1}}\right)
\end{aligned}
\]

\(\overrightarrow{C A}, ~ \overrightarrow{A B}\) 都是靠后的了， \(\overrightarrow{B B_{1}}\) 仍靠后，但没有那么靠后。\\
注转化后，最靠后的要素是 \(\overrightarrow{B B_{1}}\) ，于是更后的要素可以扡弃。下面继续处理 \(\overrightarrow{B B_{1}}=\overrightarrow{B O}+\overrightarrow{O B_{1}}\)（易处理， \(\overrightarrow{O B_{1}}\)垂直底面，尽管靠后，但无妨）

\[
\begin{aligned}
& \overrightarrow{O M}=\frac{1}{2}\left(\overrightarrow{C A}+\overrightarrow{B O}+\overrightarrow{O B_{1}}\right) \quad \overrightarrow{O T}=\frac{1}{2}\left(\overrightarrow{A B}+\overrightarrow{B O}+\overrightarrow{O B_{1}}\right)=\frac{1}{2}\left(\overrightarrow{A O}+\overrightarrow{O B_{1}}\right) \\
& |\overrightarrow{O M}|=\frac{1}{2} \sqrt{|\overrightarrow{C A}+\overrightarrow{B O}|^{2}+\left|\overrightarrow{O B_{1}}\right|^{2}}=\frac{1}{2} \sqrt{C A^{2}+B O^{2}+O B_{1}^{2}}=\frac{1}{2} \sqrt{4+3+2}=\frac{3}{2} \\
& |\overrightarrow{O T}|=\frac{1}{2} \sqrt{|A O|^{2}+\left|O B_{1}\right|^{2}}=\frac{1}{2} \sqrt{1+2}=\frac{\sqrt{3}}{2}
\end{aligned}
\]

再求 \(\angle M O T\)

\[
\cos \angle M O T=\frac{\overrightarrow{O M} \cdot \overrightarrow{O T}}{|O M| \cdot|O T|}=\frac{\frac{1}{4}\left((\overrightarrow{C A}+\overrightarrow{B O}) \cdot \overrightarrow{A O}+\left|O B_{1}\right|^{2}\right)}{\frac{3}{2} \times \frac{\sqrt{3}}{2}}=\frac{-2+2}{3 \sqrt{3}}=0
\]

\(\therefore \angle M O T=\frac{\pi}{2} \quad \therefore S=\frac{1}{2} \times \frac{3}{2} \times \frac{\sqrt{3}}{2}=\frac{3 \sqrt{3}}{8}\)

结论 找易于入手的要素关系，基本是一路顺下去，牵制关系越少越好。静态问题学会迕推，方法是尝试把靠后的要素用靠前的要素表示，只要成功就有效！怎样表示？课上讲过的公式，定理等；自己对已知条件的把握（如本题选取向量法更易进行表示，中点变成中位线来简化）。

\section*{1．4．2．2 动态问题（静态问题推广）}
静态问题的要素关系网 \(\rightarrow\) 的控制是绝对的（ \(A B C\) 唯一确定 \(O, O\) 唯一确定 \(B\) ），好比一个函数 \(f(x)\) 给你，求 \(f(1)\) ，所以它最简单。但是并不是所有的题都如此＂静态＂，即我们必须面对控制本身。

例1．19已知 \(y\) 是 \(x\) 的函数，\(y=x+1\) ，若 \(y=5\) ，求 \(x\) 。

实际上这也是一个静态问题，\(x, y\) 双变量，\(y=x+1, y=5\) 是两个已知（等式），于是减两个自由度，全面静止。我们用今天讲的内容重新审视一下这个问题：\(x \rightarrow y\) 可以作为一个控制（要素关系网），\(y=x+1\) 被添加进了 \(\rightarrow\)内，表示 \(x\) 唯一控制 \(y\) 。因此但看这个 \(x \rightarrow y\) ，是一个自由度（只加了 \(y=x+1\) 已知），但这里是 \(y\)（后面的）\(\rightarrow\) 求 \(x\)（前面的）。我们还可以问：\(y>0\) 时 \(x\) 范围（直接研究控制的函数 \(y=f(x)\) 本体），这类问题虽然动静关系不同但是都有相似性——我们必须考虑＂\(\rightarrow\)＂带来的控制关系本身。

例1．20 \(\triangle A B C\) 是等边三角形，\(D\) 为 \(\triangle A B C\) 外一点，\(B D=2, C D=1\) ，四边形 \(A B D C\) 的面积是 \(\frac{5}{4} \sqrt{3}+2\) ，求 \(\angle B D C\) 大小。\\
\includegraphics[max width=\textwidth, center]{2025_02_16_7c9aaf92bc99dc07f048g-026(1)}

我们当然可以去这样画要素关系：\(A B C \rightarrow D\) 。但先画 \(D\) ，找个 \(B D, ~ C D\) 都好画，进而画出 \(\triangle A B C\) ，于是

\[
D \xrightarrow{(1)}\left\{\begin{array}{l}
B \\
C
\end{array}\right\} \xrightarrow{(2)} \Delta A B C
\]

（1）：\(B D=2, C D=1, \triangle B D C\) 中，共 3 个自由度，去掉 2 个，还剩 1 个，正好取成 \(<0\) ．\\
（2）：等边（明确的控制）\\
第一个箭头箭头中＂\(D \rightarrow\left\{\begin{array}{l}B \\ C\end{array}\right.\) 有运动成分在里面，自由度是 1。让 \(C D\) 来控制这个自由度，那么剩下的都是顺水推舟。于是 \(\triangle A B D C\) 面积可以被 \(C D\) 表示。从 \(D\) 出发，先到 \(\triangle B D C: S_{\triangle B D C}=\frac{1}{2} \times 2 \times 1 \times \sin D\) ，再 \(\triangle A B C\) ：需知一边长度

\[
\begin{aligned}
& \because B C^{2}=4+1-2 \times 2 \cos D=5-4 \cos D \\
& \therefore S_{\triangle A D C}=\frac{\sqrt{3}}{4} B C^{2}=\frac{\sqrt{3}}{4}(5-a \cos D) \\
& \therefore S_{A B D C}=\frac{5 \sqrt{3}}{4}-\sqrt{3} \cos D+\sin D=\frac{5}{4} \sqrt{3}+2 \\
& \therefore 2=\sin D-\sqrt{3} \cos D \Rightarrow D=\frac{5}{6} \pi
\end{aligned}
\]

\(f(x)=y, y\) 已知 \(\Rightarrow\) 求 \(x\) 。这属于传统静态问题的变体，但仍算是静态问题。\\
注 另一个鲜明例子是之前那道椭圆题目，up 主（nichijou 的涂山苏苏）在 2011－10－12 出视频讲过，其中把问题归结到斜率上去，然后确定轨迹。实际上，这也是要素关系：\\
\includegraphics[max width=\textwidth, center]{2025_02_16_7c9aaf92bc99dc07f048g-026}

P （动点）是唯一自由度

把 \(P, ~ Q\) 的平行关系及牵制条件拆出来，以 \(A P\) 斜率作为表示问题的变量。大家可以结合 10 月 12 日视频进行思考（留作练习）。这道题是纯动态问题，要素关系网中的 \(P\) 是纯动态的，可以把 \(R\) 视为 \(P\) 的＂函数＂，一个 \(P\)决定一个 \(R\) 。本题研究 \(R\) 的轨迹，就是研究这个对应＂函数＂值域，

结论 这类含动态问题的题目都与要素关系网中 \(\rightarrow\) 的控制有关。有的是求 \(f(x)=y\) 之解（化回静态，实则动态）；有的研究 \(f(x)\) 的本身特征（真动态），但思想上具有一致性。

\section*{总结}
动态问题找合适的要素关系网来辅助分析，\(\rightarrow\) 尽量能确定性控制后半段（但有时 \(\rightarrow\) 后的东西有自由度也难免）。对于自由度为 0 （知 \(y\) 求 \(x\) ）的这种，建议仍然放在动态要素关系网中，选取合适的，较早的要素去表示所求，然后解方程。纯动态问题也要选好的，考前的要素来分析。\\
总之：找一个好的要素关系网（牵制少，易于搭建，易于想象）胜过一切！

\section*{1．4．3 漫游难题}
刚才我们谈到了要素关系网以及它在解动，静问题中的应用，核心是找一个好的要素关系网来达到前项控制后项的目的。所有的自由度将散播在一个要素关系网中，作为自变量控制后面的要素；大部分简单，中档题，这种思想可以让我们尽快配合动静视角找出解法。那么，对于难题，该如何处理呢？

\[
\text { 难题大概分为 }\left\{\begin{array}{l}
\text { 过程复杂 (要素关系网长); } \\
\text { 难以转化 (牵制太多); } \\
\text { 特殊性 (存在一些关键的观察, 等价转化)。 }
\end{array}\right.
\]

过程复杂的问题，要把要素关系网的构建，以及逆推来找控制关系作为重点，借助要素关系网可以减少凌乱。\\
难以转化的问题，可以尝试去掉一部分牵制（如刚才那个椭圆题，\(P, ~ Q\) 之间的牵制关系再拆一层 \(\rightarrow\) 出来），并且，此时请务必重视充要替换！一些难以转化的牵制条件，必须靠充要替换进行等价变形，具体如何等价变形，需要一定的预判能力，做题经验等，这也是希望突破难题的同学要重点关注的。一些较为特殊的问题，需要某些二级结论积累或某些关键步骤之类的，也要靠一些适当积累，但根本来讲，预判与充要替换的扎实掌握依然十分重要。简化问题，不断等价也是要素关系网难过处理复杂问题的根本。

\section*{1．4．3．1 思想试验}
思想试验：没有感觉时先做做简单情形或者静态情况的尝试，寻找突破

\section*{1．4．3．2 预判}
预判：结合要素网与元问题（后续章节会讲）事先对计算的式子结构进行预判，运筹帷幄才能决胜考场在第五章会有很多对这两个解题思维的练习

\section*{总结}
\begin{itemize}
  \item 要素关系网：反映一个题目中所有对象的层级关系（谁可以生出谁）
  \item 要素关系网 \(\rightarrow\) 可能存在动态关系 \((P\) 在 \(A B\) 上，即 \(A B \rightarrow P(\) 动 \())\)
  \item 一个题目中所有的自由度会分散到要素关系网中，在各个节点控制系统的行为
  \item 静态问题需要要素关系网中一切可以被完全确定的问题，即自由度为 0
  \item 解题时，靠前的要素来还原靠后的要素一定可以简化问题（倒推法）
  \item 切记！理解为主不要去背，原理都很好理解，从思想上认识它！
\end{itemize}

\section*{1.5 考场习惯}
本次内容是在前四节的基础上展开的，如何在考场实战中最大限度发挥其价值，让每一个认真学习阳哥提供的思维的同学，都能够在应对任何类型的高考题中拿到复合自己能力的分数。本节内容将转化型针对高中阶段的数学考试，为大家讲解考场上的思维使用，答题速度，正确率等内容。

本节内容需要读者在本章一到四节的基础上进行食用，如果哪一节还没看过的话，请先阅读后再阅读本节内容。

本章前四节，我们讲解了一系列针对高中数学题目的普适性思维范式，包括动静视角，充要替换，要素关系网等，并介绍了最大限度确保正确率的使用方法——即时检查。前述三项内容可以在任何一份高考试卷（我们主要关注全国卷和新高考全国卷）中完美覆盖到 130 分左右的容量，如果基础知识十分扎实，并熟练掌握这些思维方式，理论上将可以稳定拿到 130 分的保底分。

但是，如何才能在考场上完美发挥出这样的规划呢？真实的高考考场会受到心态，难度波动等各种因素的干扰，而且如果基础知识不够扎实，或者这些思维还不够熟练，都可能干预考场拿分的稳定性，增加不确定性，从而导致发挥失误等。

本节将会把目光聚焦考场，为观众朋友们详细解析如何在考场中，以最大的稳定性来获取高分。将会告诉大家怎样更丝滑顺畅地使用我们讲解的思维，怎样在考场上尽量提高答题速度的同时兼顾正确率。在了解了这些方法之后，还需要长时间的试错，实践和反思进行巩固提升，这些功夫绝非一朝一夕，但是在稳定性的加持下，一定会让大家慢慢进步。

\section*{1．5．1 元问题}
\section*{1．5．1．1 解题，究竟是为了什么？}
我们讲到的动静视角与要素关系网，实际上都在解剖一道题目的结构，其中，要素关系网又可以看作是动静视角的解释。绝大多数的题都是：在一些基础的要索之上，通过个断地控制关系（也就是 \(\rightarrow\) ）确定更高级的要素。这就是要素关系网的含义。

处于基础地位的要素就作为这个要素网的主自由度，它们是搭建整个题干模型的＂地基＂，是整个题目的要素结构的操纵杆，手柄。下面，我们来看两个例子，这个例子里面暂时没有涉及更多的已知条件，单纯是结构性的题干。

例 1.21\\
（1）已知 \(\triangle A B C, D\) 是 \(B C\) 的中点，\(E\) 是 \(A D\) 的中点，\(B E\) 交 \(A C\) 于 \(F\)\\
（2）已知数列 \(\left\{a_{n}\right\}\) ，首项 \(a_{1}=1, a_{n+1}=2 a_{n}\) ．

这两个例子都是题干形式的内容，可以看到，（1）是标准的要素关系网的结构，（2）则是数列的一项一项连续控制的结构，它们都可以看作是要素关系网的形式。

当要素网搭建好了之后，题干中余下的已知条件就可以向内填充了，也需要去关注题目要求什么了。这个

时候，新的已知条件就作为等量关系，去削减题目本身的自由度个数，一个等量关系减少一个自由度，最终余下的数字就是实际的自由度。自由度在第一节和第四节都在反复讲，是一道题目的系统实际蕴含的，独立的，实变量的个数。而这里的等量关系，有的是直接告诉你某些原始要素或高级要素的情况；另外一些则是在要素网中随机拉线，把两个看似无关的要素强加一个等量关系，这种我们称为牵制关系。另外，还有的已知条件并不能充要替换为等式，而是一些不等关系，它们只不过是限制了运动的范围（例如 \(0<x<1\) ），不改变自由度本身；除此之外的都是极为个别的，需要就事论事的条件（这些就十分罕见了，需把重心放到充要替换上，容易出难题，部分可能需要灵感，在我们的 130 分规划之外）。

最后，实际自由度为 0 ，则问题变为静止，所有的主自由度都是理论可求的，控制关系下的其它要素也可以顺着求出来；实际自由度为 1 ，则问题实质是一个单变量问题，最后往往要求某个量的取值范围（或者还有大小无关，位置无关等特殊情况，参见以往相应章节内容）。

以上就是我们动静视角和要素关系网的大致内容。但是，实际应用是需要找们真正落实到找思路，求解的，空有要素网和自由度，却下不去手求解问题同样不行。那么，在实际的考试中，我们真正落实到答题卡上的，是什么呢？

显然，回顾刚才的整个要素关系网，我们还会发现，唯一剩下的，有实际解题作用的，就是＂\(\rightarrow\)＂！！\\
\(\rightarrow\) 是什么？是要素关系网中的控制关系，小小的箭头包含了前向如何控制后项的全部过程。例如，取一个中点，可以是 \(\rightarrow\) ；函数的一个极值点，可以是 \(\rightarrow\) ；圆锥曲线中一些奇怪的线交出一个交点，可以是 \(\rightarrow \ldots \ldots\)

因此，\(\rightarrow\) 所代表的控制关系是更加＂元＂的关系，描述的是一些数学本身的关系，是超越所有题目的，具有极度普适性的关系。

例 1.22\\
（1）已知 \(\triangle A B C, A B=1, A C=2, \angle A=\frac{\pi}{3}\) ，求 \(B C\)\\
（2）已知 \(\triangle A B C, A B=2, A C=3, B C=4\) ，求 \(\cos A\)\\
（3）已知等比数列 \(\left\{a_{n}\right\}\) ，首项 \(a_{1}=1\) ，公比 \(q=2\) ，求 \(a_{2}\)\\
（4）已知等比数列 \(\left\{a_{n}\right\}\) ，首项 \(a_{1}=1, a_{2}=2\) ，求公比 \(q\)\\
（1）和（2），（3）和（4），它们的要素关系网可以写成一样的形式，因此 \(\rightarrow\) 也是一样的，因此，它们的＂元＂是一样的，都在做同样的操作，只不过到最后一个是解方程，一个是算数（注意余弦定理两种形式的一致性）。可能有的同学会把这个叫做方程思想，或者函数思想，这当然没问题，不过，我们继续往下走一步。

这只是个简单例子，在偏复杂的中档题和难题中，\(\rightarrow\) 代表的信息量往往更多，但是，它们的一般性都不会改变，我习惯把这种问题称为＂元问题＂。元问题的内容就是找出究竟这些要素是如何控制的，比如余弦定理的式子本身，就包含了 \(\rightarrow\) 中的全部信息；一个小小的 \(a_{2}=a_{1} q\) ，也包含了这样的全部信息。

元问题是描述控制关系的方式，而牵制条件的发散程度则更高，看似风马牛不相及的东西都能被一个牵制强行扯上关系。因此，牵制关系的分析往往比控制更难，但是这二者在解题时地位相仿，都是一个小问题，需要我们单独处理掉以便于继续解题过程，因此，他们都可以认为是解题中的＂单位＂。

只有当我们列出的式子能够包含整个控制过程的全部信息，这个式子才能正确起到充要替换的作用，才能正确的表达一个个单位，才能够让我们得以进一步往下走，完成整个题目。因此，大家可以看到的是，元问题与牵制的分析和充要替换，是解题过程中不可或缺的一步，其重要性和要素网建立，算数求解，是相当的。我们甚至可以说，绝大多数问题的处理，其实就是动静视角总览 \(\rightarrow\) 建立要素网 \(\rightarrow\) 处理元问题与牵制条件 \(\rightarrow\) 算术求解的过程，前二者在本章第 1,4 节里已经分析过，最后的求解则是算数，解方程的基本功（以及适当的即时检查）。

\section*{1．5．1．2 如何学习知识点？}
那么，剩下的中间这个，处理元问题与牵制，它的实质是什么呢？从刚才的分析中，大家想必也看到了，元问题和牵制的处理，很大程度上就是基础知识的使用！！！因此，大家也可以看到了，我们讲过的思维完美覆盖的范围是简单题，中档题和部分难题所占据的 130 分，这些题的求解在本质上是一致的都是这样的，几步走的过程，动静视角与要素网的思维范式起到高屋建瓴，指引方向的作用。剩下的基础知识应用和计算求解全部是依托于这些思维，受灵感，偶然想到之流影响很小，几乎都是在要素网的基础上，自然而然孕育出的。

而动静视角与要素网的难题，往往会在牵制上下功夫，会让一些牵制看上去不好处理，这也是区分难题的一种方法。按照高考正常的 4：4：2的比例，难题分值在30分左右，其中会有一些靠要素网复杂度堆积的题，会有一些牵制关系处理有技巧的题，还会有关键步骤不好想的题，这三者难度依次递增，我们讲到的思维对它们处理的效果也依次下降，最后大概能完美适配的占 10 分左右，这也是我告诉大家的＂130分完美适配＂的依据。

因此，一个显然的事实就是，基础知识的熟练掌握，以及计算，解方程的基本功，是解题，乃至考高分（130分及以上）的基本要求，也是刚需。这些的功夫在哪？在平时在学校的课堂学习，课后的知识巩固以及练习中，这些东西的熟练掌握是作为学生的本分，如果基础知识都不熟练，何谈考出好成绩呢？

那问题就来了，怎么就叫基础知识熟练掌握呢？是公式定理背的滚瓜烂熟吗？基本公式定理的熟记自然是刚需，但是回头看我们的那几步走，显然，基础知识的存在是为了处理元问题，因此，我们在平时学习，巩固基础知识的时候，更需要关注的是，这个基础知识能够创造出怎样的元问题？刷题千变万化，但是元问题往往花样较少，只要自己平时练习的时候多多留心，这方面肯定是能够自己解决的，就比如：等比数列的元问题就是首项＋公比可以确定整个数列的信息；余弦定理的元问题是 SAS 或 SSS 型三角形的其他要素求解；离心率的元问题是确定圆锥曲线的形状，所有大小无关参量与离心率直接挂钩．．．．．．。

类似的这些东西大家在学习或复习的时候理应去主动关注一下，不用去死记硬背，本身这些东西就与要素网搭建息息相关，大家在学知识的时候要带着要素网这样的思维，好比思维被动静视角，要素网重塑了一般，变成习惯，这样的话，知识点自然的就以这样的形式储存在你的脑子里了，使用这些思维的时候自然也就信手拈来了。这是一个大工程，需要覆盖到高中数学全部的知识，需要积跬步以至千里，不是任何投机取巧能做到的，但却是付出时间必定有收获，必定可以完成的。

\section*{1．5．2 考场如何提高解题速度}
在上一节的基础上，熟练掌握知识点之后，我们终于来到了考场，来到了做题实战中。在真正做题的时候，我们需要兼顾做题速度与正确率，这样才能保证我们的得分。这一节和下一节将分别阐述这两方面。首先是做题速度，如何提高做题速度？实际上，从大的方向来讲，一个是我们的思维敏捷度，通过科学的思维方式与自身的努力，可以让这一方面受到智商的影响大大减弱；另一个是答题习惯，这对于答题速度的影响可能远远超出你之前的想象，是需要高度重视的一点。

\section*{1．5．2．1 思维焦点与隔离意识}
我们前面谈到了解题的行为过程，这些行为主要适配于前述的 130 分容量，当然，这 130 分左右的分值是我们提高答题速度的主体对象，在答题速度这一方面，难题的影响相对较小（ 2.2 将进一步展开叙述），因此，在思维层面的提速核心就是适配我们前述的那类思维范式的题目，如何将这些题目的思维加快，用更稳定且高效的方式解决呢？

我们可以这样想：在整个做题过程中，哪些地方最有榨取时间的潜力？很明显，元问题或牵制条件的分析，以及后续计算过程是十分层次分明（因为有要素网），也是最容易提速的地方。在有了一个比较好的要素关系网之后，网中的每一个单位都需要我们进行转化，有的元问题，比如取中点，做平方之类的，或者有的牵制，比如某两条线段相等之类的，太过简单，本身也几乎不消耗我们的时间，但是每道稍微复杂一点的题总有几个比较关键的元问题或者牵制，对解题的影响很大，而且也需要花点时间去分析清楚的。

这个时候，如果我们直面庞大的要素关系网，很容易陷入凌乱，但是当我们认清楚当前要分析好哪个单位，或者，要算好哪一步的时候，我们如果能把思维重心暂时转移到这个点上，先把这个点分析好，再转而在要素网上检索下一步的内容，再重复刚才所说的过程，这样的话，我们答题的条理性和目标性都会大大提高，也会让我们尽量减少发懵而内耗掉的时间，从而提高解题效率。

也就是说，我们应当形成一种思维焦点的意识，时刻都要认清楚我们现在要做的＂小事情＂是什么，并且在认清这件事之后，立刻把思维焦点放在上面，专心搞定当前的主要矛盾，然后再去寻找下一个主要矛盾，最终完成整道题目。

例 1.23 （2022 新高考二卷）\\
已知 \(\left\{a_{n}\right\}\) 为等差数列，\(\left\{b_{n}\right\}\) 为公比为 2 的等比数列，若 \(a_{2}-b_{2}=a_{3}-b_{3}=b_{4}-a_{4}\) ，\\
求证：\(a_{1}=b_{1}\)

看完问题，要素网一目了然，\(a_{n}\) 被首项和公差控制，\(b_{n}\) 被首项控制， 3 个自由度，外带给定的牵制条件，实际自由度为 1 。这里的控制关系就是基础的等差等比公式，很简单也好处理，于是思维焦点立刻放到牵制上去，发现都是两个数列里的要素构成的等式，于是思维焦点变成：先立刻用主自由度表示这些要素，从而进行已知牵制的充要替换。

表示结束，得到

\[
a_{1}+d-2 b_{1}=a_{1}+2 d-4 b_{1}=8 b_{1}-a_{1}-3 d
\]

顺手进一步简化得到

\[
d=2 b_{1}, 2 a_{1}=12 b_{1}-5 d
\]

充要替换后立刻把思维焦点转到所证式上，立刻明白只需从刚才的式子推导 \(a_{1}=b_{1}\) ，思维焦点切换为纯代数问题，余下不表。

这个简单的小例子告诉我们做题时思维焦点切换的一般情况。答题时，迅速切割出当前要干的事情的重点 （主要矛盾），把这个问题单独隔离出来，把思维焦点放上去，利用熟悉的基础知识迅速收割掉这个东西。反之，如果做题没有这种隔离意识，以及迅速切换思维焦点的习惯，那么做题就可能面对问题出神，发懵，拖拉，导致时间慢慢浪费。

以刚才这个小问为例，读题加构建要素网大概不到 1 分钟，写出两个式子再 1 分钟，得证大概半分钟，合

计最多应在 3 分钟内完成。反之，若是读题拖拖拉拉，拖到 1 分半，写两个式子犹犹豫豫，花费 2 分钟，然后盯着式子发憘不知道接下来要干啥，再多拖 1 分钟，加起来，单这一小问就多拖了 2 分钟，若是整张卷子 20 多道题，每题都多拖 1 分钟，那可能就是＂剩 10 分钟答完＂，和＂导数来不及看＂的差距了。

因此，及时将当前主要矛盾隔离出来，把思维焦点锁定到一个个小问题上，然后气势十足地秒掉一个个小问题（甚至不少简单题和中档题的元问题和牵制都可以当默写题来做），能够有效帮助我们提高答题时的速度和专注度，提高考试时的答题效率。当然，做到这一点，前提还是知识点足够熟练，不然的话连＂默写题＂都卡壳，就不要指望能有多高的分数了。

\section*{1．5．2．2 主动找活，积极设问}
设问：遇到困境主动提问自己当前矛盾点何在，我又能干什么\\
做题时，如果全是默写题级别的元问题和牵制，那自然是好不过的，但是有的时候，我们也会遇到需要稍微动脑筋的牵制条件，而在考场的高压状况下，有可能我们会大脑短路，或者找不到处理方法（也就是通常说的 ＂找不到思路＂）。下了考场一看，其实很简单，然后懊悔不已。这种情况常见于某些牵制条件的处理，在控制条件中很罕见，而且，即使是在 130 分的规划内，有的牵制条件也可能需要大家动一动脑筋（当然，余下 20 分的难题里也可能有这种牵制出现）。那么，我们该怎么样减少这样的突如其来的灵感的左右，尽可能以稳定性较高的方式拿到较多的分数呢？

其实，解决方法十分朴实无华。我们在面对一个条件的时候，通常都要对它进行观察，里面涉及到的要素如果可以直接用主自由度处理掉，那再好不过（比如刚才的题目），但是有的牵制里面可能会隐藏有一些需要变形或者代换掉的结构，又或者，可以通过配凑进行处理的结构，等等。在简单题和中档题中，这样的结构往往藏得不深，相对比较容易识别。因此，在遇到一时卡壳的条件时，我们应当优先干的，是先对条件中的结构进行观察，去主动询问自己：对这个条件，我能干什么？

这样的设问是十分有主动性的，我们会在这样的设问下主动寻找条件中我们熟悉的结构，或者可以进行操作的结构，有时还会瞬间找到多个处理的方向。这时候，我们的心态应当是积极的，因为＂有事可做＂肯定是不小于＂无事可干，一筹莫展＂的，然后我们就要主动尝试走一走这几条路径。位于简单题和中档题位置的题目，肯定在处理上不会显得非常有技巧性或者令人叹为观止，尝试几种路子往往就能把问题简化。

例1．24（2022 全国乙）\\
已知 \(\sin C \sin (A-B)=\sin B \sin (C-A)\) ，证明： \(2 a^{2}=b^{2}+c^{2}\)

自由度和要素网都没啥好分析的，单纯的牵制条件转化问题。观察这个式子，乍一看肯定会觉得有点茫然，不知道该干啥。自问：我现在能干什么？式子中什么地方可以进行操作？这一观察，很明显能找到至少 3 条路径：把正弦用差角公式打开；利用 \(\sin C=\sin (A+B)\) ；正弦定理角化边．．．．．．在没有方向性的时候，我们可以先做一点简单的尝试，一共就这么几条路，挨个稍微试试看看？

第一条：

\[
\sin C \sin A \cos B-\sin C \sin B \cos A=\sin B \sin C \cos A-\sin B \sin A \cos C,
\]

再观察就是

\[
\sin C \sin A \cos B+\sin B \sin A \cos C=2 \sin B \sin C \cos A,
\]

好像又有点卡，所证是边相关，再自问能干什么？刚才第三条路不是说了？可以正弦化边？那就继续走走：

\[
a c \cos B+a b \cos C=2 b c \cos A
\]

看上去标致了不少，感觉像是简化了，本着＂简化即前进＂的原则，再自问我能干什么？右边挺简单的先一放，左边提个因式，剩下 \(a(c \cos B+b \cos C)\) ，还能干什么？主动寻找，就能寻找到射影定理，就成了

\[
a^{2}=2 b c \cos A
\]

已经非常简化了，和所证只差一点了，再能干啥？要证边，自然把 \(\cos A\) 变成边，就是余弦定理，至此问题结束。\\
如果把刚才的答案整理下来，给别人看，那么对方可能会感觉：为什么要这么变形？其实，对于简单题和中档题而言，主动去寻找自己能干什么，然后朝着那个方向迈步尝试，是比较容易得到令人满意的结果的，这和题目难度本身的限制也有关系。对于这类问题，积极设问带来的主动寻找目标的意识，以及找出 3 种左右的路子大胆尝试，有很大的概率把我们带入正轨。

那么，如何判断走的这条路子能不能带我们进入正轨呢？我们的原则是＂简化即前进＂，只要变形后的结果看上去茫然感更低，或者看上去更接近所求，或者更容易被基本要素拿捏（要素网的视野下的踏实感）．．．．．．．只要变形后的结果看上去给你的安全感更强，那么这个变形就一定是成功的。这种感觉（判断）有点玄学，但是相比于原来牵制的不稳定性，这种判断方式能帮助我们在茫然之中找到亮光，从而进一步逼近答案。

\section*{1．5．2．3 草稿纸的使用}
在答题速度方面，我们提到了隔离主要矛盾，迅速转移思维焦点来达到快速切割题目的方法，也为大家提到了遇到不太直白的牵制条件时利用主动设问寻找出路，以及简化原则判断试验是否有效的方法。这些方法都是针对大家考试时思维的行为的，是用来提高大家的思维敏捷度的，但是，不管什么方法都要依托于动笔演算。演算的草纸使用将会直接牵扯到答题的效率，胡乱演算，找不到演算痕迹，到处翻卷子，白白给自己增加不快，浪费时间，但是好的演算习惯却能事半功倍，即使是计算能力稍弱的同学，只要演算习惯很好，一样可以保证答题速度的稳定性。

本人在演算纸使用上的心得为：小计算用卷，大计算用纸。像前三题，算个坚式，解个方程，试算一种方法，这种步骤较短的过程，直接写在原卷上该题旁边进行即可。这样的话可以省去翻找草稿纸的时间，而且演算离题目本身距离很近，也有助于我们随时确认题干已知信息，减少抄错已知的情况出现。对于一些比较长的计算，可以在草稿纸上进行演算，比如化简一个较长的等式，但是一定要划分区域，并且不要写的太潦草。另外，答完全卷进行二轮检查的时候可以用草稿纸，除此之外尽量使用试卷本身演算。事实上，这样的话，在答题时，真正在草稿纸上完成的计算其实并不多，草稿纸本身也保持整齐，也能进一步的提高做题效率。

\section*{1．5．2．4 成为考场上的猎手}
我们已经看到了 130 分的思维层面的处理方法，但是，为了能够最大化的提高考场上的做题效率，最大程度的提高我们的做题速度，我们绝不能把思维搞僵化，搞得太过于机械。

比如，第一道选择题，拿动静视角和要素网能不能解释？当然能，但是我们本身对这类题目就相当熟悉，读完题之后就可以不假思索地动笔完成，那么这个时候，可不可以不使用思维去单独过一遍题，而直接下手秒掉这个题呢？当然可以，毕竟去分析一遍自由度也是需要消耗好几秒钟的时间的，反映到前几题上面的话，可能就是 \(30 \%\) 左右的时间占比。因此，在真实的高考考场上，大家应当以效率优先，用最舒服的方法去完成题目。

那么，我们讲到的思维的用武之地在哪呢？实际上，从刚才的分析足以看出，用特定的思维方式去分析题目，特别简单的题，效率相对偏低（因为本来就做得很快），随着难度上升，效率逐渐升高，在试卷的大部分题目

中都可以起到加快思维切入，给予自己做题踏实感与方向感的作用。如果你能在几秒钟之内看清简单题的要素网，或者自由度，那么这种从容感给考试心态带来的调节，也是不可忽视的。因此，我会给大家不断强调这些思维方式的价值。

而到了真实的考场上，一切都应当让位于效率，正确率这些偏向于结果论的事情。这个时候，积极的考试状态就是，以一个猎手的心态，以跃跃欲试，兴奋地猎杀题目的姿态，主动高效地砍瓜切莱。比如我读完一道小题，或者一个大题的第一问，我还没来得及去考虑自由度或者要素网，做题方法就已经在脑子里冒出来了，那么就别管那些了，直接上手干掉它，反而一些自由度的判断之类的可以在时间充裕的前提下帮助我们即时检查一下得出的结论合不合理，之类的。

简单的总结一下前述的关于提高答题速度（效率）的注意点。大家平时把良好的思维方式学清楚，在考试时候的绝大部分题（接近 130 分）都可以极大提高稳定性与切入速度，从而提高效率。这些习惯的培养（包括思维本身的练习，以及各种答题习惯）需要大家在平时坚持不解的努力（这一点在第四部分还会重新阐述）。当你真正踏入考场后，就应当以结果论为优先，答题时遇到直接会做的题，就上手直接做掉，以一种猎手的姿态，这样的话，效率会进一步提高，并且考场的气势打出来了，考试心态和考场状态也会有相应的提升。

\section*{1．5．3 让兼顾正确率成为习惯}
得分受到两大因素的制约——速度与正确率，下面，我们来讨论一下正确率的问题。影响正确率的主要因素是失误，失误主要分为审题错误，抄写错误和计算错误等。为了尽量高效地规避失误（也就是降低失误的概率），显然就需要我们做好两件事：正确审题与即时检查。

在审题方面，对于一般的问题，如果大家在读题的时候都是使用动静视角和要素网的话，是很难出现漏看条件的情况的，因为漏看一个就会导致自由度出现问题，导致条件太少无法做题。当然，由于不等关系不会改变自由度，因此在做题的时候，不等关系是肯定要重点关注的，很多错误都和忽视不等关系（也就是忽视定义域）有关。而且，为了防止意外发生，还是建议大家在读题的时候把所有已知条件圈一下，做题时方便随时视线上瞥查阅，用圈相对来讲比画线醒目，容易被视线抓住。

面对一些长题干的问题，尤其是概率统计，大家可以试着用读小说／科普／故事之类的心态去阅读，这样的心态下，阅读更能集中注意力，同时有助于心态的放松以及考试状态的调整。当然，在这样阅读的同时，也要记得把正常读题需要做的事情做了，也就是圈出已知和重要信息。另外，这种长题干的题目，可以多读 2－3 遍来加深印象，确保不遗漏，看错条件。

下面我们重点说一下即时检查，即时检查的作用是将出错率变为原来的平方。实际上，这在本章第三节里已经讲过，讲解了一些即时检查的小技巧。为了能在考试中实际运用即时检查，我们需要做的事情是将即时检查变成习惯，将即时检查变成正常做题的一部分，只有这样才能够最大限度的发挥即时检查的价值。

那么，怎样才能做到这一点呢？我们来介绍一下即时检查的注意事项。首先，不管是什么题，最好都要在做完时余光回瞥一下题干，看看有没有没用上或者看错的条件，因为我们之前圈了圈，所以已知都非常醒目，余光一瞥最多消耗不了几秒。其次，只要即时检查计算问题，一定要让检查中的计算和做题过程的计算发生本质性的改变，因为只有这样才能实现错误率的平方下降。

下面，我们对考试题目的即时检查进行一些具体说明，运用较多的计算方面的检查仍然参照第三节的内容，例如除法换乘法，方程验根，求出的结果带回去检验是否满足题干条件等，都是广为使用的通法。对于一些计算较多的题目，可以从中选取 2 个左右的关键步骤进行检查，例如中途求出一个函数，是否求对？中途算出一个关键要素的值，是否正确？单选题具有选项，如果是计算型选择题，而你恰好计算出这个答案，那么一般来讲会是正确的，这一点的理解会给我们带来做单选题更多的安心感。如果要检查单选题的话，可以去说明其他

选项是错的，例如集合题，可以通过几个特殊值（直接从选项里获取）是否在集合内，来判断是不是确实选你那个答案，再例如求取值范围，也可以带人端点值判断（注意：有的同学可能会希望直接这样来做选择题，并不推荐，安心感差而且出错率未必好看）。数列题标准地验证前三项来判断通项公式是否正确；立体几何通过倒回去乘来检验法向量是否正确；偏几何的问题可以通过代数，几何的不同方法来改变计算过程．．．．．．．

即时检查根据不同的题目类型，会有各种各样的方法，大家在平时做题练习中也要自己主动思考，寻找一些实用的小方法。

良好的审题习惯和即时检查习惯的养成对于我们提高正确率有着极为强大的作用，但是其养成非常需要时间的积淀。相对来讲，对于性格谨慎的同学而言，这一习惯会更容易养成。如果懒得养成即时检查的习惯，那么就要相应的承担原有的错误率带来的风险。实际上，只要掌握了那些即时检查的小方法，那么即时检查并不会耗费我们太多的时间，就相当于在实际做题的演算中额外插入一点检查演算，即使是步骤多的题目也不过是 2 处左右的额外计算，小题更不必说，有的即时检查靠目测与口算就可以完成。

\section*{1．5．4 浅谈科学提分}
\section*{1．5．4．1 时间节点与考试节奏}
刚才的内容涉及到了思维焦点，主动设问，草稿纸使用，即时检查等一系列思维与行为的习惯，这一切都是为了让我们在考试中拿到一个更好的分数。为了这个目标，还有一个十分重要的问题——考试时间。我们能否在考试中保质保量地答完试卷？为了解决这个问题，我推荐大家采用记录时间节点的方式来把控自己的考试节奏。

有关时间节点的问题在本人关于考试时间规划的杂谈中已经详细讲解过（推荐大家去那里自己了解），简单来讲，就是自己在考试中做到几个关键位置所需要的时间的期望值，例如，答完选择填空用时大约 40 分钟，做完除圆锥曲线导数还剩 35 分钟，等等。这个数字需要大家在自己的每一次考试中进行记录，通过大量考试的统计结果分析得出。与之相伴，我们还可以计算方差，从而观察自己考试节奏是否稳定。

通过认识自己的时间节点，就能了解自己在哪些题目上更容易卡壳，就能在未来更加有针对性的进行训练；就可以在培养我们之前讲到的习惯时，留心并看出时间节点的变化，从而看出这段时间的习惯养成，题目训练，反思巩固是否改善了自己的考试节奏，让自己不再慌乱，到处卡壳等。

需要注意的是，时间节点的统计需要把即时检查的时间也一并算上，这样才能准确看出即时检查在考场上是否会影响你的答题。假如真的因为即时检查而导致原本勉强做完的题做不完，并且怎么反思改进都无法改善这一问题，那么也许就不得不压缩甚至在部分题目（选自己失误率低的）放弃即时检查。

另外，如果在真实考场上时间节点波动该怎么办？这里可以提供一个 5 分钟原则，也就是如果某个位置的时间超出时间节点 5 分钟以内，那么对考试节奏的影响都比较小，不需要太过于慌张的进行追赶。另外，可以根据自己平时答完全卷的时间进行估计，比如本来能剩 20 分钟，这次选填超了 10 分钟，那么预计答完卷子就还剩 10 分钟，答得完，而且还有时间二轮检查。这样思考可以帮助我们稳定心态，当然，如果超过了 5 分钟，建议最好稍微抓紧一点，适当加快一点答题速度。

关于考试节奏，还有一个因素就是遇到不会的题卡壳怎么办？这里为大家推荐 2 分钟原则，也就是，如果这道题2分钟还没有思路（注意不是没做完，也不是还在试验中，而是一点思路和想法都没有），那么就直接跳过。

有的同学可能会认为，如果跳了好几道题，会不会耽误时间？实际上理论上是不会的，因为跳过的题只耽误了 2 分钟，还减少了实际做这道题的时间，因此时间节点的变化等于 2 分钟减去做原题号的题的时间，因此实际＂浪费＂掉的时间只有那些 2 分钟。做完别的回过头去再看也一样，只是时间的平移。当然，跳题时的时间节点也要相应的减去一些。反而在那道题目上继续耽误，极大概率耗费更多时间。2分钟原则充分反映了理性的

精神。\\
从刚才的分析可以看出，时间节点配合 2,5 分钟原则可以充分掌握我们的考试节奏，但是，考场仍可能产生很多突发情况，最常见的就是刚才说的两种，时间节点出现问题，以及跳题，虽然刚才的分析表明了处理方法，但是在紧张的考场上，心态很可能受到影响。为了减少心态的影响，就要求大家在平时就尽可能扩大考试的样本量，通过限时练的方式获取大量数据信息，让自己通过大量的实践充分信任自己的时间节点。恐惧来源于未知，但是如果我们对自己的能力情况有大量的科学的数据作为支撑，那么考试时，我们就有科学的理由对自己进行积极的心理暗示了。

\section*{1．5．4．2 试错与习惯养成}
从上一条关于时间节点的分析中，大家可以看出，大量的试验以获取数据是必要的。而我们前面所讲到的思维方式，即时检查，也都是这样。我告诉大家这些习惯是什么，原理如何，效果如何，应该怎么培养，但是真正的提高，一定是需要每一位同学自己去努力。这些东西和快节奏的学习方法吹水，华丽的二级结论欣赏不同，需要长期的努力和实践才能够有所收获，不可能第一天听完，第二天就熟练上手。我们如何努力，才能够培养出这些习惯呢？

科学的努力才能结出果实。我为大家极力推荐的方式就是对比与试错。当你拿到一道自己本应期望做出的题，却不会用前述的思维方式去分析，或者找了好几条路都没走通，又或者根本就是完全没思路的时候，怎么办？首先，尝试自己去做，然后，对照答案，尝试使用我们讲到的思维范式去分析答案的来源，这样的话（这相当于开卷考试了，难度骤降），就能得到一份符合我们所讲的思维范式的答案。然后，将自己的思维过程和这个答案进行比对。因为我所讲到的思维范式都是有理可循，逻辑清晰的，所以将自己当时的思维进行对比，找到从哪个地方开始发生的差异，然后仔细琢磨这个位置，自问为什么当时没有想到？是元问题不会分析？充要替换牵制条件的时候这条路根本就没有想过？就能够发现自己的思维上的漏洞和不足，这个不足就是之后需要反思，总结的地方。

而像即时检查与时间节点的配合，审题，草稿纸使用这些方面，则需要在长时间的练习中养成习惯。通过限时练的方式，不断调试，优化，改进这些习惯。想要改进一个时间节点？那就刷一部分题来作为试验，尝试这种方法是否有驾驭的可能；考试算错一个题而想要研究该题型的即时检查？那就刷一部分这个题型，研究一下错误率高的地方，想想这个地方怎么即时检查能又快又准，然后继续做试验，改进效果．．．．．．这就是所谓的试错，我们的所有习惯养成与学习经验，都是在大量的样本中进行分析，进行不断的微调而逐渐获取的，可以说，没有大量科学的试验和试错，就无法定位自己的不足，就无法找到科学的提升路径，就无法取得真正的进步，就只能被困在假努力的陷阻中。

对比试错是阳哥高三阶段用于提高语文和理综的手段，成效显著（可以参见部分杂谈），所以将其修改成数学适配的方案提供给大家。这些方法都是大家随时可以上手实操的，但是若要收获成效，则需要长时间的不断持续，至少要 1－2 个月甚至更长，才会慢慢有了感觉。

\section*{1．5．5 图的信息量}
图的信息：信息量尽量以直观的视觉形式呈现，降低思维容量。比如多画示意图示意表之类的，审题圈出要点方便即时回瞥检查等。

\section*{总结}
本章节内容与前四节内容的风格不同，本章节内容的重心放在较为功利的考场实战上，并且干货偏多，内容多且零散，需要观众们花费一定时间进行消化。

本节重点介绍了如何提高考场答题速度。我们介绍了使用前四节的思维方式来进行考场答题效率提升的办法，通过要素网的结构阐释了答题行为重点是分析元问题和牵制的过程，并且分别对二者的处理范式进行了说明。二者的核心都是基础知识，前者侧重于直接应用，后者难度更大，可能存在一定的技巧性。为了能够让基础知识的掌握程度提高，我为大家讲解了如何学习与复习基础知识，以适配我们讲到的思维。这种学习方式也能让大家更加深刻地理解知识本身，但是需要覆盖到所有知识，是一个大工程，需要大家在学习和高三复习中慢慢努力去完善。

为了真正提高答题速度，我们讲到了思维焦点及时跟进转移，迅速当成如同默写题一般，来处理当前主要矛盾的作战思路，以及面对牵制时的尝试与＂简化即前进＂的信念。答题效率的提升反映在时间节点上，通过长期试错带来的时间节点优化，我们可以了解到自己的答题速度有了怎样的提升。当然，考场形式千变万化，如果遇到棘手的问题，则应当遵循 2,5 分钟原则，并且相信其背后的时间节点方面的逻辑支持，这样能尽量降低心态对我们的影响。

除了答题速度以外，良好的审题习惯与科学高效的即时检查是确保正确率的关键。即时检查的原则是让计算发生本质性改变。这些因素也要记得纳人时间节点的规划与调试中。同时，为各类题型寻找合适的即时检查方法，以及养成上述的各种习惯，都需要大家在平时抓住学校训练的犬量机会，通过大量的这样的试验（典型如限时练，模拟考等），去不断对比，试错新想法，并及时进行结果回顾（也就是反思），以及优化的调试。这样才能够在几个月甚至半年的时间内逐渐找到进步的感觉，最终取得高考的胜利。

大家可以看到，第四部分讲到的东西是贯穿整个学习过程的，大家无论是在学习思维，学习知识，还是在做题练习，反思巩固中，都需要不断坚持科学理性的提分方式。不论是思维的养成，还是习惯的获得，还是分数的提高，都不存在任何捷径，只有科学与不科学，效果好与效果差的区分。那些幻想着靠迅速背点二级结论，看一两次课程，听懂一点思维，就让成绩快速飞升的人，最终只能面对理想与现实的差距，并接受进步甚微，甚至毫无进步的事实。希望大家在阅读本节讲义之后，能够对于数学的讲步问题有一个科学而清醒的认识，通过科学的方法努力，通过优化各种细节提升实力，最终在考场发挥出真实的水平。

\section*{1.6 封装命题}
事前声明\\
本节讲义是关于＂自由度＂这一概念的补充内容，分为三部分：第一部分是自由度的整体回顾，涉及到本章前五节的内容；第二部分是一个关于自由度的易错点，涉及到自由度这一概念所谈论的对象究竟是谁的问题，并衍生出关于函数思想的本质解释；第三部分是一大类在题干叙述逻辑上有共性的＂题型＂，up 主称其为＂封装命题＂。

本内容需至少吸收本章前四节后才能观看，最好是看完了 \(1 \sim 5\) 节全部，不然的话在吸收和理解上会大大折扣。

\section*{1．6．1 自由度回顾}
我们在本章第一节中引入了自由度的概念，并在第四节中通过要素关系网，对这一概念的本质进行了剖析。自由度的字面意思，是衡量某个事物自由程度的参量，所谓＂自由程度＂的严格表述是：控制某个事物所需要的独立的参数的个数，例如，正方形的面积受到其边长控制（控制关系是面积 \(=\) 边长的平方），我们就说面积的自由度是 \(1 ;\) 但我们不能说因为面积 \(=\) 边长乘周长除 4 ，就判定面积是两个自由度，因为周长和边长不是独立的参量。

那么，何为独立？这里的独立的意思是，对于选取作为控制操纵杆的参数，它们彼此之间无法列出等式关系来彼此约束。例如，方形面积 \(=\) 两个邻边，你可以随意选取它们的值；但是正方形面积 \(=\) 边长乘周长除 4 ，边长和周长之间显然存在等式关系的约束，那么就不独立。

思考：长方形面积 \(=\) 长乘宽，长和宽之间有关系：长 \(>\) 宽，这是否说明它们会影响到长方形的自由度？

在阅读完第 1,4 节讲义后，我们知道了自由度存在两幅面孔：一个是要素关系网中的自由度。列出要素关系网后，各个要素，参数都在一级一级的控制中被确立。假如暂时不看已知中的牵制关系，那么要素网中就有起始的要素（不被任何箭头所指，自身的指定就是自由的）和被控制的要素。起始的要素可能在要素网开头（题干指定一个背景），也有可能中途杀出个程咬金（中间我突然取一个点，一个新参数等），它们构成了要素网的＂操纵杆＂。其他要素就被这些起始的要素严格地控制住（起始要素定了，这些其他的要素也就顺次定了）。这种＂定＂的方式，在第五节所说，就是所谓的元问题，需要大家按照第五节的学习理念，花长时间积累，去重塑自己对所有基础知识的理解。

另一个是系统实际的自由度。当我们回到题目，把要素网以外的已知加进来，那么一些牵制关系就会卡住控制杆，让它们没法再彼此随意移动来＂肆意妄为＂地控制全局。原则是一个等式减少一个自由度，不等关系不减自由度但是约束操纵杆的移动范围。所有操纵杆的自由度的和，减去所有牵制关系所减掉的自由度，就是这道题目真实的自由度。这个自由度的用处就是判断系统是运动的还是静止的，用来给我们一个在读完题后的大方向判断，除此之外没啥用处（因为除此之外给不了我们实质的信息）。

下面，我们对自由度进行一些补充。

\section*{1．6．2＂自由度＂概念的作用对象}
1．6．2．1 作用对象与着眼点\\
首先，我们来考虑一个问题

例 1.25 给定直线 \(x+y=1 \cdots \cdots\)

其余的条件都暂时不给，请问这个例子是几个自由度？\\
也许你会说，这条直线是确定的，当然是静止的了！但是，我反问你，这里有 \(x, y\) 两个参数，一个等式，那自由度不应该是 \(2-1=1\) 吗？

与之类似的问题最近经常被乐于思考的粉丝问起，类似的问题还有：给出一个圆 \(x^{2}+y^{2}=a^{2}\) ，这个圆算几个自由度？从参数来看，应当是一个 \(a\) 作为自由度，但 \(x, y, a\) 三个变量，一个等式，自由度似乎又是 \(2 \ldots . .\).

这确实是我们以前没有重点去强调的一个问题，也确实会让大家产生误解。实际上，我们从朴素的本能出发，思考一个问题：谈论自由度是不是需要明确对象呢？是不是需要先搞清楚，我们是在讨论谁的自由度呢？

当然！这肯定是必要的，我们生活中谈论一个特征，肯定要首先明确是在谈论谁的特征。比如，我们谈论考试成绩，一定是谈论某一个人的成绩，某一个班的成绩，某一个省的成绩．．．．．．．不明确这一点，谈论成绩就没有意义。比如，谈论个人的成绩，着眼点就在这个人成绩进步退步？各科分数，失分原因？谈论班级成绩，就要关注班级平均分，整体进步退步情况，尖子生，退步生的数量……谈论一个省的高考成绩，还要考虑分数的区分度体现的好不好．．．．．．．

回到例 1.25 ，我问的问题是，这个例子是几个自由度，也就是把直线整体作为一个对象谈论，这个直线本身当然是静止的。

那么，另一种理解是从哪儿来的呢？仔细观察，会发现第二种理解的着眼点是在 \(x, y\) 参数上的，也就是说，第二种理解关注的是直线上的点 \((x, y)\) 的坐标！这个坐标被两个参数 \(x, y\) 直接控制，并且需要满足牵制关系 \(x+y=1\) ，因此自由度就是 \(2-1=1\) 。

因此，这两种理解都有它的道理，区别在于讨论自由度的具体着眼点，具体对象不同。第一个（也就是例子的自由度）是把直线整体作为讨论对象，第二个则是随便取了直线上的一个点 \((x, y)\) ，作为个体去考虑它的自由度，也就是：

例 1.26 给定直线 \(x+y=1, P(x, y)\) 是直线 \(x+y=1\) 上的任意一点．．．．．．

到了例 1.26 ，直线静止，但是它身上的 \(P\) 却有一个自由度，因此，两个问题本质是不同的问题，它们的结果就自然有了差别。

从上面的例子中，我们可以看到，研究自由度，必须要首先明确你在讨论谁的自由度，只有这样才不会在复杂的题目中产生混乱，才能够清晰地搞清楚题目的结构。

让我们回到上一个模块里讨论到的两种角度上的自由度——要素网中的自由度，以及系统整体的自由度。这两个维度的自由度，很明显，在要素网中看的那个，是分开考虑各个要素的情况，分析的着眼点在系统的结构本身，对于我们理解系统结构具有重要意义，显然更为重要。因此，之后大家在分析自由度的时候，最为推荐的思路，是在要素关系的理念下，去考虑系统中的各个主自由度，并以此去把控要素网，这也是最为贴合自由度理念的初衷的做法。虽然也可以去提前计算题目整体的自由度（也就是系统实际的自由度），但是着眼点放得非常大，放眼于整个题目，这只是一个非常粗浅的对题目全局的判断（只能告诉你题目整体最终是动还是静），是一个视野上的行为，当读题结束，思考完成，明确题目最终大致是怎样后，就可以不再去过多关注。

或者换个比喻，第一种理解是坐火车，细致进行考察；第二种理解是坐飞机，告诉你世界俯㒈是怎样的，但是无法告诉你世界中的如画风景与人间冷暖。坐飞机是有意义的，云上的风景与俯瞰的开阔感让我们胸有成竹，但后者是我们在解题时实实在在需要的，也是大家更需要重视的。

\section*{1．6．2．2 函数思想——镌刻时光的胶卷}
2.1 举出的两个例子代表例分析自由度的两种典型的对象意识——整体与个体。点动成线，线动成面，面动成体．．．．．．．点是有自由度的，但是点组成的线却是静止的，从此可以看出，同一个整体，不同的观察方法（整体考

察与个体考察），导致的自由度也不一样。这是一个很强大的地方，因为我们有希望利用这一点来增添或者削减自由度，这其实就是大家一直喜欢讲的函数思想。

例 1.27 在射线 \(O x\) 上，一个动点 \(P\) 从 \(O\) 出发，向着正方向运动，以 \(O P\) 为一边，作正方形。记 \(t=|O P|\) ，正方形的面积为 \(y\)

我们可以轻易的写下函数关系 \(y=x^{2}\) 。按照刚才说过的思维方式，如果谈论 \(P\) 点，它是一个自由度的动点，无论怎么看这都是一个运动的问题。而我们看函数 \(y=x^{2}\) ，我们既可以从函数原本的含义理解（变量之间的控制关系，动态视角），也可以从静态的角度理解（这是一条抛物线啊）。

如果我们把一个控制关系看成整体，那么它就是一个（半个）自由度为 0 的抛物线。＂抛物线＂这个事实本身便把一个动态的问题整个吸纳了进去，我们单纯看一个静止的抛物线，就是在看整个运动过程。好比儿时的一段美妙的游乐园时光，被父母用相机录了下来，以胶卷的形式保存。

反之，当我们想要回忆儿时与父母去游乐园的时光，我们就可以重新播放胶卷，将那时的点点滴滴展现出来。当我们想知道 \(x=1\) 时发生了什么，我们直接对着抛物线的那一捆＂胶卷＂去看即可：代入 \(x=1\) 时即可知道那时的面积；计算导数即可知道那时正方形面积增长的多么快．．．．．

用动静视角和要素网的观点重新审视函数思想，可以发现，它就是这么简单并且浪漫。函数思想体现的就是把一段运动过程制作成静态的＂胶卷＂，以及把一段静态的＂胶卷＂重新播放，提取信息的过程。因此，什么时候用函数思想，怎么用函数思想，这些问题的答案也＂尽在不言中＂：当我们遇到一段运动过程时，主动的尝试把这一过程提炼成一个可以静态描述的东西（可以是函数，但也未必非要是标准的函数，只要是一个静态的，可以摸得着的事物就行）；当我们想要从一段动态过程中提取信息时（这里的信息涵盖很广，主动的提取信息分析，以及被动的往上添加已知的牵制条件去翻找一个特定的信息），我们就把得到的＂胶卷＂翻出来，并带着静态的眼光去寻找，审视。

\section*{1．6．3 特殊又普遍的题型}
刚才说到的自由度的对象意识，是有关自由度的一个注意事项。下面要介绍一种直接分析自由度并不容易，但是仍然具有重要意义的题目种类——封装命题。这一种题型分类方法是从逻辑上进行的分类，具有至高无上的普遍性。

\section*{1．6．3．1 什么是封装命题}
这个词语是 up 主自己编的，反映的是数理逻辑中多变元命题的现象。\\
我们先看两个小例子：\\
例1．28 对实数 \(x\) ，若 \(x \geq 2 \ldots \ldots\) 。

例 1.29 对实数 \(x\) ，若 \(\exists y \in(0,+\infty)\) ，满足 \(y+\frac{1}{y}=x \ldots\) ．

这两个例子其实都是在说一个自由度为 1 的变量 \(x\) ，并且都在给它加上限制 \(x \geq 2\) ，但是例 1.28 是我们熟悉的，直白的限制条件，而例 1.29 似乎绕了一点。

我们下面从命题的结构上来对此进行分析。首先，让我们了解一点关于逻辑的知识。说到＂命题＂，大家一定不陌生，题目中的已知条件也可以看作命题，比如 \(x \geq 2\) 之类的。而命题又有一个十分重要的特性，就是真假性，给出命题，我们应当能说出它是真还是假，比如 \(1 \geq 2\) 为假， \(3 \geq 2\) 为真。那么，\(x \geq 2\) 是真还是假呢？

这就和 \(x\) 这个参数（变元）的取值有关了，\(x\) 取不同的值，这个命题的真假性也会改变，这个时候，我们就说一个命题里含有了变元。有变元的命题，它的真假性会受到变元的左右。我们如果用 \(P\) 表示命题，那么刚才说到的 \(x \geq 2\) 就可以用 \(P(x)\) 表示，这就是说 \(P\) 这个命题含有变元 \(x\)（或者说，含有一个自由度 \(x\) ）。

按理说，含有一个变元的命题（例如 \(P(x)\) ）应当是被这个自由度 \(x\) 控制的，至少直观上，这是一个运动的，有自由度的事情，不是一个＂静态＂的事情。但是，请看下一个例子：

例1．30 给定命题 \(Q: \forall y>3, y+\frac{1}{y}>\frac{10}{3}\)\\
这个例子乍一看，\(y\) 应当是充当一个变元的角色，\(Q\) 应当是有一个＂变元＂的，应当是有一个自由度的。但是，感觉上总有点怪怪的。 \(Q\) 它说的是，任取一个大于 3 的 \(y\) ，总是满足那个不等式，看上去像是在陈述一个既定的客观事实，既定的客观事实从感觉上来讲当然应该是静态的，不变的！这是咋回事呢？

事实上，我们类比一下上一部分讲到的＂整体与局部＂的问题，我们就可以看到，如果考虑命题 \(P(y): y \geq 2\) ，那么我们是考虑各个具体的 \(y\) ，这样的 \(y\) 充当了自由度的角色，我们可以看到，它的运动是完全自由的，不受约束的：但是在例 6 中，我们是将 \(y>3\) 作为一个整体来考虑，这时候虽然 \(y\) 看似是自由的，但是我们却不能把它直接理解为自由度，而应该把整个 \(y\) 的运动过程看作是一个整体的＂胶卷＂，例 6 中的 \(Q\) 正是考虑了这整个胶卷的性质：处处满足那个给定的不等式。

起到上述的＂封装胶卷＂的作用的，正是 \(\forall\) 这个符号，它的含义是：＂胶卷＂上处处成立；另外，在例1．29中，施加给 \(y\) 的是 \(\exists\) 这个符号，表示＂胶卷＂上至少有一处成立。我们在高中学过，这两个符号都是量词（全称量词和存在量词）因此我们可以说，量词的意义就是封装一段运动的＂时光＂，把一段运动给刻成＂胶卷＂，变成静态的整体。这和我们上一部分讲到的＂函数思想＂看似大相径庭，实则不谋而合。

被量词封装了的变元（自由度）在＂胶卷＂内部是可以运动的，而我们考察的是整个＂胶卷＂本身。逻辑学中，把例 1.28 的这种不被量词封装的变元称为自由变元，把例 6 的这种被量词封装的变元称为约束变元。而在例1．29中， x 是自由变元，\(y\) 是约束变元，\(y\) 被封装的＂胶卷＂是用来描述 \(x\) 的特性的，最终矛头指向的，是 \(x\) 。

我们就把例1．29中，描述 \(x\) 的这样的命题称为封装命题。封装命题的判断准则就是是否含有量词，被量词封装起来的约束变元应当不再看成自由度，而应该把它的运动作为整体进行考察。因此，例1．29中实际上只有一个自由度 x ，而 \(y\) 的运动则生成了一个胶卷（存在一个正数 \(y\) 来配合 \(x\) ），最终还是用来刻画 \(x\) 的性质的。

另外，被封装的约束变元也未必只有一个，可以有更多的封装情况。

例 1.31 对于 \(\forall x>0\) ，都 \(\exists y>0\) ，使得 \(x>y\)

可以看到，在这个例子里面，\(x, y\) 都是约束变元，整个命题是一个静态的客观事实，但是＂胶卷＂相当于封装了两次：首先，＂\(\exists y>0, x>y "\) 是第一层，\(y\) 被封装了起来，只看这个小命题的话，可以写成 \(\exists y>0, P(x, y)\) 的形式，这里的 \(P(x, y)\) 就代表了 \(x>y\) 这个含有两个变元的命题，这时候我们可以整体把它封装成 \(Q(x): \exists y>0, P(x, y)\) ，然后例 1.31 的原始命题就成了 \(\forall x>0, Q(x)\) ，这就是两次封装的情形。更多的约束变元的封装命题也是类似的。

封装命题广泛存在于高中数学中，如恒成立，不等式有解等都可以看成封装命题；在高考题目中，封装命题也十分常见，从简单题到压轴题都可以遇到而且难题中也不乏多个约束变元的封装命题。下面，我们讲解一下封装命题的一般的处理思维。

\section*{1．6．3．2 如何处理封装命题的一般思想}
我们先从最简单的一类封装命题——恒成立问题出发。我们在＂本章第二节＂的最后一部分中讲到了遍历法的思维。比如：

例 1.32 对于参数 \(a\) ，函数 \(f(x)=x^{2}\) ，要求 \(f(x) \geq a\) 恒成立

这样一个命题实际上就是：\(\forall x \in R, f(x) \geq a\) ，因此 \(x\) 是约束变元，被封装。这样的话，原本的自由度 \(a\) 身上还绑着一捆胶卷，更加笨重。但是这种封装命题很多情况下都是要考虑真自由度的取值情况，所以遍历法的思想就十分方便。

我们在实数 \(R\) 上活动 \(a\) 这个滑块，一边滑动一边观测命题的情况。（先随便画个 \(f\) ）如果 \(a\) 非常小，比如 \(a=-100000000\) ，那显然是没问题的，然后把 \(a\) 往上推，推的太高会发现＂胶卷＂炸了，不能普适于 \(x\) 了，所以精确调整一下，发现充要替换是 \(a \leq f(x)_{\min }\) ，就完成了。

这其实是封装命题的最常见的思维方式，也是一种不容易混乱的思维方式。通过这种遍历的意识，我们在不同的 \(a\) 的位置分开观测命题，从而避免了自由变元和约束变元交叉在一起，混淆视听的情况。

而当题目难度进一步加大时，被封装的命题可能会更加复杂，这时候，单纯的遍历的观测可能就不够用的，因为快速的遍历看到的东西可能太浅了，我们需要更加仔细的进行观测。不过，这种移动自由变元来进行观测的基本理念还是需要的，并且这种理念不论处理多么难的封装命题都是很有价值的。

我们看一个比较难的例子：

例1．33 设直线 \(l: 3 x+4 y+m=0\) ，圆 \(C: x^{2}+y^{2}-4 x+2=0\) ，若在圆 \(C\) 上存在两点 \(P, Q\) ，在直线 \(l\) 上存在点 \(M\) ，使 \(\angle P M Q=90^{\circ}\) ，则 \(m\) 的取值范围为（）

读完题会发现，真自由度只有 \(m\) ，另外三个单自由度动点 \(P, Q, M\) 都被量词 \(\exists\) 给封装了，矛头指向的还是 \(m\)。我们先按照第五节的理念，把默写步骤做完：圆是 \((x-2)^{2}+y^{2}=2, m\) 对直线 \(l\) 的控制方式是定斜率 \(\left(k=-\frac{3}{4}\right)\)的，全平面的平移。按照刚才的理念，我们需要一边尝试移动 \(m\)（根据刚才的控制方式的分析，充要替换为暂时只考虑直线平移，一边观测：什么时候能找到那样的组成直角的三个点呢？

这个时候，因为牵扯到的要素比较多，我们直接全局遍历也看不出什么来。这个时候的一般思维是：先滑动滑动直线，随便固定一个直线位置（随手一画，画个图就行），看看这时候有没有希望达成题目要求，并同步思考：什么样的直线能达成要求？实际上，这正是我们之前讲的两种思想一一冻结与设问的结果合（之后总结时会再展开）。

即使是这样，仍然有三个点，其中 \(P, Q\) 在圆上，\(M\) 在直线上。这个时候，我们很自然的改写一下命题，做一下二次封装（最外层的应当简单，这样方便冻结）：\(\exists M\) ，使得 \(\exists P, Q\) ，满足 \(\angle P M Q=90^{\circ}\) 。于是重复之前的遍历，冻结，设问：随便先点个 \(M\) ，想想什么样的 \(M\) 才能存在 \(P, Q\) 张成直角？

后面这个加粗黑体的问题，一旦被我们通过设问的方式隔离出来，我们的思维焦点只要放上去，就可以立刻看到，这是一个＂元问题＂级别的东西，只和 \(M\) 相对于圆的位置有关。下面，我们必须把这个元问题分析清楚。

什么样的点，能对着圆张出一个直角？这已经是考察大家对于解析几何的基本功了。不难看出（这时候开始在脑子里全局的遍历 \(M\) ），如果 \(M\) 在圆内或者圆上，可以随手找一组 \(P, Q\) 张一个直角出来；如果 \(M\) 在圆外，那我先象征性动一动圆上自由运动的 \(P, Q\) 点，看看张角 \(\angle P M Q\) 能是个什么情况（这时暂时 \(M\) 先放某个位置看看）。发现，\(P, Q\) 收的足够近，张角就很小；张角撑到最大就是 \(M\) 对圆的两条切线的夹角。因此，只要 \(M\) 对圆的两条切线的夹角能撑到不小于 \(90^{\circ}\) ，就能往里收一收，收出一组符合条件的 \(P, Q\)（求范围的题注意检查充要替换，发现没问题）。因此，\(M\) 位置的充要替换就是：\(M\) 在圆内或圆上或圆外（要求切线张角不小于 \(90^{\circ}\) ）。这个时候，我们已经把那个以 \(P, Q\) 为约束变元的封装命题（实质是指向 \(M\) 的命题），充要替换成了真正意义上的指向 \(M\) 的命题，因此，约束变元 \(P, Q\) 就被我们消去了。

至于后面的事情，既然一层封装已经被我们降维打击成了普通命题：直线上存在 \(M\) ，使得＂\(M\) 在圆内或圆上或圆外（要求切线张角不小于 \(90^{\circ}\) ）＂，接下来我们把思维焦点放回来，继续进行充要替换，这又是元问题级别的基本功考察。\\
＂圆内或圆上＂等价于＂直线和圆有交点＂，这个后续的充要替换已经是究级简单的基础知识；如果直线和圆没交点，那就是＂直线上存在 M ，使得 \(M\) 对圆的切线张角不小于 \(90^{\circ} "\) ，引号内的部分显然可以再试着简化成更好算的东西。什么样的点切线张角不小于 \(90^{\circ}\) ？再去随便点几个 \(M\) 看看，显然离着圆越远的点，张角越小，所以只要 \(M\) 离圆心比较近就行了。具体要近到什么程度？那肯定是算出恰好是 \(90^{\circ}\) 的那个距离，然后比这个距离小就行了。这又变成一个基础题了。

最后，我们实现了对 \(M\) 全方面的简化，接下来要做的就是处理好这些小元问题，解好这些不等式，求好直线可以允许的位置范围，然后回去求 \(m\) 即可。最终答案是 \([-16,4]\) ，具体细节留作练习。

从这个例子，我们可以看出：对于封装命题，其处理的统一宗旨就是：既然约束变元最终矛头还是指向那些真自由度，我们需要做的就是把这些带约束变元的的封装命题给充要替换为一般的普通命题。这个过程由于封装命题引入了更多的变元（多余的约束变元），我们需要把它们消除，而由于变元比较多，常用的思维方式就是遍历，冻结，设问，处理元问题的一套循环流程。用上遍历真自由度的理念，去主动移动自由变元的＂滑块＂，看看不同位置会对命题产生什么影响。有时候如果不方便直接看出，就随便先点几个真自由度的位置，然后去试着以静止的思维分析一下封装命题（就比如刚才的，随便画一个直线，看看能不能找到 \(M, P, Q\) 满足要求）。在这个过程中，主动自问（设问），什么样的直线能找到这样的 \(M, P, Q\) ？把思维焦点主动往这个方向（去尝试进行充要替换）去引导，从而抽离出真正的元问题，进而消去约束变元。

有时候题目比较难，就连约束变元也比较多，这个时候就像刚才一样进行一下二次封装，然后一层一层扒掉约束变元，重复这样的过程（需要在每消一层约束变元之后就及时切换思维焦点），迟早就能消除干净。

\section*{1．6．3．3 延伸：约束变元顺序}
我们对封装命题做一点进一步的探讨。刚才我们在例1．31与例1．33中看到了多个约束变元的情况，其中例1．33里出现了一组＂\(\exists P, Q\)＂，这按理来说应该是两个约束变元，按理说应当像我们说的那样，再去做一次冻结，比如再去冻结 \(P\) ，去观测 \(Q\) 。但是我们做题时，是直接对两个约束变元同时处理了，并没有经历＂分而治之＂的步骤。

为什么能一下子直接处理掉呢？关键点在哪里？我们来做个类比，排列与组合。排列需要区分不同事物的出现次序，但是组合则是只要这一筐子的东西一样，我就可以把它们看成一个整体，打包处理。因此，这里的关键在于次序无关：多个约束变元在命题里出现的次序不同，不影响命题的语义。这种事情有没有普遍性呢？我们再来回看一下之前举出过的例子：

例1．31 \(\forall x>0\) ，都 \(\exists y>0\) ，使得 \(x>y\)

如果例1．31 中的约束变元 \(x\) 和 \(y\) 次序无关的话，那么我们交换它们的顺序之后，命题的含义应当不变才对。我们交换一下试试：

例 \(1.34 \exists y>0\) ，使得对 \(\forall x>0\) ，都有 \(x>y\)

但是不难发现，例 1.31 是真命题，但是例 1.34 是假命题。这说明对于一般的封装命题而言，约束变元的次序是不能随意乱动的。其实可以想明白，这反而才是正常的现象。从语义上来说，约束变元出现的先后就相当于一个中，各个要素出现的先后。在控制关系的描述下，要素有初级次级的区分，当你建立了要素网之后，这个顺序就被确定了下来。但是，要素网之中也有同一级要素的说法，类似地，封装命题中按顺序出现的各个约束变元，也会有＂同级＂的说法，就像例 1.33 之中的＂\(\exists P, Q\)＂一般。这个在直观上并不难理解。

那么，什么样的约束变元是同级的呢？回看刚才的例1．31与例1．34＂，显然如果相邻两个约東变元的量词一个是＂任意＂，一个是＂存在＂，那么这时候就不能随便换序，不然会导致语义的改变。这个时候我们必须严格按照命题中约束变元出现的顺序，进行冻结研究。但是像例1．33那样的，两个相邻的约束变元的量词相同，我们就可以无视其顺序。其实所谓的 \(\exists P, Q\) ，也可以在语义的层面上改写成＂\(\exists P\) ，使得 \(\exists Q\)＂，它与＂\(\exists Q\) ，使得 \(\exists P\)＂在逻辑上是等价的。

下面，我们再来看另一个例子：\\
例1．35 \(f(x)=(x-1) e^{2 x}-\frac{4}{3} k x^{3}+k x\) ，若对 \(\forall k \in[1,2]\) ，以及 \(\forall x \in[0, k]\) ，均有 \(f(x) \leq c\) 成立，求 \(c\) 的取值范围

这道题目是又一个很好的理解约束变元次序问题的例子。这道题的命题设问中，给了两个约束变元 \(k, x\) ，它们都被全称量词 \(\forall\) 所封装。按照我们刚才讲的，这种问题的约束变元存在两种等价的次序。我们将分别为大家展示这两种次序。

如果我们直接思考这道题的话，题干先出现了 \(k\) ，然后才出现 \(x\) ，而且 \(x\) 的陈述里还与 \(k\) 有关，可以说 \(x\) 所在的胶卷本身还被 \(k\) 控制着。这时候，我们很自然地就想先冻结 \(k\) ，去研究这个关于 \(x\) 的封装命题。对 \(\forall x \in[0, k]\) ，均有 \(f(x) \leq c\) ，还是用之前的遍历法，容易看出就是求 \(f(x)\) 在 \([0, k]\) 上的最大值（它被 \(k\) 控制，记为 \(g(k)\) ），然后 \(\forall k \in[1,2]\) ，有 \(g(k) \leq c\) 。对 \(f\) 求导：

\[
f^{\prime}(x)=(2 x-1) e^{2 x}-4 k x^{2}+k=(2 x+1)(2 x-1)\left(\frac{e^{2 x}}{2 x+1}-k\right)
\]

（这样分离是基于我们之前讲的指数乘除的原则，见 up 主 2021 年 9 月 6 日视频），因此至少有一个零点是 \(\frac{1}{2}\) 。然后再令

\[
h(x)=\frac{e^{2 x}}{2 x+1}-k
\]

则

\[
h^{\prime}(x)=\frac{4 x e^{2 x}}{(2 r+1)^{2}},
\]

于是 \(h(x)\) 在 \([0, k]\) 上增。\\
这时候就得讨论一下了。 \(h(x)\) 有没有零点？与 \(\frac{1}{2}\) 关系如何？进而如何判断 \(f(x)\) 的单周性与最值？继续分析

\[
h(0)=1-k \leq 0, h(k)=\frac{e^{2 k}}{2 k+1}\left(1-\frac{k(2 k+1)}{e^{2 k}}\right),
\]

在经历类似的操作之后分析出 \(h(k)>0\) ，因此当 \(k>1\) 时，\(h(x)\) 也有零点，并且这个零点的显式表达式还写不出来。之后还要分析这个隐零点与 \(\frac{1}{2}\) 的大小关系。不过我们预判一下的话，肯定有一些区间段会导致 \(f(x)\) 先增再减再增，这时候讨论最大值会变得更加麻烦。足以看出这种解法的复杂程度。

既然这种讨论方法非常复杂，我们为何不试试将 \(k, x\) 换一下顺序呢？如果我们预判一下，先冻结 \(\forall x\) ，去遍历所谓的 \(\forall k\) 时，我们会发现，这时 \(f\) 本体是关于 \(k\) 的一次函数！最值肯定很好处理！这条路看上去十分诱人。

有同学可能会问，刚才不是说＂\(x\) 所在的胶卷被 \(k\) 控制＂吗？但是，作为量词相同的同一级约束变元，逻辑上已经保证了它们的顺序可以调换。我们把原命题改写成类似例1．33的并列语言：＂对于任意的，满足 \(k \in\) \([1,2], x \in[0, k]\) 的 \(k, x\)＂。这样看的话，\(k, x\) 彻底同级，并且 \(k, x\) 可以任意的遍历平面区域 \(k \in[1,2], x \in[0, k]\) 。既然如此，如果我们要先冻结 \(\forall x\) ，再去关注 \(k\) 的话，就得反过来让 \(k\) 所在的胶卷被 \(x\) 控制。做个思想试验来稍加讨论一下：如果 \(x\) 很小，比如 \(x=0.5\) ，那么显然 \(k\) 可以遍历 \([1,2]\) ；如果 \(x\) 较大，比如 \(x=1.5\) ，那么 \(k\) 不仅要在［ 1,2\(]\) 里，还要保证 \(k \geq x\) 。而且，显然 \(x\) 自身是可以遍历 \([0,2]\) 的。因此，如果 \(x \in[0,1]\) ，则 \(k\) 遍历 \([1,2]\) ；如果 \(x \in(1,2]\) ，则 \(k\) 遍历 \([x, 2]\) 。明确了这一点之后，我们就可以开始遍历了。

\[
f(x)=(x-1) e^{2 x}+k\left(x-\frac{4}{3} x^{3}\right)
\]

因此冻结 \(x\) 后，它关于 \(k\) 是线性函数，要么增要么威。因此需要再次讨论斜率 \(x-\frac{4}{3} x^{3}\) 。不难发现当 \(x \in\left[0, \frac{\sqrt{3}}{2}\right]\)时，\(f\) 值关于 \(k\) 递增（或常函数），于是这时 \(k=2\) 时 \(f\) 值最大，为

\[
(x+1) e^{2} x+2\left(x-\frac{4}{3} x^{3}\right)
\]

当 \(x \in\left(\frac{\sqrt{3}}{2}, 2\right]\) 时，则相反，\(k\) 越小则 \(f\) 越大，于是 \(x \in\left(\frac{\sqrt{3}}{2}, 1\right]\) 时应当让 \(k=1\) ，也就是

\[
(x-1) e^{2} x+x-\frac{4}{3} x^{3}
\]

而当 \(x \in(1,2]\) 时，\(k\) 最小也是 \(x\) ，这时候 \(f\) 值为

\[
(x-1) e^{2} x+x\left(x-\frac{4}{3} x^{3}\right)
\]

综合一下，冻结 \(x\) 时，\(f\) 关于 \(k\) 的最大值（记为 \(\varphi(x)\) ）的表达式在上面已经给出来了，即

\[
\varphi(x)=\left\{\begin{array}{l}
(x-1) e^{2 x}+2\left(x-\frac{4}{3} x^{3}\right), x \in\left[0, \frac{\sqrt{3}}{2}\right] \\
(x-1) e^{2 x}+x-\frac{4}{3} x^{3}, x \in\left(\frac{\sqrt{3}}{2}, 1\right] \\
(x-1) e^{2 x}+x\left(x-\frac{4}{3} x^{3}, x \in(1,2]\right.
\end{array}\right.
\]

现在，我们需要 \(\forall x \in[0,2]\) 都有 \(\varphi(x) \leq c\) ，也就是要求 \(\varphi(x)\) 在 \([0,2]\) 上的最大值。这是一个静态函数的分析了。\\
注意 \(x \in[0,1]\) 上的两段都是取了具体的 \(k\) 的，因此前面尝试时做的计算仍然适用。分两段计算 \(\varphi^{\prime}(x)\) ，可知第一段上（即 \(k=2) \varphi^{\prime}(x)\) 有零点 \(\frac{1}{2}\) ，还有一个零点是

\[
\frac{e^{2 x}}{2 x+1}=2
\]

之解。由于 \(h(x)\) 增，而且

\[
h\left(\frac{1}{2}\right)=\frac{e}{2}-2<0,
\]

因此经过简单的分析可知，当 \(x \in\left[0, \frac{\sqrt{3}}{2}\right]\) 时，\(\rho(x)\) 最大值要么在 \(\frac{1}{2}\) 取，要么在 \(\frac{\sqrt{3}}{2}\) 取，分别是 \(\frac{2}{3}-\frac{e}{2}\) 与 \(\left(\frac{\sqrt{3}}{2}-1\right) e^{\sqrt{3}}\) 。第二段相当于取了 \(f(x)\) 中 \(k=1\) ，因此 \(f(x)\) 单调性只与 \(h(x)=\frac{e}{2 x} 2 x+1-1\) 的零点分布有关。而显然此时有

\[
h(0)=0,
\]

因此

\[
h(x)>0, x>\frac{\sqrt{3}}{2}
\]

因此 \(f(x)\) 在 \(\left(\frac{\sqrt{3}}{2}, 1\right]\) 上增，这一段的最大值在 \(x=1\) 上取，为 \(-\frac{1}{3}\) 。这个值肯定超过了 \(\left(\frac{\sqrt{3}}{2}-1\right) e^{\sqrt{3}}\)（因为这个值也适用于第二段函数的 \(x=\frac{\sqrt{3}}{2}\) 函数值），也很容易比较出

\[
-\frac{1}{3}>\frac{2}{3}-\frac{e}{2},
\]

因此目前的最佳候选人就是 \(-\frac{1}{3}\) 。\\
再看 \(x \in(1,2]\) ，这一次得重新计算了：

\[
\varphi^{\prime}(x)=(2 x-1) e^{2 x}+2 x-\frac{16}{3} x^{3}
\]

函数不好看，但是本身我们只需考察 \(x \in(1,2]\) 的小区间，因此可能不需要做特别精细的分析就能搞定它。我们可以先代几个特殊值来直观感受一下 \(\varphi^{\prime}(x)\) 的情况。

\[
\varphi^{\prime}(1)=e^{2}-\frac{10}{3}>0,
\]

开头是正的

\[
\varphi^{\prime}\left(\frac{3}{2}\right)=2 e^{3}-15>0,
\]

而且这几个不等式似乎不是很紧。更何况这个 \(\varphi^{\prime}(x)\) 看上去就像是增的一般，因此可以考虑直接求导看看。

\[
\varphi^{\prime \prime}(x)=4 x e^{2 x}+2-16 x^{2}
\]

这玩意看上去更像增的了：

\[
\varphi^{\prime} \prime(x)=4 x\left(e^{2 x}-4 x\right)+2
\]

不难证明 \(e^{2} x-4 x>0(x>1)\) ，因此 \(\varphi^{\prime}(x)>0\) ，然后得到 \(\varphi^{\prime}(x)\) 增，于是 \(\varphi^{\prime}(x)>0, \varphi(x)\) 增，那最大的就是

\[
\varphi(2)=e^{4}-\frac{52}{3} .
\]

这就是最终答案。\\
从这道题目可以看出，交换一下平级的约束变元的顺序，计算求解会带来很不同的改变。因此，在遇到两个量词相同的约束变元并排摆放的时候，换序法是充要替换命题的一种有效方式。大家可以多加体会。

\section*{总结}
本节讲义主要对自由度的一个易错点——对象问题进行了分析，并由此引申到自由度作用对象的个体与整体性的讨论。运动的个体所包装成的＂胶卷＂就可以看成静态的对象，由此引申出的动静互化的思维方式，以及

它的一个特例—函数思想。\\
这种理解方式还与一类题型——封装命题不谋而合。封装命题就是利用量把词（ \(~\) 与 ヨ）命题的矛头指向其他真正的自由度，被封装的变元就成为了约束变元。正是这种特性，决定了封装命题解决的统一宗旨—消去约束变元，变成普通命题（进而用一般的，普遍的动静视角，要素网，元问题的思维范式处理）。由于解决时不可避免地要处理约束变元，变量势必比较多，我们就可以利用遍历与冻结的理念，辅以设问的思考方式，去变被动为主动，化特殊为一般，主动尝试去充要替换。这个过程其实和我们平时讲的思维，题目没有本质区别，思维方式除了上述注意点以外也都大差不差，无非就是把几种＂干货篇＂up 主视频列表中的思维拿来用了而已。

本节的内容由于在一定程度上涉及到了真正的逻辑学内容，所以可能有一定的抽象性。希望大家之后多花时间理解，掌握，也可以去评论区的置顶获取本文档。思维范式的形成绝非一朝一夕的功夫，包括前五节中讲到的思维也是如此。如第五节所说一般，花大量时间去实践，试错，反思，调整，才能有更多的收获，才能取得理想的进步。

\section*{第2章 高中解题思想与杂谈}
\section*{2.1 辅助思想}
2.2 具体方法

\section*{2.3 阶段性总结}
\section*{2.4 好题分享}
2．4．1 一类指数 or 对数不等式的通用思想\\
2．4．2\\
2．4．3 靠清晰的逻辑征服一切！一道不对称圆锥曲线大题的超详细思路分享\\
2．4．4 只要逻辑清晰即使是怪题你也喜欢吗？一道三角填空的超详细思路分享\\
2.5 漫游难题\\
2.6 数学联赛\\
2.7 吹水／杂谈

\section*{第3章 计算基本法—杜杆}
说到数学计算，大家常见的问题就是：式子虽然列对了，但经常算的慢，总算错；15 分大题，明明会做，因为计算失误，只拿 \(1, ~ 2\) 分。为了解决这类问题，up 主「阿不高中数学 \({ }^{1}\) 」将 10 年教学经验总结成了本章的这些内容，也是他高三模拟稳上 149 的秘密武器之一，让他成功走进北大数院。它不光适合阿不和笔者本人，也适合身为读者的你，无论你处于哪个年纪，分数段，学完都能挽回因计算失误丢掉的分！

本章内容中，我们会详细拆解计算的宏观底层思维，让你看懂别人的答案＂为什么这么算＂！教你灵活应用 ＂代数结构＂，算的又快又准！不仅如此，你还将学到：阿不的检查方法，未来会与阳哥的思维方式结合，能快速精准查错，建立计算的信心。具体怎么做呢？简单来说分为三个步骤：\\
（1）看懂＂为什么这么算＂。 很多读者遇到的计算的最大的问题就是各种看不懂，答案的计算太巧妙了，怎么想到的？这个步骤是怎么回事，是不是省略了什么？会有这些问题，都是因为你没有掌握计算的＂根＂——＂代数结构＂；我们会详细为你拆解齐次，对称，主元，换元，共轭，自由度，因式分解等＂代数结构＂，掌握了它们，你就可以明确知道：每一步的计算应该如何进行，让计有规律可循！\\
（2）达到超高的精准度，极快的计算速度。＂代数结构＂是做到这两点的秘密武器，它既可以规避无效计算，提升速度，又可以提供结构的信息帮你检查，保证正确率。为了帮你更好的掌握它，我们引入独创的＂杜杆模式＂，让大家在后续章节中一眼看到＂代数结构＂四两拨千斤的效果。我们将和你一起见证题目中巨复杂，巨困难的计算是如何被＂代数结构＂这根小小杜杆轻松解决的。\\
（3）有效检查，建立信心。 像我们开头提到的一样，很多同学在面对复杂计算时，心里没底，害怕算错，经常中途检查好几遍。为了解决这个问题，我们还会在后续章节教你检查基本法＂，有了它，你将能从＂正向角度，反向角度，结构角度＂等，来评估计算结果，让计算中的错误无处藏身。

经过 3 个训练步骤，相信你一定能彻底掌握计算的底层动机，找到速算捷径，敢算又能算对！我们将会从基础知识讲起，再在后续章节逐渐过渡到高考中难题，难题计算，所以本章内容即使你是零基础也不用怕。

\section*{3.1 主元思想}
主元思想是非常重要的恒等变形思想之一，它在学习函数，圆曲，导数时都会频繁使用。因此，我们先从这个最基本的思想出发，窥探代数结构的杠杆作用。为了讲清楚主元法，也为了提升大家的运算速度，我们先从更基本的乘法公式出来，从新的角度来理解这个公式。

\footnotetext{1高三所有模拟数学均未低于 149 ，后在北大普通学院＂复读＂一年，即狂刷数分习题，成功转专业也去了北大数院，现为杜克大学数学博士 （直博）
}\section*{3．1．1 从乘法公式开始}
\section*{3．1．1．1 俩个括号相乘}
我们在初中就已经学过了乘法公式，但是在实际运算中，很多同学运用得很不熟练。这其中的关键，就是因为在使用的时候仍旧很机械。我们知道，在化简整式的时候，需要分两步进行：\\
（1）展开括号：根据乘法公式，展开括号；\\
（2）合并同类项。\\
这是我们一开始学习的时候，都会使用的方法。例如，在计算 \((x+2)(x+3)\) 时，初学者往往会进行如下的化简：

\[
\begin{aligned}
(x+2)(x+3) & =x \cdot x+x \cdot 3+2 \cdot x+2 \cdot 3 \\
& =x^{2}+3 x+2 x+6 \\
& =x^{2}+5 x+6 .
\end{aligned}
\]

然而，对乘法公式更熟练的同学，会将（1），（2）两个步骤一起进行。怎么做？我们发现：这是两个关于 \(x\) 的一次式相乘，所以结果应该是一个二次式。那么我们只需要分别去确定二次项，一次项和常数项（零次式）分别是多少即可。\\
（1）二次项：两个括号各贡献一个 \(x\) ，得 \(x \cdot x=x^{2}\) ；\\
（2）一次项：结果中的一次式有两个来源：第一个括号贡献一个 \(x\) ，第二个括号贡献一个常数（零次项）+3 ，得到 \(3 x\) ；第一个括号贡献 +2 ，第二个括号贡献一个 \(x\) ，得到 \(2 x\) 。这样，一次项为 \(2 x+3 x=5 x\) ；\\
（3）常数项：两个括号各贡献一个常数，得 \(2 \cdot 3=6\) ．\\
这样，总的结果为 \(x^{2}+5 x+6\) 。我们不需要再严格按括号展开，而是根据＂对结果的预期＂和＂每一项的来源＂来完成展开。这是我们在使用乘法公式时最重要的基本方法，希望大家可以通过反复练习，熟练掌握。

我们再举一个更复杂一点的例子。展开 \((x-5)\left(2 x^{2}+x-4\right)\) 。\\
这是一个一次式，一个二次式的乘积，所以我们预期结果是一个三次式。这样，我们需要分别去确定三次项，二次项，一次项和常数项。\\
（1）三次项：\(x \cdot 2 x^{2}=2 x^{3}\) ；\\
（2）二次项：两个来源，第一个括号如果贡献 \(x\) ，那么第二个括号需要贡献 \(x\) ，得到 \(x^{2}\) ；第一个括号如果贡献 -5 ，那么第二个括号贡献 \(2 x^{2}\) ，得到 \(-10 x^{2} . x^{2}-10 x^{2}=-9 x^{2}\) 就是二次项；\\
（3）一次项：两个来源，第一个括号如果贡献 \(x\) ，那么第二个括号需要贡献 -4 ，得到 -4 x ；第一个括号如果贡献 -5 ，那么第二个括号贡献 \(+x\) ，得到 \(-5 x .-4 x-5 x=-9 x\) 就是一次项；\\
（4）常数项：\(-5 \cdot(-4)=20\) ．将上述结果合并，得到 \(2 x^{3}-9 x^{2}-9 x+20\) ．从而 \((x-5)\left(2 x^{2}+x-4\right)=2 x^{3}-9 x^{2}-9 x+20\) ．如何检查我们的展开结果呢？我们发现，当 \(x=5\) 时，原式为 0 。那么我们将 \(x=5\) 代人到展开式中，应该也是 0 ：

\[
2 \cdot(5)^{3}-9 \cdot 5^{2}-9 \cdot 5+20=250-225-45+20=0
\]

所以，我们的结果应该是正确的，至少＂不是错的＂。\\
总结一下，我们在展开这种只含有一个字母的括号时，要抓住两点：次数和来源。\\
1．次数：先判断结果的总次数，进而确定都有哪些项；例如，结果是四次式，那么一定由 \(x^{4}, x^{3}, x^{2}, x\) ，常数项这五部分组成；注意，有时的展开结果可能不含有某些低次项，我们可以将它的系数理解为 0 。\\
2．来源：对于每一项，确定是可以如何得到的，再去计算各个来源的系数；\\
3．将每一项结果加起来，就是最终结果。

只有两个括号相乘，且括号内项数较少（小于等于 3）时，这种方法相对更方便。如果括号内项数较多，这么算容易出错，我们按原来的方法直接展开也是可以的。

\section*{3．1．1．2 多个括号}
当我们遇到多个括号的情况时，可以两两相乘，最终化为两个括号。例如：在计算 \((x+1)(x+2)(x+3)\) 时，我们既可以像上面一样直接展开，也可以先乘前两个。\\
（1）直接展开\\
如果直接展开，我们预期结果是三次式，所以有 \(x^{3}, x^{2}, x\) 和常数项。\\
（1）\(x^{3}\) ：每个括号贡献一个 \(x, x \cdot x \cdot x=x^{3}\) ；\\
（2）\(x^{2}\) ：两个括号贡献 \(x\) ，一个括号贡献常数。这样就有三种可能，第一个括号贡献常数，就是 \(1 \cdot x \cdot x=x^{2}\) ；第二个贡献常数，就是 \(x \cdot 2 \cdot x=2 x^{2}\) ；第三个贡献常数，就是 \(x \cdot x \cdot 3=3 x^{2}\) 。和为 \(x^{2}+2 x^{2}+3 x^{2}=6 x^{2}\) ；\\
（3）\(x\) ：一个括号贡献 \(x\) ，两个括号贡献常数。这样也有三种可能，第一个括号贡献 \(x\) ，就是 \(x \cdot 2 \cdot 3=6 x\) ；第二个括号贡献 \(x\) ，就是 \(1 \cdot x \cdot 3=3 x\) ；第三个括号贡献 \(x\) ，就是 \(1 \cdot 2 \cdot x=2 x\) 。和为 \(6 x+3 x+2 x=11 x\) ．\\
（4）常数项：每个括号都贡献常数项，所以是 \(1 \cdot 2 \cdot 3=6\) ．最终结果为：\((x+1)(x+2)(x+3)=x^{3}+6 x^{2}+11 x+6\) ．\\
（2）先乘前两个\\
如果直接展开，很多同学可能会记不住项都从哪里来。这样的话，我们两个两个括号展开就行了。先去计算 \((x+1)(x+2)\) ，方法和上面一样，我就不重复了，结果是 \(x^{2}+3 x+2\) 。这样，再去计算 \((x+3)\left(x^{2}+3 x+2\right)\) 即可。\\
（1）\(x^{3}\) ：第一个括号贡献 \(x\) ，则第二个括号贡献 \(x^{2}\) ，结果为 \(x^{3}\) ；\\
（2）\(x^{2}\) ：第一个括号贡献 \(x\) ，则第二个括号贡献 \(3 x\) ，结果为 \(3 x^{2}\) ；第一个括号贡献 +3 ，则第二个括号贡献 \(x^{2}\) ，结果为 \(3 x^{2}\) ，和为 \(3 x^{2}+3 x^{2}=6 x^{2}\) ；\\
（3）\(x\) ：第一个括号贡献 \(x\) ，则第二个括号 +2 ，结果为 \(2 x\) ；第一个括号贡献 +3 ，则第二个括号贡献 \(3 x\) ，结果为 \(9 x\) ，和为 \(2 x+9 x=11 x\) ；\\
（4）常数项： \(3 \cdot 2=6\) ．\\
最终结果为：\((x+1)(x+2)(x+3)=x^{3}+6 x^{2}+11 x+6\) ．\\
因此，无论是哪种方法，都可以得到正确答案。实际应用中，可以灵活选择。推荐初学者先使用第二种方法。

\section*{3．1．2 扩展到主元思想}
在刚才对于乘法公式的使用中，我们只有一个字母 \(x\) 。但是实际应用中，我们经常会有多个字母出现在代数式中。这时，我们应该如何展开呢？

假设我们需要展开 \((a+b+2)(a+3 b-1)\) ．如果直接使用乘法公式展开，我们会得到：

\[
\begin{aligned}
(a+b+2)(a+3 b-1) & =a^{2}+3 a b-a+a b+3 b^{2}-b+2 a+6 b-2 \\
& =a^{2}+3 b^{2}+4 a b+a+5 b-2
\end{aligned}
\]

我们可以将结果按 \(a\) 的降幂排列，得到 \(a^{2}+3 b^{2}+4 a b+a+5 b-2=a^{2}+4 a b+a+3 b^{2}+5 b-2\) ．\\
我们发现，如果想要使用刚才的方法，有一定困难。因为这里有两个字母，而且次数比较凌乱，没有一个明显的规律。这样一来，我们想：如果可以把它看成一个字母的表达式，是不是就和刚才一样解决问题了呢？

这就是我们要讲的主元思想。虽然表达式中有两个字母，但是我们可以把其中一个看成主元，也就是真正

的变量，而另一个看成参数。主元充当了之前 \(x\) 的位置，是我们恒等变形的中心；而参数，虽然也是用字母代替的，我们可能不知道它的值，但是它的地位和一般的常数是一样的。它是一个＂固定＂的数，把它理解成常数就可以。这样，如果我们把 \(a\) 当作主元，把 \(b\) 当作参数，我们的式子应该写为：

\[
(a+b+2)(a+3 b-1)=[a+(b+2)][a+(3 b-1)] .
\]

这里我们给 \(b+2,3 b-1\) 加上了括号，就是为了给大家强调：\(b\) 和其他常数地位一样，都是常数，整体看 \(b+2,3 b-1\)都是常数。这样的话，我们就得到了一个关于 \(a\) 的二次式，它和 \((a+2)(a+3)\) 没有区别（当 \(b=1\) 时，就是这个式子）。因此，我们可以按刚才的方法进行展开：\\
（1）二次式：\(a^{2}\) ；\\
（2）一次式：第一个括号贡献 \(a\) ，那么第二个括号贡献 \(3 b-1\) ，得 \((3 b-1) a\) ；第一个括号贡献 \((b+2)\) ，那么第二个括号贡献 \(a\) ，得 \((b+2) a\) 。这样，一次项为 \((3 b-1) a+(b+2) a=(4 b+1) a\) 。\\
（3）常数项：\((b+2)(3 b-1)=3 b^{2}+5 b-2\) ．所以，展开式可以整理为：\((a+b+2)(a+3 b-1)=a^{2}+(4 b+1) a+3 b^{2}+5 b-2\) ．和原来得结果是相同的。

主元思想最大的作用在于帮助大家理清思路。当主元确定后，我们只需要把它看成一个只有一个未知数（一元）的展开式去展开即可。这里，我们也用到了一个至关重要的数学思想：化归，将不熟悉的问题转化为熟悉的问题。我们将多个未知数（多元）的问题，通过主元的思想，转化为了一元问题，所有关于一元问题的思路和工具就都可以派上用场了。

\[
\text { 多元 } \xlongequal{\text { 主元思想 }} \text { 一元 }
\]

\section*{3．1．3 XXXX}
\section*{3.2 对称式}
\section*{3．2．1 什么是对称式？}
\section*{3．2．1．1 基本概念}
我们在初中数学的学习中，就已经接触过对称式了。例如，对于二次方程 \(a x^{2}+b x+c=0(a \neq 0)\) ，若它的两个根为 \(x, y\) ，那么由根系关系（韦达定理）；

\[
\left\{\begin{array}{l}
x+y=-\frac{b}{a} \\
x y=\frac{c}{a}
\end{array}\right.
\]

我们发现：如果将 \(x, y\) 看成两个变量（字母），那么 \(x+y\) 和 \(x y\) 都是二元对称式：如果交换 \(x, y\) ，我们的代数式不会发生改变。

例如，若交换 \(x, y\) ，那么 \(x+y\) 将会变为 \(y+x=x+y\) ；而 \(x y\) 则会转变为 \(y x\) ，它也等于 \(x y\) 。所以，它们都是关于 \(x, y\) 的二元对称式。

我们再举一些更复杂的例子，咱们来验证它们也都是对称式。\\
（1）\(x^{3} y+x y^{3}\) 是对称式。如果我们将 \(x, y\) 交换，那么这两项分别会变为：

\[
x^{3} y \rightarrow y^{3} x, x y^{3} \rightarrow y x^{3}
\]

这样的话，\(x^{3} y+x y^{3}\) 变为 \(y^{3} x+y x^{3}=x^{3} y+x y^{3}\) ，最终得到了一样的式子！所以，它也是一个对称式。\\
（2）\(x^{2} y+2 x y^{2}\) 不是对称式。如果我们将 \(x, y\) 交换，这两项分别会变为：

\[
x^{2} y \rightarrow y^{2} x, 2 x y^{2} \rightarrow 2 y x^{2}
\]

这样的话，\(x^{2} y+2 x y^{2}\) 变为 \(y^{2} x+2 y x^{2}=2 x^{2} y+x y^{2}\) ，它和原来的 \(x^{2} y+2 x y^{2}\) 不一样。所以，它不是对称式。这里可能有的同学会有疑问：那如果 \(x=y\) ，那么这两个式子不就一样了吗？是的，但是我们要求的＂相等＂，是＂恒等＂：无论对于什么样的 \(x, y\) ，这两个式子都要相等。这里，如果 \(x=1, y=2\) ，那么

\[
x^{2} y+2 x y^{2}=10,2 x^{2} y+x y^{2}=8
\]

二者显然不等。所以，它们不恒等。\\
我们通过上面的例子发现，如果想要一个关于 \(x, y\) 的式子是对称的，你要保证：\(x\) 有什么，\(y\) 就有什么；\(y\)有什么，\(x\) 也有什么。比如，在 \(x^{3} y+x y^{3}\) 中，第一项是三个 \(x\) 配一个 \(y\) ，那么相应地，我也一定要有三个 \(y\) 配一个 \(x\) ，这样才对称；而在 \(x^{2} y+2 x y^{2}\) 中，虽然两个 \(x\) 配一个 \(y\) 和两个 \(y\) 配一个 \(x\) 都有了，但是系数不一样，那么就不是对称的。显然，这个式子更偏心 \(y\) ：有 \(y^{2}\) 的项，系数是 2 ，但是 \(x^{2}\) 的项系数仅为 1 。所以就不是对称的了。

因此，如果一个式子关于 \(x, y\) 对称，我们也通常说 \(x, y\) 地位相等：整个式子对谁都不偏心，一碗水端平，\(x\)有啥 y 都有，反过来一样。这就是对称的本质含义。

\section*{3．2．1．2 对称式的性质}
（1）封闭性\\
对称式有什么性质呢？首先，它对于加，减，乘，除都是封闭的。假如 \(f, g\) 都是对称的，那么：\(f+g, f-g, f g, \frac{f}{g}\)都是对称的。这是因为：当我们交换 \(x, y\) 的时候，\(f, g\) 本身都不发生改变。那么，关于它们四则运算的结果也就不会改变了。

特别地，常数也是对称的。因为它和 \(x, y\) 无关，所以 \(x, y\) 怎么互换都不会影响它，自然就是对称的。\\
这里补充说明一点：我们这里的例子，多为整式。但是分式，根式也可以是对称的。例如：

\[
\frac{1}{x}+\frac{1}{y}, \sqrt{x^{2}+y^{2}}
\]

都是对称的。读者可以自己验证。\\
（2）用 \(x+y, x y\) 表示对称式\\
任意关于 \(x, y\) 的二元对称多项式，都可以用 \(x+y, x y\) 来表示。例如，

\[
\begin{gathered}
x^{2}+y^{2}=(x+y)^{2}-2 x y, \\
x^{2} y+x y^{2}=x y(x+y) .
\end{gathered}
\]

如果记 \(m=x+y, n=x y\) ，那么有：

\[
x^{2}+y^{2}=m^{2}-2 n, x^{2} y+x y^{2}=m n .
\]

它们都化成了关于 \(m, n\) 的多项式。即便是复杂的式子，只要对称，也一定可以用 \(m, n\) 表示。例如：

\[
\begin{aligned}
x^{4}+y^{4} & =\left(x^{2}+y^{2}\right)^{2}-2 x^{2} y^{2} \\
& =\left(m^{2}-2 n\right)^{2}-2 n^{2} \\
& =m^{4}-4 m^{2} n+2 n^{2} .
\end{aligned}
\]

所以，我们可以用 \(m, n\) 表示所有的二元对称式。为了方便称呼，我们在这里把 \(m=x+y, n=x y\) 称为对称换元。\\
虽然我们不给出严谨的证明，但是我们利用递推的思想来验证：对于任意的正整数 \(k, x^{k}+y^{k}\) 是可以被 \(m, n\)表示的。

显然，当 \(k=1,2\) 时，我们有：

\[
x+y=m, x^{2}+y^{2}=m^{2}-2 n .
\]

那么，我们如何表示 \(x^{3}+y^{3}\) 呢？我们注意：如果去计算 \((x+y)\left(x^{2}+y^{2}\right)\) ，那么最高次项就会出现 \(x^{3}\) ，\(y^{3}\) 项，或许就可以和 \(x^{3}+y^{3}\) 建立联系。事实上，

\[
(x+y)\left(x^{2}+y^{2}\right)=x^{3}+y^{3}+x^{2} y+x y^{2}
\]

因此，

\[
\begin{aligned}
x^{3}+y^{3} & =(x+y)\left(x^{2}+y^{2}\right)-x^{2} y-x y^{2} \\
& =(x+y)\left(x^{2}+y^{2}\right)-x y(x+y) \\
& =m \cdot\left(m^{2}-2 n\right)-m n \\
& =m^{3}-3 m n .
\end{aligned}
\]

这样，我们就写出了 \(x^{3}+y^{3}\) 关于 \(m, n\) 的表达式了。那么，如何再去写 \(x^{4}+y^{4}\) 的呢？我们可以利用同样的方法：

\[
(x+y)\left(x^{3}+y^{3}\right)=x^{4}+y^{4}+x^{3} y+x y^{3} .
\]

因此，

\[
\begin{aligned}
x^{4}+y^{4} & =(x+y)\left(x^{3}+y^{3}\right)-x^{3} y-x y^{3} \\
& =(x+y)\left(x^{3}+y^{3}\right)-x y\left(x^{2}+y^{2}\right) \\
& =m\left(m^{3}-3 m n\right)-n\left(m^{2}-2 n\right) \\
& =m^{4}-3 m^{2} n-m^{2} n+2 n^{2} \\
& =m^{4}-4 m^{2} n+2 n^{2} .
\end{aligned}
\]

所以，\(x^{4}+y^{4}\) 也可以写出来。类似地，\(x^{5}+y^{5}, x^{6}+y^{6} \ldots .\). 我们可以一步一步写下去了。\\
如果更加仔细地观察刚才的计算，我们发现，在计算的时候有这么一个规律：

\[
\begin{gathered}
x^{3}+y^{3}=m \cdot\left(x^{2}+y^{2}\right)-n \cdot(x+y), \\
x^{4}+y^{4}=m \cdot\left(x^{3}+y^{3}\right)-n \cdot\left(x^{2}+y^{2}\right)
\end{gathered}
\]

依此类推，对任意的正整数 \(k\) ，都有：

\[
x^{k}+y^{k}=m \cdot\left(x^{k-1}+y^{k-1}\right)-n \cdot\left(x^{k-2}+y^{k-2}\right) .
\]

这样一来，如果我们知道 \(x^{k}-1+y^{k-1}\) 和 \(x^{k}-2+y^{k-2}\) 的表达式，我们就可以求出 \(x^{k}+y^{k}\) 的表达式。例如：知道 \(x^{3}+y^{3}, x^{4}+y^{4}\) 可以写 \(x^{5}+y^{5}\) ；知道 \(x^{4}+y^{4}, x^{5}+y^{5}\) 可以写 \(x^{6}+y^{6} \ldots \ldots\) 。这样一来，所有的 \(x^{k}+y^{k}\) 都可以用 \(m, n\) 表示了。

\section*{3．2．2 对称式与主元思想}
我们在第一个杠杆中讲了主元思想，它帮助我们把多元问题化为一元问题。那么，主元思想是如何在 \(x, y\) 的二元对称式中体现的呢？

最直接的，我们有 \(x, y\) 两个变元，可以直接使用其中一个作为主元，另一个作为参数。但是，这样做的问题是：对称性被破坏了。选取主元后，主元和参数的地位是不等价的，前者是真正的变量，后者是常数。因此，这么做效果不好。例如，我们考虑 \(x^{3} y+x y^{3}+x+y\) ，如果我们选 \(x\) 为主元，那么

\[
f(x)=y \cdot x^{3}+\left(y^{3}+1\right) x+y .
\]

除了对称性被破坏，这还是一个关于 \(x\) 的三次函数。我们对于三次函数的内容了解很少，这么选主元，并没有化简问题。

为了保护对称性，我们可以先进行对称换元，再选主元。

\[
\begin{aligned}
& x^{3} y+x y^{3}+x+y \\
& =x y\left(x^{2}+y^{2}\right)+x+y \\
& =n \cdot\left(m^{2}-2 n\right)+m \\
& =m^{2} n-2 n^{2}+m .
\end{aligned}
\]

这样一来，我们虽然无法直接看到 \(x, y\) ，但是 \(x, y\) 的对称性被 \(m, n\)＂编码＂进去了：\(m, n\) 是关于 \(x, y\) 的对称式，所以最后的表达式也关于 \(x, y\) 对称。除此之外，它还是一个二次式。无论是以 \(m\) 为主元，还是以 \(n\) 为主元，都可以得到一个二次函数：\\
（1）以 \(m\) 为主元：\(f(m)=n m^{2}+m-2 n^{2}\) ；\\
（2）以 \(n\) 为主元：\(g(n)=-2 n^{2}+m^{2} n+m\) ．\\
因此，使用对称换元，不光保护了对称性，还降低了次数，一举两得。这是使用对称换元的正确方法。\\
总结一下，当我们在处理对称式时，绝大多数情况都需要先进行对称换元，再选主元。选好后，再进行处理

\[
x, y \text { 对称式 } \stackrel{\text { 对称换元 }}{\Longrightarrow} \text { 关于 } m, n \stackrel{\text { 选主元 }}{\Longrightarrow} \text { 一元 }
\]

这样，我们又利用了主元思想将多元问题转化为了一元问题。这些代数杜杆，不光本身具有强大的威力，组合起来效果更是惊人。

选好主元后如何处理呢？我们会在后续的杜杆中再深入探讨。

\section*{3．2．3 XXXX}
\section*{3.3 齐次式}
\section*{3．3．1 什么是齐次式？}
我们展开 \((a-b)^{2}\) 时，得到展开式：

\[
(a-b)^{2}=a^{2}-2 a b+b^{2}
\]

我们发现，展开式中的所有项都是二次项：\(a^{2}, ~-2 a b, b^{2}\) 。像这样，一个多项式中的所有项都是 \(m\) 次的，我们就称它为 \(m\) 次齐次式。例如，\(a^{2}-2 a b+b^{2}\) 就是一个二次齐次式。\\
\(m\) 次齐次式最重要的性质是：当所有变量变为原来的 \(k\) 倍时，整个代数式的变为原来的 \(k^{m}\) 倍。例如，对于 \((a+b)^{2}\) 而言，若 \(a\) 变为 \(k a, b\) 变为 \(k b\) ，那么代数式变为

\[
(k a+k b)^{2}=k^{2}(a+b)^{2} .
\]

代数式是原来的 \(k^{2}\) 倍。再如，对于 \(2 a^{3}+4 a^{2} b-a b^{2}+b^{3}\) ，它是一个三次齐次式。若 \(a\) 变为 \(k a, b\) 变为 \(k b\) ，那么

代数式变为

\[
\begin{aligned}
& 2(a k)^{3}+4(a k)^{2}(b k)-(a k)(b k)^{2}+(b k)^{3} \\
= & 2 k^{3} a^{3}+4 k^{3} a^{2} b-k^{3} a b^{2}+3 k^{3} b^{3} .
\end{aligned}
\]

提公因式 \(k^{3}\) ，得到

\[
2 k^{3} a^{3}+4 k^{3} a^{2} b-k^{3} a b^{2}+3 k^{3} b^{3}=k^{3}\left(2 a^{3}+4 a^{2} b-a b^{2}+b^{3}\right) .
\]

它是原来表达式的 \(k^{3}\) 倍。\\
齐次式的这个性质决定了：当一个分式的分子，分母是同样次数的齐次式时我们可以进行比例换元。我们举一个简单的例子来说明，考虑分式

\[
\frac{2 x^{2}-x y+3 y^{2}}{y^{2}}
\]

我们的分子，分母都是二次齐次式。令 \(t=\frac{x}{y}\) ，将分子，分母同时除以 \(y^{2}\) ，得到：

\[
\frac{2 x^{2}-x y+3 y^{2}}{y^{2}}=\frac{\frac{2 x^{2}}{y^{2}}-\frac{x y}{y^{2}}+\frac{3 y^{2}}{y^{2}}}{\frac{y^{2}}{y^{2}}}=\frac{2\left(\frac{x}{y}\right)^{2}-\frac{x}{y}+3}{1}=2 t^{2}-t+3 .
\]

它转化成了一个关于 \(t\) 的二次函数。一般而言，如果分子，分母都是 \(m\) 次齐次，那么分子，分母也都会化为关于 \(t\) 的 \(m\) 次函数。在这个例子中，情况比较特殊：它的分子是二次函数，分母退化成了常数我们再举一例：对于分式

\[
\frac{x^{2}+3 x y-y^{2}}{x^{2}+y^{2}}
\]

令 \(t=\frac{x}{y}\) ，分子分母同时除以 \(y^{2}\) ，可得：

\[
\frac{x^{2}+3 x y-y^{2}}{x^{2}+y^{2}}=\frac{\frac{x^{2}}{y^{2}}+\frac{3 x y}{y^{2}}-\frac{y^{2}}{y^{2}}}{\frac{x^{2}}{y^{2}}+\frac{y^{2}}{y^{2}}}=\frac{\left(\frac{x}{y}\right)^{2}+3 \cdot \frac{x}{y}-1}{\left(\frac{x}{y}\right)^{2}+1}=\frac{t^{2}+3 t-1}{t^{2}+1}
\]

在两个例子中，我们将一个含有两个变量的表达式化简成了一个变量的表达式，完成了消元的操作。消元的思想在求取值范围，解方程时都是非常常用的。

\section*{3．3．2 齐次式与其他杜杆的联系}
我们在前两节中，着重介绍了主元思想和对称式。它们在处理代数问题中发挥的核心作用，就是去＂消元＂ ：主元将多元问题转化为一元问题；对称换元可以保留对称结构的同时，进行消元。那么，齐次式是如何起到消元作用的呢？

我们在例题中看到，如果一个分式，分子和分母都是关于 \(x, y\) 的 \(k\) 次齐次式，那么，我们可以利用比例换元 \(t=\frac{x}{y}\) ，将这个分式化为关于 \(t\) 的一个分式。例如：

\[
\frac{x^{2}+3 x y-y^{2}}{x^{2}+y^{2}}=\frac{t^{2}+3 t-1}{t^{2}+1}
\]

左侧是的分子，分母都是二次齐次式，最终得到的结果是一个分子，分母都是关于 \(t\) 的二次函数的式子。我们发现：经过比例换元，本来含有两个变量 \(x, y\) 的表达式，变为了只含有一个变量 \(t\) 的表达式。这样，所有关于一元函数的工具就都可以派上用场了。

现在的问题是：为什么形如 \(\frac{x^{2}+3 x y-y^{2}}{x^{2}+y^{2}}\) 的式子，经过比例换元后，就能达到消元的效果呢？我们考察下面两个分式，来理解一下。\\
（1）\(\frac{x^{2}+y^{2}+1}{x y}\) 无法利用比例换元消元。如果我们令 \(t=\frac{x}{y}\) ，则

\[
\frac{x^{2}+y^{2}+1}{x y}=\frac{t^{2} y^{2}+y^{2}+1}{t^{2} y^{2}} .
\]

仍然需要两个变元 \(t, y\) 来表示。这是因为：虽然分子，分母都是二次式，但是分子 \(x^{2}+y^{2}+1\) 不是齐次的：虽然 \(x^{2}+y^{2}\) 是二次齐次的，但 1 是一个零次式（当 \(x, y\) 等比例改变时，它不改变，\(k^{0}=1\) ，所以是零次）。这样，当 \(x, y\) 等比例改变时，\(x^{2}+y^{2}+1\) 就无法等比例改变了，那么比例换元也就无法起到消元的作用了。\\
（2）\(\frac{x^{2}+y^{2}}{x+y}\) 也无法利用比例换元消元。如果我令 \(t=\frac{x}{y}\) ，则

\[
\frac{x^{2}+y^{2}}{x+y}=\frac{t^{2} y^{2}+y^{2}}{t y+y}=y \cdot \frac{t^{2}+1}{t+1}
\]

虽然我们得到了部分关于 \(t\) 的函数，但是由于有 \(y\) 的存在，我们并没有完成消元。这是因为：虽然 \(x^{2}+y^{2}, x+y\)都是齐次的，但是它们次数不一样。前者是二次的，后者是一次的。当 \(x, y\) 等比例扩大的时候，分子扩大 \(k^{2}\) 倍，而分母扩大 \(k\) 倍，所以即使 \(\frac{x}{y}\) 固定，整个表达式的值也无法确定。这样，消元就失败了。

我们通过这两个例子，意识到：想要成功使用比例消元，有一个不可或缺的条件：分子，分母是同次齐次式。此时，整个表达式和常数一样，是一个零次齐次式。因为：当 \(x, y\) 同时变为原来 \(k\) 倍时，我们有：

\[
\frac{(k x)^{2}+3 k x \cdot k y-(k y)^{2}}{(k x)^{2}+(k y)^{2}}=\frac{k^{2} \cdot\left(x^{2}+3 x y-y^{2}\right)}{k^{2} \cdot\left(x^{2}+y^{2}\right)}=\frac{x^{2}+3 x y-y^{2}}{x^{2}+y^{2}}
\]

整个表达式没有发生变化。因此，即使 \(x, y\) 发生变化，但是只要比例没有发生变化，这个表达式就不发生变化。这意味着：这个式子一定由比例 \(t=\frac{x}{y}\) 决定！所以，它可以写成关于 \(t\) 的函数。事实上，所有零次齐次式都可以写成关于 \(t=\frac{x}{y}\) 的函数。我们将在＂自由度＂这一节，重新探讨这个结论。

\section*{3．3．3 平面几何中的齐次结构}
齐次结构虽然是一个代数结构，但我们初中学习的很多几何定理也是具有齐次结构的。最简单的，勾股定理：\(c^{2}=a^{2}+b^{2}\) 。我们发现：\(a^{2}, b^{2}, c^{2}\) 都是关于边长的二次式。如果我们把它变形成 \(c^{2}-a^{2}-b^{2}=0\) ，那么，当 \(a, b, c\) 变为原来的 \(k\) 倍时，原式变为原来的 \(k^{2}\) 倍。

\[
(k c)^{2}-(k a)^{2}-(k b)^{2}=k^{2}\left(c^{2}-a^{2}-b^{2}\right)=0
\]

而本来这个式子就是 0 ，所以变为原来的 \(k\) 倍后依旧是 0 。这意味着什么呢？当我们把一个直角三角形的边长变为原来的 \(k\) 倍时，它依旧是直角三角形。这是自然的：因为新的三角形和原来的三角形相似，形状不变，还是直角三角形。

我们学习的锐角三角函数，也是具有齐次结构的。在直角三角形 ABC 中，\(\angle C=90^{\circ}\) ．根据锐角三角函数的定义：

\[
\sin A=\frac{B C}{A B}, \cos A=\frac{A C}{A B} .
\]

当我们将三角形各边变为原来的 \(k\) 倍，变成三角形 \(\Delta A^{\prime} B^{\prime} C^{\prime}\) 时，角度的正弦值，余弦值都不会发生改变：

\[
\begin{aligned}
& \sin A^{\prime}=\frac{B^{\prime} C^{\prime}}{A^{\prime} B^{\prime}}=\frac{k B C}{k A B}=\frac{B C}{A B}=\sin A \\
& \cos A^{\prime}=\frac{A^{\prime} C^{\prime}}{A^{\prime} B^{\prime}}=\frac{k A C}{k A B}=\frac{A C}{A B}=\cos A
\end{aligned}
\]

因此，新角的正弦值，余弦值也没有发生改变，角度不变。这和我们相似的性质是吻合的。\\
从这个例子，我们自然可以联想到：相似三角形，平行线分线段成比例等定理，也都是具有齐次结构的。例如，当 \(\triangle A B C\) 和 \(\triangle A^{\prime} B^{\prime} C^{\prime}\) 相似时，有：

\[
\frac{A B}{A^{\prime} B^{\prime}}=\frac{A C}{A^{\prime} C^{\prime}}=\frac{B C}{B^{\prime} C^{\prime}}=k
\]

当两个三角形都等比例变为原来的 \(m\) 倍时，相似比保持不变：

\[
\frac{m A B}{m A^{\prime} B^{\prime}}=\frac{m A C}{m A^{\prime} C^{\prime}}=\frac{m B C}{m B^{\prime} C^{\prime}}=k
\]

这是因为：如果我们令 \(x=A B, y=A^{\prime} B^{\prime}\) ，那么 \(\frac{x}{y}\) 本质上是一个关于 \(x, y\) 的零次式 \((1-1=0)\) 。这样一来，当他们变为原来的 \(m\) 倍的时候，相似比为 \(k \cdot m^{0}=k\) ，保持不变。

我们再来看第三个例子：圆幂定理。以相交弦定理为例，设 \(A B C D\) 为圆内接四边形，\(A C, B D\) 交于点 \(O\) ，则 \(A O \cdot C O=B O \cdot D O\) 。两侧也都是关于边长的二次式，当我们把整个图形变为原来的 \(k\) 倍时，等式仍然成立：

\[
k A O \cdot k C O=k B O \cdot k D O
\]

因此，它仍旧是圆内接四边形。

在这些例子中，勾股定理，锐角三角函数，相似三角形，相交弦定理，它们都是描述几何图形形状定理：直角三角形，角度，三角形形状相同，圆内接四边形等。所有描述形状的定理，关于边长的表达式一定是具有齐次结构的。因为当整个图形等比例放缩的时候，形状不会发生改变，定理中的代数表达式也维持成立。如果不齐次了，有的二次，有的三次，代数式就会出现不成立的情况，定理也就失效了。

\section*{3．3．4 XXXX}
\section*{3.4 共轭根式}
\section*{3．4．1 根的处理手法}
\section*{3．4．1．1 等式中的二次根式}
由于根式的特殊性，它无法直接进行加，减：

\[
\sqrt{a}+\sqrt{b} \neq \sqrt{a+b}
\]

所以它的处理较为繁琐。很多恒等变形的难点就在于对根式的处理。可以想象：最基本的处理根式的方法就是平方。那么，每平方一次，可以消掉几个根号呢？我们以 \(x=\sqrt{2}+\sqrt{3}\) 为例，我们能不能找到一个 \(x\) 满足的方程，它的所有项系数都是整数呢？

如果我们直接将它平方，可以得到：

\[
x^{2}=2+3+2 \sqrt{6}=5+2 \sqrt{6}
\]

下一步如何处理？如果我们依旧直接平方，得到

\[
x^{4}=25+24+20 \sqrt{6}=49+20 \sqrt{6} .
\]

这样的话，根号仍未消掉。所以，为了消掉根号，我们要把根式单独放在一边。换言之，\(x^{2}=5+2 \sqrt{6}\) 应写作：

\[
2 \sqrt{6}=x^{2}-5
\]

再平方，即可得

\[
24=x^{4}-10 x^{2}+25
\]

故 \(x^{4}-10 x^{2}+1=0\) ．因此，\(x=\sqrt{2}+\sqrt{3}\) 是这个方程的一个根．\\
但如果我们仍想用 \(x^{4}=49+20 \sqrt{6}\) 这个式子，能不能也得到这个方程呢？事实上，是可以的。我们只需将两个方程中的 \(\sqrt{6}\) 消掉即可。\\
（1） \(5+2 \sqrt{6}=x^{2}\) ，所以 \(\sqrt{6}=\frac{x^{2}-5}{2}\) ；\\
（2） \(49+20 \sqrt{6}=x^{4}\) ，所以 \(\sqrt{6}=\frac{x^{4}-49}{20}\) ．\\
故 \(\frac{x^{4}-49}{20}=\frac{x^{2}-5}{2}\) ．化简得 \(x^{4}-10 x^{2}+1=0\) ．也可以得到这个方程。在这个方法中我们将 \(\sqrt{6}\) 消掉，这也是处理根号的一种常见方法。

我们考虑更复杂得式子，设 \(x=\sqrt{2}+\sqrt{3}+\sqrt{5}\) ．我们能不能找到一个 \(x\) 满足的方程，系数都是整数呢？如果我们像刚才一样直接平方，我们发现：

\[
x^{2}=2+3+5+2 \sqrt{6}+2 \sqrt{10}+2 \sqrt{15}
\]

我们发现：本来 \(x\) 中有三个根号，现在仍有三个，而且变成了不同的根号。这样直接平方，并没有减少根号的次数。要知道，我们之所以平方，关键就是为了减少根号的个数。如果根号的个数没有减少，那么我们的平方就没有意义。

怎么办呢？我们应该像刚才处理 \(\sqrt{6}\) 一样，把一个根号放在一侧，先消掉它。例如，我们将原式变形为

\[
x-\sqrt{2}=\sqrt{3}+\sqrt{5} .
\]

两边平方，得到

\[
x^{2}-2 \sqrt{2} x+2=8+2 \sqrt{15}
\]

这样一来，我们就成功处理掉了一个根号，只剩下两个根号了。仿照这个方法，我们把 \(\sqrt{15}\) 单独放在一侧：

\[
x^{2}-2 \sqrt{2} x-6=2 \sqrt{15}
\]

两边平方，得到

\[
\begin{aligned}
& x^{4}+8 x^{2}+36-4 \sqrt{2} x^{3}-12 x^{2}+24 \sqrt{2} x=60 \\
& x^{4}-4 \sqrt{2} x^{3}-4 x^{2}+24 \sqrt{2} x-24=0 .
\end{aligned}
\]

这里我们发现：为什么平方完，还是两个根号呢？根号的个数并没有减少。事实上，个数看上去没有减少，但是我们只有一类根号了：所有的根号都是 \(\sqrt{2}\) 。我们可以将视角重新改变，以 \(\sqrt{2}\) 为主元：

\[
x^{4}-4 x^{2}-24=\left(4 x^{3}-24 x\right) \sqrt{2}
\]

这样一来，再平方一次，就可以就可以处理根号了：

\[
\begin{aligned}
& x^{8}+16 x^{4}+576-8 x^{6}-48 x^{4}+192 x^{2}=2\left(16 x^{6}-192 x^{4}+576 x^{2}\right) \\
& x^{8}-40 x^{6}+352 x^{4}-960 x^{2}+576=0
\end{aligned}
\]

所以 \(x\) 满足 \(x^{8}-40 x^{6}+352 x^{4}-960 x^{2}+576=0\) ．\\
我们从这两个例子中看到：处理根号最基本的方法就是平方。但是平方也需要注意技巧：\\
1．把根号单独放在一侧再平方，才可以有效消除根号。不然，交叉项会制造出新的根号；\\
2．每平方一次，可以消除＂一类根号＂。虽然同类根号可能出现多次，但是只要是一类根号，就可以一次平方处理掉。\\
3．平衡根号的个数。如果有多于（含）三类的二次根式在同一侧，就会出现交叉相乘产生更多根式的情况。为了避免这种情况，尽可能将根号平衡在等式的两侧。

例如，我们在第一个例子中，平方了两次，因为有 \(\sqrt{2}, ~ \sqrt{3}\) 两类根号；而在第二个例子中，平方了三次，因为有 \(\sqrt{2}, ~ \sqrt{3}, ~ \sqrt{5}\) 三类根号。平方的时候，也几乎都是把根号放在一侧，才成功消掉根号。

\section*{3．4．1．2 不等式中的二次根式}
我们在等式中，可以放心大胆地平方，至少不会出现矛盾的情况。虽然可能产生增根，但是只需后续检验即可。但是，在不等式中，如果遇到根号就随意平方，会出现非常严重的错误。

解不等式 \(x<\sqrt{x^{2}-1}\) ．很多同学在初学的时候，会直接将不等式平方：

\[
\begin{aligned}
& x^{2}<x^{2}-1 \\
& 0<-1 .
\end{aligned}
\]

我们发现矛盾了。这说明，不等式没有解吗？显然，\(x=-1\) 是不等式得解：\(-1=x<\sqrt{x^{2}-1}=0\) ．那么，哪里出现了问题呢？

事实上，如果 \(x<\sqrt{x^{2}-1}\) ，并不能推出 \(x^{2}<x^{2}-1\) ．换言之，不等式两侧不能直接平方。为什么？因为根据不等式的基本性质，当 \(a<b\) 时，是不能推出 \(a^{2}<b^{2}\) 的。例如：当 \(a=-2, b=1\) 时，\(a<b\) ，但是 \(4=a^{2}>b^{2}=1\) 。反例会在 \(a<0\) 时出现。所以，我们想要平方的话，需要确保： \(0 \leq a<b\) 才可以。此时，有：

\[
0 \leq a<b \Longrightarrow a^{2}<b^{2}
\]

这样一来，我们就可以踏踏实实平方了。\\
那么，如何解含根式的不等式呢？\\
（1）首先，我们要保证根号有意义：即 \(x^{2}-1 \geq 0\) ．故 \(x \geq 1\) 或 \(x \leq-1\) ．\\
（2）其次，为了顺利平方，我们需要分类讨论。\\
（1）当 \(x \geq 0\) 时， \(0 \leq x<\sqrt{x^{2}-1}\) ．平方得 \(x^{2}<x^{2}-1\) ．矛盾！此时无解；\\
（2）当 \(x<0\) 时，\(x<0 \leq \sqrt{x^{2}-1}\) ．不等式成立．故在根号有意义的情况下，不等式的解集为 \(\{x \mid x<0\}\) ．\\
（3）将根号有意义的范围，解集取交集，得解集为 \(\{x \mid x \leq-1\}\) 。这就是不等式的解集。\\
总结一下，我们在处理含根号的不等式时，为了能够平方，需要保证不等式两侧都是正的。因此，需要对较小的一侧进行分类讨论。基本的步骤为\\
1．保证根号有意义，\\
2．根据较小一侧和 0 的关系，分类讨论；\\
3．取交集。这样，我们就可以处理好含有根号的不等式了。

\section*{3．4．2 共轭根式与有理化}
处理根式最基本方法就是平方，他可以有效消除根式。但是，平方的运算量较大，所以我们通常希望限制平方的使用次数。此时，一个非常有效的工具就应运而生：共轭根式。

形如 \(a+\sqrt{b}, a-\sqrt{b}\) 的两个根式互为共轭根式。它们的特点是：它们的乘积不含根号。根据平方差公式，我们有：

\[
(a+\sqrt{b})(a-\sqrt{b})=a^{2}-(\sqrt{b})^{2}=a^{2}-b
\]

我们发现：共轭根式不需要将整个式子平方，就可以有效消除根号。所以，在化简根号时，它是非常有效的手段。\\
我们在初中时，就使用分母有理化的方法化简根式，它的本质就是在使用共轭根式进行化简。例如：

\[
\frac{1}{2+\sqrt{3}}=\frac{2-\sqrt{3}}{(2+\sqrt{3})(2-\sqrt{3})}=2-\sqrt{3}
\]

我们给分子，分母同乘 \(2+\sqrt{3}\) 的共轭根式 \(2-\sqrt{3}\) ，进而使得根号从分母转移到分子上，再进行下一步的处理。\\
那么，为什么我们要采用分母有理化来化简根式呢？我们放在分母上不好吗？事实上，由于分子可以进行加，减操作：

\[
\frac{x+y}{z}=\frac{x}{z}+\frac{y}{z}
\]

而分母不行：

\[
\frac{x}{y+z} \neq \frac{x}{y}+\frac{x}{z} .
\]

这意味着：如果分子有很多个复杂的部分组成，我们可以将它拆开，拆成若干个简单的部分；而若分母复杂，我们无法将它拆开，处理起来相对困难。因此，我们希望分母简单，但是分子允许复杂一些。所以，通常情况下，我们会把根号从分母转移到分子上。

\section*{3．4．3 XXXX}
\section*{3.5 自由度，消元}
在第一章中我们以自由度这个概念为出发点，引出了本讲义内容的核心思想方式。而在本章节中，我们依旧用自由度来解释一些特殊的问题，你不仅能更全面的理解这个概念，也能将它作为解决计算的一支巨大的杜杆。

\section*{3．5．1 对自由度更深的理解}
自由度这个概念事实上来自于物理，它有助于我们理解系统的运行机制。那么，从数学的角度如何理解自由度呢？它如何帮助我们提高计算的准确率和速度呢？让我们从最基本的线性方程组出发来理解这个概念。

我们小学学习过鸡兔同笼问题。假设共有 10 只鸡和兔，它们共有 28 只脚。那么，各有多少只鸡和兔呢？如果我们利用方程思想来解决这个问题，我们可以设鸡，兔的数量分别是 \(x, y\) ，那么根据题目条件：

\[
\left\{\begin{array}{l}
x+y=10 \\
2 x+4 y=28
\end{array}\right.
\]

解得 \(x=6, y=4\) ．\\
我们通过基本的二元一次方程组解决了问题。但是，大家有没有想过：为什么一定只有这样一组解呢？有没有别的可能呢？我们从条件入手来理解这个问题。\\
（1）如果我们只有条件＂10 只鸡和兔＂，那么可以有很多组解。例如， 5 只鸡， 5 只兔或者 6 只鸡， 4 只兔；\\
（2）如果我们只有条件＂28 只脚＂，那么也可以有很多组解。例如， 7 只兔子， 14 只鸡， 2 只鸡和 6 只兔子，等等。

换言之，如果我们只有题中条件中的一个，那么我们是无法确定具体只数的，无论是哪个条件。但是如果是两个条件，就大不相同了：我们要在满足条件 1 的情况下，找出也满足条件 2 的情况。这样一来，一会发现，如果我们把满足条件 1 的解都列出来：

\begin{center}
\begin{tabular}{|c|c|}
\hline
鸡的数量 & 免子的数量 \\
\hline
0 & 10 \\
\hline
1 & 9 \\
\hline
2 & 8 \\
\hline
3 & 7 \\
\hline
4 & 6 \\
\hline
5 & 5 \\
\hline
6 & 4 \\
\hline
7 & 3 \\
\hline
8 & 2 \\
\hline
9 & 1 \\
\hline
10 & 0 \\
\hline
\end{tabular}
\end{center}

那么我们会得到 11 种可能。这 11 种可能，我们应该这么理解：\\
（1）当鸡的数量确定以后，兔子的数量也确定了。令鸡的数量为 \(x\) ，则兔子的数量必然为 \(10-x\) ；\\
（2）当兔子的数量确定以后，鸡的数量也确定了。令兔子的数量为 \(y\) ，则鸡的数量必然为 \(10-y\) 。\\
换言之，本来鸡，兔数量两个没有关系的量，我们经过一个条件的限制，让它们可以互相决定了！这样一来，本来＂两个＂彼此无关的量，变成了＂一个＂独立的变量了，因为一个可以决定另外一个，它们之间建立了联系。

但是问题到这，并没有解决。虽然两个变量之间建立了联系，但是还是有一个变量可以随意改变。如果我们以鸡的数量为基准，那么兔的数量取决于鸡的数量，但是鸡的数量依旧可以在 010 之间随意改变，整个系统依旧没有被确定下来。

想要确定它，就需要另外一个条件了，脚的只数总和为 48 ．当我们把所有满足条件 1 的情况列出来，并计算每种情况脚的只数时，我们可以得到如下表格：

\begin{center}
\begin{tabular}{|c|c|c|}
\hline
鸡的数量 & 免子的数量 & 脚的数量 \\
\hline
0 & 10 & 40 \\
\hline
1 & 9 & 38 \\
\hline
2 & 8 & 36 \\
\hline
3 & 7 & 34 \\
\hline
4 & 6 & 32 \\
\hline
5 & 5 & 30 \\
\hline
6 & 4 & 28 \\
\hline
7 & 3 & 26 \\
\hline
8 & 2 & 24 \\
\hline
9 & 1 & 22 \\
\hline
10 & 0 & 20 \\
\hline
\end{tabular}
\end{center}

我们发现：随着鸡的数量增多，脚的数量越来越少；反过来，随着兔子的数量增多，脚的数量越来越多。它们在 2040 之间变化。所以，我们只要找到共 28 只脚的那个组合，我们就解决了问题。表格显示：当我们有 6只鸡，4只兔子的时候，刚好有 28 只脚。这样，最终的变量也就确定了，问题就解决了。这种感觉，就像是一种特殊的密码：我们知道明文是 28 （我们数出来有 28 只脚），我们要找到它对应的暗文是多少。找到了，我们就把藏在后面的答案找到了。

回顾我们刚才对于系统的分析，我们发现：我们最终确定了两个本来无关的变量 \(x, y\) ，但其实，它们也是逐步确定的：\\
（1）根据第一个条件，我们确定了 \(x, y\) 的关系。此时，我们使用了一个限制条件，将两个无关的变量，变成了一个决定另一个，真正在变的，只有一个变量；\\
（2）根据第二个条件，我们决定了 \(x\) 。在所有可能的组合里，只有 1 个满足条件 2 。这样一来，本来在变化的 \(x\) ，也被确定了。

整个过程，像是探案一样：我们先通过一个限制条件，把 \(y\) 建立在 \(x\) 上，将两个独立变量变成一个独立变量；再利用一个限制条件，把 \(x\) 也确定，最终问题解决。两个限制条件，我们确定了两个变量，所以问题得到了回答。

在这个过程中，变量的总个数我们称为问题的维数，独立变量的个数我们称为自由度。什么叫独立变量？当我们在条件 1 的情况下寻找 \(x, y\) 时，独立变量只有一个：它确定了，整个系统就确定了，其他的变量都随着它的改变而改变。而一开始，我们不使用任何条件的时候，我们有两个独立变量，所以自由度是 2 。

我们发现，在上面的过程中，自由度，限制条件个数和维数，有着如下关系：\\
（1）当我们使用 0 个条件时，共有 2 个自由变量：\(x, y, 0+2=2\) ，是维数；\\
（2）当我们使用 1 个条件时，共有 1 个自由变量：\(x\) 或 \(y\) ，它们相互决定， \(1+1=2\) ，也是维数；\\
（3）当我们使用 2 个条件时，不再有自由变量了：\(x, y\) 均被确定。此时自由度为 \(0,2+0=2\) ，也是维数。\\
我们发现：无论何时，都成立下面的等式：

\[
\text { 自由度 }+ \text { 限制条件个数 = 维数 }
\]

在线性方程组的理论中，它叫秩－零化度定理（rank－nullity theorem）。除了这么理解，我们也可以理解成：每个限制条件都可以确定一个自由度。这样，限制条件，自由度此消彼长，维持和不变，永远等于维数。

这也就解释了，为什么当我们有两（维数）个未知数时，我们需要两个方程（限制条件）才能够解出问题。解出问题，相当于系统彻底确定，这是一个自由度为 0 的状态。根据秩－零化度定理，我们必须要求方程个数等于维数。

这样，我们就从代数结构上理解了鸡兔同笼问题。更一般地，多个未知数的线性方程组都可以这么理解。\\
这个理解虽然不能帮助我们直接把方程组解出来，但是有助于我们理解问题的运行机制。你可以问自己：这个问题到底有几个自由度？哪些变量可以完全确定这个系统？怎么设未知数更方便？只有把这些想清楚了，我们才算是完成了分析问题的第一步，也算是为我们的计算打下了坚实的基础。那么，想要从技术层面上真的解好方程组，需要做什么呢？这就需要我们的消元思想了。

\section*{3．5．2 什么是消元思想？}
我们在学习前面的几个杠杆时，我们已经意识到：为了最终求取值范围，最值，我们都需要把多元问题变为一元问题，因为我们所有熟悉的工具和方法，几乎都是针对以一元问题的。这个基本思想就是消元。事实上，这个基本思想在解方程时是最常用的。

我们从三个未知数的线性方程组来说起。例如，考虑

\[
\left\{\begin{array}{l}
x+y=2 z=9 \\
2 x-y=3 z=9 \\
3 x-2 y-z=-4
\end{array}\right.
\]

这里共有三个未知数，所以维数是 3 ；共有三个方程，所以限制条件个数也是 3 。因此，一般情况下，这个系统是确定的。那么，我们能不能仿照刚才的做法，将它进行化简呢？

出发时，我们还没有使用任何一个方程，我们处在一个自由度为 3 的状态。 \(x, y, z\) 三个未知数之间没有关系。为了让他们建立联系，我们考虑使用第一个方程。如何使用？我们将 \(x\) 用 \(y, z\) 去表示：

\[
x=9-y-2 z .
\]

使用了一个限制条件后，我们处在一个什么状态呢？我们处于自由度为 2 的状态。此时，两个独立变量是 \(y, z\) ，我们使用了一个限制条件，另一个变量 \(x\) 就是利用这个限制条件被 \(\mathrm{y}, \mathrm{z}\) 表示的。

我们将这个方程，代人到后面两个方程中，得到

\[
\left\{\begin{array}{c}
2(9-y-2 z)-y+3 z=9 \\
3(9-y-2 z)-2 y-z=-4
\end{array}\right.
\]

化简得

\[
\left\{\begin{aligned}
3 y+z & =9 \\
5 y+7 z & =31
\end{aligned}\right.
\]

这样，我们使用第一个方程，把 \(x\) 消掉，得到了一个关于 \(\mathrm{y}, \mathrm{z}\) 的方程。像这样，利用方程（限制条件），将其中一个变量用其他变量表示的方法，并将它在所有方程种替换掉，就是代人消元法。

下一步，我们想要再减少一个自由度。所以，我们考虑将第一个方程改写为现在，是一个自由度为 1 的状态：我们使用了两个限制条件，\(y\) 是独立变量，\(x\) 由 \(y, z\) 决定，\(z\) 由 \(y\) 决定，这等价于 \(x, z\) 均由 \(y\) 决定：

\[
\begin{aligned}
x & =9-y-2 z \\
& =9-y-2(9-3 y) \\
& =5 y-9 .
\end{aligned}
\]

为了最终解出方程，我们将 \(z=9-3 y\) 代人到最后一个方程，得到

\[
\begin{aligned}
& 5 y+7(9-3 y)=31 \\
& 63-16 y=31
\end{aligned}
\]

解得 \(y=2\) ．代回到之前的表达，得到

\[
x=5 y-91=1, z=9-3 y=2 .
\]

故方程的解为

\[
\left\{\begin{array}{l}
x=1 \\
y=2 \\
z=3
\end{array}\right.
\]

我们一共使用了两次代入消元法。第一步，将 \(x\) 用 \(y, z\) 表示，得到一个关于 \(y, z\) 的方程组；第二步，在这个方程组中，将 z 用 y 表示，得到一个只含有 \(y\) 的方程，解出它们。这就是代入消元法的最终目标：将方程组化为

一元方程，并求解。像刚才一样，我们回顾一下整个过程中系统的状态，可以得到如下表格：

\begin{center}
\begin{tabular}{|c|c|c|c|c|}
\hline
自由度 & \begin{tabular}{c}
使用的条 \\
件个数 \\
\end{tabular} & 独立变量 & \begin{tabular}{c}
和独立变量的关系 \\
（已使用的方程） \\
\end{tabular} & \begin{tabular}{c}
原方程组化为 \\
（未使用的方程） \\
\end{tabular} \\
\hline
3 & 0 & \(x, y, z\) & 无 & \(\left\{\begin{array}{l}x+y+2 z=9, \\ 2 x-y+3 z=9, \\ 3 x-2 y-z=-4 .\end{array}\right.\) \\
\hline
2 & 1 & \(y, z\) & \begin{tabular}{l}
\(x=9-y-2 z\) \\
\end{tabular} & \(\left\{\begin{array}{l}3 y+z=9, \\ 5 y+7 z=31 .\end{array}\right.\) \\
\hline
1 & 2 & \(y\) & \(\left\{\begin{array}{l}x=9-2 y-z \\ =5 y-9, \\ z=9-3 y .\end{array}\right.\) &  \\
\hline
0 & 3 & 无 & \(\left\{\begin{array}{l}x=5 y-9=1, \\ y=2, \\ z=9-3 y=3 .\end{array}\right.\) & \begin{tabular}{l}
解完了 \\
\end{tabular} \\
\hline
\end{tabular}
\end{center}

我们发现：自由度和限制条件个数之和保持不变，这也印证了我们秩－零化度定理。\\
当我们理解了自由度和消元思想后，我们能不能将他们应用到更广泛的问题上呢？我们之前学习的杜杆：主元，对称，齐次，共轭根式，能不能也用自由度和消元思想来解释呢？我们在例题中来回顾一下。

\section*{3．5．3 XXXX}
\section*{3.6 因式分解}
\section*{3．6．1 一元二次式}
\section*{3．6．1．1 解方程}
我们在初中就学习过因式分解，解二次方程。事实上，一元二次式的因式分解本质上就是在解二次方程。考虑二次式 \(x^{2}+b x+c\) ，如果它的因式分解结果为

\[
x^{2}+b x+c=\left(x-x_{1}\right)\left(x-x_{2}\right)
\]

那么 \(x^{2}+b x+c=0\) 等价于 \(\left(x-x_{1}\right)\left(x-x_{2}\right)=0\) ，即它的两根为 \(x_{1}, x_{2}\) 。因此，我们只需要求出方程的解，就可以完成因式分解。

根据二次方程求根公式，我们知道 \(a x^{2}+b x+c=0\) 的两根为

\[
x_{1,2}=\frac{-b \pm \sqrt{b^{2}-4 a c}}{2 a}
\]

故我们可以将 \(a x^{2}+b x+c=0\) 分解为

\[
\begin{aligned}
a x^{2}+b x+c & =a\left(x-x_{1}\right)\left(x-x_{2}\right) \\
& =a\left(x-\frac{-b+\sqrt{b^{2}-4 a c}}{2 a}\right)\left(x-\frac{-b-\sqrt{b^{2}-4 a c}}{2 a}\right) .
\end{aligned}
\]

这就完成了分解。因此，在高中的学习中，如果需要对一元二次式进行因式分解，只需利用求根公式求出两根即可。

\section*{3．6．1．2 十字相乘法}
这是我们在课内学习的最基本的二次式因式分解的方法。考虑二次式 \(x^{2}+b x+c\) ，假设我们能够将它分解成

\[
x^{2}+b x+c=\left(x+a_{1}\right)\left(x+a_{2}\right)
\]

的形式，那么将右侧展开，得到

\[
x^{2}+b x+c=x^{2}+\left(a_{1}+a_{2}\right) x+a_{1} a_{2} .
\]

由于这个式子对于任意的 \(x\) 的都成立，所以对应项相等：

\[
a_{1}+a_{2}=b, a_{1} a_{2}=c
\]

换言之，我们只需要找到这样的 \(a_{1}, a_{2}\) 就可以完成分解。\\
十字相乘法的本质在于，通过对 \(c\) 的分解进行尝试，最终找到合适的 \(a_{1}, a_{2}\) 。我们以 \(x^{2}+5 x+6\) 为例，我们需要找到 \(a_{1}, a_{2}\) 使得

\[
a_{1}+a_{2}=5, a_{1} a_{2}=6 .
\]

注意到 \(6=1 \times 6=2 \times 3\) ，我们发现：\\
（1）若 \(a_{1}=1, a_{2}=6\) ，则 \(a_{1}+a_{2}=7\) ，不符合；\\
（2）若 \(a_{1}=2, a_{2}=3\) ，则 \(a_{1}+a_{2}=5\) ，满足条件。\\
我们通过基本的尝试发现，\(a_{1}=2, a_{2}=3\) 是满足条件的，所以

\[
x^{2}+5 x+6=(x+2)(x+3)
\]

在展开计算右侧表达式的过程中，我们有如下的步骤：\\
（1）先去计算二次项：\(x \times x=x^{2}\) ；\\
（2）再去计算一次项：\(x \times 3+x \times 2=5 x\) ，\\
（3）最终计算常数项： \(2 \times 3=6\) ．\\
如果我们把这个过程更形象化地表示出来，我们可以这么表达：

\[
\begin{aligned}
& x \\
& x^{x}+2 \\
& +3
\end{aligned}
\]

（1）对应着左侧的两个 \(x\) 相乘，（3）对应着右侧的 \(+2,+3\) 相乘，而最关键的（2）对应着这样的十字相乘。因此，我们可以借助这样的十字相乘来辅助我们进行因式分解。对于第一种情况 \(a_{1}=1, a_{2}=6\) ，我们得到的十字相乘的结果为

\[
\begin{aligned}
& x+1 \\
& x \\
& \times \begin{array}{l}
+1
\end{array}
\end{aligned}
\]

这样得到的 \(x\) 的系数为 \(1+6=7\) ，不是 5 ．所以不是正确的十字相乘结果。这样一来，我们就可以总结出利用十字

相乘进行因式分解的基本步骤了：\\
（1）写出算式：

\[
\begin{aligned}
& 1 \times \begin{array}{l}
+1 \\
1
\end{array}+6
\end{aligned}
\]

我们通常会省略前面的 \(x\) ，用 1 来代替；\\
（2）十字相乘：对不同的 \(c\) 的分解进行尝试，直到十字相乘的结果能够得到正确的一次项系数。

\[
\begin{aligned}
& 1 \times+2 \\
& 1 \times+3
\end{aligned}
\]

（3）还原：将式子相乘的结果还原成因式分解：\(x^{2}+5 x+6=(x+2)(x+3)\) ．\\
这就是基本的十字相乘法分解因式，我们将会在例题中看到更多的变式。

\section*{3．6．2 因式定理}
由于二次方程的因式分解和根有着天然的联系，我们发现：如果 \(x_{1}\) 恰好是二次方程 \(x^{2}+b x+c=0\) 的根，那么 \(\left(x-x_{1}\right)\) 必然是代数式 \(x^{2}+b x+c\) 的一个因式。这是因为：由于 \(x_{1}\) 是 \(x^{2}+b x+c=0\) 的根，所以

\[
x_{1}^{2}+b x_{1}+c=0
\]

直接的计算可得：

\[
\left(x-x_{1}\right)\left(x-\frac{c}{x_{1}}\right)=x^{2}-\left(x_{1}+\frac{c}{x_{1}}\right) x+c
\]

由于 \(x_{1}^{2}+b x_{1}+c=0\) ，故

\[
b=-\frac{x_{1}^{2}+c}{x_{1}}=-\left(x_{1}+\frac{c}{x_{1}}\right) .
\]

所以

\[
\left(x-x_{1}\right)\left(x-\frac{c}{x_{1}}\right)=x^{2}-\left(x_{1}+\frac{c}{x_{1}}\right) x+c=x^{2}+b x+c .
\]

我们不光证明了 \(\left(x-x_{1}\right)\) 是它的一个因式，同时求出了另一个因式 \(\left(x-\frac{c}{x_{1}}\right)\) ．细心的同学会发现，\(x_{1} \cdot \frac{c}{x_{1}}=c\) ，这正是我们后续会学习的韦达定理。

因此，对于一个多项式 \(f(x)\) 而言，\(\left(x-x_{1}\right)\) 是 \(f(x)\) 的一个因式当且仅当 \(x=x_{1}\) 是 \(f(x)=0\) 的一个根。这就是因式定理，它为我们分解高次多项式奠定了基础：我们可以通过试根法来确定它的一次因式，再进行后续的分解。

\section*{3．6．3 一元三次式与高次式}
\section*{3．6．3．1 一元三次式}
对于三次式 \(x^{3}+a x^{2}+b x+c\) ，如果我们能够通过试根法得到一个根 \(x=x_{1}\) ，那么我们就可将三次降为二次：

\[
x^{3}+a x^{2}+b x+c=\left(x-x_{1}\right)\left(x^{2}+d x+e\right)
\]

对于二次部分，我们再采取相应的方法进行分解即可。所以，对于三次和高次多项式的分解，我们通常第一步会采用试根法。

一个试根的技巧是：对于最高次项系数为 1 的多项式（如 \(x^{3}+a x^{2}+b x+c, x^{2}+b x+c\) ），可以尝试常数项的\\
＂因数＂。例如，若常数项为 +6 ，那么可以优先尝试：

\[
-6,-3,-2,-1,1,2,3,6
\]

例如，\\
（1）对于多项式 \(g(x)=x^{3}-3 x^{2}+2 x-6\) ，我们发现：当 \(x=3\) 时，

\[
g(-3)=27-27+6-6=0 .
\]

故 \(x-3\) 是 \(f(x)\) 的一个因式；\\
（2）对于多项式 \(f(x)=x^{3}-2 x^{2}-2 x+1\) ，我们发现：当 \(x=-1\) 时，

\[
f(-1)=-1+2-1+1=0 .
\]

故 \(x-(-1)=x+1\) 是 \(f(x)\) 的一个因式。\\
所以，通过试根法试出三次式的一次因式是一个普遍的方法。现在的问题是：当我们试出这个一次因式后，剩下的二次部分应该如何求解呢？在这里，我们介绍两个基本方法：待定系数法和它的简化版本。

首先我们介绍待定系数法。以 \(f(x)=x^{3}-2 x^{2}-2 x+1\) 为例，我们已经直到 \(x+1\) 是它的一个因式了，我们设剩下的因式为 \(a x^{2}+b x+c\) ，即

\[
x^{3}-2 x^{2}-2 x+1=(x+1)\left(a x^{2}+b x+c\right) .
\]

将右侧展开，得到

\[
(x+1)\left(a x^{2}+b x+c\right)=a x^{3}+(a+b) x^{2}+(b+c) x+c .
\]

故

\[
x^{3}-2 x^{2}-2 x+1=a x^{3}+(a+b) x^{2}+(b+c) x+c .
\]

这对于任意的 \(x\) 都成立，所以对应项系数相等，我们得到

\[
\left\{\begin{array}{l}
a=1 \\
a+b=-2 \\
b+c=-2 \\
c=1
\end{array}\right.
\]

解得 \(a=1, b=-3, c=1\) ．因此

\[
x^{3}-2 x^{2}-2 x+1=(x+1)\left(x^{2}-3 x+1\right)
\]

这就完成了分解因式。像这样，设未知的系数为 \(a, b, c\) ，再通过对比系数等方法求解的系数方法就是待定系数法。

在求解的过程中，我们发现：虽然方程组有四个，但是未知数只有三个，可能会出现无解的情况。事实上，我们可以证明这个方程一定存在解：这个由因式定理保证，这里我们略去不表。

对于三次，四次等次数不是特别高的方程而言，我们可以将上述待定系数法的步骤进行简化。我们还是以 \(x^{3}-2 x^{2}-2 x+1\) 为例，设

\[
x^{3}-2 x^{2}-2 x+1=(x+1)\left(a x^{2}+b x+c\right) .
\]

我们能不能通过简单的计算，直接求出 \(a, c\) 的值呢？事实上是可以的。\\
如果我们比较两侧三次项的系数，那么左侧为 \(1 \cdot x^{3}\) ，系数为 1 ；而右侧的展开式中必然是 \(a x^{2} \cdot x=a x^{3}\) ，系

数为 \(a\) ．根据对应项相等，必然有

\[
a=1 .
\]

这样就确定了 \(a\) 的值。\\
如果我们比较常数项的系数，那么左侧为 +1 ；而右侧的展开式中必然为 \(1 c=c\) ，系数为 \(c\) ．根据对应项相等，必然有 \(c=1\) ．这样就确定了 \(c\) 的值。

因此，我们不需要列出任何方程，直接简单的比较最高次项和常数项的系数，就可以求出 \(a, c\) ．一共三个待定系数中，我们已经求出了两个，说明这个方法还是非常便捷的。

最后，对于 \(b\) ，我们可以通过两种方法求解。如果比较一次项系数，左侧为 -2 ，右侧为 \(b+c\) ，所以

\[
b+c=-2 \text {. }
\]

由于 \(c=1\) ，所以

\[
b=-2-c=-3 .
\]

这样就可以求出 \(b\) 的值。\\
如果比较二次项系数，左侧为 -2 ，右侧为 \(a+b\) ，所以

\[
a+b=-2 \text {. }
\]

由于 \(a=1\) ，所以

\[
b=-2-a=-3 .
\]

这和我们比较一次项系数得到的结果是相同的。实际应用中，大家可以采取任意一种方法。推荐读者使用一种方法，用另一种方法检验。

我们总结一下简化版的待定系数法的分解三次多项式方法。\\
（1）比较左右两侧三次项（最高次项）系数，求解 \(a\) ；\\
（2）比较左右两侧常数项系数，求解 \(c\) ；\\
（3）比较左右两侧一次项系数或二次项系数，求解 \(b\) ．\\
总结一下：分解三次多项式，我们先试根，找出一次因式；再通过简化版本的待定系数法找出剩下的二次因式，最后如有需要，将二次因式再分解即可。

\section*{3．6．3．2 高次式}
高次式和三次式的分解方法没有本质区别，都是试根配合待定系数法。对于不高于五次的多项式，不需使用大除法。在高中范围内，熟练掌握三次多项式的分解因式足以应对题目中出现的 \(99 \%\) 的计算了。在这里我们不介绍大除法的使用，因为和简化版的待定系数法相比，它的速度较慢。

\section*{3．6．4 多元}
和一元问题不同，多元的分解因式在处理时方法更多，但主要还是使用我们之前介绍过的所有代数杜杆。\\
（1）主元法：选主元，化为一元问题；\\
（2）对称，齐次：利用代数结构将问题化简；\\
（3）使用多元的因式分解公式。\\
前面的代数结构我们已经详细介绍过了，这里着重介绍重要的分解因式公式，希望大家掌握。

1．\(a^{2}-b^{2}=(a-b)(a+b)\) ；\\
2．\(a^{3}-b^{3}=(a-b)\left(a^{2}+a b+b^{2}\right)\) ；\\
3．\(a^{n}-b^{n}=(a-b)\left(a^{n-1}+a^{n-2} b+\cdots+a b^{n-2}+b^{n-1}\right)\) ；\\
4．\(a^{3}+b^{3}=(a+b)\left(a^{2}-a b+b^{2}\right)\) ；\\
5．\(a^{n}+b^{n}=(a+b)\left(a^{n-1}-a^{n-2} b+\cdots-a b^{n-2}+b^{n-2}\right)\) ，其中 \(n\) 为奇数；\\
6．\((a+b)^{3}=a^{3}+3 a^{2} b+3 a b^{2}+b^{3}\) ；\\
7．\((a-b)^{3}=a^{3}-3 a^{2} b+3 a b^{2}-b^{3}\) ；\\
8．\((a+b+c)^{2}=a^{2}+b^{2}+c^{2}+2 a b+2 a c+2 b c\) ；\\
9．\(a^{3}+b^{3}+c^{3}-3 a b c=(a+b+c)\left(a^{2}+b^{2}+c^{2}-a b-a c-b c\right)\) ．

\[
=\frac{(a+b+c)\left[(a-b)^{2}+(b-c)^{2}+(c-a)^{2}\right]}{2}
\]

这些公式是要求熟练掌握的。为了给大家更直观的感受，我们将其中立方和公式的＂右式＂分解因式，看是否能够得到左侧。这个步骤的意义在于，我们体会前面讲到的代数杜杆，是如何在分解因式这个小问题上发挥巨大功效的。

我们试着分解 \(a^{3}+3 a^{2} b+3 a b^{2}+b^{3}\) 。在这个式子中，我们可以看到非常多的代数结构。\\
（1）对称：它关于 \(a, b\) 对称；\\
（2）齐次：每一项都是关于 \(a, b\) 的三次式。\\
我们分别利用对称，齐次两个杜杆来分解因式。\\
（1）对称：如果这是一个关于 \(a, b\) 的对称式，我们可以自然地想到使用对称换元，即令 \(m=a+b, n=a b\) ．那么 \(3 a^{2} b+3 a b^{2}=3 a b(a+b)=3 m n\) 。现在只需要处理 \(a^{3}+b^{3}\) 即可。除了直接使用分解因式公式，我们还可以像第二节中一样，利用递推求解：

\[
\begin{aligned}
a^{3}+b^{3} & =(a+b)\left(a^{2}+b^{2}\right)-a b(a+b) \\
& =m\left(m^{2}-2 n\right)-m n \\
& =m^{3}-3 m n
\end{aligned}
\]

故

\[
a^{3}+3 a^{2} b+3 a b^{2}+b^{3}=m^{3}-3 m n+3 m n=m^{3}=(a+b)^{3} .
\]

这样，就完成了分解。\\
（2）齐次：如何使用齐次进行分解因式呢？齐次最大的作用在于，如果我们能够得到一个零次齐次式，那么就可以使用比例换元，将问题化为一元问题。这里，我们只有一个三次齐次式，如何化为零次齐次式呢？如果我们将它先化为分式，整体提出一个 \(b^{3}\) ，会得到什么呢？

\[
a^{3}+3 a^{2} b+3 a b^{2}+b^{3}=b^{3}\left[\left(\frac{a}{b}\right)^{3}+3\left(\frac{a}{b}\right)^{2}+3\left(\frac{a}{b}\right)+1\right] .
\]

此时，我们后边得到了一个关于 \(t=\frac{a}{b}\) 的三次多项式，它就是 \(t^{3}+3 t^{2}+3 t+1\) 。因此，我们只需要对这个一元三次式分解因式即可。我们首先尝试 1 的因数，即土 1 ．发现：当 \(t=-1\) 时，有

\[
t^{3}+3 t^{2}+3 t+1=-1+3-3+1=0
\]

故 \(t+1\) 是它的一个因式。根据待定系数法，得到：

\[
t^{3}+3 t^{2}+3 t+1=(t+1)\left(t^{2}+2 t+1\right)=(t+1)(t+1)^{2}=(t+1)^{3} .
\]

这样，就完成了因式分解。回到原题，将 \(t=\frac{a}{b}\) 代回，得到

\[
a^{3}+3 a^{2} b+3 a b^{2}+b^{3}=b^{3}\left(\frac{a}{b}+1\right)^{3}=(a+b)^{3}
\]

因此，即使是因式分解中，也可以使用代数杜杆重新理解问题。

\section*{3.7 二次式}
\section*{3．7．1 韦达定理}
设二次方程 \(a x^{2}+b x+c=0\) 有两个实根 \(x_{1}, x_{2}\) 。已经学过根系关系（韦达定理）的同学知道，我们可以利用系数表示 \(x_{1}+x_{2}, x_{1} x_{2}\) ：

\[
\left\{\begin{array}{l}
x_{1}+x_{2}=-\frac{b}{a} \\
x_{1} x_{2}=\frac{c}{a}
\end{array}\right.
\]

那么，我们如何证明这个结论呢？我们能够马上想到的方法，就是利用求根公式把他们求出来，然后代人计算。根据求根公式，

\[
x_{1}=\frac{-b+\sqrt{b^{2}-4 a c}}{2 a}, x_{2}=\frac{-b-\sqrt{b^{2}-4 a c}}{2 a}
\]

因此：

\[
x_{1}+x_{2}=\frac{-b+\sqrt{b^{2}-4 a c}}{2 a}+\frac{-b-\sqrt{b^{2}-4 a c}}{2 a}=-\frac{2 b}{2 a}=-\frac{b}{a} .
\]

和非常好验证。对于乘积，我们发现：

\[
x_{1} x_{2}=\frac{-b+\sqrt{b^{2}-4 a c}}{2 a} \cdot \frac{-b-\sqrt{b^{2}-4 a c}}{2 a}
\]

这里出现了一个我们之前反复强调的代数结构：共轭根式。事实上，\(x_{1}, x_{2}\) 的两个分子互为共轭根式，所以我们可以直接观察到平方差的结构：

\[
\begin{aligned}
\left(-b+\sqrt{b^{2}-4 a c}\right) \cdot\left(-b-\sqrt{b^{2}-4 a c}\right) & =(-b)^{2}-\left(\sqrt{b^{2}-4 a c}\right)^{2} \\
& =b^{2}-\left(b^{2}-4 a c\right) \\
& =4 a c .
\end{aligned}
\]

因此，

\[
x_{1} x_{2}=\frac{4 a c}{4 a^{2}}=\frac{c}{a}
\]

这样，也就完成了韦达定理的证明。\\
这个方法直接，但是相对计算较慢。我们能不能从因式分解的角度来理解它呢？事实上，我们可以直接利用上节课讲过的因式定理。由于 \(x_{1}, x_{2}\) 是方程的两个根，我们知道 \(x-x_{1}, x-x_{2}\) 一定是二次式 \(a x^{2}+b x+c\) 的两个因式，所以

\[
a x^{2}+b x+c=a\left(x-x_{1}\right)\left(x-x_{2}\right)
\]

这本质上相当于二次函数的两点式。将右侧的括号展开，得到

\[
a x^{2}+b x+c=a x^{2}-a\left(x_{1}+x_{2}\right) x+a x_{1} x_{2} .
\]

由于这个式子对于任意的 \(x\) 都成立，对应项的系数必然相等，所以

\[
-a\left(x_{1}+x_{2}\right)=b, a x_{1} x_{2}=c .
\]

即：

\[
\left\{\begin{array}{l}
x_{1}+x_{2}=-\frac{b}{a} \\
x_{1} x_{2}=\frac{c}{a}
\end{array}\right.
\]

这就是所谓的韦达定理，也称根系关系。韦达定理是我们最早接触＂对称＂这个代数结构的地方，它本质上就是我们之前提到的对称换元。

一个简单的韦达定理，我们就观察到了几乎之前讲过的所有代数结构：对称，共轭根式，因式分解等等。它们在我们理解所有代数定理时都扮演者举足轻重的地位，希望大家在学习基本概念和理解新问题的时候，把它们紧紧握在手中。

\section*{3．7．2 判别式}
\section*{3．7．2．1 判别式的含义}
在我们熟悉了根系关系后，能不能从它的角度来理解判别式呢？事实上，如果我们将 \(b=-a\left(x_{1}+x_{2}\right), c=a x_{1} x_{2}\)代人到判别式中，我们可以得到：

\[
\Delta=b^{2}-4 a c=a^{2}\left(x_{1}+x_{2}\right)^{2}-4 a^{2} x_{1} x_{2}=a^{2}\left(x_{1}-x_{2}\right)^{2}
\]

所以，如果方程存在两个实根，必然有 \(a^{2}\left(x_{1}-x_{2}\right)^{2} \geq 0\) ，从而 \(\Delta \geq 0\) ．这就是为什么我们可以使用判别式来判断二次方程是否有实根。

在计算二次函数零点间的距离时，我们也可以利用这个公式。两边开方，得到

\[
\left|x_{1}-x_{2}\right|=\frac{\sqrt{\Delta}}{|a|}
\]

\section*{3．7．2．2 判别式的计算}
在实际应用中，我们经常会需要计算一个具体的二次方程的判别式。可能是代字母的，可能是不代字母的。很多同学在计算时，会直接将判别式的值＂算出来＂，再去开方。事实上，由于我们经常需要使用判别式去开方，所以，在计算时要主动＂提公因式（数）＂，而不是算完再开方，这样可以减少很多计算。我们举一个简单的例子。考虑二次方程 \(9 x^{2}+36 x-16=0\) ．计算它的判别式，可以得到

\[
\Delta=36^{2}+9 \cdot 4 \cdot 16
\]

如果我们直接把所有数计算出来，那么可以得到

\[
\Delta=1296+576=1872
\]

这样的话，我们需要计算 \(36^{2} 9 \cdot 4 \cdot 16\) ，并计算 1872 开方后的结果。像这样较为复杂的四位数运算和开方运算，是相对浪费时间的。

那么，正确的做法是什么呢？正确的做法是反向操作，不去算乘法，去分解质因数，然后去提公因数。具体

而言，我们要把 \(36^{2}\) 表示成

\[
36^{2}=(4 \cdot 9)^{2}=4 \cdot 4 \cdot 9 \cdot 9
\]

我们把它写成若干个平方因子的乘积后，再和后面的数提公因数：

\[
\begin{aligned}
\Delta & =4 \cdot 4 \cdot 9 \cdot 9+9 \cdot 4 \cdot 1 \\
& =4 \cdot 9 \cdot(9 \cdot 4+16) \\
& =4 \cdot 9 \cdot 4 \cdot(9+4) \\
& =4 \cdot 9 \cdot 4 \cdot 13 .
\end{aligned}
\]

因此，

\[
\sqrt{\Delta}=\sqrt{4 \cdot 9 \cdot 4 \cdot 13}=\sqrt{4} \cdot \sqrt{9} \cdot \sqrt{4} \cdot \sqrt{13}=2 \cdot 3 \cdot 2 \cdot \sqrt{13}=12 \sqrt{13} .
\]

这样，我们甚至不需要处理大于 20 的数，就完成了一个四位数的加法计算和开方操作。对于含有分母的判别式，我们也要采取同样的策略。这将在例题 \({ }^{2}\) 中看到。

\section*{3．7．3 XXXX}
\footnotetext{2 见《高中数学积累本》
}\section*{第4章 计算基本法—实战}
\section*{4.1 函数的性质}
我们在学习函数时，除了学习它的定义和基本要素（对应关系，定义域，值域）等，我们还学习了它的有关性质。这是我们研究任何一个新的数学对象的基本方法：我们先关心它到底是什么（定义），然后关心它到底如何，怎么样（性质）。在这一节中，我们会着重讨论涉及到大量计算和代数结构的函数性质，如：奇偶性，单调性和值域。

\section*{4．1．1 奇偶性，对称性}
奇偶性是一种特殊的对称性。如果一个函数是奇函数，那么它的图像关于原点 \(O\) 对称；如果一个函数是偶函数，那么它的图像关于 \(y\) 轴 \(x=0\) 对称。从代数上讲，一个函数是奇函数当且仅当它的定义域关于原点对称，且

\[
f(x)+f(-x)=0 .
\]

一个函数是偶函数当且仅当它的定义域关于原点对称，且

\[
f(x)=f(-x)
\]

这里我们需要特别强调＂定义域关于原点对称＂。因为我们经常会遇到例子，它虽然满足后面的代数条件，但是不满足定义域对称，导致它的图像不是一个对称图形，进而是非奇非偶函数。例如：\(f(x)=x^{2}, x \geq 1\) ．

很多同学可能会有疑问：奇偶性为什么和我们的代数结构有关系呢？我们来看一个简单的例子。考虑函数

\[
f(x)=\ln \left(x+\sqrt{1+x^{2}}\right) .
\]

很多同学看到这个函数，下意识就认为它不是奇函数，也不是偶函数。为什么呢？\(y=x\) 是奇函数，\(y=\sqrt{1+x^{2}}\)是偶函数，相加得到 \(y=x+\sqrt{1+x^{2}}\) 就是非奇非偶函数了。外面还套了一层对数函数，怎么可能具有对称性呢？事实上，我们只需要检验奇函数的定义即可：

\[
\begin{aligned}
f(x)+f(-x) & =\ln \left(x+\sqrt{1+x^{2}}\right)+\ln \left(-x+\sqrt{1+x^{2}}\right) \\
& =\ln \left[\left(x+\sqrt{1+x^{2}}\right) \cdot\left(-x+\sqrt{1+x^{2}}\right)\right] \\
& =\ln \left(1+x^{2}-x^{2}\right) \\
& =\ln 1 \\
& =0 .
\end{aligned}
\]

同时，注意到 \(f(x)\) 的定义域是 \(\mathbb{R}\) ，故 \(f(x)\) 是一个奇函数。\\
相信这个结果可能会出乎很多人的意料。这么复杂的一个函数，结果居然是一个奇函数？！但是，当我们从两个角度重新理解这个计算，就会发现一切顺利成章。

首先，我们观察到内层的代根号：\(x+\sqrt{1+x^{2}}\) ．为了去检验奇偶性，我们会将 \(x\) 替换成 \(-x\) ，这样内层变为了 \(-x+\sqrt{1+(-x)^{2}}=-x+\sqrt{1+x^{2}}\) 。这个结构引起了我们的注意，马上联想到了我们在杜杆 4 中讲到的共轭根式。很多同学会有困惑：这似乎并不是共轭根式。我们的共轭根式应该是一对形如 \(a+\sqrt{b}, a-\sqrt{b}\) 的表达式，根号的

符号相反；而这里是整式项 \(x\) 前面的符号相反。事实上：

\[
-x+\sqrt{1+x^{2}}=-\left(x-\sqrt{1+x^{2}}\right)
\]

它其实是 \(x-\sqrt{1+x^{2}}\) 的相反数，而 \(x+\sqrt{1+x^{2}}, x-\sqrt{1+x^{2}}\) 是一对共轭根式，所以，我们也应该把 \(x+\sqrt{1+x^{2}}\) 和 \(-x+\sqrt{1+x^{2}}\) 看作一对共轭根式。

共轭根式有什么好的性质呢？我们在讲第二章第 4 节中刻意强调过：它可以消掉根号，把无理式变为有理式。如果将这对共轭根式相乘，我们得到：

\[
\begin{aligned}
& \left(x+\sqrt{1+x^{2}}\right)\left(-x+\sqrt{1+x^{2}}\right) \\
= & \left(\sqrt{1+x^{2}}+x\right)\left(\sqrt{1+x^{2}}-x\right) \\
= & \left(\sqrt{1+x^{2}}\right)^{2}-x^{2} \\
= & 1+x^{2}-x^{2} \\
= & 1
\end{aligned}
\]

这样，我们很自然地希望利用这个性质，来帮我们化简函数。同时，我们也意识到：\(x+\sqrt{1+x^{2}},-x+\sqrt{1+x^{2}}\) 互为倒数，即

\[
-x+\sqrt{1+x^{2}}=\frac{1}{x+\sqrt{1+x^{2}}}
\]

换言之，对于函数 \(g(x)=x+\sqrt{1+x^{2}}\) 而言，虽然它不是奇函数，也不是偶函数，但是它满足 \(g(x) g(-x)=1\) 。\\
这也就是为什么，当我们外层函数是对数函数 \(y=\ln x\) 时，复合得到的 \(f(x)\) 最终是一个奇函数了。根据对数的性质： \(\ln a+\ln b=\ln a b\) ，有

\[
\begin{aligned}
f(x)+f(-x) & =\ln \left(x+\sqrt{1+x^{2}}\right)+\ln \left(-x+\sqrt{1+x^{2}}\right) \\
& =\ln \left[\left(x+\sqrt{1+x^{2}}\right) \cdot\left(-x+\sqrt{1+x^{2}}\right)\right] \\
& =1
\end{aligned}
\]

所以，看上去不知从何而来的对数函数，帮我们把加法变成了乘法，从而承接了共轭根式的性质，最终让这个函数成为了一个奇函数。

至此，我们已经非常清晰地分析出了这个复杂函数的结构，以及它为什么成为了一个奇函数。定义域对称自不必说，从解析式上看：\\
（1）共轭根式结构：对数内层函数 \(g(x)\) 和 \(g(-x)\) 互为共轭根式。同时，满足

\[
g(-x)=\frac{1}{g(x)}
\]

它们互为倒数；\\
（2）对数结构：对数将加法转变为乘法，使得 \(f(x)+f(-x)\) 在对数中体现为 \(g(x) 。 g(-x)\) 。进而用上了（1）的性质，使得 \(f(x)+f(-x)=0\) ．用数学语言叙述，就是：

\[
\begin{aligned}
f(x)+f(-x) & =\ln g(x)+\ln g(-x)=\ln g(x)+\ln \frac{1}{g(x)} \\
& =\ln \left[g(x) \cdot \frac{1}{g(x)}\right] \\
& =\ln 1 \\
& =0 .
\end{aligned}
\]

这里的另一个关键点是： \(\ln x+\ln \frac{1}{x}=1\) 。两个数互为倒数，则对数互为相反数。\\
这样，我们就从代数结构的角度，原理的角度，理解好了为什么这个复杂的函数是一个奇函数。仅仅是一个简单的奇函数，偶函数问题，里面就藏着这么多有用的代数结构。从这里，我们也可以窥见数学学习的要义：理解代数结构，理解原理，才可以以不变应万变，看破变化莫测的表面形式，抓住底层逻辑。

关于这个问题，下面有一段我想说的话。所有这些和计算基本法本身关系不大，但是非常重要的思想和认知内容，会以【不言不语】的评论形式编辑在讲义中。它们的重要性并不亚于知识本身。它们可以帮助你调整学习方法和对于学习的认知。

\section*{【不言不语】}
如果大家去参考市面上的其他教辅或者课程，可能会有很多地方把这个性质抽象总结成：

\[
\text { "当 } g(x) \cdot g(-x)=1 \text { 时, } f(x)=\ln g(x) \text { 是奇函数" }
\]

然后给这样一个函数 \(g(x)=x+\sqrt{1+x^{2}}\) 取了一个好听的名字，叫＂对数奇函数＂。什么意思呢？意思是＂取完对数后，就是奇函数了＂。很多同学看到这个名字和二级结论，就开始记忆并利用它们解题，刚好教材和课程的编写者会为你精挑细选若干可以直接用这个东西解决的题目，然后你就像切瓜砍菜一样把它们做出来了，觉得自己学会了一个很厉害的东西，以为这就是数学学习。

但是大家可以仔细思考一个问题：我们在处理奇函数，偶函数的问题的时候，真是这么思考的吗？不是！也不需要。我们是拿着这样不本质的结论去验证一般的函数吗？不是！我们只要根据定义直接去检验 \(f(x), f(-x)\)就好了。如果相等，就是偶函数；如果互为相反数，就是奇函数。我们需要去记忆刚才被总结的结论，和知道＂对数奇函数＂这个名字吗？不需要。因为我们非常深刻地理解了它背后的原理：共轭根式带来了倒数，对数将倒数转化为相反数，最终成为奇函数。这就是理解原理强大的力量：一切都显得理所当然，而且思路清晰。既不会突然冒出来一个天外飞仙一样的二级结论，也没有任何的记忆负担，思考问题稳准狠。最终，我们既知道是什么，也知道为什么，学起来自然毫不费力。

当然，很多同学会说：我也想这样理解原理啊，但是我能力不够，做不到啊！其实，他们之所以做不到，很大一部分原因是计算能力偏弱，无法完成过程中的代数变形。很多时候，大家都知道怎么算，但是真的开始算，就束手无策了。计算基本法就是为大家解决这个障碍：当计算能力提升时，你就有能力从定义出发，从原理出发，完成中间步骤，理解所有问题。定义和原理就像数学世界的地图，而计算就像是河里的船，脚上的鞋。只知道往哪里走远远不够，你还得能走过去，划过去，这样才能够真正理解数学。

\section*{4．1．2 单调性}
除了奇偶性，单调性也是非常重要的概念。我们除了可以从图像直接判断单调性，还会从定义直接验证。设函数 \(f(x)\) 的定义域为 \(D\) 。如果一个函数满足：对于任意 \(x_{1}, x_{2} \in D, x_{1}<x_{2}\) ，\\
（1）都有 \(f\left(x_{1}\right)<f\left(x_{2}\right)\) ，那么 \(f(x)\) 是一个单调递增函数；\\
（2）都有 \(f\left(x_{1}\right)>f\left(x_{2}\right)\) ，那么 \(f(x)\) 是一个单调递减函数。\\
那么，如何利用我们的代数结构来帮助我们判断单调性呢？我们这里主要考虑多项式和有理函数的例子。\\
我们以 \(f(x)=x^{3}+x\) 为例。对于 \(x_{1}, x_{2} \in \mathbb{R}, x_{1}<x_{2}\) ，我们需要比较 \(f\left(x_{1}\right)\) 和 \(f\left(x_{2}\right)\) 的大小。因此，考虑

\[
f\left(x_{1}\right)-f\left(x_{2}\right)=x_{1}^{3}+x_{1}-x_{2}^{3}-x_{2}
\]

我们只需要考虑它和 0 的大小即可。我们在之前讲过，对于整式方程，解方程本质上就是分解因式。而这里，我

们需要＂解不等式＂：考察 \(f\left(x_{1}\right)-f\left(x_{2}\right)\) 和 0 的大小，而非相等。但这本质上是一样的！因为：只有当我们将 \(f\left(x_{1}\right)-f\left(x_{2}\right)\) 写成若干个简单因式的乘积时，我们才可以通过逐一判断每个因式的正负，来决定 \(f\left(x_{1}\right)-f\left(x_{2}\right)\)的正负。

那么，问题来了：如何将 \(x_{1}^{3}+x_{1}-x_{2}^{3}-x_{2}\) 分解因式呢？我们来观察一下结构。\\
（1）这是一个二元三次多项式。它有多个变量，无法直接使用一元的代数结构。因此，可以考虑使用主元思想；\\
（2）它不是一个对称的多项式。当我们交换 \(x_{1}, x_{2}\) 的位置时，我们得到：

\[
x_{2}^{3}+x_{2}-x_{1}^{3}-x_{1} .
\]

它和原来的式子互为相反数。所以，对称换元无法直接使用。如果我们选主元（如 \(x_{1}\) ），会得到一个三次函数：

\[
y=x_{1}^{3}+x_{1}-x_{2}^{3}-x_{2}
\]

对于三次函数，我们知道：需要通过试根的方式分解因式。那么，试哪个根呢？\\
经验丰富的同学，马上就会意识到：\(x_{1}=x_{2}\) 时，\(g\left(x_{1}\right)=0\) 。此时

\[
y=x_{2}^{3}+x_{2}-x_{2}^{3}-x_{2}=0
\]

所以，它一定可以分解出 \(x_{1}-x_{2}\) 这个因式。这样一来，我们马上得到启发：只需要进行分组分解即可：

\[
\begin{aligned}
x_{1}^{3}+x_{1}-x_{2}^{3}-x_{2} & =x_{1}^{3}-x_{2}^{3}+x_{1}-x_{2} \\
& =\left(x_{1}-x_{2}\right)\left(x_{1}^{2}+x_{1} x_{2}+x_{2}^{2}+1\right)
\end{aligned}
\]

这样，我们知道了第一个括号 \(x_{1}-x_{2}<0\) 。只需要再去判断第二个括号的正负即可。我们发现：第二个括号是一个二次式。想要判断它的正负，只需要配方即可。我们依旧需要把它看成一个关于 \(x_{1}\) 的二次函数：

\[
x_{1}^{2}+x_{2} \cdot x_{1}+x_{2}^{2}+1=\left(x_{1}+\frac{x_{2}}{2}\right)^{2}+\frac{3 x_{2}^{2}}{4}+1 \geq 1>0
\]

所以，第二个括号是正数。这样一来，一正一负相乘为负。

\[
f\left(x_{1}\right)-f\left(x_{2}\right)=\left(x_{1}-x_{2}\right)\left(x_{1}^{2}+x_{1} x_{2}+x_{2}^{2}+1\right)<0
\]

故 \(f(x)\) 是一个单调递增函数。\\
反思上面的解题思路，我们做出来的关键，就是因为我们对于代数结构有着深刻的洞察。除此之外，而对于多项式而言，这里我们想要强调一个特别的性质。假设 \(f(x)\) 是一个多项式，那么形如 \(f\left(x_{1}\right)-f\left(x_{2}\right)\) 的式子一定可以分解出 \(x_{1}-x_{2}\) 这个因式。这是因为：我们如果选 \(x_{1}\) 为主元，那么 \(y=f\left(x_{1}\right)-f\left(x_{2}\right)\) 化为一个关于 \(x_{1}\)的多项式，满足：当 \(x_{1}=x_{2}\) 时，

\[
y=f\left(x_{2}\right)-f\left(x_{2}\right)=0
\]

故 \(x_{1}=x_{2}\) 是 \(y=0\) 的一个根。根据因式定理，\(y\) 必然有 \(x_{1}-x_{2}\) 这个因式。这样一来，我们如果再面对多项式函数的时候，可以上来直接＂无脑＂分解出 \(x_{1}-x_{2}\) 这个因式，再去处理后面的结构。

如果对它感觉依旧模糊，我们可以再举一个例子：若 \(f(x)=3 x^{4}+5 x^{2}+1\) ，那么

\[
\begin{aligned}
f\left(x_{1}\right)-f\left(x_{2}\right) & =3 x_{1}^{4}+5 x_{1}^{2}+1-\left(3 x_{2}^{4}+5 x_{2}^{2}+1\right) \\
& =3\left(x_{1}^{4}-x_{2}^{4}\right)+5\left(x_{1}^{2}-x_{2}^{2}\right) \\
& =\left(x_{1}-x_{2}\right)\left(3 x_{1}^{3}+3 x_{1}^{2} x_{2}+3 x_{1} x_{2}^{2}+3 x_{2}^{3}+5 x_{1}+5 x_{2}\right) \\
& =\left(x_{1}-x_{2}\right)\left(x_{1}+x_{2}\right)\left(3 x_{1}^{2}+3 x_{2}^{2}+5\right)
\end{aligned}
\]

这和我们的观察是相符的。同时，我们发现：多项式之所以有这个性质，很大一部分来自我们非常熟悉的公式：

\[
x^{n}-y^{n}=(x-y)\left(x^{n-1}+x^{n-2} y+x^{n-3} y^{2}+\cdots+x y^{n-2}+y^{n-1}\right)
\]

我们经过分组分解后，就会产生多个这样不同的 \(x^{n}-y^{n}\) ：如 \(x_{1}^{4}-x_{2}^{4}, x_{1}^{3}-x_{2}^{3}, x_{1}^{2}-x_{2}^{2}\) 等。因此，我们又从分解因式公式本身理解了这个现象。

回顾我们整体思路，我们顺利完成它主要有以下两个要点：\\
（1）解方程 \(=\) 因式分解，解不等式也要因式分解！\\
（2）多项式 \(f\left(x_{1}\right)-f\left(x_{2}\right)\) 必然可以分解出 \(x_{1}-x_{2}\) 。这和因式定理有关，也可以理解为多组 \(x_{1}^{n}-x_{2}^{n}\) 分解后得到的结果。这样一来，我们的整体思路就非常清楚了。再遇到这样的多项式函数，我们只需要选主元，分解因式，提 \(x_{1}-x_{2}\) ，就可以很快判断单调性了。

我们发现，仅仅是一个简单的判断三次函数的单调性，在实际的计算过程中，就用到了咱们之前讲过的非常多的杜杆：主元，对称，分解因式，二次式。

\section*{【不言不语】}
在这里，我不给大家总结任何＂标准化＂的步骤和思路。这是因为：首先，问题本身可能会有变化，所以没有一套＂一定有用＂的标准化思路，只有一套高度泛化的对于代数结构的理解。每次解决问题的时候，应该通过理解问题本身，洞察代数结构，去发展出新的思路，因地制宜地解决问题。

其次，标准化思路会导致学习的懒情。背诵是永远比理解要轻松的，因为人的大脑在理解事物时会下意识地采用＂更轻松＂的做法，这可以节约能量和空间做更多的事情。因此，很多同学会下意识地去背诵公式和定理。但是，这对于数学学习无异于饮鸠止渴：背诵公式忽略了数学对象本身的特点，忽略了数学对象之间的联系，忽略了对于原理的理解。

很多同学会反驳：我学数学就是为了做题，拿一个好的分数，难道不是记住公式，＂会用＂就行了嘛？其实，难就难在＂会用＂：只有背诵，而不理解原理，依靠＂刷题＂熟悉套路，是非常低效的＂学会使用＂的方法。它既会导致大量的无意义重复劳动，也会让最终的学习成果过拟合：我只会处理这些公式和结论能用的题，一旦超出这个范围，就不会了。

但如果，我们通过理解数学对象本身，理解它们的特性，原理，和它们之间的联系学会了数学，就不会出现这个问题。即使是新的情景中，我们的概念没有变，函数依旧是函数，依靠理解就可以轻松解决问题。因此，希望大家可以通过模仿我的思路和理解，尽可能习得并发展出自己对于数学的认知。

多项式我们可以提 \(x_{1}-x_{2}\) 的性质，那么对于有理函数呢？有理函数指的是形如

\[
\frac{p(x)}{q(x)}
\]

的函数，其中 \(p(x), q(x)\) 都是关于 \(x\) 的多项式。我们 \(f(x)=\frac{x}{x^{2}+1}, x>0\) 为例。当 \(x_{1}<x_{2}\) 时，我们有：

\[
\begin{aligned}
f\left(x_{1}\right)-f\left(x_{2}\right) & =\frac{x_{1}}{x_{1}^{2}+1}-\frac{x_{2}}{x_{2}^{2}+1} \\
& =\frac{\left[x_{1}\left(x_{2}^{2}+1\right)-x_{2}\left(x_{1}^{2}+1\right)\right]}{\left(x_{1}^{2}+1\right)\left(x_{2}^{2}+1\right)} \\
& =\frac{\left(x_{1} x_{2}^{2}+x_{1}-x_{2} x_{1}^{2}-x_{2}\right)}{\left(x_{1}^{2}+1\right)\left(x_{2}^{2}+1\right)}
\end{aligned}
\]

我们发现，分子 \(y=x_{1} x_{2}^{2}+x_{1}-x_{2} x_{1}^{2}-x_{2}\) 是一个多项式。我们如果选 \(x_{1}\) 为主元，它可以写成一个关于 \(x_{1}\) 的二

次函数：

\[
y=-x_{2} \cdot x_{1}^{2}+\left(x_{2}^{2}+1\right) x_{1}-x_{2}
\]

根据刚才的经验，我们猜测 \(x_{1}=x_{2}\) 是方程的一个根。事实上，这是成立的，此时

\[
y=-x_{2}^{3}+x_{2}^{3}+x_{2}-x_{2}=0
\]

所以，它也可以分解出 \(x_{1}-x_{2}\) 这个因式。因此，

\[
y=\left(x_{1}-x_{2}\right)\left(-x_{2} x_{1}+1\right)
\]

故只需要讨论 \(-x_{1} x_{2}+1\) 的大小即可。因此，\(f(x)\) 的单调性如下：\\
（1）当 \(x>1\) 时，\(f(x)\) 单调递减。对任意 \(1<x_{1}<x_{2}\) ，有 \(-x_{2} x_{1}+1<0\) ，

\[
f\left(x_{1}\right)-f\left(x_{2}\right)=\frac{\left(x_{1}-x_{2}\right)\left(-x_{2} x_{1}+1\right)}{\left(x_{1}^{2}+1\right)\left(x_{2}^{2}+1\right)}>0
\]

故 \(f(x)\) 单调递减；\\
（2）当 \(0<x<1\) 时，\(f(x)\) 单调递增。对任意 \(0<x_{1}<x_{2}<1\) ，有 \(-x_{1} x_{2}+1>0\) ，

\[
f\left(x_{1}\right)-f\left(x_{2}\right)=\frac{\left(x_{1}-x_{2}\right)\left(-x_{2} x_{1}+1\right)}{\left(x_{1}^{2}+1\right)\left(x_{2}^{2}+1\right)}<0
\]

故 \(f(x)\) 单调递增。\\
因此，我们发现：即使是有理分式函数，也有这个性质！我们都可以在检验单调性的时候，提一个 \(x_{1}-x_{2}\)后再考虑剩下的部分。我们通过刚才的例子进行了检验，也可以直接通过运算得到这个性质：

\[
\begin{aligned}
f\left(x_{1}\right)-f\left(x_{2}\right) & =\frac{p\left(x_{1}\right)}{q\left(x_{1}\right)}-\frac{p\left(x_{2}\right)}{q\left(x_{2}\right)} \\
& =\frac{p\left(x_{1}\right) q\left(x_{2}\right)-p\left(x_{2}\right) q\left(x_{1}\right)}{q\left(x_{1}\right) q\left(x_{2}\right)}
\end{aligned}
\]

我们的分子 \(y=p\left(x_{1}\right) q\left(x_{2}\right)-p\left(x_{2}\right) q\left(x_{1}\right)\) 虽然不能写成 \(g\left(x_{1}\right)-g\left(x_{2}\right)\) 的形式，但是我们也可以通过选主元的方法，将它理解成关于 \(x_{1}\) 的多项式。而当 \(x_{1}=x_{2}\) 时，有

\[
y=p\left(x_{2}\right) q\left(x_{2}\right)-p\left(x_{2}\right) q\left(x_{2}\right)=0
\]

所以，根据因式定理，分子可以提出一个 \(x_{1}-x_{2}\) 。

\section*{4．1．3 值域}
有关函数值域的问题，我们主要讨论分子，分母都是一次函数或者二次函数的情况，且至少有一个是二次函数。例如，如下函数都在我们的讨论范围之内。

\[
\frac{x}{1+x^{2}}, \frac{x^{2}+1}{x+1}, \frac{x^{2}+2 x+1}{x^{2}+x+1}
\]

很多同学在学习这部分内容的时候，是严格依照一个思维导图去理解的。例如，当如遇到一次比二次的情况，选择将分子的一次式进行换元，化成对勾函数的形式。然而，在实际计算过程中，我们只需要抓住两个大的基本思路即可。\\
（1）集中：将变量集中到分子，分母中的一个；\\
（2）化归：将不熟悉的函数化为熟悉的函数：二次函数或对勾函数等。\\
为什么这么考虑问题呢？首先，虽然分子，分母都是我们熟悉的函数，但是它们在同时变化的时候，我们无法掌握它们的变化趋势。将所有的变量集中到其中一侧才可以有效化简问题，所以需要第一步：集中；其次，只有将不熟悉的函数化为熟悉的，才能使用我们已知的工具，所以需要第二步：化归。

我们来仔细看，这两个思路是如何统一处理所有这类问题的。对于第一个函数

\[
f_{1}(x)=\frac{x}{1+x^{2}}
\]

我们发现：分子是 \(x\) ，分母是 \(x^{2}+1\) ，将分子，分母同时除以分子可以有效化简。

\[
f_{1}(x)=\frac{1}{x+\frac{1}{x}}
\]

这样，我们同时完成了两步操作。既将所有的变量集中到了分母上，又同时将不熟悉的函数化为了熟悉的对勾函数。为什么选择除以分子？因为同时除以分母并不能化简问题：

\[
f_{1}(x)=\frac{\frac{x}{x^{2}+1}}{\frac{x^{2}+1}{x^{2}+1}}=\frac{\frac{x}{x^{2}+1}}{1}=\frac{x}{x^{2}+1}
\]

而对于第二个函数 \(f_{2}(x)=\frac{x^{2}+1}{x+1}\) ，无论是分子还是分母，都是多项式。因此，我们可以令 \(t=x+1\) ，得到

\[
f_{2}(x)=\frac{t^{2}-2 t+2}{t}=t+\frac{2}{t}-2
\]

又化成了对勾函数。之所以想到令 \(t=x+1\) ，就是因为我们希望得到一个单项式，将分母变得简单。这里，我们再一次提醒大家关于分子，分母哲学：分子，分母中，尽量把所有复杂的内容留给分子，而让分母尽可能简单。

上面两个例子中，分子，分母中都有一个是一次式，这样可以将函数化为对勾函数或它的倒数。那如果分子，分母都是二次式呢？我们能否将它化归为刚才的一次比二次的问题？也是可以的：只需要分离常数即可。例如：

\[
f_{3}(x)=\frac{x^{2}+2 x+1}{x^{2}+x+1}=1+\frac{x}{x^{2}+x+1}=1+\frac{1}{x+\frac{1}{x}+1}
\]

我们再一次化为对勾函数的倒数。这也是化归思想的自然应用：将分子的二次式化为一次式。\\
那除了这么处理，有没有别的方法呢？也有。我们观察到分子 \(x^{2}+2 x+1\) 刚好是一个完全平方式：\(x^{2}+2 x+1=\) \((x+1)^{2}\) ．这样，也可以化作一个一次式。令 \(t=x+1\) ，则

\[
f_{3}(x)=\frac{t^{2}}{t^{2}-t+1}
\]

如果我们直接将 \(t^{2}\) 除到分母上，可以得到

\[
f_{3}(x)=\frac{1}{1-\frac{1}{t}+\frac{1}{t^{2}}}
\]

分母变成了一个关于 \(\frac{1}{t}\) 的二次函数：令 \(z=\frac{1}{t}\) ，那么

\[
1-\frac{1}{t}+\frac{1}{t^{2}}=z^{2}-z+1
\]

这个二次函数的最值就决定了 \(f_{3}(x)\) 的最值。\\
学到这里，我们发现：以上的每一步计算和思考，都非常自然。

4．1．4 XXXX

\section*{4.2 不等式}
\section*{4．2．1 基本不等式（均值不等式）}
我们在教材里已经学习过了均值不等式。在这里，我们重新复习一下有关它的基本性质。设 \(a, b \geq 0\) 是非负实数，则

\[
a+b \geq 2 \sqrt{a b}
\]

我们之前所学的众多代数结构，都可以在它的身上得到体现。\\
（1）对称性：这是一个关于 \(a, b\) 对称的不等式。事实上，令 \(m=a+b, n=\sqrt{a b}\) ，那么它可以写作

\[
m \geq 2 \sqrt{n}
\]

注意，这里要求 \(a, b \geq 0\) ，我们才可以写作这个形式。如果 \(a, b\) 只是实数，则一般改写为

\[
a^{2}+b^{2} \geq 2 a b
\]

这等价于：\((a+b)^{2} \geq 4 a b\) ，即 \(m^{2} \geq 4 n\) ．\\
（2）齐次性：均值不等式的两侧都是一次齐次。利用这个特点，我们可以利用比例来表示改写这个不等式。令 \(t=\frac{a}{b}\) ，在不等式两侧同时除以 \(\sqrt{a b}\) ，可得

\[
\begin{aligned}
& \sqrt{\frac{a}{b}}+\sqrt{\frac{b}{a}} \geq 2 \\
& \sqrt{t}+\frac{1}{\sqrt{t}} \geq 2
\end{aligned}
\]

这就是我们熟悉的对勾函数。所以，对勾函数最小值的性质和均值不等式是等价的。\\
（3）二次函数：均值不等式本身，也可以看作是配方的结果。事实上，

\[
a^{2}+b^{2}-2 a b=(a-b)^{2} \geq 0
\]

因此，它也和平方数非负是等价的。我们从这里也能看出均值不等式＂基本＂的地位。\\
均值不等式除了上面两种常见形式，还有非常多其他的写法：\\
（1）\((a+b)^{2} \geq 4 a b\) 。这等价于

\[
\begin{aligned}
& a^{2}+2 a b+b^{2} \geq 4 a b \\
& a^{2}+b^{2} \geq 2 a b
\end{aligned}
\]

因此也是均值不等式。\\
（2）\(a^{2}+b^{2} \geq \frac{(a+b)^{2}}{2}\) 。这等价于

\[
\begin{aligned}
& 2 a^{2}+2 b^{2} \geq a^{2}+2 a b+b^{2} \\
& a^{2}+b^{2} \geq 2 a b
\end{aligned}
\]

有时也将它写成根号的形式：

\[
\sqrt{\frac{a^{2}+b^{2}}{2}} \geq a+b
\]

很多同学会有这样的疑问：均值不等式的变形如此之多，我们如何记忆呢？事实上，它们完全等价，不需要记忆。我们重新来理解它们。

我们可以把上面的不等式写作如下的不等式串：

\[
a^{2}+b^{2} \geq \frac{(a+b)^{2}}{2} \geq 2 a b
\]

中间的项 \(\frac{(a+b)^{2}}{2}\) 可以写作：\(\frac{a^{2}+b^{2}+2 a b}{2}\) 。事实上，它刚好是 \(a^{2}+b^{2}\) 和 \(2 a b\) 的平均值。我们根据均值不等式知道： \(a^{2}+b^{2} \geq 2 a b\) ，那么上面的不等式串就是显然的：

较大数 \(\geq\) 平均数 \(\geq\) 较小数．\\
除了利用平均数，我们还可以写成加权平均的形式。例如：当 \(a^{2}+b^{2}\) 的权重是 \(\frac{1}{3}, ~ 2 a b\) 的权重是 \({ }_{3}^{2}\) 时，可以得到：

\[
\begin{aligned}
& a^{2}+b^{2} \geq \frac{1}{3}\left(a^{2}+b^{2}\right)+\frac{2}{3} \cdot 2 a b \geq 2 a b \\
& a^{2}+b^{2} \geq \frac{a^{2}+b^{2}+4 a b}{3} \geq 2 a b
\end{aligned}
\]

这也是显然的，因为：

\[
\text { 较大数 } \geq \text { 加权平均数 } \geq \text { 较小数. }
\]

那不同的加权平均之间，如何比较呢？那就看谁＂较大数＂的权重大了。例如，

\[
\frac{a^{2}+b^{2}+2 a b}{2} \geq \frac{a^{2}+b^{2}+4 a b}{3}
\]

因为前者 \(a^{2}+b^{2}\) 的权重是 \(\frac{1}{2}\) ，而后者是 \(\frac{1}{3}\) ，所以前者大。这样，我们就可以得到一串更长的不等式：

\[
a^{2}+b^{2} \geq \frac{a^{2}+b^{2}+2 a b}{2} \geq \frac{a^{2}+b^{2}+4 a b}{3} \geq 2 a b
\]

事实上，我们可以把 \(a^{2}+b^{2}\) 看成 \(a^{2}+b^{2}\) 权重为 1 的加权平均， \(2 a b\) 看成 \(a^{2}+b^{2}\) 权重为 0 的加权平均，那么上面的不等式可以理解为：

\[
\text { 权重为 } 1 \geq \text { 权重为 } \frac{1}{2} \geq \text { 权重为 } \frac{1}{3} \geq \text { 权重为 } 0 \text {. }
\]

所以，我们可以把所有这些不等式理解成一个关于权重 \(\lambda \in[0,1]\) 的函数：

\[
f(\lambda)=\lambda\left(a^{2}+b^{2}\right)+(1-\lambda) \cdot 2 a b
\]

对于给定的 \(a \neq b, f(\lambda)\) 是一个关于 \(\lambda\) 单调递增的函数。除了我们直观地从加权平均数来理解，严格的证明也非常简单：当 \(\lambda>\mu\) 时，

\[
\begin{aligned}
f(\lambda)-f(\mu) & =(\lambda-\mu)\left(a^{2}+b^{2}\right)+(1-\lambda-1+\mu) \cdot 2 a b \\
& =(\lambda-\mu)\left(a^{2}+b^{2}-2 a b\right) \\
& =(\lambda-\mu)(a-b)^{2}>0 .
\end{aligned}
\]

因此 \(f(\lambda)\) 单调递增。

\section*{实战 2 不等式}
因此，在理解这些均值不等式的变形时，不要把它们看作没有道理和孤立的新形式。它们很多只是经过适当的加权平均得到的，只要把正确的＂权重＂找到，就可以顺利理解这些不等式。

\section*{4．2．2 利用对称性处理不等式}
我们在第一部分的总结中提到，利用对称性来消元处理不等式的方法将会在这部分内容中详细讨论。这里我们举一个具体的例子来说明。

如果像要求 \(a^{4}+b^{4}-2 a b\) 的最小值，如何求解？我们观察发现，这是一个对称的二元不等式，考虑使用对称换元。令 \(m=a+b, n=a b\) ．那么

\[
a^{4}+b^{4}=\left(a^{2}+b^{2}\right)^{2}-2 a^{2} b^{2}=\left(m^{2}-2 n\right)^{2}-2 n^{2}=m^{4}-4 n m^{2}+2 n^{2}
\]

那么原式可以写作：

\[
y=m^{4}-4 n m^{2}+2 n^{2}-2 n
\]

这样的话，我们得到了一个关于 \(m, n\) 的式子。如果我们选不同的变量为主元，会得到不同的结果：\\
（1）选 \(m^{2}\) 为主元，那么 \(y\) 是一个关于 \(m^{2}\) 的二次函数：

\[
y=m^{4}-4 n m^{2}+2 n^{2}-2 n=\left(m^{2}\right)^{2}-4 n \cdot m^{2}+2 n^{2}-2 n
\]

（2）选 \(n\) 为主元，那么 \(y\) 是一个关于 \(n\) 的二次函数：

\[
y=2 n^{2}-\left(4 m^{2}+2\right) n+m^{4}
\]

在第一种情况中，二次函数更简单，且关于 \(n\) 也是二次函数，所以我们考虑以 \(m^{2}\) 为主元。\\
这里就要用到我们心心念念的均值不等式了。根据均值不等式，有 \(m^{2} \geq 4 n\) 。因此，

\[
y=m^{2}\left(m^{2}-4 n\right)+2 n^{2}-2 n \geq 2 n^{2}-2 n
\]

我们成功地将 \(m\) 消掉，得到了一个只关于 \(n\) 的二次函数。这就是对称性和均值不等式在消元中的巨大作用。\\
最后，

\[
2 n^{2}-2 n=2\left(n-\frac{1}{2}\right)^{2}-\frac{1}{2} \geq-\frac{1}{2}
\]

故 \(y\) 的最小值为 \(-\frac{1}{2}\) ．取等条件为：\(n=\frac{1}{2} m^{2}=4 n\) ，即 \(a=b=\frac{\sqrt{2}}{2}\) 时取等。这样一来，我们就解决了一个四次表达式的最值问题。

总结一下，对于具有对称性的二元表达式，可以进行如下处理：\\
（1）对称换元，将表达式写成 \(m, n\) 的形式；（2）选 \(m\) 或 \(n\) 为主元；\\
（3）利用均值不等式：\(m^{2} \geq 4 n\)（或当 \(a, b>0\) 时，\(m \geq 2 \sqrt{n}\) ）消元；\\
（4）利用一元函数最值，求原表达式最值。\\
很多同学不能完全理解对称性的意义，以及为什么要尽可能保留对称性。这里第三步给了大家答案：只有利用对称性，我们才可以使用均值不等式（这是对称的不等式），才可以消元。所以，对称结构是在进行代数变形时需要尽量保留的。

最后，我们从自由度的角度分析一下这个题。为了得到最小值，我们需要确定取最小值时 \(a, b\) 的值。相当于，我们要得到关于 \(a, b\) 的两个条件。每一次放缩的取等条件，都提供一个限制条件。因此，需要两个独立的放缩（取等条件），才可以确定这个式子的最小值。在这道题里，两次独立的放缩是：\\
（1）均值不等式：

\[
a^{2}+b^{2} \geq 2 a b
\]

（2）二次函数：

\[
2 n^{2}-2 n=2\left(n-\frac{1}{2}\right)^{2}-\frac{1}{2} \geq-\frac{1}{2}
\]

有了这个理解，其实我们可以直接使用均值不等式：

\[
a^{4}+b^{4}-2 a b \geq 2 a^{2} b^{2}-2 a b=2\left(a b-\frac{1}{2}\right)^{2}-\frac{1}{2} \geq-\frac{1}{2}
\]

也可以得到正确答案。很多同学可能会在其他参考书上看到这样的解答，但是它没有说明背后的动机。事实上，除了＂感觉可以使用均值＂这样的做题经验，这背后还藏着我们刚才一路走来的所有分析：从问题本身的代数结构出发，通过使用我们的一个个杜杆，最终把复杂的问题化为简单的问题（均值不等式，二次函数配方等）。

\section*{4．2．3 XXXX}
\section*{4.3 三角函数，解三角形}
\section*{4．3．1 三角恒等变换}
4．3．1．1 和角公式，差角公式\\
我们在学习三角函数时，一个重要的难点就是三角恒等变换，如正弦，余弦的和角，差角公式：

\[
\begin{aligned}
& \sin (\alpha+\beta)=\sin \alpha \cos \beta+\sin \beta \cos \alpha, \\
& \sin (\alpha-\beta)=\sin \alpha \cos \beta-\sin \beta \cos \alpha, \\
& \cos (\alpha+\beta)=\cos \alpha \cos \beta-\sin \alpha \sin \beta, \\
& \cos (\alpha-\beta)=\cos \alpha \cos \beta+\sin \alpha \sin \beta
\end{aligned}
\]

虽然我们最终都会记住这四个公式，但是我们能不能从正确的角度来理解它们呢？尤其是从代数结构的角度来理解它们？\\
（1）正弦的和角，差角公式\\
事实上，我们可以非常自然地利用换元和对称来理解这几个三角函数公式。首先观察（1），（2）两个公式。对于（1），我们发现，公式的左侧是一个关于 \(\alpha, \beta\) 的对称式：当 \(\alpha, \beta\) 互换时，\(\alpha+\beta\) 保持不变， \(\sin (\alpha+\beta)\) 同理。因此，我们可以预期，等号右边应该也是关于 \(\alpha, \beta\) 的对称式。我们来验证一下：将 \(\alpha, \beta\) 交换，得

\[
\sin \beta \cos \alpha+\sin \alpha \cos \beta=\sin \alpha \cos \beta+\sin \beta \cos \alpha
\]

和原式相同。而对于公式（2），如果交换了 \(\alpha, \beta\) ，公式的左侧变为了原来的相反数：

\[
\sin (\beta-\alpha)=-\sin (\alpha-\beta)
\]

这是因为正弦函数是奇函数。因此，我们预期等号右边经过互换之后，也会变为原来的相反数：

\[
\sin \beta \cos \alpha-\sin \alpha \cos \beta=-(\sin \alpha \cos \beta-\sin \beta \cos \alpha)
\]

像这样，交换 \(\alpha, \beta\) 后，得到相反数的式子，称为反对称（skew－symmetric，antisymmetric）式。\\
除了通过对称性理解这两个公式，我们还应该注意到，它们本质上是一个公式。很多同学在理解这两个公式时，注意到了它们只差一个正负号，因此会把它们写作：

\[
\sin (\alpha \pm \beta)=\sin \alpha \cos \beta \pm \sin \beta \cos \alpha
\]

这个写法固然没有任何问题，但是我们也要理解：为什么加号成立的情况下，把中间的符号换成减号，也成立

呢？这本质上也和正弦函数的奇性，余弦函数的偶性相关。当我们把（1）中的所有 \(\beta\) 换成 \(-\beta\) 后，可以得到：

\[
\sin (\alpha+(-\beta))=\sin \alpha \cos (-\beta)+\sin (-\beta) \cos \alpha
\]

等式左侧显然变成了 \(\sin (\alpha-\beta)\) ，也就是（2）的左侧；而由于

\[
\cos (-\beta)=\cos \beta, \sin (-\beta)=-\sin \beta
\]

等式右侧变为了

\[
\sin \alpha \cos \beta-\sin \beta \cos \alpha
\]

所以，由（1）我们可以直接得到

\[
\sin (\alpha-\beta)=\sin \alpha \cos \beta-\sin \beta \cos \alpha
\]

这就是公式（2）．所以，并不是因为我们可以直接将加号换成减号，所以得到了公式（2），而是这两个公式本质上就是一个公式：将（1）中的 \(\beta\) 换为 \(-\beta\) ，就可以得到（2）．\\
（2）余弦的和角，差角公式\\
和正弦一样，我们可以通过同样的思路来理解这两个公式。显然，（3）的左右两侧都是对称的式子：交换 \(\alpha, \beta, \cos (\alpha+\beta)\) 保持不变，

\[
\cos \beta \cos \alpha-\sin \beta \sin \alpha=\cos \alpha \cos \beta-\sin \alpha \sin \beta
\]

它们依旧保持不变。对于（4），虽然 \(\alpha-\beta\) 是反对称的，但是由于余弦函数是偶函数，所以交换 \(\alpha, \beta\) 后，它依旧保持不变：

\[
\cos (\beta-\alpha)=\cos (\alpha-\beta)
\]

因此，这也是一个对称式。右侧确实是对称的：

\[
\cos \beta \cos \alpha+\sin \beta \sin \alpha=\cos \alpha \cos \beta+\sin \alpha \sin \beta
\]

将 \(\beta\) 换为 \(-\beta\) ，也可以发现（3），（4）两个公式等价：

\[
\begin{aligned}
\cos (\alpha+(-\beta)) & =\cos \alpha \cos (-\beta)-\sin \alpha \sin (-\beta) \\
& =\cos \alpha \cos \beta+\sin \alpha \sin \beta
\end{aligned}
\]

因此，（3），（4）本质上也是一个公式。\\
（3）正弦和余弦\\
我们发现正弦，余弦内部两个公式等价，这并不意外；但事实上，它们之间也是等价的。这一点可以简单地从诱导公式得到。我们从（4）出发，推（1）．将（4）中 \(\alpha\) 换为 \(\frac{\pi}{2}-\alpha\) ，得到

\[
\begin{aligned}
\cos \left(\frac{\pi}{2}-\alpha-\beta\right) & =\cos \left(\frac{\pi}{2}-\alpha\right) \cos \beta+\sin \left(\frac{\pi}{2}-\alpha\right) \sin \beta \\
& =\sin \alpha \cos \beta+\sin \beta \cos \alpha
\end{aligned}
\]

而 \(\cos \left(\frac{\pi}{2}-\alpha-\beta\right)=\sin (\alpha+\beta)\) ，所以

\[
\sin (\alpha+\beta)=\sin \alpha \cos \beta+\sin \beta \cos \alpha
\]

因此（4）可以推（1）类似的方法可以从（1）推（4），因此，这四个三角函数公式是完全等价的。

一个自然的问题是，能不能从别的角度直接推出其中一个三角函数公式呢？最简单的方法，是通过向量的内积。考虑向量 \(\vec{a}=(\cos \alpha, \sin \alpha), \vec{b}=(\cos \beta, \sin \beta)\) 。显然，\(\vec{a}\) 和 \(x\) 轴的夹角为 \(\alpha, \vec{b}\) 和 \(x\) 轴的夹角为 \(\beta\)（考虑单位圆），所以 \(\vec{a}, \vec{b}\) 的夹角是 \(\alpha-\beta\)（或 \(\beta-\alpha\) ）。根据向量内积的定义：

\[
\cos (\alpha-\beta)=\frac{\vec{a} \cdot \vec{b}}{|\vec{a}| \cdot|\vec{b}|}=\frac{\cos \alpha \cos \beta+\sin \alpha \sin \beta}{1 \cdot 1}=\cos \alpha \cos \beta+\sin \alpha \sin \beta
\]

总结一下，虽然我们学习了四个很像的三角函数公式，但是它们本质上：\\
（1）（1），（3），（4）是对称的，（2）是反对称的；\\
（2）四个都是等价的：（1），（2）差一个负号，（3），（4）差一个负号，正弦和余弦可利用诱导公式转化。\\
（3）可以从向量内积理解：（4）是两个单位向量的内积。

\section*{4．3．1．2 二倍角公式}
\section*{2．二倍角公式}
我们在学习完和角公式，差角公式后，一个自然的应用就是二倍角公式。在（1），（3）中，令 \(\beta=\alpha\) ，可得

\[
\begin{aligned}
& \sin 2 \alpha=2 \sin \alpha \cos \alpha \\
& \cos 2 \alpha=\cos ^{2} \alpha-\sin ^{2} \alpha
\end{aligned}
\]

我们发现：左侧都是二倍角的三角函数，而右侧都是关于 \(\sin \alpha, \cos \alpha\) 的二次齐次式。所以，二倍角的正弦，余弦，都可以写成一倍角的二次齐次式。

三倍角甚至是多倍角是否也有这样的规律呢？答案是肯定的。我们以三倍角为例：

\[
\begin{aligned}
\sin 3 \alpha=\sin (2 \alpha+\alpha) & =\sin 2 \alpha \cos \alpha+\sin \alpha \cos 2 \alpha \\
& =2 \sin \alpha \cos \alpha \cos \alpha+\sin \alpha\left(\cos ^{2} \alpha-\sin ^{2} \alpha\right) \\
& =3 \sin \alpha \cos ^{2} \alpha-\sin ^{3} \alpha \\
\cos 3 \alpha=\cos (2 \alpha+\alpha) & =\cos 2 \alpha \cos \alpha-\sin 2 \alpha \sin \alpha \\
& =\left(\cos ^{2} \alpha-\sin ^{2} \alpha\right) \cos \alpha-2 \sin ^{2} \alpha \cos \alpha \\
& =\cos ^{3} \alpha-3 \sin ^{2} \alpha \cos \alpha
\end{aligned}
\]

三倍角的正弦，余弦，确实都可以写作一倍角的三次齐次式。\\
很多同学记忆的三倍角公式和上述公式并不完全相同，而是

\[
\sin 3 \alpha=3 \sin \alpha-4 \sin ^{3} \alpha, \cos 3 \alpha=4 \cos ^{3} \alpha-3 \cos \alpha
\]

这个公式中，右侧不再是三次齐次式了，而是关于 \(\sin \alpha\) 或者 \(\cos \alpha\) 的三次多项式。事实上，这两个公式和上面的公式是完全一样的，差别在于多使用了一次勾股定理 \(\sin ^{2} \alpha+\cos ^{2} \alpha=1\) ：

\[
\begin{aligned}
\sin 3 \alpha & =3 \sin \alpha \cos ^{2} \alpha-\sin ^{3} \alpha=3 \sin \alpha\left(1-\sin ^{2} \alpha\right)-\sin ^{3} \alpha \\
& =3 \sin \alpha-4 \sin ^{3} \alpha \\
\cos 3 \alpha & =\cos ^{3} \alpha-3 \sin ^{2} \alpha \cos \alpha=\cos ^{3} \alpha-3\left(1-\cos ^{2} \alpha\right) \cos \alpha \\
& =4 \cos ^{3} \alpha-3 \cos \alpha
\end{aligned}
\]

所以，如果只是死记硬背公式，就会看不见真正重要的结构，导致使用时没有方向。\\
一般地，对于任意的正整数 \(n, n\) 倍角的正弦，余弦，可以写作一倍角的 \(n\) 次齐次式。如果读者已经学习过数学归纳法，我们可以直接利用它来证明。假设我们已经知道了 \((n-1)\) 倍角可以写作 \(\sin \alpha, \cos \alpha\) 的 \((n-1)\) 次

齐次式，那么

\[
\begin{aligned}
\sin n \alpha & =\sin (\alpha+(n-1) \alpha) \\
\cos n \alpha & =\sin \alpha \cos [(n-1) \alpha]+\cos \alpha \sin [(n-1) \alpha] .
\end{aligned}
\]

由于 \(\cos [(n-1) \alpha], \sin [(n-1) \alpha]\) 已经是 \((n-1)\) 次齐次式，那么

\[
\begin{aligned}
& \sin \alpha \cos [(n-1) \alpha]+\cos \alpha \sin [(n-1) \alpha] \\
& \cos \alpha \cos [(n-1) \alpha]-\sin \alpha \sin [(n-1) \alpha]
\end{aligned}
\]

必然就是 \(n\) 次齐次式了。根据数学归纳法原理，我们的结论是正确的。\\
在线下教学中，我们经常会听到类似于＂降幂扩角＂之类的口诀。它们所揭示的，就是多倍角三角函数齐次结构的特征：一倍角的二次齐次，就是二倍角的一次齐次。在实际应用中，我们可以根据具体情况选择使用一倍角或多倍角。

\section*{4．3．2 正弦定理，余弦定理}
\section*{4．3．2．1 正弦定理}
正弦定理的叙述为：在三角形 \(A B C\) 中，设 \(R\) 是三角形 \(A B C\) 外接圆的半径，则

\[
\frac{a}{\sin A}=\frac{b}{\sin B}=\frac{c}{\sin C}=2 R
\]

如果不考虑外接圆半径的部分，前面三部分同样具有齐次结构：

\[
\frac{a}{\sin A}=\frac{b}{\sin B}=\frac{c}{\sin C}
\]

它们都是关于 \(a, b, c\) 的一次齐次式，也是关于 \(\sin A, \sin B, \sin C\) 的负一次齐次式。因此，在实际使用中，可以在 \(a, b, c\) 和 \(\sin A, \sin B, \sin C\) 之间进行转化。例如，如果三角形 \(A B C\) 中，有

\[
a \sin C=b \sin A
\]

那么我们有两种思路进行化简：\\
（1）化为角：观察到它们都是关于边 \(a, b\) 的一次齐次，可以利用正弦定理化为关于角的式子。根据正弦定理，有

\[
a=2 R \cdot \sin A, b=2 R \cdot \sin B
\]

所以

\[
\begin{aligned}
2 R \cdot \sin A \sin C & =2 R \cdot \sin B \sin A \\
\sin C & =\sin B
\end{aligned}
\]

因此 \(B=C\) ，等价地，\(b=c\) 。很多老师在讲解这部分内容时直接将已知条件中的 \(a, b\) 化成 \(\sin A, \sin B\) ，而不解释原因。其原因在于正弦定理的齐次结构。\\
（2）化为边：观察到它们都是关于正弦 \(\sin A, \sin C\) 的一次齐次，可以利用正弦定理化为关于边的式子。根据正弦定理，有

\[
\sin C=\frac{c}{2 R}, \sin A=\frac{a}{2 R} .
\]

因此，

\[
a \cdot \frac{c}{2 R}=b \cdot \frac{a}{2 R}
\]

化简得 \(b=c\) ，等价地，\(B=C\) 。\\
因此，在实际应用中，如果观察到了关于 \(a, b, c\) 或 \(\sin A, \sin B, \sin C\) 的齐次式，而可以考虑使用正弦定理，将角都化为边或者将边都化为角。绝大多数情况下，边角夹杂是不好处理的，统一化为角或边更有利于化简。

\section*{4．3．2．2 余弦定理}
余弦定理的叙述为：在三角形 \(A B C\) 中，成立

\[
\begin{aligned}
& a^{2}=b^{2}+c^{2}-2 b c \cos A \\
& b^{2}=a^{2}+c^{2}-2 a c \cos B \\
& c^{2}=a^{2}+b^{2}-2 a b \cos C
\end{aligned}
\]

我们以第一个式子为例，观察它具有哪些代数结构。\\
（1）齐次性：左右都是关于 \(a, b, c\) 的二次齐次式；\\
（2）对称性：第一个是关于 \(b, c\) 的对称式。所以可以将它写作 \(b+c, b c\) 的表达式：

\[
a^{2}=(b+c)^{2}-2 b c(1+\cos A)
\]

基于以上两个结构，当我们遇到齐次式（尤其是二次齐次）和对称式时，可以考虑用余弦定理去处理。\\
余弦定理还会写作

\[
\cos A=\frac{b^{2}+c^{2}-a^{2}}{2 b c}
\]

右侧是一个关于边的零次齐次式。这和正弦定理是一样的：

\[
\sin A=\frac{a}{2 R}
\]

都是关于边长的零次齐次式。这个特点我们在齐次式这个杜杆中提到过。由于正弦，余弦是关于角度的函数，所以当边长等比例扩大或缩小时，新图形和旧图形相似，角度不变。因此，正弦，余弦如果可以写作关于边的函数，必然是零次齐次的。只有这样，当边长等比例扩大时，它们才可以保持不变。

什么时候适合使用余弦定理呢？只要有它所应具有的代数结构时，都可以考虑。\\
（1）齐次性。若三角形 \(A B C\) 中，有 \(b^{2}+c^{2}-b c=a^{2}\) ，这显然是一个二次齐次的式子。可以考虑使用余弦定理。根据余弦定理，

\[
a^{2}=b^{2}+c^{2}-2 b c \cos A=b^{2}+c^{2}-b c
\]

故 \(2 \cos A=1, A=\frac{\pi}{3}\) ．\\
注意正弦定理也可以在齐次的情况下使用，所以二者可能会同时使用。例如，2019年课标I卷的三角形大题的已知条件为：三角形 \(A B C\) 中成立

\[
(\sin B-\sin C)^{2}=\sin ^{2} A-\sin B \sin C .
\]

如果直接进行三角恒等变换，可能较难处理。我们注意到已知是关于 \(\sin A, \sin B, \sin C\) 的二次齐次式，可以考虑使用正弦定理。将 \(\sin A, \sin B, \sin C\) 用 \(\frac{a}{2 R}, \frac{b}{2 R}, \frac{c}{2 R}\) 替换，约掉所有 \(R\) ，可得

\[
(b-c)^{2}=a^{2}-b c
\]

化简得 \(b^{2}+c^{2}-b c=a^{2}\) ，化为第一种情况，可得 \(A=\frac{\pi}{3}\) 。因此，当我们有齐次式的时候可以考虑余弦定理。 （2）对称性。当已知条件是关于其中两边的对称式时，可以尝试余弦定理。例如，若三角形 \(A B C\) 的面积为 \(\frac{a^{2}+b^{2}-c^{2}}{4}\) ，如何化简？注意到这是一个二次齐次式，且关于 \(a, b\) 对称，所以可以考虑用 \(C\) 的余弦定理：\\
\(c^{2}=a^{2}+b^{2}-2 a b \cos C\) ．故

\[
S=\frac{a^{2}+b^{2}-c^{2}}{4}=\frac{1}{2} a b \cos C
\]

而根据三角形面积公式 \(S=\frac{1}{2} a b \sin C\) ，得 \(\cos C=\sin C\) 。故 \(C=\frac{\pi}{4}\) 。因此，当式子本身具有对称性时，可以考虑余弦定理。

\section*{总结}
正弦定理，余弦定理都有丰富的代数结构，在使用时可以根据已知条件变通选择：\\
（1）正弦定理：关于边，正弦都是一次齐次；\\
（2）余弦定理：关于边二次齐次，关于两边对称。\\
因此，当已知条件具有齐次结构时，既可以考虑正弦定理，也可以考虑余弦定理；当已知条件关于两边对称时，多使用余弦定理。正弦定理和余弦定理的使用不是绝对的，很多情况下，既可以使用余弦定理，也可以使用正弦定理。这揭示了这两个定理的等价性：正弦定理，余弦定理都是描述三角形内部的关系，二者可以互相推导。试试看，你能不能从余弦定理推出正弦定理，或者从正弦定理推出余弦定理呢？

\section*{4．3．3 解三角形}
解三角形，就是为了确定三角形的边，角。而三角形的边和角 \(a, b, c, A, B, C\) 共六个条件，如何才能确定下来呢？确定三角形，本质上就是从已知的边，角出发，利用 \(A+B+C=\pi\) ，正弦定理和余弦定理，把其他的边，角确定：

已知的边，角

正弦定理，余弦定理都有三个等号，所以它们相当于三个限制条件。根据自由度的知识，我们还需要 \(6-3=3\)个已知的边角，才可以确定其他所有的边角。这也就解释了，为什么所有三角形全等判定都是有三个条件的。这也是自由度思想的一种体现。

在这里，我们仅以 \(S A S\) 为例，证明它可以确定三角形。假设 \(b, c, A\) 已知。\\
（1）根据余弦定理，

\[
a^{2}=b^{2}+c^{2}-2 b c \cos A
\]

故 \(a=\sqrt{b^{2}+c^{2}-2 b c \cos A}\) ．边 \(a\) 可以确定了。\\
（2）根据正弦定理，

\[
\sin B=\frac{b}{a} \sin A, \sin C=\frac{c}{a} \sin A
\]

由于 \(a, b, c, A\) 均已知，故 \(\sin B, \sin C\) 可确定，角 \(B, C\) 即可确定。或亦可用余弦定理：

\[
\cos B=\frac{a^{2}+c^{2}-b^{2}}{2 a c}, \cos C=\frac{a^{2}+b^{2}-c^{2}}{2 a b}
\]

确定 \(\cos B, \cos C\) ，亦可确定角 \(B, C\) ．\\
所有其他的三角形判定定理，都可以通过正弦定理，余弦定理完成证明。感兴趣的同学可以参考强基计划的讲义（制作中），里面有更详细的介绍。

我们在用自由度解释解三角形时，提到＂正弦定理，余弦定理各有三个条件＂，可以确定三个自由度。但是， \(A+B+C=\pi\) 是三角形自带的条件，按理说我们一共有四个限制条件，可以确定四个自由度。那为什么还是需要三个已知的边，角才可以判定全等呢？这是因为：在证明正弦定理，余弦定理的过程中，需要使用 \(A+B+C=\pi\)这个条件。所以，这个条件隐含在正弦定理，余弦定理之中了，相当于三个限制条件。

\section*{4．3．4 XXXX}
\section*{4.4 立体几何}
\section*{4．4．1 直线和平面的确定}
\section*{直线的方向向量与平面的法向量}
我们在高中也学习过向量的内容。向量代表的是有大小，方向的量，它们有对应的坐标表示。在立体几何的学习中，我们经常会使用向量来帮助我们确定平面的和直线的方向。那么，如何从自由度的角度来理解这个方向是如何确定的呢？

空间中的直线可以直接由它的方向向量来确定。这是因为：两点确定一条直线，设 \(A\left(x_{1}, y_{1}, z_{1}\right), B\left(x_{2}, y_{2}, z_{2}\right)\)是直线 \(l\) 上的两个点，那么

\[
\overrightarrow{A B}=\left(x_{2}-x_{1}, y_{2}-y_{1}, z_{2}-z_{1}\right)
\]

就是直线的方向向量。方向确定了，只需要再知道直线上的一个点，就知道直线往哪个方向延伸，直线也就确定了。因此，直线 \(l\) 可以有如下两种不同的表示方法：\\
（1）过 \(A, B\) 两点的直线；\\
（2）过 \(A\)（或 \(B\) ）的直线，且方向向量为 \(\overrightarrow{A B}=\left(x_{2}-x_{1}, y_{2}-y_{1}, z_{2}-z_{1}\right)\) ．\\
无论是哪种方法，都需要两个条件才可以确定一个平面。换言之，一条直线 \(l\) 有两个自由度\\
我们再来看空间中的平面。如何确定空间中的平面呢？一个最基本的事实是：不共线三点确定一个平面。设 \(A\left(x_{1}, y_{1}, z_{1}\right), B\left(x_{2}, y_{2}, z_{2}\right), C\left(x_{3}, y_{3}, z_{3}\right)\) 三点不共线，那么这三个点就可以确定一个三角形，这个三角形所在平面就是这三点确定的平面。如果我们和直线类比，是否有一个类似于＂方向向量＂的向量，可以确定平面的方向呢？和直线不同，平面需要的不是方向向量，而是法向量 \(\vec{n}\) ．

平面的法向量 \(\vec{n}\) 垂直于平面中的任意一条直线，确定了它。按理说，如果想要找到平面的法向量，我们需要把平面中所有的直线都写出来，列无数个方程，才可以得到法向量。而在实际计算中，为什么只需要列两个不共线的方向就可以了呢？这就是＂平面向量基本定理＂的功劳。

设 \(\vec{n}\) 是平面 \(A B C\) 的一个法向量。那么 \(\vec{n}\) 满足：

\[
\vec{n} \cdot \overrightarrow{A B}=0, \vec{n} \cdot \overrightarrow{A C}=0
\]

为什么我们不需要再单独要求 \(\vec{n} \cdot \overrightarrow{B C}=0\) 了呢？因为根据前面两个式子，我们可以直接推出这个式子：

\[
\vec{n} \cdot \overrightarrow{B C}=\vec{n} \cdot(\overrightarrow{A C}-\overrightarrow{A B})=\vec{n} \cdot \overrightarrow{A C}-\vec{n} \cdot \overrightarrow{A B}=0-0=0
\]

我们看到，因为 \(\overrightarrow{\boldsymbol{B C}}\) 可以被 \(\overrightarrow{\boldsymbol{A B}}, \overrightarrow{\boldsymbol{A C}}\) 表示，所以根据向量点乘的分配律，这个条件已经蕴含在前面两个式子中了。对于这个平面中的其他任何一个向量 \(\vec{a}\) ，根据＂平面向量基本定理＂，存在唯一的实数 \(\lambda, \mu\) 使得

\[
\vec{a}=\lambda \cdot \overrightarrow{A B}+\mu \cdot \overrightarrow{A C}
\]

因此，有

\[
\vec{n} \cdot \vec{a}=\vec{n} \cdot(\lambda \cdot \overrightarrow{A B}+\mu \cdot \overrightarrow{A C})=\lambda \cdot \vec{n} \cdot \overrightarrow{A B}+\mu \cdot \vec{n} \cdot \overrightarrow{A C}=0+0=0
\]

因此，对于这个平面内的任意一个向量 \(\vec{a}\) ，都有 \(\vec{n} \perp \vec{a}\) ，故 \(\vec{n}\) 是这个平面的法向量。\\
从直观上讲，当平面的法向量确定之后，这个平面的＂倾斜角度＂也就确定了。那么只要再给出平面所过的一个点，就可以确定平面了。因此，我们有如下两种方式来确定一个平面：\\
（1）不共线三点 \(A, B, C\) 确定一个平面；\\
（2）过 \(A\)（或 \(B\) ，或 \(C\) ），且 \(\vec{n}\) 是这个平面的法向量。\\
明显（1）中有三个限制条件，但是为什么（2）只有两个限制条件？平面的自由度到底是 3 还是 2 ？\\
事实上，＂\(\vec{n}\) 为法向量＂这个条件相当于两个限制条件。从前面我们对于 \(\vec{n}\) 的说明即可看出：\(\vec{n}\) 是平面的法向量，当且仅当它和平面中两个不共线向量垂直。在这个例子中，它可以是＂\(\vec{n} \cdot \overrightarrow{A B}=0, \vec{n} \cdot \overrightarrow{A C}=0 "\) ，也可以是＂ \(\vec{n} \cdot \overrightarrow{A B}=0, \vec{n} \cdot \overrightarrow{B C}=0 "\) ，还可以是 \(" \vec{n} \cdot \overrightarrow{B C}=0, \vec{n} \cdot \overrightarrow{A C}=0\) ．＂因此，给出法向量相当提供了两个限制条件。无论是 （1），（2），都有三个限制条件。因此，平面的自由度为 3 。

事实上，还有一种确定平面的方法：两条相交直线确定一个平面。我们可以从方向向量的角度来理解。设这两条相交平面交于 \(A\) ，方向向量分别为 \(\overrightarrow{a_{1}}, \overrightarrow{a_{2}}\) ．那么，这两条相交直线确定的平面满足：\\
（1）它经过 \(A\) ；\\
（2）法向量 \(\vec{n}\) 垂直于 \(\overrightarrow{a_{1}}\) ；\\
（3）法向量 \(\vec{n}\) 垂直于 \(\overrightarrow{a_{2}}\) 。\\
也是三个限制条件！因此，我们理解了：当给出两条相交直线时，它们确定一个平面的方式是给出交点和平面上两个不共线的向量。后者确定了平面方向（法向量的方向），再加上一个平面所过的点，平面就确定了。因此，这三种方式都可以确定一个平面：\\
（1）不共线三点；\\
（2）两条相交直线；\\
（3）一点和一个法向量。\\
其实，＂不共线三点＂可以很自然地理解成两条相交直线：直线 \(A B, A C\) 。所以，（1），（2）本质上相同。\\
最后，我们来回答一个问题：为什么确定法向量的方向，需要两个条件呢？事实上，我们可以从多个角度理解这个问题。\\
（1）平面向量基本定理：两个不共线向量可以表示所有这个平面中的向量，所以只要和两个垂直即可；\\
（2）自由度：法向量 \(\vec{n}=(x, y, z)\) 有三个未知数。想要确定它的方向，只需要确定它们个坐标的比值，即 \(x: y: z\) ．而根据齐次结构的理论，确定 \(n\) 个未知数的比例，需要 \((n-1)\) 个齐次式。所以，我们只需要两个齐次条件即可。那两个齐次条件呢？就是＂和两个不共线向量垂直＂。为什么它们是齐次条件呢？设 \(\overrightarrow{a_{1}}=\left(x_{1}, y_{1}, z_{1}\right), \overrightarrow{a_{2}}=\left(x_{2}, y_{2}, z_{2}\right)\) ，因此

\[
\begin{aligned}
& \vec{n} \cdot \overrightarrow{a_{1}}=x_{1} x+y_{1} y+z_{1} z=0 \\
& \vec{n} \cdot \overrightarrow{a_{2}}=x_{2} x+y_{2} y+z_{2} z=0
\end{aligned}
\]

它们都是关于 \(x, y, z\) 的一次齐次式！这就确定了法向量的方向。因此，我们要知道：在三维空间中，当一个向量和两个已知向量垂直，它的方向就确定了。

在立体几何的建系计算中，为了计算方便，我们往往直接写出一个法向量。想要写出一个法向量，就需要

补充一个非齐次条件，比如 \(x=1\) 。这样，一个非齐次条件和两个齐次条件，就可以确定法向量的三个坐标分量了。这也体现了咱们自由度理论在立体几何中的基本应用。

\section*{4．4．2 法向量的计算}
在建系计算中，经常会需要计算给定平面的法向量。传统的方法是直接解如下方程组：

\[
\begin{aligned}
& \vec{n} \cdot \overrightarrow{a_{1}}=x_{1} x+y_{1} y+z_{1} z=0 \\
& \vec{n} \cdot \overrightarrow{a_{2}}=x_{2} x+y_{2} y+z_{2} z=0
\end{aligned}
\]

但这个计算往往需要消元。有没有更简单的方法计算法向量呢？\\
事实上，如果给定的两个向量中，有向量的某分量为 0 ，我们可以很快地写出可能的法向量。例如，对于向量 \(\vec{a}=(0,1, \sqrt{3})\) ，所有和它垂直的向量一定可以写作 \((x,-\sqrt{3}, 1)\) 的形式。如果利用这个性质，就可以很快写出法向量可能的形式了，并进一步求解。

例如，若平面内两直线的方向向量为 \(\vec{a}=(0,1, \sqrt{3}), \vec{b}=(-1,3 \sqrt{3}, 3)\) ，则可直接设平面法向量为 \(\vec{n}=\) \((x,-\sqrt{3}, 1)\) ．此时，\(\vec{n} \perp \vec{a}\) 总成立，故只需满足

\[
\vec{n} \cdot \vec{b}=-x+9+3=12-x=0
\]

即可。解得 \(x=12\) ．故 \(\vec{n}=(12,-\sqrt{3}, 1)\) 是一个法向量。\\
在一些情况下，可能没有分量是 0 ．此时该如何处理呢？有两种处理方法：例如，若平面内两直线的方向向量为 \(\vec{a}=(1,-2,4), \vec{b}=(2,3,1)\) ，如何求法向量？\\
（1）对已知向量进行线性组合。考虑

\[
\vec{c}=\vec{b}-2 \vec{a}=(2,3,1)-(2,-4,8)=(0,7,-7)
\]

故只需让 \(\vec{n}\) 和 \(\vec{a} \vec{c}\) 垂直即可。这是因为：\(\vec{c}\) 也是在平面内和 \(\vec{a}\) 不共线的向量，所以和 \(\vec{a}, \vec{b}\) 等价于和 \(\vec{a} \vec{c}\) 垂直。由于 \(\vec{c}\) 的 \(x\) 分量为 0 ，故法向量必然具有 \((x, 7,7)\) 的形式。当我们遇到这种有大于 1 的公约数的情况时，可以＂约分＂：约掉 7 ，设法向量为 \(\vec{n}=(x, 1,1)\) 。这是因为，我们只关心方向，不关心大小，同乘，同除不影响方向。这样一来，可以直接得

\[
\vec{n} \cdot \vec{a}=x-2+4=x+2=0
\]

故 \(x=-2\) ，平面的一个法向量为 \((-2,1,1)\) ．\\
（2）可以直接设一个分量为 1 。这是因为，我们最重要确定法向量需要一个非齐次条件，可以直接在一开始就使用这个非其次条件。令 \(\vec{n}=(1, y, z)\) ，则

\[
\begin{aligned}
& \vec{n} \cdot \vec{a}=1-2 y+4 z=0 \\
& \vec{n} \cdot \vec{b}=2+3 y+z=0
\end{aligned}
\]

解方程即可。解得 \(y=z=-\frac{1}{2}\) ．\\
这两种方法可以根据实际情况选择，各有优劣。

\section*{4．4．3 平行于垂直的判定和性质}
如何从方向向量，法向量的角度来理解平行，垂直的判定定理和性质定理呢？首先，我们从平行出发。

\section*{4．4．3．1 平行}
（1）线面平的判定定理：若 \(a / / b\) ，且 \(a \not \subset, b \subset \alpha\) ，则 \(a / / \alpha\) ．\\
（2）面面平的判定定理：若 \(a, b \subset \alpha, a \cap b \neq \emptyset\) ，且 \(a / / \beta, b / / \beta\) ，则 \(\alpha / / \beta\) ．\\
（3）线面平的性质定理：若 \(a / / \alpha, a \subset \beta, \alpha \cap \beta=b\) ，则 \(a / / b\) 。\\
（4）面面平的性质定理：若 \(\alpha / / \beta, \gamma \cap \alpha=a, \gamma \cap \beta=b\) ，则 \(a / / b\) 。\\
我们首先观察到了如下两个基本规律：\\
（1）＂\(X X\)＂判定定理，以＂\(X X "\) 作为结论；＂\(X X\)＂性质定理，以＂\(X X\)＂作为条件。判定的意思是我们要最终判断出这个结果，所以＂\(x X^{\prime}\)＂作为结论出现；性质的意思是我们关心这个情况下它有哪些性质，所以＂XX＂作为条件出现。\\
（2）判定定理＂升维＂，性质定理＂降维＂。直线是二维对象，平面则是三维对象。判定定理都是从低维关系推导出高维关系，而性质定理都是从高维关系推导出低维关系。这也解释了人教版教材上给出的定理之间的关系图：\\
\includegraphics[max width=\textwidth, center]{2025_02_16_7c9aaf92bc99dc07f048g-094}

如何用向量的语言来描述平行关系呢？显然线线平，线面平是很好描述的，只要方向向量或法向量平行即可：

\[
\begin{aligned}
& l_{1} / / l_{2} \Leftrightarrow \overrightarrow{a_{1}} / / \overrightarrow{a_{2}}, \\
& \alpha / / \beta \Leftrightarrow \overrightarrow{n_{1}} / / \overrightarrow{n_{2}} .
\end{aligned}
\]

对于线面平而言，只需要直线的方向向量和平面的法向量垂直即可。可以想象：这条直线可以通过平移成为这个平面上的一条直线，所以它的方向向量必然和平面的法向量垂直。因此，

\[
a / / \alpha \Leftrightarrow \overrightarrow{a_{1}} \perp \vec{n}, a \not \subset \alpha
\]

有了这些条件，我们的判定定理，性质定理就变得非常自然：\\
（1）线面平判定定理：由于 \(a\) 和 \(b\) 平行，所以它们方向向量平行。设 \(\vec{n}\) 是平面的法向量，则

\[
\vec{a} \cdot \vec{n}=\lambda \vec{b} \cdot \vec{n}=0
\]

又 \(a \not \subset \alpha\) ，所以 \(a / / \alpha\) ．\\
（2）面面平判定定理：设 \(\overrightarrow{n_{1}}, \overrightarrow{n_{2}}\) 分别为 \(\alpha, \beta\) 的法向量，由于 \(a / / \beta, b / / \beta\) ，因此

\[
\vec{a} \cdot \overrightarrow{n_{2}}=0, \vec{b} \cdot \overrightarrow{n_{2}}=0
\]

由于 \(\vec{a}, \vec{b}\) 不共线，所以根据自由度的理论， \(\overrightarrow{n_{2}}\) 的方向就确定了。又因为 \(\vec{a} \cdot \overrightarrow{n_{1}}=\vec{b} \cdot \overrightarrow{n_{1}}=0\) ，所以 \(\overrightarrow{n_{1}}\) 和 \(\overrightarrow{n_{2}}\) 方向相同，所以 \(\overrightarrow{n_{1}} / / \overrightarrow{n_{2}}\) ，故 \(\alpha / / \beta\) ．\\
（3）线面平性质定理：设 \(\alpha, \beta\) 的法向量为 \(\overrightarrow{n_{1}}, \overrightarrow{n_{2}}\) ，由于 \(b=\alpha \cap \beta\) ，所以

\[
\vec{b} \cdot \overrightarrow{n_{1}}=0, \vec{b} \cdot \overrightarrow{n_{2}}=0
\]

根据自由度理论，\(\vec{b}\) 方向确定了。而 \(a \subset \beta\) ，所以 \(\vec{a} \cdot \overrightarrow{n_{2}}=0\) ．又 \(a / / \alpha\) ，所以 \(\vec{a} \cdot \overrightarrow{n_{1}}=0\) ．因此，对于 \(\vec{a}\) ，也成立

\[
\vec{a} \cdot \overrightarrow{n_{1}}=0, \vec{a} \cdot \overrightarrow{n_{2}}=0
\]

根据自由度理论，\(\vec{a}\) 方向也确定了，且必然和 \(\vec{b}\) 方向相同。所以 \(a / / b\) ．\\
（4）面面平性质定理：由于 \(\alpha, \beta\) 平行，所以可设 \(\vec{n}\) 都是它们的法向量。设 \(\gamma\) 的法向量分别为 \(\overrightarrow{n_{1}}\) 。由于 \(a=\alpha \cap \gamma\)\\
，所以

\[
\vec{a} \cdot \overrightarrow{n_{1}}=\vec{a} \cdot \vec{n}=0
\]

由于 \(b=\beta \cap \gamma\) ，所以

\[
\vec{b} \cdot \overrightarrow{n_{1}}=\vec{b} \cdot \vec{n}=0
\]

所以 \(a / / b\) ．\\
我们发现，所有这些平行的判定和性质定理，本质上都是利用和两个不共线向量垂直，进而确定方向，导出平行的。

\section*{4．4．3．2 垂直}
（1）线面垂的判定定理：若 \(a, b \subset \alpha, a \cap b \neq \emptyset\) ，且 \(l \perp a, l \perp b\) ，则 \(l \perp \alpha\) ；\\
（2）面面垂的判定定理：若 \(n \perp \alpha, n \subset \beta\) ，则 \(\alpha \perp \beta\) ；\\
（3）线面垂的性质定理：若 \(a \perp \alpha, b \perp \alpha\) ，则 \(a / / b\) ；\\
（4）面面垂的性质定理：若 \(\alpha \perp \beta, l=\alpha \cap \beta, n \subset \alpha\) ，且 \(n \perp l\) ，则 \(n \perp \beta\) ．\\
和平行类似，我们首先用向量刻画垂直关系。\\
（1）\(a \perp b \Leftrightarrow \vec{a} \cdot \vec{b}=0\) ；\\
（2）\(a \perp \alpha \Leftrightarrow \vec{a} / / \vec{n}\) ，其中 \(\vec{n}\) 是平面 \(\alpha\) 的法向量；\\
（3）\(\alpha \perp \beta \Leftrightarrow \overrightarrow{n_{1}} \perp \overrightarrow{n_{2}}\) ，其中 \(\overrightarrow{n_{1}}, \overrightarrow{n_{2}}\) 分别是平面 \(\alpha, \beta\) 的法向量。\\
有了这些关系，有关垂直的定理也就显而易见了：\\
（1）线面垂的判定定理：由 \(\vec{l} \cdot \vec{a}=\vec{l} \cdot \vec{b}=0\) 直接可得；\\
（2）面面垂的判定定理：设 \(\overrightarrow{n_{1}}\) 为 \(\beta\) 的法向量，故 \(\overrightarrow{n_{1}} \perp \overrightarrow{n_{2}}\) ．故 \(\alpha \perp \beta\) ；\\
（3）线面垂的性质定理：由 \(a \perp \alpha, b \perp \alpha\) 知 \(a / / n, b / / n\) ，故 \(a / / b\) ；\\
（4）面面垂的性质定理：设 \(\overrightarrow{n_{1}}, \overrightarrow{n_{2}}\) 分别为 \(\alpha, \beta\) 的法向量，则 \(\overrightarrow{n_{1}} \cdot \overrightarrow{n_{2}}=0\) ．又 \(l=\alpha \cap\)\\
\(\beta\) ，所以 \(\vec{l} \cdot \overrightarrow{n_{1}}=\vec{l} \cdot \overrightarrow{n_{2}}=0\) ．又 \(n \perp l\) ，所以 \(\vec{n} \perp \vec{l}=0\) ．又 \(n \subset \alpha\) ，所以 \(\vec{n} \cdot \overrightarrow{n_{1}}=0\) ．\\
注意到 \(\overrightarrow{n_{2}}\) 满足

\[
\overrightarrow{n_{1}} \cdot \overrightarrow{n_{2}}=\vec{l} \cdot \overrightarrow{n_{2}}=0
\]

而 \(\vec{n}\) 满足

\[
\overrightarrow{n_{1}} \cdot \vec{n}=\vec{l} \cdot \vec{n}=0
\]

所以 \(\vec{n}\) 和 \(\overrightarrow{n_{2}}\) 共线，故 \(n \perp \beta\) ．

\section*{4．4．4 XXXX}
\section*{4.5 数列}
数列本质上就是函数。我们可以将数列中的每一项 \(a_{n}\) 理解成自变量为 \(n\) ，函数值为 \(a_{n}\) 的函数 \(f\) ，即 \(a_{n}=f(n)\) 。所有的数列通项公式都可以这么理解。例如，等差数列 \(a_{n}=a_{1}+(n-1) d=d n+a_{1}-d\) ，它是一个关于 \(n\) 的一次函数；等比数列 \(a_{n}=a_{1} q^{n-1}\) ，可以理解成关于 \(n\) 的指数函数。因此，求数列的通项公式就是求函数的解析式。

我们首先以等差数列和等比数列为例，来理解这部分内容。

\section*{4．5．1 等差数列与一次函数}
等差数列 \(a_{n}=a_{1}+(n-1) d\) 可以看作关于 \(n\) 的一次函数。

\[
a_{n}=a_{1}+(n-1) d=d n+\left(a_{1}-d\right)=f(n)
\]

可以看到，这是一个斜率为 \(d\) ，截距为 \(a_{1}-d\) 的一次函数。\\
虽然数列只能在 \(n\) 为正整数的时候取值，但是我们可以自然地将它延拓到所有实数上。例如，对于首项为 1 ，公差为 2 的等差数列，它的通项公式为

\[
a_{n}=1+2(n-1)=2 n-1
\]

我们可以自然地去计算

\[
a_{2.5}=2 \cdot 2.5-1=4, a_{\frac{4}{3}}=2 \cdot \frac{4}{3}-1=\frac{5}{3}
\]

这样的话，数列本身就和一次函数没有任何区别了。数列的每一项，都可以理解成一次函数上的一个点。例如， \(a_{2.5}=4\) 可以理解成直线 \(y=2 x-1\) 上的点 \((2.5,4), a_{\frac{4}{3}}\) 可以理解成直线上的点 \(\left(\frac{4}{3}, \frac{5}{3}\right)\) 。所有等差数列 \(a_{n}=2 n-1\)的项，都可以理解成直线 \(y=2 x-1\) 上的点 \(A_{n}(n, 2 n-1)\) ．

等差数列还有一个重要的结构，就是齐次结构。数列中的所有项，都可以看作关于 \(a_{1}, d\) 的一次齐次式。例如：

\[
a_{3}=a_{1}+2 d, a_{6}=a_{1}+5 d
\]

因此，所有关于数列的项的齐次式都可以转化成 \(a_{1}, d\) 的齐次式。例如：\\
（1） \(2 a_{7}=a_{5}+2 a_{3}\) ，可以写作

\[
\begin{aligned}
& 2 a_{1}+12 d=a_{1}+4 d+2 a_{1}+4 d \\
& a_{1}=4 d
\end{aligned}
\]

进而，可以求得 \(a_{1}, d\) 的比例。\\
（2）\(a_{4}^{2}=2 a_{1} a_{8}\) ，是关于项的二次齐次式，那么我们可以得到：

\[
\begin{aligned}
& \left(a_{1}+3 d\right)^{2}=2 a_{1} \cdot\left(a_{1}+7 d\right) \\
& a_{1}^{2}+8 a_{1} d-9 d^{2}=0 \\
& \left(a_{1}-d\right)\left(a_{1}+9 d\right)=0
\end{aligned}
\]

故 \(a_{1}=d\) 或 \(a_{1}=-9 d\) ．也可以求出 \(a_{1}, d\) 的比例。\\
这和我们在自由度中获得的理解是一致的：对于有两个自由度的问题，只需要一个齐次式，就可以求出比例。因此，所有关于等差数列的项的齐次条件，都可以确定 \(a_{1}, d\) 的比例。当我们再看到类似的条件时，就可以＂有底气＂地处理它。

\section*{4．5．2 利用齐次，线性快速计算}
想要确定一个等差数列，本质上需要两个条件：公差 \(d\) 和首项 \(a_{1}\) 。这和它的图像是吻合的：它可以理解为一次函数（直线），而两点确定一条直线，所以只需要知道直线上的两点（等差数列中的两项），即可确定等差数列。

假设我们知道 \(a_{m}=a, a_{n}=b\) 的值，如何求出首项 \(a_{1}\) 和公差 \(d\) 呢？如果直接从通项公式理解：

\[
\begin{aligned}
& a_{1}+(m-1) d=a \\
& a_{1}+(n-1) d=b
\end{aligned}
\]

我们可以直接使用消元的方法求解 \(a_{1}, d\) ．两式相减，得

\[
\begin{aligned}
& (m-n) d=a-b \\
& d=\frac{a-b}{m-n}
\end{aligned}
\]

代回第一个式子，得

\[
\begin{aligned}
a_{1}=a-(m-1) d & =a-(m-1) \cdot \frac{a-b}{m-n} \\
& =\frac{a m-a n-a m+b m-a+b}{m-n} \\
& =\frac{b m-a n-a+b}{m-n} .
\end{aligned}
\]

重新改写为 \(a_{m}, a_{n}\) 的式子，可以得到：

\[
d=\frac{a_{m}-a_{n}}{m-n}, a_{1}=\frac{(m+1) a_{n}-(n+1) a_{m}}{m-n} .
\]

因此，我们代数上验证了自由度的观察。同时，我们发现这两个式子是关于 \(m, n\) 对称的式子。这也是合理的，因为 \(m, n\) 地位平等。

这个公式又复杂，又不好记，如何快速计算呢？事实上，我们不需要记忆上面这个式子来求 \(a_{1}, d\) ，而可以从更直观的角度理解。回到直线的理解，我们可以将 \(a_{m}, a_{n}\) 理解成函数值。那么，求通项公式，等价于求直线的方程。已知直线过 \(\left(m, a_{m}\right),\left(n, a_{n}\right)\) 两点，那么这条直线的斜率为

\[
k=\frac{a_{m}-a_{n}}{m-n}
\]

而等差数列的通项公式 \(a_{n}=d n+\left(a_{1}-d\right)=f(n)\) ，将 \(f(n)\) 就是斜率为 \(d\) ，截距为 \(a_{1}-d\) 的直线方程：

\[
y=d x+\left(a_{1}-d\right)
\]

因为，公差就是斜率

\[
d=k=\frac{a_{m}-a_{n}}{m-n}
\]

接下来再去求截距，即可求出等差数列的通项。\\
在实际应用中，只需要理解＂公差就是斜率＂，就可以顺利求出等差数列的通项，不需要刻意记忆任何内容。这是因为：实际应用中都是相对简单的数字，不是复杂的字母，直接计算远比死记硬背要方便得多。例如，已知 \(a_{10}=3, a_{20}=8\) ，那么可以理解为经过 \((10,3),(20,8)\) 的直线。它的斜率为

\[
d=\frac{8-3}{20-10}=\frac{5}{10}=\frac{1}{2}
\]

故首项为

\[
a_{1}=a_{10}-9 d=3-\frac{9}{2}=-\frac{3}{2}
\]

故数列的通项公式为 \(a_{n}=-\frac{3}{2}+(n-1) \cdot \frac{1}{2}=\frac{n}{2}-2\) ．\\
刚才的方法，可以已知数列中的两项，快速求出数列的通项。求出通项，就可以求数列中的任意一项了。但是，如果我们要求数列中指定的一项，有没有更好的方法呢？例如，我们已知 \(a_{1}=7, a_{4}=13\) ，如何求 \(a_{2}\) 呢？

按着刚才的方法，我们首先可以求出斜率：

\[
d=\frac{a_{4}-a_{1}}{4-1}=\frac{13-7}{4-1}=2
\]

那么 \(a_{2}=a_{1}+d=7+2=9\) ．我们用了两步计算完成。但事实上，如果你理解了＂\(a_{2}\) 是 \(a_{1}, a_{4}\) 靠近 \(a_{1}\) 的三等分点＂，我们可以一步得到答案：

\[
a_{2}=\frac{2}{3} a_{1}+\frac{1}{3} a_{4}=\frac{2}{3} \cdot 7+\frac{1}{3} \cdot 13=\frac{27}{3}=9
\]

为什么＂\(a_{2}\) 是 \(a_{1}, a_{4}\) 靠近 \(a_{1}\) 的三等分点＂呢？还是回到我们直线的理解：\(A\left(1, a_{1}\right), B\left(2, a_{2}\right), C\left(4, a_{4}\right)\) 是直线上共线的三点。观察到它们的横坐标 \(1,2,4\) ，刚好满足 \((2-1)=\frac{1}{3} \cdot(4-1)\) 。所以，\(B\) 是 \(A C\) 的定比分点，满足

\[
\overrightarrow{A B}=\frac{1}{3} \overrightarrow{A C}
\]

那么，根据向量的坐标表示，可得

\[
\left(1, a_{2}-a_{1}\right)=\frac{1}{3} \cdot\left(3, a_{4}-a_{1}\right)=\left(1, \frac{a_{4}-a_{1}}{3}\right)
\]

故

\[
a_{2}-a_{1}=\frac{a_{4}-a_{1}}{3}=\frac{2}{3} a_{1}+\frac{1}{3} a_{4}
\]

在实际应用中，我们只要有关于定比分点的理解，就可以直接书写：\\
（1）比例之和一定为 1 ，所以 \(a_{2}\) 一定可以写作

\[
a_{2}=t a_{1}+(1-t) a_{4}
\]

的形式；\\
（2） 2 是 1,4 之间的一个三等分点，如果不考虑远近，\(t\) 必然是 \(\frac{1}{3}\) 或 \(\frac{2}{3}\) ；\\
（3） 2 离 1 ，比 2 离 4 近。离得近，意味着更像，所以权重更大，所以 \(t=\frac{2}{3}\) ．故

\[
a_{2}=\frac{2}{3} a_{1}+\frac{1}{3} a_{4}
\]

我们再举一个更复杂的例子：用 \(a_{3}, a_{13}\) 表示 \(a_{7}\) 。我们发现： \(7-3=4,13-7=6,4: 6=2: 3\) ，所以 7 是 3,13的一个五等分点，靠近 3 的第二个五等分点。靠近它，权重大，故

\[
a_{7}=\frac{3}{5} a_{3}+\frac{2}{5} a_{13}
\]

可以用通项公式验证：\(\frac{3}{5} a_{3}+\frac{2}{5} a_{13}=\frac{3}{5} \cdot\left(a_{1}+2 d\right)+\frac{2}{5} \cdot\left(a_{1}+12 d\right)=a_{1}+\left(\frac{6}{5}+\frac{24}{5}\right) d=a_{1}+6 d=a_{7}\) ．因此，当已知数列中的两项时，可以直接利用定比分点求解数列中的任意一项。事实上，定比分点中点概念的推广，我们刚才的求法也是等差中项的推广。

除了在 \(a_{n}\) 中可以利用这个观察，在 \(S_{n}\) 中也可以。我们知道，等差数列求和是利用＂首尾配对＂的思路求解的。例如：

\[
\begin{aligned}
& S_{9}=a_{1}+a_{2}+\cdots+a_{9} \\
& 2 S_{9}=\left(a_{1}+a_{9}\right)+\left(a_{2}+a_{8}\right)+\cdots+\left(a_{8}+a_{2}\right)+\left(a_{9}+a_{1}\right)
\end{aligned}
\]

我们发现，每一对括号内的项的下角标之和均为 10 ，所以，可以利用刚才的观察： 5 是 \(1,92,83,74,6\) 的中点，

所以

\[
\begin{aligned}
& \frac{1}{2} a_{1}+\frac{1}{2} a_{9}=a_{5} \\
& \frac{1}{2} a_{2}+\frac{1}{2} a_{8}=a_{5} \\
& \frac{1}{2} a_{3}+\frac{1}{2} a_{7}=a_{5} \\
& \frac{1}{2} a_{4}+\frac{1}{2} a_{6}=a_{5}
\end{aligned}
\]

因此，每个括号都是 \(2 a_{5}\) ，从 \(a_{1}\) 到 \(a_{9}\) 共 9 个括号，所以

\[
2 S_{9}=9 \cdot 2 a_{5}=18 a_{5}
\]

故 \(S_{9}=9 a_{5}\) ．事实上，我们发现，\(a_{5}\) 正好是这些项的平均数：

\[
a_{5}=\frac{S_{9}}{9}
\]

一般地，前 \(n\) 项的平均数是 \(\frac{a_{1}+a_{n}}{2}=a_{\underline{n+1}}^{2}\) ，所以

\[
S_{n}=n \cdot a_{\frac{n+1}{2}}^{2}
\]

这样一来，当我们知道前若干项的和，本质上也就是知道了数列中的某一项。

\section*{4．5．3 等比数列与指数函数}
等比数列的通项公式为 \(a_{n}=a_{1} \cdot q^{n-1}\) ，我们可以将它理解成指数函数。有没有可能像等差数列一样，利用＂定比分点＂和齐次结构来理解等比数列呢？

事实上，只需要取一个对数即可。假设 \(a_{1}, q\) 都是正数，那么

\[
\ln a_{n}=\ln \left(a_{1} \cdot q^{n-1}\right)=\ln a_{1}+(n-1) \cdot \ln q
\]

因此，等比数列的对数 \(\left\{\ln a_{n}\right\}\) 可以看作是首项为 \(\ln a_{1}\) ，公差为 \(\ln q\) 的等差数列。这样一来，所有关于等差数列＂定比分点＂都可以平移到等比数列来。

例如，\(a_{1}, a_{4}\) 如何确定 \(a_{2}\) ？由于 \(\ln a_{n}\) 是等差数列，所以

\[
\begin{aligned}
& \ln a_{2}=\frac{2}{3} \cdot \ln a_{1}+\frac{1}{3} \cdot \ln a_{4} \\
& a_{2}^{3}=a_{1}^{2} a_{4} .
\end{aligned}
\]

故 \(a_{2}^{3}=a_{1}^{2} a_{4}\) 。我们同样在这里观察到了齐次结构：左右两侧都是关于项的三次齐次。\\
注意到，我们假设了 \(a_{1}, q\) 都是正数，才自然地写成了对数形式。如果不是正数，是否仍旧可以这么书写？事实上，只要将定比分点中的分母乘到两侧作为幂次即可。例如，\(a_{1}=1, a_{4}=-2, a_{7}=4\) 成等比数列，由于 \(a_{4}<0\) ，所以

\[
\ln a_{4}=\frac{1}{2} \ln a_{1}+\frac{1}{2} \ln a_{7}
\]

虽然无意义，但是可以理解成

\[
\begin{aligned}
& \ln a_{4}^{2}=\ln a_{1}+\ln a_{7} \\
& a_{4}^{2}=a_{1} a_{7} .
\end{aligned}
\]

即 \(a_{4}^{2}=a_{1} a_{7}\) ．所有在等差数列中得到的关于齐次，权重的理解，都可以像这样迁移到等比数列。

\section*{4．5．4 XXXX}
\section*{4.6 概率与统计}
\section*{4．6．1 平均数，方差}
无论是概率论还是统计中，我们都会遇到对于统计量或数学期望的计算：例如平均数（数学期望），方差，标准差，极差等。它们作为反映数据或随机变量的指标，可以量化地体现趋势。根据它们不同的含义，也有其不同的性质。

对于平均数（数学期望）而言，我们有如下性质：

\section*{4．6．1．1 平均数（数学期望）}
对于数据 \(x_{1}, x_{2}, \ldots x_{n}\) ，设 \(\bar{x}\) 是它们的平均数，若所有的数据变为原来的 \(a\) 倍，那么 \(\bar{x}\) 变为 \(a \bar{x}\) 。如何从代数结构的角度来理解这些性质呢？根据平均数的定义，我们有：

\[
\bar{x}=\frac{1}{n} \cdot\left(x_{1}+x_{2}+\cdots+x_{n}\right)
\]

如果我们将 \(\bar{x}\) 看成关于 \(x_{1}, x_{2}, \ldots x_{n}\) 的函数，那么它是一个一次齐次的函数：所有的项 \(\frac{1}{n} x_{1}, \frac{1}{n} x_{2}, \ldots \frac{1}{n} x_{n}\) 都是一次的。这样一来，当自变量变为原来的 \(a\) 倍时，函数值也会变为原来的 \(a^{1}=a\) 倍。因此，我们要明确地意识到，平均数是一个＂一次＂的量。

随机变量的数学期望也是一次式。设随机变量 \(X\) 的分布列为

\begin{center}
\begin{tabular}{|c|c|c|c|c|c|}
\(x\) & \(x_{1}\) & \(x_{2}\) & \(x_{3}\) & \(\ldots\) & \(x_{n}\) \\
\hline
\(\mathbb{P}(X=x)\) & \(p_{1}\) & \(p_{2}\) & \(p_{3}\) & \(\ldots\) & \(p_{n}\) \\
\hline
\end{tabular}
\end{center}

那么根据定义，有

\[
\mathbb{E}(X)=p_{1} x_{1}+p_{2} x_{2}+\cdots+p_{n} x_{n}
\]

我们发现，和平均数一样， \(\mathbb{E}(X)\) 也是一个关于 \(x_{1}, x_{2}, \ldots, x_{n}\) 的一次齐次式。因此，我们有

\[
\mathbb{E}(a X)=a \mathbb{E}(X)
\]

平均数和数学期望除了这个性质，还有一个关于数据平移的性质。对于平均数而言，若所有的数据变为 \(x_{i}+a\) ，那么平均数也从 \(\bar{x}\) 变为 \(\bar{x}+a\) ；对于数学期望而言，\(Y=X+a\) 作为一个新的随机变量，它的期望比 \(X\) 的大 \(a\) ：

\[
\mathbb{E}(Y)=\mathbb{E}(X+a)=\mathbb{E}(X)+a
\]

这意味着：平均数和数学期望是随着数据的平移一同平移的。这从直观上非常好理解：如果全班成绩都加 1 分，那么全班的平均分也会加 1 分。这个性质和刚才的性质，共同构成了期望的＂线性性＂的一部分。这两个式子可以写成一个式子或性质：\\
（1）平均数：设 \(\bar{x}\) 是数据 \(x_{1}, x_{2}, \ldots, x_{n}\) 的平均数，若所有的 \(x_{i}\) 变为 \(a x_{i}+b, i=1,2, \ldots, n\) ，其中 \(a, b\) 为常数，那么 \(\bar{x}\) 变为 \(a \bar{x}+b\) ；\\
（2）数学期望：设 \(X\) 为随机变量，\(a, b\) 为常数，那么

\[
\mathbb{E}(a X+b)=a \mathbb{E}(X)+b
\]

我们发现：当数据进行了线性变换后，平均数也随着需要进行同样的变换。这个性质，可以帮我们快速得到平

均数和数学期望。

\section*{4．6．1．2 方差}
设 \(x_{1}, x_{2}, \ldots, x_{n}\) 的方差为 \(s^{2}\) 。当我们进行同样的操作时，\(s^{2}\) 又会如何变化呢？根据定义，我们知道

\[
s^{2}=\frac{1}{n}\left[\left(x_{1}-\bar{x}\right)^{2}+\left(x_{2}-\bar{x}\right)^{2}+\cdots+\left(x_{n}-\bar{x}\right)^{2}\right]
\]

如果还是将 \(s^{2}\) 看成一个关于 \(x_{1}, x_{2}, \ldots, x_{n}\) 的函数，那么它是几次齐次的呢？事实上，这是一个二次齐次函数，因为： \(\bar{x}\) 是一个一次齐次，所有的 \(x_{i}-\bar{x}\) 也是一次齐次；平方完 \(\left(x_{i}-\bar{x}\right)^{2}\) 则是二次齐次。因此，\(s^{2}\) 表达式中的每一项那都是二次齐次，它也是一个二次齐次式。这样的话，若所有的 \(x_{i}\) 都变为 \(a x_{i}, s^{2}\) 就会变为 \(a^{2} s^{2}\) 。

数学期望道理相同。根据定义，随机变量 \(X\) 的方差为

\[
D(X)=p_{1}\left(x_{1}-\mathbb{E}(X)\right)^{2}+p_{2}\left(x_{2}-\mathbb{E}(X)\right)^{2}+\cdots+p_{n}\left(x_{n}-\mathbb{E}(X)\right)^{2}
\]

注意到 \(\mathbb{E}(X)\) 是一个一次式，所以 \(x_{i}-\mathbb{E}(X)\) 也是一次式，平方完 \(\left(x_{i}-\mathbb{E}(X)\right)^{2}\) 是二次式。故 \(D(X)\) 中的每一项都是二次式，\(D(X)\) 也是一个二次式。因此，随机变量 \(Y=a X\) 的方差为

\[
D(Y)=D\left(a X^{2}\right)=a^{2} D(X)
\]

以上是我们对数据进行放缩后，方差的改变。那么平移呢？事实上，平移方差不变。我们还是用成绩来直观理解。当所有人的成绩都加 1 分后，平均分也加了 1 分，那么＂每个人和平均分的差距＂并没有发生改变。用数学语言说，就是 \(x_{i}-\bar{x}\)（或 \(x_{i}-\mathbb{E}(X)\) ）没有变化。这样一来，方差自然也不发生改变。

把以上两个性质总结成一个式子或一个性质，就是：\\
（1）设 \(s^{2}\) 是 \(x_{1}, x_{2}, \ldots, x_{n}\) 的方差，若所有的数据 \(x_{i}\) 都变为 \(a x_{i}+b\) ，那么方差变为 \(a^{2} s^{2}\) ；\\
（2）设随机变量 \(Y=a X+b\) ，其中 \(X\) 为随机变量，\(a, b\) 为常数。则

\[
D(Y)=D(a X+b)=a^{2} D(X)
\]

这个性质是由方差的二次齐次性和平均数平移的性质得到的。

\section*{4．6．2 成对数据处理}
我们依旧从代数结构的角度来考虑数据点的变化（平移和放缩）将会引起统计参数的哪些变化。

\section*{4．6．2．1 最小二乘法}
考虑成对数据 \(\left(x_{1}, y_{1}\right),\left(x_{2}, y_{2}\right), \ldots,\left(x_{n}, y_{n}\right)\) 。设最小二乘法得到的最佳逼近直线为 \(y=\hat{b} x+\hat{a}\) 。那么有

\[
\hat{b}=\frac{\sum_{i=1}^{n}\left(x_{i}-\bar{x}\right)\left(y_{i}-\bar{y}\right)}{\sum_{i=1}^{n}\left(x_{i}-\bar{x}\right)^{2}}, \hat{a}=\bar{y}-\hat{b} \bar{x}
\]

其中 \(\bar{x}, \bar{y}\) 分别是 \(x_{i}(i=1,2, \ldots, n), y_{i}(i=1,2, \ldots, n)\) 的平均数。系数 \(\hat{b}, \hat{a}\) 分别是关于 \(x_{i}, y_{i}\) 的几次齐次呢？\\
对于 \(\hat{b}\) ，由于平均数是一次齐次式，所以 \(x_{i}-\bar{x}, y_{i}-\bar{y}\) 都是一次齐次，所以 \(\left(x_{i}-\bar{x}\right)\left(y_{i}-\bar{y}\right)\) 是二次齐次。同理，分母是二次齐次，两个二次齐次的式子相除，得到的是零次齐次。故 \(\hat{b}\) 是一个关于 \(x_{i}, y_{i}\) 的零次齐次式。这也符合它的几何含义：它在直线方程中，相当于斜率。斜率是一个零次齐次的量，因为

\[
k=\frac{y_{2}-y_{1}}{x_{2}-x_{1}}
\]

对于 \(\hat{a}\) ，由于 \(\hat{b}\) 是零次齐次，所以 \(\hat{b} \bar{x}\) 是一次齐次式，所以 \(\hat{a}=\bar{y}-\hat{b} \bar{x}\) 是一次齐次。这也和几何意义相符。在直线方程中，\(\hat{a}\) 代表的是截距，它是一个一次齐次的量，因为它是直线和 \(y\) 轴交点的纵坐标。

以上讨论的都是关于 \(x_{i}, y_{i}\) 这 \(2 n\) 个变量的次数。有时候，我们也会需要讨论对于单独 \(x_{i}\) 或 \(y_{i}\) 的次数。\\
对于 \(\hat{b}\) ，它的分子是关于 \(x_{i}\) 的一次式，分母是关于 \(x_{i}\) 的二次式，所以 \(\hat{b}\) 是关于 \(x_{i}\) 的负一次式；它的分子是关于 \(y_{i}\) 的一次式，分母与 \(y_{i}\) 无关，所以 \(\hat{b}\) 是关于 \(y_{i}\) 的一次式。

对于 \(\hat{a} \bar{y}\) 与 \(x_{i}\) 无关，所以它是关于 \(x_{i}\) 的零次式；\(\hat{b}\) 是关于 \(x_{i}\) 的负一次式， \(\bar{x}\) 是关于 \(x_{i}\) 的一次式，所以 \(\hat{b} \cdot \bar{x}\)是关于 \(x_{i}\) 的零次式。因此，\(\hat{a}\) 是关于 \(x_{i}\) 的零次式。而 \(\bar{y} \hat{b}\) 。 \(\bar{x}\) 则都是关于 \(y_{i}\) 的一次式，所以 \(\hat{a}\) 是关于 \(y_{i}\) 的一次式。

我们将所有这些关于次数的信息进行总结：

\begin{center}
\begin{tabular}{|c|c|c|c|}
\hline
变量 & \(x_{i}, y_{i}\) & \(x_{i}\) & \(y_{i}\) \\
\hline
\(\hat{a}\) & 1 & 0 & 1 \\
\hline
\(\hat{b}\) & 0 & -1 & 1 \\
\hline
\end{tabular}
\end{center}

有了这些理解，我们很快地判断当 \(x_{i}, y_{i}\) 发生变化时，\(\hat{b}, \hat{a}\) 如何变化。\\
（1）当 \(x_{i}\) 变为原来的 \(a\) 倍时，\(\hat{b}\) 关于 \(x_{i}\) 是 -1 次，所以变为 \(\frac{1}{a} \cdot \hat{b} ; \hat{a}\) 关于 \(x_{i}\) 是 0 次，所以不变。从几何的角度，也可以理解：\(x\) 在斜率分母上，所以 \(x\) 变为 \(a x\) 时，斜率 \((\hat{b})\) 变为原来的 \(\frac{1}{a}\) 。当横坐标放缩时，截距不改变，所以 \(\hat{a}\) 不变。\\
（2）当 \(y_{i}\) 变为原来的 \(a\) 倍时，\(\hat{b}\) 的关于 \(y_{i}\) 是 1 次，所以变为 \(a \cdot \hat{b} ; \hat{a}\) 关于 \(y_{i}\) 是 1 次，所以变为 \(a \cdot \hat{a}\) 。从几何的角度理解：\(y\) 在斜率的分子上，所以 \(y\) 变为 \(a y\) 时，斜率变为原来的 \(a\) 倍。当纵坐标放缩时，截距直接随着其发生同样的改变，所以 \(\hat{a}\) 变为原来的 \(a\) 倍。

再利用平均数关于平移的性质，我们可以得到：\\
（1）当 \(x_{i}\) 都变为 \(x_{i}+a\) 时，由于 \(\hat{b}\) 中 \(x_{i}\) 都以 \(x_{i}-\bar{x}\) 的形式出现，所以不发生改变；而随着变为 \(\bar{x}+a\) 因此 \(\hat{a}\) 变为

\[
\bar{y}-\hat{b} \cdot(\bar{x}+a)=\bar{y}-\hat{b} \cdot \bar{x}-\hat{b} \cdot a=\hat{a}-\hat{b} \cdot a
\]

即 \(\hat{a}\) 要减少 \(\hat{b} \cdot a\) 。这从图像平移角度也可以理解：\(x_{i}\) 都变为 \(x_{i}+a\) ，相当于图像向右平移 \(a\) 个单位。平移不改变斜率，所以 \(\hat{b}\) 不变；若正相关，则截距也会相应减少，以斜率为速度减少，所以减少 \(\hat{b} \cdot a\) 。\\
（2）若 \(y_{i}\) 变为 \(y_{i}+a\) ，注意到 \(\hat{b}\) 中的 \(y_{i}\) 也都是以 \(y_{i}-\bar{y}\) 的形式出现的，所以不发生改变；由于 \(\bar{y}\) 变为 \(\bar{y}+a\) ，所以 \(\hat{a}\) 变为

\[
\bar{y}+a-\hat{b} \cdot \bar{x}=\hat{a}+a
\]

故增加了 \(a\) 。从图像平移的角度看：\(y_{i}\) 都变为 \(y_{i}+a\) ，相当于图像向上平移 \(a\) 个单位。平移不改变斜率，所以 \(\hat{b}\) 不变；向上平移，截距跟着发生同样的变化，所以增加 \(a\) 。

注意，以上所有的讨论都是建立在所有的 \(x_{i}\) 或所有的 \(y_{i}\) 一起发生同样的改变后，产生的结果。如果只是改变其中的若干个 \(x_{i}, y_{i}\) ，或 \(x_{i}\) 之间发生了不一样的改变，那么对于 \(\hat{a}, \hat{b}\) 的影响是无法确定的。因此，考虑这些问题时要注意适用范围。

\section*{4．6．2．2 相关系数}
成对数据 \(\left(x_{1}, y_{1}\right),\left(x_{2}, y_{2}\right), \ldots,\left(x_{n}, y_{n}\right)\) 的相关系数为

\[
r=\frac{\sum_{i=1}^{n}\left(x_{i}-\bar{x}\right) \cdot\left(y_{i}-\bar{y}\right)}{\sqrt{\sum_{i=1}^{n}\left(x_{i}-\bar{x}\right)^{2} \cdot \sum_{i=1}^{n}\left(y_{i}-\bar{y}\right)^{2}}}
\]

首先，可以观察到 \(r\) 是一个齐次式。分子关于 \(x_{i}, y_{i}\) 分别都是 1 次，关于所有的 \(x_{i}, y_{i}\) 这 \(2 n\) 个变量是 2 次。分母也一样，所以 \(r\) 是一个零次齐次式。

如果我们知道了 \(r\) 是齐次式，有没有方法不去计算具体次数，直接看出来它是零次的呢？事实上，根据柯西不等式，我们知道 \(r \in[-1,1]\) 。换句话说，\(r\) 是有界的。然而，只有零次齐次的量是有界的。所有其他次数的齐次量都是无界的！例如，我们考虑一次齐次的量，那么当所有自变量变为原来的 \(10^{100}\) 倍时，代数式的值也会跟着变为 \(10^{100}\) ，它是无界的。因此，当我们知道齐次式有界时，它必然是零次齐次。除了相关系数，正弦函数，余弦函数也具有同样的性质。

\section*{4．6．3 统计数据计算}
除了统计量的代数结构，我们还会遇到具体的统计数据的计算。例如，有如下一组数据，我们想要求出它们的平均数：

\[
189.3,189.9,191.1,190.2,190.0,189.7,188.8,189.9,190.2,190.9
\]

很多同学直接就将这 10 个数据相加，再除以 10 去计算。然而，考试时时间有限，这么算需要算多次的多位数加法，非常不划算。一个正确且更快的方法是：我们去看这些数据和＂某个基准＂的差距。我们观察到，这些数据都和 190 非常接近，所以我们可以把它们都和 190 进行比较：

\[
\begin{aligned}
& 189.3=190.0-0.7,189.9=190.0-0.1 \\
& 191.1=190.0+1.1,190.2=190.0+0.2 \\
& 190.0=190.0+0.0,189.7=190.0-0.3 \\
& 188.8=190.0-1.2,189.9=190.0-0.1 \\
& 190.2=190.0+0.2,190.0=190.0+0.9
\end{aligned}
\]

我们只需要去计算这些数据和 190.0 的差距的和，就可以得到平均数了。注意到：

\[
-0.7-0.1+1.1+0.2+0.0-0.3-1.2-0.1+0.2+0.9=0
\]

所以它们的平均数就是 190．0．这是因为：这 10 个数的和，和 \(190 \times 10\) 的差距，就是 0.0 。所以它们的平均数就是 190.0 这样的话，我们就不需要重复四位数加减法，只需要简单做这种小数计算即可。

这个方法其实利用了期望的线性性。当对所有数据选取了一个基准 \(a\) 后，我们只需要关心每个数据和 \(a\) 的差距，就可以得到平均数。写作数学表达式就是：

\[
\mathbb{E}(X)=a+\mathbb{E}(X-a)
\]

在这个例子中，\(a=190.0\) 。\\
除了可以利用平移的性质，还可以利用一次的性质。例如，如果所有的数据都是两位小数，计算上有一定困难：

\[
17.12,17.09,16.99,17.11,17.02,17.21,17.13,17.08,17.00,17.25
\]

我们可以把它们都化成整数或者一位小数。例如，如果将它们都化作整数，相当于给每个数据都乘 100 ：

\[
1712,1709,1699,1711,1702,1721,1713,1708,1700,1725 .
\]

我们选 1700 为基准，误差的和为

\[
12+9-1+11+2+21+13+8+0+25=100
\]

那么，平均每个数获得了 \(\frac{100}{10}=10\) ，所以平均数为 1710 。注意，这里不能直接将 100 加回到 1700 ，而应去计算 100 平均给这十个数据是多少。这是因为： 100 是所有数据 \(X-1700\) 的和，而不是平均数，我们要计算的是\\
\(\mathbb{E}(X-1700)\) ，所以还需要除以 10 。因此，乘 100 后，平均数为 1710 。那么原来数据的平均数为 17.10 。这里，我们则用到了

\[
\mathbb{E}(a X)=a \mathbb{E}(X)
\]

这个性质。

\section*{4．6．4 XXXX}
\section*{4.7 导数}
\section*{4．7．1 求导法则的理解}
我们在学习求导法则时，学过乘法法则和除法法则：

\[
\begin{aligned}
& (f(x) g(x))^{\prime}=f(x) g^{\prime}(x)+f^{\prime}(x) g(x) \\
& \left(\frac{f(x)}{g(x)}\right)^{\prime}=\frac{f^{\prime}(x) g(x)-f(x) g^{\prime}(x)}{g^{2}(x)}
\end{aligned}
\]

事实上，这两个求导法则是一个法则，只需要借助复合函数求导法则即可。这是因为：我们可以将除以一个数，看成乘一个数的倒数，即：

\[
\frac{f(x)}{g(x)}=f(x) \cdot \frac{1}{g(x)}
\]

因此，

\[
\left(\frac{f(x)}{g(x)}\right)^{\prime}=f^{\prime}(x) \cdot \frac{1}{g(x)}+f(x) \cdot\left(\frac{1}{g(x)}\right)^{\prime}
\]

根据复合函数求导法，\(y=\frac{1}{g(x)}\) 可以看作 \(y=\frac{1}{u}, u=g(x)\) 的复合，所以

\[
\left(\frac{1}{g(x)}\right)^{\prime}=y_{u}^{\prime}(u) \cdot u_{x}^{\prime}(x)=-\frac{1}{g^{2}(x)} \cdot g^{\prime}(x)=-\frac{g^{\prime}(x)}{g^{2}(x)}
\]

故

\[
f^{\prime}(x) \cdot \frac{1}{g(x)}+f(x) \cdot\left(\frac{1}{g(x)}\right)^{\prime}=\frac{f^{\prime}(x)}{g(x)}-\frac{f(x) g^{\prime}(x)}{g^{2}(x)}=\frac{f^{\prime}(x) g(x)-f(x) g^{\prime}(x)}{g^{2}(x)}
\]

这样一来，对除法法则的理解也就更为透彻，也解释了为什么除法法则的分子上有一个负号：它来自 \(y=\frac{1}{u}\) 的导数 \(-\frac{1}{u^{2}}\) 。知道了原理，就再也不会记混淆了。甚至，只需要知道乘法法则和复合函数求导法则，就可以推导除法法则了。

乘法法则和除法法则的等价性，其实和四则运算中乘法，除法的等价性是一回事。我们都是将＂除以一个数＂看成＂乘这个数的倒数＂。在实际运算中，我们也可以利用这个方法加快我们的运算。

例如，我们经常会需要计算形如 \(\frac{f(x)}{e^{x}}\) 形式的导数。这里 \(f(x)\) 可能是一个多项式，可能是一个二次函数，或者是其他非指数函数。例如，\(f(x)=x^{2}+3 x+1\) ．那么

\[
\begin{aligned}
\left(\frac{x^{2}+3 x+1}{e^{x}}\right)^{\prime} & =\frac{(2 x+3) e^{x}-\left(x^{2}+3 x+1\right) e^{x}}{e^{2 x}}=\frac{\left(-x^{2}-x+2\right) e^{x}}{e^{2 x}} \\
& =\frac{-x^{2}-x+2}{e^{x}}
\end{aligned}
\]

在使用除法法则时，我们最后还需要多进行一次约分。但如果直接使用乘法法则，则免去了这次约分。我们把除

以 \(e^{x}\) 看作乘 \(e^{-x}\) ，那么：

\[
\left(\left(x^{2}+3 x+1\right) e^{-x}\right)^{\prime}=(2 x+3) e^{-x}+\left(x^{2}+3 x+1\right)\left(e^{-x}\right)^{\prime}=\left(-x^{2}-x+2\right) e^{-x}
\]

可以直接写出结果。\\
除了乘法法则和除法法则的等价性，对数函数，指数函数的求导法则也可以有它自己的理解。我们知道，以 \(a\) 为底的指数函数，对数函数的导数为：

\[
\begin{aligned}
& \left(a^{x}\right)^{\prime}=\ln a \cdot a^{x} \\
& \left(\log _{a} x\right)^{\prime}=\frac{1}{\ln a} \cdot \frac{1}{x}
\end{aligned}
\]

对于 \(a \neq e\) 的情况，我们使用较少，考试的时候经常记混。那如何解决这个问题呢？只需要将它化成 \(a=e\) 的情况即可。\\
（1）对于指数函数，注意到 \(a=e^{\ln a}\) ，所以

\[
a^{x}=\left(e^{\ln a}\right)^{x}=e^{x \ln a}
\]

由于 \(\ln a\) 是一个常数，根据复合函数求导法则，有：

\[
\left(a^{x}\right)^{\prime}=\left(e^{x \ln a}\right)^{\prime}=\ln a \cdot e^{x \ln a}=\ln a \cdot a^{x}
\]

（2）对于对数函数，根据换底公式，我们把它换成自然对数：

\[
\log _{a} x=\frac{\ln x}{\ln a}
\]

又因为 \(\ln a\) 是一个常数，所以

\[
\left(\log _{a} x\right)^{\prime}=\left(\frac{\ln x}{\ln a}\right)^{\prime}=\frac{1}{\ln a} \cdot \frac{1}{x}
\]

这样一来，我们只需要知道 \(e^{x}, \ln x\) 的求导法则，就可以求出所有指数函数，对数函数的导数了，进一步减少了死记硬背的负担。同时，我们也感受到了数学知识之间紧密的联系：乘法和除法，指数和对数，不同底数间的转化，它们可以在导数中发挥它们自己的作用。

我们之所以洞察到了这些导数知识之间的联系，就是因为我们使用了第三章讲过的＂化归＂的思想。化归指的是：将不熟悉的问题转化为熟悉的问题。例如，人脑是先学习乘法，后学习除法的，且天然对乘法更亲切。所以，我们通常选择将除法向除法转化；再例如，以 \(e\) 为底的指数，对数是相对常用的，其他的指数和对数如果也能向其转化，就可以获得新的理解。我们尽可能将更多的知识都通过＂化归＂建立新的联系，来加深理解。

\section*{4．7．2 积分因子}
在模拟卷中，曾经非常流行的一种题型（也被称为＂构造＂），是给一个有关函数 \(f(x)\) 和它导数 \(f^{\prime}(x)\) 的一个不等式，例如：

\[
f(x)+f^{\prime}(x)>0
\]

根据这个，进行推理。我们知道，如果我们得到的信息是

\[
f^{\prime}(x)>0
\]

那么可以利用单调性求解；但实际上，还有 \(f(x)\) 这一项，无法求解。因此，我们考虑给不等式乘一个新的函数 \(g(x)\) ，使得它刚好可以写成某个新的函数 \(F(x)\) 的导数的形式。这样，就可以利用单调性求解了。

一些教辅的参考答案会给出这样的做法：＂令 \(F(x)=e^{x} f(x)\) ．则

\[
F^{\prime}(x)=\left(f(x)+f^{\prime}(x)\right) e^{x}>0
\]

故 \(F(x)\) 单调递增，故 \(\cdots\) 然而，问题是，我们怎么从题目的不等式，想到这样的构造呢？\\
问题的关键，就是要找到这个 \(e^{x}\) 。事实上，我们假设不知道我们要乘的是 \(e^{x}\) ，而是一个新的函数 \(g(x)\) 。不妨设它是正的（或至少在某一区间上是正的），那么在不等式两侧同时乘 \(g(x)\) ，可以得到：

\[
g(x) f(x)+g(x) f^{\prime}(x)>0
\]

为了让它刚好是某个函数 \(F(x)\) 的导数，我们马上就联想到了导数的乘法法则。如果我们的不等式刚好是

\[
g^{\prime}(x) f(x)+g(x) f^{\prime}(x)>0
\]

的形式，那么它就是 \(f(x) g(x)\) 的导数了！将原来的不等式和乘法法则对比，发现：\(f(x)\) 的系数不同，一个是 \(g(x)\) ，一个是 \(g^{\prime}(x)\) 。因此，我们大胆假设：如果 \(g(x)=g^{\prime}(x)\) ，问题就解决了！这样一来：

\[
g(x) f(x)+g(x) f^{\prime}(x)=g^{\prime}(x) f(x)+g(x) f^{\prime}(x)=(g(x) f(x))^{\prime}>0
\]

那么，什么样的函数满足 \(g(x)=g^{\prime}(x)\) 呢？根据我们的经验，\(g(x)=e^{x}\) 满足。这就是为什么想到去构造 \(F(x)=\) \(e^{x} f(x)\) 。

我们再考察一个相对复杂的例子。若对于 \(x \in \mathbb{R}\) ，均有

\[
2 x f(x)+f^{\prime}(x)>0
\]

应如何构造函数？我们还是想在不等式两侧乘函数 \(g(x)>0\) ：

\[
2 x g(x) f(x)+g(x) f^{\prime}(x)>0
\]

和导数的乘法法则作比较：

\[
g^{\prime}(x) f(x)+g(x) f^{\prime}(x)>0
\]

所以，我们只需要要求：

\[
\begin{aligned}
& 2 x g(x)=g^{\prime}(x) \\
& \frac{g^{\prime}(x)}{g(x)}-2 x=0
\end{aligned}
\]

即可。那么，什么样的函数满足这个条件呢？我们观察到，左侧刚好可以写成一个函数的导数：

\[
\frac{g^{\prime}(x)}{g(x)}=(\ln g(x))^{\prime}, 2 x=\left(x^{2}\right)^{\prime}
\]

故

\[
\begin{aligned}
& \left(\ln g(x)-x^{2}\right)^{\prime}=0 \\
& \left(\ln \left[g(x) e^{-x^{2}}\right]\right)^{\prime}=0
\end{aligned}
\]

什么函数的导数是 0 呢？只有常数函数的导数为 0 。因此， \(\ln \left[g(x) e^{-x^{2}}\right]\) 是一个常数，\(g(x) e^{-x^{2}}\) 也是一个常数。所以可以选择 \(g(x)=e^{x^{2}}\) ．

我们来验证一下。令 \(F(x)=e^{x^{2}} f(x)\) ．那么

\[
F^{\prime}(x)=2 x e^{x^{2}} f(x)+e^{x^{2}} f^{\prime}(x)=e^{x^{2}}\left(2 x f(x)+f^{\prime}(x)\right)>0
\]

故 \(F(x)\) 单调递增。我们成功构造出了函数。\\
事实上，这里的运算已经涉及到微分方程的内容，我们不再展开，但是希望读者理解这里进行推理的指导

思想：把不熟悉的表达式，向导数的乘法法转化，这本质上也是一种化归。转化的过程中，需要待定系数，在左右两侧乘一个未知的函数 \(g(x)\) ，根据乘法法则求解它。我们左右两侧乘的因子 \(g(x)\) ，就称作积分因子。在实际备考中，不需要浪费过多时间准备它。

\section*{4.7 .3 消元思想的应用}
我们在前面学习过的所有代数结构，处理导数时都会发挥它们相应的作用。在这里，我们借助极值点偏移的一道经典例题来体会。

例 4.1 设 \(f(x)=e^{x}-a x\) ．若 \(f(x)=0\) 有两个实根 \(x_{1}, x_{2}\) ，求证：\(x_{1}+x_{2}>2\) ．\\
解 由于 \(f(x)=0\) 的两根为 \(x_{1}, x_{2}\) ，所以

\[
\begin{aligned}
& e^{x_{1}}-a x_{1}=0 \\
& e^{x_{2}}-a x_{2}=0
\end{aligned}
\]

要证明的是 \(x_{1}+x_{2}\) ，和 \(a\) 没有直接关系，所以考虑消掉 \(a\) ．注意到

\[
e^{x_{1}}=a x_{1}, e^{x_{2}}=a x_{2}
\]

所以，可以将两个式子作比，得到

\[
\frac{e^{x_{1}}}{e^{x_{2}}}=\frac{x_{1}}{x_{2}}
\]

令 \(t=\frac{x_{1}}{x_{2}}\) ，那么 \(e^{x_{1}-x_{2}}=t\) ．故 \(x_{1}-x_{2}=\ln t\) ．和 \(x_{1}=t x_{2}\) 联立，可得

\[
x_{2}=\frac{\ln t}{t-1}, x_{1}=\frac{t \ln t}{t-1}
\]

因此，只需要证明：

\[
x_{1}+x_{2}=\frac{(t+1) \ln t}{t-1}>2
\]

即可。为了摆脱分式，我们考虑函数 \(F(t)=(t+1) \ln t-2 t+2\) ．求导得

\[
\begin{aligned}
& F^{\prime}(t)=\frac{t+1}{t}+\ln t-2=\frac{1}{t}+\ln t-1 \\
& F^{\prime \prime}(t)=\frac{1}{t}-\frac{1}{t^{2}}=\frac{t-1}{t^{2}}
\end{aligned}
\]

因此，\(F^{\prime}(t)\) 在 \(t=1\) 时取最小值，故 \(F^{\prime}(t) \geq F^{\prime}(1)=0\) ．故 \(F^{\prime}(t) \geq 0, F(t)\) 单调递增。故：\\
（1）当 \(t>1\) 时，\(F(t)>F(1)=0\) ．故

\[
\begin{aligned}
& (t+1) \ln t>2(t-1) \\
& \frac{(t+1) \ln t}{t-1}>2
\end{aligned}
\]

（2）当 \(t<1\) 时，\(F(t)<F(1)=0\) ．故

\[
\begin{aligned}
& (t+1) \ln t<2(t-1) \\
& \frac{(t+1) \ln t}{t-1}>2
\end{aligned}
\]

由于 \(x_{1} \neq x_{2}\) ，所以 \(t=e^{x_{1}-x_{2}} \neq 1\) ．故不需要考虑 \(t=1\) 的情况。故命题得证。

这里面涉及到了一个非常巧妙的技巧：比例换元。我们令 \(\frac{x_{1}}{x_{2}}=t\) ，然后一切都顺理成章的解决了。但问题时，为什么我们要这么处理呢？不这么处理，有没有别的方法？

首先，我们从自由度的角度分析这个问题。思考：\(x_{1}, x_{2}\) 这两个实数根，是由谁决定的呢？很显然，当参数 \(a\)确定后，\(x_{1}, x_{2}\) 就都确定了。因此，整个系统其实只有一个自由度。这个自由度可以用 \(a\) 表示，也可能可以用别的字母表示。

我们更清晰地理解这个系统，我们能不能让参数 \(a\) 变起来？当它和变量混杂在一起时，直接变化会影响函数图像。所以，我们先进行参变分离：

\[
f(x)=0 \Rightarrow \frac{1}{a}=\frac{x}{e^{x}}
\]

换言之，\(x_{1}, x_{2}\) 就是函数 \(g(x)=\frac{x}{e^{x}}\) 和直线 \(y=\frac{1}{a}\) 的两个交点的横坐标。那么，我们让 \(a\) 动起来，自然就会看到 \(x_{1}, x_{2}\) 也会发生改变。

在改变的过程中，我们接着思考：既然这个系统只有一个自由度，我们还能用哪个变量来表示它呢？根据刚才的解答，我们知道：比例换元 \(t=\frac{x_{1}}{x_{2}}\) 最终成功表示了 \(x_{1}, x_{2}\) 。因此，\(t\) 也可以作为这个自由度。

为什么呢？不妨设 \(x_{1}<x_{2}\) 。仔细观察，我们会发现，当 \(a\) 变化时，比例 \(\frac{x_{1}}{x_{2}}\) 也在跟随着发生给变化。而且，如果这个比例确定了，\(a\) 的值也会被确定。换言之，如果我们知道了两个横坐标的比值，也可以确定这个系统！我们还会发现，这个比值 \(t\) 随着 \(a\) 的增大而减小。这个单调性，更保证了 \(t\) 作为自由度的合理性。

从几何上，我们理解了 \(t\) 作为自由度的合理性。从代数上是否也能看出来呢？事实上，根据

\[
e^{x_{1}-x_{2}}=\frac{x_{1}}{x_{2}}
\]

令 \(t=\frac{x_{1}}{x_{2}}\) 虽然是自然的，但是我们也能从齐次结构来理解。如果两边同时取对数，就可以得到

\[
x_{1}-x_{2}=\ln \left(\frac{x_{1}}{x_{2}}\right)
\]

左侧 \(x_{1}-x_{2}\) 是一个一次齐次式，而 \(\frac{x_{1}}{x_{2}}\) 是零次齐次，所以 \(\ln \left(\frac{x_{1}}{x_{2}}\right)\) 也是。左右两侧次数不齐，所以是非齐次式。对于有两个变量的系统，我们只需要一个齐次式，一个非齐次式就可以表示所有的变量了。因此，我们只需要确定 \(t=\frac{x_{1}}{x_{2}}\) ，问题就解决了。所以，从齐次结构角度分析，令 \(t=\frac{x_{1}}{x_{2}}\) 也是合理的。

除了利用比值换元，我们还有其他方法吗？既然我们注意到了这是只有一个自由度的系统，我们是否可以选择 \(x_{1}, x_{2}\) 中之一，作为变量去表示其他的量呢？如果我们再去仔细考察结构会发现，\(x_{1}+x_{2}\) 是一个对称式。因此，可以不妨设 \(x_{1}<x_{2}\) 。一般而言，我们是不轻易破坏对称结构的。但是在这道题中，我们可以大胆尝试，将它变形为

\[
x_{2}<2-x_{1}
\]

从几何上看，\(x_{1}\) 确实可以确定 \(x_{2}\) 。当小根确定后，我们作 \(x\) 轴的垂线，平行线，就可以自然找到 \(x_{2}\) 的位置。所以，也可以选 \(x_{1}\) 或 \(x_{2}\) 作为自由度。

要知道，虽然 \(x_{2}\) 是可以用 \(x_{1}\) 表示的，但是我们无法写出一个显式的表达式，即 \(x_{2}=h\left(x_{1}\right)\) 的形式。如果写不出来，就达不到消元的目的，无法得到最终结果。

因此，为了消元，我们必须要借助于 \(x_{1}, x_{2}\) 的＂定义＂，或者说来源：它们是 \(\frac{x}{e^{x}}=\frac{1}{a}\) 的零点。令 \(g(x)=\frac{x}{e^{x}}\) ，所以

\[
g\left(x_{1}\right)=g\left(x_{2}\right)
\]

即

\[
\frac{x_{1}}{e^{x_{1}}}=\frac{x_{2}}{e^{x_{2}}}
\]

虽然我们无法直接将 \(x_{2}\) 用 \(x_{1}\) 表示，但是如果我们的 \(x_{2}\) 都能够以 \(g\left(x_{2}\right)\) 的形式出现，问题也就解决了。

为此，我们分析 \(g(x)\) 的性态。

\[
g^{\prime}(x)=\left(x e^{-x}\right)=e^{-x}-x e^{-x}=(1-x) e^{-x}
\]

故 \(g(x)\) 的极大值点为 1 ，在 \((-\infty, 1)\) 上单增，在 \((1,+\infty)\) 上单减。由于 \(g\left(x_{1}\right)=g\left(x_{2}\right)\) ，所以必然有

\[
x_{1}<1<x_{2}
\]

因此，

\[
2-x_{1}>1
\]

所以， \(2-x_{1}, x_{2}\) 都在 \(g(x)\) 的单调递减区间上。\\
如果两个横坐标都在单调递减区间上，我们可以怎么比较它们的大小呢？只要比较函数值就可以了。这样，我们就把 \(x_{2}\) 成功以 \(g\left(x_{2}\right)\) 的形式进行处理：

\[
2-x_{1}<x_{2} \Leftrightarrow g\left(2-x_{1}\right)>g\left(x_{2}\right)=g\left(x_{1}\right)
\]

故只需证明：

\[
g\left(2-x_{1}\right)>g\left(x_{1}\right)
\]

显然，我们的 \(x_{1}\) 是可以取遍 \((0,1)\) 的，所以我们可以考虑函数

\[
G(x)=g(2-x)-g(x), x \in(0,1)
\]

只需证明 \(G(x)>0\) 即可。\\
到这里，我们已经成功消元。只需要研究 \(G(x)\) 的性态即可。虽然我们需要求导，但是甚至不需要写出 \(G(x)\)的表达式。这是由求导法则决定的：

\[
\begin{aligned}
& G^{\prime}(x)=(g(2-x))^{\prime}-g^{\prime}(x) \\
& (g(2-x))^{\prime}=-g^{\prime}(2-x)
\end{aligned}
\]

所以：

\[
\begin{aligned}
G^{\prime}(x) & =-g^{\prime}(2-x)-g(x) \\
& =-(x-1) e^{x-2}-(1-x) e^{-x} \\
& =(1-x)\left(e^{x-2}-e^{-x}\right) \\
& =(1-x) e^{-x}\left(e^{2 x-2}-1\right)
\end{aligned}
\]

由于 \(x \in(0,1)\) ，所以

\[
1-x>0, e^{-x}>0, e^{2 x-2}<e^{0}=1
\]

所以 \(G^{\prime}(x)<0, G(x)\) 在 \((0,1)\) 上单调递减，所以

\[
G(x)>G(1)=g(1)-g(1)=0 .
\]

问题也解决了。\\
第二个方法，从自由度的角度，选取了 \(x_{1}\) 为自由度。而第一个方法，选取了 \(t\) 为自由度。到这里，这个题目的所有结构，我们都研究清楚了：\\
（1）自由度：所有字母 \(x_{1}, x_{2}, t, a\) 都可以作为自由度，决定这个系统。其中，选择 \(t\) ，会得到方法一；选择 \(x_{1}\) 或 \(x_{2}\) ，获得到方法二；（问题：是否可以选择 \(a\) ）\\
（2）消元思想：无论是方法一还是方法二，关键的指导思想都是消元。我们只有将所求变为一个单变量函数的不

等式，才有可能借助高中阶段关于函数的工具完成证明。无论是 \(t=\frac{x_{1}}{x_{2}}\) ，并用 \(t\) 表示 \(x_{1}, x_{2}\) ，还是用不等关系 \(x_{1}>2-x_{2}\) ，都是在消元。因此，当我们分析好自由度后，就可以选取适当的自由变量去消元了；\\
（3）齐次结构：\(t=\frac{x_{1}}{x_{2}}\) 之所以可以表达 \(x_{1}, x_{2}\) ，就是因为

\[
e^{x_{1}-x_{2}}=\frac{x_{1}}{x_{2}}
\]

是一个非齐次结构。我们通过确定比值（零次齐次），两个变量也就确定了，消元完成；\\
（4）对称结构：所求 \(x_{1}+x_{2}\) 是一个对称式。但是这里，我们毫不犹豫地破坏了对称式。这是因为：\(x_{1}+x_{2}\) 作为对称结构的意义是，去表示其他对称式。这里并没有对称式需要我们表示，所以可以直接破坏掉。更重要的原因是，我们需要 \(x_{1}\) 或 \(x_{2}\) 分离开，这样才能借助它们＂零点＂的定义，和函数产生联系。如果总是 \(x_{1}+x_{2}\) 一起出现，那么就用不到 \(f\left(x_{1}\right)=f\left(x_{2}\right)=0\) 这个条件了。利用上这个条件，我们也完成了消元。

\section*{4．7．4 XXXX}
\section*{4.8 圆雉曲线}
\section*{4．8．1 离心率问题}
有一类常见的小题，是求解离心率。求解的时候，我们会遇到关于 \(a, b, c\) 的方程。这些方程有哪些特点呢？需要几个这样的方程才能够确定离心率呢？

以椭圆为例，根据定义，有 \(a^{2}=b^{2}+c^{2}\) 。这是一个关于 \(a, b, c\) 的二次齐次式。确定离心率，就是确定 \(\frac{c}{a}\)的值。当确定了 \(\frac{c}{a}\) 的值，就等价于确定了 \(a, b, c\) 三者的比例。这是因为：

\[
\frac{b^{2}}{a^{2}}=\frac{a^{2}-c^{2}}{a^{2}}=1-\frac{c^{2}}{a^{2}}=1-e^{2}
\]

故 \(\frac{b}{a}=\sqrt{1-e^{2}}\) 。因此，\(\frac{b}{a}, \frac{c}{a}\) 和其他任意的 \(a, b, c\) 三者中两者的比值，都可以确定 \(a, b, c\) 三者的比例。从自由度的角度也很好理解：如果确定了 \(a, b, c\) 中任意两者的比值，相当于获得了一个齐次条件；\(a, b, c\) 默认成立 \(a^{2}=b^{2}+c^{2}\) ，所以已经有两个齐次条件了。根据我们自由度的理论，两个齐次条件就可确定三个未知数的比例关系。因此，我们只需要知道任意两者的比例即可。

这样一来，我们在处理离心率小题时，只需要得到一个关于 \(a, b, c\) 的比例关系即可。这个比例关系往往来自于 \(a, b, c\) 的齐次式。例如，当我们得到如下关系式时，都可以确定离心率：\\
（1）\(a^{2}-3 a c+2 c^{2}=0\) 。注意到这是一个关于 \(a, c\) 的二次齐次式，可以直接分解因式或写成关于 \(e=\frac{c}{a}\) 的方程。离心率中 \(a\) 在分母，所以左右两侧同时除以 \(a^{2}\) ，可得：

\[
\begin{aligned}
& 1-\frac{3 a c}{a^{2}}+\frac{2 c^{2}}{a^{2}}=0 \\
& 2 e^{2}-3 e+1=0 \\
& (2 e-1)(e-1)=0
\end{aligned}
\]

解得 \(e_{1}=1, e_{2}=\frac{1}{2}\) ．由于椭圆离心率 \(e<1\) ，所以 \(e=\frac{1}{2}\) ．\\
（2）\(b^{2}=a c\) 。在（1）中，我们已知的是只关于 \(a, c\) 的齐次方程，因此直接除即可；此处，我们得到的是关于\\
\(a, b, c\) 的方程，所以需要将 \(b\) 消元再处理。根据 \(b^{2}=a^{2}-c^{2}\) ，有

\[
\begin{aligned}
& a^{2}-c^{2}=a c \\
& a^{2}-a c-c^{2}=0
\end{aligned}
\]

这同样是一个关于 \(a, c\) 的二次齐次式，左右两侧同时除以 \(a^{2}\) ，可得

\[
\begin{aligned}
& 1-\frac{c}{a}-\frac{c^{2}}{a^{2}}=0 \\
& e^{2}+e-1=0
\end{aligned}
\]

解得 \(e_{1}=\frac{-1-\sqrt{5}}{2}, e_{2}=\frac{-1+\sqrt{5}}{2}\) ．由于 \(e>0\) ，所以 \(e=\frac{\sqrt{5}-1}{2}\) ．\\
（3） \(2 \sqrt{3} c^{2}=a b\) 。这同样是关于 \(a, b, c\) 的二次齐次式，所以需要将 \(b\) 消元。但是，这里只有 \(b\) ，所以可以平方后再消元。

\[
\begin{aligned}
& 12 c^{4}=a^{2} b^{2} \\
& 12 c^{4}=a^{2}\left(a^{2}-c^{2}\right) \\
& a^{4}-a^{2} c^{2}-12 c^{4}=0
\end{aligned}
\]

这是一个关于 \(a, c\) 的四次齐次式。很多同学在这里担心，可能无法求解四次方程；但事实上不用担心，我们两边除以 \(a^{4}\) ，可得

\[
\begin{aligned}
& 1-\frac{c^{2}}{a^{2}}-\frac{12 c^{4}}{a^{4}}=0 \\
& 12 e^{4}+e^{2}-1=0
\end{aligned}
\]

它虽然是关于 \(e\) 的四次方程，但它是＂假四次，真二次＂：它没有 \(e^{3}, e\) 这些奇数次项，所以我们可以将它看成关于 \(e^{2}\) 的二次方程。分解因式得

\[
\left(4 e^{2}-1\right)\left(3 e^{2}+1\right)=0
\]

由于 \(3 e^{2}+1>0\) ，所以 \(4 e^{2}-1=0\) ．故 \(e=\frac{1}{2}\) ．\\
可以看到，当我们掌握了齐次结构后，这些方程都会迎刃而解。所以，只需要将题目中的几何条件翻译成合适的齐次方程即可。

\section*{4．8．2 焦点三角形}
很多解析几何中的题目都来自焦点三角形。焦点三角形指的是以圆雉曲线上一点 \(P\) 和两个焦点 \(F_{1}, F_{2}\) 组成的三角形。有非常多关于焦点三角形的二级结论，但是我们在理解这个问题的时候，可以从新的角度出发来理解。

当我们椭圆给定时，例如 \(\frac{x^{2}}{4}+\frac{y^{2}}{3}=1\) ，那么这个焦点三角形只有一个 \(P\) 点是自由的。 \(P(x, y)\) 点有两个坐标，同时它在椭圆上，所以本质上只有一个自由度。因此，想要确定这个焦点三角形，我们只需要一个额外条件即可。

当我们需要求解椭圆时（即 \(a, b\) 均未知），那么我们需要更多的条件。\\
（1）如果 \(a, b\) 中知道一个，那么算上点 \(P\) 坐标，共有两个自由度。此时，需要两个条件；\\
（2）如果 \(a, b\) 均未知，那么算上 \(P\) 点坐标，共有三个自由度。此时，需要三个条件。\\
我们从如下几个例子，来体会自由度在这些问题中的应用。

4．8．2．1 已知两个的情形（圆雉曲线确定）\\
已知椭圆 \(\frac{x^{2}}{4}+y^{2}=1\) ．若 \(\angle F_{1} P F_{2}=90^{\circ}\) ，且 \(\left|P F_{1}\right|>\left|P F_{2}\right|\) ，求 \(P\) 点坐标。由于粗圆确定，所以这个问题只有一个自由度。确定它，就是需要题目中的直角条件。注意，这里 \(\left|P F_{1}\right|>\left|P F_{2}\right|\) 虽然也是一个条件，但是它不会带来＂方程＂，只带来＂不等式＂，所以无法消除自由度。那么，如何翻译 \(\angle F_{1} P F_{2}=90^{\circ}\) 这个条件呢？\\
（1）勾股定理：根据勾股定理，

\[
\left|P F_{1}\right|^{2}+\left|P F_{2}\right|^{2}=\left|F_{1} F_{2}\right|^{2}
\]

根据已知条件，\(\left|F_{1} F_{2}\right|=2 c=2 \sqrt{3}\) ，所以

\[
\left|P F_{1}\right|^{2}+\left|P F_{2}\right|^{2}=12
\]

为了确定这个三角形，我们只需要解出 \(\left|P F_{1}\right|,\left|P F_{2}\right|\) 的长度即可。为了化简这个表达式，令 \(\left|P F_{1}\right|=m,\left|P F_{2}\right|=n\) ，所以

\[
m^{2}+n^{2}=12
\]

观察到这是一个关于 \(m, n\) 的对称式，所以我们想到，也可以再引入一个关于 \(m, n\) 的对称式。这就是椭圆的第一定义：

\[
m+n=\left|P F_{1}\right|+\left|P F_{2}\right|=2 a=4
\]

所以，\(m, n\) 满足

\[
\left\{\begin{aligned}
m+n & =4 \\
m^{2}+n^{2} & =12
\end{aligned}\right.
\]

由于 \(m=\left|P F_{1}\right|>\left|P F_{2}\right|=n\) ，解这个二次方程，可得

\[
m=2+\sqrt{2}, n=2-\sqrt{2}
\]

根据三角形面积公式（或勾股定理）：

\[
S_{\Delta P F_{1} F_{2}}=\frac{1}{2} m n=\frac{1}{2}\left|y_{P}\right| \cdot\left|F_{1} F_{2}\right|
\]

故

\[
\left|y_{P}\right|=\frac{m n}{\left|F_{1} F_{2}\right|}=\frac{2}{2 \sqrt{3}}=\frac{\sqrt{3}}{3}
\]

故根据椭圆方程，

\[
x_{P}^{2}=4\left(1-y_{P}^{2}\right)=\frac{8}{3}
\]

又因为 \(\left|P F_{1}\right|>\left|P F_{2}\right|\) ，所以 \(x_{P}>0\) ．故 \(x_{P}=\frac{2 \sqrt{6}}{3}\) ．因此 \(P\) 点坐标为 \(\left(\frac{2 \sqrt{6}}{3}, \frac{\sqrt{3}}{3}\right)\) 或 \(\left(\frac{2 \sqrt{6}}{3},-\frac{\sqrt{3}}{3}\right)\)\\
事实上，如果我们直接观察到了 \(\left|y_{P}\right|=\frac{m n}{\left|F_{1} F_{2}\right|}\) ，甚至不需要求解 \(m, n\) 具体的值。这是因为，我们已知 \(m+n, m^{2}+n^{2}\) 这两个关于 \(m, n\) 的对称式，那么就可以确定所有 \(m, n\) 的对称式：

\[
m n=\frac{(m+n)^{2}-\left(m^{2}+n^{2}\right)}{2}=2
\]

（2）向量：我们可以利用 \(\overrightarrow{P F_{1}} \cdot \overrightarrow{P F_{2}}=0\) 来翻译。设 \(P(x, y)\) ，根据 \(F_{1}(-\sqrt{3}, 0), F_{2}(\sqrt{3}, 0)\) ，可以得到

\[
\overrightarrow{P F_{1}} \cdot \overrightarrow{P F_{2}}=(x+\sqrt{3})(x-\sqrt{3})+y^{2}=x^{2}+y^{2}-3=0
\]

又因为 \(P(x, y)\) 在椭圆上，所以

\[
\frac{x^{2}}{4}+y^{2}=1
\]

这样，我们就得到了关于 \(x^{2}, y^{2}\) 的一元二次方程组：

\[
\left\{\begin{array}{c}
x^{2}+y^{2}=3 \\
\frac{x^{2}}{4}+y^{2}=1
\end{array}\right.
\]

消元解得 \(x^{2}=\frac{8}{3}, y^{2}=\frac{1}{3}\) ．因为 \(\left|P F_{1}\right|>\left|P F_{2}\right|\) ，所以 \(P\) 在 \(y\) 轴右侧，所以 \(x>0\) ．故 \(P\) 点坐标为 \(\left(\frac{2 \sqrt{6}}{3}, \frac{\sqrt{3}}{3}\right)\) 或 \(\left(\frac{2 \sqrt{6}}{3},-\frac{\sqrt{3}}{3}\right)\) ．\\
（3）斜率：设 \(P(x, y)\) ，则 \(k_{P F_{1}} \cdot k_{P F_{2}}=-1\) ．因此

\[
\begin{aligned}
& \frac{y}{x+\sqrt{3}} \cdot \frac{y}{x-\sqrt{3}}=-1 \\
& \frac{y^{2}}{x^{2}-3}=-1 \\
& x^{2}+y^{2}-3=0
\end{aligned}
\]

这和（3）中得到的方程不谋而合。我们要意识到，用斜率和向量翻译直角是完全等价的。\\
（4）利用几何意义：由于 \(P F_{1} F_{2}\) 是直角三角形，所以斜边中线等于斜边一半。因此

\[
|O P|=\frac{1}{2}\left|F_{1} F_{2}\right|=\sqrt{3}
\]

因此，\(P(x, y)\) 在以 \(O\) 为圆心，\(\sqrt{3}\) 为半径的圆上。所以

\[
x^{2}+y^{2}-3=0
\]

这和（2），（3）中的得到的方程不谋而合。

\section*{4．8．2．2 已知一个参数的情形}
已知椭圆 \(\frac{x^{2}}{a^{2}}+y^{2}=1\) ．若 \(\angle F_{1} P F_{2}=90^{\circ}\) ，且 \(\left|P F_{1}\right|=(3+2 \sqrt{2})\left|P F_{2}\right|\) ．求 \(P\) 点坐标。在这个情况中，椭圆中有一个参数不确定，所以一共有两个自由度。为了确定两个自由度，需要两个限制条件或方程。根据刚才的处理，我们需要另一个方程。如果我们数末知数个数，会发现这道题中有三个未知数：\(x, y, a\) 。那这个问题到底有几个自由度呢？事实上，如果认为只有两个自由度，是因为我们考虑了 \(P\) 在粗圆上。算上这个限制，有两个自由度；如果不考虑，就有三个自由度。所以算上＂\(P\) 在椭圆上＂，我们一共需要三个方程。\\
（1）勾股定理：同样设 \(\left|P F_{1}\right|=m,\left|P F_{2}\right|=n\) ．根据已知条件，我们可以列出三个方程：\\
（i）\(P\) 在椭圆上，所以 \(m+n=\left|P F_{1}\right|+\left|P F_{2}\right|=2 a\) ；\\
（ii）\(\left|P F_{1}\right|=(3+2 \sqrt{2})\left|P F_{2}\right|\) ，所以 \(m=(3+2 \sqrt{2}) n\) ；\\
（iii）三角形 \(P F_{1} F_{2}\) 是直角三角形，所以

\[
m^{2}+n^{2}=\left|P F_{1}\right|^{2}+\left|P F_{2}\right|^{2}=\left|F_{1} F_{2}\right|^{2}=(2 c)^{2}=4 c^{2}=4\left(a^{2}-1\right)=4 a^{2}-4
\]

因此，我们得到了关于 \(m, n, a\) 的三个方程：

\[
\left\{\begin{array}{c}
m+n=2 a \\
m=(3+2 \sqrt{2}) n \\
m^{2}+n^{2}=4 a^{2}-4
\end{array}\right.
\]

这个三元方程如何求解？观察这组方程的结构， 13 两个方程是对称的， \(1, ~ 2\) 两个方程是齐次的。我们考虑先利用对称性。

\[
m n=\frac{(m+n)^{2}-\left(m^{2}+n^{2}\right)}{2}=\frac{4 a^{2}-\left(4 a^{2}-4\right)}{2}=2
\]

因此，我们可以替换掉第三个方程：

\[
\left\{\begin{array}{c}
m+n=2 a \\
m=(3+2 \sqrt{2}) n \\
m n=2
\end{array}\right.
\]

【方法一】如果观察前两个方程，把 \(a\) 当作参数，那么我们得到了一个关于 \(m, n\) 的非齐次方程和一个齐次方程，所以可以用 \(a\) 来表示 \(m, n\) ．将 \(m=(3+2 \sqrt{2}) n\) 代入第一个方程，得

\[
\begin{aligned}
& (4+2 \sqrt{2}) n=2 a \\
& n=\frac{2 a}{4+2 \sqrt{2}}=\frac{a}{2+\sqrt{2}}=\frac{(2-\sqrt{2}) a}{2}
\end{aligned}
\]

而

\[
m=2 a-n=2 a-\frac{(2-\sqrt{2}) a}{2}=\frac{(2+\sqrt{2}) a}{2}
\]

代入第三个方程，得

\[
m n=\frac{(2-\sqrt{2}) a}{2} \cdot \frac{(2+\sqrt{2}) a}{2}=\frac{1}{2} a^{2}=2
\]

故 \(a^{2}=4\) ．即 \(a=2\) ．所以解为

\[
\left\{\begin{array}{l}
m=2+\sqrt{2} \\
n=2-\sqrt{2} \\
a=2
\end{array}\right.
\]

【方法二】事实上，后两个方程可以直接确定 \(m, n\) 的值。将 \(m=(3+2 \sqrt{2}) n\) 代人，得

\[
m n=(3+2 \sqrt{2}) n^{2}=2
\]

故

\[
n^{2}=\frac{2}{3+2 \sqrt{2}}=2(3-2 \sqrt{2})
\]

如果对 \(3-2 \sqrt{2}=(\sqrt{2}-1)^{2}\) 非常熟悉，就可以直接得到答案：

\[
n=\sqrt{2} \cdot(\sqrt{2}-1)=2-\sqrt{2}
\]

从而

\[
m=\frac{2}{n}=\frac{2}{2-\sqrt{2}}=\frac{\sqrt{2}}{\sqrt{2}-1}=\sqrt{2}(\sqrt{2}+1)=2+\sqrt{2}
\]

故

\[
2 a=m+n=4
\]

所以 \(a=2\) ．\\
上述两个方法都可以得到 \(m, n, a\) 的值，但是都需要较高的对根式处理的能力。求解出 \(m, n\) 后，求 \(P\) 点坐标的过程和（1）一样，这里不再赘述。\\
（2）向量，斜率，几何意义：这三个方法本质相同，都可以得到

\[
x^{2}+y^{2}=c^{2}=a^{2}-1
\]

问题的关键是如何翻译 \(\left|P F_{1}\right|=(3+2 \sqrt{2})\left|P F_{2}\right|\) 这个条件。如果直接用两点间距离公式翻译，将会非常复杂。所以，我们要另辟蹊径。

如果我们重新分析条件，我们会发现，其实题目给了三个已知条件：\\
（i）方程中短半轴长为 1 ，这是一个非齐次条件；\\
（ii）\(\angle F_{1} P F_{2}=90^{\circ}\) ，可以理解为一个关于边长的齐次条件；\\
（iii）\(\left|P F_{1}\right|=(3+2 \sqrt{2})\left|P F_{2}\right|\) ，也可以理解为一个关于边长的齐次条件。\\
在我们对于离心率的处理中，我们已经知道：只需要两个关于边长的齐次条件，就可以确定离心率。所以，我们可以通过（ii），（iii）两个方程确定离心率。确定离心率后，通过（i）确定椭圆方程，这样问题就转化为＂椭圆方程＂已知的情况了。

事实上，（ii），（iii）两个条件直接确定了三角形 \(P F_{1} F_{2}\) 的形状，它是一个较大锐角正切值为 \(3+2 \sqrt{2}\) 的直角三角形。根据这个条件，我们可以把三边比例都表示出来。设 \(\left|P F_{2}\right|=n\) ，那么 \(\left|P F_{1}\right|=(3+2 \sqrt{2}) n\) ．

\[
\left|F_{1} F_{2}\right|^{2}=\left|P F_{1}\right|^{2}+\left|P F_{2}\right|^{2}=n^{2}+(9+8+12 \sqrt{2}) n^{2}=(18+12 \sqrt{2}) n^{2}
\]

注意到

\[
18+12 \sqrt{2}=6(3+2 \sqrt{2})=6 \cdot(\sqrt{2}+1)^{2}
\]

所以

\[
\left|F_{1} F_{2}\right|^{2}=6(\sqrt{2}+1)^{2} n^{2}
\]

所以

\[
2 c=\left|F_{1} F_{2}\right|=\sqrt{6} \cdot(\sqrt{2}+1)=(2 \sqrt{3}+\sqrt{6}) n
\]

能否用 \(n\) 也来表示 \(a\) 呢？根据椭圆的第一定义，

\[
2 a=\left|P F_{1}\right|+\left|P F_{2}\right|=(3+2 \sqrt{2}+1) n=(4+2 \sqrt{2}) n
\]

因此

\[
e=\frac{2 c}{2 a}=\frac{c}{a}=\frac{2 \sqrt{3}+\sqrt{6}}{4+2 \sqrt{2}}
\]

如果直接化简这个分数，就需要大量的根式运算。因此，我们要利用计算 \(c\) 时的表达式：\(\sqrt{6} \cdot(\sqrt{2}+1)=(2 \sqrt{3}+\sqrt{6})\)

\[
e=\frac{2 \sqrt{3}+\sqrt{6}}{4+2 \sqrt{2}}=\frac{\sqrt{6}(\sqrt{2}+1)}{4+2 \sqrt{2}}=\frac{\sqrt{6}(\sqrt{2}+1)}{2 \sqrt{2}(\sqrt{2}+1)}=\frac{\sqrt{6}}{2 \sqrt{2}}=\frac{\sqrt{3}}{2} .
\]

注意到 \(b=1\) ，所以

\[
1=a^{2}-c^{2}=\frac{4}{3} c^{2}-c^{2}=\frac{1}{3} c^{2}
\]

故 \(c=\sqrt{3}\) ，同时 \(a=\frac{c}{e}=2\) 。后续步骤和 1 中一样。\\
事实上，我们如果重新观察勾股定理中【方法一】的解方程的方法，它本质上也是在确定三角形的形状。在刚才的方法中，我们根据勾股定理，用 \(n\) 去表示 \(c\) ，再根据第一定义，用 \(n\) 表示 \(a\) 。这个方法的弊端就是需要开根号，需要对根号非常敏锐才行；而【方法一】则利用题目条件，第一定义，用 \(a\) 来表示 \(m, n\) ，同时利用对称性，将勾股定理巧妙地转化为 \(m n\) 的形式，让共轭根式派上了用场。这样，就规避了平方的问题。事实上，方法一就算使用 \(m^{2}+n^{2}=4 a^{2}-4\) ，也可以很快得到答案。\\
（3）椭圆有两个参数未知的情况。\\
例如，对于椭圆

\[
\frac{x^{2}}{a^{2}}+\frac{y^{2}}{b^{2}}=1(a>b>0)
\]

其中 \(P\) 是椭圆上一点。这个问题就有三个自由度（如果考虑 \(P\) 在椭圆上）。和（2）中情况类似，最终会得到四个未知数 \((a, b, m, n\) 或 \(a, b, x, y)\) 和四个方程。我们需要根据题目给出的三个条件和＂\(P(x, y)\) 在椭圆 \(\frac{x^{2}}{a^{2}}+\frac{y^{2}}{b^{2}}=1\)上＂来求解。

设 \(P\) 是椭圆 \(\frac{x^{2}}{a^{2}}+\frac{y^{2}}{b^{2}}=1(a>b>0)\) 上一点，且 \(\angle F_{1} P F_{2}=90^{\circ},\left|P F_{1}\right|=(3+2 \sqrt{2})\left|P F_{2}\right|, S_{\triangle P F_{1} F_{2}}=1\) 。求 \(P\)点坐标。这个问题是在（2）的基础之上，将 \(b=1\) 的条件去掉，换为面积这个条件。我们依旧可以逐条分析这些条件：\\
（i）\(\angle F_{1} P F_{2}=90^{\circ}\) ，关于边的齐次条件；\\
（ii）\(\left|P F_{1}\right|=(3+2 \sqrt{2})\left|P F_{2}\right|\) ，关于边的齐次条件；\\
（iii）\(P\) 在椭圆上，\(\frac{x^{2}}{a^{2}}+\frac{y^{2}}{b^{2}}=1\) ，这也是齐次条件。除了它可以看作 \(x, y, a, b\) 的齐次条件，它也可以看作边的齐次条件：根据椭圆的第一定义，\(\left|P F_{1}\right|+\left|P F_{2}\right|=2 a\) ．所以是边的齐次条件；\\
（iv）\(S_{\triangle P F_{1} F_{2}}=1\) ，唯一的非齐次条件。\\
根据自由度，齐次的理论，我们知道：三个齐次条件和一个非齐次条件，可以帮我们确定四个变量，即 \(a, b, x, y\)或 \(a, b, m, n\) 。事实上，（i）～（iii）三个齐次条件可以确定整个图形的形状。根据 2 中的讨论，我们已经求出了椭圆的离心率 \(e=\frac{\sqrt{3}}{2}\) 以及 \(m, n\) 和 \(a, b\) 的关系：

\[
m=\frac{2+\sqrt{2}}{2} a, n=\frac{2-\sqrt{2}}{2} a .
\]

那么，根据 \(S_{\triangle P F_{1} F_{2}}=1\) ，可得：

\[
S_{\triangle P F_{1} F_{2}}=\frac{1}{2} m n=\frac{1}{2} \cdot \frac{2+\sqrt{2}}{2} \cdot \frac{2-\sqrt{2}}{2} a^{2}=\frac{1}{4} a^{2}=1 .
\]

故 \(a=2, b=1, m=2+\sqrt{2}, n=2-\sqrt{2}\) ．后续步骤和 1 一样，\(P\) 的坐标为 \(\left(\frac{2 \sqrt{6}}{3}, \frac{\sqrt{3}}{3}\right)\) 或 \(\left(\frac{2 \sqrt{6}}{3},-\frac{\sqrt{3}}{3}\right)\) ．

\section*{圆曲计算的基本思路}
对于所有需要联立的题目，我们基本可以把圆雉曲线的计算分为如下三个大的步骤：\\
（1）设未知数：决定完全以点出发计算，还是以线出发计算；\\
（2）联立：将直线和圆雉曲线的方程联立，得到关于未知数的方程；\\
（3）翻译：根据题目条件，用关键点的坐标将所有需要表示的量表示出来；\\
（4）化简：对翻译好的表达式进行恒等变形，完成证明或求取值范围。\\
虽然看上去解决圆雉曲线的题目看上去只需要这四个步骤，但每个步骤中，需要注意的关键点都非常多。因此，想要在一节内将它的底层逻辑讲清是十分困难的。在这一节中，我们解决其中最本质，最关键的内容。其他内容，留待后续有机会进行深入的讲解。

这四个步骤中，和计算关系较大的是（2）～（4）三个步骤。

\section*{4．8．2．3 联立}
在实际运算中，我们经常需要将直线方程和圆雉曲线方程进行联立。在计算时，有哪些需要注意的点呢？\\
我们以一个简单的计算为例。设直线 \(y=k x+b\) 和椭圆 \(\frac{x^{2}}{3}+\frac{y^{2}}{4}=1\) 交于 \(A, B\) 两点。那么，它们的坐标满足如下方程组：

\[
\left\{\begin{array}{l}
y=k x+b \\
\frac{x^{2}}{3}+\frac{y^{2}}{4}=1
\end{array}\right.
\]

此时，为了消元，我们可以直接将直线 \(y=k x+b\) 代入到椭圆方程中，得到：

\[
\frac{x^{2}}{3}+\frac{(k x+b)^{2}}{4}=1
\]

在这里，很多同学可能会直接展开计算。在计算前，我们还是要着重观察它的结构：这个式子本质上是分式，有分子，分母两部分组成。如果我们不通分，直接计算，会需要反复计算很多分数的加，减法。因此，在这里，我们通常选择直接\\
通分：

\[
4 x^{2}+3(k x+b)^{2}=12
\]

这里有一个易错点：通分（两边同乘 12）后，很多同学会忘记在右侧也乘分母，导致还是写1．希望同学们在这里多停下来检查一下。

当我们得到了整式后，再整理即可。由于我们关心点的横坐标，所以这是一个关于 \(x\) 的二次方程，我们以 \(x\) 为主元进行整理。这样就得到

\[
\left(3 k^{2}+4\right) x^{2}+6 k b x+3 b^{2}-12=0
\]

这个二次方程整体较简单，我们再来看另一个情况。设直线 \(l\) 经过点 \(A(2,1)\) ，且斜率为 \(k\) ，它和椭圆 \(\frac{x^{2}}{3}+\frac{y^{2}}{4}=1\) 交于 \(B, C\) 两点。首先，直线 \(l\) 的方程为

\[
y=k(x-2)+1
\]

因此，\(B, C\) 两点满足的方程组为

\[
\left\{\begin{array}{l}
y=k(x-2)+1 \\
\frac{x^{2}}{3}+\frac{y^{2}}{4}=1
\end{array}\right.
\]

将 \(y=k(x-2)+1\) 代人后，得到

\[
\begin{aligned}
& \frac{x^{2}}{3}+\frac{(k(x-2)+1)^{2}}{4}=1 \\
& 4 x^{2}+3[k(x-2)+1]^{2}=12
\end{aligned}
\]

接下来，如何展开括号呢？如果我们不去观察代数结构，有如下两种直接的展开方法：

方法（1）：直接利用三项的平方展开：

\[
\begin{aligned}
{[k(x-2)+1]^{2} } & =(k x-2 k+1)^{2}=(k x)^{2}+4 k^{2}+1-4 k^{2} x+2 k x-4 k \\
& =k^{2} x^{2}+\left(2 k-4 k^{2}\right) x+4 k^{2}-4 k+1
\end{aligned}
\]

这样一来，\(x\) 满足的二次方程为：

\[
\begin{aligned}
& 4 x^{2}+3\left[k^{2} x^{2}+\left(2 k-4 k^{2}\right) x+4 k^{2}-4 k+1\right]=12 \\
& \left(3 k^{2}+4\right) x^{2}+\left(6 k-12 k^{2}\right) x+12 k^{2}-12 k-9=0
\end{aligned}
\]

方法（2）：将 \(k(x-2), 1\) 看作两项，分两次展开括号：

\[
\begin{aligned}
{[k(x-2)+1]^{2} } & =k^{2}(x-2)^{2}+2 k(x-2)+1 \\
& =k^{2} x^{2}-4 k^{2} x+4 k^{2}+2 k x-4 k+1 \\
& =k^{2} x^{2}+\left(2 k-4 k^{2}\right) x+4 k^{2}-4 k+1
\end{aligned}
\]

这样一来，\(x\) 满足的二次方程为：

\[
\begin{aligned}
& 4 x^{2}+3\left[k^{2} x^{2}+\left(2 k-4 k^{2}\right) x+4 k^{2}-4 k+1\right]=12 \\
& \left(3 k^{2}+4\right) x^{2}+\left(6 k-12 k^{2}\right) x+12 k^{2}-12 k-9=0
\end{aligned}
\]

虽然这两种方法都可以得到正确答案，但是它们都没有利用好问题本身的代数结构。在这里，我们要充分利用好主元，参数的地位的差别。\\
方法（3）：由于我们选 \(x\) 为主元，所以，\(k\) 和其他的常数一样，它被看作是固定的值。因此，\(k(x-2)+1\) 是一个未被完全化简的形式。正确的形式为：

\[
k(x-2)+1=k x+(1-2 k)
\]

这样一来，我们把 \(1-2 k\) 看作一个整体，它就是常数项。再去展开时，只需用到两项的完全平方公式，计算起来非常方便：

\[
\begin{aligned}
{[k(x-2)+1]^{2} } & =[k x+(1-2 k)]^{2} \\
& =k^{2} x^{2}+\left(2 k-4 k^{2}\right) x+4 k^{2}-4 k+1
\end{aligned}
\]

\(x\) 满足的二次方程为：

\[
\begin{aligned}
& 4 x^{2}+3\left[k^{2} x^{2}+\left(2 k-4 k^{2}\right) x+4 k^{2}-4 k+1\right]=12 \\
& \left(3 k^{2}+4\right) x^{2}+\left(6 k-12 k^{2}\right) x+12 k^{2}-12 k-9=0
\end{aligned}
\]

比较起来，这三个方法，（3）最快：它只使用了一次二项完全平方公式的展开。（1）使用了三项完全平方公式，项数很多，整理起来很凌乱；（2）则使用了两次完全平方公式，复杂度较高。因此，采用方法（3）最快。

从几何意义上讲，（3）也更加合理：我们本质上是找到了直线的 \(y\) 轴截距 \(b=1-2 k\) ，把直线方程化成了 \(y=k x+b\) 的形式进行展开。这和前一种情况是完全相同的，也是化归思想的一种体现。

从这个例子，我们充分体会到了代数结构发挥的杜杆作用。如果不顾结构，盲目展开，就会白算；利用好基本的主元，参数的思想，就可以以较快的速度完成运算。

\section*{4．8．2．4 翻译}
在翻译的时候，我们经常需要用坐标去表示长度，角度，内积等。在多数情况下，这些量都是关于点的坐标 \(x_{1}, x_{2}\) 是对称的。这样一来，就可以利用韦达定理进行化简。因此，我们最终的目标都是将所有这些量用关

于 \(x_{1}, x_{2}\) 的对称式进行表示。我们先从最简单的内容出发求解。设 \(A\left(x_{1}, y_{1}\right), B\left(x_{2}, y_{2}\right)\) 是直线 \(y=k x+b\) 和某二次曲线的交点。\\
（1）\(A B\) 长度：我们可以用关于 \(x_{1}, x_{2}\) 的对称式表示 \(|A B|\) ．根据两点间距离公式，有

\[
|A B|=\sqrt{\left(x_{2}-x_{1}\right)^{2}+\left(y_{2}-y_{1}\right)^{2}}
\]

注意到，这里是关于四个字母 \(x_{1}, x_{2}, y_{1}, y_{2}\) 的表达式，我们可以利用方程进行消元。显然，利用直线方程更好，因为它次数更低。

\[
\begin{aligned}
|A B| & =\sqrt{\left(x_{2}-x_{1}\right)^{2}+\left(k x_{1}+b-k x_{2}-b\right)^{2}} \\
& =\sqrt{\left(x_{2}-x_{1}\right)^{2}+k^{2}\left(x_{2}-x_{1}\right)^{2}} \\
& =\sqrt{\left(k^{2}+1\right)\left(x_{1}-x_{2}\right)^{2}} \\
& =\sqrt{k^{2}+1} \cdot\left|x_{1}-x_{2}\right|
\end{aligned}
\]

一些教辅会将这个公式称为＂弦长公式＂。事实上，它就是两点间距离公式的等价形式，不需要 \(A B\) 构成某二次曲线的弦长。

在这个公式中，出现了 \(\left|x_{1}-x_{2}\right|\) 。首先，观察到这是一个对称式：交换 \(x_{1}, x_{2}\) 后，式子不发生改变：

\[
\left|x_{2}-x_{1}\right|=\left|x_{1}-x_{2}\right|
\]

那么，如何处理这个式子呢？\\
（2）对于初学的同学：根据杜杆 2 的内容，我们知道：对称多项式可以用 \(x_{1}+x_{2}, x_{1} x_{2}\) 表示。因此，我们通常不把 \(\left(x_{1}-x_{2}\right)^{2}\) 的根号开出来，而是还原成：

\[
|A B|=\sqrt{k^{2}+1} \cdot \sqrt{\left(x_{1}-x_{2}\right)^{2}}
\]

的形式。并将它写成

\[
\begin{aligned}
|A B| & =\sqrt{k^{2}+1} \cdot \sqrt{x_{1}^{2}-2 x_{1} x_{2}+x_{2}^{2}} \\
& =\sqrt{k^{2}+1} \cdot \sqrt{\left(x_{1}+x_{2}\right)^{2}-4 x_{1} x_{2}}
\end{aligned}
\]

这样一来，我们还可以使用韦达定理，直接代入 \(x_{1}+x_{2}, x_{1} x_{2}\) 的表达式。再代入韦达定理即可。\\
（2）对于熟练的同学：根据杜杆 7 的内容，我们知道：两根之差和判别式有着非常深人的联系。设 \(x_{1}, x_{2}\) 是方程 \(a x^{2}+b x+c=0\) 的两根。那么：

\[
x_{1,2}=\frac{-b \pm \sqrt{b^{2}-4 a c}}{2 a}
\]

因此，

\[
\begin{aligned}
\left|x_{1}-x_{2}\right| & =\left|\frac{-b+\sqrt{b^{2}-4 a c}}{2 a}-\frac{-b-\sqrt{b^{2}-4 a c}}{2 a}\right| \\
& =\left|\frac{2 \sqrt{b^{2}-4 a c}}{2 a}\right| \\
& =\frac{\sqrt{\Delta}}{|a|}
\end{aligned}
\]

因此，如果已经计算好判别式，也可以直接代入：

\[
|A B|=\sqrt{k^{2}+1} \cdot \frac{\sqrt{\Delta}}{|a|}
\]

这里我们需要特别强调：此处的 \(a\) ，不是椭圆，双曲线方程中的 \(a\) ，而是联立后，\(x_{1}, x_{2}\) 满足的二次方程的二次项系数。例如，若联立后的二次方程为

\[
\left(3 k^{2}+4\right) x^{2}+6 k b x+3 b^{2}-12=0
\]

那么 \(a=3 k^{2}+4\) ，

\[
|A B|=\sqrt{k^{2}+1} \cdot \frac{\sqrt{\Delta}}{\left|3 k^{2}+4\right|}
\]

（2）其他对称式：我们以长度为例，展示了如何翻译。事实上，我们还会遇到很多其他的量，它们也都可以用 \(x_{1}, x_{2}\) 的对称式表示。在处理这些量的时候，一个大原则，就是＂利用直线方程消元，保护对称性＂。

例如，设 \(P(1,2)\) ，如何翻译 \(\overrightarrow{P A} \cdot \overrightarrow{P B}=2\) ？根据向量的坐标表示，

\[
\overrightarrow{P A}=\left(x_{1}-1, y_{1}-2\right), \overrightarrow{P B}=\left(x_{2}-1, y_{2}-2\right)
\]

因此，

\[
\begin{aligned}
\overrightarrow{P A} \cdot \overrightarrow{P B} & =\left(x_{1}-1\right)\left(x_{2}-1\right)+\left(y_{1}-2\right)\left(y_{2}-2\right) \\
& =x_{1} x_{2}-\left(x_{1}+x_{2}\right)+1+y_{1} y_{2}-2\left(y_{1}+y_{2}\right)-4
\end{aligned}
\]

显然，这是关于 \(x_{1}, x_{2}\) 和 \(y_{1}, y_{2}\) 都对称的式子。为了后面直接使用韦达定理，我们需要将 \(y_{1} y_{2}, y_{1}+y_{2}\) 也都用 \(x_{1} x_{2}, x_{1}+x_{2}\) 表示。这就需要使用直线方程消元了。\\
（1）\(y_{1} y_{2}\) 可以写作：

\[
\begin{aligned}
y_{1} y_{2} & =\left(k x_{1}+b\right)\left(k x_{2}+b\right) \\
& =k^{2} x_{1} x_{2}+k b\left(x_{1}+x_{2}\right)+b^{2}
\end{aligned}
\]

（2）\(y_{1}+y_{2}\) 可以写作：

\[
y_{1}+y_{2}=k x_{1}+b+k x_{2}+b=k\left(x_{1}+x_{2}\right)+2 b .
\]

这样，就将所有关于 \(y\) 的对称式转化为关于 \(x\) 的对称式了。最后，代人即可：

\[
\begin{aligned}
\overrightarrow{P A} \cdot \overrightarrow{P B}= & x_{1} x_{2}-\left(x_{1}+x_{2}\right)+1+y_{1} y_{2}-2\left(y_{1}+y_{2}\right)-4 \\
= & x_{1} x_{2}-\left(x_{1}+x_{2}\right)+1+k^{2} x_{1} x_{2}+k b\left(x_{1}+x_{2}\right)+b^{2} \\
& -2\left[k\left(x_{1}+x_{2}\right)+2 b\right]+4 \\
= & \left(k^{2}+1\right) x_{1} x_{2}+(k b-2 k-1)\left(x_{1}+x_{2}\right)+b^{2}-4 b+5
\end{aligned}
\]

这样，我们就维持了表达式本身的对称结构。我们在整理时，也要以 \(x_{1} x_{2}, x_{1}+x_{2}\) 为主元整理。两个主要原因：\\
一是因为它是对称式，所以可以写作它们的表达式；二是在下一步中，我们需要利用韦达定理化简。\\
为了给大家更充足的理解，我们这里提供更多的例子，大家可以试着进行化简：\\
（3）\(x_{1} y_{2}+x_{2} y_{1}\) ：它可以写作：

\[
\begin{aligned}
x_{1} y_{2}+x_{2} y_{1} & =x_{1}\left(k x_{2}+b\right)+x_{2}\left(k x_{1}+b\right) \\
& =2 k x_{1} x_{2}+b\left(x_{1}+x_{2}\right)
\end{aligned}
\]

（4）\(\frac{y_{1}}{x_{1}-1}+\frac{y_{2}}{x_{2}-1}\) ：它的几何意义是：设 \(Q(1,0)\) ，那么 \(A Q, B Q\) 两条直线的斜率之和为

\[
k_{A Q}+k_{B Q}=\frac{y_{1}}{x_{1}-1}+\frac{y_{2}}{x_{2}-1}
\]

利用直线方程消元，可得：

\[
\frac{y_{1}}{x_{1}-1}+\frac{y_{2}}{x_{2}-1}=\frac{k x_{1}+b}{x_{1}-1}+\frac{k x_{2}+b}{x_{2}-1}
\]

在这里，我们既可以选择直接通分，也可以选择分离常数。\\
（i）直接通分：

\[
\begin{aligned}
\frac{k x_{1}+b}{x_{1}-1}+\frac{k x_{2}+b}{x_{2}-1} & =\frac{\left[\left(k x_{1}+b\right)\left(x_{2}-1\right)+\left(k x_{2}+b\right)\left(x_{1}-1\right)\right]}{\left(x_{1}-1\right)\left(x_{2}-1\right)} \\
& =\frac{k x_{1} x_{2}-k x_{1}+b x_{2}-b+k x_{1} x_{2}-k x_{2}+b x_{1}-b}{x_{1} x_{2}-\left(x_{1}+x_{2}\right)+1} \\
& =\frac{2 k x_{1} x_{2}+(b-k)\left(x_{1}+x_{2}\right)-2 b}{x_{1} x_{2}-\left(x_{1}+x_{2}\right)+1}
\end{aligned}
\]

（ii）分离常数：

\[
\begin{aligned}
\frac{k x_{1}+b}{x_{1}-1}+\frac{k x_{2}+b}{x_{2}-1} & =k+\frac{b+k}{x_{1}-1}+k+\frac{b+k}{x_{2}-1} \\
& =2 k+(k+b) \cdot\left(\frac{1}{x_{1}-1}+\frac{1}{x_{2}-1}\right) \\
& =2 k+\frac{(k+b)\left(x_{1}+x_{2}\right)-2(k+b)}{x_{1} x_{2}-\left(x_{1}+x_{2}\right)+1}
\end{aligned}
\]

我们发现，分离常数计算量明显较小。因为通分的操作非常直接，不需要过多系数的计算。\\
（3）反对称式：在实际应用中，我们还会遇到反对称式。反对称的意思是：当交换 \(x_{1}, x_{2}\) 时，代数式的值变为原来的相反数。例如，最简单的 \(x_{1}-x_{2}\) 就是反对称的。对于这种反对称式，我们可以用平方的方法将它变为对称式。例如，

\[
\left(x_{1}-x_{2}\right)^{2}=\left(x_{1}+x_{2}\right)^{2}-4 x_{1} x_{2}
\]

在实际应用中，我们还会遇到非对称式。非对称式的处理需要利用韦达定理，在这里我们暂时略去不表，在化简中进行更深人的讨论。

在翻译的时候，我们如何判断某个量用 \(x_{1}, x_{2}\) 表达式完后一定是对称式呢？我们只需要交换 \(A, B\) 两个字母后会发生什么样的变化即可。例如，\(|A B|\) 交换 \(A, B\) 不发生改变，它变为 \(|B A|=|A B| ; \overrightarrow{P A} \cdot \overrightarrow{P B}\) 交换 \(A, B\) 后，得到 \(\overrightarrow{P B} \cdot \overrightarrow{P A}=\overrightarrow{P A} \cdot \overrightarrow{P B}\) ．所以它也是对称的。

\section*{4．8．2．5 化简}
化简时，我们通常需要充分利用韦达定理和式子的对称性。在后续的运算中，我们都假设 \(A\left(x_{1}, y_{1}\right), B\left(x_{2}, y_{2}\right)\)是直线 \(y=k x+b\) 和椭圆 \(\frac{x^{2}}{3}+\frac{y^{2}}{4}=1\) 的两个交点。那么，它们的坐标满足如下方程组：

\[
\left\{\begin{array}{l}
y=k x+b \\
\frac{x^{2}}{3}+\frac{y^{2}}{4}=1
\end{array}\right.
\]

联立得：\(x_{1}, x_{2}\) 满足方程

\[
\left(3 k^{2}+4\right) x^{2}+6 k b x+3 b^{2}-12=0
\]

根据韦达定理，

\[
\left\{\begin{array}{c}
x_{1}+x_{2}=-\frac{6 k b}{3 k^{2}+4} \\
x_{1} x_{2}=\frac{3 b^{2}-12}{3 k^{2}+4}
\end{array}\right.
\]

我们可以用 \(x_{1}+x_{2}, x_{1} x_{2}\) 表示所有关于 \(x_{1}, x_{2}\)（包括 \(y_{1}, y_{2}\) ）的对称式。\\
在化简时，我们需要直接代入，但是根据问题的具体结构，有时可以采用更快的运算方法。\\
（1）定值问题：对于定值问题，处理方法相对直接，只需要将韦达定理代入即可。设直线 \(A B\) 经过点 \(M\left(\frac{1}{2},-3\right)\) ，设 \(Q\left(\frac{3}{2}, 1\right)\) ，求证 \(k_{A Q}+k_{B Q}\) 为定值。根据 2 中运算，

\[
\frac{k x_{1}+b-1}{x_{1}-\frac{3}{2}}+\frac{k x_{2}+b-1}{x_{2}-\frac{3}{2}}=\frac{2 k x_{1} x_{2}+\left(b-\frac{3 k}{2}-1\right)\left(x_{1}+x_{2}\right)-3 b+3}{x_{1} x_{2}-\frac{3}{2}\left(x_{1}+x_{2}\right)+\frac{9}{4}}
\]

如果我们直接代入韦达定理，就需要书写非常多的分式：

\[
\begin{aligned}
& \frac{2 k x_{1} x_{2}+\left(b-\frac{3 k}{2}-1\right)\left(x_{1}+x_{2}\right)-3 b+3}{x_{1} x_{2}-\frac{3}{2}\left(x_{1}+x_{2}\right)+\frac{9}{4}} \\
= & \frac{2 k \cdot \frac{3 b^{2}-12}{3 k^{2}+4}+\left(b-\frac{3 k}{2}-1\right) \cdot\left(-\frac{6 k b}{3 k^{2}+4}\right)+3-3 b}{\frac{3 b^{2}-12}{3 k^{2}+4}-\frac{3}{2} \cdot\left(-\frac{6 k b}{3 k^{2}+4}\right)+\frac{9}{4}}
\end{aligned}
\]

去分母，得

\[
\begin{aligned}
& =\frac{2 k \cdot\left(3 b^{2}-12\right)+\left(b-\frac{3 k}{2}-1\right) \cdot(-6 k b)+(3-3 b) \cdot\left(3 k^{2}+4\right)}{3 b^{2}-12+9 k b+\frac{9}{4} \cdot\left(3 k^{2}+4\right)} \\
& =\frac{-24 k+6 k b+9 k^{2}+12-12 b}{3 b^{2}+\frac{27 k^{2}}{4}+9 k b-3}
\end{aligned}
\]

由于所有系数都是 3 的倍数，可以约分：

\[
\frac{-24 k+6 k b+9 k^{2}+12-12 b}{3 b^{2}+\frac{27 k^{2}}{4}+9 k b-3}=\frac{-8 k+2 k b+3 k^{2}+4-4 b}{b^{2}+\frac{9}{4} k^{2}+3 k b-1}
\]

事实上，我们在代入韦达定理的时候，可以适当地合并代入和去分母两个步骤。如果将这个去分母的表达式和代入前含有 \(x_{1} x_{2}\) 的表达式做对比，我们发现：只需要将 \(x_{1} x_{2}\) 换成 \(3 b^{2}-12\) ，将 \(x_{1}+x_{2}\) 换成 \(-6 k b\) ，给常数项乘 \(\left(3 k^{2}+4\right)\) 即可。

由于直线经过点 \(M\left(\frac{1}{2},-3\right)\) ，所以 \(\frac{1}{2} k+b=-3\) 。在这个条件下，我们只需要证明所求的式子是定值即可。我们非常自然地想到消元。在这里，\(k, b\) 次数相同，都是一次式，所以消哪个均可。为了避免分数，我们可以消 \(k\) ：由 \(\frac{1}{2} k+b=-3\) 得

\[
k=-2 b-6
\]

故

\[
\begin{aligned}
& \frac{-8 k+2 k b+3 k^{2}+4-4 b}{b^{2}+\frac{9}{4} k^{2}+3 k b-1} \\
= & \frac{8(2 b+6)-2 b(2 b+6)+3(2 b+6)^{2}+4-4 b}{b^{2}+\frac{9}{4}(2 b+6)^{2}-3 b(2 b+6)-1} \\
= & \frac{8 b^{2}+72 b+160}{4 b^{2}+36 b+80} \\
= & 2 .
\end{aligned}
\]

故为定值，\(k_{A Q}+k_{B Q}=2\) 。\\
（2）定点问题：设 \(Q\left(\frac{3}{2}, 1\right)\) ，且 \(k_{A Q}+k_{B Q}=2\) ，求证直线 \(A B\) 过定点。那么，根据 2 中的运算，

\[
\begin{aligned}
& \frac{k x_{1}+b-1}{x_{1}-\frac{3}{2}}+\frac{k x_{2}+b-1}{x_{2}-\frac{3}{2}}=\frac{2 k x_{1} x_{2}+\left(b-\frac{3 k}{2}-1\right)\left(x_{1}+x_{2}\right)-3 b+3}{x_{1} x_{2}-\frac{3}{2}\left(x_{1}+x_{2}\right)+\frac{9}{4}}=2 \\
& (2 k-2) x_{1} x_{2}+\left(b-\frac{3 k}{2}+2\right)\left(x_{1}+x_{2}\right)-3 b-\frac{3}{2}=0
\end{aligned}
\]

如果我们直接将韦达定理代入，就会需要反复书写很多分式：

\[
(2 k-2) \cdot \frac{3 b^{2}-12}{3 k^{2}+4}+\left(b-\frac{3 k}{2}+2\right)\left(-\frac{6 k b}{3 k^{2}+4}\right)-3 b-\frac{3}{2}=0
\]

去分母，得：

\[
(2 k-2) \cdot\left(3 b^{2}-12\right)+\left(b-\frac{3 k}{2}+2\right)(-6 k b)-\left(3 b+\frac{3}{2}\right)\left(3 k^{2}+4\right)=0
\]

在这里，如果想要加快运算，可以适当地将这两步合并。如果将这个去分母的表达式和代人前含有 \(x_{1} x_{2}\) 的表达式做对比，我们发现：只需要将 \(x_{1} x_{2}\) 换成 \(3 b^{2}-12\) ，将 \(x_{1}+x_{2}\) 换成 \(-6 k b\) ，给常数项乘 \(\left(3 k^{2}+4\right)\) 即可。接下来，只要将括号展开即可。

\[
\begin{aligned}
& (2 k-2) \cdot\left(3 b^{2}-12\right)+\left(b-\frac{3}{2} k+2\right)(-6 k b)-\left(3 b+\frac{3}{2}\right)\left(3 k^{2}+4\right)=0 \\
& 6 k b^{2}-24 k-6 b^{2}+24-6 k b^{2}+9 k^{2} b-12 k b-9 k^{2} b-12 b-\frac{9}{2} k^{2}-6=0 \\
& -6 b^{2}-\frac{9}{2} k^{2}-12 k b-24 k-12 b+18=0
\end{aligned}
\]

观察到所有系数都以 3 为约数，可以直接约分，同时乘 2 去掉分数：

\[
-4 b^{2}-3 k^{2}-8 k b-16 k-8 b+12=0
\]

同时，我们发现：大部分项的系数都是负数，所以可以调整符号。

\[
4 b^{2}+3 k^{2}+8 k b+16 k+8 b-12=0
\]

接下来，如何处理这个二次式呢？我们来观察这个式子的结构。\\
（1）它是一个二元二次式，没有对称结构；\\
（2）题目要求证明直线 \(y=k x+b\) 经过定点。换言之，我们要证明存在一个点 \(\left(x_{0}, y_{0}\right)\) ，它永远在直线上。例如，如果 \((-2,1)\) 在直线上，那么

\[
\begin{aligned}
& 1=-2 k+b \\
& -2 k+b-1=0
\end{aligned}
\]

换言之，我们要从这个二元二次式中，找到一个关于 \(k, b\) 的一次因式。因此，这意味着我们需要因式分解。\\
对于二元二次式的因式分解，我们在杜杆中也已经阐述过。这里我们再复习一下。由于问题本身是二元的，我们更熟悉一元式子的因式分解，所以通常先选主元。这里选择 \(k, b\) 哪个为主元均可，因为次数均为 2 次。我们选 \(k\) 为主元，整理得

\[
3 k^{2}+(8 b+16) k+4 b^{2}+8 b-12=0
\]

我们有两种基本方法：\\
（1）直接分解：这需要我们首先对常数项 \(4 b^{2}+8 b-12\) 分解。它可以写作：

\[
4 b^{2}+8 b-12=(b-1)(4 b+12)
\]

因此，原式化为

\[
3 k^{2}+(8 b+16) k+(b-1)(4 b+12)=0
\]

如果我们还想继续分解，仍需要十字相乘。此时，我们可以去考察 \(b\) 的系数。注意到 \(4 b+12=4(b+3)\) ，所以它还有其他可能的分解形式：

\[
(b-1)(4 b+12)=(2 b-2)(2 b+6)
\]

选择哪个？我们需要比较 \(b\) 的系数。我们发现，前者无法得到 \(8 b+16\) 的 8 ，而后者可以： \(3 \times 2+1 \times 2=8\) 。所以，可以完成十字相乘：

\[
\begin{aligned}
& 3 \\
& \times \begin{array}{r}
2 b-2 \\
2 b+6
\end{array}
\end{aligned}
\]

十字相乘，就可以得到 \(3(2 b+6)+2 b-2=8 b+16\) ．完成了分解。故

\[
(3 k+2 b-2)(k+2 b+6)=0
\]

（2）求根公式：分解因式本质上就是在解方程。这是我们在杜杆中反复提到的理解。因此，可以用求根公式的方法求解。

\[
\Delta=(8 b+16)^{2}-12(b-1)(4 b+12)
\]

计算时，要贯彻我们在杜杆 7 中讲到的＂提公因数＂的方法：

\[
\begin{aligned}
& (8 b+16)^{2}-12(b-1)(4 b+12) \\
= & 4^{2} \cdot(2 b+4)^{2}-12 \cdot 4(b-1)(b+3) \\
= & 4^{2} \cdot\left[(2 b+4)^{2}-3(b-1)(b+3)\right] \\
= & 4^{2} \cdot\left(4 b^{2}+16 b+16-3 b^{2}-6 b+9\right) \\
= & 4^{2} \cdot\left(b^{2}+10 b+25\right) \\
= & 4^{2} \cdot(b+5)^{2}=(4 b+20)^{2}
\end{aligned}
\]

因此，\(\Delta\) 是一个完全平方式。所以，

\[
\begin{aligned}
k_{1,2} & =\frac{-(8 b+16) \pm \sqrt{\Delta}}{6}=\frac{-(8 b+16) \pm(4 b+20)}{6} \\
& =\frac{-4 b-8 \pm(2 b+10)}{3}
\end{aligned}
\]

故

\[
\begin{aligned}
& k_{1}=\frac{-4 b-8-2 b-10}{3}=-2 b-6 \\
& k_{2}=\frac{-4 b-8+2 b+10}{3}=-\frac{2}{3} b+\frac{2}{3}
\end{aligned}
\]

原式可分解为

\[
\begin{aligned}
& 3(k+2 b+6)\left(k+\frac{2}{3} b-\frac{2}{3}\right)=0 \\
& (k+2 b+6)(3 k+2 b-2)=0
\end{aligned}
\]

那么，应该是哪个式子为 0 呢？事实上，直线 \(A B\) 一定不会经过 \(Q\left(\frac{3}{2}, 1\right)\) 点。否则 \(k_{A Q}=k_{B Q}\) ，无法保证\\
\(k_{A Q}+k_{B Q}=2\) 恒成立。因此，

\[
\begin{aligned}
& k \cdot \frac{3}{2}+b \neq 1 \\
& 3 k+2 b-2 \neq 0
\end{aligned}
\]

因此，第一个括号不为 0 。所以 \(k+2 b+6=0\) ．它经过哪个顶点呢？我们依然可以从两个角度理解：\\
（1）消元：\(k=-2 b-6\) 代入直线方程，得

\[
y=-(2 b+6) x+b
\]

由于直线 \(A B\) 过定点，所以我们需要找到一组 \((x, y)\) ，使得它对任意的 \(b\) 都成立。因此，我们再以 \(b\) 为主元整理，得到一个关于 \(b\) 的一次方程：

\[
y+6 x=(1-2 x) b
\]

这个方程有无数组解，所以

\[
1-2 x=0, y+6 x=0
\]

解得

\[
x=\frac{1}{2}, y=-3 .
\]

所以定点坐标为 \(M\left(\frac{1}{2},-3\right)\) ．\\
（2）和直线方程作比较。我们观察到直线方程是 \(y=k x+b\) 的形式，其中 \(b\) 的次数为 1 。所以，我们可以将 \(k+2 b+6=0\) 改写成：

\[
k \cdot \frac{1}{2}+b=-3
\]

将它和

\[
k \cdot x+b=y
\]

作对比，显然有 \(x=\frac{1}{2} y=-3\) 。因此，定点坐标为 \(\left(\frac{1}{2},-3\right)\) 。这个方法可以最快得到结果，但是仍应按照（1）书写过程。

我们看到，求根公式的方法有其独特的优势：当我们解出 \(k_{1}, k_{2}\) 后，可以自然地代入直线方程消元。所以，如果看不出分解因式，强烈推荐利用求根公式。这是一个稳定得分，没有偶然因素的方法。这种＂通法＂是我们的底线，不能不会。

最后，我们来统一分析一下这两个题目本身的几何结构。首先，我们发现，如下两个命题是完全等价的：\\
（1）\(k_{Q A}+k_{Q B}=2\) ，即斜率之和为定值；\\
（2）直线 \(A B\) 经过点 \(M\) ，即直线 \(A B\) 过定点。\\
因此，斜率之和为定值和过定点是同一个问题。这个问题有几个自由度呢？看上去有两个动点 \(A, B\)（或一条动直线），它有两个自由度；但是无论是＂过定点＂还是＂\(k_{Q A}+k_{Q B}\)＂，都限制住了一个条件，因此它只有一个自由度。注意，看上去 \(A, B\) 两点共有四个坐标，是 4 个自由度，但由于 \(A, B\) 均在椭圆上，所以本质上有 \(4-2=2\)个自由度。

因此，想要确定这个系统，我们只需要给出一个自由变量即可。那么，有哪些方式来给出自由变量呢？\\
（1） 2 个变量 \textbackslash ＆一个限制：若（1）成立，可以设 \(y=k x+b\) ，同时受 \(k_{Q A}+k_{Q B}=2\) 约束，等价于一个自由度。这就是上文的方法；\\
（2） 1 个变量：若（1）成立，可以设 \(k_{Q A}=k\) ，那么 \(k_{Q B}=2-k \cdot A, B\) 坐标均可用 \(k\) 表示，直线 \(A B\) 方程也可以用 \(k\) 表示，这也只有一个自由度；\\
（3） 2 个变量\textbackslash ＆一个限制：若（1）成立，可以设 \(y=k x+b\) ，同时受 \(A B\) 过 \(M\) 约束，即 \(\frac{1}{2} k+b=-3\) 。这也是上文的方法；\\
（4） 1 个变量：若（2）成立，由于直线过 \(M\) 点，可以采用点斜式：

\[
y+3=k\left(x-\frac{1}{2}\right)
\]

事实上，这和（3）等价。\\
分析自由变量的选取，本质上就是在设未知数。好的未知数取法可以加快计算速度。这里，我们之所以选择了（1），（3）两种设法，有多个因素的考量：\\
（1）如果选择（2），（4），那么联立相对复杂。这是我们在1中看到的情况；\\
（2）如果在解决定点问题时选择（2），求解定点的步骤也相对复杂。\\
为了让大家对这个方法有实感，这里我们也尝试着对它计算：\\
假设 \(k_{A Q}+k_{B Q}=2\) ，设 \(k_{A Q}=k, k_{B Q}=2-k\) ．我们可以用 \(k\) 来表示 \(A, B\) 两点的坐标。直线 \(A Q\) 的方程为

\[
\begin{aligned}
& y-1=k\left(x-\frac{3}{2}\right) \\
& y=k x+1-\frac{3}{2} k
\end{aligned}
\]

和椭圆方程联立，得

\[
\begin{aligned}
& 4 x^{2}+3\left(k x+1-\frac{3}{2} k\right)^{2}=12 \\
& \left(3 k^{2}+4\right) x^{2}+6\left(1-\frac{3}{2} k\right) k x+3\left(1-\frac{3}{2} k\right)^{2}-12=0 \\
& \left(3 k^{2}+4\right) x^{2}+\left(6 k-9 k^{2}\right) x+\frac{27}{4} k^{2}-9 k-9=0
\end{aligned}
\]

为了求 \(A\) 点坐标，看上去需要求解二次方程，可能会出根号，但其实不会。这是因为 \(Q\) 在直线上，所以 \(x=\frac{3}{2}\) 是方程的根。因此，根据韦达定理：

\[
x_{A} \cdot \frac{3}{2}=\frac{\frac{27}{4} k^{2}-9 k-9}{3 k^{2}+4}
\]

故

\[
x_{A}=\frac{\frac{9}{2} k^{2}-6 k-6}{3 k^{2}+4}
\]

故

\[
y_{A}=k\left(x_{A}-\frac{3}{2}\right)+1=k \cdot\left(\frac{-6 k-12}{3 k^{2}+4}\right)+1=\frac{-3 k^{2}-12 k+4}{3 k^{2}+4}
\]

故

\[
A\left(\frac{\frac{9}{2} k^{2}-6 k-6}{3 k^{2}+4}, \frac{-3 k^{2}-12 k+4}{3 k^{2}+4}\right)
\]

在计算点 \(B\) 的坐标时，只需要将 \(A\) 中所有 \(k\) 替换为 \(2-k\) 即可。这是因为：在计算\\
\(A, B\) 两点的坐标时，唯一不同的地方就是斜率。前者是 \(k\) ，后者是 \(2-k\) 。因此，我们只需要将所有的 \(k\) 替换成\\
\(2-k\) 即可。即

\[
\begin{gathered}
x_{B}=\frac{\frac{9}{2}(2-k)^{2}-6(2-k)-6}{3(k-2)^{2}+4}=\frac{\frac{9}{2}\left(k^{2}-4 k+4\right)-12+6 k-6}{3 k^{2}-12 k+12+4} \\
=\frac{\frac{9}{2} k^{2}-12 k}{3 k^{2}-12 k+16} . \\
y_{B}=\frac{-3(2-k)^{2}-12(2-k)+4}{3(2-k)^{2}+4}=\frac{-3 k^{2}+24 k-32}{3 k^{2}-12 k+16} .
\end{gathered}
\]

因此

\[
B\left(\frac{\frac{9}{2} k^{2}-12 k}{3 k^{2}-12 k+16}, \frac{-3 k^{2}+24 k-32}{3 k^{2}-12 k+16}\right)
\]

我们只需要求 \(A B\) 的方程即可。但是从 \(A, B\) 两点的坐标可以看出，它们的表达式非常复杂，所以直接求解直线方程几乎是不可能的，需要另辟蹊径。事实上，如果我们能提前猜出 \(A B\) 经过点 \(M\left(\frac{1}{2},-3\right)\) ，就可以通过斜率相等来证明 \(A B\) 经过 \(M\) 。经过计算得

\[
\begin{aligned}
& K_{A M}=\frac{y_{A}-y_{M}}{x_{1}-x_{M}}=\frac{\frac{-3 k^{2}-12 k+4}{3 k^{2}+4}+3}{\frac{\frac{9}{2}-6 k-6}{3 k^{2}+4}-\frac{1}{2}}=\frac{6 k^{2}-12 k+16}{3 k^{2}-6 k-8} \\
& K_{B M}=\frac{y_{B}-y_{M}}{x_{B}-x_{M}}=\frac{\frac{-3 k^{2}+24 k-32}{3 k^{2}-12 k+16}+3}{\frac{\frac{9}{2}-12 k}{3 k^{2}-12 k+16}-\frac{1}{2}}=\frac{6 k^{2}-12 k+16}{3 k^{2}-6 k-8}
\end{aligned}
\]

因此，\(k_{A M}=k_{B M}\) ，即 \(A, B, M\) 三点共线。\\
这一步的计算相对直接，但是需要知道提前猜出定点 \(M\left(\frac{1}{2},-3\right)\) 的坐标，这需要代入两条特殊直线。这个基本方法将会在后面涉及。

这里我们看到，如果直接消元计算，计算量非常大。这是因为，本身点 \(Q, M\) 都不是特殊点（在 \(x\) 轴或者 \(y\) 轴上）。同时，因为同样的原因，我们无法利用对称性较快地猜出定点位置，所以实战中不推荐使用这个方法。

\section*{第5章 选填破题基本法}
中档以上难度的选填，不仅容易耗时极长，而且是失分大头，错一点， 5 分就没了。做题时，找不到最快，最巧的思路，无法做到一下就出答案，这是因为我们没有用选填该用的方式训练。选填的命题逻辑和大题其实完全不同，大题考察＂综合能力＂，丢分是因为＂做不对＂；选题考察＂破题能力＂，丢分是因为＂想不到＂。而在平时的练习中，我们从未对＂如何想到＂做刻意训练，所以失分严重。

那么该怎么刻意训练＂如何想到＂呢？\\
事实上，如果刻意训练去使用一些视角，你也可以获得极强的选填破题能力：极端遍历法与临界状态，基底，化归思想，动态图像，思想试验等等。为了让更多同学能学会他们，我们进行了系统的总结，在前两章的内容的基础上，本章将对上述内容进行系统的讲解。

讲解中，我们首先会配合足量的《真题》，《原创题》和《模拟题》做练习，带大家逐一熟悉与掌握每一个解题视角。此外我们会用这些新视角，带你重学一遍高中数学，例如：有了动态图像，在函数，向量，圆曲问题里，看似需要复杂推导的二级结论，其实可以被一眼看出来；学会了极端遍历法，在导数，多变量，不等式问题中，你可以生动体会极限情况，快速得到取等条件，边界值。

当你随我学完这一遍，会看到：原本看起来没有关系的知识，背后是相同的数学思想；原来数学可以不是死的，不需要死记硬背，而是活的，有道理的。而我们不只是学会了独立的知识，更是学会了真正的数学！另外，在本章的末尾，还有几套高考真题供大家检验学习效果，我也会展示我真实的解题过程，方便大家进一步体会解题基本法是如何使用的。

最后，我想说一下本章内容的难度，如果你的分数在 90 分以下，建议优先看完第四章的内容。

\section*{5.1 思想试验}
更考验思维能力的＂思想试验＂请结合索引阅读章节＜科学的思维方式解题＞中的相关内容。

\section*{5．1．1 思想试验如何破题？}
面对一个全新的，陌生的问题，最快破解这个问题的方法之一，就是数学实验。我们以 \(2022, ~ 2023\) 两年的新高考真题为例。

\section*{例 5.1 「 2022 年新高考 1 卷」}
已知函数 \(f^{\prime}(x)\) 及其导函数 \(f^{\prime}(x)\) 的定义域均为 \(\mathbb{R}\) ．记 \(g(x)=f^{\prime}(x)\) ，若 \(f\left(\frac{3}{2}-2 x\right), g(2+x)\) 均为偶函数，则\\
\(\qquad\) —。\\
A ，\(f(0)=0\)\\
B，\(g\left(-\frac{1}{2}\right)=0\)\\
C,\(f(-1)=f(4)\)\\
D，\(g(-1)=g(2)\)

导函数有奇偶性，这是相对陌生的条件。我们如何处理呢？\\
（1）

\[
f\left(\frac{3}{2}-2 x\right) \text { 是偶函数, 意味着 } f(x) \text { 本身是一个轴对称函数; }
\]

（2）

\[
g(x+2) \text { 也是偶函数, 说明 } f^{\prime}(x) \text { 也具有轴对称性。 }
\]

什么样的函数，\(f(x), f^{\prime}(x)\) 都具有轴对称性呢？我们熟悉的几大类函数中，除了常函数 \((f(x)=1\) 等），只有三角函数具有这个特点：对于 \(f(x)=\sin x\)\\
（1）\(f(x)=\sin x\) 关于 \(x=k \pi+\frac{\pi}{2}(k \in \mathbb{Z})\) 对称，\\
（2）\(f^{\prime}(x)=\cos x\) ，本身是一个偶函数。\\
因此，我们猜测，\(f(x)\) 可能可以是三角函数，或者和它差一个常数\\
三角函数本身具有周期性。我们看到 \(\frac{3}{2}, ~ 2\) ，马上想到，这个函数可能是一个以 2 为周期的函数。这样的三角函数可以是谁？显然 \(f(x)=\sin \pi x\) 是满足条件的。我们来验证一下，它是否满足题意：当 \(f(x)=\sin \pi x\) 时，\\
（1）是偶函数：

\[
f\left(\frac{3}{2}-2 x\right)=\sin \left(\frac{3}{2} \pi-2 \pi x\right)=-\sin \left(2 \pi x-\frac{\pi}{2}\right)=\cos (2 \pi x)
\]

（2）\(f^{\prime}(x)=\pi \cos (\pi x)\) ，所以 \(g(x+2)=\pi \cos (\pi(x+2))=\pi \cos (\pi x+2 \pi)=\pi \cos (\pi x+2 \pi) \pi \cos (\pi x)\) ，也是偶函数。它确实符合题意！\\
所以，我们可以利用 \(f(x)=\sin \pi x\) 来验证选项。为了让我们的答案更具普适性，我们可以考虑 \(f(x)=\sin \pi x+1\) 。\\
选项 A ：

\[
f(0)=\sin 0+1=1, \text { 错误; }
\]

选项 B：

\[
g\left(-\frac{1}{2}\right)=\pi \cos \left(-\frac{1}{2} \pi x\right)=0, \text { 可能正确; }
\]

选项 C ：

\[
f(-1)=0, f(4)=0 \text {, 可能正确; }
\]

选项 D：

\[
g(-1)=\pi \cos (-\pi)=-\pi, g(2)=\pi \cos (2 \pi)=\pi, \text { 错误。 }
\]

由于这是一个多选题，所以正确答案为 BC．\\
很多同学可能会有疑问：这难道不就是特值法吗？和我们代入具体的数，有什么区别吗？最大的区别在于，面对一个陌生的，新情境的题目，需要准确地判断应该代入什么例子，以及哪个例子是最合适的。这才是举例子的难点。以刚才的题目为例，我们并不是一股脑把各种函数代进去胡乱尝试，而是对题目条件有了相对准确的翻译。最关键的，是如下两个判断：\\
（1）\(f(x), f^{\prime}(x)\) 都有对称轴。什么样的函数具有对称轴呢？我们高中学过的几大函数：幂函数，指数函数，对数函数，三角函数中，只有三角函数和常函数具有这样的性质；\\
（2）上下平移一个函数图像，不改变对称轴。如果 \(f(x)\) 满足条件，那么 \(f(x)+100, f(x)-20\) ，也都满足条件。所以，我们不可能得到关于函数具体取值的信息。如果只有第一个判断，而代人了 \(f(x)=\sin \pi x\) ，可能会误选 A 。因此，第二个判断也是非常关键的，它帮助我们找到最合适的例子：\(f(x)=\sin \pi x+1\) ，补一个常数。

事实上，面对困难问题，举例子最大的作用，就是让一个不可做的问题变成一个可做的问题，让不熟悉变为熟悉。我们之所以被一个复杂的问题，往往就是因为我们没看到藏在它背后更简单的问题。再举一例，来体会这一思路：

例 5.2 已知 \(|k| \geq 1\) ，求 \(c=\sqrt{k^{2}+1} \cdot|x-k|+\sqrt{1+\frac{1}{k^{2}}} \cdot\left|x+\frac{1}{k}\right|\) 的最小值。

这是一个复杂的二元函数，有两个变量 \(x, k\) 。这是坏消息。但好消息是，这两个变量之间，没有额外的限制条件。换言之，它们都是自由变化的。我们在高中阶段，处理只有一个变量的问题又非常多的手段和工具：求导，不等式等。如果能够将问题化为一元函数，就可以解决。为此，那我们可以举一些简单的例子：让 \(x\) 或 \(k\) 固定，看看另一个变量怎么动：\\
（1）当 \(x=1\) 时，

\[
c=\sqrt{k^{2}+1} \cdot|1-k|+\sqrt{1+\frac{1}{k^{2}}} \cdot\left|1+\frac{1}{k}\right| .
\]

这仍是一个关于 \(k\) 的复杂函数，不好处理；\\
（2）当 \(k=1\) 时，

\[
c=\sqrt{2}|x-1|+\sqrt{2}|x+1|=\sqrt{2} \cdot(|x-1|+|x+1|) .
\]

后者是两个绝对值。根据绝对值的性质，我们可以直接看到

\[
|x-1|+|x+1| \geq 2
\]

因此 \(c \geq 2 \sqrt{2}\) ．我们成功求出了 \(k=1\) 时的最小值。注意，这不是 \(c\) 的最小值，是＂\(k=1\) 时＂，\(c\) 的最小值；\\
（3）当 \(k=2\) 时，

\[
\begin{aligned}
c & =\sqrt{5} \cdot|x-2|+\frac{\sqrt{5}}{2} \cdot\left|x+\frac{1}{2}\right| \\
& =\frac{\sqrt{5}}{2} \cdot\left(2|x-2|+\left|x+\frac{1}{2}\right|\right)
\end{aligned}
\]

又一次化成了绝对值函数。记 \(f(x)=2|x-2|+\left|x+\frac{1}{2}\right|\) 。和 \(k=1\) 时的区别在于，此时这两个绝对值前面的系数不同。但我们仍可以利用几何性质翻译：数轴上 \(x\) 对应的点为 \(A, 2\) 对应的点为 \(B,-\frac{1}{2}\) 对应的点为 \(C\) ，那么

\[
2|x-2|+\left|x+\frac{1}{2}\right|=2|A B|+|A C|
\]

如何求这个的最小值？显然，\(A\) 应该在 \(B, C\) 之间时取最小。此时

\[
2|A B|+|A C|=|A B|+(|A B|+|A C|)=|A B|+|B C|=|A B|+\frac{5}{2} .
\]

为了让它最小，必然需要 \(A, B\) 重合。所以最小值为 \(\frac{5}{2}\) ．\\
如果直接写不等式，也可以得到结果：直接利用 \(|x-2| \geq 0\) 即可，即

\[
2|x-2|+\left|x+\frac{1}{2}\right| \geq|x-2|+\left|x+\frac{1}{2}\right| \geq\left|2-\left(-\frac{1}{2}\right)\right|=\frac{5}{2}
\]

当我们有了这三个例子后，已经知道：对于给定的 \(k, c\) 会化为关于 \(x\) 的绝对值函数；利用几何意义或绝对值不等式，就可以得到答案。对于一般的 \(k\) ，我们也可以利用这个思路进行求解。

由于 \(|k| \geq 1\) ，所以

\[
\sqrt{k^{2}+1} \geq \sqrt{1+\frac{1}{k^{2}}}
\]

所以第一个绝对值的系数大。这样一来：

\[
\begin{aligned}
c & =\sqrt{k^{2}+1} \cdot|x-k|+\sqrt{1+\frac{1}{k^{2}}} \cdot\left|x+\frac{1}{k}\right| \\
& \geq \sqrt{1+\frac{1}{k^{2}}} \cdot\left(|x-k|+\left|x+\frac{1}{k}\right|\right) \\
& \geq \sqrt{1+\frac{1}{k^{2}}} \cdot\left|k+\frac{1}{k}\right|=\frac{\sqrt{k^{2}+1}}{|k|} \cdot \frac{k^{2}+1}{|k|}=\frac{\left(\sqrt{k^{2}+1}\right)^{3}}{k^{2}} .
\end{aligned}
\]

这样，就得到了一个关于 \(k\) 的函数，可以求导求最小值。如果担心求导复杂，可以先换元：令 \(\sqrt{k^{2}+1}=m\) ，则

\[
c \geq \frac{m^{3}}{m^{2}-1} .
\]

其中 \(m \geq 1\) ．记 \(g(m)=\frac{m^{3}}{m^{2}-1}\) ，则

\[
g^{\prime}(m)=\frac{3 m^{2}\left(m^{2}-1\right)-m^{3} \cdot 2 m}{\left(m^{2}-1\right)^{2}}=\frac{m^{4}-3 m^{2}}{\left(m^{2}-1\right)^{2}}
\]

令 \(g^{\prime}(m)=0\) ，由于 \(m^{2} \geq 1\) ，则 \(m^{2}=3\) ，所以 \(g(m)\) 在 \([1, \sqrt{3})\) 上单调递增，在 \((\sqrt{3},+\infty)\) 上单调递减，所以

\[
g(m) \geq g(\sqrt{3})=\frac{3 \sqrt{3}}{2} .
\]

这就是 \(c\) 的最小值。

在解决向量，三角形的问题时，我们也经常使用举例子的方法。对于一般的三角形，如果问题不好处理，我们可以选择特殊的三角形。

例 5.3 （2022 新高考 1 卷）在三角形 \(A B C\) 中，点 \(D\) 在边 \(A B\) 上，\(B D=2 D A\) ，记 \(\overrightarrow{C A}=\) \(\qquad\) －\\
A， \(3 \vec{m}-2 \vec{n}\)\\
B ，\(-2 \vec{m}+3 \vec{n}\)\\
C， \(3 \vec{m}+2 \vec{n}\)\\
D， \(2 \vec{m}+3 \vec{n}\)

题目本身并不难，如何快速得到结果？题目要求我们用 \(\bar{C} A, \overline{C D}\) 表示 \(\bar{C} B\) ，如果让为 \(C\) 原点， \(\overrightarrow{C A}, \overrightarrow{C D}\) 刚好就是平面直角坐标系中的基底，我们只需求出 \(B\) 点坐标即可。\\
因此，考虑以 \(C\) 为原点，建立平面直角坐标系，并令 \(A(1,0), D(0,1)\) ．\\
由于 \(D\) 在边 \(A B\) 上，且 \(B D=2 A D\) ，所以

\[
\overrightarrow{D B}=2 \overrightarrow{A D}=(-2,2)
\]

所以

\[
\overrightarrow{C B}=\overrightarrow{C D}+\overrightarrow{D B}=(-2,3)=-2 \vec{m}+3 \vec{n}
\]

也可以直接利用平面几何关系：过 \(B\) 作 \(x\) 轴垂线，垂足为 \(E\) ，那么 \(B E=D E=2\) ．所以 \(B(-2,3)\) ．所以 \(\vec{C} B=\) \((-2,3)=-2 \vec{m}+3 \vec{n}\) ．选 B．

在这个题中，我们相当于举了一个特殊的例子：让三角形 \(C A D\) 刚好是等腰直角三角形。为什么这样也可以得到正确的结果呢？这体现了我们举例子时运用的一个重要原理：维持题目本身的＂核心结构＂不受破坏。这个题目中，只要比例关系，共线关系不变，那么向量表出也不会发生变化，所以我们可以任意选择特殊形状的三角形，来方便运算。这个方法在复杂的题目中，威力更大。

\section*{5．1．2 初步建立知识体系}
反思刚才举例的过程。我们能够快速找到三角函数的例子，同样离不开我们对这部分知识本身的熟练掌握（类比第一章的元问题）。事实上，不光是破题，想要掌握这些抽象的知识，也需要大量的例子。高中数学之所以困难，最大的痛点就在于其抽象性。在初中阶段，我们学习二次方程，平面几何，整体而言非常直观，所以理解起来相对容易。但是到高中，我们需要面对大量的抽象概念：函数的定义域，对应关系，值域，繁多的三角函数公式，复杂的递推数列，导数等等。想要理解好这些知识，不光需要抽象思维。大量的例子加以辅助是非常有必要的。

我们仍以三角函数为例。三角函数有大量相关的公式，我们怎么才能不死记硬背地就能理解他们呢？我们学习过三角函数的和角公式：

\[
\sin (\alpha+\beta)=\sin \alpha \cos \beta+\sin \beta \cos \alpha .
\]

但看这个公式，我们可能看不出什么感觉，甚至会怀疑，它可能是错的。所以，为了让大家确信它＂应该是对的＂，我们先代人几个具体的数看一看：\\
（1）令 \(\beta=0\) ，左式变为 \(\sin \alpha\) ，右式变为

\[
\sin \alpha \cos 0+\sin 0 \cos \alpha=\sin \alpha
\]

成立；\\
（2）令 \(\beta=-\alpha\) ，左式变为 \(\sin (\alpha-\alpha)=0\) ，右式变为

\[
\sin \alpha \cos (-\alpha)+\sin (-\alpha) \cos \alpha=\sin \alpha \cos \alpha-\sin \alpha \cos \alpha=0 .
\]

成立。\\
到这里，我们至少不会怀疑它＂那么离谱＂了。但是它的正确性，还需要更多的体会。我们代入一些具体的数看看：\\
（3）令 \(\beta=\pi\) ，左式变为 \(\sin (\alpha+\pi)\) ，右式变为

\[
\sin \alpha \cos \pi+\sin \pi \cos \alpha=-\sin \alpha .
\]

所以，应有 \(\sin (\alpha+\pi)=-\sin \alpha\) ．哦！这就是关于 \(\pi\) 的诱导公式！\\
（4）令 \(\beta=-\frac{\pi}{2}\) ，左式变为 \(\sin \left(\alpha-\frac{\pi}{2}\right)\) ，右式变为

\[
\sin \alpha \cos \left(-\frac{\pi}{2}\right)+\cos \alpha \sin \left(-\frac{\pi}{2}\right)=-\cos \alpha .
\]

所以，应有 \(\sin \left(\alpha-\frac{\pi}{2}\right)=-\cos \alpha\) ．哦！这就是关于 \(\frac{\pi}{2}\) 的诱导公式！\\
（5）令 \(\beta=\alpha\) ，左式变为 \(\sin (\alpha+\alpha)=\sin 2 \alpha\) ，右式变为

\[
\sin \alpha \cos \alpha+\cos \alpha \sin \alpha=2 \sin \alpha \cos \alpha
\]

所以，应有

\[
\sin 2 \alpha=2 \sin \alpha \cos \alpha .
\]

哦！这就是二倍角公式！\\
（6）令 \(\beta=\frac{\pi}{3}\) ，左式变为 \(\sin \left(\alpha+\frac{\pi}{3}\right)\) ，右式变为

\[
\sin \alpha \cos \frac{\pi}{3}+\cos \alpha \sin \frac{\pi}{2}=\frac{1}{2} \sin \alpha+\frac{\sqrt{3}}{2} \cos \alpha
\]

所以，应有

\[
\sin \left(\alpha+\frac{\pi}{3}\right)=\frac{1}{2} \sin \alpha+\frac{\sqrt{3}}{2} \cos \alpha .
\]

哦！这就是辅助角公式！根据辅助角公式；令 \(a=\frac{1}{2}, b=\frac{\sqrt{3}}{2}\) ，那么

\[
\sin \varphi=\frac{b}{\sqrt{a^{2}+b^{2}}}=\frac{\sqrt{3}}{3}, \cos \varphi=\frac{a}{\sqrt{a^{2}+b^{2}}}=\frac{1}{2},
\]

取 \(\varphi=\frac{\pi}{3}\) ，那么

\[
\frac{1}{2} \sin \alpha+\frac{\sqrt{3}}{2} \cos \alpha=a \sin \alpha+b \cos \alpha=\sqrt{a^{2}+b^{2}} \sin (\alpha+\varphi)=\sin \left(\alpha+\frac{\pi}{3}\right) .
\]

这就是我们刚才得到的公式。想想看，你会证明辅助角公式了吗？\\
（7）\(\beta=-\frac{\alpha}{2}\) ，左式为 \(\sin \left(\alpha-\frac{\alpha}{2}\right)=\sin \frac{\alpha}{2}\) ，右式变为

\[
\sin \alpha \cos \frac{\alpha}{2}+\sin \left(-\frac{\alpha}{2}\right) \cos \alpha=\sin \alpha \cos \frac{\alpha}{2}-\sin \frac{\alpha}{2} \cos \alpha .
\]

因此， \(\sin \frac{\alpha}{2}=\sin \alpha \cos \frac{\alpha}{2}-\sin \frac{\alpha}{2} \cos \alpha\) 。我们把它整理成关于 \(\sin \frac{\alpha}{2}, \cos \frac{\alpha}{2}\) 的式子：

\[
(1+\cos \alpha) \sin \frac{\alpha}{2}=\sin \alpha \cos \frac{\alpha}{2}
\]

两边同时除以 \(\cos \frac{\alpha}{2}\) ，得

\[
\tan \frac{\alpha}{2}=\frac{\sin \alpha}{1+\cos \alpha} .
\]

哦！这就是半角公式；\\
（8）令 \(\beta=\frac{\pi}{2}-\alpha\) ，则左式化为 \(\sin \left(\alpha+\frac{\pi}{2}-\alpha\right)=\sin \left(\frac{\pi}{2}\right)=1\) ．右式化为

\[
\sin \alpha \cos \left(\frac{\pi}{2}-\alpha\right)+\sin \left(\frac{\pi}{2}-\alpha\right) \cos \alpha=\sin ^{2} \alpha+\cos ^{2} \alpha
\]

所以，应有 \(\sin ^{2} \alpha+\cos ^{2} \alpha=1\) 。哦！这就是勾股定理啊！勾股定理，居然也可以从和角公式出发得到！\\
仅仅一个和角公式，我们通过代入一些简单的值，就推出了海量的结论！诱导公式，二倍角公式，辅助角公式，半角公式，甚至是勾股定理，全都可以通过它来理解。所以，我们学习了那么多三角函数公式，其实都是一个公式：和角公式。只需要记住这一个公式，就掌握了所有公式。即使考场上无法马上回忆起，也可以快速推导。不仅如此，我们还把所有的公式，通过和角公式串联起来：\\
（1）诱导公式：让 \(\beta=\frac{\pi}{2}, \pi,-\frac{\pi}{2},-\pi\) 等，即可得到；\\
（2）二倍角公式：让 \(\beta=\alpha\) 即可得到；\\
（3）辅助角公式：和和角公式互为＂逆运算＂，可以用和角公式直接证明；\\
（4）半角公式：让 \(\beta=-\frac{\alpha}{2}\) 即可得到；\\
（5）勾股定理：让 \(\beta=\frac{\pi}{2}-\alpha\) 。\\
试试看，让 \(\beta\) 取不同的值，你还能得到哪些等式？\\
我们重新理解和角公式。它真正想要解决的问题，当角度相加时，三角函数是如何相加的。对于一般的实数而言，\(\alpha+\beta\) 直接计算即可；但当它们并非特殊角，我们只知道 \(\alpha, \beta\) 的三角函数值时，如何求 \(\alpha+\beta\) 的三角函数值呢？这就是这个公式的作用。它告诉我们，不能简单地将正弦函数相加：

\[
\sin (\alpha+\beta) \neq \sin \alpha+\sin \beta
\]

需要给每个正弦函数，配一个＂权重＂：

\[
\sin (\alpha+\beta)=\cos \beta \cdot \sin \alpha+\cos \alpha \cdot \sin \beta .
\]

而这个权重，就是另外一个角的余弦值。\\
这也是为什么，和角公式又名加法公式。它的地位，相当于三角函数的＂加法＂。我们所有的三角函数公式，都可以理解成加法，减法：\\
（1）诱导公式：告诉你 \(\alpha+\frac{\pi}{2}, \alpha+\pi\) 怎么相加；\\
（2）二倍角公式：告诉你 \(\alpha+\alpha\) 怎么相加，\\
（3）半角公式：告诉你 \(\alpha+\left(-\frac{\alpha}{2}\right)\) 怎么相加。

\section*{5．1．3 思想试验在选填中的运用}
XXX

例5．4（2022 新高考 II 卷）已知函数 \(f(x)\) 的定义域为 \(\mathbb{R}\) ，且 \(f(x+y)+f(x-y)=f(x) f(y), f(1)=1\) ，则 \(\sum_{k=1}^{22} f(k)=\)\\
\(\qquad\) －。

【方法一】观察这个表达式，哪个函数可能具有这样的性质？\\
（1）指数函数：指数函数满足 \(e^{x} \cdot e^{y}=e^{x+y}\) ，不会出现 \(x-y\) 项；\\
（2）对数函数：对数函数对乘积 \(\ln x \ln y\) 没有特别的运算规律；\\
（3）幂函数：如果是 \(f(x)=x^{\alpha}\) ，那么左侧应该是 \(\alpha\) 次，右侧是 \(2 \alpha\) 次，不成立；\\
（4）三角函数：何时会出现 \(\sin x \sin y, \cos x \cos y\) 这样的式子？在余弦的和角，差角公式中会见到，所以我们可以考虑尝试 \(f(x)=\cos x\) ．发现：

\[
\begin{gathered}
\cos (x+y)+\cos (x-y)=\cos x \cos y-\sin x \sin y+\cos x \cos y+\sin x \sin y \\
=2 \cos x \cos y
\end{gathered}
\]

然而，差了一个倍数，所以我们要进行修正。注意到左侧是关于 \(f\) 的一次式，右侧是关于 \(f\) 的二次式，可以利用这个次数的差来调整：将左右两侧同时乘 2 ，可得

\[
2 \cos (x+y)+2 \cos (x-y)=(2 \cos x) \cdot(2 \cos y)
\]

所以 \(f(x)=2 \cos x\) 符合题意。\\
一般地，对 \(x\) 修改倍数，仍满足：\(f(x)=2 \cos \omega x\) 均符合第一个条件。由于 \(f(1)=1\) ，所以 \(\cos \omega=\frac{1}{2}\) ．可令 \(\omega=\frac{\pi}{3}\) ．所以 \(f(x)=2 \cos \left(\frac{\pi}{3} x\right)\) 满足题意。可以直接利用它计算结果：

\[
\sum_{k=1}^{22} f(k)=2 \cos \frac{\pi}{3}+2 \cos \frac{2 \pi}{3}+\ldots+2 \cos \left(\frac{22 \pi}{3}\right)
\]

计算时，可以适当利用周期性：\(f(x)\) 的最小正周期为 6 ，

\[
\begin{aligned}
& f(1)=2 \cos \frac{\pi}{3}=1, f(2)=2 \cos \frac{2 \pi}{3}=-1, f(3)=2 \cos \frac{3 \pi}{3}=-2 \\
& f(4)=2 \cos \frac{4 \pi}{3}=-1, f(5)=2 \cos \frac{5 \pi}{3}=1, f(6)=2 \cos \frac{6 \pi}{3}=2
\end{aligned}
\]

所以必然有

\[
f(1)+f(2)+\ldots+f(6)=\frac{1}{2}-\frac{1}{2}-1-\frac{1}{2}+\frac{1}{2}+1=0
\]

所以前 24 项和为 0 ，故

\[
\sum_{k=1}^{22} f(k)=-f(23)-f(24)=-f(5)-f(6)=-3
\]

【方法二】由于我们只关心 \(f(x)\) 在整数处的取值，所以我们可以令 \(y=1\) ，将问题化成一个数列问题。此时

\[
f(x+1)+f(x-1)=f(x)
\]

所以，对于任意的正整数 \(n\) ，均有

\[
f(n+1)+f(n-1)=f(n)
\]

如果觉得数列的记号看起来更舒服，可以令 \(a_{n}=f(n)\) ，那么

\[
a_{n+1}+a_{n-1}=a_{n}
\]

且我们知道 \(a_{1}=1\) ．想要求出 \(a_{2}\) ，我们仍需要知道 \(a_{0}\) 值。令 \(x=1, y=0\) ，得 \(f(0)=2\) ．所以 \(a_{0}=2\) ．这样一来，就可以写出数列中每一项的值并计算。

\[
\begin{gathered}
a_{2}=a_{1}-a_{0}=-1 \\
a_{3}=a_{2}-a_{1}=-2 \\
a_{4}=a_{3}-a_{2}=-1 \\
a_{5}=a_{4}-a_{3}=1 \\
a_{6}=a_{5}-a_{4}=2 \\
a_{7}=a_{6}-a_{5}=1 \\
a_{8}=a_{7}-a_{6}=-1
\end{gathered}
\]

我们发现，\(a_{6}=a_{0}, a_{7}=a_{1}\) ，由于这个数列是二阶的（ \(a_{n}\) 只和它的前两项有关），所以，后面的项都是 6 个一循环。所以

\[
\begin{aligned}
a_{1}+a_{2}+\ldots+a_{22} & =\left(a_{1}+a_{2}+\ldots+a_{24}\right)-a_{23}-a_{24} \\
& =4\left(a_{1}+a_{2}+\ldots+a_{6}\right)-a_{5}-a_{6} \\
& =4 \cdot(2+1-1-2-1+1)-3 \\
& =-3
\end{aligned}
\]

【注】我们也可以从数列本身的性质出发证明 \(a_{n}\) 是周期为 6 的数列。由

\[
\begin{gathered}
a_{n+2}+a_{n}=a_{n+1} \\
a_{n+3}+a_{n+1}=a_{n+2}
\end{gathered}
\]

相加，得

\[
a_{n+3}+a_{n}=0
\]

即 \(a_{n+3}=-a_{n}\) ．换言之，每往后推三项，变为相反数。那么往后推六项，就变回原来的项：

\[
a_{n+6}=-a_{n+3}=a_{n}
\]

所以它以 6 为周期。

\section*{5．1．4 大题中的思想试验}
\section*{XXX}
例 5.5 （2021 年新高考 1 卷）在平面直角坐标系 \(x O y\) 中，\(F_{1}(-\sqrt{17}, 0), F_{2}(\sqrt{17}, 0),\left|M F_{1}\right|-\left|M F_{2}\right|=2\) ．点 \(M\) 的

轨迹为 \(C\) 。\\
（1）求 \(C\) 的方程；\\
（2）设点 \(T\) 在直线 \(x=\frac{1}{2}\) 上，过 \(T\) 的两条直线分别交 \(C\) 于 \(A, B\) 两点和 \(P, Q\) 两点。且 \(|T A| \cdot|T B|=|T P| \cdot|T Q|\) ．求直线 \(A B\) ，直线 \(P Q\) 斜率之和。

对于（1），根据题意，\(C\) 为双曲线的右支。且 \(2=2 a, c=\sqrt{17}\) ．所以 \(a=1, c=\sqrt{17}, b=4\) ，所以 \(C\) 的方程为

\[
x^{2}-\frac{y^{2}}{16}=1(x>0) .
\]

（2）可以考虑特殊情况。当 \(T\) 恰好是 \(\left(\frac{1}{2}, 0\right)\) ，且 \(A B, P Q\) 关于 \(x\) 轴对称时，符合题意。\\
这是因为：整个图形时完全对称的，即 \(T A=T P, T B=T Q\) ．此时，\(k_{A B}+k_{P Q}=0\) ．\\
故答案应为 0 ．设 \(T\left(\frac{1}{2}, y_{0}\right)\) ．设 \(A B, P Q\) 的斜率为 \(k_{1}, k_{2}\) ，那么：

\[
\begin{aligned}
& l_{A B}: y=k_{1}\left(x-\frac{1}{2}\right)+y_{0} \\
& l_{P Q}: y=k_{2}\left(x-\frac{1}{2}\right)+y_{0}
\end{aligned}
\]

和双曲线方程联立，得

\[
\begin{gathered}
x^{2}-16 y^{2}=16 \\
x^{2}-16\left(k x+y_{0}-\frac{1}{2} k\right)^{2}=16 \\
\left(1-16 k^{2}\right) x^{2}+16 k\left(k-2 y_{0}\right) x-16\left(y_{0}-\frac{1}{2} k\right)^{2}-16=0
\end{gathered}
\]

由于 \(x_{1}, x_{2}>1\) ，那么

\[
\begin{aligned}
|T A| \cdot|T B| & =\sqrt{k_{1}^{2}+1} \cdot\left(x_{1}-\frac{1}{2}\right) \cdot \sqrt{k_{1}^{2}+1} \cdot\left(x_{2}-\frac{1}{2}\right) \\
& =\left(k_{1}^{2}+1\right)\left(x_{1} x_{2}-\frac{1}{2}\left(x_{1}+x_{2}\right)+\frac{1}{4}\right) \\
& =\frac{k_{1}^{2}+1}{1-16 k_{1}^{2}} \cdot\left(-16\left(y_{0}-\frac{1}{2} k\right)^{2}-16+8 k\left(k-2 y_{0}\right)+\frac{1}{4}-4 k^{2}\right) \\
& =\frac{k_{1}^{2}+1}{1-16 k_{1}^{2}} \cdot\left(-4 k^{2}+16 k y_{0}-16 y_{0}^{2}-16 k y_{0}+8 k^{2}+\frac{1}{4}-4 k^{2}\right) \\
& =\frac{k_{1}^{2}+1}{1-16 k_{1}^{2}} \cdot\left(\frac{1}{4}-16 y_{0}^{2}\right)
\end{aligned}
\]

同理，

\[
|T P| \cdot|T Q|=\frac{k_{2}^{2}+1}{1-16 k_{2}^{2}} \cdot\left(\frac{1}{4}-16 y_{0}^{2}\right)
\]

由于 \(|T A| \cdot|T B|=|T P| \cdot|T Q|\) ，所以

\[
\frac{k_{1}^{2}+1}{1-16 k_{1}^{2}}=\frac{k_{2}^{2}+1}{1-16 k_{2}^{2}}
\]

化简得

\[
\begin{gathered}
-16 k_{1}^{2} k_{2}^{2}+k_{1}^{2}-16 k_{2}^{2}+1=-16 k_{1}^{2} k_{2}^{2}+k_{2}^{2}-16 k_{1}^{2}+1 \\
17 k_{1}^{2}=17 k_{2}^{2} \\
17\left(k_{1}-k_{2}\right)\left(k_{1}+k_{2}\right)=0
\end{gathered}
\]

所以 \(k_{1}+k_{2}=0\) ．

\section*{总结：举例子为什么这么有用？}
举例子，看似大家都会使用的基本方法，到底为什么能够发挥这么大的功效呢？\\
（1）它将不可做的问题，变成可做的问题。举例子让复杂，抽象的情况，变得简单，具体。例如：复杂的二元函数问题，通过举例子，变成了简单的一元函数绝对值问题。一个不会的问题，背后往往藏着一个更简单，你依旧不会的问题。只有把这个更简单的问题解决，才能够解决现在正在面对的问题。\\
（2）虽非严谨证明，但是体会一个新问题的必要准备。举例子，可以非常好地帮助你找到解决问题的感觉：这个条件意味着什么？是一个很松的条件，还是很苛刻的条件？它可能和我学过的什么知识产生关联？这些都可以通过具体的例子来体会。体会好这些感觉后，我们会对问题本身做一些判断：某某结论应成立；某某应不成立，有反例，等等。有了这些判断，才有了严格的证明。这是＂破题＂非常关键的步骤。

\section*{5.2 遍历与临界}
\section*{5．2．1 极端遍历法}
我们在举例子时，有一些特殊情况是我们经常需要考虑的。例如，多个变量的问题中，我们看当他们都相等的时候会发生什么：如均值不等式，取等条件就是各个变量均相等。除了这种对称的情况，我们还会考虑极端情况。什么是极端情况？就是让变量取到最大或者最小的情况。

\section*{5．2．1．1 利用极端遍历法画函数草图}
我们在学习函数时，学习了分子，分母是一次函数或者二次函数的情形。例如：\(y=\frac{1}{x(x-1)}\) ．在刚刚学习它的时候，难免会觉得它非常抽象：它既不是我们熟悉的二次函数，一次函数，也不是我们刚学会的有理分式函数，那么它到底＂长什么样＂呢？

极端遍历法，可以有效帮助你＂定性＂地画出这个函数的图像。考察一个函数，我们首先关心的是定义域。为了让 \(y\) 有意义，要求分母 \(x(x-1) \neq 0\) 。所以定义域为 \(\{x \mid x \neq 0\) 且 \(x \neq 1\}\) 。定义域被分成了三段：\((\infty, 0) \cup(0,1) \cup(1,+\infty)\) ．

对于这个函数而言，什么是极端状态呢？显然，当 \(x\)＂特别大＂或者 \(x\)＂特别小＂的时候，都是极端状态。 （1）当 \(x\)＂特别大＂的时候，比如，我们让 \(\mathrm{x}=1000\) ，那么，

\[
x(x-1)=1000 \times(1000-1) \approx 100000=10^{6} .
\]

这是一个非常的大数了！所以，\(y=\frac{1}{x(x-1)}\) 会变得非常大，那么 \(y\) 就会是一个非常小的数。虽然它非常小，但它仍是一个正数。所以它的图像，就会像滑滑梯一样，慢慢滑向 \(x\) 轴，但永远在 \(x\) 轴的上方。此时，\(x\) 轴（直线 \(y=0\) ）就是它的一条渐近线。\\
（2）当 \(x\)＂特别小＂的时候。这里的小，指的是在数轴上特别靠左：即，它是一个绝对值很大的负数。比如，我们让 \(x=-1000\) ，那么，

\[
x(x-1)=-1000 \times(-1000-1) \approx 10^{6} .
\]

它仍然是一个非常的大的数。所以，和 O—样，当 \(x\)＂特别小＂时，它的图像像一个滑梯。不过，此时的滑梯向左滑向 \(x\) 轴，仍在 \(x\) 轴上方。

这样，我们通过简单地代入了两个极端的数，确定了远处 y 的图像。那么，中间的图像，长什么样呢？

由于定义域被分成了三段，我们依次考虑。\\
（3）\(x>1\) 的情况。 \(x \rightarrow+\infty\) 的情况我们已经考虑清楚了，另一个极端情况，就是当 \(x\) 离 1 非常近的时候。虽然 \(x\) 不能取 1 ，但是它离 1 要多近，有多近，所以我们，可以考虑 \(x=1.001\) ，即 \(x=1+10^{-3}\) 。此时，

\[
x(x-1)=\left(1+10^{-3}\right) \times 10^{-3}=0.001001 \approx 10^{-3}
\]

此时，

\[
y=\frac{1}{x(x-1)} \approx 1000
\]

它也是一个非常大的数。所以，当 \(x\) 离 1 特别近的时候，\(y\) 会是一个特别大的正数。因此，\(y\) 的图像会变得非常陡。根据单调性，\(x(x-1)>0\) 在 \((1,+\infty)\) 上单调递增，所以 \(y=\frac{1}{x(x-1)}\) 单调递减。因此，在 \((1,+\infty)\) 上，\(y\) 的图像，就像一个从特别高的地方滑下来的滑梯，最后越来越缓，离 \(x\) 轴越来越近。\\
（4）\(x \in(0,1)\) 的情况。我们考虑当 \(x\) 离 0,1 特别近的时候的情况。若 \(x\) 离 1 很近，例如 \(x=0.999=1-10^{-3}\) 那么

\[
x(x-1)=0.999 \times(-0.001) \approx-0.001
\]

千万不要小看这个绝对值很小的负数：它直接导致 \(y\) 变成一个绝对值很大的负数：

\[
y=\frac{1}{x(x-1)} \approx-1000
\]

所以，当 \(x\) 离 1 很近时，会迅速递减，像一个很陡峭的＂峭壁＂。类似地，当 \(x\) 离 0 很近时，例如 \(x=0.001\) ，那么

\[
x(x-1)=0.001 \times(-0.999) \approx-0.001
\]

它也会让 \(y\) 变成一个绝对值很大的负数：

\[
y=\frac{1}{x(x-1)} \approx-1000
\]

因此，在 0 一侧，也会是这样的＂峭壁＂。\\
（5）\(x \in(-\infty, 0)\) 的情况。我们已经考虑了 \(x \rightarrow-\infty\) 的情况，只需考虑 \(x\) 离 0 较近的情况。如果 \(x\) 是一个特别小的负数，例如 \(x=-0.001\) ，那么

\[
x(x-1)=-0.001 \times(-1.001) \approx 0.001
\]

因此，

\[
y=\frac{1}{x(x-1)} \approx 1000
\]

这样，\(y\) 会变成很大的正数。所以，这一侧会像一个很陡峭的坡，一下子爬上去\\
当我们考虑了这些极端情况后，我们发现：\(y\) 的图像应为：\\
\includegraphics[max width=\textwidth, center]{2025_02_16_7c9aaf92bc99dc07f048g-138}

事实上，分母 \(y_{1}=x(x-1)\) 是一个二次函数，它的对称轴是 \(x=\frac{1}{2}\) ；

\[
y_{1}=x^{2}-x=\left(x-\frac{1}{2}\right)^{2}-\frac{1}{4} .
\]

所以，整个函数图像关于 \(x=\frac{1}{2}\) 对称。这也和我们前面的分析一致。\\
当我们有了极端遍历法这个视角，我们可以理解好必修一中学到的各种各样的函数。指数函数 \(e^{x}\) ，幂函数 \(x^{n}(n>0)\) ，对数函数 \(\ln x\) ，当 \(x\) 趋向于无穷的时候，有什么关系呢？根据教材的图示，当 \(x \rightarrow+\infty\) 时，\(e^{x}\) 增长远快于 \(x^{n}(n>0), x^{n}\) 的增长远快于 \(\ln x\) 。下图中，绿色的线是 \(y=e^{x}\) ，橘色是 \(y=x^{2}\) ，红色是 \(y=x\) ，蓝色是 \(y=\ln x\) 。\\
\includegraphics[max width=\textwidth, center]{2025_02_16_7c9aaf92bc99dc07f048g-139}

而当 \(x \rightarrow-\infty\) 时，会发生这样的情况。\\
（1）对于指数函数 \(y=e^{x}\) ：如果 \(x=-1000\) ，那么

\[
y=e^{-1000}=\frac{1}{e^{1000}} \approx 5.1 \times 10^{-435}
\]

是一个非常小的正数。所以它会像滑滑梯一样，滑向 \(x\) 轴；\\
（2）对于幂函数 \(y=x^{n}\) ，需要分类讨论。若 \(n\) 为偶数，负负得正，我们会得到一个巨大的正数，所以 \(y\) 会趋近于正无穷；若 \(n\) 为奇数，则会趋近于负无穷。

当理清楚他们之间的增长速度关系，就可以有效地画出很复杂的函数的草图：\\
（1）\(f(x)=e^{x}-3 x\) 。当 \(x \rightarrow+\infty\) 时，\(e^{x}\) 增长远快于 \(3 x\) ，所以 \(f(x) \rightarrow+\infty ; x \rightarrow-\infty\) 时，\(e^{x} \rightarrow 0\) ，而 \(-3 x \rightarrow+\infty\) ，所以 \(f(x) \rightarrow+\infty\) ；\\
（2）\(f(x)=e^{x}-x^{2}\) 。当 \(x \rightarrow+\infty\) 时，\(e^{x}\) 增长远快于 \(x^{2}\) ，所以 \(f(x) \rightarrow+\infty\) ；当 \(x \rightarrow-\infty\) 时，\(e^{x} \rightarrow 0\) ，而 \(-x^{2} \rightarrow-\infty\) ，所以 \(f(x) \rightarrow+\infty\) ；\\
（3）\(f(x)=x-\ln x\) ．当 \(x \rightarrow+\infty\) 时，\(x\) 增长速度远快于 \(\ln x\) ，所以 \(f(x) \rightarrow+\infty\) ；当 \(x \rightarrow 0\) 时， \(\ln x \rightarrow-\infty\) ，所以 \(f(x) \rightarrow+\infty\)

极端遍历法，还可以让你直接得到零点等信息，我们会在例题中看到它的具体应用。

\section*{5．2．1．2 几何问题中的极端遍历法}
设 \(A(-1,0)\) ．动点 \(M\) 在射线 \(y=x(x>0)\)（不含端点）上运动。过 \(M\) 作 \(x\) 轴垂线，垂足为 \(N\) ．如果我们想要求 \(\frac{|A M|}{|M N|}\) 的取值范围，应该如何处理呢？\\
\includegraphics[max width=\textwidth, center]{2025_02_16_7c9aaf92bc99dc07f048g-140(1)}

最简单的方法，我们可以把 \(M\) 点的坐标设出来。设 \(M(x, x)\) ．那么

\[
\begin{aligned}
|A M|=\sqrt{(x+1)^{2}+x^{2}} & =\sqrt{2 x^{2}+2 x+1}, \\
|M N| & =x .
\end{aligned}
\]

所以

\[
\frac{|A M|}{|M N|}=\frac{\sqrt{2 x^{2}+2 x+1}}{x} .
\]

为了求这个函数的最小值，我们可以利用第二章中的内容，将它化作 \(t=\frac{1}{x}\) 的函数：

\[
\frac{\sqrt{2 x^{2}+2 x+1}}{x}=\sqrt{\frac{1}{x^{2}}+\frac{2}{x}+2}=\sqrt{\left(\frac{1}{x}+1\right)^{2}+1} .
\]

由于 \(x>0\) ，所以 \(\frac{1}{x}>0, \frac{1}{x}+1>1\) ，所以

\[
\sqrt{\left(\frac{1}{x}+1\right)^{2}+1}>\sqrt{2}
\]

而当 \(x \rightarrow 0\) 时，\(\frac{1}{x} \rightarrow+\infty\) ，所以 \(\frac{|A M|}{|M N|} \rightarrow+\infty\) ．所以 \(\frac{|A M|}{|M N|}\) 的取值范围是 \((\sqrt{2},+\infty)\)\\
利用代数方法和简单的极端遍历法，我们得到了答案。但是利用极端遍历法，事实上可以直接从几何上得到答案。由于 \(M\) 在射线上运动，所以 \(M\) 可以无限靠近 \(O\) ，可以无限远离。考虑这两种极端情况：\\
（1）当 \(M\) 无限接近 \(O\) 时，\(|M N| \rightarrow 0,|A M| \rightarrow 1\) ，所以 \(\frac{|A M|}{|M N|} \rightarrow+\infty\) ．换言之，\(\frac{|A M|}{|M N|}\) 要多大有多大；\\
\includegraphics[max width=\textwidth, center]{2025_02_16_7c9aaf92bc99dc07f048g-140}\\
（2）当 \(M\) 无限远离 \(O\) 时，即 \(x \rightarrow+\infty\) 时，三角形 \(A M N\) 会无限接近一个等腰三角形：\(\angle M A N\) 会无限接近 \(45^{\circ}\) ，但小于 \(45^{\circ}\) ，因为

\[
\tan \angle M A N=\frac{|M N|}{|A N|}=\frac{x}{x+1} \rightarrow 1
\]

所以，\(\frac{|A M|}{|M N|}\) 会无限接近 \(\sqrt{2}\) ，但永远比 \(\sqrt{2}\) 大。\\
\includegraphics[max width=\textwidth, center]{2025_02_16_7c9aaf92bc99dc07f048g-141}

因此，\(\frac{|A M|}{|M N|}\) 可以取遍 \((\sqrt{2},+\infty)\) 的所有实数，它的取值范围是 \((\sqrt{2},+\infty)\) 。\\
我们考虑更复杂的情形：\(A(-1,0)\) ，此时 \(M\) 在双曲线的一支 \(C: x^{2}-y^{2}=1(x>0)\) 上运动，且不在 \(x\) 轴上。过 \(M\) 作 \(x\) 轴垂线，垂足为 \(N\) ，求 \(\frac{|A M|}{|M N|}\) 的取值范围。\\
\includegraphics[max width=\textwidth, center]{2025_02_16_7c9aaf92bc99dc07f048g-141(1)}

如果利用极端遍历法，仍可以直接得到答案。设 \(M(x, y)\) ，\\
（1）当 \(M\) 无西安接近 \(x\) 轴的时，\(|M N| \rightarrow 0\) ，而 \(|A M| \rightarrow 2\) ，所以 \(\frac{|A M|}{|M N|} \rightarrow+\infty\) ，要多大又多大；\\
\includegraphics[max width=\textwidth, center]{2025_02_16_7c9aaf92bc99dc07f048g-141(2)}\\
（2）当 \(M\) 无限远离 \(x\) 轴时，由于 \(C\) 的渐近线 \(y= \pm x\) ，所以当 \(x\) 极大时，在计算 \(\frac{|A M|}{|M N|}\) 时，我们可以认为 \(M\) 在渐近线 \(y= \pm x\) 上。这就化成了一开始直线的情形：当 \(x \rightarrow+\infty\)\\
时，

\[
\tan \angle M A N=\frac{x}{x+1} \rightarrow 1
\]

所以，\(\angle M A N \rightarrow 45^{\circ}, \frac{|A M|}{|M N|} \rightarrow \sqrt{2}\) ．同时，由于 \(\angle M A N\) 永远小于 \(45^{\circ}\) ，所以 \(\frac{|A M|}{|M N|}>\sqrt{2}\) ，但可以取到任意一个离 \(\sqrt{2}\) 非常近的数。\\
\includegraphics[max width=\textwidth, center]{2025_02_16_7c9aaf92bc99dc07f048g-142}

综合（1），（2），\(\frac{|A M|}{|M N|} \in(\sqrt{2},+\infty)\) 。无需计算，也直接得到了答案。\\
如果不使用极端遍历法，而使用代数方法，需利用双曲线的方程。

\[
\begin{gathered}
|A M|=\sqrt{(x+1)^{2}+y^{2}}=\sqrt{x^{2}+2 x+1+x^{2}-1}=\sqrt{2 x^{2}+2 x} \\
|M N|=|y|=\sqrt{x^{2}-1}
\end{gathered}
\]

所以

\[
\frac{|A M|}{|M N|}=\sqrt{\frac{2 x^{2}+2 x}{x^{2}-1}}=\sqrt{2} \cdot \sqrt{\frac{x^{2}+x}{x^{2}-1}}
\]

我们只需要求 \(\sqrt{\frac{x^{2}+x}{x^{2}-1}}\) 在 \((1,+\infty)\) 上的值域即可。注意到分子，分母都有 \((x+1)\) 为因式，所以

\[
\sqrt{\frac{x^{2}+x}{x^{2}-1}}=\sqrt{\frac{x(x+1)}{(x+1)(x-1)}}=\sqrt{\frac{x}{x-1}}=\sqrt{1+\frac{1}{x-1}}
\]

由于 \(x-1\) 在 \((1,+\infty)\) 上可以取遍所有的正实数，所以 \(\frac{1}{x-1} \in(0,+\infty)\) ，故

\[
\sqrt{1+\frac{1}{x-1}} \in(1,+\infty)
\]

故

\[
\frac{|A M|}{|M N|}=\sqrt{2} \cdot \sqrt{1+\frac{1}{x-1}} \in(\sqrt{2},+\infty)
\]

这个方法也可以得到正确答案，但需要一定计算。\\
这里我们需要特别注意的是：极端遍历法，只能够看到部分。想要看到全部的情况，必要时仍需利用代数方法。不过，极端遍历法的方法也是严格，严谨的。这是因为：\\
（1）整个运动过程连续。它既可以取到比 \(\sqrt{2}\) 大，且＂要多近，有多近＂的数，又可以取到＂要多大，有多大＂的数，所以可以取到所有 \((\sqrt{2},+\infty)\) 的数；\\
（2）\(\frac{|A M|}{|M N|}>\sqrt{2}\) 恒成立，所以取不到 \((-\infty, \sqrt{2}]\) 内的任何数。\\
所以它的取值范围就是 \((\sqrt{2},+\infty)\) 。但是对于一般的情况，极端遍历法未必能够直接得到最值。我们将在例题中看到。

\section*{5．2．1．3 比大小}
极端遍历法，除了可以看 \(x \rightarrow+\infty\) 时的趋势，还可以看 \(x \rightarrow 0\) 或者其他具体数时的趋势。我们以 2022 年新高考 1 卷的题目为例：

例 5.6 【2022 新高考 1 卷】设 \(a=0.1 \mathrm{e}^{0.1}, b=\frac{1}{9}, c=-\ln 0.9\) ，则（ ）\\
A，\(a<b<c\)\\
B，\(c<b<a\)\\
C，\(c<a<b\)\\
D，\(a<c<b\)

学习过泰勒展开的同学知道：这个题，可以用泰勒展开＂出结果＂。但是我们在学习它的时候，可能没有完全理解它的原理，对它的印象停留在它是＂高等数学＂的一部分，理解它需要额外的知识。事实上，我们完全可以利用极端遍历法来理解。

我们首先估计这三个数。他们应该都在 0 附近：\(-\ln 0.9 \approx-\ln 1=0, \frac{1}{9}=0.11 \ldots \approx 0,0.1 e^{0.1} \approx 0\) 。同时，显然 \(a\) 可以看作 \(f(x)=x e^{x}\) 在 \(x=0.1\) 处的取值，即 \(a=f(0.1)\) 。我们能否相应地构造其他函数呢？构造的要点是，利用 \(x=0.1\) ，来表示 \(b, c\) ．例如：

\[
b=\frac{1}{9}=\frac{\frac{1}{10}}{\frac{9}{10}}=\frac{x}{1-x}, c=-\ln 0.9=-\ln (1-x)
\]

记 \(g(x)=\frac{x}{1-x}, h(x)=-\ln (1-x)\) 。我们只需要比较 \(f(x), g(x), h(x)\) 在 \(x=0.1\) 处的取值即可。\\
这里的难点在：我们并不知道 \(f(x)-g(x), f(x)-h(x), h(x)-g(x)\) 的正负。如果化为函数问题求解，可能还需要证明导数的正负号，较为繁琐。此时，我们观察到：

\[
f(0)=g(0)=h(0)=0
\]

由于 \(x=0\) ．说明：\(f(x), g(x), h(x)\) 这三个函数起点相同。如果我们把 \(x=0.1\) 看作时间，\(f(x), g(x), h(x)\) 看作 \(a, b, c\) 三者的距离（位移），那么我们想比较的是：当时间经过 0.1 个单位时，\(a, b, c\) 分别处在什么位置。谁跑得远，谁大！

由于 0.1 个单位并不是一个很长的时间单位，所以，我们可以直接比较它们在起点处 \(x=0\) 的速度。在这个较短的时间里，速度变化也不会很大，那么 \(f(0.1), g(0.1), h(0.1)\) 和它们的速度几乎成正比。速度，就是它们在 \(x=0\) 处的导数：

\[
f(0.1) \approx 0.1 \cdot f^{\prime}(0), g(0.1) \approx 0.1 \cdot g^{\prime}(0), h(0.1) \approx 0.1 \cdot h^{\prime}(0)
\]

这就是极端遍历法：我们让时间间隔 \(x\) 足够小，看看一开始到底谁快谁慢。经过很短的时间积累，远近应该直接由速度大小决定。

然而，经过计算发现：

\[
\begin{gathered}
f^{\prime}(x)=(x+1) e^{x}, f^{\prime}(0)=1 \\
g^{\prime}(x)=\left(1+\frac{1}{1-x}\right)=\frac{1}{(x-1)^{2}}, g^{\prime}(0)=1 \\
h^{\prime}(x)=(-\ln (1-x))^{\prime}=\frac{1}{1-x}, h^{\prime}(0)=1
\end{gathered}
\]

这说明，它们一开始的速度完全一样\\
那么怎么办？速度一样，我们就只能比较更精细的指标：加速度。想象两个人赛跑：起点一样，速度一样，但是甲比乙的加速度大，那么甲在一开始就会领先乙。所以，加速度越大，最后位移也会越大。加速度，就是速度的导数，也就是二阶导数。

\[
\begin{aligned}
f^{\prime \prime}(x) & =(x+2) e^{x}, f^{\prime \prime}(0)=2 \\
g^{\prime \prime}(x) & =-\frac{2}{(x-1)^{3}}, g^{\prime \prime}(0)=2 \\
h^{\prime \prime}(x) & =\frac{1}{(x-1)^{2}}, h^{\prime \prime}(0)=1
\end{aligned}
\]

终于出现了分别 \(!f^{\prime \prime}(0)=g^{\prime \prime}(0)=2\) 较大，\(h^{\prime \prime}(0)=1\) 较小。所以：\(a, b\) 加速度仍一样，大于 \(c\) 。所以 \(c\) 最小。\\
那么，如何比较 \(a, b\) 呢？我们只能看：加速度，谁涨得快了。加速度涨得快，速度就会快，最终位移大。相当

于：＂加速度＂的增长，累积到速度上；＂速度＂的增长，累积到位移上。涨得快，跑得远。加速度的增长，就要看加速度的导数，也就是三阶导数。

\[
\begin{aligned}
& f^{\prime \prime \prime}(x)=(x+3) e^{x}, f^{\prime \prime \prime}(0)=3 \\
& g^{\prime \prime \prime}(x)=\frac{6}{(x-1)^{4}}, g^{\prime \prime \prime}(0)=6
\end{aligned}
\]

由于 \(6>3\) ，所以 \(b\) 加速度涨得比 \(a\) 快，所以 \(b>a\) ．因此：

\[
b>a>c
\]

选 C．\\
虽然这个方法并不是严格的证明（它可以被严格化），但是给了我们正确的答案。这里极端遍历法的逻辑，本质上就是导数的定义。我们再看这个式子：

\[
f(0.1) \approx 0.1 \cdot f^{\prime}(0)
\]

让 \(0.1=\Delta x\) ，它就是：

\[
\begin{gathered}
f(\Delta x)-f(0) \approx \Delta x \cdot f^{\prime}(0) \\
\frac{f(\Delta x)-f(0)}{\Delta x} \approx f^{\prime}(0)
\end{gathered}
\]

这就是我们导数的＂定义＂。所以，刚才的比较本质上是：用切线斜率（起点的瞬时速度），替代了割线斜率（总的平均速度）。由于时间很短，这两个相差不大。在初始速度的差距是 1 这个量级时， 0.1 的时间间隔不足以改变大小关系，所以直接用初始速度比较即可。

有了这个方法后，我们再也不需要死记硬背泰勒展开公式，不需要背诵拗口的系数。因为我们理解了它的底层原理：\\
（1）高阶导数是低阶导数的变化率，变化率积累，最终直接反映在了原函数上。\\
（2）起点（原函数）相同，比速度（一阶导数）；速度（一阶导数）相同，比加速度（二阶导数）；加速度（二阶导数）相同，比＂加速度的加速度＂（三阶导数）……直到最后分出胜负。

泰勒展开中的系数和 \(x^{n}\) ，只是为了体现导数是如何影响原函数的，类似于＂比例系数＂。不同的函数，这个比例系数是没区别的。所以比大小时，不需要可以背诵，只需要比较导数大小即可。

\section*{5．2．2 遍历＂临界点＂}
举一个物理中的例子：标准大气压下，水沸腾的温度是 100 摄氏度。所以， 100 摄氏度对于水的物理状态，就是一个临界点：高于它，会变成气态；低于它，仍是液态。像这样，在该点附近的两侧，整个系统处于不同状态的点，就是临界点。

我们也可以以简单的不等式来理解临界点。对于不等式 \(x^{2}-3 x+2<0\) ，它的临界点是 \(x_{1}=1, x_{2}=2\) ，因为它们刚好是 \(x^{2}-3 x+2=0\) 的零点。此时：\\
（1）\(x_{1}\) 附近左侧的实数，不在解集中，\(x_{1}\) 附近右侧的实数，在解集中；\\
（2）\(x_{2}\) 附近左侧的实数，在解集中，\(x_{2}\) 附近右侧的实数，不在解集中。\\
对于均值不等式 \(a+b \geq 2 \sqrt{a b}\) 而言，\(a=b\) 则是一个临界点或者临界状态。它使得不等式刚好取到等号，而其他的状态下，\(a+b>2 \sqrt{a b}\) ．

临界点之所以重要，就是因为它经常被用作刻画系统的性质。例如，一个参数的取值范围 \([a, b]\) 的两个端点 \(a, b\) ，就是临界点：\(a\) 附近左侧的实数，不满足题意，而其附近右侧的实数，满足题意；\(b\) 附近左侧的实数，满

足题意，而其附近右侧的实数，不满足题意。如果我们能够准确地把临界点找到，就可以直接找到参数的取值范围。我们以 2022 年浙江省高考真题问题为例，看看如何利用临界状态得到答案。

例 5.7 （2022 浙江）已知 \(a, b \in \mathbb{R}\) ，若对任意实数 \(x\) ，都成立 \(a|x-b|+|x-4|-|2 x-5| \geq 0\) 成立，则 \(\qquad\)。\\
A，\(a \leq 1, b \geq 3\)\\
B ，\(a \leq 1, b \leq 3\)\\
C，\(a \geq 1, b \geq 3\)\\
D，\(a \geq 1, b \leq 3\)

我们之前考虑的极端遍历法，在这个题目中仍然是适用的。\\
（1）考虑 \(x\) 足够大的情况。无论 \(b\) 为何值，都取比 \(b\) 大的 \(x\) 。同时，也要求 \(x>4,2 x>5\) ．此时，不等式化为：

\[
\begin{gathered}
a(x-b)+x-4-(2 x-5) \geq 0 \\
(a-1) x+1-a b \geq 0
\end{gathered}
\]

左侧可以看作一条直线的方程。只有当斜率非负时，它才可能在 \(x\) 足够大的时候，一直大于等于 0 ．因此，必然有 \(a-1 \geq 0\) 。即 \(a \geq 1\) ；\\
（2）\(a\) 越大，就会有越多的 \(|x-b| \geq 0\) 在大于等于号的左侧，不等式越容易成立。所以，可以考虑最＂苛刻＂的情况，即 \(a=1\) ．在这个情况下，我们看是否能得到 \(b\) 的范围；\\
（3）当 \(a=1\) 时，选项要求我们判断 \(b\) 和 3 的关系。显然，可以将不等式化为

\[
|x-b|+|x-4|-2\left|x-\frac{5}{2}\right| \geq 0
\]

我们将 \(A\left(\frac{5}{2}\right), B(3), C(4)\) 三个点标注在数轴上，令 \(O(x)\) 对应的实数为 \(x\) 。\\
\includegraphics[max width=\textwidth, center]{2025_02_16_7c9aaf92bc99dc07f048g-145}

我们考虑＂临界状态＂\(b=3\) ，看看不等式发生了什么。当 \(b=3\) 时，不等式转化为：

\[
|O C|+|O B| \geq 2|O A|
\]

考虑＂最不利＂的情况（这也是一种极端遍历法）：当 \(3 \leq x \leq 4\) 时，

\[
|O C|+|O B|=1
\]

取到了最小值。然而，此时

\[
2|O A| \geq 2 \cdot \frac{1}{2}=1
\]

不等式刚好为 0 ，达到了临界状态。当 \(x\) 再向右偏一点时，如 \(x=3.1\) ，就会产生矛盾。此时，

\[
\begin{gathered}
|O C|+|O B|=1 \\
2|O A|=2\left|x-\frac{5}{2}\right|=2 \cdot 0.6=1.2>|O B|+|O C|
\end{gathered}
\]

矛盾！究其原因，是因为当 \(x\) 比 3 大一点时，\(|O B|+|O C|=|x-3|+|x-4|=|B C|\) 的值为 1 保持不变，但是 \(\left|x-\frac{5}{2}\right|=|O A|\) 在增大，所以就从 \(x=3\) 时取等的情况，导致了矛盾。

为了挽救这个矛盾，我们得让 \(b\) 比 3 小：这样才能让 \(|b-4|\) 的距离更大，同时 \(\left|b-\frac{5}{2}\right|\) 距离变小，这样当 \(x=b\) 时，矛盾不会产生。所以 \(b \leq 3\) 。正确答案为 D．

有同学可能对最后一步判断产生疑问：想让 \(|b-4|\) 变大，也可以让 \(b\) 大于 4 取更大的实数。但此时，\(\left|b-\frac{5}{2}\right|\)

也会随着变大，并没有阻止矛盾的产生：若 \(b>4\) ，当 \(x=b\) 时，

\[
|x-b|+|x-4|-2\left|x-\frac{5}{2}\right|=b-4-2\left(b-\frac{5}{2}\right)=1-b<0
\]

如果利用我们第一节课学过的举例子方法，能否也得到正确答案呢？也可以。根据绝对值不等式，当 \(b=1\)时，刚好有

\[
x-b+x-4=2 x-5
\]

所以，根据绝对值不等式，

\[
|x-1|+|x-4| \geq|x-1+x-4|=|2 x-5|
\]

\(b=1\) 符合题目条件。事实上，当 \(a=1\) 时，\(b\) 只可能为 1 ．这是因为：\\
（1）当 \(x=b\) 时，我们要求：

\[
|b-4| \geq|2 b-5|
\]

两边平方得

\[
\begin{gathered}
b^{2}-8 b+16 \geq 4 b^{2}-20 b+25 \\
3 b^{2}-12 b+9 \leq 0 \\
3(b-1)(b-3) \leq 0
\end{gathered}
\]

所以 \(1 \leq b \leq 3\) ．\\
（2）当 \(x=4\) 时，我们要求：

\[
|b-4| \geq 3
\]

所以 \(b \geq 7\) 或 \(b \leq 1\) ．\\
由于（1），（2）都得成立，所以 \(b=1\) ．为什么代入这两个值？因为它们是让 \(|x-b|+|x-4|\) 取最小值 \(|b-4|\) 的极端情况，它们往往也是取等的情况。

我们再来看一个临界点的例子。

例 5.8 （2021 上海）已知实数 \(x_{1}, y_{1}, x_{2}, y_{2}, x_{3}, y_{3}\) 同时满足：\\
（1）\(x_{1}<y_{1}, x_{2}<y_{2}, x_{3}<y_{3}\)\\
（2）\(x_{1}+y_{1}=x_{2}+y_{2}=x_{3}+y_{3}\)\\
（3）\(x_{1} y_{1}+x_{3} y_{3}=2 x_{2} y_{2}>0\) ．\\
则下列选项中，一定成立的有：\\
A， \(2 x_{2}<x_{1}+x_{3}\)\\
B， \(2 x_{2}>x_{1}+x_{3}\)\\
C，\(x_{2}^{2}<x_{1} x_{3}\)\\
D，\(x_{2}^{2}>x_{1} x_{3}\)

【方法一】面积\\
我们如何处理题目中的条件呢？显然，可以将 \(x, y\) 都看作点的坐标。令 \(A\left(x_{1}, y_{1}\right), B\left(x_{2}, y_{2}\right), C\left(x_{3}, y_{3}\right)\) ．\\
（2），（3）两个条件如何翻译？\\
（2）：直线 \(A, B, C\) 三点共线。它们都在直线 \(x+y=k\) 上，其中 \(k=x_{1}+y_{1}=x_{2}+y_{2}=x_{3}+y_{3}\) ．不妨假设 \(k>0\) ．\\
（3）：\(x y\) 的形式，让我们联想到面积。过 \(A, B, C\) 分别作坐标轴的垂线，\(x\) 轴上的垂足为 \(A_{1}, B_{1}, C_{1}, y\) 轴上的垂足

为 \(A_{2}, B_{2}, C_{2}\) 。我们假设 \(A, B, C\) 都在第一象限，那么（3）等价于：

\[
S_{O A_{1} A A_{2}}+S_{O C_{1} C C_{2}}=2 S_{O B_{1} B B_{2}}
\]

这三个矩形有这样的面积关系。我们考虑 \(A B\) 选项：需要考虑 \(x_{2}\) 和 \(\frac{x_{1}+x_{3}}{2}\) 的大小，也就是 \(B\) 和 \(A C\) 中点的位置关系。我们取出 \(A C\) 的中点 \(D\) ，和另外三个点一样作出现，垂足为 \(D_{1}, D_{2}\) 。我们来比较面积大小。

如图所示，这几组垂线，将这些矩形分成了六个部分。其中

\[
\begin{gathered}
S_{O A_{1} A A_{2}}=x_{1} y_{1}=S_{1}+S_{2}+S_{4} \\
S_{O D_{1} D D_{2}}=x_{2} y_{2}=S_{2}+S_{3}+S_{4}+S_{5} \\
S_{O C_{1} C C_{2}}=x_{3} y_{3}=S_{4}+S_{5}+S_{6} .
\end{gathered}
\]

\begin{center}
\includegraphics[max width=\textwidth]{2025_02_16_7c9aaf92bc99dc07f048g-147}
\end{center}

那么

\[
S_{O A_{1} A A_{2}}+S_{O C_{1} C C_{2}}=S_{1}+S_{2}+2 S_{4}+S_{5}+S_{6}
\]

由于 \(D\) 是 \(A C\) 中点，所以 \(S_{1}=S_{2}, S_{5}=S_{6}\) ，所以

\[
S_{O A_{1} A A_{2}}+S_{O C_{1} C C_{2}}=2 S_{2}+2 S_{4}+2 S_{5}=2\left(S_{2}+S_{4}+S_{5}\right)
\]

而

\[
2 S_{O D_{1} D D_{2}}=2\left(S_{2}+S_{3}+S_{4}+S_{5}\right)>2\left(S_{2}+S_{4}+S_{5}\right)=S_{O A_{1} A A_{2}}+S_{O C_{1} C C_{2}}
\]

所以

\[
2 S_{O D_{1} D D_{2}}>S_{O A_{1} A A_{2}}+S_{O C_{1} C C_{2}}=2 S_{O B_{1} B B_{2}}
\]

也就是 \(S_{O D_{1} D D_{2}}>S_{O B_{1} B B_{2}}\) ．又因为 \(x_{2}<y_{2}\) ，为了让面积更小，\(x_{B}\) 应小于 \(x_{D}\) ，即 \(x_{2}<\frac{x_{1}+x_{3}}{2}\) ．故选 A．我们这里只考虑了 \(A, B, C\) 都在第一象限，以及 \(k>0\) 的情况。类似的，可以讨论其他的情况，在例题和作业中我们会处理。

\section*{【方法二】均值不等式}
我们还可以怎么翻译条件呢？条件（2）意味着，每组 \(\left(x_{i}, y_{1}\right)\) 的平均数是相等的。我们记为 \(m\) ，即

\[
m=\frac{x_{1}+y_{1}}{2}=\frac{x_{2}+y_{2}}{2}=\frac{x_{3}+y_{3}}{2}
\]

根据条件（1），我们知道：\(x_{1}, x_{2}, x_{3}\) 在 \(m\) 的左侧，\(y_{1}, y_{2}, y_{3}\) 在 \(m\) 的右侧。同时，\(x_{i}, y_{i}\) 到平均数 \(m\) 的距离相等。

分别记作 \(a, b, c\) ：

\[
\begin{aligned}
a & :=m-x_{1}=y_{1}-m>0 \\
b & :=m-x_{2}=y_{2}-m>0 \\
c & :=m-x_{3}=y_{3}-m>0
\end{aligned}
\]

\begin{center}
\includegraphics[max width=\textwidth]{2025_02_16_7c9aaf92bc99dc07f048g-148}
\end{center}

那么，条件（3）可以化简为非常简单的式子：

\[
\begin{aligned}
& x_{1} y_{1}=(m-a)(m+a)=m^{2}-a^{2} \\
& x_{2} y_{2}=(m-b)(m+b)=m^{2}-b^{2} \\
& x_{3} y_{3}=(m-c)(m+c)=m^{2}-c^{2}
\end{aligned}
\]

因此，条件（3）等价于： \(2 b^{2}=a^{2}+c^{2}\) 。看选项：\\
AB：比较 \(2 x_{2}\) 和 \(x_{1}+x_{3}\) 的关系，即

\[
\begin{gathered}
2 x_{2}=2 m-2 b \\
x_{1}+x_{3}=m-a+m-c=2 m-(a+c)
\end{gathered}
\]

即比较 \(2 b\) 和 \(a+c\) 的关系。根据均值不等式：

\[
(a+c)^{2}=a^{2}+c^{2}+2 a c \leq 2 a^{2}+2 c^{2}=4 b^{2}=(2 b)^{2}
\]

又因为 \(a, b, c\) 均为正数，所以 \(a+c<2 b\) ．所以

\[
2 x_{2}=2 m-2 b<2 m-a-c=x_{1}+x_{3} .
\]

故 A 正确，B 错误；\\
C，D：比较 \(x_{2}^{2}\) 和 \(x_{1} x_{3}\) 的关系，这个表达式相对复杂：

\[
\begin{gathered}
x_{2}^{2}=m^{2}-2 b m+b^{2} \\
x_{1} x_{3}=(m-a)(m-c)=m^{2}-(a+c) m+a c
\end{gathered}
\]

因此：

\[
\begin{aligned}
x_{2}^{2}-x_{1} x_{3} & =(a+c-2 b) m+b^{2}-a c \\
& =(a+c-2 b) m+\frac{a^{2}+c^{2}}{2}-a c \\
& =(a+c-2 b) m+\frac{(a-c)^{2}}{2}
\end{aligned}
\]

虽然 \(\frac{(a-c)^{2}}{2}>0\) ，但是 \(a+c<2 b\) ，所以 \(x_{2}^{2}-x_{1} x_{3}\) 的正负不确定。我们可以考虑极端情况：\\
（1）当 \(m \rightarrow+\infty\) 时，\((a+c-2 b) m \rightarrow-\infty\) ，此时 \(x_{2}^{2}<x_{1} x_{3}\) ；\\
（2）当 \(m \rightarrow-\infty\) 时，\((a+c-2 b) m \rightarrow+\infty\) ，此时 \(x_{2}^{2}>x_{1} x_{3}\) ．\\
故 C，D 均不恒成立。

\section*{如何理解和应用？}
极端遍历法和临界状态，可以理解为整个变化情况中的特殊情况。极端遍历法，考虑的是＂最＂的情况：最大，最小，最多，最少，最利，最不利，等等；临界状态，考虑的是＂相变＂的点：有解 vs 无解， 1 个解 vs 2 个解，中点等等。它们之所以重要，就是因为这些情况下得到的参数范围，经常会是参数最终范围的一部分。除此之外，利用导数的性质去考察函数的变化，也是一种极端遍历法：我们利用 \(f(\Delta x) \approx f(0)+f^{\prime}(0) \cdot \Delta x\)（对很小的 \(\Delta x)\) 来估计 \(f(\Delta x)\) ．这样一来，问题就可以化归为 \(f(x)\) 在 \(x=0\) 处的导数。

极端和临界状态，也是我们在做思想试验时经常会单独拿出来考虑的状态。它有助于我们体会哪个条件，变量，会导致系统产生质变。

\section*{5．2．3 XXXX}
\section*{5.3 再论自由度}
\section*{5．3．1 自由度回顾}
我们前面章节中，已经讨论过自由度。我们简单回顾一下关于它的基本定义。想要确定一个系统，我们需要确定系统的一些参数。例如，确定二次函数 \(y=a x^{2}+b x+c\) ，我们需要确定 \(a, b, c\) 三个参数。因此，这个系统 （二次函数）有三个自由度。在第二章中，我们只讨论了多元一次方程（线性方程组）的自由度。这一节中，我们将用它解决其他部分的内容。

例 5.9 （2023 乙卷）已知圆的半径为 1 ，直线 \(P A\) 与圆 \(O\) 相切于 \(A\) 点，直线 \(P B\) 与圆 \(O\) 交于 \(B, C\) 两点，\(D\) 为 \(B, C\) 中点。若 \(P O=\sqrt{2}\) ，则 \(\overrightarrow{P A} \cdot \overrightarrow{P D}\) 的最大值为 \(\qquad\)。\\
A，\(\frac{1+\sqrt{2}}{2}\)\\
B，\(\frac{1+2 \sqrt{2}}{2}\)\\
C． \(1+\sqrt{2}\)\\
D， \(2+\sqrt{2}\)

这是一个直线和圆且结合向量的问题。显然，我们可以建系处理。以 \(O\) 为圆心建立平面直角坐标系，那么圆 \(O\) 的方程为 \(x^{2}+y^{2}=1\) 。由于整个图形旋转不变（旋转整个图形，不改变最大值），所以不妨设 \(A(0,1)\) 。由于 \(P A\) 和圆 \(O\) 相切，所以\\
\(P A \perp O A\) ．故 \(P\) 点纵坐标也是 1．又因为 \(|P O|=\sqrt{2}\) ，所以 \(P(1,1)\) 或 \((-1,1)\) ．不妨设 \(P(1,1)\)（另一侧的情况关于 \(y\) 轴对称）。

我们已经有了 \(P\) 点坐标，所以 \(\overrightarrow{P A}=(-1,0)\) ．我们只需要表示 \(D\) 点坐标，就可以求取值范围。终于，我们等到自由度登场了：这个系统，有几个自由度呢？显然，如果确定了直线 \(B C\) ，那么整个系统就确定了。而直线 \(B C\)经过点 \(P\) ，所以只能绕着 \(P\) 点旋转。因此，它只有 1 个自由度。那么，我们设哪个变量作为自由度呢？\\
（1）斜率。这是最直接的方式，那么直线 \(B C\) 的方程为

\[
\begin{gathered}
y-1=k(x-1) \\
y=k x+1-k
\end{gathered}
\]

为了求 \(D\) 点坐标，只需联立即可：将 \(y=k x+1-k\) 代人 \(x^{2}+y^{2}=1\) ，得

\[
\begin{gathered}
x^{2}+(k x+1-k)^{2}=1 \\
\left(k^{2}+1\right) x+2 k(1-k) x+k^{2}-2 k=0
\end{gathered}
\]

所以

\[
x_{B}+x_{C}=\frac{2 k(k-1)}{k^{2}+1}=\frac{2 k^{2}-2 k}{k^{2}+1}
\]

由于 \(D\) 为 \(B, C\) 中点，所以

\[
x_{D}=\frac{x_{B}+x_{C}}{2}=\frac{k^{2}-k}{k^{2}+1}
\]

又因为 \(D\) 在直线 \(y=k(x-1)+1\) 上，所以

\[
y_{D}=k\left(x_{D}-1\right)+1=k \cdot \frac{-k-1}{k^{2}+1}+1=\frac{1-k}{k^{2}+1}
\]

故 \(D\left(\frac{k^{2}-k}{k^{2}+1}, \frac{1-k}{k^{2}+1}\right)\) ．所以

\[
\overrightarrow{P D}=\left(\frac{-k-1}{k^{2}+1}, \frac{-k-k^{2}}{k^{2}+1}\right)
\]

因此

\[
\overrightarrow{P A} \cdot \overrightarrow{P D}=\frac{k+1}{k^{2}+1}
\]

只需求 \(f(k)=\frac{k+1}{k^{2}+1}\) 的最大值即可。注意，这里斜率 \(k\) 有取值范围：显然，过 \(P\) 的两条切线分别为 \(x=1, y=1\) ．所以，斜率 \(k>0\) 即可保证有两个交点 \(B, C\) 。因此，需要求 \(f(k)\) 在 \((0,+\infty)\) 的最大值即可。

可以利用均值不等式：令 \(k+1=t\) ，那么

\[
f(k)=\frac{t}{(t-1)^{2}+1}=\frac{t}{t^{2}-2 t+2}=\frac{1}{t+\frac{2}{t}-2}
\]

由于 \(t \in(1,+\infty)\) ，根据均值不等式：\(t+\frac{2}{t} \geq 2 \sqrt{2}\) ，当且仅当 \(t=\sqrt{2}\) 时取等。因此，

\[
f(k) \leq \frac{1}{2-\sqrt{2}}=\frac{\sqrt{2}+1}{2} .
\]

故 \(\overrightarrow{P A} \cdot \overrightarrow{P D}\) 的最大值为 \(\frac{\sqrt{2}+1}{2}\) ．\\
这个方法虽然奏效，但是计算繁琐。根据我们的经验，往往直线和圆的题目，是不需要联立直线的。虽然我们知道，联立是保底的方法和通法，肯定能够得到答案，但能否更好地利用圆的几何性质和系统的自由度呢？

事实上，在圆中，设角度往往是可以考虑的方法。这是因为：题目中有切线，中点，会有大量的垂直（中点也会带来垂直，因为圆中有＂垂径定理＂）。一旦有垂直，就可以＂导角＂，就可以用角度来确定系统。设哪个角度？最自然的方法，就是设和直线 \(B C\) 的倾斜角相关的角度。显然，它可以作为系统的自由度：倾斜角确定，直线，斜率也就确定了。\\
（2）设 \(\angle A P D=\theta\) ．那么 \(\theta \in\left(0, \frac{\pi}{2}\right)\) ．那么

\[
\overrightarrow{P A} \cdot \overrightarrow{P D}=|\overrightarrow{P A}| \cdot|\overrightarrow{P D}| \cdot \cos \theta
\]

由于 \(D\) 为 \(B C\) 中点，根据垂径定理，有 \(\angle O P D=\frac{\pi}{2}\) ．所以

\[
|\overrightarrow{P D}|=|\overrightarrow{P O}| \cdot \cos \angle O P D=\sqrt{2} \cos \left(\theta-\frac{\pi}{4}\right)
\]

所以

\[
\overrightarrow{P A} \cdot \overrightarrow{P D}=\sqrt{2} \cos \theta \cos \left(\theta-\frac{\pi}{4}\right)
\]

我们只要处理这个三角函数即可。有两种方法：\\
（i）直接展开：

\[
\sqrt{2} \cos \theta \cos \left(\theta-\frac{\pi}{4}\right)=\sqrt{2} \cos \theta\left(\cos \theta \cdot \frac{\sqrt{2}}{2}+\sin \theta \cdot \frac{\sqrt{2}}{2}\right)=\cos ^{2} \theta+\cos \theta \sin \theta .
\]

根据二倍角公式：

\[
\cos ^{2} \theta+\cos \theta \sin \theta=\frac{\cos 2 \theta+1}{2}+\frac{\sin 2 \theta}{2}=\frac{\sqrt{2}}{2} \sin \left(2 \theta+\frac{\pi}{4}\right)+\frac{1}{2} .
\]

当且仅当 \(\sin \left(2 \theta+\frac{\pi}{4}\right)=1\) 时取最大值，此时 \(\theta=\frac{\pi}{8}, \overrightarrow{P A} \cdot \overrightarrow{P D}=\frac{\sqrt{2}+1}{2}\) ．\\
（ii）注意到

\[
\cos \left(\theta-\frac{\pi}{4}\right)=\cos \left(\frac{\pi}{4}-\theta\right),
\]

所以两个余弦内的角度和为定值：

\[
\theta+\frac{\pi}{4}-\theta=\frac{\pi}{4}
\]

所以，可以考虑使用积化和差：

\[
\cos \theta \cos \left(\theta-\frac{\pi}{4}\right)=\cos \theta \cos \left(\frac{\pi}{4}-\theta\right)=\frac{1}{2}\left(\cos \frac{\pi}{4}+\cos \left(\frac{\pi}{4}-2 \theta\right)\right)
\]

故

\[
\overrightarrow{P A} \cdot \overrightarrow{P D}=\sqrt{2} \cos \theta \cos \left(\theta-\frac{\pi}{4}\right)=\frac{1}{2}+\frac{\sqrt{2}}{2} \cos \left(\frac{\pi}{4}-2 \theta\right) \leq \frac{\sqrt{2}+1}{2}
\]

当且仅当 \(\frac{\pi}{4}-2 \theta=0\) 时取等，即 \(\theta=\frac{\pi}{8}\) 。\\
我们还有没有其他设自由度的方法呢？事实上，（1），（2）两种方法都聚焦角度，它们帮助我们确定了直线：它是过 \(P(1,1)\) 且倾斜角为 \(\theta\)（或斜率为 \(k\) ）的直线。除了＂定点＋角度＂，还可以通过＂两点＂来确定直线。所以我们也可以设点的坐标。图中有三个动点：\(B, C, D\) ，前两者在圆上动，可以用参数方程等方法来表示。

图中有三个动点：\(B, C, D\) ，前两者在圆上动，可以用参数方程等方法来表示。但这个设法并不直接：因为我们要求的是 \(\overrightarrow{P A} \cdot \overrightarrow{P D}\) 的范围，而不是 \(\overrightarrow{P A} \cdot \overrightarrow{P B}\) 或 \(\overrightarrow{P A}\) ．\\
\(\overrightarrow{P C}\) 的范围。所以，如果能够写出 \(D\) 点坐标，可以一步到位。\\
如何设 \(D\) 点坐标呢？如果我们直接设 \(D\left(x_{0}, y_{0}\right)\) ，会和我们系统的自由度产生分歧：\(x_{0}, y_{0}\) 如果都是自由变量，那系统就会有两个自由度。而根据前述分析，我们知道其实只有一个，所以 \(x_{0}, y_{0}\) 一定有联系。换言之，\(D\)点的轨迹是我们可以确定的。通过（2）中的推理，我们发现：根据垂径定理，\(\angle P D O=\frac{\pi}{2}\) 。所以，\(D\) 在以 \(O P\)为直径的圆（弧）\(O_{1}\) 上（夹在 \(y=1, x=1\) 间的部分）！它的半径为 \(r_{O_{1}}=\frac{\sqrt{2}}{2}\) ．这就把我们的问题重新化简了：点 \(D\) 在圆 \(O_{1}\) 夹在 \(y=1, x=1\) 之间的圆弧上运动，求 \(\overrightarrow{P A} \cdot \overrightarrow{P D}\) 的最大值：

为了求最大值，不再需要严格做代数推理，只要利用数量积的几何意义。过 \(D\) 作 \(P A\) 垂线，垂足为 \(D_{1}\) 。那么

\[
\overrightarrow{P A} \cdot \overrightarrow{P D}=\left|P D_{1}\right| \cdot|P A|=\left|P D_{1}\right|
\]

故只需要求 \(\left|P D_{1}\right|\) 的最大值即可。而

\[
\left|P D_{1}\right|=1-x_{D}
\]

所以只要看 \(D\) 点横坐标的最小值即可。显然，在 \(O_{1} D\) 平行于 \(P A\) ，且 \(D\) 在 \(O_{1}\) 左侧时取最大值。此时

\[
\left|P D_{1}\right|=\frac{1}{2}|P A|+r_{O_{1}}=\frac{1+\sqrt{2}}{2}
\]

故最大值为 \(\frac{1+\sqrt{2}}{2}\) ．

我们总结此题思路，便发现：所有不同的未知变量的设法，都源自于＂直线＂的自由度。\\
（1）直线上一定点 \(P(1,1)\) ，和直线的＂方向＂，可以确定直线：\(k\) 是直线斜率，\(\angle A P D=\theta\) 是倾斜角；\\
（2）两点确定一条直线：\(P B, P C, P D\) 均可确定。但 \(B, C\) 只是单位圆上的点，和 \(D\) 并无直接关系。所以我们想找到 \(D\) 点轨迹，并以此确定直线即可。

\section*{5．3．2 用自由度重学高中数学}
直线的自由度我们搞清楚了，那么高中数学里其他对象的自由度呢？每个自由度又是如何影响系统的呢？我们从函数开始，一个个来体会其中的作用。

\section*{5．3．2．1 三角函数}
对于三角函数 \(y=A \sin (\omega x+\varphi)(A>0, \omega>0,|\varphi| \leq \pi)\) ，这里共有 \(A, \omega, \varphi\) 三个未知变量，所以它的自由度为 3．每个变量是如何决定系统的呢？\\
（1）\(A\) ：振幅。它决定了函数的最大值和最小值：\(y\) 的最大值为 \(A\) ，最小值为 \(-A\) 。也是振动的幅度。\\
（2）\(\omega\) ：角速度。顾名思义，就是角度的变化率。每当 \(x\) 变化一个单位时，正弦内的角度就变化 \(\omega\) 个单位。当 \(\omega\)越大，振动越快，图像变化剧烈；当 \(\omega\) 越小，振动越慢，图像变化缓慢。我们可以想象，当 \(\omega\) 从 0 变化到 \(+\infty\)时，整个函数图像就像一个弹簧一样：从完全平坦（一条平行于 \(x\) 轴的直线，\(y=A \sin \varphi\) ），慢慢压缩，到振动极快的状态。\\
（3）\(\varphi\) ：初相。当 \(x=0\) 时的角度值。它起到的作用就是水平方向平移图像。

\section*{5．3．2．2 平面向量，空间向量}
我们学过三角函数后，还学了平面向量。我们知道，向量是有＂大小＂和＂方向＂的量。在平面中，方向的变化只有一个维度，所以平面向量有 2 个自由度。除了这么确定一个向量，还可以通过坐标表示确定：\(\vec{a}=(x, y)\) ，只要把 \(x, y\) 两个坐标确定即可，所以也是 2 个自由度。事实上，一个系统的自由度个数是内蕴的：它不受表示方法影响。采用不同的表示方法，会得到同样个数的自由度。

类似地，空间向量则有 3 个自由度。一则，\(\vec{a}=(x, y, z)\) ，我们需要三个坐标才可以表示；二则，它的＂大小＂（1 个自由度）和＂方向＂（ 2 个自由度）共 3 个自由度。怎么理解呢？将 \(\vec{a}\) 的起点放到平面直角坐标系的原点处，终点为 \(A\) 。设 \(A\) 在 \(x O y\) 平面上的投影为 \(A^{\prime}\) 。那么， \(\overrightarrow{O A}\) 由以下两个角度完全确定：\\
（1） \(\overrightarrow{O A}\) 和 \(x\) 轴正方向的夹角 \(\theta\) ．它决定了 \(\vec{a}\) 在 \(x O y\) 平面上投影的方向；\\
（2）\(\varphi=\angle A O A^{\prime}\) 。它决定了 \(\vec{a}\) 在坚直方向上的仰角。\\
换言之，确定了投影的方向以及仰角，整个方向就确定了。所以方向有两个自由度，大小有一个自由度，整个向量有 3 个自由度。用向量长度 \(r\) ，以及角度 \(\theta, \varphi\) 来表示一个点的坐标系，叫做球坐标系。一般地，\(n\) 维空间中的向量则有 \(n\) 个自由度。

\section*{5．3．2．3 中心在原点，对称轴为 \(x\) 轴，\(y\) 轴的圆雉曲线}
我们在第一章动静视角中也举过类似的例子：\\
对于椭圆

\[
\frac{x^{2}}{a^{2}}+\frac{y^{2}}{b^{2}}=1
\]

和双曲线

\[
\frac{x^{2}}{a^{2}}-\frac{y^{2}}{b^{2}}=1, \frac{y^{2}}{a^{2}}-\frac{x^{2}}{b^{2}}=1,
\]

都只需要确定两个参数 \(a, b\) 即可确定方程。从几何图形上看：\\
（1）椭圆：\(a\) 代表短半轴长，\(b\) 代表长半轴长；\\
（2）双曲线：\(a\) 代表实半轴长，\(b\) 代表虚半轴长。\\
从定义上看：\\
（1）椭圆：到两定点 \(F_{1}, F_{2}\) 距离之和为定值的点所形成的轨迹。即

\[
\left|P F_{1}\right|+\left|P F_{2}\right|=2 a
\]

我们只要确定两定点间的距离，以及＂距离之和＂ \(2 a\) ，即可确定。两定点称作焦点，它们之间的距离称作焦距 （＝ \(2 c)\) ；＂距离之和＂恰好就是长轴长；\\
（2）双曲线：到两定点距离之差绝对值为定值的点所形成的轨迹。即

\[
\left\|P F_{1}|-| P F_{2}\right\|=2 a
\]

我们只要确定两定点间的距离，以及＂距离之差的绝对值＂ \(2 a\) ，即可确定。两定点称作焦点，它们之间的距离称作焦距 \((=2 c) ; "\) 距离之差＂恰好就是长轴长。还可以怎么确定一个圆雉曲线呢？\\
（1）确定离心率。当离心率确定了，整个图形的形状就确定了。这也是很多求离心率题目的关键：本质上就是要确定＂形状＂，单位长度不重要。例如，当椭圆离心率确定了，那么长短轴的比例 \(\frac{a}{b}\) 就确定了，所以形状确定；当双曲线离心率确定了，实半轴，虚半轴的比例也就确定了，所以形状确定。等价地，如果能确定两条渐近线的夹角或斜率 \(\frac{b}{a}\) ，也可以确定形状。\\
（2）再确定＂单位长度＂。例如，确定长半轴，短半轴中的一个，辅以离心率，自然确定了整个椭圆。这个道理就是＂大题计算的代数杜杆＂里讲的＂齐次＂：确定 \(e\) ，等价于确定比例 \(\frac{a}{b}\) ，再得到一个非齐次条件（如 \(a=4\) ），就可以确定系统了。

\section*{5．3．2．4 数列}
对于等差数列和等比数列而言，都只有两个自由度。从参数的角度理解，等差数列的通项公式为

\[
a_{n}=a_{1}+(n-1) d
\]

其中 \(a_{1}\) 为首项，\(d\) 为公差。只要确定了 \(a_{1}, d\) ，整个数列就确定了。从含义上讲，确定首项就是确定从哪里开始种树。公差 \(d\) ，就是告诉我们每两棵树的间隔是多少。由于等差，相邻两棵树之间的距离是保持不变的，所以整排数就确定了；

对于等比数列，同样的道理：\(a_{n}=a_{1} \cdot q^{n-1}\) 。它和等差数列道理完全相同。我们可以通过取对数理解：当 \(a_{1}, q>0\) 时，

\[
\ln a_{n}=\ln a_{1}+(n-1) \cdot \ln q .
\]

所以 \(\left\{\ln a_{n}\right\}\) 是以 \(\ln a_{1}\) 为首项， \(\ln q\) 为公差的等差数列。\\
在后面的例题中，我们还将分析一些初中学过的平面几何图形的自由度。

\section*{总结：劣定，适定和超定}
我们通过自由度确定系统，往往都是在＂限制条件个数＂等于＂自由度个数＂的情况下确定的。例如，对于多元一次方程组，一般情况下，方程的个数等于未知数的个数，那么我们就可以得到唯一的解。这是＂适定＂的情况。

类似地，如果限制条件个数少于自由度个数，那就是＂劣定＂。这意味着，系统中还有可以完全自由变化的量。例如，圆雉曲线中的定值，定点问题，往往就是劣定的：对于任意的斜率 \(k\) ，直线 \(l\) 都经过某一定点，斜率 \(k\) 是完全自由变的；当 \(P\) 在曲线上动的时候， \(\overrightarrow{P A} \cdot \overrightarrow{P B}\) 都是定值，\(P\) 也是完全自由变化的。\\
限制条件个数多于自由度个数，那就是＂超定＂。这时，可能会出现矛盾的情况。比如：方程组

\[
\left\{\begin{array}{l}
x+y=1 \\
2 x-y=2 \\
x-y=4
\end{array}\right.
\]

前两个方程已经可以完全确定 \(x, y\) 了：\(x=1, y=0\) ．但又因为有第三个方程，导致方程组无解，即前两个方程和第三个方程矛盾。这是超定的情况。

只有在少数情况下，劣定，超定的系统是可以被确定的。例如，对于实数 \(x, y\) ，若满足 \(x^{2}+y^{2}=0\) ，则必然有 \(x=y=0\) 。劣定的系统可以被确定，往往都是因为取到了边界情况。超定系统可以被确定，则是因为多出的条件和其他条件是相容的（不矛盾的）：如，方程组

\[
\left\{\begin{array}{l}
x+y=1 \\
2 x-y=2 \\
x-y=1
\end{array}\right.
\]

也是适定的，\(x=1, y=0\) 是解。\\
我们前面说所有系统可以被＂确定＂，不代表它是＂唯一确定＂的。＂确定＂的意思指的是，本来有无数多或＂不可数＂多种情况，但最终被限制在有限种，这就是确定。例如，确定了三角形的两边和一角（不是夹角），可能会有两个三角形满足条件。虽然系统不是唯一确定的，但是它也确定了：从原来的＂无穷多种＂可能，坍缩成现在的两种。这也是确定。

当我们学习过自由度后，去判断劣定，适定和超定，往往是非常重要的一步。\\
（1）对于适定的系统：如果要求某自由度 \(x\) ，我们只要找到一个关于它的方程即可。不需要直接求 \(x\) ，只要是一个和它有关的方程，就能够把它确定；\\
（2）对于劣定的系统：往往都是有＂恰好＂能求的理由。对于劣定系统，如果还能求出答案，一定是定值，定点类的。这时，可以先代人特殊的数，把答案求出来，再去验证或严格证明。

自由度如何帮助我们理解一个新的问题呢？除了它能够帮助我们判断系统是否确定，还给我们提供了研究问题的一个重要视角：要把每个影响系统的因素＂分开＂看，解构系统（之后用第一章讲到的要素关系网，重构系统）。一个系统之所以复杂，是因为我们无法完全掌控它的所有要素。它的运动情况过于复杂，导致我们看不清到底它是怎么被确定的。多因素耦合的时候，想不清楚，就要固定其他因素，每次只看一个因素。这就像做实验时的＂控制变量法＂：其他因素完全一样，看其中一个因素是如何影响系统的。所有因素都这么考察后，整个系统也就考察清楚了。

5．3．3 XXXX

\section*{5.4 动态图像}
\section*{5．4．1 什么是动态图像？}
我们在前三节中，分别学习了：\\
（1）如何做思想试验和举例子（第一节）；\\
（2）什么叫高明，有用的例子（极端情况，临界状态，第二节）；\\
（3）一个系统，需要多少个自由度才可以确定（自由度，第三节）。\\
在讨论自由度的时候，关键是知道整个系统有哪些量是可以自由变化的。我们想要体会清楚每个量的作用，可以控制其他的变量，只改变这个单一变量，去看它是如何影响系统的。这样，我们就可以知道每个自由变量对系统的影响，从而把握住整个系统。

在使用自由度去分析问题的时候，我们注意到：往往不能只考虑临界状态和极端情况，还要把状态之间是怎么转化的想清楚。例如，随着 \(A, \omega, \varphi\) 的改变，\(f(x)=A \sin (\omega x+\varphi)\) 是连续变化的，而不是一个个孤立的情况。这意味着：如果我们只是举例子，考虑极端情况，是不可能充分挖掘出整个系统的所有信息的。因此，我们要把这个连续变化也刻画清楚，这就是＂动态图像＂。

动态图像，简单说，就是要想象这个变量是怎么连续变化的。我们以 2022 年天津卷的填空压轴题为例，来体会动态图像思想应用。

例5．10（2022 天津）设 \(a \in \mathbb{R}\) 。对任意实数 \(x\) ，用 \(f(x)\) 表示 \(|x|-2, x^{2}-a x+3 a-5\) 中的较小者。若函数 \(f(x)\)至少有三个零点，则实数 \(a\) 的取值范围是 \(\qquad\)。

原来，九省联考中出现的多个函数中的较小者，较大者，并非新题，在以前的高考题中就已经有所涉及（事实上，人教版教材中，它作为例题出现）。如何分析这个题目呢？\\
（1）做思想试验。我们可以代人特别的 \(a\) ，看看 \(f(x)\) 的表达式或者图像。我们以 \(a=1\) 为例，此时 \(f(x)\) 表示 \(|x|-2\) 和 \(x^{2}-x-2\) 中的较小者。我们把两个函数图像 \(y=|x|-2, y=x^{2}-x-2\) 画出，得到：\\
\includegraphics[max width=\textwidth, center]{2025_02_16_7c9aaf92bc99dc07f048g-155(1)}

那么 \(f(x)\) 的图像是哪段呢？从图像上看，它应该是每段函数图像中，在下方的拼接而成得到的，即：\\
\includegraphics[max width=\textwidth, center]{2025_02_16_7c9aaf92bc99dc07f048g-155}

当 \(x \leq 0\) 时，\(|x|-2\) 较小；当 \(0<x \leq 2\) 时，\(x^{2}-x-2\) 较小；当 \(x>2\) 时，\(|x|-2\) 较小。所以 \(f(x)\) 的解析式为

\[
f(x)=\left\{\begin{array}{c}
|x|-2, x \leq 0 \\
x^{2}-x-2,0<x \leq 2 \\
|x|-2, x>2
\end{array}\right.
\]

显然，\(f(x)\) 不符合题意。虽然如此，但是我们对问题本身有了一个基本的把握：\\
（1）只要把二次函数 \(x^{2}-a x+3 a-5\) 的图像和 \(|x|-2\) 的图像画出来，就可以找到 \(f(x)\) 的图像，进而得到零点，不需要生硬地进行代数讨论；\\
（2）\(f(x)\) 的零点，来自 \(f_{1}(x)=|x|-2\) 和 \(f_{2}(x)=x^{2}-a x+3 a-5\) 的零点。之所以会只有三个零点，就是因为零点有损耗：当 \(f_{1}(x)=0\) 时，如果 \(f_{2}(x)<0\) ，就会有 \(f(x)=f_{2}(x) \neq 0\) 。从而损失了一个零点。这样，\(f_{1}(x), f_{2}(x)\) 本身最多有四个零点，经过损耗才会得到题目要求的三个零点。\\
（2）临界状态：我们不太能够马上看出来哪个是临界状态，因为缺乏对问题的把握，所以后面再看；\\
（3）自由度：这是一个自由度为 1 的系统。只要 \(a\) 确定，那么 \(f(x)\) 就确定了。 \(a\) 通过影响二次函数来影响 \(f(x)\) ，所以看 \(f_{2}(x)\) 即可。显然，\(a\) 和 \(f_{2}(x)\) 的一次项系数和常数项系数有关，所以它影响对称轴和方位。当 \(a\) 变化时， \(f_{2}(x)\) 是如何变化的呢？\\
（1）先看看能否消掉参数 \(a\) ，得到 \(f_{2}(x)\) 过的定点。根据第二章与第四章中的内容，选 \(a\) 为主元：

\[
f_{2}(x)=(3-x) a+x^{2}-5
\]

所以 \(f_{2}(3)=4\) ．所以，\(f_{2}(x)\) 恒过 \((3,4)\) ．\\
（2）看顶点位置。事实上，由于 \(f_{2}(x)\) 的二次项系数已经确定，那么 \(f_{2}(x)\) 的形状就确定了。所以，只要确定它顶点的位置，\(f_{2}(x)\) 的图像就确定了。配方得

\[
\begin{aligned}
& f_{2}(x)=x^{2}-a x+3 a-5 \\
& \quad=\left(x-\frac{a}{2}\right)^{2}-\frac{a^{2}}{4}+3 a-5
\end{aligned}
\]

所以 \(f_{2}(x)\) 的顶点坐标为 \(\left(\frac{a}{2},-\frac{a^{2}}{4}+3 a-5\right)\) ．显然，纵坐标是横坐标的二次函数。令 \(\frac{a}{2}=t\) ，则 \(f_{2}(x)=(x-t)^{2}-\) \(t^{2}+6 t-5\) ．故顶点 \(M\left(t,-t^{2}+6 t-5\right)\) ．换言之，\(f_{2}(x)\) 的顶点在 \(y=-x^{2}+6 x-5\) 这个抛物线上！\\
有了（1），（2）两个观察，我们就可以清楚地画出 \(f_{2}(x)\) 的运动轨迹了。它是一个顶点在 \(y=-x^{2}+6 x-5\) ，且二次项系数为 1 的二次函数。视频中我们展示了相应的动画，这里给大家截取不同的 \(a, f_{2}(x)\) 的图像：\\
\includegraphics[max width=\textwidth, center]{2025_02_16_7c9aaf92bc99dc07f048g-156(1)}\\
\(a=2\)\\
\includegraphics[max width=\textwidth, center]{2025_02_16_7c9aaf92bc99dc07f048g-156}\\
\(a=4\)

图中红色虚线就是 \(f_{2}(x)\) 的图像，橘色线是 \(y=-x^{2}+6 x-5\) ．\\
这样，我们就成功建立起 \(f_{2}(x)\) 和 \(f(x)\) 的＂动态图像＂。系统被我们完全掌握清楚了，有一种＂把整个宇宙都握在手心＂的感觉。我们根据题目要求去移动 \(a\) 即可。\\
（1）\(f(x)\) 有至少三个零点，那么至少 \(f_{2}(x)\) 要有零点。橘色抛物线和 \(x\) 轴的交点分别为 \((1,0)\) 和 \((5,0)\) ，所以 \(\frac{a}{2} \geq 5\) 或 \(\frac{a}{2} \leq 1\) ，即 \(a \geq 10\) 或 \(a \leq 2\) 。\\
（2）分两种情况讨论：\\
\includegraphics[max width=\textwidth, center]{2025_02_16_7c9aaf92bc99dc07f048g-157}\\
（1）当 \(a \geq 10\) 时，根据动态图像，\(|x|-2\) 的图像一定在 \(f_{2}(x)\) 的零点上方。这样，\(f_{2}(x)\) 的零点，一定是 \(f(x)\) 的零点；同时，\(f_{2}(x)\) 的图像也一定在 \(|x|-2\) 的零点上方，所以 \(x=2, x=-2\) 也都是 \(f(x)\) 的零点。这样，\(f(x)\) 就至少有三个零点，满足题意；\\
（2）当 \(a \leq 2\) 时，根据动态图像，可以画出 \(f(x)\) 的图像。但是这个做法反而舍近求远：我们其实并不关心 \(f(x)\)的图像，只关心零点的个数。所以我们只要看 \(f_{2}(x), f_{1}(x)\) 二者零点的位置关系即可。当 \(a=2\) 时，\(f_{2}(x) \geq f_{1}(x)\) ，那么 \(f(x)=f_{1}(x)=|x|-2, f(x)\) 只有两个零点，不符合题意。所以考虑 \(a<2\) 的情况即可，此时 \(f_{2}(x)\) 有两个不同的零点 \(x_{1}<x_{2}\) 。那么，有哪些可能的位置关系呢？\\
（i）\(-2<x_{1}<x_{2}<2\) ：此时，\(x_{1}, x_{2}\) 不是 \(f(x)\) 的零点，因为 \(\left|x_{1}\right|-2<0=f_{2}\left(x_{1}\right),\left|x_{2}\right|-2<0=f_{2}\left(x_{2}\right)\) ，所以 \(f(x)\)至多两个零点；\\
（ii）\(-2<x_{1}<2<x_{2}\) ：此时，\(x_{1}, 2\) 不是 \(f(x)\) 的零点。由于 \(\left|x_{1}\right|-2<0=f_{2}\left(x_{1}\right)\) ，所以 \(f\left(x_{1}\right)=f_{1}\left(x_{1}\right)<0\) ；由于 2 介于二次函数的两根之间，所以 \(f_{2}(2)<0\) 。故 \(f(2)=f_{2}(2)<0\) ．\\
（iii）\(x_{1}<-2<x_{2}<2\) ：此时，\(x_{2},-2\) 不是 \(f(x)\) 的零点，理由同（ii）．\\
（iv）\(x_{1}<-2<2<x_{2}\) ：此时 \(-2,2\) 都不是 \(f(x)\) 的零点，因为 \(f(-2)<0, f(2)<0\) 。所以，\(a \leq 2\) 不符合题意。综上所述，\(a \geq 10\) 。事实上，我们可以把上述分类讨论进行化简：根据 \(f_{1}(x), f_{2}(x)\) 的性质，如果存在根的交叉，那么就会损失零点：\\
（1）若 \(x_{1}<x<x_{2}\) ，则 \(f_{2}(x)<0\) ．那么 \(f(x) \leq f_{2}(x)<0\) ，所以 \(x_{1}, x_{2}\) 之间没有 \(f(x)\) 的零点；\\
（2）若 \(-2<x<2\) ，则 \(f_{1}(x)<0\) ．那么 \(f(x) \leq f_{1}(x)<0\) ，所以 \(-2,2\) 之间没有 \(f(x)\) 的零点。\\
所以，必须有 \(x_{1}<x_{2}<-2<2\) 或 \(-2<2<x_{1}<x_{2}\) 。前者不满足题意，后者就是 \(a \geq 10\) 的情况。\\
我们来总结一下，使用动态图像的关键思路：\\
（1）首先通过思想试验，临界状态，极端遍历法等方法，体会系统本身的性质，得到初步的观察。这也是破题的第一步；\\
（2）和（1）并列的，可以同时分析自由度。通过思想试验的方法，理解每个自由度是如何影响系统的；\\
（3）当完成了（1），（2）两步后，我们让＂图像＂动起来，形成＂动态图像＂。随着自由变量发生改变，系统处于不同的状态，根据题意分析即可。生成动态图像的关键，是要把动点或者动曲线的＂轨迹＂找到。

上题的动态图像，是动曲线的轨迹。更常见的一类，比如 2023 年乙卷 12 题（第三节的例题），是研究动点的轨迹。即使是确定的系统（自由度为 0 ），也可以靠动态图像来理解。例如下面的解三角形题目：

例 5.11 在（凸）四边形 \(A B C D\) 中，\(A B=6, B C=4, \angle A B C=120^{\circ}, A D=2 B D, S_{\triangle B C D}=2 \sqrt{3}\) ，求 \(A D\) 的长度。

我们首先来分析系统的自由度。四边形有五个自由度，题目中有五个条件，所以整个系统是确定的。但是否唯一，不确定。前面三个条件，确定了三角形 \(A B C\)（ \(S A S\) 判定定理）。所以，只要再确定三角形 \(B C D\) ，整个图形就确定了。因此，问题变成了解三角形的问题。

固定三角形 \(A B C\) 。后面两个条件，确定了 \(D\) 点的位置。如何确定？只要找到＂轨迹＂交点即可：我们先把 \(D\) 看作在整个平面内自由运动的动点。根据这两个条件，找到它的轨迹，轨迹的交点，就是 \(D\) 的位置。\\
（1）\(A D=2 B D\) ，说明 \(D\) 在阿氏圆上。取 \(A B\) 靠近 \(B\) 的三等分点 \(F\) ，延长 \(A B\) 至 \(G\) 使得 \(A G=2 B G\) ．那么 \(D\)在以 \(F G\) 为直径的圆上。我们可以从以下两个角度理解：\\
（1）代数角度：以 \(A B\) 中点为坐标原点，如图建立平面直角坐标系。那么 \(A(-3,0), B(3,0)\) ．此时 \(F(1,0), G(9,0)\) ．设 \(D(x, y)\) ，那么

\[
|A D|=\sqrt{(x+3)^{2}+y^{2}},|B D|=\sqrt{(x-3)^{2}+y^{2}}
\]

根据题目条件，得

\[
\sqrt{(x+3)^{2}+y^{2}}=2 \sqrt{(x-3)^{2}+y^{2}}
\]

平方得

\[
\begin{gathered}
x^{2}+6 x+9+y^{2}=4 x^{2}-24 x+36+4 y^{2} \\
3 x^{2}+3 y^{2}-30 x+27=0 \\
x^{2}+y^{2}-10 x+9=0 \\
(x-5)^{2}+y^{2}=16
\end{gathered}
\]

故它是以 \((5,0)\) 为圆心，\(r=4\) 为半径的圆，确实是以 \(F G\) 为直径的圆。\\
（2）几何角度：取 \(F G\) 中点 \(M\) ，那么 \(F M=G M=4\) 。以 \(M\) 为圆心作半径为 4 的圆。注意到 \(|M B|=2,|M A|=8\) ，对于圆上任意一点 \(P\) ，有

\[
|P M|^{2}=|M A| \cdot|M B|=16
\]

所以 \(\triangle M P B \sim \triangle M A P\) ．所以

\[
\frac{|P A|}{|P B|}=\frac{|P M|}{|M B|}=2
\]

所以，点 \(D\) 在圆 \(M\) 上运动。\\
（2）\(S_{\triangle B C D}=2 \sqrt{3}\) ．设 \(D\) 到 \(B C\) 的距离为 \(h\) ，那么

\[
S_{\triangle B C D}=|B C| \cdot \frac{h}{2}=2 h=2 \sqrt{3}
\]

所以 \(h=\sqrt{3}\) ．故 \(D\) 在距 \(B C\) 为 \(\sqrt{3}\) 上的两条直线 \(l_{1}, l_{2}\) 上运动。\\
因此，\(D\) 既在圆 \(M\) 上动，也在 \(l_{1} \cup l_{2}\) 上动。这两个轨迹的交点，就是 \(D\) 点的位置。我们把整个图像画出来：\\
\includegraphics[max width=\textwidth, center]{2025_02_16_7c9aaf92bc99dc07f048g-158}

图中 \(D_{1}, D_{2}, D_{3}, D_{4}\) 都是满足条件的。如果在坐标系中看：由于 \(\angle A B C=120^{\circ}, B C=\)

4，所以 \(C(5,2 \sqrt{3})\) ，两条直线 \(l_{1}, l_{2}\) 和直线 \(B C: y=\sqrt{3}(x-3)\) 平行，且距离为 \(\sqrt{3}\) ，那么它们的方程分别为

\[
l_{1}: y=\sqrt{3}(x-1), l_{2}: y=\sqrt{3}(x-5)
\]

由于题目条件要求 \(A B C D\) 是凸四边形，所以只有 \(D_{1}\) 是符合条件的。注意到 \(\angle D F M=60^{\circ}, F M=D_{1} M=4\) ，所以 \(\triangle F D_{1} M\) 是一个等边三角形。所以 \(D_{1} F=4\) ，故三角形 \(A D_{1} F\) 是顶角为 \(120^{\circ}\) 的等腰三角形，所以 \(A D_{1}=4 \sqrt{3}\)

以上的方法是动态图像的方法。如果我们只去分析自由度，能不能得到一个代数的解法呢？事实上，我们在三角形 \(A B C\) 的基础上，只需要将剩下两个条件用两个自由度翻译好即可。如何选择自由度呢？题目条件求 \(A D\) ，我们可以自然地设 \(B D=x\) ．想要确定 \(D\) 点位置，还需要知道 \(\angle D B C\) 或者 \(\angle A B D\) ．设 \(\angle D B C=\alpha\) ．那么 \(\angle A B D=120^{\circ}-\alpha\) ．将两个条件翻译好即可：\\
（1）\(A D=2 B D\) ．那么 \(A D=2 x\) ．由于 \(\angle A B D=120^{\circ}-\alpha\) ，根据余弦定理：

\[
|A D|^{2}=|A B|^{2}+|B D|^{2}-2|A D| \cdot|B D| \cdot \cos \angle A B D
\]

故

\[
x^{2}+36-12 x \cos \left(120^{\circ}-\alpha\right)=4 x^{2}
\]

（2）\(S_{\triangle B C D}=2 \sqrt{3}\) ．根据三角形面积公式，

\[
S_{\triangle B C D}=\frac{1}{2}|B D| \cdot|B C| \cdot \sin \angle C B D=\frac{1}{2} \cdot x \cdot 4 \cdot \sin \alpha=2 x \sin \alpha=2 \sqrt{3}
\]

所以

\[
x \sin \alpha=\sqrt{3} .
\]

问题变成了解方程：

\[
\left\{\begin{array}{c}
x^{2}+36-12 x \cos \left(120^{\circ}-\alpha\right)=4 x^{2} \\
x \sin \alpha=\sqrt{3}
\end{array}\right.
\]

两个未知数，两个方程，系统确定，但是如何求解呢？首先将余弦展开：

\[
\begin{gathered}
3 x^{2}=36-12 x \cdot\left(-\frac{1}{2} \cos \alpha+\frac{\sqrt{3}}{2} \sin \alpha\right) \\
x^{2}=12-4\left(-\frac{1}{2} \cos \alpha+\frac{\sqrt{3}}{2} \sin \alpha\right) \\
x^{2}=12+2 x \cos \alpha-2 \sqrt{3} x \sin \alpha
\end{gathered}
\]

这样一来，可以代入 \(x \sin \alpha=\sqrt{3}\) ，得

\[
x^{2}=6+2 x \cos \alpha
\]

所以

\[
x \cos \alpha=\frac{x^{2}}{2}-3
\]

有了 \(x \sin \alpha, x \cos \alpha\) ，我们其实可以用 \(x\) 来表示 \(\sin \alpha, \cos \alpha\) ：

\[
\sin \alpha=\frac{\sqrt{3}}{x}, \cos \alpha=\frac{x}{2}-\frac{3}{x}
\]

到这里有多种处理手段，既可以消 \(x\) ，也可以消 \(\alpha\) ：\\
（1）消 \(\alpha\) ：利用 \(\cos ^{2} \alpha+\sin ^{2} \alpha=1\) ，得

\[
\begin{gathered}
\left(\frac{x}{2}-\frac{3}{x}\right)^{2}+\left(\frac{\sqrt{3}}{x}\right)^{2}=1 \\
\frac{x^{2}}{4}+\frac{12}{x^{2}}-3=1 \\
\frac{x^{4}}{4}-4 x^{2}+12=0 \\
\left(\frac{x^{2}}{2}-4\right)^{2}=4
\end{gathered}
\]

由于 \(x^{2}>0\) ，所以 \(\frac{x^{2}}{2}=6\) 或 2 ，即 \(x=2 \sqrt{3}\) 或 \(2 . x=2\) 要舍弃，因为此时 \(A D=4, B D=2\) ，而 \(A B=6\) ，所以 \(D\)在 \(A B\) 上，不构成四边形（对应方法一中的 \(D_{2}\) 的情况）。所以 \(x=2 \sqrt{3}, A D=2 x=4 \sqrt{3}\) ．\\
（2）消 \(x\) ：将 \(x=\frac{\sqrt{3}}{\sin \alpha}\) 代人 \(\cos \alpha\) 的表达式，得

\[
\begin{aligned}
& \cos \alpha=\frac{\sqrt{3}}{2 \sin \alpha}-\sqrt{3} \sin \alpha \\
& \sin \alpha \cos \alpha=\frac{\sqrt{3}}{2}-\sqrt{3} \sin ^{2} \alpha \\
& \sin \alpha \cos \alpha+\sqrt{3} \sin ^{2} \alpha=\frac{\sqrt{3}}{2}
\end{aligned}
\]

左侧都是 \(\sin \alpha, \cos \alpha\) 的一次式。根据第二章和第四章中的内容，它可以化为二倍角。所以

\[
\begin{gathered}
\frac{1}{2} \sin 2 \alpha+\sqrt{3} \cdot \frac{1-\cos 2 \alpha}{2}=\frac{\sqrt{3}}{2} \\
\sin 2 \alpha=\sqrt{3} \cos 2 \alpha \\
\tan 2 \alpha=\sqrt{3}
\end{gathered}
\]

所以 \(2 \alpha=\frac{\pi}{3}\) 或 \(\frac{4 \pi}{3}\) ，即 \(\alpha=\frac{\pi}{6}\) 或 \(\frac{2 \pi}{3}\) 。后者舍，此时 \(D\) 在 \(A B\) 上（对应 \(D_{2}\) ），所以 \(\alpha=\frac{\pi}{6}\) ， \(x=\frac{\sqrt{3}}{\sin \alpha}=2 \sqrt{3}\) ．故 \(|A D|=2 x=4 \sqrt{3}\) ．

\section*{5．4．2 控制变量法}
我们在举例时，往往都是固定其他所有变量，只让我们关心的自由度产生变化。这样就可以确定它对系统的影响。例如，某物理量是关于它单调递增还是单调递减。这个思路在处理很多问题是都是一种通用的方法。我们以一道＂简单＂的立体几何为例，来看控制变量法是如何起效的。

例 5.12 四面体 \(A B C D\) 的外接球半径为 \(R=2\) ．已知 \(A B=\sqrt{3}, \angle B C A=60^{\circ}\) ，求 \(A B C D\) 体积的最大值。

四面体本身有 6 个自由度，题目中有 3 个条件，故这是一个有 3 个自由度的系统。那么，这三个自由度分别是什么呢？我们可以这么确定：首先固定 \(A, B\) 两点，\(A B=\sqrt{3}\) ．分别来确定 \(C, D\) 两点的位置。\\
（1）由于 \(\angle B C A=60^{\circ}\) ，所以 \(C\) 在圆上运动。事实上，根据正弦定理，设三角形 \(A B C\) 的外接圆半径为 \(r\) ．那么

\[
2 r=\frac{A B}{\sin \angle B C A}=\frac{\sqrt{3}}{\frac{\sqrt{3}}{2}}=2
\]

故 \(r=1\) ．所以，\(C\) 在一个以 \(r=1\) 为半径的圆，圆心记为 \(O^{\prime}\) ．圆是一个一维图形（自由度为 \(1, O^{\prime} C\) 和 \(A B\) 的夹角确定了 \(C\) 的位置），\(C\) 有一个自由度；\\
（2）\(D\) 在球上运动。根据第三节中的内容，想要确定给定球面上的一个点，要确定投影的方向和仰角。所以有两个自由度。或理解成：球面本身是一个二维图形，所以 \(D\) 有两个自由度。\\
因此，我们可以分别固定 \(C, D\) ，看看另一个点满足什么条件时，体积取最大值。\\
（1）假设 \(A, B, D\) 点固定，\(C\) 点动。此时，平面 \(A B C\) 的位置也是确定的。这是因为 \(A B C\) 的外接圆半径为 \(r=1\) ，所以 \(O\) 到平面 \(A B C\) 的距离为

\[
d=\sqrt{R^{2}-r^{2}}=\sqrt{3}
\]

所以，整个平面也确定了。那么 \(D\) 到底面 \(A B C\) 的距离也就确定了。体积最大，相当于 \(\triangle A B C\) 面积最大。我们把圆 \(O^{\prime}\) 画出来：（原书中也无图片）

显然，当 \(C\) 在北极点的时候，三角形 \(A B C\) 面积最大。此时刚好是一个正三角形，面积为 \(S=\frac{\sqrt{3}}{4} \cdot(\sqrt{3})^{2}=\frac{3 \sqrt{3}}{4}\) 。 （2）假设 \(A, B, C\) 点固定，\(D\) 点动。此时，平面 \(A B C\) 的位置确定，\(O\) 到它们的距离 \(d\) 也确定，\(d=\sqrt{3}\) 。显然，\(D\) 在整个球的北极点时，体积最大。此时 \(D O \perp\) 底面 \(A B C\) 于 \(O^{\prime}\) 。此时四面体的高为

\[
h=O D+d=\sqrt{3}+2
\]

综合（1），（2），可得：

\[
V=\frac{1}{3} S h \leq \frac{1}{3} \cdot \frac{3 \sqrt{3}}{4} \cdot(2+\sqrt{3})=\frac{3+2 \sqrt{3}}{4}
\]

总结一下。控制变量一般是这样使用的：\\
（1）只变化一个变量，固定其他所有变量。看关心的量关于该变量是如何变化的；\\
（2）确保多个变量一起动的时候，每个都能取到（1）中相应的最值。这一步是非常重要的，否则会出现错误的放缩。\\
事实上，控制变量法就是我们在第二章或第四章里讲过的累次极值。

\section*{总结}
动态图像的精髓在于，把每个变量的完整连续变化过程想清楚。想要做到这一点，我们可以：\\
（1）确定自由度后，分别考虑每一个自由度。这需要我们使用控制变量法。\\
（2）为了确定每一个自由度的连续变化过程，可以灵活使用我们之前的所有方法。例如，思想试验，举高明的例子；考虑极端情况和临界状态等等。\\
这四节的内容综合运用，可以帮助大家系统地处理任何一个新的问题。这个方法论是普适的：无论是什么样的新问题，都可以用它理出思路；给定一个系统，通过这样的研究范式得到的观察，可以回答有关这个系统所有的问题。这是真正具有强大泛化性，高效且＂锋利＂的解题哲学。

\section*{5．4．3 XXXX}
\section*{5.5 数形结合}
\section*{5．5．1 数形结合的本质}
数形结合是一个我们耳熟能详的数学思想。它告诉我们：一个问题，可以分别从代数角度和几何角度去观察。两个角度，往往会有不同的收获。然而，在实际应用中，大部分此类题目都在加强学生的这种印象：数形结

合，其实是＂给一个看上去复杂的代数表达式，一个精巧的几何含义，然后题目就秒杀了＂。这种＂秒杀＂的快感是大部分人无法拒绝的：又能快速把题做对，把分拿到手，还可以在别人面前露一手，岂不两全其美？

这种印象到底对不对呢？我们的核心是希望建立起一个属于自己的＂破题＂系统。这个破题系统稳定而高效，兵来将挡，水来土掩。我们不想依靠秒杀，依靠＂巧妙＂的方法，得过且过：这道题能用，结果别的题又不能用了。当以这种标准来理解＂数形结合＂，我们才能够看到它的真面目。我们以 2022 年天津高考为例（天津高考是《选填破题基本法》的常客，因为他的小题质量还是非常不错的！即便是全国卷的同学，参考天津卷也是非常有帮助的）来看看：

例 5.13 （2022 天津）在三角形 \(\triangle A B C\) 中， \(\overrightarrow{C A}=\vec{a}, \overrightarrow{C B}=\vec{b}, D\) 是 \(A C\) 中点， \(\overrightarrow{C B}=2 \overrightarrow{B E}\) ．若 \(\overrightarrow{A B} \perp \overrightarrow{D E}\) ，则 \(\angle A C B\) 的最大值为 \(\qquad\) ＿．

看到这个题目，我相信很多同学第一反应，是直接利用向量翻译题目条件，然后求角度。题目问 \(\angle A C B\) 的最大值，又给出了向量，我们自然想到利用向量的内积来表示 \(\cos \angle A C B\) ，以求最值。事实上，

\[
\cos \angle A C B=\frac{\vec{a} \cdot \vec{b}}{|\vec{a}| \cdot|\vec{b}|}
\]

我们只要根据题目条件，得到 \(\vec{a}, \vec{b}\) 的等式即可。注意到

\[
\begin{gathered}
\overrightarrow{B E}=\frac{1}{2} \overrightarrow{C B}=\frac{1}{2} \vec{b} \\
\overrightarrow{D B}=\overrightarrow{D A}+\overrightarrow{A B}=\frac{1}{2} \overrightarrow{C A}+(\overrightarrow{C B}-\overrightarrow{C A})=\overrightarrow{C B}-\frac{1}{2} \overrightarrow{C A}=\vec{b}-\frac{1}{2} \vec{a} \\
\overrightarrow{D E}=\overrightarrow{D B}+\overrightarrow{B E}=\vec{b}-\frac{1}{2} \vec{a}+\frac{1}{2} \vec{b}=\frac{3}{2} \vec{b}-\frac{1}{2} \vec{a} \\
\overrightarrow{A B}=\overrightarrow{C B}-\overrightarrow{C A}=\vec{b}-\vec{a}
\end{gathered}
\]

所以

\[
\overrightarrow{A B} \cdot \overrightarrow{D E}=0 \Rightarrow\left(\frac{3}{2} \vec{b}-\frac{1}{2} \vec{a}\right) \cdot(\vec{b}-\vec{a})=0
\]

即

\[
\frac{3}{2}|\vec{b}|^{2}-2 \vec{a} \cdot \vec{b}+\frac{1}{2}|\vec{a}|^{2}=0
\]

这样一来，我们就可以用 \(|\vec{a}|,|\vec{b}|\) 来表示 \(\vec{a} \cdot \vec{b}\) ，以化简 \(\cos \angle A C B\) ：

\[
\cos \angle A C B=\frac{\frac{3}{4}|\vec{b}|^{2}+\frac{1}{4}|\vec{a}|^{2}}{|\vec{a}| \cdot|\vec{b}|}
\]

注意到，这是一个关于 \(|\vec{a}|,|\vec{b}|\) 的零次齐次式（第二章与第四章又出现了！），分子，分母都是二次，所以整个分式是零次的。如果看着没感觉，我们简单换个元：令 \(x=|\vec{a}|, y=|\vec{b}|\) ，那么

\[
\cos \angle A C B=\frac{\frac{3}{4} y^{2}+\frac{1}{4} x^{2}}{x y}=\frac{3 y^{2}+x^{2}}{4 x y}
\]

处理它，就有非常多的方法：\\
（1）根据齐次结构，它一定可以写作 \(t=\frac{x}{y}\) 的表达式。事实上，

\[
\frac{3 y^{2}+x^{2}}{4 x y}=\frac{1}{4} t+\frac{3}{4 t}
\]

这是一个关于 \(t\) 的对勾函数。根据均值不等式，

\[
\frac{1}{4} t+\frac{3}{4 t} \geq 2 \sqrt{\frac{1}{4} t \cdot \frac{3}{4 t}}=\frac{2 \sqrt{3}}{4}=\frac{\sqrt{3}}{2}
\]

所以 \(\cos \angle A C B \geq \frac{\sqrt{3}}{2}\) ．\\
（2）直接用均值不等式，其实也很快：由于 \(x, y \geq 0\) ，所以

\[
3 y^{2}+x^{2} \geq 2 \sqrt{3 y^{2} \cdot x^{2}}=2 \sqrt{3} x y
\]

所以

\[
\frac{3 y^{2}+x^{2}}{4 x y} \geq \frac{2 \sqrt{3} x y}{4 x y}=\frac{\sqrt{3}}{2}
\]

综上，\(\angle A C B\) 最大为 \(\frac{\pi}{6}\) 。\\
这是一个相对比较直接的代数方法。这个方法的关键，在于通过垂直条件，建立起了 \(\vec{a} \cdot \vec{b}\) 和 \(|\vec{a}|,|\vec{b}|\) 的联系。\\
从图像上，如何理解这个条件呢？事实上，如果有心的同学会发现，我们想画出一个相对准确的图是比较困难的。如果画一个一般的三角形 \(A B C\) ，那么 \(A B, D E\) 的夹角大概率是一个锐角，而且是一个很小的锐角，很难是 \(90^{\circ}\) 。如下图所示：\\
\includegraphics[max width=\textwidth, center]{2025_02_16_7c9aaf92bc99dc07f048g-163}

那么，如何画出一个相对标准的图呢？事实上，之所以 \(A B, D E\) 的角度无法控制，是因为它们的交点不是我们作图中已知的点。例如，\(C A, C B\) 的夹角很好控制，因为它们就交于点 \(C, C\) 是我们作图中基本的点。因此，想要控制好 \(A B, D E\) 的夹角，我们可以利用平行来转移其中一条线段，让它们有一个公共的基本点。所以，我们先画出比较标准的 \(90^{\circ}\) ，再找到其他点。

先做 \(A B \perp A F\) ．相当于：过 \(A\) 作 \(D E\) 的平行线交 \(B C\) 于 \(F\) ．\\
\includegraphics[max width=\textwidth, center]{2025_02_16_7c9aaf92bc99dc07f048g-163(1)}

由于 \(A B \perp D E\) ，所以 \(A F \perp A B\) ．同时，由于 \(D\) 是 \(A C\) 中点，所以 \(E\) 也是 \(C F\) 中点。而由于 \(\overrightarrow{C B}=2 \overrightarrow{B E}\) ，所以 \(B\) 是 \(C E\) 上靠近 \(E\) 的三等分点。由于 \(E\) 是 \(C F\) 中点，所以 \(C B, B E, E F\) 这三段的比例为 \(|C B|:|B E|:|E F|=2: 1: 3\) ．\\
所以，\(E\) 是 \(B F\) 上靠近 \(B\) 的四等分点，我们可以取出 \(E\) 点。\\
\includegraphics[max width=\textwidth, center]{2025_02_16_7c9aaf92bc99dc07f048g-164}

又由于 \(E\) 应为 \(C F\) 中点，所以可以倍长 \(F E\) ，得到 \(C\) 。再取 \(A C\) 中点，就得到 \(D\) 。这样，一个相对标准的图就画出来了。\\
\includegraphics[max width=\textwidth, center]{2025_02_16_7c9aaf92bc99dc07f048g-164(1)}

到这里，很多同学可能会有疑问：我们画一个示意图，难道不行吗？为什么要画这么精准的一个图？它对我们做题真的有帮助吗？答案是有的：这就要用我们上一节动态图像的内容来理解。题目中垂直的条件在几何上是非常难处理的，经过平行转移，我们直接确定了 \(A\) 点的轨迹：它是以 \(B F\) 为直径的圆上一点。 \(F\) 是怎么来的呢？哦！是我们倍长 \(D E\) 得到的。所以，这道题我们可以重新改写：\\
＂\(A\) 是以 \(B F\) 为直径的圆 \(O\) 上一点，\(C\) 在射线 \(F B\) 上，满足 \(C B=\frac{1}{2} B F\) 。求 \(\angle A C B\) 的最大值。＂\\
那什么时候最大呢？显然是 \(C A\) 和圆 \(O\) 相切的时候最大。此时，\(\angle C A O=90^{\circ}\) 。由于 \(C B=\frac{1}{2} B F\) ，所以 \(C O=B F=2 O A\) ．此时 \(\angle A C B=30^{\circ}\) ．因此，\(\angle A C B\) 的最大值为 \(\frac{\pi}{6}\) 。

这个几何方法非常神奇！只是精准地把图形画出来，问题就得以解决，没有任何复杂的计算。这也是几何方法常常具有的巧妙之处，它充分尊重了问题本身的结构。如何尊重的？找到动态图像和轨迹，找到关键的自由度，一下就把极端情况找到，得到答案。

我们来从极端情况，自由度，动态图像的角度，来理解这个方法。\\
（1）整个系统是一个自由度为 1 的系统。我们从题目条件分析：三角形 \(A B C\) 本身有三个自由度，由于我们只关心形状，应有 2 个自由度；又规定了 \(E D \perp A B \quad(D, E\) 点本身都由三角形 \(A B C\) 决定的，不影响自由度），所以应有 1 个自由度。\\
（2）同时，通过精准的画图，\(D, E\) 两点被我们成功消掉，化成了更重要的 \(F\) 点。这样，题目中只有四个点 \(A, B, C, F\) ．\\
（3）自由度为 1 ，就需要找出动点轨迹。事实上，\(A\) 点的轨迹（圆 \(O\) ）是一个阿氏圆。注意到 \(|C B|=2|B E|,|C F|=\) \(2|F E|\) ，所以 \(B F\) 是这个阿氏圆的直径，圆 \(O\) 上的点 \(A\) 均满足 \(|C A|=2|A E|\) ！不过，解决这个题目并不需要阿氏圆的具体知识，只需知道＂直径所对圆周角是 \(90^{\circ}{ }^{\circ}\) 。\\
（4）极端情况：相切时角度最大。\\
上面这个例题，体现了几何方法的优点。那么代数方法有什么优点呢？它的优点看上去并不如几何方法那么吸引人，但其实是一种非常＂厚重，朴素＂的有点：那就是稳定，且接近通法。什么意思？如果我们能够把几

何问题转化为一个可计算的代数问题，那么只要把它＂按部就班＂地算出来就好，不太需要我们有太多的＂奇思妙想＂。这也是解析几何诞生的主要动机。

\section*{5．5．2 什么是＂等值线＂？}
我们在前面的章节，已经出现过等值线的思想，但是没有过系统的梳理和介绍。作为数形结合的具体应用，在这里详细介绍它的概念和使用方法。

在物理，地理等学科中，我们都接触过＂等某线＂的讲法。在物理学中，描述电场时，就会用到＂等势线＂：一条等势线上的点，势能相等。在地理学中，＂等高线＂上的点的海拔相同。从物理上看，它其实是重力势能的＂等势线＂。在数学中，我们可以引入同样的概念，来描述一个轨迹，使得在这个轨迹上任意一点 \(P\) ，关于它的某个＂势能＂的值 \(h(P)\) 是定值。这就是＂等值线＂。

我们学过的哪些轨迹，是等值线呢？

\section*{5．5．2．1 直线是等值线}
考虑 \(x+y=1\) ．对于这条直线上的点 \(P(x, y)\) ，哪个值 \(h(P)\) 是不变的呢？从直线方程就可以直接看出来：

\[
h(P)=x+y=1
\]

所以，\(l: x+y=1\) 是 \(h(P)=x+y\) 的等值线。\\
从几何上看，它还和向量的分解有关。令 \(A(0,1), B(1,0)\) 。根据平面向量基本定理，存在唯一的实数 \(\lambda, \mu\) ，使得

\[
\overrightarrow{O P}=\lambda \overrightarrow{O A}+\mu \overrightarrow{O B}
\]

定义新的势能 \(h(P)=\lambda+\mu\) 。由 \(P(x, y)\) 知 \(\mu=x, \lambda=y\) ．对于直线上 \(l\) 上的点，它关于 \(\overrightarrow{O A}, \overrightarrow{O B}\) 两个向量的表示的系数之和 \(h(P)=\lambda+\mu\) 也是定值。反过来，对于空间上一点 \(P\) ，若 \(\overrightarrow{O P}=\lambda \overrightarrow{O A}+\mu \overrightarrow{O B}\) ，且 \(h(P)=\lambda+\mu=1\) ，则 \(P\)在直线 \(x+y=1\) 上。所以，\(l: x+y=1\) 是 \(h(P)=\lambda+\mu\) 的等值线。

还是考虑这个势能 \(h(P)=\lambda+\mu\) ．那么 \(h(P)=2\) 的等值线，是哪条直线呢？这意味着，\(P\) 满足：

\[
\overrightarrow{O P}=\lambda \overrightarrow{O A}+(2-\lambda) \overrightarrow{O B}
\]

所以 \(P(\lambda, 2-\lambda)\) 。我们发现，它的横纵坐标之和为 \(\lambda+2-\lambda=2\) 。所以，\(P\) 必然在 \(x+y=2\) 上。反过来，也成立。所以，\(h(P)=2\) 对应的等值线，是 \(l_{2}: x+y=2\) 。一般地，对于任意实数 \(k, h(P)=k\) 对应的等值线，是 \(x+y=k\) 。

从这个例子，我们可以做一个简单的推广。还是考虑 \(h(P)=\lambda+\mu\) ，但：\\
（1）考虑平面上任意的不共线三点 \(O, A, B, h(P)=1\) 的等值线，就是直线 \(A B\) 。这是我们在学习向量时，学过的结论。\\
（2）考虑任意的势能值 \(k, h(P)=k\) 的等值线其实是和 \(A B\) 平行的直线。如何理解？只要进行一个＂位似变换＂即可。我们以 \(k=2\) 为例。倍长 \(O A, O B\) ，得到 \(A^{\prime}, B^{\prime}\) 。那么

\[
\overrightarrow{O A}=2 \overrightarrow{O A}, \overrightarrow{O B}=2 \overrightarrow{O B}
\]

这样一来，若 \(h(P)=2\) ，则

\[
\overrightarrow{O P}=\lambda \overrightarrow{O A}+\mu \overrightarrow{O B}=\frac{\lambda}{2} \overrightarrow{O A^{\prime}}+\frac{\mu}{2} \overrightarrow{O B^{\prime}}
\]

此时，

\[
\frac{\lambda}{2}+\frac{\mu}{2}=\frac{h(P)}{2}=1
\]

所以 \(P\) 必然在直线 \(A^{\prime} B^{\prime}\) 上。所以，\(h(P)=2\) 的等值线，就是直线 \(A^{\prime} B^{\prime}\) 。此时，\(A B\) 刚好是 \(A^{\prime} B^{\prime}\) 的中位线。一般地，\(h(P)=k\) 的等值线，只要作 \(A B\) 的平行线 \(l\) ，交直线 \(O A, O B\) 于 \(A^{\prime}, B^{\prime}\) ，使得

\[
\overrightarrow{O A^{\prime}}=k \overrightarrow{O A}, \overrightarrow{O B^{\prime}}=k \overrightarrow{O B}
\]

即可。所以，如果想要求势能值 \(k\) ，看方向和比例 \(\frac{\left|O A^{\prime}\right|}{|O A|}, \frac{\left|O B^{\prime}\right|}{|O B|}\) ，或者 \(\frac{\left|A^{\prime} B^{\prime}\right|}{|A B|}\) 都可以。\\
当我们有了这个观点后，很多有关向量表示的题目，可以直接通过数形结合得到结果。例如2009年安徽高考题：

例 5.14 （2009 安徽）给定单位向量 \(\overrightarrow{O A}, \overrightarrow{O B}\) ，它们的夹角为 \(120^{\circ}\) 。点 \(C\) 在以 \(O\) 为圆心的圆弧 \(A B\)（劣弧）上运动。若 \(\overrightarrow{O C}=x \overrightarrow{O A}+y \overrightarrow{O B}(x, y \in \mathbb{R})\) ，求 \(x+y\) 的最大值。

我们来比较建系和等值线两个不同的方法。\\
（1）代数方法：建系。由于 \(\angle A O B=120^{\circ}\) ，不妨设 \(A\left(-\frac{1}{2}, \frac{\sqrt{3}}{2}\right), B(1,0)\) ．那么

\[
\overrightarrow{O P}=x \cdot \overrightarrow{O A}+y \cdot \overrightarrow{O B}=\left(y-\frac{1}{2} x, \frac{\sqrt{3}}{2} x\right)
\]

由 \(|\overrightarrow{O P}|=1\) 知

\[
\left(y-\frac{1}{2} x\right)^{2}+\left(\frac{\sqrt{3}}{2} x\right)^{2}=x^{2}-x y+y^{2}=1
\]

在这个条件下，求 \(x+y\) 的最大值。方法较多，可以用均值不等式：根据均值不等式，有

\[
1=x^{2}-x y+y^{2} \geq 2 x y-x y=x y
\]

而

\[
(x+y)^{2}=x^{2}+2 x y+y^{2}=1+3 x y \leq 4 .
\]

所以 \(x+y \leq 2\) ．当且仅当 \(x=y=1\) 时取等。\\
（2）几何方法：等值线。求 \(x+y\) 的最大值，我们只要看等值线 \(h(P)=x+y\) 是哪条直线，看圆弧上的哪个点势能最大。我们知道，\(h(P)\) 的等值线是平行于 \(A B\) 的直线，因此，最大值所在的等值线 \(l\) ，一定是在平行于 \(A B\) ，且和圆弧相切的直线。设 \(l\) 和直线 \(O A, O B\) 交于 \(A^{\prime}, B^{\prime}\) 。它的势能值 \(h(P)=k\) ，就是 \(\frac{\left|O A^{\prime}\right|}{|O A|}\) 。根据相似三角形的性质，我们除了看这个比例，还可以看：\(O\) 到 \(A^{\prime} B^{\prime}\) 的距离 \(d^{\prime}\) ，以及 \(O\) 到 \(A B\) 的距离 \(d\) 。因为

\[
|k|=\frac{\left|O A^{\prime}\right|}{|O A|}=\frac{\left|O B^{\prime}\right|}{|O B|}=\frac{\left|A^{\prime} B^{\prime}\right|}{|A B|}=\frac{d^{\prime}}{d}
\]

为什么想到看这个比例？因为切线：由于 \(A^{\prime} B^{\prime}\) 和圆 \(O\) 相切，所以 \(d^{\prime}=1\) 。在三角形\\
\(O A B\) 中，它是一个 \(1,1, \sqrt{3}\) 的顶角为 \(120^{\circ}\) 的等腰三角形。所以 \(d=\frac{1}{2}\) ．故取最大值时的 \(k=\frac{d^{\prime}}{d}=2\) ．故最大值为 2 。

很明显，等值线的方法计算量少，出答案快。这是因为它直接找到了临界状态，省去了代数运算。

\section*{5．5．2．2 圆是等值线}
我们可以从两个角度来理解圆是等值线。\\
（1）同弧所对圆周角相等。例如，如果给定 \(A, B\) 两点，对于点 \(P\) ，考虑势能 \(h(P)=\angle A P B\) 。那么，\(h(P)\) 的等势线是（两段）圆弧。以 \(A B\) 为边作正三角形 \(A B C\) ，那么优弧 \(A C B\) 上的点 \(P\) 满足 \(h(P)=\angle A C B=\frac{\pi}{3}\) 。所以，优

弧 \(A C B\) 就是 \(h(P)=\frac{\pi}{3}\) 的等值线（的一部分，另一部分是这个圆弧关于 \(A B\) 对称的圆弧）。这个理解在前面的题目中也反复出现，我们省去例题；\\
（2）圆上的点到圆心距离相等。这是圆的定义。所以，当我们出现距离问题时，可以考虑用圆来理解。以如下问题为例：

例 5.15 设 \(A\) 是曲线 \(g(x)=e^{x}-1\) 上一点。设 \(B(1,-1)\) ，求 \(|A B|\) 的最小值。

我们还是比较代数方法和几何方法。\\
（1）代数方法：直接计算。设 \(A\left(x, e^{x}-1\right)\) ．那么

\[
|A B|^{2}=(x-1)^{2}+\left(e^{x}-1+1\right)^{2}=e^{2 x}+x^{2}-2 x+1
\]

只需求 \(f(x)=|A B|^{2}=e^{2 x}+x^{2}-2 x+1\) 的最小值。求导得

\[
f^{\prime}(x)=2 e^{2 x}+2 x-2=2\left(e^{2 x}+x-1\right)
\]

显然 \(f^{\prime}(0)=0\) ，且 \(f^{\prime}(x)\) 单调递增（ \(e^{2 x}\) 单增，\(x\) 也单增）。所以，\(f^{\prime}(x)\) 在 \((-\infty, 0)\) 上正，在 \((0,+\infty)\) 上负。故 \(f(x)\) 在 \((-\infty, 0)\) 上单调递减，在 \((0,+\infty)\) 上单调递增，所以

\[
f(x) \geq f(0)=2
\]

所以 \(|A B| \geq \sqrt{2}\) ，当且仅当 \(x=0\) 时取到。\\
（2）几何方法：我们以 \(B\) 为圆心作圆。圆上的点到 \(B\) 点距离相等。随着这个圆扩大，当它和曲线 \(y=e^{x}-1\) 产生第一个交点 \(A\) 时，此时的 \(|A B|\) 就是距离的最小值。何时产生第一个交点？应该是＂相切＂的时候。但是，我们没有学过曲线和曲线的相切。事实上，两条曲线在某点处相切，意味着它们在该点处的切线是同一条。设此时 \(A\left(x_{0}, e^{x_{0}}-1\right)\) ．圆的切线垂直于半径，所以切线斜率为

\[
k=-\frac{1}{k_{A B}}=\frac{1-x_{0}}{e^{x_{0}}}
\]

同时：

\[
k=g^{\prime}\left(x_{0}\right)=e^{x_{0}}
\]

所以

\[
\begin{gathered}
\frac{1-x_{0}}{e^{x_{0}}}=e^{x_{0}} \\
e^{2 x_{0}}+x_{0}-1=0
\end{gathered}
\]

我们惊奇地发现，这就是代数方法中 \(f^{\prime}(x)\) 的一半！所以，代数方法中 \(f^{\prime}(x)=0\) ，除了理解成距离取最小值，还可以理解圆和曲线相切。

由于 \(h(x)=e^{2 x}+x-1\) 单调递增，且 \(h(0)=0\) ，所以 \(x_{0}=0\) ．所以最小值在 \(A(0,0)\) 处取到，距离为 \(|A B|=\sqrt{2}\) ．在例题中，我们会接着看到其他等值线。

5．5．3 XXXX

\section*{5.6 基底思想}
\section*{5．6．1 XXXX}
在高中数学中，我们学习过的＂最抽象＂的两个定理，莫过于两个向量的＂基本定理＂：平面向量基本定理和空间向量基本定理。它们的叙述拗口，字母繁多，而且逻辑关系难以理清，所以很多同学理解起来有非常大的困难。然而，这两个定理背后所体现的＂基底＂思想，其实是向量（甚至是线性代数）中最重要的思想之一。它不光在数学上重要，在解题上也极其重要，我们以一道题目为例来体会它的使用。

例 5.16 如图，平行六面体 \(A B C D-A_{1} B_{1} C_{1} D_{1}\) 的底面 \(A B C D\) 是菱形，且 \(\angle C_{1} C B=\angle C_{1} C D=\angle B C D=\frac{\pi}{3}, C D=\) \(C C_{1}\) ．求证：\(C A_{1} \perp\) 平面 \(C_{1} B D\) ．

乍一看，这是一个有高度对称性的图形。底面是菱形，还有非常多的 \(60^{\circ}\) 角，建系应该很简单。连接 \(A C, B D\) ，根据菱形的性质，\(A C \perp B D\) ，所以可以以它们的交点 \(O\) 为坐标原点建系，以 \(\overrightarrow{O D}, \overrightarrow{O A}\) 方向分别为 \(x, y\) 轴正方向。 \(x, y\) 轴有了，\(z\) 轴在哪里呢？

事实上，这里的 \(z\) 轴并不好确定，因为没有直接和底面垂直的直线。建系的方法下，只能利用 \(\angle C_{1} C D=\) \(\angle C_{1} C B=\frac{\pi}{3}\) 来写出 \(C_{1}\) 的坐标，进而得到所有点的坐标。设菱形的边长为 2 。设 \(C_{1}(x, y, z)\) ．其他需要的点的坐标为：

\[
D(1,0,0), C(0,-\sqrt{3}, 0), B(-1,0,0)
\]

（1）\(C C_{1}=2\) ，所以 \(x^{2}+(y+\sqrt{3})^{2}+z^{2}=4\) ，这是第一个方程；\\
（2）\(\angle C_{1} C D=\frac{\pi}{3}\) ．而 \(C C_{1}=C D=2\) ，所以

\[
\begin{gathered}
\overrightarrow{C C_{1}} \cdot \overrightarrow{C D}=\cos \angle C_{1} C D \cdot\left|C C_{1}\right| \cdot|C D|=2 \\
\overrightarrow{C C_{1}}=(x, y+\sqrt{3}, z), \overrightarrow{C D}=(1, \sqrt{3}, 0)
\end{gathered}
\]

所以

\[
\overrightarrow{C C_{1}} \cdot \overrightarrow{C D}=x+\sqrt{3} y+3=2
\]

所以 \(x+\sqrt{3} y=-1\) ．\\
（3）\(\angle C_{1} C B=\frac{\pi}{3}\) ．和（2）同理，得 \(\overrightarrow{C C_{1}} \cdot \overrightarrow{C B}=2\) ，而 \(\overrightarrow{C B}=(-1, \sqrt{3}, 0)\) ，所以

\[
-x+\sqrt{3} y=-1
\]

这样，我们得到了关于 \(x, y, z\) 的三个方程，可以确定 \(C_{1}\) 点坐标了。由于 \(x+\sqrt{3} y=-x+\sqrt{3} y=-1\) ，所以 \(x=0, y=-\frac{\sqrt{3}}{3}\) ．代入回（1）中方程，有

\[
\begin{gathered}
\frac{4}{3}+z^{2}=4 \\
z^{2}=\frac{8}{3}
\end{gathered}
\]

我们不妨设 \(z>0\)（另一种情况完全对称），所以 \(z=\frac{2 \sqrt{6}}{3}\) ．所以 \(C_{1}\left(0,-\frac{\sqrt{3}}{3}, \frac{2 \sqrt{6}}{3}\right)\) ．那么

\[
\overrightarrow{C_{1} A_{1}}=\overrightarrow{C A}=(0,2 \sqrt{3}, 0)
\]

所以 \(A_{1}\left(0, \frac{5 \sqrt{3}}{3}, \frac{2 \sqrt{6}}{3}\right)\) ．那么 \(\overrightarrow{C A_{1}}=\left(0, \frac{8 \sqrt{3}}{3}, \frac{2 \sqrt{6}}{3}\right)\) ．为了证明 \(\overrightarrow{C A_{1}} \perp\) 平面 \(C_{1} B D\) ，只需证明：\\
（1） \(\overrightarrow{C A_{1}} \cdot \overrightarrow{B D}=0 \cdot \overrightarrow{B D}=(2,0,0)\) ，确实有 \(\overrightarrow{C A_{1}} \cdot \overrightarrow{B D}=0 \times 2+\frac{8 \sqrt{3}}{3} \times 0+\frac{2 \sqrt{6}}{3} \times 0=0\) ；\\
（2） \(\overrightarrow{C A_{1}} \cdot \overrightarrow{O C_{1}}=0 \cdot \overrightarrow{O C_{1}}=\left(0,-\frac{\sqrt{3}}{3}, \frac{2 \sqrt{6}}{3}\right)\) ，所以

\[
\overrightarrow{C A_{1}} \cdot \overrightarrow{O C_{1}}=\left(-\frac{\sqrt{3}}{3}\right) \cdot\left(\frac{8 \sqrt{3}}{3}\right)+\left(\frac{2 \sqrt{6}}{3}\right)^{2}=-\frac{8}{3}+\frac{8}{3}=0
\]

成立。\\
\(C A_{1}\) 垂直于平面 \(B C_{1} D\) 内两条相交直线（ \(B D, O C_{1}\) 交于 \(O\) ），根据线面垂判定定理，\(C A_{1}\) 垂直于平面 \(B C_{1} D\) ．虽然我们顺利地用建系的方法把问题解决了，但是计算量不小，而且建系都没建＂全＂，连 \(z\) 轴是哪条直线都没搞清楚。根据题目条件，国脚地翻译得到 \(C_{1}\) 坐标的三个方程，得到了答案。那有没有不用建系，更快的方法呢？相信很多读者从标题中可以得到答案：应该是用向量吧？怎么用呢？在这之前，我们先来理解：建系的方法，本质上是在做一件什么事？事实上，向量的三个坐标，分别代表从坐标原点出发，应该走多少个单位。例如，\(\vec{a}=(1,2,3)\) ，意味着：先沿着 \(x\) 轴正方向走 1 个单位，再沿着 \(y\) 轴正方向走 2 个单位，再沿着 \(z\) 轴正方向走 3 个单位。这样得到的向量，就是 \(\vec{a}\) ．如果我们设 \(x, y, z\) 轴正方向的单位向量分别为 \(\overrightarrow{e_{x}}=(1,0,0), \overrightarrow{e_{y}}=(0,1,0), \overrightarrow{e_{z}}=(0,0,1)\) ，那么

\[
\vec{a}=1 \cdot \overrightarrow{e_{x}}+2 \cdot \overrightarrow{e_{y}}+3 \cdot \overrightarrow{e_{z}}
\]

所以，我们建系写向量坐标，本质上就是使用＂空间向量基本定理＂：用 \(\overrightarrow{e_{x}}, \overrightarrow{e_{y}}, \overrightarrow{e_{z}}\) 这组基底，来表示空间中的所有向量，进而表示点的坐标。

那既然是选择基底去表示向量，我们一定要选择这样两两垂直的向量吗？不一定。之所以选择它们，是因为计算＂向量的内积＂方便：

\[
\overrightarrow{e_{x}} \cdot \overrightarrow{e_{y}}=\overrightarrow{e_{x}} \cdot \overrightarrow{e_{z}}=\overrightarrow{e_{y}} \cdot \overrightarrow{e_{z}}=0
\]

所有交叉项都是 0 ，所以只要对应坐标相乘即可。但是，这并不是必需的：如果题目中不好找到两两垂直的向量，找到不共面的，有其他好性质的向量，也可以作为基底。

对于这道题而言，一个天生的基底，就是 \(\vec{a}=\overrightarrow{C D}, \vec{b}=\overrightarrow{C C_{1}}, \vec{c}=\overrightarrow{C B}\) 。为什么它们比两两垂直的更合适呢？\\
（1）根据空间向量基本定理，这三个向量可以表示所有向量，且表示方法唯一，所以它们是基底；\\
（2）它们两两夹角虽然不是 \(90^{\circ}\) ，但是是 \(60^{\circ}\) ，计算内积也并不复杂。延续之前的设法，设 \(|C D|=\left|C C_{1}\right|=|C B|=2\) ，那么

\[
\vec{a} \cdot \vec{b}=\vec{a} \cdot \vec{c}=\vec{b} \cdot \vec{c}=2
\]

这样，计算内积也并不复杂；\\
（3）容易表示需要的向量。 \(C D, C C_{1}, C B\) 刚好是这个平行六面体的三条相交的棱，那么表示其他以平行六面体的顶点为端点的向量是非常容易的。\\
综合这三个考虑，我们可以选 \(\vec{a}=\overrightarrow{C D}, \vec{b}=\overrightarrow{C C_{1}}, \vec{c}=\overrightarrow{C B}\) 为基底去表示所有向量。

由于 \(\overrightarrow{C A_{1}}\) 是体对角线，所以

\[
\begin{gathered}
\overrightarrow{C A_{1}}=\overrightarrow{C D}+\overrightarrow{D A}+\overrightarrow{A A_{1}}=\vec{a}+\vec{b}+\vec{c} \\
\overrightarrow{B D}=\overrightarrow{C D}-\overrightarrow{C B}=\vec{a}-\vec{c} \\
\overrightarrow{B C_{1}}=\overrightarrow{C C_{1}}-\overrightarrow{C B}=\vec{b}-\vec{c}
\end{gathered}
\]

只需验证：\\
（1） \(\overrightarrow{C A_{1}} \cdot \overrightarrow{B D}=0\) ．注意到

\[
\begin{gathered}
\overrightarrow{C A_{1}} \cdot \overrightarrow{B D}=(\vec{a}+\vec{b}+\vec{c}) \cdot(\vec{a}-\vec{c}) \\
=\vec{b} \cdot(\vec{a}-\vec{c})+(\vec{a}+\vec{c}) \cdot(\vec{a}-\vec{c}) \\
=2-2+|\vec{a}|^{2}-|\vec{c}|^{2} \\
=0
\end{gathered}
\]

成立！\\
（2） \(\overrightarrow{C A_{1}} \cdot \overrightarrow{B C_{1}}=0\) ．注意到

\[
\begin{aligned}
\overrightarrow{C A_{1}} \cdot \overrightarrow{B D} & =(\vec{a}+\vec{b}+\vec{c}) \cdot(\vec{b}-\vec{c}) \\
& =\vec{a} \cdot(\vec{b}-\vec{c})+(\vec{b}+\vec{c}) \cdot(\vec{b}-\vec{c}) \\
& =2-2+|\vec{b}|^{2}-|\vec{c}|^{2} \\
& =0 .
\end{aligned}
\]

成立！\\
这样，我们就不用求任何点的坐标，直接利用内积验证垂直，就完成了证明。\\
基底思想不光在向量中有应用，在我们学习过的其他模块中也有应用。\\
（1）在学习对数时，我们经常会遇到这样的题目：设 \(a=\lg 2 b=\lg 3\) ，用 \(a b\) 来表示 \(\lg 150\) 。利用对数的性质，

\[
\lg 150=\lg 3 \times 2 \times 5^{2}=\lg 3+\lg 2+2 \lg 5
\]

由于 \(\lg 2+\lg 5=\lg 10=1\) ，所以 \(\lg 5=1-\lg 2=1-a\) 。所以

\[
\lg 150=b+a+2(1-a)=2-a+b
\]

我们发现： \(\lg 2 \lg 3\) 是＂独立＂的，彼此之间无法＂表示 \("\) ；但是 \(\lg 2, \lg 5\) 不是独立的，因为它们的和是 1 ．所以，我们可以用 \(a, b\) 去表示所有因子只有 \(2,3,5\) 的整数的对数。再复杂一点：

\[
\lg 675=\lg 27+\lg 25=3 \lg 3+2 \lg 5=3 b-2 a+2
\]

（2）学习三角函数时，我们可以用 \(\sin x, \cos x\) 去表示所有形如 \(f(x)=A \sin (x+\varphi)\) 的函数。根据辅助角公式，

\[
f(x)=(A \cos \varphi) \sin x+(A \sin \varphi) \cos x
\]

如果我们把形如 \(f(x)=A \sin (x+\varphi)\) 的函数看成一个集合（空间），那么 \(\sin x, \cos x\) 就是＂基底＂。\\
（3）学习函数时，所有定义在 \(\mathbb{R}\) 上的函数，都可以表示为奇函数，偶函数的和。这个分解可以＂显式＂地写出：对于 \(f(x)\) ，令

\[
f_{e}(x)=\frac{f(x)+f(-x)}{2}, f_{o}(x)=\frac{f(x)-f(-x)}{2}
\]

那么

\[
\begin{aligned}
& f_{e}(-x)=\frac{f(-x)+f(x)}{2}=f_{e}(x) \\
& f_{o}(-x)=\frac{f(-x)-f(x)}{2}=-f_{o}(x)
\end{aligned}
\]

所以 \(f_{e}(x), f_{o}(x)\) 分别是偶函数，奇函数。显然，有

\[
f_{e}(x)+f_{o}(x)=f(x)
\]

所以，\(f(x)\) 写成了一个奇函数和一个偶函数的和。想一想，这个表示是不是唯一的？为什么？\\
事实上，（2）中的写成正弦函数，余弦函数的做法，就是将形如 \(f(x)=A \sin (x+\varphi)\) 写成了奇函数（正弦），偶函数（余弦）的部分：

\[
\begin{aligned}
& f_{e}(x)=\frac{f(x)+f(-x)}{2}=\frac{A \sin (x+\varphi)+A \sin (-x+\varphi)}{2}=(A \sin \varphi) \cos x \\
& f_{o}(x)=\frac{f(x)-f(-x)}{2}=\frac{A \sin (x+\varphi)-A \sin (-x+\varphi)}{2}=(A \cos \varphi) \sin x
\end{aligned}
\]

（4）所有二元对称式，都可以用 \(x+y, x y\)（和常数）表示。这个在第二章和第四章中我们详细讨论过，这里省略；

我们发现，基底思想，本质上就是在＂表示＂，而且是＂最小表示＂。我们试图用个数最少的量，来表示所有需要的东西：用 \(\lg 2, \lg 3\) 表示对数，不需要 \(\lg 5\) ；用 \(\sin x, \cos x\) 表示对数，不需要其他的；用奇函数，偶函数表示所有定义在 \(\mathbb{R}\) 上的函数；二元对称式，用 \(x+y, x y\) 可以表示所有的变量。

这个思想我们在哪里见过，讨论过呢？其实就是＂自由度＂，＂消元＂。只不过，这里不再是单纯的代数式的表示，而是函数，向量等更复杂内容的表示了。但无论是消元还是基底思想，本质都是用最少的东西表示所有需要的内容。

\section*{5．6．2 XXXX}
\section*{5.7 分类讨论}
\section*{5．7．1 分类讨论的时机}
分类讨论是非常重要的数学思想。我们在学习绝对值的时候，就开始接触分类讨论了。然而，在实际运用中，很多同学对分类的时机，标准和动机把握都不够准确。这一章我们主要以计数原理（排列组合）的题目为例，来讲清楚个中问题。

如果用一句话来概括分类的时机，就是：当遇到了无法确定的因素，使推理无法进行下去时，就要分类讨论。我们以 2023 年新高考 1 卷的题目为例：

例 5.17 （2023 年新高考 1 卷）某学校开设了 4 门体育类选修课和 4 门艺术类选修课，学生需要从这 8 门课中选修 2 门或 3 门课，并且每类课至少选修 1 门，则不同的选课方案共 \(\qquad\)有种。

显然， 2 门课和 3 门课是完全不一样的情况。想要进行下一步的分类讨论，我们必须要确定到底是 2 门课还是 3 门课。这就是分类讨论的动机。\\
（1）若选择 2 门课，则 1 门体育， 1 门艺术。根据乘法原理，共有 \(\mathrm{C}_{4}^{1} \cdot \mathrm{C}_{4}^{1}=16\) 种不同的选法；\\
（2）若选择 3 门课，则 1 门体育， 2 门艺术或者 2 门体育， 1 门艺术。这里仍需要分类讨论：\\
（1）若 1 门体育， 2 门艺术，根据乘法原理，共有 \(C_{4}^{1} \cdot C_{4}^{2}=24\) 种；\\
（2）若 2 门体育， 1 门艺术，根据乘法原理，共有 \(\mathrm{C}_{4}^{2} \cdot \mathrm{C}_{4}^{1}=24\) 种。\\
根据加法原理，共有 \(16+24+24=64\) 种。\\
题目本身不难，但是直接直观地体现了分类讨论的时机。分类讨论后，我们需要应用相应的计数原理求方法数。从名字上就可以看出什么时候需要分类讨论了：分类加法原理，分步乘法原理。所有我们分类讨论后得到的不同种数，都要用加法原理相加；所有使用加法原理的地方，本质上就是在分类讨论。所以，分类讨论也并非什么＂高超＂的技巧，只是在你进行逻辑推理，面对＂分岔路＂的情况时，把每个分岔路走一遍的过程。

\section*{5．7．2 原则和目标}
＂不重不漏＂是分类讨论的目标，既不能出现一个情况算两遍，也不能丢掉任何一种情况。换句话说，在我们沿着逻辑路线走的时候，既要保证每个分岔口都走下去看看，保证每个都被算过（不漏），也要保证每个分岔口的结果不能被计算两遍（不重）。那我们怎么确保这一点呢？这就是分类讨论时要遵守一个最基本的原则，就是按＂序＂讨论。最常见的序有 2 种：\\
（1）大小顺序。例如，如果可以选择 \(1^{\sim} 3\) 门课程，那么我们从小到大讨论 1 门， 2 门， 3 门的情况。切忌跳跃性地讨论，一会看 2 门的情况，还没看完就去算 3 门\\
的情况，想到哪是哪，这样肯定会出现重复或者漏解的情况。在 2023 年新高考 1 卷中，我们就行了两次大小顺序的分类讨论：课程的数量和体育课，美术课的数量。\\
（2）方向顺序。例如，照相排队有从左到右的顺序。甲有特殊的照相要求（不能两头，不能和某人相邻）时，可以按着甲的位置从左到右讨论。\\
基本上，绝大多数的题目，只要按着这两个序中的其中一种分类，就可以保证不重不漏了。\\
我们再举两个经典的例题，来大家体会这两个顺序和不重不漏的原则。

例 5.18 三个正整数的和为 8 ，不考虑顺序的情况下，有多少种可能？

注意，这个和将 8 个篮球分给三个同学，每个人至少有一个不一样：三个同学不同，有区别；但是这里三个正整数没有顺序（即 134 和 341 是一种情况）。所以，不能直接使用插板法。\\
（1）按最大数来分类。设最大数为 \(x\) ，另外两个数为 \(y, z\) 。由于三个数都是正整数，所以 \(x \leq 8-1-1=6\) ．又因为 \(x\) 是最大数，所以

\[
8=x+y+z \leq 3 x
\]

故 \(x \geq 3\) ．所以 \(x=3,4,5,6\) 共四种情况\\
（1）当 \(x=6\) 时，另外两个数只能为 1 ，故 1 种情况；\\
（2）当 \(x=5\) 时，另外两个数和为 3 ．而 \(3=1+2\) ，故 1 种情况；\\
（3）当 \(x=4\) 时，另外两个数和为 4 ．而 \(4=1+3=2+2\) ，故 2 种情况；\\
（4）当 \(x=3\) 时，另外两个数和为 5 。虽然 \(5=1+4=2+3\) ，但是 \(1+4\) 的情况再（3）中已经被计算过了， \(1,3,4\)三个数的最大数为 4 ．所以只有 \(8=2+3+3\) 这 1 种情况。\\
根据加法原理，共有 \(1+1+2+1=5\) 种。\\
（2）按最小数来分类。设最小数为 \(y\) 。由于均为正整数，所以 \(y \geq 1\) ；又因为 \(y\) 是最小数，所以

\[
8=x+y+z \geq 3 y
\]

故 \(y \leq 2\) ．所以 \(y=1,2\) 共两种情况。\\
（1）当 \(y=1\) 时，另外两个数和为 7 ，而 \(7=1+6=2+5=3+4\) ，共 3 种可能；\\
（2）当 \(y=2\) 时，另外两个数和为 6 ，且都不小于 2 ，所以 \(6=2+4=3+3\) ，共 3 种可能。根据加法原理，共有 \(2+3=5\) 种。

为什么一定要按＂最大数＂，＂最小数＂来分类呢？能不能只找其中一个数来分类呢？答案是不行。这是因为：由于三个数没有顺序，所以＂只找一个数＂会出现重复的情况。＂其中一个数是 1 ＂，并没有指定具体是哪个数，相当于＂这三个数里有 1 ＂。那和其他的情况＂其中一个数是 \(2^{\prime \prime}\) 之间，有重合，违背不重不漏的原则。

在这个题目中，为什么考虑最小数比考虑最大数的分类种数要少很多呢？事实上，同样的题目，将 8 换成 \(n\) ，那么最大数的满足 \(\frac{n}{3} \leq x \leq n-2\) ，而最小数 \(y\) 满足 \(1 \leq y \leq \frac{3}{n}\) 。所以 \(x\) 的范围差不多是 \(\frac{2}{3} n\) 这个量级，而 \(y\) 的范围差不多是 \(\frac{1}{3} n\) 这个量级，后者是前者的一半。所以，当有多个分类标准时，也应适当考虑分类的数量。

例 5.19 六人排队拍照，甲不在最左边，且和乙不相邻，有多少种不同的排队方法？

设从左到右的位置编号为 \(1^{\sim} 6\) 。甲不在两端，所以甲只能在 \(2,3,4,5,6\) 号位。这几个位置中，唯一和其他位置不一样的是 6 号位。因为只有 5 号位和 6 号位相邻，其他的位置都有两个位置相邻。所以我们可以按这个标准分类。\\
（1）若甲在 \(2,3,4,5\) 号位中的一个，那么乙只能去剩下位置中的 3 个。如，甲若在 2 号位，那么乙只能去 \(4,5,6\)号位。但无论是 \(2,3,4,5\) 中的哪一个，乙都有 3 种可能。当甲，乙确定后，剩下四个人有 \(A_{4}^{4}\) 种排法，所以总数为 \(4 \times 3 \times A_{4}^{4}=288\) 种；\\
（2）若甲在 6 号位，那么乙能去剩下位置中的 4 个。剩下的四个人有 \(A_{4}^{4}\) 种排法，所以总数为 \(4 \times A_{4}^{4}=96\) 种。故一共有 \(96+288=384\) 种。

这里我们使用了＂对称性＂：甲在 \(2,3,4,5\) 号位是等价的。所以，这四种情况对应的方法种数一样，我们不用四种都列出来再相加，只要算其中一种，最后乘 4 即可。这也是 \(4 \times 3 \times A_{4}^{4}\) 中 4 的含义。因此，对称性可以有效帮助我们减少分类讨论的种数，将部分加法原理转化为乘法原理。

如果分类讨论的逻辑比较复杂，我们可以借助树状图来分析讨论的情况。检查的时候，只要检查每一个分叉点，我们是否考虑了所有情况，且不同的情况之间是否有重叠。没有的话，就是一次有效的分岔，只要遍历所有分岔的种数，就是最后总的方法数。

\section*{5.7 .3 染色问题}
染色问题经常需要分类讨论。但是分类的标准有时候模糊不清，大部分答案只给出了最后分类的结果，但是分类的动机和原理并没有讲清楚。所以我们以两道例题为例，来看看染色问题的分类应该如何进行。

例5．20（2008 高考）如图，一环形花坛分成 \(A, B, C, D\) 四块，现有 4 种不同的花供选种，要求在每块里种 1 种花，且相邻的 2 块种不同的花，则不同的种法总数为\\
A， 96\\
B ， 84\\
C， 60\\
D， 48\\
\includegraphics[max width=\textwidth, center]{2025_02_16_7c9aaf92bc99dc07f048g-174}

我们可以首先根据选用的花的种数来分类。这是按多少，大小分类。\\
（1）若只使用 2 种花，那么必然是 \(A, C\) 同色，\(B, D\) 同色（且和 \(A, C\) 不同色）。根据分步乘法原理，共有 \(4 \times 3=12\)种；\\
（2）若使用三种花，一定是＂要不然 \(A, C\) 同色，要不然 \(B, D\) 同色＂。这是因为：若 \(A, C\) 不同色，那么 \(B, D\) 一定是第三种颜色，必然同色。所以我们只要考虑 \(A, C\) 不同色，\(B, D\) 同色的情况，然后将种数乘 2 即可。此时，\(A\) 有四种染色方法，\(C\) 有三种，\(B, D\) 有两种，即 \(4 \times 3 \times 2=24\) 种，故有 \(24 \times 2=48\) 种；\\
（3）若使用四种花，那么必然四部分的颜色各不相同。那么共有 \(A_{4}^{4}=24\) 种。\\
根据加法原理，合计 \(12+48+24=84\) 种。\\
虽然这个方法得到了答案，但是它并不自然，也并不普适。这题的图形结果简单，所以当我们用三种颜色填涂时，同色关系是确定的。当遇到更复杂的结构，比如五个花围，六个花諫时，只是去对颜色分类，并没有解决问题。同时，如果有更多的颜色供选择（如 8 个颜色），这个分类将会非常复杂。

这个方法不够自然则在于：我们拿到一个染色问题，＂直觉＂上是先从某一块区域开始，染染看，看能不能直接通过乘法原理得到结果。如果不行，我们再看是不是需要分类讨论。例如，对于排成一条直线的 4 块花東，如果用 4 种颜色染色，只要求相邻不同，那么\\
（1）从第一块开始染色即可，共有 4 种方法；\\
（2）第二块不能和第一块同色，所以 3 种；\\
（3）第三块不能和第二块同色，所以 3 种；\\
（4）第四块不能和第三块同色，所以 3 种。\\
因此，共有 \(4 \times 3 \times 3 \times 3=108\) 种。这是直接使用分步乘法原理的方法，也是最自然的方法。\\
那这个方法，能不能适用于环形的花東中呢？我们来简单尝试一下。\\
（1）从 \(A\) 开始涂色，那么共有 4 种方法；\\
（2）接着涂 \(B\) ，由于颜色和 \(A\) 不同，所以有 3 种方法；\\
（3）接着涂 \(C\) ，由于颜色和 \(B\) 不同，所以有 3 种方法；\\
（4）接着涂 \(D\) 。问题来了：它的颜色和 \(A, C\) 都不相同！但是 \(A, C\) 并不相邻，颜色可能相同，可能不同，需要分类讨论。这个例子，就非常好地说明了分类讨论的动机：由于推理无法进行（不知道颜色是否相同），所以要对两种情况分别讨论。由于 \(A, C\) 颜色是否相同，也直接影响了（1）～（3）的涂色，所以我们重新使用乘法原理：\\
（i）假设 \(A, C\) 颜色相同。那么 \(A, C\) 共有 4 种颜色可以选择；由于 \(B, D\) 都和 \(A, C\) 不同，所以各有 3 种可能。根据乘法原理，共有 \(4 \times 3 \times 3=36\) 种可能；\\
（ii）假设 \(A, C\) 颜色不同。那么 \(A, C\) 共有 \(4 \times 3=12\) 种颜色可以选择；那么 \(B, D\) 都只有 2 种选择。根据乘法原理，共有 \(12 \times 2 \times 2=48\) 种可能。\\
根据加法原理，共有 \(36+48=84\) 种不同的方法。\\
所以，我们在答案中看到的，要对第一块和倒数第二块的颜色进行分类讨论，道理在：最后一块和第一块，倒数第二块的颜色不同。所以第一块和倒数第二块的颜色是否相同，决定了最后一块有多少种选择。

在这个题目的基础上，我们看一个更复杂的问题：

例 5.21 用四种不同的颜色给如图所示的六块区域 \(A, B, C, D, E, F\) 涂色，要求相邻区域涂不同颜色，则涂色方法的总数是（）\\
A， 120\\
B， 72\\
C，48\\
D， 24\\
\includegraphics[max width=\textwidth, center]{2025_02_16_7c9aaf92bc99dc07f048g-175}

虽然整个图形较为复杂，但是基本形状仍是一个环形，但是有五个部分，和中间的部分。中间的部分是特殊的部分，它占用的颜色其他部分都不能占用。所以，我们可以先给中间涂色，再按上一题的方法给圆环涂色。\\
（1）\(E\) 有 4 种选择；\\
（2）剩下的五个部分，可以用三种颜色。仿照上题，我们从 \(A\) 开始，并对 \(B C\) 是否同色进行分类讨论。\\
（1）若 \(B, C\) 同色，共 3 种可能。那么 \(A\) 有 2 种可能，剩下的 \(D, F\) 不能和 \(B C\) 同色，不能和 \(E\) 同色，所以也只有两种颜色可选。同时，彼此还不能同色，所以无论 \(D\) 选哪种，\(F\) 只能选剩下的最后一种，所以共 2 种可能。这一类共 \(3 \times 2 \times 2=12\) 种；\\
（2）若 \(B, C\) 不同色，共 \(3 \times 2=6\) 种可能。那么 \(A\) 只有 1 种可能。此时 \(D\) 有两种颜色可选，一种和 \(C\) 同色，一种是和 \(B, E, C\) 都不同色的最后一种。 \(D\) 是否和 \(C\) 同色，会影响 \(F\) 的选择，所以也要分类讨论：\\
（i）若 \(D, C\) 同色，那么 \(F\) 有两种选择；\\
（ii）若 \(D, C\) 不同色，那么 \(F\) 有一种选择。\\
故 \(F\) 有 3 种选择。所以这类里共有 \(6 \times 1 \times(2+1)=18\) 种。\\
根据加法原理，圆环共有 \(12+18=30\) 种不同的涂法。那么共有 \(4 \times 30=120\) 种不同的方法。\\
这里，我们把整个涂色步骤看成了＂\(E\)＂和＂圆环＂两个部分。也可以把它们都在 \(B, C\) 同色以及 \(B, C\) 不同色两类里考虑：根据上面的讨论，\\
（1）\(B, C\) 同色，那么 \(E\) 有 4 种，圆环有 12 种，所以 \(4 \times 12=48\) 种；\\
（2）\(B, C\) 不同色，那么 \(E\) 有 4 种，圆环有 18 种，所以 \(4 \times 18=72\) 种。\\
根据加法原理， \(48+72=120\) 种。两种计算路径的关系，就是乘法分配律的等式两边：

\[
4 \times(12+18)=4 \times 12+4 \times 18
\]

虽然两种分类方法都得到了答案，但是第二种分类方法具有普适性。假设用 3 种颜色给首尾相连（每一块只和其中的两块相邻）的 \(n\) 块区域染色，要求相邻不同色。根据刚才的讨论，当 \(n=5\) 时，共有 30 种染色方法。对于一般的 \(n\) ，答案是多少呢？

设 \(a_{n}\) 为 \(n\) 块首尾相连的区域染色的方法种数，那么 \(a_{5}=30\) ．同样的方法，可以得到：\\
（1）当 \(n=3\) 时，共 \(\mathrm{A}_{3}^{3}=6\) 种颜色；\\
（2）当 \(n=4\) 时，从 \(A\) 开始填涂，考虑 \(B, D\) 是否同色。\\
（i）若 \(B, D\) 同色，\(A\) 有三种方法，\(B, D\) 有两种方法，\(C\) 有两种方法，根据乘法原理，共 \(3 \times 2 \times 2=12\) 种；\\
（ii）若 \(B, D\) 不同色，\(A\) 有三种方法，\(B, D\) 有 \(A_{2}^{2}=2\) 种方法，\(C\) 有一种方法，根据乘法原理，共 \(3 \times 2 \times 1=6\) 种。因此，共 \(12+6=18\) 种方法；\\
（3）对于一般的 \(n\) ，从 \(A_{1}\) 开始染色，最后一块是 \(A_{n}\) 。我们虽然无法具体分析每一块的涂法数量，但是可以写出地

推关系。\\
（i）若 \(A_{2}, A_{n}\) 同色，给这个圆环染色，等价于：给一个 \((n-2)\) 块的圆环染色，然后将第一块臂开（变成 \(A_{2}, A_{n}\) ），中间插入 \(A_{1}\) 。由于 \(A_{1}\) 和 \(A_{n}, A_{2}\) 颜色不同，所以共有 \(2 \times a_{n-2}=2 a_{n-2}\) 种方法；\\
（ii）若 \(A_{2}, A_{n}\) 不同色，给这个圆环染色，等价于：给一个 \((n-1)\) 块的圆环染色，然后在第一块和最后一块之间插入新的一块。这一块的颜色是确定的：必须和 \(A_{n}, A_{2}\) 不同色。所以共有 \(a_{n-1}\) 种方法。根据加法原理，

\[
a_{n}=a_{n-1}+2 a_{n-2}
\]

我们用 \(n=5\) 进行验证：\(a_{5}=30, a_{4}=18, a_{3}=6\) ，确实成立 \(a_{5}=a_{4}+2 a_{3}\) ．你是否能写出这个数列的通项公式？

\section*{总结}
我们以排列组合为例，说明了分类讨论的关键点：在因为＂某一环节有若干种情况＂导致推理无法进行时，要进行分类讨论。所以，分类讨论只是一个＂自然＂的选择：我想接着推理，但是我推不下去了，只有分成若干类确定的情况，我才能继续推理，所以我分而治之。这种自然的感觉，其实是数学中非常常见的。高中阶段我们学习过的所有方法，定理，都有这种自然的感觉。如果你暂时还没有建立起来，说明这里可能有一些你没有完全搞清楚的地方。

计数问题要求不重不漏。为了达到这个标准，必须按＂序＂进行分类。常见的＂序＂，有大小和方向。必要时，我们也要指定大小和方向（比如讨论＂最大的数＂的大小等）。

\section*{5．7．4 XXXX}
\section*{5.8 充要与化归}
\section*{5．8．1 再论充要替换}
充要替换，指的是将研究的问题＂等价地＂转化为另外一个好研究的问题。例如，如果想要判断一个四边形是平行四边形，我们可以证明：＂一组对边平行且相等＂；但如果是在解析几何中，线段长度需要利用勾股定理计算，较为复杂，我们可以把问题转化为＂对角线互相平分＂。这样，只需验证两条对角线的中点重合，就可以证明平行四边形了。

为什么这个思想重要？因为在复杂的题目中，我们要处理很多个条件。只有＂充要替换＂，才可以保证我们＂精确＂地翻译每个条件。既不遗漏信息，也不认为增加题目中没有的条件。举一个简单的例子，来说明充要替换的意义。

考察二次方程 \(x^{2}-2 a x-a+2=0\) 。\\
（1）它有两个大于 0 的实根，求实数 \(a\) 的取值范围；\\
（2）它有两个大于 \(\frac{1}{2}\) 的实根，求实数 \(a\) 的取值范围。\\
很多同学在处理这类问题时，都是＂凭感觉＂列方程。处理（1）时，发现这是一个两根问题，我们自然地想到了＂韦达定理＂。所以：

\[
\left\{\begin{array}{l}
\Delta=4 a^{2}+4 a-8>0 \\
x_{1}+x_{2}=2 a>0 \\
x_{1} x_{2}=2-a>0
\end{array}\right.
\]

联立解得 \(1<a<2\) ．如法炮制，对于（2），有：

\[
\left\{\begin{array}{r}
\Delta=4 a^{2}+4 a-8>0 \\
x_{1}+x_{2}=2 a>1 \\
x_{1} x_{2}=2-a>\frac{1}{4}
\end{array}\right.
\]

联立解得 \(1<a<\frac{7}{4}\) 。同样的方法，结果（1）得到了正确答案，但是（2）却得到了错误答案。\\
然后自己再看答案，发现答案中，（2）列的式子不一样，但是不理解为什么要多列某个式子（或者修改某个式子），这个疑问就从高一留到了高三。事实上，这里（2）的解答，问题就在于没有进行＂充要替换＂。

在（1）中，我们将命题 \(p:\)＂该方程有两个正实根＂，转化成了第一组的三个式子（命题 \(q\) ）。事实上，命题 \(p\) 和命题 \(q\) 是等价的。这是因为：\\
（1）\(p \Rightarrow q\) ：当方程有两个正实根时，肯定有判别式 \(\Delta>0\) ，所以第一个式子成立；由于 \(x_{1}, x_{2}>0\) ，所以 \(x_{1}+x_{2}>0, x_{1} x_{2}>0\) ．因此，第二个，第三个也都成立，所以 \(p\) 是 \(q\) 的充分条件；\\
（2）\(q \Rightarrow p\) ：当这三个不等式成立时，\(\Delta>0\) ，所以方程有两个不同的实根 \(x_{1}, x_{2}\) ；根据韦达定理，

\[
x_{1}+x_{2}=2 a>0, x_{1} x_{2}=2-a>0 .
\]

这两个条件，可以推出 \(x_{1}>0, x_{2}>0\) 。这是因为：当 \(x_{1} x_{2}>0\) 时，意味着它们两个同号，或者 \(x_{1}, x_{2}>0\) ，或者 \(x_{1}, x_{2}<0\) ．然而，\(x_{1}+x_{2}>0\) ，所以只能是 \(x_{1}, x_{2}>0\) ．\\
因此，方程有两个不同的实根，\(p\) 也是 \(q\) 的必要条件。\\
因此，\(p, q\) 是充要的，解 \(q\) 中的不等式即可得到 \(a\) 的取值范围。\\
然而，同样的等价性，在（2）中却不成立了。我们记命题 \(p_{1}\) ：＂该方程有两个大于 \(\frac{1}{2}\) 的实根＂，命题 \(q_{1}\) 是第二组的三个不等式。\\
（1）\(p_{1} \Rightarrow q_{1}\) ：当方程的两根都大于 \(\frac{1}{2}\) 时，首先 \(\Delta>0\) ，其次 \(x_{1}+x_{2}>1 x_{1} x_{2}>\frac{1}{4}\) 均成立，所以 \(q_{1}\) 成立；\\
（2）然而，\(q_{1}\) 推不出 \(p_{1}\) 。换言之，当第二组三个不等式成立时，无法保证两个根都大于 \(\frac{1}{2}\) 。这是为什么呢？\(\Delta>0\)毋庸置疑，问题出在了后两个不等式上。当实数 \(x_{1}, x_{2}\) 满足 \(x_{1}+x_{2}>1, x_{1} x_{2}>\frac{1}{4}\) 时，能否推出 \(x_{1}, x_{2}\) 都大于 \(\frac{1}{2}\)呢？不可以！举一个简单的反例：当 \(x_{1}=1, x_{2}=\frac{1}{3}\) 时，

\[
x_{1}+x_{2}=\frac{4}{3}>1, x_{1} x_{2}=\frac{1}{3}>\frac{1}{4}
\]

这两个不等式都成立。但是 \(x_{2}=\frac{1}{3}<\frac{1}{2}\) ．所以，这不是一个充要的转化！\\
事实上，\(q_{1}\) 只是 \(p_{1}\) 的必要条件。所以这样解得的范围，包含了正确的范围。那么，如何得到正确的不等式组呢？我们可以借鉴（1）的思路：

\[
x_{1}>\frac{1}{2}, x_{2}>\frac{1}{2} . \Leftrightarrow x_{1}-\frac{1}{2}>0, x_{2}-\frac{1}{2}>0
\]

令 \(t_{1}=x_{1}-\frac{1}{2}, t_{2}=x_{2}-\frac{1}{2}\) ，只要保证 \(t_{1}, t_{2}\) 是正数即可。这等价于

\[
t_{1}+t_{2}>0, t_{1} t_{2}>0
\]

所以，

\[
x_{1}>\frac{1}{2}, x_{2}>\frac{1}{2} . \Leftrightarrow\left\{\begin{array} { c } 
{ t _ { 1 } + t _ { 2 } > 0 } \\
{ t _ { 1 } t _ { 2 } > 0 }
\end{array} \Leftrightarrow \left\{\begin{array}{r}
x_{1}+x_{2}>1 \\
\left(x_{1}-\frac{1}{2}\right)\left(x_{2}-\frac{1}{2}\right)>0
\end{array}\right.\right.
\]

再加上 \(\Delta>0\) ，这样得到的不等式组，才是和 \(p_{1}\) 等价的。这完成了充要替换。\\
为什么我们还是希望写成 \(x_{1}, x_{2}\) 的对称式呢？根据第二章与第四章的内容，\(\left(x_{1}-\frac{1}{2}\right)\left(x_{2}-\frac{1}{2}\right)\) 是一个对称式，

可以用 \(x_{1}+x_{2}, x_{1} x_{2}\) 表示。这样韦达定理就可以派上用场，整个不等式可以化为关于参数 \(a\) 的不等式。\\
命题 \(p_{1}\) 等价于：

\[
\left\{\begin{array}{c}
\Delta=4 a^{2}+4 a-8>0 \\
x_{1}+x_{2}=2 a>1 \\
\left(x_{1}-\frac{1}{2}\right)\left(x_{2}-\frac{1}{2}\right)=x_{1} x_{2}-\frac{1}{2}\left(x_{1}+x_{2}\right)+\frac{1}{4}=\frac{9}{4}-2 a>0
\end{array}\right.
\]

解得 \(1<a<\frac{9}{8}\) ．这才是正确的答案。\\
充要替换，不仅仅帮助我们得到了正确的答案，还给予了我们一种＂完全掌握＂的自信。当我们将题目条件做＂充要替换＂时，每一步都是非常扎实， \(100 \%\)\\
正确的。这样，做的过程中，就不需要我们担心条件翻译是否有漏洞，是否漏解，多解了。\\
接下来，我们将举若干个不同的例子，来帮助大家体会不同模块中充要替换的意义。

【例 1】第二章中，我们阐释了如何解根式方程。其中，经常会出现＂增根＂的情况。虽然只要最后统一检验增根即可规避错误，但是我们可以从＂充要替换＂的角度，来理解＂增根＂的含义。考虑方程

\[
x+1=\sqrt{x+7}
\]

两边平方，得

\[
(x+1)^{2}=x+7
\]

移项，分解因式，得

\[
\begin{gathered}
x^{2}+x-6=0 \\
(x+3)(x-2)=0
\end{gathered}
\]

故 \(x_{1}=-3, x_{2}=2\) ．代人检验：\\
（1）当 \(x=-3\) 时，\(-2 \neq 2\) ，舍；\\
（2）当 \(x=2\) 时， \(3=3\) ．成立。\\
故方程的解为 \(x=2\) ．这里为什么产生增根了呢？就是因为将方程左右两侧平方，并不是一个充要的转化。我们记命题 \(p: " x+1=\sqrt{x+7} "\) ，命题 \(q: "(x+1)^{2}=x+7 "\) ．事实上，命题 \(q\) 只是命题 \(p\) 的必要条件，而非充分条件：

\[
(x+1)^{2}=x+7 \Leftrightarrow x+1=\sqrt{x+7} \text { 或 } x+1=-\sqrt{x+7} \text {. }
\]

命题 \(q\) 对应着两种情况，命题 \(p\) 只对应着其中的一中。当我们平方后，得到的两个解，分别是这两个方程的根。前者是 \(p\) 的解，后者是增根。

如果我们写得再＂紧凑＂些，这个逻辑就是：\(a^{2}=b^{2}\) 只是 \(a=b\) 的必要条件，而非充分条件，因为

\[
a^{2}=b^{2} \Leftrightarrow a=b \text { 或 } a=-b \text {. }
\]

后面的命题中，前者是方程的解，后者是增根。

【例 2】在导数的题目中，很多同学喜欢使用＂放缩＂的方法。这个方法本身没有问题，但是要意识到：这不是一个＂充要替换＂的方法，而是证明了一个更强的结论。例如，为了证明 \(e^{x} \geq \ln x+2\) ，可以使用两次不等式：\\
（1）由于 \(e^{x} \geq x+1\) ，所以只需证明 \(x+1 \geq \ln x+2\) 即可。这等价于 \(x \geq \ln x+1\)\\
（2）由于 \(x \geq \ln x+1\) ，命题成立。

写成一串等式，相当于：

\[
e^{x} \geq x+1 \geq \ln x+2
\]

完成了证明。这个方法本身没有问题，但是从几何上看，是我们在两个函数图像之间，插入了一个新的函数图像，＂刚好＂介于它们之间，从而完成了证明。换言之，我们证明的命题 \(p: A \geq B, B \geq C\) ，而题目要求证明的是命题 \(q: A \geq C\) ．前者是后者的充分条件，但不是必要条件。这是放缩法的本质。

这决定了，放缩法本身是非常容易失败的。同样是 \(e^{x} \geq x+1\) ，我们只要稍\\
微调整，就不再奏效。若要证明 \(e^{x} \geq e^{\frac{x}{2}}-1\) ，若仍使用这个方法，就是希望证明：

\[
e^{x} \geq x+1 \geq e^{\frac{x}{2}}-1
\]

然而，第二个不等式是错误的：令 \(f(x)=e^{\frac{x}{2}}-x-2\) ．它等价于证明 \(f(x) \leq 0\) ．\\
然而 \(f(4)=e^{2}-6>0\) ．所以这个方法失效了。原因是：我们强行将 \(x+1\) 插人到这两个函数之间，但它和 \(e^{\frac{x}{2}}-1\)互有胜负（对不同的 \(x\) ，有的 \(x+1\) 大，有的 \(e^{\frac{x}{2}}-1\) 大），\(x+1 \geq e^{\frac{x}{2}}-1\) 不成立。

事实上，一个简单的均值不等式就可以完成证明：\(e^{x}+1 \geq 2 \sqrt{e^{x}}=e^{\frac{x}{2}}\) ．所以 \(e^{x} \geq e^{\frac{x}{2}}-1\) ．\\
到这里，有的同学会说：那我只要也学会用均值不等式，就可以了呀！这没问题，多学几个不等式自然不错，但是不等式有无数多个，全都记住并不现实。另外，我们也很难保证，考场上看到题目，马上反应出＂正确＂的不等式。最终往往都会变成，拿着一大串钥匙，一把把去试；而且还不知道，这里面有没有正确的钥匙，只能＂赌＂。

所以，强烈推荐大家，把大量导数相关的不等式的学习，放到次要的位置上。因为它是不充要的转化，成功率低；相反，求导，分析单调性，画图象的方法，应该首先掌握，因为它是充要的方法。它可以帮助你掌握几乎所有你需要的函数的信息，要什么，有什么。这个方法掌握好后，再去熟悉必要的，常见的不等式即可

【例 3】圆雉曲线的小题中，也经常需要做充要替换，翻译几何条件。我们以 2024 届福建省检的圆曲小题为例：\\
例5．22 已知双曲线 \(\frac{x^{2}}{a^{2}}-\frac{y^{2}}{b^{2}}=1(a>0, b>0)\) 的左焦点为 \(F\) ，过 \(F\) 的直线 \(l\) 交圆 \(x^{2}+y^{2}=a^{2}\) 于 \(A, B\) 两点，交 \(C\)的右支于 \(P\) ．若 \(|A F|=|B P|,|P F|=2|A B|\) ，则 \(C\) 的离心率为

我们首先分析题目条件。所有条件都是齐次条件，所求是离心率，所以我们只关心图形的形状。这里的关键是如何将 \(|A F|=|B P|\) 这个条件进行充要转换。由于 \(A, B\) 都在圆 \(O\) 上，所以 \(O A=O B, \angle O A B=\angle O B A\) ．这个对称性让我们马上注意到，\(F O\) 应该也等于 \(P O\) ．这可以由全等得到：在三角形 \(F A O\) 和三角形 \(P B O\) 中，

\[
\left\{\begin{aligned}
O A & =O B \\
\angle F A O & =\angle P B O \\
A F & =B P
\end{aligned}\right.
\]

所以 \(F O=P O\) ．取双曲线的右焦点 \(F_{2}\) ，那么

\[
F_{2} O=F O=P O
\]

所以 \(P\) 在以 \(F F_{2}\) 为直径的圆上，所以 \(\angle F P F_{2}=90^{\circ}\) 。因此，第一个条件帮助我们确定了 \(P\) 点的位置。它告诉我们，虽然图和圆 \(x^{2}+y^{2}=a^{2}\) 有关，但是 \(P\) 点是以 \(F F_{2}\) 为直径的圆和双曲线的交点，它是＂确定＂的。这里确定的含义是指：由双曲线本身的参数决定。换言之，每个双曲线，都有这样的一个（其实是对称的四个，不妨考虑其中一个）点 \(P\) ：它在双曲线上，且 \(\angle F_{1} P F_{2}=90^{\circ}\) 。

题目所求的双曲线，它刚好满足第二个齐次条件：\(|P F|=2|A B|\) ．可以这么理解：对于形状（离心率）的双曲线，这样作出来的 \(|P F|,|A B|\) 的比值是不一样的。我们要找的是，比值刚好是 2 的那个。所以，我们将利用第二个条件列方程，确定离心率。

这里可以考虑不同的策略：\\
（1）由于 \(P\) 点确定，可以求出它的坐标（用 \(a, b\) 表示），然后也可以写出 \(A, B\) 的坐标，利用它们在圆上列方程。但这会出现根号，甚至是根号和有理式混合，所以不优先考虑；\\
（2）利用圆的对称性。过 \(O\) 作 \(A B\) 垂线，垂足为 \(D\) ．由垂径定理，\(A D=B D\) ．我们设 \(A D=B D=y\) ，那么 \(A B=\) \(2 y, P F=2 A B=4 y\) 。故 \(F A=B P=y\) 。为了避免根号，再设 \(O D=x\) 。由于 \(F A=B P\) ，所以 \(D\) 也是 \(P F\) 中点。所以 \(O D\) 是三角形 \(F P F_{2}\) 的中位线。故 \(P F_{2}=2 O D=2 x\) ．这样，所有需要的边长，都已确定。

显然第二种方法相对简单。虽然方程较多，但是可以最大程度避免根号，计算量小。现在共有 \(x, y, a, c\)（或者 \(x, y, a, b)\) 四个未知数。根据第二章与第四章中自由度和齐次的理论，我们想要三个齐次的方程。这三个分别是：\\
（1）\(P\) 在双曲线上，所以 \(P F-P F_{2}=2 a\) ．即 \(4 y-2 x=2 a\) ，即 \(2 y-x=a\) ；\\
（2）\(|O P|=c\) 。根据勾股定理，\(c^{2}=x^{2}+4 y^{2}\) ．\\
（3）\(|O A|=a\) ．根据勾股定理，\(a^{2}=x^{2}+y^{2}\) ．\\
四个未知数，三个齐次方程，可以确定比例。相信很多同学这里会产生疑问：为什么不是使用 \(|P F|=2|A B|,|A F|=\) \(|P B|\) 来列方程呢？这是因为，这两个条件已经被用作表示所有边长，和确定 \(P\) 的位置。那如何保证，我们现在得到的方程，就是对题目条件充分必要的转化呢？从两个角度理解：\\
（1）从内蕴的结构上看，我们列出了给四个未知数，列出了三个无关的方程，可以确定离心率 \(e\) 了。因为它就是其中两个未知数 \(a, c\) 的比值。题目本身，就是需要我们确定它，所以我们没有人为增加和减少条件；\\
（2）从每一步的充要替换上看，题目条件等价于如下若干个表述：首先，\(|P B|=|A F|\) 和 \(|O A|=|O B|\) 帮我们确定了 \(P\) 的位置，它 \(P\) 是双曲线和 \(x^{2}+y^{2}=c^{2}\) 的一个交点，这是第一个条件的充要替换；第二个条件，可以有如下两种理解：\\
（i）\(A, B\) 是 \(P F\) 和圆 \(x^{2}+y^{2}=a^{2}\) 的交点，满足 \(|P F|=2|A B|\) ；\\
（ii）\(A, B\) 是 \(P F\) 的两个不是中点的四等分点，且 \(|O A|=|O B|\) ．\\
这两种表述是等价的。我们列方程时，采用的是（ii）的表述：先利用四等分点，用 \(x, y\) 表示所有边长，然后用 \(|O A|=|O B|=a\) 列方程。

因此，问题化归为解如下方程组：

\[
\left\{\begin{array}{c}
2 y-x=a \\
x^{2}+4 y^{2}=c^{2} \\
x^{2}+y^{2}=a^{2}
\end{array}\right.
\]

多个未知数，需要消元。显然第一个方程可以用来消元，因为没有平方。既可以选择消 \(x\) ，也可以选择消 \(a\) 。这里展示消 \(a\) 的做法：将 \(a=2 y-x\) 代人第三个方程，得

\[
\begin{array}{r}
x^{2}+y^{2}=x^{2}-4 x y+4 y^{2} \\
4 x y=3 y^{2}
\end{array}
\]

故 \(4 x=3 y\) ．不妨设 \(x=3, y=4\) ，那么 \(a=5, c=\sqrt{x^{2}+4 y^{2}}=\sqrt{73}\) ．所以 \(e=\frac{\sqrt{73}}{5}\) ．\\
从这个题中，我们看出：充要替换在对于复杂题目也很有帮助。必要时，要去检验每个条件的翻译和转化，是否是充要的。既可以单纯从逻辑上检查，可以从自由度等结构本身内蕴的特点出发检验。

\section*{5．8．2 提升注意力的方法}
对于大部分练习充分的同学来说，处理熟悉的问题不再是难点。这部分不是我们这一节的重点（如果这里你需要更多的帮助，可以试试看阿不的＂大合集＂）。我们的重点是：到底怎么能够在拿到新问题的时候，我们就＂反应＂过来，＂意识到＂，或者说＂注意到＂，它就是我们已经练了几百遍的老问题呢？如何能够突破题目的束缚，做一个＂注意力惊人＂的人呢？

在揭晓谜底之前，我们以＂数列求通项＂为例，来看看到底如何运用化归思想。很多同学在学这部分内容的时候非常痛苦，因为题型和方法实在是太多了：一会是等比，等差，一会是差比数列，待定系数，还有什么错位相减，裂项相消，分组求和，倒序相加等等。学完这一节，一定要意识到：所有的数列，都是＂披着狼皮的羊＂，看上去凶猛，难对付，其实是你在教科书上见过不知道多少次的温顺的绵羊（其实一共只有两只，等差和等比）。你要做的，就是扯下它们的狼皮。

我们以接下来 12 个常见的求通项为例，看每一个题目，到底如何化归为等差数列和等比数列。我们要知道，所有求通项公式的题目，都是从＂递推公式＂出发，去求＂通项公式＂。什么叫＂递推公式＂？就是揭示相邻若干项的关系的公式，它可以＂递进地推导＂出每一项。如果想要求通项，就要把相邻两项的关系（递推）积累起来，最终得到一般的公式（通项）。这就产生了一般的累加，累乘的方法：前者把相邻两项的差累计，得到 \(a_{n}-a_{1}\) ，进而得到 \(a_{n}\) ；后者把相邻两项的比累计，得到 \(\frac{a_{n}}{a_{1}}\) ，进而得到 \(a_{n}\) 。

这两种方法分别适用于哪些数列呢？在我们学习等差，等比数列，推导它们的通项公式时，我们知道：累加法可以推导等差数列的通项（相邻两项差为定值），累乘法则可以推导等比数列的通项（相邻两项比为定值）。因此，所有告诉我们相邻两项差的式子，或者形如 \(a_{n+1}-a_{n}=\ldots\) 的式子，都可以考虑累加；所有告诉我们相邻两项商的式子，或者形如 \(\frac{a_{n+1}}{a_{n}}=\ldots\) 的式子，都可以考虑累乘。我们来看几个简单的例子。

【例1】已知数列 \(\left\{a_{n}\right\}\) 满足 \(a_{1}=1, a_{n+1}-a_{n}=n\) ，求数列 \(\left\{a_{n}\right\}\) 的通项公式；\\
递推公式告诉我们相邻两项的差，可以考虑累加。累加的时候，最好直接从 \(a_{n}-a_{n-1}\) 开始累加，这样求出来的就是 \(a_{n}\) ．

\[
\begin{gathered}
a_{n}-a_{n-1}=n \\
a_{n-1}-a_{n-2}=n-1 \\
a_{n-2}-a_{n-3}=n-2 \\
\cdots \\
a_{2}-a_{1}=1
\end{gathered}
\]

累加得

\[
a_{n}-a_{1}=1+2+\ldots+n=\frac{n(n+1)}{2}
\]

故对所有 \(n \geq 2\) ，都有 \(a_{n}=a_{1}+\frac{n(n+1)}{2}=\frac{n^{2}}{2}+\frac{n}{2}\) ．所以通项公式为：

\[
a_{n}=\left\{\begin{array}{c}
1, n=1 \\
\frac{n^{2}}{2}+\frac{n}{2}+1, n \geq 2
\end{array}\right.
\]

【例 2 】已知数列 \(\left\{a_{n}\right\}\) 满足 \(a_{1}=1, \frac{a_{n+1}}{n+1}-\frac{a_{n}}{n}=2\) ，求数列 \(\left\{a_{n}\right\}\) 的通项公式；

递推公式告诉我们的不是＂相邻两项的差＂，但是＂\(\frac{a_{n+1}}{n+1}, \frac{a_{n}}{n} "\) 的差。它是新数列相邻两项的差。令 \(b_{n}=\frac{a_{n}}{n}\) ，那么

\[
b_{n+1}-b_{n}=2
\]

所以 \(\left\{b_{n}\right\}\) 是一个等差数列，\(b_{1}=\frac{1}{1}=1\) ，所以 \(b_{n}=2 n-1\) ．故

\[
a_{n}=n b_{n}=2 n^{2}-n
\]

因此，如果不是 \(a_{n}\) 的相邻差，我们可以看是不是＂新数列＂的相邻两项差。通过化为一个新数列进行操作。这里我们为什么没有累加呢？因为 \(\left\{b_{n}\right\}\) 是等差数列，不用再累加一遍求它的通项，可以直接用等差数列通项公式。因此，我们可以这么理解累加：它更像是一个基本思路，而不是一个技巧。我们要尽量把所有表达式，往 \(a_{n+1}-a_{n}=\ldots\) 靠近。这个 \(a_{n}\) 可以是原来的数列，可以是新的数列。根据右侧的表达式，再看是否需要累加，还是直接就是一个等差数列。我们给这样的数列一个更一般的名字，叫＂差数列＂，意思就是知道相邻两项差的数列。

【例 3】已知数列 \(\left\{a_{n}\right\}\) 满足 \(a_{1}=1,(n+1) a_{n+1}-n a_{n}=\left(\frac{1}{2}\right)^{n}\) ，求数列 \(\left\{a_{n}\right\}\) 的通项公式；这也是一个差数列。我们知道的不是 \(a_{n+1}, a_{n}\) 的差，而是 \(b_{n}=n a_{n}\) 的相邻两项差：

\[
b_{n+1}-b_{n}=\left(\frac{1}{2}\right)^{n}
\]

故

\[
\begin{aligned}
b_{n}-b_{n-1} & =\left(\frac{1}{2}\right)^{n-1} \\
b_{n-1}-b_{n-2} & =\left(\frac{1}{2}\right)^{n-2} \\
\ldots & \\
b_{2}-b_{1} & =\left(\frac{1}{2}\right)^{1}
\end{aligned}
\]

累加得

\[
b_{n}-b_{1}=1-\left(\frac{1}{2}\right)^{n-1}
\]

而 \(b_{1}=1 \cdot a_{1}=1\) ，所以当 \(n \geq 2\) 时，有

\[
b_{n}=2-\left(\frac{1}{2}\right)^{n-1}
\]

这对 \(n=1\) 也成立，所以 \(b_{n}=2-\left(\frac{1}{2}\right)^{n-1}\) ．

【例 4】已知数列 \(\left\{a_{n}\right\}\) 满足 \(a_{1}=1, a_{n+1}-2 a_{n}=1\) ，求数列 \(\left\{a_{n}\right\}\) 的通项公式；这个式子不再像前面的几个例子一样，直接是 \(a_{n+1}, a_{n}\) 的差，而是 \(a_{n+1}\) 和 \(2 a_{n}\) 的差。怎么办呢？这里就要我们尽可能去化归。这个式子，既不是等差数列的通项公式，也不是等比数列的通项公式。但是，差得都不多：\\
（1）等差数列：如果改为 \(a_{n+1}-a_{n}=1\) ，就是等差数列，可以直接累加。但是，这里也不是不能累加，需要做个小操作。累加的精髓，其实和裂项一样，可以做到前后相消。从 \(a_{n}-2 a_{n-1}=1\) 出发，为了将 \(2 a_{n}\) 消掉，应该给下一个式子乘2：

\[
\begin{gathered}
2\left(a_{n-1}-2 a_{n-2}\right)=2 \\
2 a_{n-1}-4 a_{n-2}=2 .
\end{gathered}
\]

这样，\(a_{n-1}\) 就顺利消掉了：

\[
a_{n}-4 a_{n-2}=3
\]

如何消掉 \(a_{n-1}\) 呢？给下一个式子乘 4 即可。依次类推：

\[
\begin{gathered}
a_{n}-2 a_{n-1}=1 \\
2 a_{n-1}-4 a_{n-2}=2 \\
4 a_{n-2}-8 a_{n-3}=4 \\
8 a_{n-3}-16 a_{n-4}=8
\end{gathered}
\]

最后一个式子是多少呢？我们来观察前面的几个式子。如果写成 \(2^{n}\) 的形式，其实是：

\[
\begin{gathered}
2^{0} a_{n}-2^{1} a_{n-1}=2^{0}, \\
2^{1} a_{n-1}-2^{2} a_{n-2}=2^{1}, \\
2^{2} a_{n-2}-2^{3} a_{n-3}=2^{2}, \\
2^{3} a_{n-3}-2^{4} a_{n-4}=2^{3}, \\
2^{4} a_{n-4}-2^{5} a_{n-5}=2^{4},
\end{gathered}
\]

我们发现，每一项的角标和 2 的指数，和都是 \(n\) ；同时，右侧常数和被减数系数相同。所以，最后一个应该是：

\[
2^{n-2} a_{2}-2^{n-1} a_{1}=2^{n-2}
\]

累加后，中间的项全都消掉，得到

\[
a_{n}-2^{n-1} a_{1}=1+2+\ldots+2^{n-2}=2^{n-1}-1
\]

由 \(a_{1}=1\) ，得

\[
a_{n}=2^{n-1}+2^{n-1}-1=2^{n}-1
\]

这对于所有 \(n \geq 2\) 都成立。同时，对 \(n=1\) 也成立。所以通项为 \(a_{n}=2^{n}-1\) 。\\
（2）等比数列。如果等式的右侧是 0 ，那它就是一个标准的等比数列。然而，这里等式的右侧不是 0 ，所以我们需要进行适当地调整，让它变成一个等比数列。我们将 \(a_{n}\) 挪到等式右边，然后同时 +1 ，那么

\[
a_{n+1}+1=2 a_{n}+2=2\left(a_{n}+1\right)
\]

所以数列 \(\left\{a_{n}+1\right\}\) 是一个等比数列。因此，

\[
a_{n}+1=2^{n-1}\left(a_{1}+1\right)=2^{n}
\]

故 \(a_{n}=2^{n}-1\) ．这个神奇的数 1 是怎么得到的呢？事实上，我们希望构造出一个新的数列 \(\left\{a_{n}+x\right\}\) ，使得它是等比数列。注意到 \(a_{n}\) 的系数是 2 ，所以我们希望

\[
a_{n+1}+x=2\left(a_{n}+x\right)
\]

展开得

\[
a_{n+1}=2 a_{n}+x
\]

和原来的通项公式比较，得 \(x=1\) ．关于这里更细致的讲解，在魔法书中有介绍，这里按下不表。\\
这两个方法的核心思想，都是化归。当我们注意到这个数列递推式和等差，等比数列的相似性时，可以考虑把它们向这两个数列转化。事实上，向等差数列（或者一般的差数列）转化，本质上就是用累加法；向等比数列\\
（或者一般的比数列）转化，本质上就是用累乘法。因此，求通项的题目绝大多数都可以化归成这两类。

【例 5】已知数列 \(\left\{a_{n}\right\}\) 满足 \(a_{1}=1, a_{n+1}-2 a_{n}=3^{n}\) ，求数列 \(\left\{a_{n}\right\}\) 的通项公式；\\
（1）累加法：和（4）类似的思路：

\[
\begin{aligned}
& a_{n}-2 a_{n-1}=3^{n-1} \\
& 2 a_{n-1}-4 a_{n-2}=2^{1} \cdot 3^{n-2} \\
& 4 a_{n-2}-8 a_{n-3}=2^{2} \cdot 3^{n-3} \\
& \cdots \\
& 2^{n-2} a_{2}-2^{n-1} a_{1}=2^{n-2} \cdot 3^{1}
\end{aligned}
\]

累加得

\[
a_{n}=2^{n-1}+2^{n-2} \cdot 3^{1}+2^{n-3} \cdot 3^{2}+\ldots+2 \cdot 3^{n-2}+3^{n-1}
\]

根据第二章或第四章中的内容，它刚好可以看作 \(a=2, b=3\) 时的

\[
a^{n-1}+a^{n-2} b+a^{n-3} b^{2}+\ldots+a^{2} b^{n-3}+a b^{n-2}+b^{n-1}=\frac{a^{n}-b^{n}}{a-b}=3^{n}-2^{n}
\]

故 \(a_{n}=3^{n}-2^{n}\) 。这对 \(n \geq 2\) 都成立。同时，它也对 \(n=1\) 成立，所以 \(a_{n}=3^{n}-2^{n}\) 。（2）化为差数列：我们发现 \(a_{n+1}, a_{n}\) 的系数差 2 ，能不能进行一个操作，变成 \(b_{n+1}-b_{n}\) 的形式呢？我们联想到＂等比数列＂：以 2 为公比的等比数列，满足 \(a_{n+1}-2 a_{n}=0\) ，右侧是 0 ，和这个式子有相似的地方。所以，我们两侧同时除以一个公比为 2 的等比数列， \(2^{n+1}\) ，得

\[
\frac{a_{n+1}}{2^{n+1}}-\frac{a_{n}}{2^{n}}=\frac{3^{n}}{2^{n+1}}
\]

这样的话，\(\frac{a_{n+1}}{2^{n+1}}, \frac{a_{n}}{2^{n}}\) 冈好好构成了 \(b_{n+1}, b_{n}\) 的形式。令 \(b_{n}=\frac{a_{n}}{2^{n}}\) ，那么

\[
\begin{aligned}
& b_{n+1}-b_{n}=\frac{3^{n}}{2^{n+1}} \\
& b_{n}-b_{n-1}=\frac{3^{n-1}}{2^{n}}
\end{aligned}
\]

累加得

\[
b_{n}-b_{1}=\frac{3^{n-1}}{2^{n}}+\frac{3^{n-2}}{2^{n-1}}+\ldots+\frac{3^{1}}{2^{2}}
\]

注意到 \(b_{1}=\frac{1}{2}\) ，所以

\[
b_{n}=\frac{3^{n-1}}{2^{n}}+\frac{3^{n-2}}{2^{n-1}}+\ldots+\frac{3^{1}}{2^{2}}+\frac{3^{0}}{2^{1}}
\]

右侧是一个等比数列。它可以写作：

\[
\frac{3^{n-1}}{2^{n}}+\frac{3^{n-2}}{2^{n-1}}+\ldots+\frac{3^{1}}{2^{2}}+\frac{3^{0}}{2^{1}}=\frac{1}{2}\left(1+\frac{3}{2}+\ldots+\left(\frac{3}{2}\right)^{n-1}\right)=\frac{1}{2} \cdot \frac{\left(\frac{3}{2}\right)^{n}-1}{\frac{3}{2}-1}=\left(\frac{3}{2}\right)^{n}-1
\]

所以

\[
a_{n}=2^{n} b_{n}=3^{n}-2^{n}
\]

这对 \(n \geq 2\) 都成立。经检验，对 \(n=1\) 也成立。所以 \(a_{n}=3^{n}-2^{n}(n \geq 1)\) ．\\
事实上，这个方法也可以用于例4．

【例 6】已知数列 \(\left\{a_{n}\right\}\) 满足 \(a_{1}=1, a_{n+1}-2 a_{n}=n\) ，求数列 \(\left\{a_{n}\right\}\) 的通项公式。

累加法和化为差数列的方法都是可行的。这里我们使用化为差数列的方法。两边同时除以 \(2^{n+1}\) ，得

\[
\frac{a_{n+1}}{2^{n+1}}-\frac{a_{n}}{2^{n}}=\frac{n}{2^{n+1}}
\]

令 \(a_{n}=\frac{b_{n}}{2^{n}}\) ，则

\[
b_{n}-b_{n-1}=\frac{n-1}{2^{n}}
\]

累加得

\[
b_{n}=\frac{n-1}{2^{n}}+\frac{n-2}{2^{n-1}}+\ldots+\frac{1}{2^{2}}+\frac{1}{2}
\]

这是一个差比数列（等差数列和等比数列的乘积），求和方法和等比数列一项，可以＂错位相减＂：

\[
2 b_{n}=1+\frac{1}{2^{1}}+\ldots+\frac{n-2}{2^{n-2}}+\frac{n-1}{2^{n-1}}
\]

将 \(2 b_{n}\) 和 \(b_{n}\) 同分母的对齐相减，得

\[
b_{n}=\frac{1}{2^{n-1}}+\frac{1}{2^{n-2}}+\ldots+\frac{1}{2^{2}}+\frac{1}{2}-\frac{n-1}{2^{n-1}}=1-\frac{1}{2^{n}}-\frac{n-1}{2^{n-1}}=1-\frac{2 n-1}{2^{n}}
\]

故

\[
a_{n}=2^{n} b_{n}=2^{n}-(2 n-1)=2^{n}-2 n+1
\]

这对 \(n \geq 2\) 都成立。经检验，对 \(n=1\) 也成立，所以 \(a_{n}=2^{n}-2 n+1\) ．

【例 7】已知数列 \(\left\{a_{n}\right\}\) 满足 \(a_{1}=1, n a_{n+1}-(n+1) a_{n}=2\) ，求数列 \(\left\{a_{n}\right\}\) 的通项公式。\\
对于这个式子，我们注意到 \(a_{n+1}, a_{n}\) 的系数也不相同。我们试图两边同时乘一个因子，使得等式左侧化为 \(b_{n+1}-b_{n}\)的形式。如果想要写成这样的形式，最终一定化为：

\[
g(n+1) a_{n+1}-g(n) a_{n}=\ldots
\]

的式子。其中 \(g(n)\) 是一个关于 \(n\) 的函数。这样，令 \(b_{n}=g(n) \cdot a_{n}\) 即可。所以，我们的目标就是要找到合适的 \(g(n)\) 。

由于我们两侧都是同时乘一个因子，所以 \(a_{n+1}, a_{n}\) 系数的比例是不改变的：即

\[
\frac{g(n+1)}{n}=\frac{g(n)}{n+1}
\]

可以这么理解这个式子：我们在原来的递推式左右两侧同时乘一个因子 \(f(n)\) ，那么

\[
n f(n) \cdot a_{n+1}-(n+1) f(n) \cdot a_{n}=2 f(n)
\]

和 \(g(n)\) 的表达式比较，有

\[
\begin{gathered}
g(n+1)=n f(n) \\
g(n)=(n+1) f(n)
\end{gathered}
\]

消掉 \(f(n)\) ，就有

\[
\frac{g(n+1)}{n}=\frac{g(n)}{n+1}
\]

什么样的 \(g(n)\) 满足条件呢？我们可以用累乘法求解，但事实上，它可以变形为

\[
(n+1) g(n+1)=n g(n)
\]

所以，\(n g(n)\) 是常数。不妨令 \(g(1)=1\) ，那么 \(g(n)=\frac{1}{n}\) 符合题意。所以，我们希望最终式子变形为：

\[
\frac{a_{n+1}}{n+1}-\frac{a_{n}}{n}=\ldots
\]

比较系数会发现，只要在递推式的左右两侧同时除以 \(n(n+1)\) 即可：

\[
\frac{a_{n+1}}{n+1}-\frac{a_{n}}{n}=\frac{2}{n(n+1)}
\]

右侧可以裂项，即

\[
\frac{a_{n+1}}{n+1}-\frac{a_{n}}{n}=\frac{2}{n}-\frac{2}{n+1}
\]

可以变形为

\[
\frac{a_{n+1}}{n+1}+\frac{2}{n+1}=\frac{a_{n}}{n}+\frac{2}{n}
\]

令 \(b_{n}=\frac{a_{n}}{n}+\frac{2}{n}\) ，那么 \(b_{n+1}=b_{n}\) ．所以 \(\left\{b_{n}\right\}\) 是一个常数列。当 \(n=1\) 时，\(b_{1}=3\) ．所以

\[
\frac{a_{n}}{n}+\frac{2}{n}=3
\]

即 \(a_{n}=3 n-2\) ．\\
很多同学在理解形如＂\(n a_{n+1}-(n+1) a_{n}=\ldots\)＂的递推式时，是完全通过记忆进行的。简单，常用的靠记忆可以解决，但复杂的就只能通过原理进行推导。所有这些变形的方法都是在左右两侧同时乘一个因子 \(f(n)\) ，使得最终可以化作 \(b_{n+1}-b_{n}\) 的形式。为什么要化成这样的形式？因为我们最终要使用累加法，只有 \(b_{n+1}-b_{n}\) 的形式才可以前后抵消。这个例子中我们没有直接使用累加法，因为我们观察到了裂项；但如果不裂项，也可以对

\[
\frac{a_{n+1}}{n+1}-\frac{a_{n}}{n}=\frac{2}{n(n+1)}
\]

直接累加，得到答案。\\
以上所有的例子，都是通过累加法得到的。使用累加法，只要完成两个步骤即可：\\
（1）构造适当的 \(\left\{b_{n}\right\}\) ，将递推公式化为 \(b_{n+1}-b_{n}=\ldots\) 的形式。这一步往往是代数变形要求较高的一步，但是它也有一般方法（例 7 中＂待定系数＂的方法）；\\
（2）累加。\\
类似地，我们来看累乘法如何应用。它相对简单，因为 \(a_{n+1}, a_{n}\) 的比例，往往不需要额外的构造，就可以得到答案。

【例 8】已知数列 \(\left\{a_{n}\right\}\) 满足 \(a_{1}=1, a_{n+1}=\left(1+\frac{1}{n}\right) a_{n}\) ，求数列 \(\left\{a_{n}\right\}\) 的通项公式；根据递推公式，有

\[
\begin{aligned}
& \frac{a_{n}}{a_{n-1}}=\frac{n}{n-1} \\
& \frac{a_{n-1}}{a_{n-2}}=\frac{n-1}{n-2}
\end{aligned}
\]

\[
\frac{a_{2}}{a_{1}}=\frac{2}{1}
\]

累乘得

\[
\frac{a_{n}}{a_{1}}=n
\]

故 \(a_{n}=n(n \geq 2)\) 。经检验，对 \(n=1\) 也成立，所以 \(a_{n}=n\) ，对 \(n \geq 1\) 都成立。事实上，不累乘可以直接看出＂定

值＂：

\[
\frac{a_{n+1}}{n+1}=\frac{a_{n}}{n}
\]

这意味着，数列 \(b_{n}=\frac{a_{n}}{n}\) 是常数列，因为它满足 \(b_{n+1}=b_{n}\) ．所以，\(b_{n}=b_{1}=1\) ．故 \(a_{n}=n\) ．这个情况在例 7 中也出现了，所以，当我们能明显看到＂常数列＂或＂定值＂时，可以直接构造对应的数列。

【例 9】已知数列 \(\left\{a_{n}\right\}\) 满足 \(a_{1}=1, a_{n+1}=2\left(1+\frac{1}{n}\right) a_{n}\) ，求数列 \(\left\{a_{n}\right\}\) 的通项公式；直接累乘：

\[
\begin{aligned}
& \frac{a_{n}}{a_{n-1}}=2 \cdot \frac{n}{n-1} \\
& \frac{a_{n-1}}{a_{n-2}}=2 \cdot \frac{n-1}{n-2} \\
& \ldots \\
& \frac{a_{2}}{a_{1}}=2 \cdot \frac{2}{1}
\end{aligned}
\]

故

\[
\frac{a_{n}}{a_{1}}=2^{n-1} \cdot n
\]

故 \(a_{n}=2^{n-1} n(n \geq 2)\) ．经检验，对 \(n=1\) 也成立，所以 \(a_{n}=2^{n-1} n\) ．如果参考例 9 的思路，可以发现：\(b_{n}=\frac{a_{n}}{n}\) 是一个等比数列，因为

\[
\frac{a_{n+1}}{n+1}=2 \cdot \frac{a_{n}}{n}
\]

即

\[
b_{n+1}=2 b_{n} .
\]

而 \(b_{1}=1\) ，所以 \(b_{n}=2^{n-1}\) ，故 \(a_{n}=2^{n-1} n\) ．\\
【例 10】已知数列 \(\left\{a_{n}\right\}\) 满足 \(a_{1}=1, \frac{a_{n+1}}{a_{n}}=\frac{n+2}{n}\) ，求数列 \(\left\{a_{n}\right\}\) 的通项公式；直接累乘，得

\[
\begin{gathered}
\frac{a_{n}}{a_{n-1}}=\frac{n+1}{n-1} \\
\frac{a_{n-1}}{a_{n-2}}=\frac{n}{n-2} \\
\ldots \\
\frac{a_{2}}{a_{1}}=\frac{3}{1}
\end{gathered}
\]

累乘得

\[
\frac{a_{n}}{a_{1}}=\frac{(n+1) \cdot n \cdot \ldots 3}{(n-1) \cdot(n-2) \cdot \ldots 1}=\frac{(n+1) n}{2}
\]

所以 \(a_{n}=\frac{(n+1) n}{2}\) ．这对 \(n \geq 2\) 都成立。经检验，对 \(n=1\) 也成立，所以 \(a_{n}=\frac{(n+1) n}{2}\) ．如果能直接观察到定值也可以：

\[
\frac{a_{n+1}}{a_{n}}=\frac{n+2}{n}=\frac{(n+2)(n+1)}{(n+1) n}
\]

所以

\[
\frac{a_{n+1}}{(n+2)(n+1)}=\frac{a_{n}}{(n+1) n}
\]

令 \(b_{n}=\frac{a_{n}}{(n+1) n}\) ．那么 \(b_{n+1}=b_{n}\) ．所以

\[
b_{n}=b_{1}=\frac{1}{2 \cdot 1}=\frac{1}{2}
\]

故 \(a_{n}=(n+1) n b_{n}=\frac{(n+1) n}{2}\) ．

【例 11】已知数列 \(\left\{a_{n}\right\}\) 满足 \(a_{1}=1, a_{2}=2, \frac{a_{n+2}}{a_{n}}=\frac{n}{n+4}\) ，求数列 \(\left\{a_{n}\right\}\) 的通项公式；\\
我们发现，递推公式不再是 \(a_{n+1}, a_{n}\) 的关系，而是 \(a_{n+2}, a_{n}\) 的关系。我们可以想象一下它们的＂决定关系＂：\\
（1）当 \(n=1\) 时，可以求 \(a_{3}\) ；当 \(n=3\) 时，可以求 \(a_{5} ; \cdots\) 因此，奇数项依次决定；\\
（2）当 \(n=2\) 时，可以求 \(a_{4}\) ；当 \(n=4\) 时，可以求 \(a_{6} ; \cdots\) 因此，偶数项依次决定。因此，这个数列本质上可以被看成两个数列：奇数项组成的数列和偶数项组成的数列。因为它们内部就完全决定了自己，彼此之间无关。\\
（1）先看奇数项。奇数项可以写成 \(a_{2 k+1}\) 的形式。令 \(b_{k}=a_{2 k-1}\) ，那么 \(b_{k}\) 其实就是 \(\left\{a_{n}\right\}\) 中第 \(k\) 个奇数项。为了求 \(b_{k}, b_{k-1}\) 的关系，我们让 \(2 k-1=n+2\) ，即 \(n=2 k-3\) 。这样代人到递推公式中，就会出现 \(b_{k}=a_{2 k-1}\) 了。即：

\[
\frac{a_{2 k-1}}{a_{2 k-3}}=\frac{2 k-3}{2 k+1}
\]

左侧刚好是

\[
\frac{b_{k}}{b_{k-1}}=\frac{2 k-3}{2 k+1}
\]

故

\[
\begin{gathered}
\frac{b_{k-1}}{b_{k-2}}=\frac{2 k-5}{2 k-1}, \\
\ldots \\
\frac{b_{2}}{b_{1}}=\frac{1}{5} .
\end{gathered}
\]

累乘得

\[
\frac{b_{k}}{b_{1}}=\frac{(2 k-3) \cdot(2 k-5) \cdot \ldots \cdot 1}{(2 k+1) \cdot(2 k-1) \cdot \ldots \cdot 5}=\frac{3}{(2 k+1)(2 k-1)} .
\]

注意到 \(b_{1}=a_{1}=1\) ，所以

\[
b_{k}=\frac{3}{(2 k+1)(2 k-1)}
\]

这对 \(k \geq 2\) 都成立。经检验，对 \(k=1\) 也成立，所以 \(b_{k}=\frac{3}{(2 k+1)(2 k-1)}\) 。我们需要将它写为 \(a_{n}\) 的形式。由于 \(a_{2 k-1}=b_{k}\) ，所以

\[
a_{2 k-1}=\frac{3}{(2 k+1)(2 k-1)}
\]

令 \(2 k-1=n\) 那么

\[
a_{n}=\frac{3}{n(n+2)}
\]

这是奇数列的情况；\\
（2）再看偶数项。令 \(c_{k}=a_{2 k}\) ，那么 \(c_{k}\) 就是 \(\left\{a_{n}\right\}\) 中的第 \(k\) 个偶数项。令 \(2 k=n+2\) ，那么 \(n=2 k-2\) ，代入递推公式，得

\[
\frac{a_{2 k}}{a_{2 k-2}}=\frac{2 k-2}{2 k+2}=\frac{k-1}{k+1}
\]

即

\[
\frac{c_{k}}{c_{k-1}}=\frac{k-1}{k+1}
\]

可以累乘，但也可以观察到＂常数列＂：

\[
\begin{aligned}
\frac{c_{k}}{c_{k-1}} & =\frac{k(k-1)}{(k+1) k} \\
(k+1) k c_{k} & =k(k-1) c_{k-1}
\end{aligned}
\]

这意味着，\(k(k+1) c_{k}\) 是定值。而 \(c_{1}=a_{2}=2\) ，所以

\[
\begin{gathered}
k(k+1) c_{k}=2 c_{1}=4 \\
c_{k}=\frac{4}{(k+1) k}
\end{gathered}
\]

即

\[
a_{2 k}=\frac{4}{(k+1) k}
\]

令 \(n=2 k\) ，那么

\[
a_{n}=\frac{4}{\frac{n}{2}\left(\frac{n}{2}+1\right)}=\frac{16}{n(n+2)}
\]

因此，数列 \(\left\{a_{n}\right\}\) 的通项公式为：

\[
a_{n}=\left\{\begin{array}{l}
\frac{3}{n(n+2)}, n \text { 为奇数, } \\
\frac{16}{n(n+2)}, n \text { 为偶数. }
\end{array}\right.
\]

【例 12】已知数列 \(\left\{a_{n}\right\}\) 满足 \(a_{1}=1, a_{n+1}=-a_{n}+n\) ，求数列 \(\left\{a_{n}\right\}\) 的通项公式；\\
【方法一】化为差数列。我们注意到 \(a_{n}, a_{n+1}\) 的系数差一个 -1 ，所以可以用 \((-1)^{n}\) 进行调整。在递推公式两侧同乘 \((-1)^{n+1}\) ，得

\[
(-1)^{n+1} a_{n+1}=(-1)^{n} a_{n}+(-1)^{n+1} n
\]

累加得

\[
(-1)^{n} a_{n}=(-1)^{1} a_{1}+(-1)^{n}(n-1)+(-1)^{n-1}(n-2)+\ldots(-1)^{2} \cdot 1
\]

后面的求和部分，仍可以视作差比数列求和。所以记 \(S=(-1)^{n}(n-1)+(-1)^{n-1}(n-2)+\ldots(-1)^{2} \cdot 1\) ，那么

\[
-S=(-1)^{n+1}(n-1)+(-1)^{n}(n-2)+\ldots+(-1)^{4} \cdot 2+(-1)^{3} \cdot 1
\]

和 \(S\) 错位相减，得

\[
\begin{gathered}
2 S=(-1)^{n+2}(n-1)+(-1)^{n}+(-1)^{n-1}+\ldots+(-1)^{3}+(-1)^{2} \\
=(n-1)(-1)^{n}+\frac{(-1)^{n}+1}{2}
\end{gathered}
\]

故 \(S=\frac{1}{4}+\frac{2 n-1}{4} \cdot(-1)^{n}\) ．所以

\[
(-1)^{n} a_{n}=-1+S=\frac{2 n-1}{4} \cdot(-1)^{n}-\frac{3}{4}
\]

即

\[
a_{n}=\frac{2 n-1}{4}+\frac{3}{4} \cdot(-1)^{n-1}
\]

这对 \(n \geq 2\) 成立。经检验，\(n=1\) 也成立。所以 \(a_{n}=\frac{2 n-1}{4}+\frac{3}{4} \cdot(-1)^{n-1}\) 。\\
【方法二】注意到 \(a_{n}\) 的系数是 -1 ，我们可以写作

\[
a_{n+1}+a_{n}=n
\]

这意味着，相邻两项的和是有关系的。这样一来，我们可以通过思想试验来考察通项公式。\\
（1）当 \(n=1\) 时，\(a_{1}+a_{2}=1\) ．所以 \(a_{2}=0\) ；\\
（2）当 \(n=2\) 时，\(a_{2}+a_{3}=2\) ．所以 \(a_{3}=2\) ；\\
（3）当 \(n=3\) 时，\(a_{3}+a_{4}=3\) 。所以 \(a_{4}=1\) ；\\
（4）当 \(n=4\) 时，\(a_{4}+a_{5}=4\) ．所以 \(a_{5}=3\) ；\\
（5）当 \(n=5\) 时，\(a_{5}+a_{6}=5\) ．所以 \(a_{6}=2\) ．\\
我们发现：奇数项刚好是 \(1,2,3, \ldots\) ，而偶数项刚好是 \(0,1,2, \ldots\) 所以，它们各自是公差为 1 的等差数列。为什么？可以对递推公式作差得到：

\[
\begin{gathered}
a_{n}+a_{n+1}=n \\
a_{n+1}+a_{n+2}=n+1
\end{gathered}
\]

相减得

\[
a_{n+2}-a_{n}=1
\]

故隔项是公差为 1 的等差数列。由于 \(a_{1}=1, a_{2}=0\) ，所以通项公式为：

\[
a_{n}=\left\{\begin{array}{l}
\frac{n+1}{2}, n \text { 为奇数. } \\
\frac{n-2}{2}, n \text { 为偶数. }
\end{array}\right.
\]

如何快速写出这个公式呢？\\
（1）无论 \(n\) 是奇数偶数，都是关于 \(n\) 的一次函数 \(a_{n}=k n+b\) ；\\
（2）由于隔项等差，所以下角标每增加 2 ，项会增加 1 ．所以 \(k=\frac{1}{2}\) ；\\
（3）根据 \(a_{1}=1\) ，得奇数项时 \(a_{n}=\frac{n+1}{2}\) ；根据 \(a_{2}=0\) ，得奇数项时 \(a_{n}=\frac{n-2}{2}\) 。相当于：前者是斜率为 \(\frac{1}{2}\) ，过\\
\((1,1)\) 的直线；后者是斜率为 \(\frac{1}{2}\) ，过 \((2,0)\) 的直线。这个理解在第二章和第六章中有详细介绍。\\
【注】两个方法的通项公式是相同。\\
（1）当 \(n\) 为奇数时，

\[
a_{n}=\frac{2 n-1}{4}+\frac{3}{4} \cdot(-1)^{n-1}=\frac{2 n-1}{4}+\frac{3}{4}=\frac{n+1}{2} .
\]

（2）当 \(n\) 为偶数时，

\[
a_{n}=\frac{2 n-1}{4}+\frac{3}{4} \cdot(-1)^{n-1}=\frac{2 n-1}{4}-\frac{3}{4}=\frac{n-2}{2} .
\]

方法一计算复杂，但是不需要分奇偶讨论；方法二更直观简单。\\
以上 12 个例子充分说明了：绝大多数没有额外提示的求通项问题，都可以化作＂差数列＂去累加，或者＂比数列＂去累乘。它的难点，不在于最后的累加和累乘，而在于把式子变形成 \(a_{n+1}-a_{n}=\ldots\) 或者 \(\frac{a_{n+1}}{a_{n}}\) 的形式。这里的代数技巧，本质上就是待定系数法。原理在例 7 中已经展示，我们会在后续的例题中加以大量的练习进行巩固。

这样一来，大部分求通项问题，不是累加，就是累乘，只有这两类。这充分体现了化归思想：我们把新的问题，都向经典问题转化。都把它们转化为差数列或者比数列。这样，我们学过的方法就可以派上用场。当你再遇到一个新的数列求通项的问题时，先问自己：它可以向差数列转化，还是可以向比数列转化？哪个更接近？它和之前接触过的哪类问题结构上有相似的地方？

我们再举两个简单的例子说明，化归思想在理解知识和做题中的应用。

【例 1】数列项数。很多同学在数等差数列的项数时，会遇到问题。总是记不住 \(+1,-1\) 的公式，导致数错。事实上，我们只要记住一些基本的等差数列的项数，就可以将所有不同的等差数列都化作基本的数列，从而得到项数。首先，我们知道， \(1,2,3, \ldots, n\) ，共 \(n\) 个数字。利用这个，我们可以求很多等差数列的项数：\\
（1） \(2,4,6, \ldots, 2024\) ，共 1012 项。把这个数列的每一项除以 2 ，变成

\[
1,2,3, \ldots, 1012
\]

我们没有增加或者减少项数，所以项数不变。从 1 1012，应 1012 项。所以原来也有 1012 项；\\
（2） \(1,3,5,7, \ldots, 2023\) ，也有 1012 项。我们把这个数列的每一项 +1 ，变成

\[
2,4,6,8, \ldots, 2024
\]

项数不变，所以只要看这个新的数列的项数。根据（1），它是 1012 项。问题解决。\\
（3） \(2,5,8,11, \ldots, 2024\) ，有 675 项。注意到这个数列的公差为 3 ，如果我们能过化为首项是 3 的数列，又可以＂故技重施＂。给每一项 +1 ，得到

\[
3,6,9, \ldots, 2025
\]

项数不变；再给每一项除以 3 ，得

\[
1,2,3, \ldots, 675
\]

所以共 675 项；\\
（4）\(-13,-5,3, \ldots, 2019\) ，有 255 项。注意到这个数列的公差为 8 ，所以我们化为首项是 8 的数列：给每一项 +21 ，得

\[
8,16, \ldots, 2040
\]

项数不变；再给每一项除以 8 ，得

\[
1,2, \ldots, 255
\]

所以共 255 项。\\
因此，等差数列的项数公式不需要记。遇到新的等差数列，经过适当的操作，保证项数不变，化为 \(1,2,3, \ldots, n\) 即可。

【例 2】内切球问题。很多同学学习内切球问题的时候，缺乏思路。事实上，一个化归思想足以解决问题：将所有立体几何问题，化归为平面几何问题。考虑正四棱台 \(A B C D-A_{1} B_{1} C_{1} D_{1}\) ，上底面边长 \(A B=2\) ，下底面边长 \(A_{1} B_{1}=4\) ．若它有内切球，求它的高。

我们把它的所有切点标出：根据对称性，上下底面，和其中两个不相邻的侧面的切点，应该在同一平面内。这个平面，截正四棱台得到一个等腰梯形，记为 \(E F G H\) ，其中：\(E F=2, G H=4, F G=H E, E F / / G H\) ．这个梯形的高，就是正四棱台的高。

内切球的球心 \(O\) ，就在等腰梯形的对称轴上。记 \(E F, G H\) 的中点为 \(J, K\) ，那么 \(O\) 在 \(J K\) 中点上。设梯形的高

为 \(h\) ，那么 \(O\) 的半径 \(R=\frac{h}{2}\) 。过 \(O\) 作 \(F H, E G\) 的垂线，垂足为 \(L, M\) ，那么 \(O L=O M=R=\frac{h}{2}\) ．这样，问题划归为一个平面几何问题。

根据全等三角形的判定 \((H L)\) ，三角形 \(O J F\) 和三角形 \(O L F\) 全等；三角形 \(O L G\) 和三角形 \(O K G\) 全等。所以 \(F G=F L+L G=F J+G K=3\) ．这样，等腰梯形的腰也确定了，高就可以确定了。过 \(F\) 作 \(H G\) 的垂线，垂足为 \(N\) ，那么 \(G N=1\) ，所以

\[
h=\sqrt{G F^{2}-G N^{2}}=2 \sqrt{2} .
\]

内切球半径也自然得到了：\(R=\frac{h}{2}=\sqrt{2}\) ．\\
所有的内切球问题，只要把由切点，球心组成的平面找到，解相应的四边形或者三角形，就可以解决问题。

\section*{总结}
感悟：比充要替换更本质的思想是化归！\\
充要替换，也是一种化归。它在把问题变得＂可做＂的同时，保证了转化后的问题和原问题是等价的。它更重要的意义在于，我们处理问题时，要建立起＂检查等价性＂的意识：每一步操作，是否都维持了原问题的信息？如果不是，需要增补什么？事实上，我们很难在处理所有问题时时刻刻都保证充要替换。但有了这个意识，我们就可以适时调整，使得最终得到的结果是正确的。

所有化归的操作，最终目标只有一个：都是为了把问题变得更＂简单＂。例如，将平面几何问题放到坐标系中，用解析几何来处理；将递推数列化为和等比，等差类似的数列。这样，把一个不可做的问题，变成一个可做的，熟悉的问题。只有这样，我们才可以用已经会的知识和工具，来解决新问题。所以说，化归的能力，是处理新问题最关键的能力之一。

\section*{5．8．3 XXXX}
\section*{5.9 证明方法}
很多读者看到这个标题可能感到奇怪：这是一本关于小题的书，为什么要讲大题才需要的证明方法呢？事实上，证明方法背后所体现的＂做题逻辑＂是非常重要的，这一点与大题，小题无关。例如，我们会讲到＂不妨设＂这个技巧，虽然它经常在大题出现，但是小题中也会见到；＂反证法＂所体现的＂正难则反＂的思路，其实在求取值范围，计数原理中也会反复见到。这一节中，我们会重点梳理这些大题，小题都会用到的证明技巧。

另外，对于追求选填满分的同学，证明方法也是必不可少的。当我们通过前面的技巧完成破题后，为了保证正确率，至少需要一定程度的＂证明＂。尤其是对于认知风格更偏序列型的同学（擅长计算，语言文字等复杂的序列处理），尝试去证明是非常有意义的。

\section*{5．9．1 不妨设：不失一般性}
不妨设，这里的妨指的是＂妨碍＂的意思。它和＂设＂的区别在于，＂设＂可以是没有任何理由的假设。例如，我们假设今天的温度是 20 摄氏度，然后进行讨论。这个假设不需要任何先决条件，即使今天的温度是零下二十度，也可以假设这个并进行讨论。

但不妨设，简单说，就是要求你的假设，没有妨碍到问题的讨论。换句话说，这个假设，没有对问题产生结构性的改变。什么叫结构性的改变？例如，如果已知 \(a+2 b=3\) ，但是我们假设 \(a \geq b\) ，这就会产生结构性的改

变。本来 \(a=0, b=\frac{3}{2}\) 满足题意，但是 \(a \geq b\) 把它排除在外了。这样，就可能漏掉需要讨论的情况，产生错误。另外，做出＂不妨设＂这个假设的前提，一定是它会方便我们问题的讨论和解决。如果只会让讨论更麻烦，就没必要这样做了。所以，我们在讨论以下例子时，要重点关注这两个特点：\\
（1）＂不妨设＂后，不会产生结构性改变；\\
（2）＂不妨设＂后，问题更简单了。\\
那什么样的假设是这样的呢？比较常见的，有如下几种。

\section*{1．对称性}
我们在导数问题中，经常会遇到类似 \(\frac{f\left(x_{1}\right)-f\left(x_{2}\right)}{x_{1}-x_{2}}>1\left(x_{1} \neq x_{2}\right)\) 这样的表述。部分题目中会给出已知条件 \(x_{1}>x_{2}\) ，然后推出

\[
\begin{aligned}
& f\left(x_{1}\right)-f\left(x_{2}\right)>x_{1}-x_{2} \\
& f\left(x_{1}\right)-x_{1}>f\left(x_{2}\right)-x_{2}
\end{aligned}
\]

这意味着，对 \(g(x)=f(x)-x\) ，有：当 \(x_{1}>x_{2}\) 时，\(g\left(x_{1}\right)>g\left(x_{2}\right)\) ．所以 \(g(x)\) 单调递增，故

\[
\frac{f\left(x_{1}\right)-f\left(x_{2}\right)}{x_{1}-x_{2}}>1\left(x_{1} \neq x_{2}\right) \Rightarrow g(x) \text { 单调递增。 }
\]

但我们是否真的需要这个已知条件呢？直观上我们会认为需要，如果 \(x_{1}<x_{2}\) ，那么不等号会反号，可能会推出相反的结论。出人意料的是，即使不添加 \(x_{1}>x_{2}\) 这个条件，我们也可以得到＂\(g(x)\) 单调递增＂这个结论。

我们在第七章，第八章讲过，解决问题需要不重不漏，需要尽可能充要替换。所以，我们把 \(x_{1} x_{2}\) 的两种大小关系都讨论一遍，同时对不等式作恒等变形，看能得到什么结论。\\
（1）当 \(x_{1}>x_{2}\) 时，和之前的情况类似，可以推出 \(g(x)\) 单调递增；\\
（2）当 \(x_{1}<x_{2}\) 时，除了之前乘 \(x_{1}-x_{2}\) 然后调整不等式方向，我们还可以改写分式，让分子，分母同时变为相反数：

\[
\frac{f\left(x_{1}\right)-f\left(x_{2}\right)}{x_{1}-x_{2}}=\frac{f\left(x_{2}\right)-f\left(x_{1}\right)}{x_{2}-x_{1}}>1
\]

此时 \(x_{2}-x_{1}>0\) ，那么两边同乘 \(x_{2}-x_{1}\) ，不等号方向不变，得

\[
\begin{aligned}
& f\left(x_{2}\right)-f\left(x_{1}\right)>x_{2}-x_{1} \\
& f\left(x_{2}\right)-x_{2}>f\left(x_{1}\right)-x_{1}
\end{aligned}
\]

令 \(g(x)=f(x)-x\) ，那么 \(g\left(x_{2}\right)>g\left(x_{1}\right)\) ．这意味着：当 \(x_{2}>x_{1}\) 时，\(g\left(x_{2}\right)>g\left(x_{1}\right)\) ，所以 \(g(x)\) 单调递增。因此，无论（1），（2）哪种情况，\(g(x)\) 都是单调递增的。所以，\(g(x)\) 单调递增。

原来，即使题目中不给出 \(x_{1}>x_{2}\) 这个条件，我们也可以得到这个结论！究其原因，是因为这里 \(x_{1}, x_{2}\) 地位完全相同，是对称的。我们在未来第四章的杜杆实战中讲过，\(f\left(x_{1}\right)-f\left(x_{2}\right), x_{1}-x_{2}\) 这样的式子，是反对称的：交换 \(x_{1}, x_{2}\) 后，式子变为原来的相反数：

\[
\begin{aligned}
f\left(x_{2}\right)-f\left(x_{1}\right) & =-\left(f\left(x_{1}\right)-f\left(x_{2}\right)\right) \\
x_{2}-x_{1} & =-\left(x_{1}-x_{2}\right)
\end{aligned}
\]

但是，当两个反对称的式子相乘，相除时，就会得到对称的式子：对 \(\frac{f\left(x_{1}\right)-f\left(x_{2}\right)}{x_{1}-x_{2}}\) 交换 \(x_{1}, x_{2}\) ，得到，它和原式相等：

\[
\frac{f\left(x_{2}\right)-f\left(x_{1}\right)}{x_{2}-x_{1}}=\frac{f\left(x_{1}\right)-f\left(x_{2}\right)}{x_{1}-x_{2}}
\]

所以，\(x_{1}, x_{2}\) 这两个变量完全对称。\\
当我们遇到地位完全相等的两个变量时，可以不妨假设它们的大小关系。这是因为：对于地位完全相等的

两个变量，\(x_{1}, x_{2}\) 两个记号只代表了＂这两个数＂，并没有指定谁是哪个。题中的不等式其实是在说：＂随便给我两个实数 \(x_{1}, x_{2}\)＂，我这个神奇的函数 \(f(x)\) ，它都满足 \(\frac{f\left(x_{1}\right)-f\left(x_{2}\right)}{x_{1}-x_{2}}>1\) 。假设这两个数分别是 1,3 ，那么无论是 \(x_{1}=1, x_{2}=3\) ，还是 \(x_{1}=3, x_{2}=1\) ，这个式子都可以写成：

\[
\frac{f(3)-f(1)}{3-1}>1
\]

所以一般地，无论 \(x_{1}, x_{2}\) 谁大谁小，我们都可以把这个式子写成

\[
\frac{f(\text { 大 })-f(\text { 小 })}{\text { 大 }- \text { 小 }}>1
\]

这样的话，我们就直接讨论 \(x_{1}\) 是大数，\(x_{2}\) 是小数的情况就好了。所以，我们可以不妨假设 \(x_{1}>x_{2}\) ．\\
另一种理解方式是：对于 \(x_{2}>x_{1}\) 的情况，我们只是把所有的 \(x_{2}\) 都替换成了 \(x_{1}, ~ x_{1}\) 都替换成了 \(x_{2}\) 。论证的道理是完全相同的（同理），那么只讨论一种就够了。所以，我们可以不妨设第一种情况成立，第二种情况道理相同，就不用再单独处理了。所以，这也方便了问题的讨论。

\section*{2．齐次性}
在证明部分不等式的时候，经常会出现假设某个变量或者某个式子为 1 的情况。例如，证明不等式

\[
x^{2}+y^{2}+z^{2} \geq x y+x z+y z
\]

可以有很多方法：例如，观察到等式两侧都是关于 \(x, y, z\) 的二次齐次式，所以可以两边同时除以 \(z^{2}\) ；由于它是关于 \(x, y, z\) 都对称，可以利用对称换元；可以选主元等等。但如下＂不妨设＂的方法，也是一个正确的方法：

不妨设 \(z=1\) ．那么原不等式化为

\[
x^{2}+y^{2}+1 \geq x y+x+y
\]

选 \(x\) 为主元，整理得

\[
x^{2}-(y+1) x+y^{2}-y+1 \geq 0
\]

为了证明这个关于 \(x\) 的二次不等式成立，可以考虑：\\
（1）配方：

\[
\begin{aligned}
x^{2}-(y+1) x+y^{2}-y+1= & \left(x-\frac{y+1}{2}\right)^{2}-\frac{(y+1)^{2}}{4}+y^{2}-y+1 \\
& =\left(x-\frac{y+1}{2}\right)^{2}+\frac{3}{4} y^{2}-\frac{3}{2} y+\frac{3}{4} \\
& =\left(x-\frac{y+1}{2}\right)^{2}+\frac{3}{4}(y-1)^{2} \geq 0
\end{aligned}
\]

故原不等式成立；\\
（2）判别式：等价于判别式非负，即

\[
\Delta=(y+1)^{2}-4\left(y^{2}-y+1\right)=-3 y^{2}+6 y-3=-3(y-1)^{2} \leq 0
\]

成立。\\
故原不等式成立。\\
这个方法看上去非常荒谬：怎么可以只考虑 \(z=1\) 的情况，就说明不等式是成立的呢？事实上，这里我们利用了关键的代数结构：齐次性。这个不等式两侧都是关于 \(x, y, z\) 的二次齐次式，这意味着：如果等比例改变 \(x, y, z\) ，那么不等号是不会发生改变的。例如：\\
（1）假设 \(x=3, y=2, z=1\) ，那么

\[
\begin{aligned}
& x^{2}+y^{2}+z^{2}=1+4+9=14 \\
& x y+x z+y z=2+3+6=11
\end{aligned}
\]

\(14 \geq 11\) ，所以左大于右；\\
（2）如果等比例扩大若干倍，会发生什么？假设 \(x=30, y=20, z=10\) ，那么

\[
\begin{aligned}
& x^{2}+y^{2}+z^{2}=100+400+900=1400 \\
& x y+x z+y z=200+300+600=1100
\end{aligned}
\]

\(1400 \geq 1100\) ，所以左大于右。\\
对于其他的比例，大家可以直接验算，不等号的方向不会改变（相当于给不等式两侧同乘 \(k^{2}, k^{2}\) 非负，不等号方向不变）。

既然等比例改变的时候，不等号方向是不改变的，那么我们只要检验：对每个不同的比例 \(x: y: z\) 中的一组，不等号是否成立。如果成立，那么和这组数等比例的 \(x, y, z\) ，也一定成立。那么我们检验每个不同的比例中的哪一组呢？可以检验 \(Z=1\) 的那一组，它是最简单，直接的。这就是＂不妨设＂的道理。

既然观察到齐次，根据第二章的思想，我们可以考虑不等式两侧同时除以 \(z^{2}\) ，利用比例来表示这个式子。得：

\[
\left(\frac{x}{z}\right)^{2}+\left(\frac{y}{z}\right)^{2}+1 \geq \frac{x}{z}+\frac{y}{z}+\frac{x}{z} \cdot \frac{y}{z}
\]

这个式子和我们假定 \(z=1\) 是完全一样的：做一个简单的换元，令 \(x^{\prime}=\frac{x}{z}, y^{\prime}=\frac{y}{z}\) ，原不等式等价于

\[
\left(x^{\prime}\right)^{2}+\left(y^{\prime}\right)^{2}+1 \geq x^{\prime}+y^{\prime}+x^{\prime} y^{\prime} .
\]

所以，不妨设 \(z=1\) ，其实就是在不等式两侧同时除以 \(z^{2}\) 。本来是 \(x: y: z\) 的三个数，变成了 \(\left(\frac{x}{z}\right):\left(\frac{y}{z}\right): 1\) 的三个数，比例不变，而 \(z=1\) 。

这个假设的意义在于，它帮我们把一个三元问题化归为了一个二元问题。变量个数越少，越好解决，这也是化归思想的体现。

总结一下：当我们遇到有对称，齐次结构的式子时，必要时可以添加假设。对称的式子可以假设大小关系，齐次的式子可以添加一个非齐次条件。

\section*{5．9．2 正难则反}
如果一个问题的正面非常难解决，或者反面很容易，都可以考虑从反面去考虑。如下若干个小问题，都非常适合从反面考虑来解决。\\
（1）扔 6 个均匀的正方体鹃子，至少有一个大于等于 5 的概率是多少？如果我们直接考虑＂至少有一个大于等于 \(5 "\) ，需要分类讨论且计算量较大。对于一个骰子而言，大于等于 5 的概率为 \(\frac{2}{6}=\frac{1}{3}\) ，那么概率为

\[
p=C_{6}^{6}\left(\frac{1}{3}\right)^{6}+C_{6}^{5}\left(\frac{1}{3}\right)^{5}\left(\frac{2}{3}\right)+C_{6}^{4}\left(\frac{1}{3}\right)^{4}\left(\frac{2}{3}\right)^{2}+C_{6}^{3}\left(\frac{1}{3}\right)^{3}\left(\frac{2}{3}\right)^{3}+C_{6}^{2}\left(\frac{1}{3}\right)^{2}\left(\frac{2}{3}\right)^{4}+C_{6}^{1}\left(\frac{1}{3}\right)\left(\frac{2}{3}\right)^{5}
\]

如果从反面考虑，＂至少有一个大于等于 5 ＂的反面是＂所有都小于 5 ＂，而一个股子小于 5 的概率为 \(\frac{4}{6}=\frac{2}{3}\) ，所以

\[
1-p=\left(\frac{2}{3}\right)^{6}=\frac{64}{729}
\]

故 \(p=1-\frac{64}{729}=\frac{665}{729}\) ．\\
如果你担心计算结果错误，还可以简单验证一下。只要把＂至少有一个大于等于 5 ＂的概率和＂所有都小于 5 ＂的概率相加，看是不是 1 即可。事实上，这就是二项式定理＂反过来＂写：

\[
C_{6}^{6}\left(\frac{1}{3}\right)^{6}+C_{6}^{5}\left(\frac{1}{3}\right)^{5}\left(\frac{2}{3}\right)+C_{6}^{4}\left(\frac{1}{3}\right)^{4}\left(\frac{2}{3}\right)^{2}+C_{6}^{3}\left(\frac{1}{3}\right)^{3}\left(\frac{2}{3}\right)^{3}+C_{6}^{2}\left(\frac{1}{3}\right)^{2}\left(\frac{2}{3}\right)^{4}+C_{6}^{1}\left(\frac{1}{3}\right)\left(\frac{2}{3}\right)^{5}+\left(\frac{2}{3}\right)^{6}=\left(\frac{1}{3}+\frac{2}{3}\right)^{6}=1
\]

（2）丢两个均匀的正方体骰子，第一个点数比第二个点数大的概率是多少？直接枚举，并不复杂：所有可能情况有 \(6^{2}=36\) 种，其中第一个比第二个大的有：\\
（1）第一个为 \(2:(2,1)\) ；\\
（2）第一个为 \(3:(3,1),(3,2)\) ；\\
（3）第一个为 \(4:(4,1),(4,2),(4,3)\) ；\\
（4）第一个为 \(5:(5,1),(5,2),(5,3),(5,4)\) ；\\
（5）第一个为 \(6:(6,1),(6,2),(6,3),(6,4),(6,5)\) ．\\
共 15 种，故概率为 \(\frac{15}{36}=\frac{5}{12}\) 。\\
如果从＂反面＂考虑，可以利用对称性。所有可能的情况无非三种：第一个大，相等，第二个大。根据对称性，第一个大的概率必然和第二个大的概率是一样的：只要将第一个大的样本点，交换点数，就变成了第二个大的样本点。所以，只要减去相等的概率再除以 2 就是答案。相等的概率显然是 \(\frac{6}{36}=\frac{1}{6}\) ，所以答案为 \(\frac{1-\frac{1}{6}}{2}=\frac{5}{12}\) （3）扔 12 枚硬币，正面比反面多的概率是多少？如果（2）中＂枚举＂计算量也不大，这里的计算量就不太能接受了。如果直接枚举（正面考虑），那么概率为

\[
p=C_{12}^{12}\left(\frac{1}{2}\right)^{12}+C_{12}^{11}\left(\frac{1}{2}\right)^{12}+C_{12}^{10}\left(\frac{1}{2}\right)^{12}+C_{12}^{9}\left(\frac{1}{2}\right)^{12}+C_{12}^{8}\left(\frac{1}{2}\right)^{12}+C_{12}^{7}\left(\frac{1}{2}\right)^{12}
\]

这样需要计算大量的二项式系数。\\
如果从＂反面＂考虑，仍利用对称性。所有可能的情况无非三种：正面多，一样多，反面多。根据对称性，正面多的概率和反面多的概率是相等的：只要将所有硬币翻面，正面多的情况就变成反面多的情况了。所以，只要减去一样多的情况再除以 2 ，就是答案。故

\[
p=\frac{1-C_{12}^{6}\left(\frac{1}{2}\right)^{12}}{2}=\frac{793}{2048}
\]

这里的对称性，其实是二项式系数的对称性。如果考虑反面个数多的情况，列出来的式子是：

\[
p^{\prime}=C_{12}^{0}\left(\frac{1}{2}\right)^{12}+C_{12}^{1}\left(\frac{1}{2}\right)^{12}+C_{12}^{2}\left(\frac{1}{2}\right)^{12}+C_{12}^{3}\left(\frac{1}{2}\right)^{12}+C_{12}^{4}\left(\frac{1}{2}\right)^{12}+C_{12}^{5}\left(\frac{1}{2}\right)^{12}
\]

而 \(C_{12}^{k}=C_{12}^{12-k}\) ，所以 \(p=p^{\prime}\) ．

除了在概率，排列组合中，求取值范围的时候也会用到正难则反。当我们看到＂至少一个＂这样的描述时，都可以从反面考虑，这样就可以把问题变得简单清晰明了。反证法，本质上也是一种正难则反：直接证明比较复杂，所以考虑反面。

\section*{5．9．3 同理}
字面意思就是＂道理相同＂。道理相同，只是替换了一个数字，那么可以省略论述的道理，直接出结果。\\
（1）圆雉曲线的计算中，经常出现同理。最简单，当我们计算完关于 \(\left(x_{1}, y_{1}\right)\) 的式子后，得到了结果；如果想要对 \(\left(x_{2}, y_{2}\right)\) 进行相同的计算，那么只要将 \(x_{1}\) 替换为 \(x_{2}\) ，将 \(y_{1}\) 替换为 \(y_{2}\) 即可。

更复杂的情况：假设过 \((1,0)\) ，斜率为 \(k\) 的直线截椭圆的弦长为

\[
\sqrt{1+k^{2}} \cdot \frac{|2 k|}{2 k^{2}+1}
\]

那同样是过 \((1,0)\) ，斜率为 \(-\frac{1}{k}\) 的直线截椭圆的弦长是多少？注意，这里唯一的区别在于斜率从 \(k\) 变为了 \(-\frac{1}{k}\) ，所以只要将结果中的 \(k\) 都替换为 \(-\frac{1}{k}\) 即可：

\[
\sqrt{1+\frac{1}{k^{2}}} \cdot \frac{\left|2 \cdot\left(-\frac{1}{k}\right)\right|}{\frac{2}{k^{2}}+1}=\sqrt{\frac{1+k^{2}}{k^{2}}} \cdot \frac{|2 k|}{k^{2}+2}=\sqrt{1+k^{2}} \cdot \frac{2}{k^{2}+2}
\]

事实上，我们可以把弦长看成 \(k\) 的函数：

\[
f(k)=\sqrt{1+k^{2}} \cdot \frac{|2 k|}{2 k^{2}+1} .
\]

从解析式或者函数的定义上看，它都是函数：对于给定的 \(k\) ，都存在一个确定的弦长和它对应。因此，\(-\frac{1}{k}\) 的情况，就是一个变量替换（或者说是函数的复合），计算 \(f\left(-\frac{1}{k}\right)\) 。\\
（2）不妨设，很多时候也是同理。有些不妨设的情况，就是因为道理相同，所以我们只论证一遍。例如，当 \(a, b, c\)地位完全相同时，我们不妨设 \(a\) 是最大数。这是因为 \(a, b, c\) 哪个最大，道理都一样，只不过是换了个字母。

\section*{5．9．4 XXXX}
\section*{第6章 科学的思维方式解题}
本章为各模块阳哥的分享解题思想的内容，也有自己结合前面各章节的内容对一些好题的讲解。选题以难题，综合题为主，锻炼具体题目解题思想，达到一题抵百题的效果。

\section*{6.1 排列组合与计数问题}
排列组合计数问题在高考中是一类独立的题型，在概率大题中有时会出现，同时小题里一定会有一道。通常来讲这个小题的难度不会很高。这样的小题要么是计数，要么是求概率，不过本质上大多数都和计数有关。

之前不少观众也希望我讲一点排列组合问题，从我的观察来看，大家的问题主要有：分类讨论不清楚，不知道该怎么分类；会出现讨论或计数时有重复或有遗漏的情况；稍微复杂点，思维就容易乱，等等。

本节内容将主要围绕计数问题展开，主要讲以下几个内容：怎样的分类讨论是正确的；怎么用对称性简化讨论；怎么用树状图辅助思考。

注 本节内容所有观点均为 up 主（nichijou 的涂山苏苏）本人的做题经验与心得。写出来单纯是因为，多一种做题习惯，多一种思考方式，必定是没有坏处的。是他试做几个排列组合题并同步观察自己的习惯之后总结的，大家看个乐呵就行。

\section*{6．1．1 计数问题的题目特征}
计数问题是在某些限制下计算满足某个问题的所有情况的个数，问题本身由背景和牵制两部分组成。例如，把 \(1, ~ 2, ~ 3, ~ 4,5\) 排成一列， 1 不在首位的情况数量。问题背景就是＂ \(1,2,3,4,5\) 排成一列＂，牵制就是＂1不在首位＂。

题目类型：纯静态问题，我们要计的数都是具有确定性的。通过题目，我们能明确知道我们要计算一个什么的数量。这意味着，计数问题实际上如同案板上的鱼肉，唾手可得！

题目类别：应用题，通常具有一个实际背景，题干叙述往往是自然语言。这意味着，把题目要求计的内容先变成具有确定性的数学模型是关键！

常用方法：分类讨论，反向思考（取补集）。分类讨论是先把鱼剁成鱼头，鱼身等部分，再分开各个击破；反向思考是在一个全集内，考虑不满足计数要求的情况数量，再用全集内的情况数量作差得到结果。

\section*{6．1．2 分类讨论——对立事件}
作为计数问题，唯一的原则就是不重复，不遗漏。因此在选择计数的时候，只能做充要替换，也就是寻找原计数要求的等价变形！！

因此，这个时候，如果我们要进行分类讨论的话，必须必须必须，分的几类是对立的事件！

\section*{定义 6.1}
若一族事件 \(A_{1}, A_{2} \cdots A_{n}(n \geq 2)\) 满足，任意两个不能同时发生，也就是 \(\forall i, j=1,2 \cdots n, i \neq j, A_{i}, A_{j}\) 均不可能同时发生；且必定有一个事件发生，也就是 \(\exists k=1,2 \cdots n, A_{k}\) 发生，那么称 \(A_{1}, A_{2} \cdots A_{n}\) 是彼此对立的。

面对一道计数问题，如果我们希望分类讨论，分成一个个小情况 \(A_{1}, A_{2} \cdots A_{n}\) 来分开计数，最后再加起来，那么一定一定一定，要保证 \(A_{1}, A_{2} \cdots A_{n}\) 彼此对立！

例 6.1 安排 5 个人到 \(a, b, c\) 三家宾馆入住，每家宾馆至少一个人入住，且 \(A, B\) 两人都在 \(a\) 宾馆入住，则安排种数为（）

做题第一步：翻译成数学模型：\(A, B, C, D, E\) 扔到三个筐 \(a, b, c\) 里，每个筐里面都要有东西，\(A, B\) 都扔到 \(a\)筐。（注意：不擅长计数的同学，请先老老实实想明白要计数的具体内容，也就是数学模型，可以用集合语言来想，最准确）

分类讨论的划分：选择的方法就特别多，基本上只要划分的对，就能做！比如可以这样划分：\(a\) 只有 \(A, B ; a\)还有新的 1 个；\(a\) 还有新的 2 个；\(a\) 还有新的 3 个（后两种情况后续可以轻易排除）。还可以这样划分：\(a, b, c\) 的人数分别是 \(311,221,212 \ldots \ldots\).

划分完之后，特别特别特别重要的一步是：即时检查！！其他题目我不强求大家随时做好即时检查，但是排列组合的分类讨论，请务必及时检查一下！！检查内容就是，你讨论划分的这几类是不是彼此对立的。

比如刚才第一个划分，\(a\) 还有几个？肯定同一个情况 \(a\) 不会既 2 个又 3 个，所以肯定只能发生一次；每一种情况 \(a\) 都对应有入住的人数，所以必定发生一个，于是检查完毕。这就是标准的检查对立性的方法——设问。

自己问问自己，是不是两两不交叉？是不是每个情况都要对应一个你的划分？自己问问自己，检查效果比什么都好！

另外，关于对立事件的分类讨论，我想再提两点：首先，一道题只要你能成功地把问题分解成 2 个以上的彼此对立的事件，你接下来只要专心地分开看，每一个小情况拿出来单独分析，最后合起来就行了。因为只要分开了小情况，基本上每个小情况的种类数都会比最后的答案小很多，实际上已经起到了简化问题的作用。尽管可能方法不是最简便的，但是本身高考中的计数问题也不会难的离谱，所以放心地去试着做一下吧。

另外，有的时候，某一个小情况的讨论，可能还需要继续再分类，这个时候，我们就对这种小情况重复刚才讲到的对立事件分类讨论，继续研究即可。当然，除了分对立事件讨论以外，还有一个重要的思想就是反向思考。比如刚才提的小例子， \(1 \sim 5\) 排列， 1 不在首位，我们就可以计算 1 在首位的情况数，然后用 \(1 \sim 5\) 的排列减去刚才的数字，就是 \(5!-4!\) 。实际上，这也是一种对立事件，因为如果我们把 \(1 \sim 5\) 排列看作一个计数对象的话，＂ 1在首位＂和＂1不在首位＂就是一组对立事件，所以＂1在首位＂+ ＂ 1 不在首位＂\(=\)＂ \(1 \sim 5\) 排列＂。因此反向思考也属于对立事件划分的一种。这两大思想基本上可以处理所有的排列组合计数问题了。

\section*{6．1．3 对称计数技巧一 双射思想}
当分类或者反向思考结束之后，我们常常会得到接下来需要计数的好几个小事件。此时，经常会遇到的一种可能是，我们分类的某一些事件，它们看上去差不多，而且，包含的情况数目是相等的。遇到这种事情，我们就可以少计算一部分事件包含的情况数，进而简化问题。

但是，做出这样的判断到底是否足够让我们放心？如何正确利用这种对称性来帮助我们提高讨论计数的效率？

我们可以这样想，不同的分类出来的事件包含的情况可以分别看成一些集合，如果根据对称性，我们想判

定某两个集合的元素个数相等，实际上它们的元素也应该在这种对称性之下彼此对应。比如一开始问大家的问题， \(1,2,3,4,5\) 排成一列， 1 不在首位，我们可以像刚才一样分类讨论 1 分别在第二，三，四，五位的情况。（当然可以直接 \(5!-4!\) 做，我只是拿来做个例子）

比如我们考虑 1 在第二位， 1 在第三位，这两个大类分别的情况数，对称性的直觉让我们感受到，这两大类里的情况数量应该是一样的，严格的理性告诉我们确实如此，因为如果按照 \(2,3,4,5\) 的排列顺序的话，这两个大类里的元素可以建立一一对应。

比如 21345 和 23145 可以对应（排列顺序都是 2345）， 31254 和 32154 可以对应（排列顺序都是 3254），因此这两个大类的所有情况可以这样来一一对应起来。

自然地，情况数量也就相同了。同理可以想到， 1 的位置为基准分出的四个大类，各自的元素个数都相等，因此我们只需要计算任何一个大类的元素个数即可，每一个的话都是 4 ！，一共就是 \(4 \times 4\) ！。

也就是说，当我们把计数对象经过分类，变成了 \(A_{1}, A_{2} \cdots A_{n}\) 这些彼此对立的事件后，如果对于其中的某些 \(A_{i}, A_{j}\) ，我们通过问题本身的＂对称性＂构造出二者中的元素之间的一一对应，那么 \(A_{i}, A_{j}\) 的元素个数必定相等。这种思考问题的方法我称之为双射思想（这比起对称这种命名，更体现一一对应在逻辑上对这种思考方式正确性的保障）。

对于双射思想，我的建议是，做题的时候，只要遇到对称性相关的，似乎可以确保二者数量相等的时候，都用双射思想稍微思考一下：能不能根据对称性，给出这两大类的一一对应？最多花你 10 秒时间，但是能让你的讨论更踏实。毕竟不少同学可能做排列组合总是心里不踏实，觉得是不是漏了，重复了？但是一一对应这种逻辑上的保障能够让我们有一种安全感。

下面我们回到例题 1：\(A, B, C, D, E\) 扔到三个筐 \(a, b, c\) 里，每个筐里面都要有东西，\(A, B\) 都扔到 \(a\) 筐。我们给出了这样划分：\(a\) 只有 \(A, B ; a\) 还有新的 1 个；\(a\) 还有新的 2 个；\(a\) 还有新的 3 个。其中后两种情况容易排除，因为后两种情况会导致 \(a\) 中不少于 4 ，那么 \(b, c\) 至少空一个，不符合题意。因此只需要讨论 \(a\) 只有 \(A, B ; a\) 还有新的 1 个。下面就要独立地看这两部分。

如果 \(a\) 只有 \(A, B\) ，那么 \(C, D, E\) 随便往 \(b, c\) 里扔，要求每个筐都有人。我们可以继续讨论，\(b 2 c 1\) 或 \(b 1 \mathrm{c} 2\)（是对立的吧，自己想想）。发现这俩数量一样，都是3（直接用双射思想也行，两个分别算一下也不费时间），合起来， 6 。如果 \(a\) 还有新的一个，新的可以是 \(C, D, E\) ，发现是对立的，并且是对称的（可以思考一下这个一一对应怎么造）。于是算出 \(a\) 是 \(A, B, C\) 的情况，乘 3 就行；这时候 \(D, E\) 随机扔到 \(b, c\) ，显然就 2 个扔法。于是 \(3 \times 2=6\)

合起来就是 \(6+6=12\) 。

\[
\text { 树状图示例: }\left\{\begin{array}{l}
a \text { 只有 } A, B\left\{\begin{array}{l}
b 2 c 1: 3 \\
b 1 c 2
\end{array}\right. \\
a \text { 有一个新的 }\left\{\begin{array}{l}
C: 2 \\
D \\
E
\end{array}\right.
\end{array}\right.
\]

我们再换个做法，来证明随便分类都能做。我们还讲了划分方式：\(a, b, c\) 的人数分别是 \(311,221,212\) 。不难看出后两者对称，一一对应就是 \(b, c\) 互换。

第一个， 311 ，相当于 \(a\) 有 \(A, B\) 前提下，剩下三个 \(C D E\) 排列一下就插进来，于是总数为 \(A_{3}^{3}=6\) ；第二个，\(b\)

2 cl ，数量是 \(C_{3}^{2}=3\) 。最后合起来： \(6+2 \times 3=12\) 。

\[
\text { 树状图示例: }\left\{\begin{array}{l}
a 3 b 1 c 1: A_{3}^{3}=6 \\
\left\{\begin{array}{l}
a 2 b 2 c 1: C_{3}^{2}=3 \\
a 2 b 1 c 2
\end{array}\right.
\end{array}\right.
\]

可以看到，上述讨论过程非常清晰，并且在对立事件和双射思想的加持下毫无难度，可以稳收。如果思考比较慢，还可以画树状图来辅助思考（这一部分看视频讲解）。

从中可以看到，对立事件＋即时检查确保好的分类，并时不时用反向思考和双射思想扫视一下，看看哪里有简化的希望；讨论复杂一点就配合树状图直观观察，基本是可以稳通大部分计数问题的。

注意一下，树状图的绘制可以直接在试卷上做，省时间，还可以随时检查已知条件是否看错。

\section*{6．1．4 综合应用}
我们再看一个稍微复杂一点点的问题。

例 6.2 有 3 个男孩， 3 个女孩， 1 个老师站成一排，老师站中间，有且仅有 2 个女孩相邻，则不同站法种数为（）

这个翻译简单，就是个排列问题。感觉这个直接反向思考不太好啊，于是我随手写下一个：＂？？？师？？？＂，然后开始思考，如何正面去分类讨论？

很直白的想法就是，两个一起的女孩在哪？如果位点是 1234567 的话，自然就有 \(12,23,56,67\) 这 4 种（是对立的，没错）马上拿对称和反向思考扫一下，发现只要做一个整个排列的镜像对称，就可以给 12 这种和 67 这种建立一一对应， 23 和 56 同理，于是我写下树状图：

\[
\text { 树状图示例: }\left\{\begin{array}{l}
\text { 一起的女孩在 }\left\{\begin{array}{l}
12 \\
67
\end{array}\right. \\
\text { 一起的女孩在 }\left\{\begin{array}{l}
23 \\
56
\end{array}\right.
\end{array}\right.
\]

12 这种，另一个女孩在哪？现在顺序是＂12？师？？？＂，那就有 \(5,6,7\) 三个地方了。这三个地方数量一样？双射一下：只要把原来男孩按顺序填进去，女孩换到其它位点，就是一一对应，所以是一样的！

23 这种，另一个女孩在哪？现在顺序是＂？ 23 师？？？＂，那就有 \(5,6,7\) 三个地方了。和刚才一样，这仁情况都一样的数量！

\[
\text { 树状图进化: }\left\{\begin{array}{l}
\text { 一起的女孩在 }\left\{\begin{array}{l}
12\{3 \times " \text { 女孩在 } 125 " \\
67
\end{array}\right. \\
\text { 一起的女孩在 }\left\{\begin{array}{l}
23\{3 \times \text { 女孩在 } 235 \\
56
\end{array}\right.
\end{array}\right.
\]

OK，那么女孩在 125 ，男孩在 367 ，一共就 \(A_{3}^{3} A_{3}^{3}=36\) ；女孩在 235，男孩在 147 ，一共也是 \(A_{3}^{3} A_{3}^{3}=36\) 。最后就是 \(36 \times 3 \times 2+36 \times 3 \times 2=432\) 。

可以看到，讨论上没有太大的难点，全程重复刚才讲的思想。\\
最后，我们讲一道难题，供有能力的观众拔高：

例 6.3 （2013 年全国高中数学联赛第 8 题）已知数列 \(\left\{a_{n}\right\}\) 共有 9 项，其中，\(a_{1}=a_{9}=1\) ，且对每个 \(i \in\{1,2, \cdots, 8\}\) ，均有

\[
\frac{a_{i+1}}{a_{i}} \in\left\{2,1,-\frac{1}{2}\right\}
\]

则这样的数列的个数为（）

题目让我们计算数列个数，因此我们要处理好问题的等价变形。首先，要把问题转化成我们熟悉的背景 + 牵制的形式。这里把 \(a_{1}=1\) 和

\[
\frac{a_{i+1}}{a_{i}} \in\left\{2,1,-\frac{1}{2}\right\}
\]

作为题目背景，把 \(a_{9}=1\) 作为牵制更让人感觉舒服，相当于在那样的数列里筛选满足 \(a_{9}=1\) 的数列。\\
那么，如何去找？观察已知，\(a_{1}=1\) 必然是基底，要去找 \(a_{9}\) ，要从一堆

\[
\frac{a_{i+1}}{a_{i}}
\]

出发的话，那必然只能是全乘起来了！从 \(a_{1}\) 到 \(a_{9}\) 一共要乘 8 个 \((i=1,2 \cdots 8)\) ，因此这 8 个乘起来就是 1 （终点是\\
1 ）。注意！！这是一个等价！！\\
因此，我们肯定要考察这 8 个数了，原命题等价于这 8 个数乘起来是 1 ，而且，用上我们讲过的双射思想，每一组这样按顺序排列的 8 个数，和一个题干中的 9 项数列，也是一一对应的！给你这按顺序排列的 8 个比值，唯一决定一个 9 项数列；给你一个 9 项数列，也唯一决定一组这样按顺序排列的 8 个比值。于是，我们可以只研究这样按顺序排列的 8 个比值就行！这 8 个比值也相当于一个数列啊，那按照数列的语言来说吧。每一项都要从

\[
\left\{2,1,-\frac{1}{2}\right\}
\]

里选，最后乘积还是 1 ，很明显了，要考虑各个数出现的个数。\\
如果 2 出现次数为 \(m,-\frac{1}{2}\) 出现次数为 \(n\) ，那么

\[
2^{m}\left(-\frac{1}{2}\right)^{n}=1
\]

于是整理一下就是

\[
(-1)^{n} 2^{m-n}=1,
\]

显然 \(n\) 为偶数，\(m=n\) 。\\
那想一下，怎样就行？也就是 2 和 \(-\frac{1}{2}\) 出现同样的次数，还都是偶数次，那它们公共的出现次数只能是 \(0,2,4\)了。就分成了 3 个情况，不难看出这 3 个情况是互相对立的。

俩都出现 0 次，全是 1 ，就一种。\\
俩都出现两次，相当于一个 8 个数的排列，选俩是 2 ，选俩 \(-\frac{1}{2}\) ，其余都是 1 ，这不就 \(C_{8}^{2} C_{6}^{2}\) 吗？\\
俩都出现 4 次，那就 \(C_{8}^{4} C_{4}^{4}\) 呗！合起来就是

\[
1+C_{8}^{2} C_{6}^{2}+C_{8}^{4} C_{4}^{4}=491
\]

\section*{总结}
排列组合计数问题其实并不难，只要大家掌握好了思想方法。我们总结一下拿到排列组合题目的做题流程： （1）首先，翻译一下你要干什么，把题干翻译成纯粹的排列组合模型。然后，将你要计数的问题进行分类讨论，分类的标准必须必须必须是彼此对立的事件。然后将你分类了的一个个小事件进行计数，如果还不好记，就继续分类。分类的时候不要忘了通过自问自答（每种情况都能匹配到一个？我分的各种事件互不相交？来即时检查分类的准确性！（2）同时，手握两把宝剑——反向思考和双射思想，时刻都留个心眼，哪里适合就随时拿来使用。好比 NPC 给了你藏宝图和两把武器，让你完成寻宝任务，你照着藏宝图走，顺便两把武器来清怪，开路。（3）如果讨论比较复杂，容易混乱，别忘了画个讨论用树状图，好比开启天眼模式，方便思考。

通常高考真题中的排列组合不会特别难，正常按照上述路径走都能走通。如果碰上比较难的，就要努力去充要替换，把问题尽量等价成方便计数的情形。另外，虽然我们今天没有遇到，但是，假如真的碰到了怎么讨论都讨论不明白的情况，还是考虑一下换个分类方法吧。通常结合牵制条件都能找到好的分类方法。祝大家今后排列组合问题不再出错！

\section*{6.2 数列难题}
事前声明\\
本次内容为 up 主结合粉丝私信提问，所提胨出的数列问题的较为普遍的思维方式与解题理念。阅读所需的准备为学习过高中课内全部数列问题，最好看过第一章前五节的内容，对本人一直以来的讲解实践的思维方式 （动静视角，控制理念等）较为熟悉。

本次选题均为粉丝私信提问的题目，整体难度会比较高，大部分题目属于压轴难度（甚至好几道能达到高考的最高难度了），因此本节内容可能需要大家多用心消化。

数列是高中数学中极为特殊的模块，有两大特点：一是非常容易产生超难题，难度可以毫无上限。数列不仅是高考的一大组成部分，在高中数学竞赛里面更是极为活跃，联赛的压轴小题经常可以看见数列的身影；另一个是很重视代数变形与计算，可以说是高中数学所有部分里面最重视技巧的模块了。幸运的是，除了浙江卷这种喜欢竞赛风格和超难题的卷子以外，像全国卷地区的数列题并没有那么恐怖，整体仍然是思维方式和逻辑底力起主要作用。但是即便如此，它仍然会对计算与分析有比较高的要求。

本次数列问题的经验分享，将会以此前的＂思维方式与逻辑能力＂章节为基础，采用习题课的方式，通过题目来展示高考数列问题的思维方式。需要注意的是，由于数列本身难度无上限的特点，许多同学在平时练习（例如部分模拟题）中都可能会遇到一些难度接近竞赛的数列题，其可能使用一些非常依赖灵感的处理手段。这些题目如果遇到，建议通过此前讲的 2,5 分钟原则进行取舍。

\section*{6．2．1 数列有几个自由度？}
\section*{6．2．1．1 无穷自由度}
我们在最开始学习数列的时候就会学到，数列是定义在一些正整数（自然数）上的函数。有的数列只有有限项，比如，\(a_{n}=n\) ，取 \(n=1,2 \cdots 100\) ；更多的是定义在整个正整数集上的数列，也就是无穷项数列。

如果我问：一个数列有几个自由度，你该如何回答？\\
似乎有点难以直接回答，我们先看几个特殊例子。首先，一个等差数列有几个自由度？那显然是两个。首项和公差这两个＂主自由度＂就能定死整个等差数列，它们也彼此独立；一个等比数列呢？也是两个，主自由度选为首项和公比。这些都感觉是以前的老问题了。那我再问：一个 100 项的数列，它有几个自由度呢？这时候它就单纯是一个数列罢了，没别的已知条件，那只能说它的自由度是 100 ，每一项都是一个自由度。

那我再问：考虑一个数列 \(\left\{a_{n}\right\}, n \in N^{+}\)，它有几个自由度呢？\\
每一项都是自由度．．．．．．那它就有无穷项，所以只能说，它有无穷个自由度。这听上去感觉很不爽，但我们必须要接受这个事情，也就是一般的无穷项数列的自由度其实是无穷的，准确的来说，是＂正整数个＂。我们似乎第一次遇到无穷自由度这样的事情，但其实除了数列（正整数集上的函数），还有一个更常见的对象也是这样，那就是函数！我考虑一个定义在 \(R\) 上的函数 \(f(x)\) ，不给其他任何信息，那么这个 \(f(x)\) 在每一个点的取值都可以看作一个自由度，所以合起来就有＂实数个自由度＂！这样的函数我们其实有一个名字，叫做抽象函数，就是别的信息没有，单纯就一个函数，所以抽象函数的自由度也是无穷的。这就是我们在高中接触到的自由度为无穷的对象。

无穷自由度和有穷自由度的对象，它们的性质差异会很大，大学数学的很多内容都是研究无穷维对象的，像微积分，极限理论等就是典型，继续深人会涉及到函数空间，泛函分析等分析学的领域，感兴趣的同学可以去自行了解。

\section*{6．2．1．2 静止化方法}
回到正题。一般数列有无穷自由度，但是我们做题的时候，像我们平时做的等差等比，自由度却只有 2 。这是为什么呢？原本那么多自由度都去哪了？

要理解这个问题，我们需要回到等差，等比的原始定义上去。以等差数列为例，什么是等差数列？相邻项之差为定值（公差）。如果我们直接把 \(a_{1}\) 和 \(d\) 都定死，那么相当于先消灭首项的自由度，那么后面的项的自由度都是怎么消失的呢？稍微想想就能明白，这是因为 \(a_{n}+1-a_{n}=d\) 这个关系是对所有 \(n\) 成立的，于是 \(a_{1}\) 靠这个定出 \(a_{2}\) ，然后 \(a_{2}\) 再靠这个定出 \(a_{3} \ldots \ldots\). 如此下去，前项不断的定出后项，像多米诺骨牌一样滚过去，最后整个数列都变成静止了。

这里面的关键在哪呢？就在于它的静止化是靠的这种多米诺骨牌的方法：给出一个递推关系，这个递推关系本身就起到了静止化的作用，相当于拉出了一个要素网：

\[
a_{1} \rightarrow a_{2} \rightarrow \cdots a_{n} \rightarrow \cdots
\]

因此，通过明确递推关系，可以把数列的无穷自由度直接削减到很小，归结到递推关系和首项本身上面来。也有同学会说，那么我把通项公式直接写出来，这不是更加简单的静止化吗？话虽不假，有通项公式固然是最好，但是实际做题中，有很多数列是写不出显式的通项公式的。一个很简单的例子，我随便写一个递推

\[
a_{n+1}=e^{a_{n}}-a_{n}
\]

通项就写不出来了。而恰恰是这种通项不好写，或者即使写出来也很丑陃的数列，更容易考察难题。因此，递推关系的静止化是更加基本的，也是数列问题独有的。即便是像刚才的等差，等比，当初在推导它们通项的时候也是用的递推关系来推导的。

\section*{6．2．2 静态分析的一般理念}
想要静止化一个数列，其实需要的东西并不多，只要题目中给出的关系配合初始值能够把各项按照多米诺骨牌的方式各个定死，那么这个数列就是静态的。绝大部分数列问题都是静态的，那些很难的问题也不例外。但是，如何处理题目给出的这个关系，往往是最为重要的一个部分。

\section*{6．2．2．1 数列大题}
首先来说大题。高考的数列大题基本都是比小题简单的，因为作为考察写步骤和计算的大题，势必需要涉及到和具体的通项相关的内容。因此，所考查到的数列也基本是可以写出显式的通项公式的数列。这种数列显然比那些写不出通项公式的货色要让人舒服很多。

既然是这种问题，研究起来就肯定是如何去一步步把数列给静止化了。最简单的问题莫过于给你通项清晰的数列了，比如已知 \(\left\{a_{n}\right\}\) 是等差，\(\left\{b_{n}\right\}\) 是等比这种，本身就有明确的主自由度，剩下的就结合题目给定的牵制关系去算即可。这类问题和我们以前那些通过建立要素网解决的题目没有任何差别。

特殊一点的就是不知道这个数列长啥样（初始自由度是无穷），但通过关系来起到多米诺骨牌的静止化，从而可以研究的那种。这时候，控制的意识仍然十分重要，因为所有的信息都在控制的关系中。同时，充要替换则是重中之重。我们来先看一个例子，通过例子来感受数列大题的思维方式。

例 6．4 记 \(S_{n}\) 为数列 \(\left\{a_{n}\right\}\) 的前 \(n\) 项和，\(b_{n}\) 为数列 \(\left\{S_{n}\right\}\) 的前 \(n\) 项积，已知 \(\frac{2}{S_{n}}+\frac{1}{b_{n}}=2\)\\
（1）证明：数列 \(\left\{b_{n}\right\}\) 是等差数列\\
（2）求 \(\left\{a_{n}\right\}\) 的通项公式

这是全国乙卷数列大题，难度在高考数列大题中算是偏上的，正好用来讲解一下数列问题的思维方式。首先，我们要尝试建立要素关系网，不过由于数列本身的自由度比较复杂，我们不方便直接用以往的自由度算法去处理。我们要回到更根本的东西上面去，也就是——为什么要建立要素关系网？要素网的核心目的是控制，是着清楚一个系统中的要素彼此之间是什么关系，一个东西如何决定另一个东西。所以，只要能认清彼此之间谁决定谁就够了，这就是要素网的理念。

我们来看条件，\(\left\{S_{n}\right\}\) 和 \(\left\{a_{n}\right\}\) 是什么关系？是彼此控制的关系，已知 \(\left\{a_{n}\right\}\) ，我们就相当于知道了 \(\left\{S_{n}\right\}\)（直接求和），已知 \(\left\{S_{n}\right\}\) 我也可以知道 \(\left\{a_{n}\right\}\)（基本的元问题）。注意这里大家不需要去考虑具体怎么写通项表达式，只需要关注控制是否成立（好比我们以前做其他题目也没特别关心某个控制的具体函数表达式能不能写出来一样，这里不要陷人追求显式表达的误区）。随后又有个前 \(n\) 项积，道理是一样的（加减变乘除罢了）。也就是说，\(\left\{a_{n}\right\}\) 和 \(\left\{S_{n}\right\}\) 还有 \(\left\{b_{n}\right\}\) ，它们三者携带的信息量是完全一样的，彼此可以互相控制（看到这，请自行反思一下自己学前 \(n\)项和这个知识点的时候有没有自己去这么思考过，我们第五期里已经强调过按照控制和元问题的理念去学习知识点了）。

既然如此，现在唯一的已知条件

\[
\frac{2}{S_{n}}+\frac{1}{b_{n}}=2
\]

我在处理它的时候，就可以完全不考虑 \(\left\{a_{n}\right\}\)（里头都没 \(\left\{a_{n}\right\}\) 我急着考虑它干嘛），只把后两个数列拿出来关注，等后面问题问到 \(\left\{a_{n}\right\}\) 再说。那么，这个已知能干嘛？我们这相当于知道了 \(\left\{S_{n}\right\}\) 和它的前 \(n\) 项积的一个等式。注意

这个等式是对所有 \(n\) 成立的。它能不能静止化数列（因为彼此等价，所以不需要具体指出哪个数列）呢？\\
稍微想想就知道，当然可以。取 \(n=1\) 算出第一项，取 \(n=2\) 结合第一项就能算出第二项．．．．．．．肯定是可以静止化的。从这里可以看出递推关系有可能长的不是那种标准的

\[
x_{n+1}=? ? ? x_{n}
\]

这种结构，只要能起到多米诺骨牌效应的都算。退一万步讲，第二问都直接问你通项公式了，这不就告诉你能静止化了？不就相当于揭了这一家三口的底裤吗？反正这题都是要我们静止化来写通项的。

但是能静止化到一个什么程度呢？需要具体问题具体分析。这问题第一问让你说明 \(b_{n}\) 怎么怎么样，说明我们需要把式子给充要替换为纯粹 \(b_{n}\) 相关的东西。怎么就充要替换呢？原则很简单，看你替换完之后的那个式子是否仍然具有＂静止化 \(b_{n}\)＂的效力。

比如这题，一看式子里，有 \(b_{n}\) 又有 \(S_{n}, b_{n}\) 是我要研究的先留着，\(S_{n}\) 我要用 \(\left\{b_{n}\right\}\) 的信息把它控制住（刚才分析过三个数列信息量相同，这种控制的习惯和理念全程不要丢）。于是隔离问题：怎么用 \(\left\{b_{n}\right\}\)（前 \(n\) 项的积）的信息控制 \(S_{n}\) ？隔离完之后马上自问一下，然后就想起来这个元问题的处理方法了：

\[
S_{n}=\frac{b_{n}}{b_{n-1}},
\]

直接代入得到

\[
\frac{2 b_{n-1}+1}{b_{n}}=2,
\]

因为有一个 \(n=1\) 要单独讨论，所以再补加

\[
b_{1}=\frac{3}{2}
\]

我知道你已经看出这是个等差数列了，但是你先别急，先看看这个关系配合能不能实现 \(\left\{b_{n}\right\}\) 静止化？稍微看几个小的 \(n\) 的递推效果就能看出来确实可以，这样就没问题了，确实可以继续往下化简。这样做的目的是确保在逻辑上这类问题你学会了，别的题可能没等差这么好的形式。数列静止化需要的就是一个首项和一个关系实现多米诺骨牌效应，只要你变形完之后，这两样静止化需要的原料还在，那就保持了数列信息的完整刻画，你的变形才能保持信息的＂不流失不蒸发＂，也就是我们第二期讲到的充要替换（回顾第二期讲解时就强调过信（息的完整性）。

到这里其实 \(\left\{b_{n}\right\}\) 的显式已经好算了，那三个数列又是彼此决定的，所以第二问就是送的，只要三者之间彼此的转化方式很清楚（数列的基础知识点，只要你按照元问题的理解方式消化了）就毫无威胁，就不多说了，大家自己去做一下就好，最后答案是

\[
S_{n}=\frac{n+2}{n+1}, \quad a_{n}=\left\{\begin{array}{l}
\frac{3}{2}, n=1 \\
-\frac{1}{n(n+1)}, n \geq 2
\end{array}\right.
\]

从这道题可以看出，做数列大题的时候，思维方式和以往的题没有很大区别，还是控制意识（要素网），静态化的理念，元问题的做题方式在起作用，只不过虽然药没换，但汤确实换了，换成数列特有的静止化和控制方式了。只要熟悉了数列本身特有的这种思维方式，高考真题的数列大题是不成问题的，这道题已经算偏难偏灵巧的了，但是从上面的分析过程来看，稳定性完全可以达到 \(100 \%\) 。

\section*{6．2．2．2 静态把握与整体认知}
刚才的题目还算比较接近我们以往讲过的思维方式。但实际上，很多数列难题考察的更多是对一个描述性的数列的静态把握能力。这类问题在高考考场上是非常吃香的，也容易出难题。

我们来看一个例子。

例 6.5 已知数列 \(1,1,2,1,2,1,2,4,1,2,4,8,1,2,4,8,16 \cdots\) 其中第一项是 \(2^{0}\) ，接下来的两项是 \(2^{0}\) ， \(2^{1}\) ，再接下来的三项是 \(2^{0}, 2^{1}, 2^{2}\) ，以此类推。则满足下列条件的最小正整数 \(N: N>100\) 且该数列的前 \(N\) 项和为 2 的整数幂，问 \(N\)是 \(\qquad\) —。

这个问题也是高考真题，其中的数列一看就是静止的，但是叙述内容比较冗余，其项是一组一组的出现，不过规律性是很清晰的。这种数列我们很容易对它形成一个直观的认知，但是如果想用显式表达式去写通项，那会很累，不好下手。而且完全没必要。我们只需要想办法把它用尽可能舒服的形式展现出来，便于我们分析计算即可。因此，工程师的意识将在这类问题中十分重要，你必须想到一个舒服的看清整个数列的方法，也就是，如果让你把这个数列摆在我眼前，让我一目了然，你应该怎么去摆放。

这里因为是一组一组的，所以最舒服的摆放方法肯定不是一条横线摆放，而是一组一组的摆放，也就是摆成三角堆的样子：

\[
\begin{aligned}
& 2^{0} \\
& 2^{0}, 2^{1} \\
& 2^{0}, 2^{1}, 2^{2} \\
& 2^{0}, 2^{1}, 2^{2}, 2^{3}
\end{aligned}
\]

这样的话，求和也会很舒适。虽然我们关心是具体哪一项，但是只要我们有办法定位出这个三角堆里面那个点的位置，剩下的就把后续问题单独隔离出来做即可。控制的理念告诉我们，直接对三角堆下手，省去很多不必要的麻烦。

接下来，既然是要研究求和，那肯定就要做一些求和计算。对于这种能够清楚看清整个数列全貌的东西，计算时就一定要遵循这种整体的感知。比如，这里都是一堆一堆的，我计算也要以堆为单位，各堆计算完之后再考虑更精细的堆内具体位置的计算。不难看出 \(n>100\) 那必须至少加满第 13 堆（这就是 91 项），而且第 14 堆至少要加到第 10 个才行。\\
第一堆求和

\[
2^{0}=2^{1}-1,
\]

第二堆求和

\[
2^{0}+2^{1}=2^{2}-1 \ldots \ldots
\]

第 \(k\) 堆求和那自然就是 \(2^{k}-1\)（我个人的喜好：不需要急着写一堆带着未定指标 \(k\) 的抽象算式，写几个初始项，再自然地写下一般通式，有利于我更清晰地看清整个数列）。因此前 \(k\) 堆加起来就是 \(\left(2^{1}-1\right)+\left(2^{2}-1\right)+\cdots+\left(2^{k}-1\right)=\)\\
\(2^{k+1}-2-k\) 。

\[
\begin{aligned}
& 2^{0} \rightarrow 2^{1}-1 \\
& 2^{0}, 2^{1} \rightarrow 2^{2}-1 \\
& 2^{0}, 2^{1}, 2^{2} \rightarrow 2^{3}-1 \\
& 2^{0}, 2^{1}, 2^{2}, 2^{3} \rightarrow 2^{4}-1 \\
& \ldots \\
& 2^{0}, 2^{1} \cdots 2^{k-2}, 2^{k-1} \rightarrow 2^{k}-1 \\
& 2^{0}, 2^{1} \cdots 2^{t-1} ?
\end{aligned}
\]

打个比方说我现在的最小值就恰好加满第 \(k\) 堆，然后我该顺着摸下一层（这种加一加，整体思考一下，然后往上不断摸的思维方式对于静态数列问题是比较重要的），同时一边摸一边试着对已知做充要替换，试着去摸我要的那个 \(k\) 。我现在是要 2 的整数幂，我得摸到啥样的 \(k\) ？我就停下来想，假如我触摸到最后一个完整堆之后，我还需要再加多少？（思维焦点放到＂加多少＂上面来）有没有可能不需要多加了？那就正好是

\[
2^{k+1}-2-k,
\]

它得是 2 的整数幂，但这不可能，指数到了 13 差不多的位置之后，已经远远比 \(2+k\) 这种货色大了，那这个结果会很接近 \(2^{k+1}\) ，减掉一点，离 \(2^{k}\) 还远着呢（充要替换成数学语言就是 \(2^{k+1}-2-k>2^{k}\) ），不行不行！所以肯定是往上再摸几项达到一个 2 的幂次（这个时候自然出现了两个问题，具体再摸几项，达到 2 的多少幂次）。

再摸几项？这是关于确定具体位置的，是关键信息，需要更多信息来支撑它才行，所以往后稍稍。先想达到 2 的多少幂次？（思维焦点放过来）到下一堆（第 \(k+1\) 堆）之后先慢慢往前摸，初始的几项很小，都是 \(1,2,4\) 这种玩意（静态数列的直观把控）, \(2+\mathrm{k}\) 会慢慢被抵消掉，好像摸不了几项就有希望全部抵消！那就是 \(2^{k+1}\) 了不成？不一定，没准再多摸一些，摸到 \(2^{k+2}\) 甚至再大了呢？（对这些数的大小，以及层的大小进行直观的评估）再摸摸看看，这差距可不得了，你即便一直摸，摸完整个这第 \(k+1\) 层，充其量也就能到 \(2^{k}+2-2-(k+1)\) ，我们刚才不是说假设就摸完整的第 \(k\) 层吗？你都摸完下一层了还没到，那只能说明必须是 \(2^{k}+1\) 了。所以你的最小值 \(N\) 如果是包含完整的 \(k\) 层，那么第 \(k+1\) 层只能摸一点点，摸到正好抵消 \(2^{k}+1-2-k\) 里面的 \(-2-k\) ，就正好结束。这样的话，第二个问题，多少幂次，就得到了解答。

上一段的最后一句话非常直观的展示了摸的要求。自问一下是充要替换吗？是的，刚才的一通分析都是双向考虑的，没有造成逻辑上的信息遗漏。那就好，抛弃一波，我们现在只需要看上一段最后一句话了。现在就要摸完 \(k\) 层以后再摸 \(t\) 个， \(1 \leq t \leq k\) 且第 \(k+1\) 层的前 \(t\) 个加起来正好是 \(2+k\) ，这不就相当于再摸几项的刻画了吗！接下来显然继续充要替换，那这就等价于 \(2^{\prime}-1=2+k\) ，也就等价于 \(3+\mathrm{k}\) 是 2 的幂次。（现在信息都被等价完了，而且看上去很适合计算）我们希望 \(k\) 尽量小，这样需要摸的层就尽量少，\(N\) 就小，还要保证 \(k \geq 13\) ，那最小的就是 \(3+k=16\) ，然后再摸 \(t\) 项，\(t=4\) ，可是这样的话总项数还是不够 100 ，差了一点；那我们就不得不取下一个 2 的幂次： \(3+\mathrm{k}=32\) ，这时候 \(2^{\prime}=32, t=5\) ，总共的项数为 \((1+2+\cdots 29)+5\) ，题目就做完了。

从这个题可以看出，题目本身就是一个纯静态的数列分析问题。虽然通项不好写（也没意义），但是整个数列的外貌是栩栩如生的展现在我们面前的。这种问题（包括一些奇偶分开的数列，另一些类似的奇怪描述的数列）都是一个路子：首先，你需要用好工程师的理念（这里的工程师与我们以前讲过的要素网建立的工程师有点差别，后者是一个梳理主自由度控制系统的过程，是造机械，是重工业；前者是将一个复杂庞大的静态系统给整理清楚，使之井井有条，是收拾房间，是服务业，不过现在似乎也有了打扫房间的机器人了），把数列以一个好算的形

式整理利索，然后分析题目研究的要素（求和，具体哪一项之类的）。注意这里的分析要思维与计算同时进行，当你发现一个东西想不好想的时候，就把能算的要素算出来看看。数列的很多信息是藏在式子里面的，在边想，边摸，边算的过程中，你才能让对整个数列的认知越来越清晰，进而便于对题目要求进行充要替换。

我们再看一个更难的例子。

例 6.6 数列 \(\left\{a_{n}\right\}\) 共 \(M\) 项 \((M>5)\) ，对于任意正整数 \(k \leq M, a_{k}+a_{M+1-k}=0\) ，当 \(n \leq \frac{M}{2}\) 时，\(a_{n}=\frac{1}{2^{n}}\) ，记 \(\left\{a_{n}\right\}\) 的前 \(n\) 项和为 \(S_{n}\) ，则下列说法中正确的有（多选）

A．若恒有 \(S_{n} \leq \frac{1023}{1024}\) ，则 \(M \leq 20\)

B．\(\left\{a_{n}\right\}\) 中可能存在连续 5 项构成等差数列

C．对任意小于 \(M\) 的正整数 \(p, q\) ，存在正整数 \(i, j\) ，使得 \(a_{i}+a_{j}=S_{p}-S_{q}\)

D．对 \(\left\{a_{n}\right\}\) 中任意一项 \(a_{r}\) ，必存在 \(a_{s}, a_{t}(s \neq t)\) ，使得 \(a_{r}, a_{s}, a_{t}\) 按一定顺序排列可以构成等差数列

老规矩，先看看数列是否静止。 \(M\) 是一个未知的对象，但是如果 \(M\) 定了，那么整个数列就定了吗？已知是两个，一个直接给了近半数的数列项，另一个我们来做一下试验。 \(k\) 是个被封装的，相当于我们可以随便跑它，随便带几个 \(k\) 就可以看到 \(a_{1}+a_{M}=0, a_{2}+a_{M-1}=0 \cdots\) 就是说这个数列是两头反过来的。这样的话确实，定了 \(M\) 之后，数列就完全定下来了。可以取几个 \(M\) 看看，\(M=6\) 就是

\[
\frac{1}{2}, \frac{1}{4}, \frac{1}{8},-\frac{1}{8},-\frac{1}{4},-\frac{1}{2}
\]

而 \(M=7\) 的话，已知是说 \(n \leq 3.5\)（3）时是等比，\(n=5,6,7\) 时顺势翻过来，那 \(n=4\) 呢？刚才的＂两头反＂的直观理解立刻就让我们看到中间 \(a_{4}=0\)（相当于取 \(k=4\) 代人）。那这个数列该如何工程师一般去理解？其实就是先写一段等比数列（这一段至少要写到 3 项），然后要么直接对称翻过去（ \(M\) 是偶数），要么后面加个 0 再翻过去（ \(M\) 是奇数），这个理解已经很舒服了。下面就可以开始分析选项了。

A 选项看的是求和，那我们把等比外加（可能添 0 ）翻转这一排拿过来看，很明显，你从头一直往后摸着加，把等比数列加完达到最大（可能添个 0 ），然后就开始慢慢把加的项又减去了，所以 \(S_{n}\) 最大就是

\[
\frac{1}{2}+\cdots+\frac{1}{2^{n}}=1-\frac{1}{2^{n}} \leq \frac{1023}{1024}
\]

这就得到 \(n \leq 10\) ，也就是正的等比部分最多加到 10 个（或者再添个 0 ），然后一下子把那最多 10 个的项翻过去，所以 \(M\) 最多是 21 个。如果你漏掉了 0 这一项的可能的话，回头看看这一段，添 0 一直作为括号后缀而放在脑子里，这就是时刻有着充要替换和逻辑完整性的习惯，已经从经验分享第二期就开始强调的东西了。

B 选项，可能存在连续 5 项等差，那我把刚才的直观拿出来，等比对折翻转（可能添 0 ），我来看看里头有没有连续 5 项等差。很明显，要分成奇数偶数讨论了。这 5 项有可能在哪出现？整体是一个对称结构（工程师意识搭建起来的直观），两侧都是等比下降，不断折半，感觉折半就不太可能等差啊，比如从 \(\frac{1}{2^{t}}\) 折到 \(\frac{1}{2^{t+1}}\) ，下一项想要是等差就得是 0 了，不能再折了。 0 是个必要条件，这么一看， 0 的话，那只有考虑 \(M\) 是奇数，\(\frac{1}{2^{t+1}}\) 后接 0 ，

只有这样了，再往后是

\[
-\frac{1}{2^{t+1}},-\frac{1}{2^{t}},
\]

正好是等差！那没问题了，当然是可能有。

C 和 D 是重头戏，我们看一下如何进行思维的分析。 C 是标准的封装命题结构，所有的参数 \(p, q, i, j\) 全被封装，所以矛头还是指向数列本身。我们按照研究封装命题的一般思路去走，最外面的一层是任意 \(p, q\) ，我们先随便点一组最外层的 \(p, q\) ，然后看里面那一层：＂存在 \(i, j\) 使得 \(a_{i}+a_{j}=S_{p}-S_{q}\)＂，里面这一层非常简单，就是一个选两项给加一下的问题，所以问题关键在于我们随便点的这个 \(S_{p}-S_{q}\) 到底是多少，所以问题归结为计算 \(S_{p}-S_{q}\)的形式。 \(p, q\) 是任意小于 \(M\) 的，其余什么信息都没有，我们自然要设问：如何计算这个玩意？

很明显，这是一个经典的元问题，其实就是连续若干项相加。但是稍微一遍历 \(p, q\) ，就会理解肯定需要先讨论一下二者大小关系：如果 \(p=q\) ，那就是 0 ，（思维焦点过来）显然随便找一组互为翻转的 \(a_{i}\) 和 \(a_{j}\) 即可；\(p>q\)和 \(p<q\) 是一样的，因为这相当于把 \(S_{p}-S_{q}\) 符号换了一下，如果 \(p>q\) 能找到 \(i, j\) ，那么相应的 \(p<q\) 就是 \(-\left(S_{q}-S_{p}\right)\) ，把之前的构造里的 \(i, j\) 都给翻转到另一侧即可（工程师意识带给你的，数列整体的翻转的直观，要一直放在脑子里）。所以我们只需要考虑 \(p>q\) ，也就是 \(a_{q+1}+\cdots+a_{p}\) ，也就是第 \(M\) 项之前，不包含第一项的任意一段截出来求和（这是充要替换，我们拿出来，抛弃前面的 \(S_{p}-S_{q}\) ）。

把整个数列铺开放在眼前，心里想着我要截取连续的一段来求和，求出表达式来再看能不能找到 \(i, j\) 。截的方式就跟上一道题一样，一边摸一边看着往上加。比如从 \(\frac{1}{2^{t}}\) 开始往后截，截到 \(\frac{1}{2^{t+p}}\)（瞎写两个数），那求和就是一段正项等比数列；再往后截可能截到负的，那就从 \(\frac{1}{2^{n}}\) 开始往回走，被那些负的又给抵消掉，这样抵消就相当于原本摸的时候少摸一点，还是相当于一段等比数列加起来。如果一直抵消到 \(-\frac{1}{2^{t}}\) ，那啥也没了，不就是 0 呗！继续往下截，正项全部抵消完，变成负的，那也是负的一段等比数列相加呗！刚才分析 \(p<q\) 还分析过变成负的无所谓，因为数列本体就可以让你选择 \(i, j\) 时把符号颠倒过去（从这又看到先前整体静态的印象有多重要）。所以我完全可以只考虑一段正等比数列的求和：比如考虑从 \(\frac{1}{2^{t}}\) 加到 \(\frac{1}{2^{t+p}}\)（这里 \(t \geq 2, t+p \leq n\) ），加起来是

\[
\frac{1}{2^{t-1}}-\frac{1}{2^{t+p}}
\]

你看，正好是数列中的某两项求和吧！这不就找到了！所以 C 是对的。

D 也是一样的，封装的结构和 C 差不多，先随便点一个 \(a_{r}\) ，然后看看能不能选出制造等差的 \(s \neq t\) 。那还是去遍历一下 \(a_{r}\) 看看，如果选的 \(a_{r}=0\) ，那就显然能找到。之后根据数列整体的翻转特性，只需要看

\[
a_{r}=\frac{1}{2^{m}},
\]

那这时候能选出合适的 \(s \neq t\) 吗？那就拿过整个数列来找找，如果里面有 0 ，那选 0 和 \(-\frac{1}{2^{m}}\) 就完事，不然的话，我们试着讨论一下（这时候要讨论的数列不含 \(0, a_{r}=\frac{1}{2^{m}}\) 是其中一员）。

还是一样的，我们就顺着去摸就行，因为三项的等差数列是两个自由度，其中已知一项，剩一个自由度，所以先试着找一个 \(s\) ，然后就能基本锁定最后一项的数值了，再去筛选这些数值是不是在数列里即可（实际上你也不可能同时摸 \(s\) 和 \(t\) ，按照冻结的思维理念肯定也是先去摸其中一个再看另一个）。比如我先往后摸，可能下一项是 \(\frac{1}{2^{m+1}}\) ，也可能直接这就是等比数列这正的一半的末尾（那下一项就是 \(-\frac{1}{2^{m}}\) ）。如果是 \(\frac{1}{2^{m+1}}\) ，那么 \(a_{t}\) 能是啥？把数列铺开然后算一下，就会发现 \(a_{t}\) 要么是 0 （现在没有），要么是（中间插不上，只好插在左边）

\[
\frac{2}{2^{m}}-\frac{1}{2^{m+1}}=\frac{3}{2^{m+1}}
\]

没这一项。那继续摸，如果继续摸到 \(\frac{1}{2^{m+p}}(p \geq 2)\) 的话，你会发现，如果 \(a_{t}\) 插在右边，那么 \(a_{t}\) 的值就会是

\[
\frac{2}{2^{m+p}}-\frac{1}{2^{m}}=\frac{2-2^{p}}{2^{m+p}}
\]

所以只要 \(p=2\) 就有

\[
\frac{1}{2^{m}}, \frac{1}{2^{m+2}},-\frac{1}{2^{m+1}}
\]

是等差数列，也就是说，只要往后摸两项还是正的，那么往后摸两项就能选出合适的 \(a_{t}\) 来！但是刚才也说过了，如果往后还没摸到两正项就进了负区域，怎么办？那往回摸两项也行

\[
\left(\frac{1}{2^{m-2}}, \frac{1}{2^{m}},-\frac{1}{2^{m-1}}\right)
\]

！那如果往回也摸不到两项呢？那就相当于往前的正部分最多一项，往后的正部分最多也一项，那只有一种可能了，就是 \(M=6\) ，考虑 \(a_{r}=\frac{1}{4}\) ，这时候一共就这么几项，直接构造就行了！\(a_{s}\) 一共六种可能，挨个数一遍！最后发现，取

\[
\frac{1}{4},-\frac{1}{8},-\frac{1}{2}
\]

即可。证明完毕！

这道题比刚才的例 2 还要难，需要分析的东西很多，但是从我绝大部分分析都是文字这一点来看，在脑子里过的时间是主要的，真正计算量反而很小，因此这道题谁的逻辑基本功越扎实谁做的就越清楚，就越快。用到的思维方式和刚才的例 2 是一样的，首先把自己作为一个工程师，形成一个对数列的整体印象，知道一段等比能完全定死整个数列即可。然后边分析边计算，计算应该辅佐思考，遇到思考不动的东西再算算看看。多去用遍历的思维方式去看问题（也就是摸项，沿着 \(n\) 去摸，边摸边观察）。大概需要注意的就是这些，另外我们以前讲过的那些思维方式，隔离，思维焦点，设问等等都会有用。如果平时对我们讲过的思维方式已经比较熟悉了，那么这类数列问题看视频的讲解熟悉一下其特有的思考方式，就可以了。另外，由于思考的过程往往比较长，所以可以把一些静态直观（比如例 2 的三角堆，例 3 的翻转）都给画出来摆在试卷显眼的位置，这样方便大脑不断强化它们。

\section*{6．2．3 递推数列问题}
刚才讲过的静态数列分析一般都是我们熟悉的数列或者整体直观比较好把握的数列，它们比较方便我们进行整体的计算和分析。但是第一部分里也说过。更多的数列我们是很难像例 2 ，例 3 这样，直接看透它们整体的样貌的。往往它们只有一个递推关系用来满足静止化。这种问题的难度往往更高，也更难用一般化的手段处理，常常需要特事特办，需要一定的做题量积累，甚至需要一定的反应和灵感。这类题目就更多的面向冲击 140 分以上的考生了。不过，养成一些良好的思考习惯和充要替换习惯也会有助于我们更好的处理它们。我们通过一些例子来展示一些好的思维方式和习惯以提高解题稳定性，但即便如此，这类题目也无法保证稳定性达到 \(100 \%\) ，同学们量力而行。

例 6.7 数列 \(\left\{a_{n}\right\}\) 满足 \(a_{1}=a, a_{n+1}=3 a_{n}-a_{n}^{2}-1\) ，则下列说法正确的有（多选）

A．若 \((a-1)(a-2) \neq 0\) ，则数列 \(\left\{a_{n}\right\}\) 单调递减

B．若存在无穷多个正整数 \(n\) ，使得 \(a_{n}+1=a_{n}\) ，则 \(a=1\)

C．当 \((a-1)(a-2)>0\) 时，\(\left\{a_{n}\right\}\) 不存在最小值

D．当 \(a=3\) 时,\(\frac{1}{a_{1}-2}+\frac{1}{a_{2}-2}+\cdots \frac{1}{a_{n}-2} \in\left(\frac{1}{2}, 1\right]\)

题目的已知条件可以看到，数列的递推关系是明朗的，所以只要首项一个自由度即可定死整个数列。先来看 A ，给了一个首项的范围，说这时候数列单调递减。是不是单调递减我们肯定看不出来，因为只有一个递推，数列长啥样我们也不清楚。遇到这种情况我们第五期是怎么说的？应该主动问自己我们能干什么。我们能干的只有一个，就是思考递减怎么处理。（元问题隔离出来）怎么就递减？那当然是恒有 \(a_{n+1}<a_{n}\) ，这个是前后项的关系，根据控制关系，可以充要替换成 \(3 a_{n}-a_{n}^{2}-1<a_{n}\) ，这是一个关于 \(a_{n}\) 的不等式，本着第五节的＂简化即胜利＂的原则，肯定要把它求解掉，解得．．．．．．．怎么解得 \(a_{n} \neq 1\) ？？那我就懂了， A 是说，只要首项不是 1 或 2 ，那么整个 \(\left\{a_{n}\right\}\) 没有一项等于 1 。这可比原本的要简单，直接落实到项上。

那么，事实确实如此吗？我们判断整个 \(\left\{a_{n}\right\}\) 咋样肯定不好判断，所以玩个逻辑游戏，考虑它的逆否命题 （或者说用反证法）：如果有某一项是 1 ，看看它首项是不是只能是 1 或者 2 。这个好办，因为某一项是 1 的话，可以往回摸，比如 \(a_{n}=1\) ，那解方程得到 \(a_{n}-1\) 是 1 或 2 ；如果 \(a_{n}-1=1\) ，跟刚才一样了，如果 \(a_{n-1}=2\) ，那么 \(a_{n-2}^{2}-3 a_{n-2}+3=0\) ，无解！当然，如果你倒到最开头了，取 \(a_{1}=2\) 当然没问题，但只要没倒到头，那就必须选择 1 继续往回倒．．．．．．．那自然 A 就是对的了。可以看到， A 选项的分析过程一直是在充要替换，其实只要大家在分析时保持了严谨性，并且有明显简化的步骤（比如这里解出 \(a_{n} \neq 1\) ），那么根据第五期的做题原则，大家就坚定的往下尝试即可。

B 选项，老规矩，充要替换，\(a_{n+1}=a_{n} \Leftrightarrow a_{n}=1\) ，刚才刚刚分析过，有 1 就必须一路的 1 倒回去，开头（首项）可以是 1 也可以是 2 ，这个时候当然是满足无穷个 \(n\) 的条件（ 1 往前走也一直是 1 ），所以 \(a=2\) 也可能，B 错误。

C 选项，当那个范围的时候，最小值不存在．．．．．．．又不会了，那就自问，怎么就最小值不存在了？什么样的数列找不到最小值？想想常见的，等差等比，等差如果递减，就没有最小值；等比如果递减也没有．．．．．．．不管怎么说好像只要跟＂减＂扯上关系就容易没有哎。这个减我好像在哪里刚见过．．．．．．． A 不就是吗？？\((a-1)(a-2) \neq 0\) 的时候它就递减，那现在当然也满足这个条件啊，一个数列严格单调递减，它怎么可能有最小值呢？C 当然对。

D 是最难的选项。碰到这种陌生数列的求和问题的时候，记住，我们只有两个方法，一个是裂项，或者说把求和想办法给求出来，把很多项变成很少项；另一个是放缩，具有较高的技巧性和经验依赖性。有的题目可能需要二者混用。这道题考虑的是

\[
\left\{\frac{1}{a_{n}-2}\right\}
\]

的求和，看外表完全没啥思路，这个时候和大家分享一个好习惯，就是尽量把数列整理成看着比较舒服的形式，哪怕舒服一点都比原来的强。比如这里 \(a_{n}-2\) 看着就很别扭，所以我换元，令 \(b_{n}=a_{n}-2\) ，这样的话求和的数列就变成了 \(\frac{1}{b}\) ，这样看着能赏心悦目一些。这样的话 \(\left\{a_{n}\right\}\) 和 \(\left\{b_{n}\right\}\) 彼此控制，当然也是静止的。首项 \(b_{1}=1\) ，我们需要把递推关系也迁移过来从而静止化 \(\left\{b_{n}\right\}\) ，这个计算很简单，大家自行练习。最后得到的是

\[
b_{n+1}=-1-b_{n}-b_{n}^{2}
\]

到这之后就要选方法了，尝试直接裂项把求和给算出来变成好看的一两项还是考虑放缩？好像现在看也没啥可放缩的。那就试试裂项吧。裂项式子的唯一可能来源就是那个递推关系，我们尝试观察。首先肯定需要凑出

倒数形式，也就是 \(\frac{1}{b_{n}}\) 。直接倒过来变成

\[
\frac{1}{b_{n+1}}=-\frac{1}{1+b_{n}+b_{n}^{2}}
\]

吗？右边底下没法因式分解，干啥都干不了，走不通；还有其他能凑出 \(\frac{1}{b_{n}}\) 的途径吗？右边有两个带 \(b_{n}\) 的部分单独拿出一个来弄？那另一个和 \(b_{n+1}\) 缠在一起，更乱了；我看那里有个 \(b_{n}+b_{n}^{2}\) ，把 \(b_{n}\) 整体提出来，就有个 \(b_{n}\left(1+b_{n}\right)\)了，然后还出来一个 \(b_{n}+1\) ，这下出来一组相似的结构：\(b_{n}+1\) 了（别忘了裂项的目的就是把 \(\frac{1}{b_{n}}\) 裂成两个相同结构的差），有门！于是尝试整理：

\[
\frac{1}{b_{n+1}+1}=-\frac{1}{b_{n}\left(b_{n}+1\right)}=-\left(\frac{1}{b_{n}}-\frac{1}{b_{n}+1}\right),
\]

得

\[
\frac{1}{b_{n}}=\frac{1}{b_{n}+1}-\frac{1}{b_{n+1}+1},
\]

这样就完成了裂项，这是巨大的简化，肯定我们走的路是对的。于是原命题等价为

\[
\frac{1}{b_{1}+1}-\frac{1}{b_{n+1}+1} \in\left(\frac{1}{2}, 1\right],
\]

把数值代人，进一步替换为

\[
b_{n}+1 \leq-3 \Leftrightarrow a_{n+1} \leq-1
\]

到这已经简化很多了，我们继续分析。还是老样子，往前摸摸试试。如果希望 \(a_{n+1} \leq-1\) ，解不等式（充要替换）得到 \(a_{n} \leq 0\) 或 \(a_{n} \geq 3\) 。这不正好，第一项是 3 ，第二项可以算出来就是 -1 ，然后 -1 还在那个区间里，所以再摸一项还是不大于 -1 ，继续往下还在这个区间里，于是下一项还不大于－1．．．．．．以此类推，可知后面的全不大于 -1 ，还是往前不断摸的这种思维方式证明了 D 选项成立。

这样的题目是很难保证 \(100 \%\) 做出的，只能通过科学的思维方式来尽量提高稳定性。这道题给我们体现出来很多曾经讲过的思维方式的应用，除了隔离，分解出元问题，充要替换维持信息量等，还有一些应对难题的思维方式。比如，＂简化即胜利＂判断操作是否可行的理念，反面思考，也就是反证法（这种思考方式越是难题越要重视）以及数列特化的遍历法，也就是摸项的思维方式。这种摸项遍历对于静态数列都有用（包括这种没法一眼望穿的），这是因为数列这一对象的特殊性（ \(n\) 可以滑动）。上篇里面我们讲那些静态数列分析问题时候，几乎每个题都在用，现在仍然在用，其重要性可见一斑。

我们再来看一道更棘手的问题。

例 \(6.8 \triangle A_{n} B_{n} C_{n}\) 是直角三角形，\(A_{n}\) 是直角，角 \(A_{n}, B_{n}, C_{n}\) 所对边分别是 \(a_{n}, b_{n}, c_{n}\) ，三角形面积为 \(S_{n}\) 。若

\[
b_{1}=4, c_{1}=3, b_{n+1}^{2}=\frac{a_{n+1}^{2}+c_{n}^{2}}{3}, c_{n+1}^{2}=\frac{a_{n+1}^{2}+b_{n}^{2}}{3}
\]

则（多选）\\
A．\(\left\{S_{2 n}\right\}\) 是递增数列\\
B．\(\left\{S_{2 n-1}\right\}\) 是递减数列\\
C．\(\left\{b_{n}-c_{n}\right\}\) 存在最大项\\
D．\(\left\{b_{n}-c_{n}\right\}\) 存在最小项

这道题较为困难，已知条件上来就错综复杂，而且涉及到的数列还很多。继续沿用我们的思维方式。首先 \(b_{1}, c_{1}\) 都已知，代表第一个三角形静止，然后令 \(n=1\) ，代人递推得到两个方程，这里相当于方程中的未知数只有\\
\(a_{2}, b_{2}, c_{2}\) ，而直角三角形两个自由度，这方程足以定死第二个三角形．．．．．以此往复，很容易看出来已知条件能够完全实现静止化。

既然如此，我们知道数列是静止的，那就好办了。首先，从我们之前讲的思维方式来看，已知条件并不舒服，因为已知里面前项控制后项的关系不明朗，所以我们肯定是要先手动把控制关系给整明朗。我们之前说过，标准的控制关系应当是＂\(x_{n+1}=? ? ? x_{n}\)＂形式的（上一道例 4 就是标准的这种形式），好比函数关系一样，后项完全被前项表示。这道题目很明显，是可以做到这一点的（刚才试着判断是否静止的时候代入 \(n=1\) 时就能发现，就是在解一个方程组罢了），实际上所有靠递推关系实现静止化的数列都应该能做到这一点。这道题的这种需求就比较明显，不这么做的话我们就看不清楚前项如何控制的后项。

怎么操作呢？控制就是前项决定后项，那我们要做的就是把 \(a_{n}, b_{n}, c_{n}\) 全当已知（静止，相当于已知前项），然后求后项，也就是解方程组

\[
\left\{\begin{array}{l}
b_{n+1}^{2}=\frac{a_{n+1}^{2}+c_{n}^{2}}{3} \\
c_{n+1}^{2}=\frac{a_{n+1}^{2}+b_{n}^{2}}{3}
\end{array}\right.
\]

注意直角三角形两个自由度，求解边的时候需要勾股定理条件，我们选取主自由度为 \(b_{n+1}, c_{n+1}\)（很好理解，画图法告诉我们直角三角形最舒服的主自由度就是它们，而且四个选项都围着他们转），先求这俩。因此把 \(a_{n+1}^{2}=\) \(b_{n+1}^{2}+c_{n+1}^{2}\) 换进去，就变成了关于 \(b_{n+1}, c_{n+1}\) 的方程组

\[
\left\{\begin{array} { l } 
{ b _ { n + 1 } ^ { 2 } = \frac { b _ { n + 1 } ^ { 2 } + c _ { n + 1 } ^ { 2 } + c _ { n } ^ { 2 } } { 3 } } \\
{ c _ { n + 1 } ^ { 2 } = \frac { b _ { n + 1 } ^ { 2 } + c _ { n + 1 } ^ { 2 } + b _ { n } ^ { 2 } } { 3 } }
\end{array} , \text { 解得 } \left\{\begin{array}{l}
b_{n+1}^{2}=\frac{b_{n}^{2}+2 c_{n}^{2}}{3} \\
c_{n+1}^{2}=\frac{2 b_{n}^{2}+c_{n}^{2}}{3}
\end{array}\right.\right.
\]

顺便我们还可以把 \(a_{n+1}\) 算出来，控制的意识下很简单，就直接用 \(b_{n+1}, c_{n+1}\) 控制它：加一下得到 \(a_{n+1}^{2}=b_{n}^{2}+\) \(c_{n}^{2} \ldots \ldots\) 我马上就注意到，这不就是相当于 \(a_{n+1}=a_{n}\) 吗？这不就说明 \(a_{n}\) 是个常数列？这真是意外之喜！不过这个＂意外之喜＂也没有那么唐突，基本就是在我们写控制关系的时候顺理成章出来的，由此，再次可见控制的重要性。

注意一个做题经验，当你遇到这种＂意外之喜＂的时候，不要轻易放过，好东西肯定不可能徒有其表，肯定能派上用场。有了 \(a_{n}\) 定值，上面的结果我们还可以进一步简化。

\[
\left\{\begin{array}{l}
b_{n+1}^{2}=\frac{b_{n}^{2}+2 c_{n}^{2}}{3} \\
c_{n+1}^{2}=\frac{2 b_{n}^{2}+c_{n}^{2}}{3}
\end{array}\right.
\]

这玩意是我们费劲算出来的，配合首项就承载了数列的全部信息，而且是直白的控制，肯定是好的，我们下面肯定要对它操作分析。 \(a_{n}\) 定值相当于

\[
b_{n}^{2}+c_{n}^{2}=a_{n}^{2}=25
\]

这是我们在控制 \(a_{n}\) 时得到的结果，相当于把两组式子相加。两组式子的全部信息肯定没有彻底包含在这样一个和式里面，因为与两组式子具有同等逻辑效力（可充要替换）的话仍需要两组式子（早就讲过的自由度的观点）。

因此，想要这个 \(b_{n}^{2}+c_{n}^{2}=25\) 发挥作用，必须把它和原本两组式子再配合，发生新的化学反应。比如我把

\[
\left\{\begin{array}{l}
b_{n+1}^{2}=\frac{b_{n}^{2}+2 c_{n}^{2}}{3} \\
b_{n}^{2}+c_{n}^{2}=25
\end{array}\right.
\]

放一起，接下来怎么搞？这么单独把这俩放一块就很明显能看到，把 \(c_{n}^{2}=25-b_{n}^{2}\) 往里带就得到

\[
b_{n+1}^{2}=\frac{50-b_{n}^{2}}{3}
\]

\(c\) 居然直接给干没了！那同理，我还可以类似得到 \(c_{n}^{2}\) 的控制关系

\[
c_{n+1}^{2}=\frac{50-c_{n}^{2}}{3}
\]

这下控制看的更清楚了，两个数列被分开了，各自控制各自的，这就很好（形式还一致，很美观）。不过这一步有点灵感的成分，如果没看出来的话，这两个数列的控制就仍然纠缠在一起。我们下面讲两种接下来的走向，一种是能看出来的，一种是没看出来的。

首先说一下看出来的，也就是用

\[
\left\{\begin{array}{l}
b_{n+1}^{2}=\frac{50-b_{n}^{2}}{3} \\
c_{n+1}^{2}=\frac{50-c_{n}^{2}}{3}
\end{array}\right.
\]

继续往下做的（这当然足以静止化，所以信息量完全保留了）。我们接下来还能干什么？把两个数列分开隔离看看，一分开就很容易发现 \(\left\{b_{n}^{2}\right\}\) 和 \(\left\{c_{n}^{2}\right\}\) 都是标准的

\[
x_{n}+1=p x_{n}+q
\]

形式的数列，都是有求通项的方法的，你当然可以去直接算通项出来，都是程序性的过程（默写型步骤），看到这里基本就知道这题结束了。

A 和 B 都看的是面积分奇偶的单调性，我们当然可以算出两个数列然后去直接搞。这里再讲一种算法。因为这里我们手的东西全都是边的平方，干脆也直接充要的看面积的平方，再去掉多余的系数，就看 \(\left\{b_{n}^{2} c_{n}^{2}\right\}\) 得了。令 \(x_{n}=b_{n}^{2} c_{n}^{2}\) ，那么一样，用控制关系去处理：

\[
x_{n+1}=\frac{50-b_{n}^{2}}{3} \frac{50-c_{n}^{2}}{3}=\frac{1250+x_{n}}{9},
\]

首项 \(x_{1}=144\) 这仍是 \(x_{n+1}=p x_{n}+q\) 形式（定式套路，解方程 \(p x+q=x\) 得到 \(x_{0}\) ，则 \(\left\{x_{n}-x_{0}\right\}\) 是等比数列），直接配凑一下去算得了：

\[
x_{n+1}-\frac{625}{4}=\frac{1}{9}\left(x_{n}-\frac{625}{4}\right),
\]

于是很容易算出

\[
x_{n}-\frac{625}{4}=\left(-\frac{49}{4}\right) \cdot\left(\frac{1}{9}\right)^{n-1},
\]

单调性的关键在后面这个指数部分，容易看出它是增的，所以面积应该是一直递增的才对（ AB 这奇偶分了个寂寞）。

C 和 D 都在考虑 \(\left\{b_{n}-c_{n}\right\}\) ，这玩意有没有最值呢？这玩意当然也是静止的，而且 \(b_{n}\) 和 \(c_{n}\) 的关系早就有了，平方和是定值，所以 \(b\) 越大 \(c\) 越小，\(b-c\) 也就越大，这是单调的。所以 \(b_{n}-c_{n}\) 若能取得最值，相应的 \(b_{n}\) 也能

取得最值，反之亦然，这是一个基本的单调性的思考。所以我们只需要关注 \(b_{n}\) ，或者说，\(b_{n}^{2}\) 。那么，干脆还是跟刚才一样，直接算得了，还省事。

\[
b_{n+1}^{2}-\frac{25}{2}=-\frac{1}{3}\left(b_{n}^{2}-\frac{25}{2}\right),
\]

然后算出

\[
b_{n}^{2}=\frac{25}{2}+\frac{7}{2} \cdot\left(-\frac{1}{3}\right)^{n-1}
\]

这是分奇偶的上下震荡的数列，指数部分会越来越靠近 0 ，稍微画一下图就能看出来它确实两个最值都存在（就是震荡最开始的部分）。

上面的做法基本都是默写步骤，所以把这两个数列给解开确实给我们省了一些事。我们再讲一下，如果一开始没想到的话，那仍然用

\[
\left\{\begin{array}{l}
b_{n+1}^{2}=\frac{b_{n}^{2}+2 c_{n}^{2}}{3} \\
c_{n+1}^{2}=\frac{2 b_{n}^{2}+c_{n}^{2}}{3}
\end{array}\right.
\]

来做。 A 和 B 的面积，肯定还是要考虑

\[
x_{n}=b_{n}^{2} c_{n}^{2}
\]

乘一下得到

\[
x_{n+1}=\frac{5 x_{n}+2\left(b_{n}^{4}+c_{n}^{4}\right)}{9}
\]

这样的话，要研究的 \(x_{n}, x_{n+1}\) 肯定要留着，于是思维焦点就自然地放到那个多余的 \(b_{n}^{4}+c_{n}^{4}\) 上面去。我们说了，好东西有 \(x_{n}=b_{n}^{2} c_{n}^{2}\) ，是个乘积形式，还有什么好东西？那当然是 \(b_{n}^{2}+c_{n}^{2}=25\) 定值，那么两个东西 \(\left(b_{n}^{2}\right.\) 和 \(\left.c_{n}^{2}\right)\) 乘积是好东西，和也是好东西，我现在想研究它们的平方和（就是四次方），这么一隔离，一想，就是个经典的问题了：

\[
b_{n}^{4}+c_{n}^{4}=\left(b_{n}^{2}+c_{n}^{2}\right)^{2}-2 b_{n}^{2} c_{n}^{2},
\]

本质是

\[
x^{2}+y^{2}=(x+y)^{2}-2 x y,
\]

初中生都会。这样 \(b_{n}^{4}+c_{n}^{4}\) 也成了好东西了，全给带换进去算出来就是之前得到的

\[
x_{n+1}=\frac{1250+x_{n}}{9},
\]

然后又回去了。\\
至于 C 和 D ，因为现在两个数列纠缠在一起，所以不太容易直接想到单独分析 \(b_{n}^{2}\) 。我们知道 \(b_{n}^{2}+c_{n}^{2}=25\) ，这玩意是我们加出来的，刚才就分析过了信息全在那个递推关系里，这种意外之喜只是我们从递推关系里面榨取出来的，信息量肯定是有所下滑的，所以原递推关系不能不看（这种思维方式与逻辑是我们一直都在强调的，大家应该做题有这样的意识）。加一下出来一个关系，如果能再补一个关系就好了，这样，两个关系就没有信息量的流失。两个式子，已经加过了，那么再减一下，不就行了？两个式子一加一减得到两个新式子本身就保持了充要（二对二，两组式子可以互推，类比方程组加减消元）。再加上递推关系的两个式子

\[
\left\{\begin{array}{l}
b_{n+1}^{2}=\frac{b_{n}^{2}+2 c_{n}^{2}}{3} \\
c_{n+1}^{2}=\frac{2 b_{n}^{2}+c_{n}^{2}}{3}
\end{array}\right.
\]

的结构本身就很对称，一加一减也看着舒心，于是作差：

\[
b_{n+1}^{2}-c_{n+1}^{2}=\frac{c_{n}^{2}-b_{n}^{2}}{3}
\]

这一看，这不是能直接把 \(b_{n}^{2}-c_{n}^{2}\) 看成一个整体直接计算吗！！这下我直接计算出

\[
b_{n}^{2}-c_{n}^{2}=7 \cdot\left(-\frac{1}{3}\right)^{n-1}
\]

把原本两式相加的 \(b_{n}^{2}+c_{n}^{2}=25\) 配合一下，直接就把 \(b_{n}^{2}\) 和 \(c_{n}^{2}\) 全算出来了，这相当于把静止化做到尽头了，又回归到刚才的方法了。那剩下的就没啥难度了。

这道题有些地方需要一定的灵感，但是做题的宗旨是不变的，一个是搞清楚控制关系，把控制关系给规范化之后，很容易就能够看出来 \(b_{n}^{2}+c_{n}^{2}\) 是定值，因为本身算 \(a_{n}^{2}=b_{n}^{2}+c_{n}^{2}\) 就是顺手一算。其实，原题目题干的两个已知相加相减也能够得到这个结论，也能够直接算出 \(b_{n}^{2}\) 和 \(c_{n}^{2}\) ，适合反应能力比较强的同学，大家可以去尝试一下。如果直接想到加减比较困难的话，老老实实通过思维方式的帮助，先去整理控制，在控制 \(a_{n}\) 的时候顺便直接看出来相加定值，然后充要替换的习惯会比较自然的把你导向信息量完整性的考虑，进而发现这些数列甚至可以写出通项来，这样的话稳定性会高一些。但是，这些发现仍然具有一定的偶然性，有见招拆招的色彩，不是一眼能够看到底的，所以难度仍然存在，难以保证考试就一定能处理完美。另外，这个题目本身也考察了 \(x_{n+1}=p x_{n}+q\) 型数列的处理方法，虽属于偏二级结论的东西，但是并不难，并且也不偏门，很多地方都能看到。直接算出通项那肯定是最舒服的，不然的话又得像例 4 一样去分析单调性。同学们可以尝试一下。

\section*{6．2．4 抽象函数问题}
说完数列，我们再简单说一下抽象函数。在第一部分里说过，抽象函数也是自由度无穷的对象，处理它的思维方式和数列非常接近，都是考虑通过类似＂前项控制后项＂的递推关系来达到控制的效果。

我们首先举一个很简单的例子：

\[
f(x)=f(-x)
\]

这是一个抽象函数，已知条件告诉我们它在 \(y\) 轴两侧是完全对称的，所以 \(x \geq 0\) 的部分和 \(x \leq 0\) 控制。 \(x \geq 0\) 部分的函数值可以任意取，取完一半就直接控制了另一半；或者也可以说 \(x \leq 0\) 的部分任取，再控制函数整体。这就是抽象函数的控制方法：用一部分的函数值去控制其他部分的函数值，和数列的＂首项靠递推控制全体＂的模式是不是有点像？而且，选取控制用的部分也不一定是死板的，像刚才的偶函数，选非负部分和非正部分都一样，是不是又和主自由度的想法有点类似？再比如

\[
f(x)=f(x+1)
\]

这个就是周期为 1 的函数，任意选一段长度为 1 的区间然后一段一段平移就可以控制整个 \(R\) 上的函数值，这个任选的一段区间就可以看成以前的＂主自由度＂，控制的起点。

我们还可以看一些更复杂的例子，比如

\[
f(x)=f(-x), f(1+x)+f(1-x)=0,
\]

这里有了两个条件，我们需要仔细一点分析。它是偶函数，且关于 \((1,0)\) 中心对称，两个条件交织在一起，似乎有点不好处理。我们先别急，去用类似摸项法的手段，去摸一摸实数轴，看看它长什么样。这下肯定不能随便取 \(x \geq 0\) 的函数值了，我们先试着取离对称轴，对称中心近的数，比如这里涉及到的轴和中心是 0 和 1 ，我们先随便取一段 \([0,1]\) 上的函数看看，既然是中心对称，那得保证满足 \(f(1)=0\) ，其他的好像暂时看不出来。我们把 \([0,1]\)

给先沿着这些轴翻转一下，看看确定了 \([0,1]\) 上的函数值之后，有多少部分能够得到控制。\\
首先沿着对称轴翻一下，［－1，1］就全都填满了；然后再沿着对称中心翻一下，［－1，3］全满了，再沿着对称轴翻一下，\([-3,3]\) 全知道了，再沿着对称中心翻一下，\([-3,5]\) 全知道了．．．．．．．我发现了，每翻转一下，这个区间都在变大，就是关于 0 和 1 不断进行着对称的疆域开拓。这早晚会把整个实数轴全盖满，所以 \([0,1]\) 上的函数值确实就足以决定整个 \(R\) 上的函数值。这是一个很神奇的事情，因为原本一个条件 \(f(x)=f(-x)\) 需半个实数轴才能控制整个实数轴，但是现在多一条件，就仅需一小段区间就能不断开拓疆域，四两拨千斤，控制整个实轴。当然还有另一种理解方式，就是我们知道

\[
f(x)=f(-x), f(1+x)+f(1-x)=0
\]

这俩条件可推出 \(f\) 是周期函数，

\[
f(x)=f(-x)=-f(2-x)=-f(x-2)=-f(4-x)=f(x-4),
\]

所以从周期的角度讲，长度为 4 的区间就能够定死整个函数，但是推导周期函数的时候丢失了更精细的函数信息，比如既然长度 4 足够，那［ 0,2 ］再偶函数翻折就是 4 ，所以长度 2 也够；［ 0,1 ］中心对称翻折就是［ 0,2 ］，那长度 1 也够。

做了这样一些游戏之后，相信大家也对抽象函数的控制方式有了一个直观的了解，我们来看一道问题。

例 \(6.9 f(x)\) 定义域为 \(R\) ，满足 \(f(x+1)+f(x-1)=f(2022), f(1-2 x)=f(2 x+5)\) ，若 \(f\left(\frac{1}{2}\right)=\frac{1}{2}\) ，则 \(f(2022)=\) \(\ldots ; \sum_{k=1}^{100} k f\left(k-\frac{1}{2}\right)=\) \(\qquad\)

这道题在抽象函数问题中，算中等偏难的了。需要注意的是，从题目来看，我们是要求一些具体的函数值，所以大家不需要过于纠结在这里需要多长的区间能定死整个函数，只要那些已知条件配合控制关系能够定死我们需要的部分即可。我们先来具体看一下问题。第一个条件是说 \(f\) 每相差 2 就有一个规律：和为定值，这个定值还恰好是 \(f(2022)\) ，后面还提问了这个数是多少，所以肯定是可知的，所以长度为 2 的区间就能够起到控制效力。还有一个条件，是一个轴对称条件，相当于

\[
f(5+x)=f(1-x)
\]

几何直观告诉我们对称轴是 \(x=3\)（或者直接推也能推出来）。又有一个轴对称，所以沿着对称轴把刚才长度为 2的区间折半，就可以知道长度为 1 的区间就够了。这个题，思考控制区间的长度比刚才的例子简单多了，所以直接考虑了一下。

现在我要求 \(f(2022)\) ，已知又有一个 \(f\left(\frac{1}{2}\right)\) ，是要用这个推过去求 \(f(2022)\) 吗？细想一下发现不对，这个函数是指定［2，3］上的函数值，然后对称成［2，4］上的函数，然后长度 2 利用第一个条件开拓到整个 \(R\) 上。设 \(k=f(2022)\) ，那么每相差 2 ，函数值就要关于 \(\frac{k}{2}\) 做一次对称，是这样控制的。那么 \(f\left(\frac{1}{2}\right)\) 能决定多少函数值？它是从 \(f\left(\frac{5}{2}\right)\) 出发，出来 \(f\left(\frac{7}{2}\right)\) ，然后翻到［0，2］上，变成 \(f\left(\frac{1}{2}\right)\) 和 \(f\left(\frac{3}{2}\right)\) ，也们当然是可以互相决定的。所以 \(f\left(\frac{1}{2}\right)\) 涉及到的位置都是些 \(\frac{2 n-1}{2}\) 这样的货色管不到，\(f(2022)\) 。这也同时说明 \(f\left(\frac{1}{2}\right)\) 这个已知条件和 \(k=f(2022)\) 的数值完全没有关系，所以我们做第一个填空就可以完全不管这个条件，肯定只需要那两个递推关系配合就能把第一个空搞出来。

那么怎么求 \(k\) 呢？肯定是前两个条件组合起来有猫娬。我们来思考一下其中的控制。

\[
f(x+1)+f(x-1)=k,
\]

每差 2 就要折一次，具体折多少暂且不知道，但是一个巧合是 \(f(2022)\) 恰好等于 \(k\) ，所以肯定要关注与 2022 彼

此控制的那些数，也就是偶数（相差 2 嘛）。所以就能看出来：假设 \(f(0)=t\) ，则

\[
f(2)=k-t
\]

则

\[
f(4)=k-(k-t)=t
\]

回去了，然后

\[
f(6)=k-t, f(8)=t \cdots
\]

可以看到，偶数的函数值就是 \(t, k-t, t, k-t \cdots\) 循环，所以到了 2022 ，应该是多少？ 4 的倍数都是 \(t\) ，所以

\[
f(2020)=t, f(2022)=k-t=k,
\]

所以 \(t=0\) ，循环就是 \(0, k, 0, k \cdots\) 第一个条件就告诉我们这么多，还有第二个条件呢！关于 \(x=3\) 对称！所以 \(f(2)=f(4)\) ，这不就是 \(k=0\) 吗！这不就完事了！由此又可见控制理念对做题稳定性的意义，几乎不需要算，就看看关系就直接得到了结果。

现在我们知道了，\(f\) 每加 2 就要变号。第二个空关注的函数位点都是 \(\frac{2 n-1}{2}\) 这些玩意，我们刚刚分析过 \(f\left(\frac{1}{2}\right)\)正好能把它们全部决定！所以这些值都是可知的！我们还是从 \(f\left(\frac{5}{2}\right)\) 出发，假设 \(f\left(\frac{5}{2}\right)=s\) ，那么 \(f\left(\frac{7}{2}\right)=s\) ，减 2再翻过去可得

\[
f\left(\frac{3}{2}\right)=f\left(\frac{1}{2}\right)=-s=\frac{1}{2}
\]

往后翻折也一样，都是两组两组一样的，然后交替变号所以从 \(f\left(\frac{1}{2}\right)\) 开始，分别是

\[
\frac{1}{2}, \frac{1}{2},-\frac{1}{2},-\frac{1}{2}, \frac{1}{2}, \frac{1}{2},-\frac{1}{2},-\frac{1}{2} \cdots
\]

现在第二个空所需的全部数值都静止了，代人写出来就是

\[
\frac{1}{2}(1+2)-\frac{1}{2}(3+4)+\frac{1}{2}(5+6) \cdots-\frac{1}{2}(99+100)
\]

求这个当然很简单，有 114514 种办法，分成等差数列，两两合并云云，比如我两两合并，

\[
\frac{1}{2}(1+2)-\frac{1}{2}(3+4)=-2
\]

能合并出 25 组来，最后就是 -50 了。

从这道题可以看出，抽象函数问题搞清楚如何控制之后，几乎就是顺水推舟，一路顺风就做下去了。大家以往在做这种题的时候，可能会写一大堆抽象推导，但可以看出，我没怎么写那些就过去了。这是我个人的做题风格，就是能直观就尽量直观，像这种偏静态的问题，直观去看控制，直接看具体数值有时候会更容易看清楚规律。当然，抽象推导在很多时候更有利于帮助我们直接得到规律，比如这里的

\[
f(x+1)+f(x-1)=k,
\]

再写一个

\[
f(x+1)+f(x+3)=k,
\]

一减就能得到周期函数

\[
f(x-1)=f(x+3)
\]

不同的做法有不同的优点，因为我一直跟大家讲控制的理念，所以我做什么题都会带有这样的影子，数列这种

问题大家可以自行选择适合自己的方式。

\section*{总结}
数列问题具有较为独特的控制模式和思维方式。数列这一对象本身可能具有无穷的自由度，但是我们在高中遇到的绝大部分数列都是静态的或者自由度很低的，这是因为数列往往被递推的关系给静止化了。因此，认识数列就必须从递推人手，按顺序从下往上去摸。这也决定了摸项的思维方式是重中之重。摸项就是一种特殊的遍历思维，把指标 \(n\) 作为滑块往上去摸，一边摸一边分析我们需要的东西。这种摸项的思维在我们上下两篇的绝大部分题目中都起到了很有价值的效果，帮助我们更好的感知数列，寻找下一步的出路。

数列问题的大题和小题差异很大。大题主要考察传统的等差，等比数列和衍生问题，比较简单的大题就是传统的等差等比的考查，自由度非常固定，与我们以前讲解的自由度，要素网的思维方式没有区别。更难一点的数列大题则可能涉及到更加陌生的数列形式，这种问题对我们的思维就提出了更高的要求，需要我们将经验分享 1－5 期的思维方式形成习惯，并将其和数列特有的静止化方法结合起来，同时做好隔离和元问题处理。这样的话，数列的大题基本就不成问题。

数列的小题经常在压轴难题位置出现，最为常见的就是静态数列分析。题目给出一个静态描述的数列，通过题干的叙述，我们可以将其完全静止化，呈现在我们眼前。这样的问题需要我们在讲过的基本的思维方式之上做好摸项的遍历，一边摸一边根据题目要求去塑造它的性质（也就是充要替换）。计算和思考是并行的，当单纯的思考提供不了更多信息时，就尝试用计算来做充要替换。这类静态的数列问题就突出考查一个静态分析能力和逻辑基本功，只要大家对我们讲过的思维方式很熟练，并熟悉了其特有的思维方式，就问题不大。

更难的数列小题常常以递推的形式出现，很难直接使整个数列一览无余，这时候把握好递推带来的控制关系就十分重要。递推关系包含了我们需要的全部信息量，而且广泛使用于做题的充要替换过程中。从例4和例 5 就可以看到，利用这些递推关系做好基本的充要替换，摸项这些工作，有助于我们获取到需要的信息。但是，这类问题的难度是没有上限的，像放缩，裂项之类的也是有可能出现，会对做题经验和临场反应有一定的要求。实在搞不定这些的话也没关系，少这几分不影响追求 130 和 140 。

控制的理念是十分深刻的，它背后的价值观就是事物会被如何决定。想清楚一个事物如何生成，会对理解这个事物，做出题目都有帮助。包括最后的抽象函数部分也体现了这个道理。希望本视频能够加深您对控制理念的理解，今后学习时至少能够拿下那些用思维可以完全取胜的数列问题，比如数列大题，静态分析小题等。

\section*{6.3 三角函数难题}
\section*{6．3．1 二级结论怎么学习？}
视频日期：2023年8月5日

题目：\\
在 \(\triangle A B C\) 中，\(\frac{\pi}{6} \leq C \leq B \leq A\) ，则 \(\cos \frac{A}{2} \sin \frac{B}{2} \cos \frac{C}{2}\) 的最小值为 \(\qquad\) ，最大值为 \(\qquad\)

这是一道较难的三角函数小题，相比于高考，更适合放在竞赛里。通过这道题，简单来谈谈 up 主对于二级结论的认识，并进一步强化一下大家对冻结法的理解。

题目考察的是三角形，三角形本体是三个自由度。但是不论是已知还是所求都只涉及角度（也就是形状），与大小无关。因此，题目涉及到的只有两个形状自由度。而已知的唯一条件 \(\frac{\pi}{6} \leq C \leq B \leq A\) 也是不等式的关系，只对自由度范围有约束，不产生实质的等量关系，因此不减少自由度。所以这个题目是一个双自由度的问题。既然是双自由度，自然就要使用我们的老朋友——冻结变量法了。先冻结一个自由度，把它融入题目背景，再研究在这个自由度背景下其他自由度的行为（这样的话，一次性研究的变量就少了，更加简单），求出这时候的最值。这个最值会和当前背景绑定，也就会直接被冻掉的自由度控制。这样的话就把其他自由度和冻结的自由度捆起来了。再解冻原本被冻结的自由度，这时候剩下的自由度就减少了，继续寻找最值即可。

这道题的主自由度选取没啥好说的，肯定就是三角形的角。看所求 \(\cos \frac{A}{2} \sin \frac{B}{2} \cos \frac{C}{2}\) ，我们需要先选取一个冻结的变量。选取谁呢？\(A, B, C\) 似乎都可以。从对称的角度来看，只有 \(B\) 前面是正弦，另外俩都是余弦，因此可以考虑选取 \(B\) 冻结一下试试（当然，你选取其他角冻结也是可以的）。冻结 \(B\) ，把 \(B\) 当作已知。冻结的时候要留意一下其取值范围，这样才能更好的理解 \(B\) 作为＂背景＂的信息，帮助我们后续分析。 \(B\) 都有可能取到什么值呢？显然下限 \(\frac{\pi}{6}\) 是能取到的。有没有可能更大？设问一下什么样的 \(B\) 可能满足条件，不难发现其实就是对于这个 \(B\) ，存在符合 \(\frac{\pi}{6} \leq C \leq B \leq A\) 的三角形罢了，也就是希望找另外两个角 \(A, C\) 夹住 \(B\) 。于是尝试寻找：

\[
\frac{\pi}{6} \leq \pi-A-B \leq B \leq A \Leftrightarrow \max \{B, \pi-2 B\} \leq A \leq \frac{5 \pi}{6}-B
\]

所以能找到 \(A\) 就等价于

\[
\max \{B, \pi-2 B\} \leq \frac{5 \pi}{6}-B \Leftrightarrow \frac{\pi}{6} \leq B \leq \frac{5 \pi}{12}
\]

所以冻结住的 \(B\) 就一定在 \(\left[\frac{\pi}{6}, \frac{5 \pi}{12}\right]\) 里，而且因为是充要替换，所以 \(B\) 可以遍历这个范围。冻结 \(B\) 后，暂时只需要关注 \(\cos \frac{A}{2} \cos \frac{C}{2}\) 。注意 \(\sin \frac{B}{2}>0\) ，因此要求最终的最小值，就要先求 \(\cos \frac{A}{2} \cos \frac{C}{2}\) 的最小值；最大值同理。\\
\(\cos \frac{A}{2} \cos \frac{C}{2}\) 怎么研究呢？常规高中方法就比较难处理。这就涉及到一个竞赛常用的二级结论：积化和差。有现成公式

\[
\cos \alpha \cdot \cos \beta=\frac{1}{2}[\cos (\alpha+\beta)+\cos (\alpha-\beta)] .
\]

这样的话上式直接变成

\[
\frac{1}{2}\left(\cos \frac{A+C}{2}+\cos \frac{A-C}{2}\right)=\frac{1}{2}\left(\sin \frac{B}{2}+\cos \frac{A-C}{2}\right)
\]

显然此时就只需要研究 \(A-C\) 这个直接被主自由度们控制的东西了。

说到这里，简单谈一谈 up 主关于二级结论的理解。二级结论是什么？我认为可以这样说：二级结论是超出考纲范围的，无法在大题中直接使用的，补充或拔高课内知识的知识点。既然如此，二级结论同样是知识点，因此自然的，它的学习也适用于经验分享第五期中的＂元问题＂的学习方法：搞清一个二级结论（知识点）对应的控制结构本质。例如：当你学习到上面积化和差公式，你该怎么从控制结构上去理解？那你就想，等式左边是什么？是三角函数乘积。它是怎么被控制的？是被两个角直接控制。等式右边是什么？是两角和与两角差的三角函数和。它是怎么被控制的？是先两角控制出一个和与一个差，再去控制三角函数值。因此，它的控制结构就是：两个角实现的控制，可以变成用两角和与两角差去控制，这在自由度上是平衡的，很好。那再从关系网落实到元问题上：具体是怎么控制的呢？左边是乘积，右边是和差，形式很简洁，所以才叫积化和差。所以我就懂了，上面说的一个改变控制方式的行为，可以把三角函数乘积与三角函数和差这两种计算给互相转化。所以它的元问题就是：直接控制与变成和差再控制（控制结构），可以实现乘积与和差之间的转化（元问题的形式），转化简洁优美，便于改换计算形式（作用）。这样我们就把它的元问题理解清楚了。

使用二级结论也是一样的。像刚才的题目，就是在 \(B\) 确定的情况下研究 \(\cos \frac{A}{2} \cos \frac{C}{2}\) 的取值情况。此时分析问题结构：只有一个自由度，体现在 \(A+C\) 为定值上面，因此围绕该牵制的信息（天然的求和）联想到使用转化成和差来实现控制的积化和差公式就不困难了。这也是一般的元问题调用（使用知识点）的思路：提取当前问题结构，匹配相应的元问题。因此，二级结论的学习和使用无非就是大一号的元问题学习罢了。再举几个例子：递推数列的特征根法无非就是一种特定形式的递推控制：极点极线无非就是一种特定的点控制线，线控制点，不过这套控制之下有一系列定点定直线等不变量问题；仿射无非就是除了大小无关的另一类更深刻的无关自由度．．．．．．．不过，二级结论自有其高于一般知识点的特殊之处，有些二级结论常常是特化针对某一类不太容易想出解法的控制结构。所以学习一些二级结论，有助于我们了解更多种类的控制结构，如果考试考察到了这类控制结构，就可以大大节省做题时间。

虽说如此，但是所有的二级结论都值得学习吗？也未必，有些二级结论越挖越多，过于繁琐或者过于难记，过于执着这样的二级结论反而容易走火人魔，耽误思维的培养和核心知识的学习。如何选择好的二级结论学习？这里给大家两点建议。第一，如果一个二级结论的证明很简单，本身的控制结构十分简洁清晰，具有美感，那么推荐学习。例如刚才使用到的积化和差公式 \(\cos \alpha \cdot \cos \beta=\frac{1}{2}[\cos (\alpha+\beta)+\cos (\alpha-\beta)]\) ，刚才分析了它的控制结构十分简洁干练，很容易看清。而且，它的证明也非常简单，配一组 \(\pm \frac{1}{2} \sin \alpha \cdot \sin \beta\) 合成右边即可，两行即可完成。对于这种证明不困难的二级结论，推荐了解如何快速将其步骤写成证明，这样的话就可以用于大题之中了。第二，如果一个二级结论揭示的控制结构的处理方法比较单一，而且不好想，那么也推荐学习。例如极点极线，要计算一个点对圆雉曲线的切点弦，有一套特定的推理方法，几乎不需要怎么计算，但是寻常很难想到；圆锥曲线统一极坐标，直接用角度控制到焦点的弦长，结论表达式短小精悍，充满美感。这样的二级结论掌握的话，做题时带来的简化具有不可替代性，同样推荐学习。反之，如果一个二级结论形式十分复杂繁琐，记忆成本太高还容易记错；或者做题时带来不了实质性的简化，或者简化效果很低（比如正常算也就慢个一两分钟），这样的二级结论就很无趣，没有数学美感，不推荐学习。Bilibili 上有很多关于二级结论的视频，其中不乏高质量视频也有不少很无趣的二级结论；有的学校老师也会进行二级结论的补充。具体什么样的二级结论是哪种类型，请大家自行甄别，自主思考能力是很重要的。

回到原题目。我们说到 \(\cos \frac{A}{2} \cos \frac{C}{2}=\frac{1}{2}\left(\sin \frac{B}{2}+\cos \frac{A-C}{2}\right)\) ，因此只需要关注 \(A-C\) 的取值范围，然后在这个范围里选让 \(\cos \frac{A-C}{2}\) 最大或最小的情形。下面就要看唯一的阳制 \(\frac{\pi}{6} \leq C \leq B \leq A\) ，很明显 \(0 \leq A-C<\frac{\pi}{2}\) ，所以希望原式小，就得 \(A-C\) 尽量大，反之亦然。而因为 \(A+C\) 是定值，所以一个小另一个就大，因此其实 \(A-C\)尽量大就是 \(A\) 尽量大 \(A-C\) 尽量小就是 \(A\) 尽量小。

那么 \(A\) 能大到（小到）什么程度呢？隔离出这个小问题思考，这不就是分析 \(A\) 的取值范围吗？在固定 \(B\) 的前提下，已知不就是关于 \(A\) 的不等式组吗？那就对原式直接充要替换呗！\(B\) 现在是背景，当静态已知，然后我们把 \(C\) 消去：\(\frac{\pi}{6} \leq \pi-A-B \leq B \leq A\) ，前面算过了，这等价于 \(\max \{B, \pi-2 B\} \leq A \leq \frac{5 \pi}{6}-B\) ，且 \(B \in\left[\frac{\pi}{6}, \frac{5 \pi}{12}\right]\) 是可以保证一定存在 \(\iota\) 的。于是放心大胆下结论：\(A\) 最大就是 \(\frac{5 \pi 6}{-} B\) ，最小就是 \(\max \{B, \pi-2 B\}\) 。\\
\(A\) 最大值比较简洁，我们先研究它，刚才说过，这对应原式取最小值的情形。现在 \(A\) 被 \(B\) 彻底控制，于是整个系统都被 \(B\) 控制，代回去把原式彻底控制出来即可。你可以代入 \(\frac{1}{2}\left(\sin \frac{B}{2}+\cos \frac{A-C}{2}\right)\) 的式子里慢慢展开整理，也可以直接代进 \(\cos \frac{A}{2} \sin \frac{B}{2} \cos \frac{C}{2}\) 里，预判下来都能稳定计算出来。以前者为例：

\[
\cos \frac{A-C}{2}=\cos \left(\frac{\pi}{3}-\frac{B}{2}\right)=\frac{1}{2} \cos \frac{B}{2}+\frac{\sqrt{3}}{2} \sin \frac{B}{2},
\]

于是原式为 \(\frac{1}{2} \sin \frac{B}{2}\left(\frac{1}{2} \cos \frac{B}{2}+\left(1+\frac{\sqrt{3}}{2}\right) \sin \frac{B}{2}\right)=\frac{1}{8} \sin B+\left(\frac{1}{2}+\frac{\sqrt{3}}{4}\right) \frac{1-\cos B}{2}\) 。这个函数显然在 \(B \in\left[\frac{\pi}{6}, \frac{5 \pi}{12}\right]\) 上单调递增（正弦增余弦减），于是最终的最小值就是在 \(B=\frac{\pi}{6}\) 时取，代进去就算出原式最小值是 \(\frac{1}{8}\) 。

下面考虑 \(A\) 的最小值，这对应原式的最大值。这需要进一步讨论：如果 \(B \in\left[\frac{\pi}{3}, \frac{5 \pi}{12}\right]\) ，则此时固定 \(B\) 之下的最大值在 \(A=B\) 时取，对应最大值 \(\cos \frac{B}{2} \sin \frac{B}{2} \cos \frac{\pi-2 B}{2}=\frac{1}{2} \sin ^{2} B\) ；如果 \(B \in\left[\frac{\pi}{6}, \frac{\pi}{3}\right]\) ，则此时固定 \(B\) 之下的最大值在 \(A=\pi-2 B\) 时取，此时有 \(C=B\) ，于是显然最大值仍然是 \(\frac{1}{2} \sin ^{2} B\) 。有趣！这说明 \(B \in\left[\frac{\pi}{6}, \frac{5 \pi}{12}\right]\) 时，固定 \(B\)之下的最大值都是一样的表达式：\(\frac{1}{2} \sin ^{2} B\) 。这时候单调性显然是增的，于是最后的最大值就在 \(B=\frac{5 \pi}{12}\) 时取这时 \(\sin \frac{5 \pi}{12}=\frac{\sqrt{6}+\sqrt{2}}{4}\) ，于是最终得数易算出是 \(\frac{2+\sqrt{3}}{8}\)

\section*{总结：}
本期视频讲到的题目就是普通的冻结法的题目，大家再好好回顾一下冻结法解题的流程。本期视频想给大家强调两个重要事情。第一个是注意多关注各个自由度的取值范围，比如一开始我们就计算了 \(B\) 的范围。不管你在何处，关注自由度的取值范围都是非常重要的事情，因为只有明确了取值范围，才能使用遍历法，才能清楚地看到系统的结构。不然的话，就只有范围混乱，遍历失效，做不清楚题的份。第二个是关于二级结论的学习与使用。它们与普通的元问题知识点没有太大区别，使用起来也是一样的。但是学习二级结论的态度要端正，取其精华，去其糟粕。这我们也在前面强调了，大家可以回去翻文档。只有正确科学地，有选择地学习二级结论，才能让二级结论更好地服务于我们的解题，帮助我们冲刺高分。

\section*{课后练习}
1．积化和差与和差化积有许多公式，大家可以在网上轻松地找到它们。尝试用视频类似的方法，用不超过三行的过程完成每个公式的证明。然后盖住这些公式，只剩下等号左边部分，尝试不死记硬背，利用刚才的证明思路去直接心算拼凑出这些公式的等号右边部分。

2．我们在做这道题的时候，选择了冻结 \(B\) 。请你尝试冻结 \(A\) 与冻结 \(C\) ，用视频内的思路，完成这道题目的求解。这样的话，算上视频内的方法，一共有三种冻结方式，于是就相当于有三种方法求最大值，也有三种方法求最小值。你认为这两个最值分别用哪一种方法最简单？

\section*{6．3．2＂灵感＂也扎根于思维}
视频日期：2023年8月29日

\section*{题目：}
已知在 \(\triangle A B C\) 中，角 \(A, B, C\) 所对的边分别是 \(a, b, c, B C\) 边上的高长度为 \(h\) ，且满足 \(a+h=b+c\) ：\\
（1）若 \(a=\frac{3}{2} h\) ，求 \(\cos A\) 的值\\
（2）当 \(A\) 最小时，求 \(\tan A\) 的值

这道三角函数题目难度偏大，此前已有不少同学提问过，但考虑到这道题风格与高考题有一点差别，所以一直没讲。现在本年高考结束，且不久前又有同学提问，所以拿出来讲一下。整道题从头到尾都是大小无关的：把三角形保持形状的放大缩小，所有的已知仍然保持，所有的所求或所证结论不变。所以大小这个自由度可以无视，换而言之，我们可以自主指定一个大小，从而题目条件白多一个，自由度白减一个。原始题干中仅有一个已知条件：\(a+h=b+c\) ，但是这个条件比较怪异，我们暂时不知道该如何利用，也就没什么可以先做的事情，只能暂时到这。所以直接看具体的问号。

第一问又给了一个条件 \(a=\frac{3}{2} h\) ，配合原题干里已有的那个条件，自由度减 2 。三角形本身只有三个自由度，再把大小自由度减掉，直接静止了。而所求的 \(\cos A\) 也是直接被三角形的要素 \(A\) 控制，所以思路就很明确了：把三角形的静止化完成，然后去输出这个角 \(A\) 的信息即可。

按照惯例，我们先来给三角形规定一个大小。由于已知 \(a=\frac{3}{2} h\) ，因此我们自然可以不妨假设 \(h=2, a=3\)（当然你选取其他大小也可以）。然后题目还有一个已知：\(b+c=5\) 。把这些要素合起来看：有两个都是三角形基础要素，另一个 \(h\) 算个次级要素。这种次级要素没法直接研究，肯定得先用初级要素控制出来，于是自然提炼元问题：怎么控制 \(B C\) 边上的高？这个元问题的处理就比较单一：先算面积，然后已知底 \(a=3\) ，就有了 \(h\) 。而面积显然可以轻易被基础要素控制：考虑到 \(b+c=5\) 中的 \(b, c\) 对称，且题目最后要求 \(A\) ，所以我们用 \(A, b, c\) 来表示面积：

\[
\frac{1}{2} b c \sin A=\frac{1}{2} a h \Leftrightarrow b c \sin A=6 .
\]

这样的话三个方程，理论上的静止化就完成了，接下来只需要求解 \(A\) 。方程组需要三个主元，自然就选 \(A, b, c\) ，然后把 \(a=3\) 也用 \(A, b, c\) 表示出来（列牵制）：

\[
9=b^{2}+c^{2}-2 b c \cos A
\]

接下来就是纯粹解方程了。为了求 \(A\) ，需要消去 \(b, c\) ，而这些式子都关于 \(b, c\) 对称，所以消元方式就很简单了：

\[
b^{2}+c^{2}+2 b c=25
\]

然后有

\[
9=25-2 b c-2 b c \cos A=25-2 \cdot \frac{6}{\sin A}(1+\cos A)
\]

式子只剩下 \(A\) 了，目标完成。再计算得 \(\frac{1+\cos A}{\sin A}=\frac{4}{3}\) ，进而可以解出 \(\cos A=\frac{7}{25}, \sin A=\frac{24}{25}\) ．

下面来看第二问。第二问无非就是少了第一问的条件，研究对象与第一问一样，都是 \(A\) 。由于少了一个等式，所以第二问在大小无关后仍然有一个自由度，所以这时候，选取什么量作为主自由度呢？这是本题最大的难点。选主自由度需要保证可计算性，不能选了之后发现控制式太复杂，无法研究，甚至根本连控制式的显式表达式都列不出来。那么，这道题我们应该如何选取主自由度，才能满足前面的要求呢？比如，我选择 \(A, b, c\) 作为主自由度，你会发现，结合第一问的过程，\(h=\frac{b c \sin A}{a}\) ，代进之后 \(a+h=b+c\) 之后，式子极为混乱，没法计算。那么，该如何选取主自由度呢？

其实，这道题的第一问已经给了我们提示。第一问用 \(a\) 和 \(h\) 作为已知，顺利地把 \(A\) 用方便计算的式子拿捏得死死的，这说明选取 \(a\) 和 \(h\) 作为主自由度（去掉大小正好是 1 个自由度）可以很好地拿捏 \(A\) 。那么，这么选取主自由度，它的控制图景（要素关系网）是怎样的呢？我们不妨设 \(a=1\) ，然后固定下一个 \(h\) ，这时候相当于要找使得 \(b+c=1+h\) 的三角形。既然我们把 \(h\) 选为了主自由度，那就让控制围绕这个 \(h\) 来。高是 \(h\) ，就是把 \(B, C\) 静止地摆在地上，然后拿一根长度 \(h\) 的小棍横向移动。那么应该滑动到哪里呢？自然就是当恰好满足 \(b+c=1+h\) 时，就卡出一个符合题意的三角形。在此基础上，变化 \(h\) ，每个 \(h\) 之下都按照上述方式搞出三角形，我们就得到了 \(h\)

对题目的控制方式。那么自然的，我们现在需要确定 \(h\) 的取值范围（定义域）以及 \(h\) 对 \(A\) 的控制。后者的表达式我们根据第一问的经验可以相信比较好写出来，而前者在列表达式的时候把相应的式子成立的约束写出来即可。于是目标明确，只需要开算即可。

有的同学可能会问：从第一问的线索去思考，这是怎么想到的呢？其实大家回忆一下我们之前讲的一种思维方式：思想试验（出处为 2023年3月5日的视频），大家就会发现，其实本质是一个东西。思想试验就是先预设某个自由度静止，在分析这时的状况（自由度少，更简单）时发现问题本质，进而将本质提取出来，用于分析刚才自由度运动时（自由度多，更复杂）的控制。在这样的思维方式之下，加上两个问号除了差一个条件外，其余问题核心都一样的特性，想要意识到第一问就是第二问的一个思想试验，就不觉得困难了。这也给大家一个启示：思维方式只有转变成习惯，才能在无形之中给予你力量，只是机械地把思维方式当模板背下来去套用，是无法发挥出其十成效果的。只有在学习思维方式的时候，多去想＂这是如何思考的＂，多去让大脑顺从这种内在的思考行为，才能把思维方式变成自己的习惯，无需刻意回忆便能潜移默化地自如应用。

回到题目。按照我们之前讲过的习惯，选好主自由度之后要先考察其取值范围（定义域）。这样才符合我们遍历法的思维方式。假定 \(a=1\) 时，对于什么样的 \(h\) ，才能控制出符合题意的三角形呢？我们把思维焦点暂时聚焦到这个问题上。换个说法就是：什么样的 \(h\) ，存在符合原题的三角形。有了结构层面的控制，才有具体的计算。这时候的控制图景是如何的呢？刚才说过，先把 \(B C\) 静止放在那，然后取一条与 \(B C\) 平行，垂直距离为 \(h\) 的直线 \(l\) 。考虑在 \(l\) 上滑动的点，让它恰好满足已知的 \(b+c=1+h\) 即可。也就是说，只要滑动过程中存在满足 \(b+c=1+h\)的点即可。这个条件提炼成控制，就是到两个定点的距离和。因此这个元问题的处理方式就很明确了。两条道路：一个是关注 \(b+c=1+h\) 的角度，距离和是定值，等价于说，这个点要在以 \(B C\) 为两个焦点，长轴长 \(1+h\) 的椭圆上；另一个则是关注点所在的滑动轨迹一一直线，定直线 \(l\) 上一动点到两定点 \(B, C\) 的距离和，这正是初中难度的基本问题一一将军饮马。这两条路都可以让我们搞出 \(h\) 的取值范围。以后者的将军饮马为例，经过对称，我们可以发现 \(l\) 上动点到 \(B, C\) 距离和有一个最小值 \(\sqrt{1+4 h^{2}}\) ，但取值上不封顶。因此＂能取到 \(1+\mathrm{h}\)＂，就等价于

\[
\sqrt{1+4 h^{2}} \leq 1+h
\]

解出范围是 \(0<h \leq \frac{2}{3}\) 。求范围红束，我们把思维焦点切换回原题。\\
完全模仿第一问的计算，由刚才的 \(a=1\) ，于是有

\[
\frac{1}{2} b c \sin A=\frac{1}{2} a h \Leftrightarrow b c \sin A=h .
\]

然后把 \(a=1\) 变成 \(1=b^{2}+c^{2}-2 b c \cos A\) ，再加上

\[
b+c=1+h \Leftrightarrow b^{2}+c^{2}+2 b c=(h+1)^{2}
\]

我们可得

\[
1=(h+1)^{2}-2 b c(1+\cos A)=(h+1)^{2}-\frac{2 h}{\sin A}(1+\cos A)
\]

于是

\[
\frac{1+\cos A}{\sin A}=1+\frac{h}{2}
\]

这里提醒大家一下，我们说 \(m\) 能控制 \(n\) 是指理论上有了 \(m\) 就能定下 \(n\) ，我们列出的式子只要能起到这个效果即可，不一定非要强行写成 \(n=f(m)\) 的形式。比如上面的，显然给一个 \(h\) ，我们就能写出一个关于 \(A\) 的方程，这就可以了。

控制式和取值范围都有了，于是我们得到

\[
1<\frac{1+\cos A}{\sin A} \leq \frac{4}{3}
\]

我们要找 \(A\) 最小的时候，而当 \(A\) 增大时 \(\cos A\) 减小， \(\sin A\) 变大，于是 \(\frac{1+\cos A}{\sin A}\) 变小，因此 \(A\) 最小就是 \(\frac{1+\cos A}{\sin A}\)虽大，也就是 \(\frac{4}{3}\) 。这正是第一问的式子。第一问算过

\[
\cos A=\frac{7}{25}, \sin A=\frac{24}{25}
\]

因此自然的此时 \(\tan A=\frac{24}{7}\) ，这就是最终答案

\section*{总结与提升：}
本期视频选取的例题难度较大，主要难度体现在思维方式上。首先要通过大小无关与规定大小来减少自由度，更好地看清本质，这我们强调得够多了，不再赘述。主要是要素关系网的建构与主自由度的选取。第一问是一个热身，仅通过题设与所求就能轻松选取可计算性强的主自由度，然后就是基本的代数处理。而第二问与第一问在结构上的差异就是把明显的主自由度选取提示给去掉了。要看到这一层，对于自由度与要素结构的把握是不可缺少的。随后，思想试验的思维方式让我们联想到照搬第一问的设参数方式，只不过这里留下一个主自由度 \(h\) 。选完主自由度之后，要构建一个好的要素关系网。本题讲解时选用了画直线加卡 \(b+c\) 的方式，实际上如果选择先画椭圆，再拿高度 \(h\) 的小棍去卡，也是可以的。这样理解的话，求 \(h\) 的范围时就会自然地考虑写椭圆方程计算，计算量稍微大一点，但仍然是一条可行的道路。

整道题无处不在强调对结构把握的重要性。本题最大的难点就是想到使用第一问的方法去处理第二问，而要想到这一层，则需要思想试验的思维方式成为自己的一种潜意识，这对思维的要求很高。本题的难度个人认为应当已经超过了高考，作为高中数学联赛的填空题第 7 题难度应该刚好。题目虽难，但是只要仍然在考试题的难度范围之内，其思考都应是有迹可循的。重视思维方式的学习，多通过科学的思维方式反思如何思考才能走上正轨，做出题目，并反作用于自己的思维反思，进而提升自己，才能不断进步。

\section*{课后练习：}
本题在确定 \(h\) 的取值范围时，还给出了以 \(b+c=1+h\) 为控制，以＂高度 \(h\)＂为牵制的一种要素结构。请使用这种要素结构来计算 \(h\) 的取值范围，并利用＂焦点三角形面积公式＂这个二级结论（不知道的话可以百度）来完成整个题目的求解。再结合 \(u p\) 主 2023 年 8 月 5 日发布的视频中二级结论的学习及应用理念，你认为在求解该题的时候如何思考，才能想到这种基于椭圆的解法。

\section*{6．3．3 大小自由度更深的理解}
题目：\\
记 \(\triangle A B C\) 的内角 \(A, B, C\) 所对的边分别为 \(a, b, c\) ，已知 \(\frac{a-\cos B}{a-\cos C}=\frac{\sin B}{\sin C}\)\\
（1）若 \(b \neq c\) ，证明：\(a^{2}=b+c\)\\
（2）若 \(B=2 C\) ，证明： \(2 c>b>\frac{2}{3}\)

先看第一问。三角形三个自由度，题干的已知给出来之后还剩俩，这个已知相当于一个牵制等式。第一问让我们证明的也是一个等式，从自由度的角度考虑的话，明显这是让我们去做充要替换，所以思维焦点放到对式子的代数变形上面去。

自然先交叉相乘去掉分式：

\[
(a-\cos B) \sin C=(a-\cos C) \sin B
\]

接下来的整理方式很明显：由于所证全是边，自然把角全搞成边（注意这里的必胜逻辑，全变成边是充要替换，逻辑上信息量不丧失，而三边的等式牵制会成为三个独立自由度的多项式，而逻辑上多项式互推只有等价变形和因式分解，基本功罢了），先变成

\[
(a-\cos B) c=(a-\cos C) b
\]

于是

\[
a(c-b)=c \cos B-b \cos C=c \frac{a^{2}+c^{2}-b^{2}}{2 a c}-b \frac{a^{2}+b^{2}-c^{2}}{2 a b}=\frac{c^{2}-b^{2}}{a}
\]

变成了纯这的多项式，而所证也是纯边的多项式，所以上面这个式子必定能直接推过去（做题方向选取的信心来源正是逻辑）。接下来就很简单了，显然有公因式 \(c-b\) ，约去得到 \(a=\frac{c+b}{a}\) ，这正是我们所证的。因此，充要替换的理念可以迅速帮助我们找到第一问的解题思路。

关键是第二问。 \(B=2 C\) 说明一定是 \(b \neq c\) ，所以原式的信息量全集中在 \(a^{2}=b+c\) 上（注意第一问去分式的时候得到的等式得要求 \(a \neq \cos C\) ，但如果不然，则 \(a-\cos C=0\) ，进而有 \(a-\cos B=0\) ，这又会得到 \(B=C\) 。所以一个 \(B=2 C\) 的已知就说明 \(a^{2}=b+c\) 和原题的式子完全等价，信息量完全不丢失）。而且 \(a^{2}=b+c\) 多么简洁漂亮，从第五期的简化原则来看，一定是用这个来做题。所以我们彻底抛弃原题条件，就从 \(a^{2}=b+c\) 和 \(B=2 C\) 出发去推导所证式子。

下面才是重头戏。所证式子有两个，都是不等式，其中左边的那个 \(2 c>b\) 是大小无关的，只和三角形形状有关，而右边那个 \(b>\frac{2}{3}\) 稍微有点摸不着头脑。肯定要先试着分析已知，已知也很有趣。B＝2C 也是只衡量形状的，它直接把大小以外的两个自由度消去 1 ，说明在这个条件的基础之上，一个参量就能定死形状。很显然，把这个主自由度给 \(C\) 最合适，\(C\) 控制 \(B\) ，进而控制整个三角形的三个内角，形状就定下来了。而第二个已知却不是大小无关的了，\(a^{2}=b+c\) 不齐次。但是我们依然要用控制的基本理念去思考（注意控制，要素关系网要内化成思维范式，内化成习惯才能用出最大效果）。题目本身就是两个式子，实际自由度为 1 ，而先前 \(C\) 把形状控制了，所以这个条件刚好可以把三角形静止下来。那这个 \(a^{2}=b+c\) 无疑就是把最后一个自由度，也就是大小给定下来的！换而言之，\(C\) 决定了 \(a, b, c\) 三者的比例，因此实际的三边就是在这个比例之上乘一个数（大小自由度的实际理解就是这样，当初讲大小无关时渗透的思维也是这样的）。把 \(C\) 冻住（先定下形状）的话，就可以设 \(a=a_{0} k, b=b_{0} k, c=c_{0} k\) ，其中的 \(a_{0}, b_{0}, c_{0}\) 就是被 \(C\) 控制的比例（跟解法向量那样选一组喜欢的即可），\(k\) 就是决定大小的系数。那么，利用非齐次的已知就可以解出 \(k=\frac{b_{0}+c_{0}}{a_{0}^{2}}\) 来，从而实现静止化。

梳理一下刚才控制的思维过程，你可以发现，其实只要已知中有大小无关的条件，那么就可以把主自由度先分配给形状，把形状控制消楚之后，剩下的已知自然就能去求解真正的大小，可以说是把大小无关的思维方式发挥到了极致。

有了这样的思维指南，我们做题自然也会是砍瓜切菜。第一个不等式 \(2 c>b\) 本身就和大小无关，所以最后一步靠 \(k\) 定下大小的已知 \(a^{2}=b+c\) 在这里没用！所以从逻辑上就可以预判到 \(B=2 C\) 的条件就能证出这个不等式（注意这个不等式隐含了定义域 \(C \in\left(0, \frac{\pi}{3}\right.\) ）因为 \(B+C<\pi\) ）。思路自然也很简单了，大小无关时唯一的主自由

度就是 \(C\) ，那就用 \(C\) 去计算

\[
\frac{b}{c}=\frac{\sin B}{\sin C}=\frac{\sin 2 C}{\sin C}=2 \cos C
\]

由 \(\cos C<1\) 即得证。\\
第二个 \(b>\frac{2}{3}\) 的不等式其实就是一个衡量大小的。刚才就分析过了，\(C\) 定死形状，加一个已知就能定死大小。那自然就要用主自由度 \(C\) 去计算大小。因为是求 \(b\) ，那干脆就把边的比例都用 \(b\) 作为标定。刚才算了

\[
\frac{b}{c}=2 \cos C \Rightarrow c=\frac{b}{2 \cos C},
\]

还差个 \(a\) ，随便拿正张弦定理算算（ \(C\) 做主自由度）：

\[
\frac{a}{b}=\frac{\sin A}{\sin B}=\frac{\sin (B+C)}{\sin B}=\frac{\sin 3 C}{\sin 2 C} \Rightarrow a=\frac{b \sin 3 C}{\sin 2 C}
\]

这里得化简一下（本题最大的难点）：

\[
\frac{\sin 3 C}{\sin 2 C}=\frac{\sin 2 C \cos C+\cos 2 C \sin C}{2 \sin C \cos C}=\frac{2 \cos ^{2} C+\cos 2 C}{2 \cos C}
\]

（主要是约去一层 \(\sin C\) ），最后

\[
a=b \frac{4 \cos ^{2} C-1}{2 \cos C}
\]

现在代人控制大小的已知，让 \(C\) 定出 \(b\) ：

\[
b^{2}\left(\frac{4 \cos ^{2} C-1}{2 \cos C}\right)^{2}=b\left(1+\frac{1}{2 \cos C}\right)
\]

把能终的都约去，得

\[
b=\frac{2}{\cos } C(2 \cos C+1)(2 \cos C-1)^{2}
\]

注意 \(C\) 遍历 \(\left(0, \frac{\pi}{3}\right)\) ，令 \(t=2 \cos C \in(1,2)\) 最后求证的式子就是

\[
\frac{t}{(t+1)(t-1)^{2}}>\frac{2}{3}
\]

一路充要替换下来，到这已经变成很基本的不等式问题了，接下来怎么做都行了。所证式等价于 \(3 t>2(t+\) 1）\((t-1)^{2} \Leftrightarrow 2 t^{3}-2 t^{2}-5 t+2<0\) ，最无脑的方法是三次函数求导当然也可以因式分解成 \((t-2)\left(2 t^{2}+2 t-1\right)<\) \(0 \Leftrightarrow 2 t^{2}+2 t-1>0\) ，更加简单。

做一个总结。这道题在解三角形大题里难度算是数一数二的。最想向大家传递的是大小自由度的问题。我们之前强调过很多次大小无关，实际上，任何几何问题都可以把大小和形状分开来看，大小占 1 个自由度，其余的自由度都是形状的自由度。有一些题目的已知可能不全都是大小无关的，但是有一部分是大小无关的（判断方法和以前讲的一样），这一部分带来的牵制只作用于形状，不作用于大小，如果分开看的话，会对整个问题的要素关系网，对这个系统的运作规律有一个更深的理解，会更好的帮助我们去建构整个问题，选取方法计算。前期花了很长时间带大家去分析这个题目的控制关系也是希望大家对我们将过的思维方式有着更进一步的理解。

\section*{思考题}
1．在证明 \(2 c>b\) 时，我们说过已知 \(a^{2}=b+c\) 在这里没用。但是如果我换一个已知，比如说 \(a^{2}=2 b-3 c\) ，你会发现里面隐含着 \(2 b>3 c\) 的要求。请证明：若 \(2 b>3 c\) ，则 \(C\) 将无法遍历区间 \(\left(0, \frac{\pi}{3}\right)\) ，这说明这种定下大小的已知也可能影响形状。请你思考，原题中的已知条件 \(a^{2}=b+c\) 是否会额外附加对形状的影响（不要去做细致的计算，不然思维判断就本末倒置了）？换而言之，请你不通过细致计算来判断，是否对于 \(\forall C \in\left(0, \frac{\pi}{3}\right)\) 所定下来的形状，

由 \(a^{2}=b+c\) 都能够解出唯一的 \(\triangle A B C\) ？（提示：看看当时关于 \(k\) 的分析）\\
2．我们知道 \(b+c>a\) ，请你利用这个不等式，充分利用第二问的已知条件，想一个简单办法，不经过复杂计算证明 \(c>\frac{1}{3}\) ，你认为这种巧妙方法是否能用来证明 \(b>\frac{2}{3}\) ？

\section*{6．3．4 从一道小题中挖出大小无关}
视频日期：2023年9月29日

题目：\\
已知在 \(\triangle A B C\) 中，角 \(A, B, C\) 所对的边分别是 \(a, b, c\) ，且 \(\frac{c}{a}=\frac{\sqrt{3}}{3} \sin B+\cos B\) 。若 \(D\) 为 \(B C\) 的中点，且 \(A D=2 \sqrt{3}\) ，则 \(a\) 的最小值为 \(\qquad\)。

大小无关问题我们一直在反复讲。实际上，大小无关的理念就是把几何问题的大小和形状分开看。因为控制形状所需要的自由度少于原题，而大小则是纯粹的比例，很好处理，所以问题的难度会降低，更易于我们入手解题。理念是指引我们前行的灯塔，只有好的理念才能让我们朝着胜利越来越近。正是因为大小与形状分开看的理念很有价值，所以我们才会反复去讲解。除了最基本的大小无关题目（所有的已知与所求全部大小无关），我们还在2023年1月30日的视频内讲解了一道进化版的＂大小无关＂问题，在那里把纯粹的大小无关给进化为大小形状分开看的思考方式。而这道题则会通过大小与形状的有趣处理，让大家更深刻地认识到这一理念在分解问题要素关系，筛选解题策略上的价值。先来读题。整个系统的背景简明，就是三角形再加个中点。但是给出的已知信息很有趣，第一个已知 \(\frac{c}{a}=\frac{\sqrt{3}}{3} \sin B+\cos B\) 是纯粹的大小无关，只与形状有关。而第二个条件是给了一条线段的长度，是个很纯的大小条件。最后所求则是大小条件。已知与所求都是大小形状泾渭分明，这很明显的提醒我们，把大小和形状分开来看待是一个不错的选择。如果我们把大小与形状拆成要素网的两个分支，那么形状是两个自由度，被已知减掉一个，所以本题的形状是由单自由度控制的。在形状确定的前提下，给出一个大小条件就能把大小也定死，从而确定系统。这就是这个题目在形状与大小的视角下的一个要素结构。

既然这个题可以把大小和形状分开看，这就说明可以把它们分开处理。因此，我们就有了一个大致的解题方向：首先专心处理形状部分，找到用纯粹的形状自由度定死形状的式子：然后在拿捏住形状之后，把大小求解出来（因为是求解，所以解完之后的式子里也应该只剩形状参数），用这个求解出来的大小去刻画最终所求。我们先来看形状部分。已知条件中涉及到 \(a, B, c\) 三个要素，自然的SAS 型控制。其中 \(B\) 是纯形状参数；\(a, c\) 都是大小参数，但是 \(\frac{c}{a}\) 就是大小无关的形状参数，记 \(\frac{c}{a}\) 为 \(k\) 。于是已知就是 \(k=\frac{\sqrt{3}}{3} \sin B+\cos B\) 。很明显，这是一个形状参数 \(B\) 控制另一个形状参数 \(k\) 的式子！！因此，自然的就选 \(B\) 作为主自由度，直接控制消楚了两个形状参数，整个形状被 \(B\) 牢牢㨘在手里，对应的元问题就是 SAS 型控制结构。选好主自由度后，我们多次强调，要注意主自由度的取值范围。这里稍微遍历一下就知道显然需要 \(k>0\) ，也就是 \(\frac{\sqrt{3}}{3} \sin B+\cos B>0\) ，即 \(B \in\left(0, \frac{2}{3} \pi\right)\) 。

形状的控制清晰了，那么就该从要素网向外输出信息了。该去算什么呢？自然就要从大小的已知里找。已知 \(A D=2 \sqrt{3}\) ，那么要控制的对象自然就是 \(A D\) ，也就是中线。刚才的元问题是 SAS 型控制，那么 SAS 怎么控制中线？这是基础知识了，自然是用 SAS 刻画的向量基底直接控制。用 \(c, a\) 对应的向量 \(\overrightarrow{B A}, \overrightarrow{B C}\) 去表示 \(A D\) ，不难写出 \(\overrightarrow{A D}=\frac{1}{2} \overrightarrow{B C}-\overrightarrow{B A}\) 。然后平方就得到 \(|\overrightarrow{A D}|^{2}=\frac{1}{4} a^{2}+c^{2}-a c \cos B\) 。接下来，我们贯彻理念，把大小与形状分解开（大家注意这是关键步骤）。选取 \(a\) 作为大小的指标（因为这是所求），这样用形状主自由度 \(B\) 加上大小主自由度 \(a\) 就可

以控制全局了。下面开始计算：先把 \(c\) 控制，自然就是 \(c=a k=a\left(\frac{\sqrt{3}}{3} \sin B+\cos B\right)\) ，带进去就实现了 \(a, B \rightarrow A D\)的控制。加上 \(A D\) 长度的牵制导到 \(12=a^{2}\left[\frac{1}{4}+\left(\frac{\sqrt{3}}{3} \sin B+\cos B\right)^{2}-\cos B\left(\frac{\sqrt{3}}{3} \sin B+\cos B\right)\right]\) 。可以看到，这个式子实现了本题用形状参数 \(B\) 定下大小的效果：输入一个主自由度 \(B\) 的值，就吐出一个大小 \(a\) 。最后要研究的就是 \(a\) ，于是直接把 \(a\) 分到左边，右边展开化简，表示成辅助角公式形式，得到 \(\frac{12}{a^{2}}=\frac{\sqrt{3}}{6} \sin 2 B-\frac{1}{6} \cos 2 B+\frac{5}{12}\) 。用辅助角公式得到右边的最大值为 \(\frac{3}{4}\) ，等号成立时 \(B=\frac{\pi}{3}\) ，可以取到。于是进一步得到 \(a_{\mathrm{min}}=4\) 。

题目虽然用这种方法做完了，但是我想跟大家分享的东西还没完。这道题刚才我们在选取主自由度时用到的核心思想就是分割大小形状，体现了大小形状分开看的思想。事实上，这种思想还可以有别的形态的变化。对于本题及其他类似的问题，我给大家分享一种玩弄大小与形状的特殊手段，通过这种手段，大家可以把要素关系网等价变形成纯控制，无牵制的形态。这样的操作虽然不直接减少解题运算量，但要素网结构的精简有助于我们更好地进行思考。

我们继续回头看这个题，刚才我们说到 \(\sin B+\cos B\) 是纯大小无关的条件，而另一个已知和所求都是一次的大小条件。我们以前做的不少题都是从题干到所求全程大小无关，而这道题并不是这样。有没有办法，把它化归到纯粹的大小无关问题呢？

我们来做这样一个操作：所求的 \(a\) 给它充要替换成 \(\frac{a}{A D} \cdot 2 \sqrt{3}\) 。因为有 \(A D=2 \sqrt{3}\) 的已知，这个长度对于本题是既定事实，所以这个转化确实是充要替换。这么一写，我们就发现，所求变成大小无关的形式了！这时，我们思考：做了这个转化之后，\(A D=2 \sqrt{3}\) 这个条件还有用吗？我把这个条件拿掉，可以吗？这个题的形状本身是一个主自由度在控制而 \(A D=2 \sqrt{3}\) 只是规定了一个大小罢了，不对形状有任何干涉。刚才我们已经把所求改写成大小无关形式了，所以原题刻意规定一个大小已经没有意义了，不管你规定 \(A D=2 \sqrt{3}\) 还是 \(A D=4 \sqrt{3}\) 还是 \(A D=114514\) 都不会影响最后的答案了，所以直接把这个条件去掉也不会有任何问题了！因此题目就被我们改造成了：

例 6.10 已知 \(\frac{c}{a}=\frac{\sqrt{3}}{3} \sin B+\cos B\) 。若 \(D\) 为 \(B C\) 的中点，则 \(\frac{a}{A D} \cdot 2 \sqrt{3}\) 最小值为 \(\qquad\) －

从刚才的分析中我们可以看到，这个新题的答案应当是与旧题的答案一样的。那么这个改造有什么好处呢？我们现在抛弃旧题，只看这个新题。不难发现新题是一个纯粹大小无关的，两个自由度的题，舍去大小自由度后，实际自由度为 1 ，与旧题一致。但现在这是纯大小无关的题了，我们是有统一的一种处理办法的：不妨设一个大小。现在我们不妨设 \(a=1\) ，于是已知变成 \(c=\frac{\sqrt{3}}{3} \sin B+\cos B\) ，仍然是刚才的 \(a, B, c\) 的信息，不过我们看得比之前清晰得多：\(a=1\) 是确定的线段长，给一个 \(B\) 吐一个 \(c, B\) 直接彻底拿捏 \(\operatorname{SAS}\) 控制，一路下来纯粹的控制，思路比刚才的牵制顺畅得多！所求也等价于求 \(A D\) 的最大值，这时就相当于 \(B\) 按照 SAS 控制全局，让我们直接输出 \(A D\) 罢了。这个思维上的简化太大了，直接从 \(a, B, c\) 控制 \(A D\) 长，基础题中的基础题，写出控制式 \(|\overline{A D}|^{2}=\frac{1}{4} a^{2}+c^{2}-a c \cos L\) 是完全不需要思考的过程！剩下的就是代人化简得到 \(|\overline{A D}|^{2}=\frac{\sqrt{3}}{6} \sin 2 B-\frac{1}{6} \cos 2 B+\frac{5}{12}\) ，求出 \(A D\) 的最大值是 \(\frac{\sqrt{3}}{2}\) ，然后带回 \(\frac{a}{A D} \cdot 2 \sqrt{3}\) 即可。思路相比原题而言，只剩纯粹的控制过程，完全没有任何选取主自由度的难度，也不需要去思考什么形状参数 \(k\) ，因为对大小与形状的处理已经在题目改写的过程中全部搞定了。

第二种解法的思想其实就是化归：把陌生问题转化为熟悉问题。题目里有明显的区分大小与形状的条件的时候，大家就要瞪起眼来，注意这题很可能是要分开处理形状与大小。而对于这种可以把所求改写成大小无关的情况，则有希望把问题充要改写为大小无关的新问题，而由于是充要替换，因此这样两个问题的答案应当是一样的。从刚才这道题大家可以看出，计算总量没有实质性的改变，判断定义域，化简三角函数式，求最值一样不少，处理的元问题也基本一样，但是在思维层面带来的简化却是实质性的。第一种解法我们不仅要考虑牵制转化，还要主动分离大小与形状，但第二种解法则只剩下一条单线的控制，要素关系得到大大简化，几乎不怎么需要思考，可以提笔就写。因此，这样的转化节省的是思维时间和思维容量，有助于我们更快看清问题本质，更快完成解题。希望大家通过本次视频，更深刻地认识到大小与形状的视角在几何问题中的重要性，更好地将其应用于解题中。

\section*{课后练习：}
冰雪聪明的 Cirno 同学在看完本视频后，又遇到了这样一道问题：\\
己知在 \(\triangle A B C\) 中，角 \(A, B, C\) 所对的边分别是 \(a, b, c\) ，且 \(B-C=\frac{\pi}{2}, a=\cos B\) 。则 \(a+\frac{\sqrt{3} b}{\cos B} \cdot \sin A\) 的取值范围是 \(\qquad\) －

冰雪聪明的 Cirno 同学立刻意识到这道题类似于我们刚才讲的问题，于是她先改写了所求，并将其等价变形为大小无关的 \(\cos B+\sqrt{3} \sin B\) 。随后她抛弃已知的 \(a=\cos B\) ，选取 \(B\) 为主自由度，并通过 \(B-C=\frac{\pi}{3}\) 求出了 \(B\) 的取值范围是 \(\left(\frac{\pi}{3}, \frac{2 \pi}{3}\right)\) ，进而求出题目答案为 \((1,2)\) 。可是当她翻看答案的时候，却发现答案是 \((\sqrt{3}, 2)\) 。她百思不得其解错因在哪。\\
（1）Cirno 同学是如何将原所求等价变形成 \(\cos B+\sqrt{3} \sin B\) 的呢？请结合视频的思想与方法，补充完整这一过程。\\
（2）你能帮帮 Cirno 同学，告诉她错因在哪吗？这种错误的产生是视频里强调的哪个注意点没做好导致的呢？我们应当如何规避这种错误？

\section*{6．3．5 看似设问奇怪，实则还是思维}
视频日期：2024年6月23日

\section*{题目：}
在 \(\triangle A B C\) 中，角 \(A, B, C\) 所对的边分别是 \(a, b, c\) 。已知 \(\triangle A B C\) 的外接圆半径为 1 ，且满足 \(\sin 2 A+\sin 2 B+\sin 2 C=\) 1，则 \(\triangle A B C\) 的面积为 \(\qquad\) ；当 \(A\) 取得最大值时，\(a^{4}-8 a\) 的值为 \(\qquad\) －

这道题目比较有趣，设问方式也比较新颖。我们来分析一下。首先已知给了两个牵制条件，而放在具有 3 个自由度的三角形系统内，应当起到减少两个自由度的作用，因此系统应当还有 1 个自由度。但是上来第一个空就问我们一个定值问题。因此这道题要么是不变量问题，要么是临界问题。第二个空又提到了最大值，因此应该不是临界问题（否则系统卡死，问最大值就没啥意思了）。所以这道题大概率第一个空就是不变量问题。然后第二个空自然就是单自由度求最值的问题。

动静视角的分析清楚了，我们来建立要素关系网。但是上来就让人不爽。两个牵制很明显一个管大小，一个

管形状。 \(\sin 2 A+\sin 2 B+\sin 2 C=1\) 把形状减掉 1 个自由度，然后外接圆半径固定大小。但是用外接圆半径固定大小可真不常见。平时我们大多拿某条边的长度来固定大小，这样可以先白嫖两个静态点。可是把外接圆半径给定出来就有一定的困难了，画图也似乎找不大到抓手。已知的角度牵制看上去也没啥好的信息。这时候该怎么办呢？

我们讲三个方法。第一个方法，既然是不变量问题求最值，那我把答案试出来不就行了。反正一个自由度余地很大，应该是可以挺随便地取特殊值。比如 \(A=\frac{\pi}{2}\) ，带回去得到另一组解 \(15^{\circ}, 75^{\circ}, ~ 90^{\circ}\) ，再计算也行。

第二个方法，既然题目的已知是一个管大小，一个管形状，根据我们2023年9月29日视频中的方法，我们可以做一个强行的大小无关转换来抓出问题本质：把所求的 \(S\) 改写为

\[
\frac{S}{R^{2}}=\frac{a b \sin C}{2 R^{2}}=2 \sin A \sin B \sin C
\]

抛弃外接圆的条件。因此问题进一步等价为：在 \(\triangle A B C\) 中，已知 \(\sin 2 A+\sin 2 B+\sin 2 C=1\) ，求 \(2 \sin A \sin B \sin C\)的值。上述视频内的方法告诉我们：这个问题改写是充要替换，不丢失信息量，一定是可解的。因此我们把思维焦点放过来，现在是纯角度推导的时间。这里需要一点二级结论储备，我们用和差化积公式：

\[
\sin 2 A+\sin 2 B=2 \sin (A+B) \cos (A-B)
\]

这里 \(\sin (A+B)=\sin C\) ，因此代人原式左侧整理得到 \(2 \sin C[\cos (A-B)+\cos C]\) 。再把 \(\cos C=-\cos (A+B)\)换进去，化简得到 \(4 \sin A \sin B \sin C\) 。于是 \(4 \sin A \sin B \sin C=1\) ，这就得到 \(\frac{S}{R^{2}}=\frac{1}{2}\) ，即本题答案。

第三个方法比较巧妙。既然它现在要我们用外接圆半径定大小，那我就画一个半径为 1 的圆，把三角形放进去。比如我先画个锐角 \(\triangle A B C\) 示意图，圆心取为 \(O\) 。一个巧妙的观察是 \(\angle A O B=2 C\) ，再加上半径为 1 的圆这个背景，联系到角度的牵制条件和面积的所求，我就注意到

\[
S_{\triangle A O B}=\frac{1}{2}|A O| \cdot|O B| \cdot \sin \angle A O B=\frac{1}{2} \sin 2 C .
\]

正好，这三个小三角形拼起来就是大三角形 \(\triangle A B C\) 了！因此大的面积

\[
S=\frac{1}{2}(\sin 2 A+\sin 2 B+\sin 2 C)=\frac{1}{2} .
\]

就这么做完了！！不过需要注意刚才实际上只讨论了锐角三角形的情况，为了论证严谨还需要补充直角和钝角三角形的情况（留作练习），不过这已经足够说明这个几何解法的巧妙了。

不管怎么说，三种方法对应了三种不同的思考，但得到第一个空的答案并不困难。在开始第二个空之前，我们要再强调一下之前讲到过的有关充要替换的一个重要方法，就是自由度调节。这道题总自由度是 1 ，信息量等同于三角形系统外加两个牵制，但是选取哪两个牵制呢？现在有三个等式：外接圆半径，三角函数式，面积值。从信息量的角度上来说，任选两个牵制都可以贡献同样的信息量，因此我们直觉上相信这是个充要替换（在细微的定义域讨论之后）。事实上确实如此，如果我们选取外接圆半径 \(R=1\) 和面积值 \(S=\frac{1}{2}\) 作为这个题的已知条件，把那个鬼畜的三角函数式扔掉，回顾刚才第一个空的计算过程，显然三角函数式是可以重新推出来的。因此这个充要替换是合理的，而且 \(S=\frac{1}{2}\) 显然在构建要素网的时候会比那一坨三角函数舒服得多。自由度调节是充要替换的重要方法之一，在难题中的重要性不亚于遍历法，大家一定要注意。

下面，我们用 \(R=1\) 和 \(S=\frac{1}{2}\) 作为新已知（信息量与原题一致），来看 \(A\) 何时取最大值。首先配合画图法来试着选主自由度。综合已知，目前唯一能画的就是：先画个半径为 1 的圆。外接圆半径涉及到的核心元问题是正弦定理，也就是 \(\frac{a}{\sin A}=2\) 。显然，在外接圆里画一条边（弦）看着顺很多，比强行画个 \(A\) 角舒服多了。于是画完圆之后，我们再画一条弦 \(B C\)（它占据一个主自由度 \(a\) ）。跟踪定义域：显然必须 \(0<a \leq 2\) 。不难看出，弦 \(B C\)画出来（固定这个主自由度）之后，让点 \(A\) 在圆上跑，则角 \(A\) 只能取两个固定值。

已经把主自由度给 \(B C\) 了，因此 \(A\) 不能再自由跑了。为了把 \(A\) 的位置确定下来，就需要第二个牵制：\(S=\frac{1}{2}\) 上场了。面积的元问题最常见的有三种：SAS 型控制，底乘高除 2 ，以及坐标系中的分割加减法。现在放在外接圆里跑，SAS 控制没有抓手，但是刚好一条底边 \(B C\) 拿来做主自由度了，因此自然想到底乘高除 2 的路子。设 \(B C\)边上的高线长度为 \(h\) ，则

\[
\frac{1}{2} a h=S=\frac{1}{2} \Leftrightarrow h=\frac{1}{a} .
\]

这下要素网就清楚了：先把主自由度给 \(a \in(0,2]\) ，在固定的圆里画出弦长 \(a\) ，然后平行于该弦，距离为 \(h=\frac{1}{a}\) 作两条平行线，和圆的交点就是合法的 \(A\) 点。这个要素网的每一步都是常见的控制，因此我们信心满满，这个网肯定能解出答案来！

开始尝试解题。 \(a\) 对于系统的控制我们还不完全清楚，主要是这两条平行线的位置还没研究明白。因此可以搞点思想试验先直观看看。定义域是 \(a \in(0,2]\) ，不妨先看看 \(a=2\) 的情形：\(B C\) 是直径，平行线到 \(B C\) 的距离是 \(\frac{1}{2}\) ，显然是一个确定的直角三角形，此时 \(A\) 当然是直角。我们要求 \(A\) 的最大值，没想到起步就是直角，前途无量啊。然后我们尝试用一下遍历法，稍微缩小一点 \(a\) 的值。比如取 \(a=1.99\) ，则平行线到 \(B C\) 的距离是 \(\frac{1}{1.99}\) ，约等于 \(\frac{1}{2}\) 但略大。这时 \(B C\) 相比直径稍微往下推了一点，而两条平行线相对 \(B C\) 参考系，则向外分开了一点，这导致 \(B C\) 下方的平行线往下推了两小步。这时显然两条平行线与外接圆仍然有 4 个交点，但是对应的 \(\angle B A C\) 却是两个锐角和两个钝角，它们显然刚好互补，正是 \(\sin A=\frac{a}{2}\) 的两个解。我们是要求 \(\angle B A C\) 的最大值啊，怎么可能要锐角，肯定是取下面的针角。然后继续往下推 \(B C\)（对应 \(a\) 继续减小），两条平行线继续与外接圆有 4 个交点，不过 \(B C\) 下方的平行线却离 \(B C\) 越来越远，随时有跌到外接圆下方的危险。但是，往下推 \(B C\) 的收益高啊， \(\sin A=\frac{a}{2}\) 的钝角解是越来越大的，我们还得硬着头皮往下推 \(B C\) 。 \(\frac{1}{a}\) 在 \(a\) 下降到 0 时会趋于正无穷，因此终有一个时刻，下方的那条平行线会刚好和圆底部相切。这时候再往下推一点，这条平行线就和圆没有交点了，我们的钝角解也就消失了，只剩锐角解了。再往下推 \(B C\) ，针角解也回不来了。我们要求的是 \(A\) 的最大值了，怎么能允许这种事情发生！！因此，这就告诉我们，\(A\) 的最大值就刚好是在 \(B C\) 下方的平行线与外接圆底部相切时取到，这正是钝角解能达到的最大值，自然也是 \(A\) 的最大值！！

显然，这已经给出了 \(A\) 取最大值时的静止化条件。我们画出图来，以 \(a\) 为主自由度即可轻松列出方程：

\[
\left(1-\frac{1}{a}\right)^{2}+\left(\frac{a}{2}\right)^{2}=1
\]

稍加化简即得 \(a^{4}-8 a=-4\) ，问题解决。

做个总结。这道题的两个空质量都不错。第一个空考察的是对于大小与形状的理解，这个我们之前为大家重点讲过了。投机取巧的话可以猜答案，靠一点灵感还能找到几何解法，但不管怎么说，搞出答案还是不难的。紧接着是一个很重要的充要替换，基于对自由度的理解。这个充要替换服从＂简化即胜利＂的原则，为我们构建第二问的要素关系网打好了基础。随后第二问就是一个动静视角的基本理念：单自由度求最值，关键是找出最值点的静止化条件，然后用好思想试验探索与遍历法的思想稳步求解即可。

\section*{课后练习：}
（1）使用充要条件的定义，验证我们为第二个空所作的条件替换确实是充要替换。\\
（2）\(B C\) 上方的那条直线如果推的太高会怎么样呢？你能算出本题 \(a\) 的取值范围吗？你能算出 \(A\) 取最小值时的系统状态吗？（写出 \(a\) 需要满足的等式或不等式即可）

\section*{6．3．6 趣味性题目纯思维历程分享}
\section*{题目：}
对于 \(\theta>0\) ，若存在实数 \(\varphi\) ，使得对于 \(\forall n \in N^{*}\) ，都有 \(\cos (n \theta+\varphi)<\frac{\sqrt{3}}{2}\) ，则 \(\theta\) 的最小值为 \(\qquad\) ．

题目来自上海春考，是一道颇有难度的趣味性题目。虽然难度很大（我也不要求大家非得自己做出来），但是完全听懂讲解，理解其中的逻辑，绝对是一次非常好的逻辑练习。

这道题的题干就比较绕，里面带了两个量词，也就是我们在＂经验分享番外篇＂中讲到的封装命题，还是两层，我们需要从内侧向外侧扒，从而理解清楚问题的内涵。最内层的是

\[
" \text { 对于 } \forall n \in N^{*} \text {, 都有 } \cos (n \theta+\varphi)<\frac{\sqrt{3}}{2} ",
\]

也就是相当于从 \(\varphi\) 开始不断往上添 \(\theta\) ，其中的余弦值都要在 \(\left[-1, \frac{\sqrt{3}}{2}\right)\) 内，换言之，\(n \theta+\varphi\) 一定在某个范围内，即

\[
\left(\frac{\pi}{6}+2 k \pi, \frac{11 \pi}{6}+2 k \pi\right)
\]

但是我们很快就发现一个问题，这样一个充要替换是非常难处理的。且不说外面还套了一层封装，你不断加 \(\theta\) ，还要求时刻控制着 \(n \theta+\varphi\) 在什么位置，这是一个相当大的工作量。而且，这个范围

\[
\left(\frac{\pi}{6}+2 k \pi, \frac{11 \pi}{6}+2 k \pi\right)(k \in \mathbb{Z})
\]

也很怪异，是实数轴上一段一段均匀分布的开区间段。你要控制 \(n \theta+\varphi\) 时刻都跳在这里面，光是想象一下就头大，完全没有方向感。

我们一直以来讲过的思维方式，包括动静视角，要素关系网，隔离等等，很多都是为了尽可能让我们直观地看到问题的结构，清楚接下来该干什么的。换而言之，就是可视化。但这个题，在扒掉一层封装之后，仍然非常混沌，毫无入手想法。我们在经验分享第五期中说过，简化即胜利，做题需要时不时出现一些简化的行为来让我们保持继续做下去的动力，但是刚才的转化无疑是失败的，一点简化都没看出来。于是，我们不禁怀疑，刚才的充要替换真的能舒服地把题做下去吗？我们不得不重新审视这个转化。转化唯一需要的知识就是三角函数，我们开始思考，既然这个转化很臭，那有没有其他转化方式，能至少哪怕一点点，让人感到舒服呢？

这个时候，我尝试回归根本，回忆一下：三角函数从哪里来？我学过什么三角函数值的知识？三角函数的定义来源于单位圆上角终边的坐标，随后出现了诱导公式，恒等变换等一系列关系。作为函数，它本身的函数的性质也有不少。刚才我们做的求值把整个问题弄的很混沌，那么还有什么路子可以尝试？看看刚才列举的东西，不难注意到其中有一条＂定义＂，就是单位圆上的点坐标。通过这个方式确定余弦值的路子是，先确定角的终边，再直接读取横坐标。所以 \(n \theta+\varphi\) 所在点要一直在 \(x<\frac{\sqrt{3}}{2}\) 的区域内，也就是，除了最右边的张角 \(\frac{\pi}{3}\) 的闭圆弧以外的任何地方。换而言之，\(n \theta+\varphi\) 始终不能踏进这段闭圆弧。这个理解让我们瞬间感觉很舒服，因为它具有极强的可视性，整个问题的几何直观仿佛立刻呈现在我们眼前。这无非就是说 \(\varphi\) 不断往上加一个一个一个的 \(\theta\) 角，都始终踏不进刚才那段闭圆弧。这个问题立刻就变得好像可以着手思考了一样，让人不自觉的感到舒服。（这种几何直观带来方向感的感受，是一种数学直觉，仿佛在告诉我们继续沿这条路走下去一定能有好的结果。当难度达到一定高度时，直觉不得不参与进来，但这题的直觉仍然与＂简化即胜利＂原则有较强的关联）

本着＂简化即胜利＂的原则，我们继续往下走。这一层的直观感受已经很强了，我们再往下扒掉一层封装。我们需要存在一个 \(\varphi\) ，使得往上加 \(\theta\) 角，都进不去闭圆弧。这种封装结构，我们自然还是用番外篇讲过的处理封

装结构的思维方式去看，遍历加冻结加设问。我们摸一摸 \(\varphi\) ，随便停下来看看，自问一下这样的 \(\varphi\) 满不满足？比如我在单位圆上随便找一点 \(\varphi\) ，往上加几段 \(\theta\) 看看，加一加看看能不能进入刚才那段闭圆弧内？这样一加就知道了，比如一开始 \(\varphi=\frac{5 \pi}{3}\) ，再加一下就是 \(\frac{5 \pi}{3}+\theta\) ，如果这就直接踏进去了，那就完蛋了，所以不能让 \(\frac{5 \pi}{3}+\theta\) 进去，也就是说，如果把闭圆弧 \(\left[-\frac{\pi}{6}, \frac{\pi}{6}\right]\) 往回转（顺时针转）\(\theta\) 角，得到一个新的闭圆弧，那么 \(\frac{5 \pi}{3}\) 不能在这里头。那我 \(\frac{5 \pi}{3}\) 加俩 \(\theta\) 呢？那就是说 \(\frac{5 \pi}{3}+2 \theta\) 也不能踏进去，换而言之，把闭圆弧 \(\left[-\frac{\pi}{6}, ~ \frac{\pi}{6}\right]\) 往回转 \(2 \theta\) 角得到一个新的闭圆弧，\(\frac{5 \pi}{3}\) 也不能在里头。

这样一冻结，一摸索，再设问一下，那 \(\varphi\) 究竟能站在哪？我们就立刻明白了。把闭圆弧 \(\left[-\frac{\pi}{6}, \frac{\pi}{6}\right]\) 往回不断转 \(\theta\) ，转出一大堆（无限个）闭圆弧，你的 \(\varphi\) 必须精准身朵开这些闭圆弧才行。那如果要存在 \(\varphi\) ，自然就是说，这一大堆闭圆弧盖不满整个单位圆！反之显然也对，只要没盖满，就从漏网之鱼里面随便挑一个 \(\varphi\) 就行了。所以这当然是充要替换，问题等价为，求 \(\theta\) 的最小值，使得闭圆弧 \(\left[-\frac{\pi}{6}, \frac{\pi}{6}\right]\) 不断往回转 \(\theta\) 得到一堆闭圆弧盖不满单位圆。

现在替换出来的这个问题里面，没有封装了，量词全被消去了，而且几何直观仍然很强，所以我们抛弃原来的题干，转而看刚才这个等价表述。从位置无关的角度来看，其实就是说一段张角 60 度的闭圆弧往回转 \(\theta\) 角的故事。我们很快就会发现，因为 \(n\) 是无穷多个，所以稍微转一会，这一堆闭圆弧所能盖到的区域似乎就会变得很＂多＂。比如说，如果 \(\theta\) 比较小，那转着转着就铺满了。但是，也容易看出来，有一些 \(\theta\) 是满足条件的，比如 \(\theta=2 \pi, ~ 一\)转转一圈，再或者 \(\theta=\pi\) ，一转转半圈，两次就开始循环了。从这些直观与试验中，我们可以得到的初步直观是，你不停的转，如果最后不能重合，形成一个周期，那么随着越转越多，早晚会被盖满的。

如何来严格表述这种直观呢？现在假设我们转五次，连带一开始的 \(\left[-\frac{\pi}{6}, \frac{\pi}{6}\right]\) ，总共形成了六段闭圆弧，每段的张角都是 \(\frac{\pi}{3}\) 。这六段闭圆弧加起来的长度和已经是一整个圆周了，所以必然会有其中两段出现相交。随便打个比方，比如初始的第 0 段和第 4 段有交点。那么，如果这两段不是完全重合，那么它们会拼成一个更长的圆弧。也就是说，每转 \(4 \theta\) ，新圆弧就会和老圆弧拼上，拉长覆盖的占比。那么，第4段再转 \(4 \theta\) 到第 8 段，这两段又会把刚才的＂更长的圆弧＂继续拉长（注意第 0 段和第 4 段的相对位置与第 4 段和第 8 段的相对位置一样），并且每 \(4 \theta\) 带来的拉长距离都是一样的。这样的话，迟早会拉长到占满整个单位圆，这与我们的期望是矛盾的。

因此，如果在转的过程中，某两段出现了相交，那必须让这两段完全重合，否则不重合且相交的两段就可以类似的不断拉长，直至铺满整个单位圆。因此，在转的过程中，要么任何两段都不相交，只要相交就一定刚好重合，进入周期循环。这就告诉我们，在闭圆弧不断旋转 \(\theta\) 角，产生更多的圆弧段的时候，我们的视觉观察就是，一段一段的分离的闭圆弧罢了（出现重合时进人了周期循环，视觉上完全重合的就一样了），而且显然这样观察到的闭圆弧就是很少几个（刚才说过了，六个闭圆弧就一定会出现相交，从而产生重合，进入周期循环）。因此，场上实际观察到的闭圆弧最多也就 5 个罢了，其余的全都是重合。

既然没几个，那么完全就可以枚举了。场上能观察到的分离的圆弧段可能有 \(1,2,3,4,5\) 个，如果就 1 个，那就是一次转整数圈；如果是 2 个，那一次就转半圈（加整数圈）；如果是 3 个，那一次就转 120 度，或者 240 度之类的；如果是 4 个，那就一次转 90 度，或 270 度之类的；如果是 5 个，那就一次转 \(\frac{360}{5}=72\) 度之类的。都能枚举了，那自然找最小值也就很轻松了。很明显，观察到的圆弧越多，那么一次转动的角度就可以更小，场上最多存在 5 个，那最小的转角就是 \(\frac{2 \pi}{5}\) 了，这就是最后的答案。

简单总结一下。这道题是一道思维难度很高的问题，也充满趣味性。最开始需要想到把问题搬到单位圆上面去。这需要一定的联想能力（重要的是回归根本的意识），然后处理封装命题，扔掉量词，变成一个直观性很强的问题。后续的处理需要先抛弃原题，把思维焦点放到新的简化问题——闭圆弧转圈上，然后论证清楚其中的逻辑，这并不容易，需要非常好的逻辑思考能力。这些逻辑上的难点让这道题的难度达到了一个较高的水平。想清楚这个

提到论证逻辑，本身就是一种思维和逻辑的练习。

\section*{思考题}
1．请结合讲解中的讨论，尝试写出题目中 \(\theta\) 的所有可能取值构成的集合\\
2．仿照视频中论证 \(\theta\) 角取值的逻辑，结合下面的提示，完成下题的证明：

\section*{题目：}
对于 \(x \in \mathbb{R}\) ，定义高斯函数 \([x]\) 为不超过 \(x\) 的最大整数，例如 \([\sqrt{2}]=1,[-3.5]=-4\) ，再定义函数 \(\{x\}=x-[x]\) ，称为 \(x\) 的小数部分。不难看出 \(\{x\} \in[0,1), \forall x \in \mathbb{R}\) 。\\
对于 \([0,1)\) 的非空子集 \(A\) ，称 \(A\) 在 \([0,1)\) 上稠密，若对 \(\forall x \in[0,1)\) 以及 \(\forall \varepsilon>0\) ，均 \(\exists y \in A\) ，满足 \(|x-y|<\varepsilon\)（直观上就是说 \(A\) 中的元素可以以任意小的尺度接近 \([0,1)\) 中的任何元素）。比如 \([0,1)\) 中的全部有理数构成的集合就在 \([0,1)\) 上稠密。

现在，对于 \(\alpha>0\) ，我们定义集合 \(A_{\alpha}=\left\{\{n \alpha\} \mid n \in N^{*}\right\}\) ，显然 \(A_{\alpha} \subseteq[0,1)\) 。求证：\(A_{\alpha}\) 在 \([0,1)\) 上稠密当且仅当 \(\alpha\)是无理数

提示：\\
（1）考虑集合 \(B_{\alpha}=\left\{2 \pi x \mid x \in A_{\alpha}\right\}\) ，证明：把 \(B_{\alpha}\) 中的元素看作角，把这些角的终边对应的点标在单位圆上，那么这些点就是所有形如 \(2 \pi n \alpha\left(n \in N^{*}\right)\) 的角度对应的点，把这些点构成的集合记为 \(P_{\alpha}\)\\
（2）证明：如果 \(\alpha\) 是有理数，则 \(P_{\alpha}\) 只有有限个点构成，从而不可能 F 密\\
（3）证明：如果 \(\alpha\) 是无理数，则所有的 \(2 \pi n \alpha\left(n \in N^{*}\right)\) 对应的点两两不同（提示：利用反证法，想想看如果存在两个相同的点，会发生什么）\\
（4）借助第三问的结论，并模仿视频中原题讲解的一些论述思路，证明：对 \(\forall \theta \in(0, \pi), P_{\alpha}\) 中存在两点 \(2 \pi n_{1} \alpha\) 与 \(2 \pi n_{2} \alpha\left(n_{1}<n_{2}\right)\) ，使得这两点在单位圆上对圆心的张角小于 \(\theta\)（也就是说 \(P_{\alpha}\) 中存在离的任意近的点）\\
（5）借助第四问的结论，并体会视频原题讲解中提到的相对位置一说，论述：如果 \(\alpha\) 是无理数，则 \(A_{\alpha}\) 在 \([0,1)\) 上稠密

\section*{6．3．7 XXXX}
\section*{6.4 新定义问题}
\section*{6．4．1 九省联考压轴题如何拿高分？}
视频日期：2024年1月20日

观前提醒：本视频是对一道近期热度较高的数学压轴题的讲解视频。讲解时用到的思维方式和思想都是 up主往期重点视频中重点强调的思想。视频中用到的核心思想为＂高中数学经验分享系列＂中分析系统结构与要素关系的思想，以及＂数学的干货＂视频列表中出现过的分析法（投稿日期为2023年2月5日）。不过视频中使用时会争取让还未观看过这些视频的同学也能够理解这些思想的精髓。为了尽量还原高中生临场时的思考，视频内将不引入带余除法，同余，剩余系，费马小定理等数论名词，对背后蕴含的概念也不加深入，主要突出考场

思维引导与拿分策略。感兴趣的同学可进一步观看其他 up 主对本题的讲解与深挖．

\section*{题目（为2024年九省联考数学压轴题）：}
离散对数在密码学中有重要的应用。设 \(p\) 素数，集合 \(X=\{1,2 \cdots p-1\}\) ，若 \(m \in N, u, v \in X\) ，则记 \(u \otimes v\) 为 \(u \cdot v\) 除以 \(p\) 的余数，\(u^{m, \otimes}\) 为 \(u^{m}\) 除以 \(p\) 的余数；现设 \(a \in X\) ，满足 \(1, a, a^{2, \otimes} \ldots a^{p-2, \otimes}\) 两两不同，若 \(a^{n, \otimes}=b(n \in\{0,1 \cdots p-2\})\) ，则称 \(n\) 是 \(b\) 的以 \(a\) 为底的离散对数，记 \(n=\log (p)_{a} b\) 。\\
（1）若 \(p=11, a=2\) ，求 \(a^{p-1, \otimes}\) ；\\
（2）对 \(m_{1}, m_{2} \in\{0,1 \cdots p-2\}\) ，记 \(m_{1} \oplus m_{2}\) 为 \(m_{1}+m_{2}\) 除以 \(p-1\) 的余数（当 \(m_{1}+m_{2}\) 能被 \(p-1\) 整除时，规定 \(\left.m_{1} \oplus m_{2}=0\right)\) 。证明： \(\log (p)_{a}(b \otimes c)=\log (p)_{a} b \oplus \log (p)_{a} c\) ，其中 \(b, c \in X\)\\
（3）已知 \(n=\log (p)_{a} b\) ，对 \(x \in X, k \in\{1,2 \cdots p-2\}, y_{1}=a^{k, \otimes}, y_{2}=x \otimes b^{k, \otimes}\) ，证明：\(x=y_{2} \otimes y_{1}^{n(p-2), \otimes}\)

拿到题目，乍一看非常新颖，各种乱七八糟的定义，还出来余数了，高中都从来没练过，怎么办？其实，不管什么题目，拿到手时，先对题目讲了一个什么事进行梳理，把题目描述的体系理解清楚，把题目中出现的各个要素生成的先后顺序搞清楚，都是很有意义的。这就是结构为上，优先解剖题目结构，然后通过系统的结构来思考如何解题的思维方式。比如这道题，我们先看题干，先给了一个素数 \(p\) ，有了 \(p\) 就可以定义集合 \(X\) ，这些构成了题目的基本背景。然后，在这个系统里出现了两个运算：\(u \otimes v\) 与 \(u^{m, \otimes}\) 。虽然长得挺别致，但是它就是我们传统认知的乘法与乘幕再取个除 \(p\) 的余数罢了。以上就是由一个素数 \(p\) 生出的背景系统。

接着，题目给出了一个 \(a, a\) 是几？我们不知道，但是它有一个很强的限制（翻译成我们熟悉的人话）：\(a\) 取从 0 到 \(p-2\) 这些幂次，得到 \(p-1\) 个数，它们除以 \(p\) 的余数两两不同。天哪！一个数除 \(p\) 一共才可能出几个余数啊，不就 0 到 \(p-1\) 这 \(p\) 个可能吗？现在直接占掉 \(p-1\) 个了！就比全占少了一个数！可能有的同学会好奇，少了谁，能知道吗？我们先不急着思考这个问题，先继续看。若 \(a^{n, \otimes}=b\) ，也就是刚才这一堆余数互不相同的幂次 \(a^{n, \otimes}(n \in\{0,1 \cdots p-2\})\) ，我把余数记为 \(b\) ，然后就定义出来了离散对数 \(n=\log (p)_{a} b\) ，长得还挺像大家以前学的对数的。因此，题目的系统就清晰了：先是 \(p\) 决定的一个静态背景，然后在赋予了一个很牛掱的 \(a\) 之后，就可以定义一种新型对数。这种新型对数我们可以发现，\(a^{n, \otimes=b}\) 实现的是 \(n\) 对 \(b\) 的控制，但是反过来，写成 \(n=\log (p)_{a} b\)的时候，就是 \(b\) 对 \(n\) 的控制。这说明在给出这个牛羒 \(a\) 之后，\(b\) 和 \(n\) 可以建立一个互相控制。Up 主一直以来都给大家强调，要多关注题目中各个要素的控制关系。这下倒好，题目直接把控制式甩在你脸上，明摆着直说＂\(b\)和 \(n\) 在现在的余数语境下也可以建立相对控制，元问题就是题目定义的取余数版指数与离散对数＂。

但是，up 主经常给大家强调，控制肯定得关注定义域。我现在是把这 \(n-1\) 个余数全拿出来，它们的集合就是 \(b\) 的定义域。这个定义域是啥啊？没说啊？能不能猜出来呢？以前我们学过，对数定义域不含 0 。现在呢？0能不能作为某个 \(b\) 的取值呢？那等价说法就是能不能存在某个 \(n\) ，使得 \(a^{n}\) 是 \(p\) 的倍数呢？这怎么可能？？你本身 \(a\) 就比 \(p\) 小，\(p\) 又是素数，所以 \(a\) 肯定没 \(p\) 这个因数，那你不断做多少乘法也变不出新的因数啊。所以这下懂了，这堆 \(a^{n}\) 除 \(p\) 的余数没 0 ，正好是 \(1 \cdots p-1\) 这 \(p-1\) 个数！！这下控制彻底通了，这堆余数互不相同，并起来刚好就是 \(1 \cdots p-1\) ，然后反向对应回去，取一个 \(b \in\{1 \cdots p-1\}=X\) 反回来找：取多少的幂次（即 \(a^{n}\) ），对应的余数是 \(b\) 啊？这个幂次 \(n\) 就是离散对数。

题目的要素结构关系分析完之后，我们来解题（大家可以看到，up 主的做题习惯就是先把要素结构分析清

楚，分析要素结构的基本原则在上面的讲解中也体现出来了）。第一问是一个纯静态的问题。所谓＂静态＂就是各个要素都有确定的位置，没有任何可以变化（运动）的参数。基本背景 \(p\) 和牛羒神数 \(a\) 都告诉你了，你只需要计算即可，没有任何难度，算就行。 \(a^{p-1}=2^{10}=1024\) ，除一下 11 即可，答案是 1 。

下面看第二问。又给了一个 \(m_{1} \oplus m_{2}\) ，表示加完除 \(p-1\) 的余数（注意这里是 \(p-1\) ，不是 \(p\) 了。Up 主在多个套卷视频中多次强调审题要画圈，把要点圈出来，你做到了吗）。然后让我们证明对于 \(b, c \in X\) ，有 \(\log (p){ }_{a}(b \otimes c)=\) \(\log (p)_{a} b \oplus \log (p){ }_{a} c\) 。如何解读这个问号？仍然秉持我们讲过的分析结构的原则，我们先来分析一下这个问号的系统。首先，题干原有的背景 \(p\) 与一个给定的元问题仍然继承。但是这时 \(b, c \in X\) 是任意给的，所以我们可以把它们看成两个独立的，运动的参量，他们就像操纵杆一样控制了整个系统。这和我们在经验分享第四期讲到的主自由度的理念不谋而合，只不过这里这两个参数只能取有限个值，与以前的连续运动不一样。但是，思想是相通的。

然后该解决问题了。所证式呢？式子里有三个对数，它们反映了 \(b\) 的指对相互控制，\(c\) 的指对相互控制（显然在这个小问）以及 \(b \otimes c \in X\) 的指对相对控制。式子的左右两侧都是被操纵杆 \(b, c\) 控制的东西，并且它们的语义都很明确——给定了 \(b, c\) ，这俩都是确定的数。如何研究具体问题？那就要把控制关系对应到元问题上去了。刚才我们分析了，这道题是题目现场给你一个元问题——＂ \(0 \cdots p-2\) 当幂次与 \(1 \cdots p-1\) 取对数之间可以一一对应互相控制＂，而且离散对数这个对象我们目前只知道这一个元问题，所以肯定就要用它。既然是互相控制，而且显然取对数的控制我们看不清楚，但反过来的求幂次控制我们看得很清楚，因此方向就很明确了：我们需要把幂次转回指数去，这样才能让问题变得可以处理。操纵杆 \(b, c\) ，我给转成 \(x, y\) ，满足 \(a^{\mathrm{x}}\) 除 \(p\) 余 \(b, a^{\mathrm{y}}\) 除 \(p\) 余 \(c\)。等式右边就成了 \(x \oplus y\) ，那左边怎么处理？现在我把操纵杆给了 \(x, y\) ，由它们控制 \(b, c\) ，然后再让 \(b, c\) 去控制出更复杂的 \(\log (p)_{a}(b \otimes c)\) 吗？这个控制出来，你也很难研究，更别说它是不是和 \(x \oplus y\) 相等了。那怎么办？我们之前讲过，在碰到存在双向控制的时候，颠倒一下要素关系网，让两个要素的控制关系反过来，常常有新的发现。我们故技重施：我不搞对数了，我只需要证明 \(a^{\mathrm{x} \oplus \mathrm{y}}\) 除 \(p\) 的余数是 \(b \otimes c\) 即可。因此，现在我只需要证明：\(a^{x}\) 除 \(p\) 余 \(b, a^{\mathrm{y} \Theta \mathrm{y}}\) 除可推出 \(a^{x \oplus y}\) 除 \(p\) 余 \(b \otimes c\) 。

怎么证明呢？\(a^{\mathrm{x} \oplus \mathrm{y}}\) 看者像个指数加法。 \(x \oplus y\) 是啥？相加，然后除掉 \(p-1\) 取余数，和真正的加法差了个余数。我们遍历一下 \(x\) 和 \(y\) 看看：都是从 0 到 \(p-2\) 里随机选数，相加就是从 0 长到 \(2(p-2)\) 。那这 \(x \oplus y\) 不就要么是 \(x+y\) ，要么就是 \(x+y-(p-1)\) 吗？真就和真正的加法差不多。所以要么就是算 \(a^{\mathrm{x}+\mathrm{y}}\) 的余数，要么就是 \(a^{\mathrm{x}+\mathrm{y}-(\mathrm{p}-1)}\)的余数。事已至此，\(a^{\mathrm{x}+\mathrm{y}}\) 是绕不开了，必须得乘一下。既然要乘 \(a^{\mathrm{x}}\) 和 \(a^{\mathrm{y}}\) ，就得把它们都放等号左边。我们目前的要求是 \(a^{x}\) 除 \(p\) 余 \(b, a^{y}\) 除 \(p\) 余 \(c\) ，怎么写成式子，还要把 \(a^{\mathrm{x}}\) 和 \(a^{\mathrm{y}}\) 都放等式一侧？除法，余数相关的元问题是啥？小学就学了，被除数 \(\div\) 除数＝商 \(\ldots\) 余数，等价于被除数 \(=\) 除数 \(\times\) 商 + 余数。所以设 \(a^{\mathrm{x}}\) 和 \(a^{\mathrm{y}}\) 除 \(p\) 的商分别是 \(m, n\) ，则 \(a^{x}=m p+b, a^{y}=n p+c\) 。这下可以乘了，\(a^{x+y}=a^{x} \cdot a^{y}=(m p+b)(n p+c)=m n p^{2}+(m c+b n) p+b c\) 。那么，这个式子能不能看出 \(a^{x}+y\) 除 \(p\) 余多少？这不显然吗？\(m n p^{2}+(m c+b n) p\) 除 \(p\) 都没了，剩个 \(b c\) 接着除，那就剩下 \(b \otimes c\) 呗！所以直接证完了一半。

还有剩下的一半：若 \(x+y \geq p-1\) ，那么现在需要考虑 \(a^{\mathrm{x}+\mathrm{y}-(\mathrm{p}-1)}\) 。但是这个时候，刚才证明的 \(a^{x}+y, \otimes=b \otimes c\)仍然成立。如果我们能证明 \(a^{x}+y-(p-1)\) 和 \(a^{x}+y\) 除 \(p\) 余数一样的话，那不就完事了吗？事实上，这个灵感我们稍微一细想，就能明白它的合理之处。我们还是继续遍历一下，如果 \(x+y\) 比较大，那么这时候它可以从 \(p-1\)遍历到 \(2(p-2)\) ，也就是接近 \(p-1\) 的两倍的位置。这已经基本把从 \(p-1\) 再往上数一轮 \(p-1\) 个数字的范围涵盖到了。所以这基本相当于证明：＂任取 \(k \in X, a^{k}\) 与 \(a^{k+(p-1)}\) 除 \(p\) 余数一样＂了。这是一个很一般的结果，而且刚才这一思考也能明白，这个结果肯定得是对的。而且关键的一点是，这个命题不依赖于 \(b, c\) 了，只是 \(a\) 的一个属

性，这意味着接下来我们可以把思维焦点放到这个命题上来。我们证明了它，就是证明了原命题。如何证明呢？注意 \(a^{k+(p-1)}=a^{k} \cdot a^{p-1}\) ，这又是一个乘积的形式，刚才我们刚从乘积的余数那里过来，所以故技重施：\(a^{k} \cdot a^{p-1}\)除 \(p\) 的余数是 \(a^{k, \otimes} \otimes a^{p-1, \otimes}\) 。那我们要证什么？不就是它等于 \(a^{k, \otimes}\) 吗？而 \(k\) 是在很大的范围内随便取的，所以要想证明乘一个 \(a^{p-1, \otimes}\) 不影响的话，只有一种可能：\(a^{p-1, \otimes}=1\) 。反过来，只要 \(a^{p-1, \otimes}=1\) ，刚才的逻辑就都通了。这是一个充要条件，信息量完全不流失，所以只要原题目正确，那么这个结论也一定正确，大家完全可以放心大胆去尝试证它。\\
\(a^{p-1, \otimes}=1\) 这个命题已经彻底和原题目没关系了，就是一个单纯的 \(a\) 的固有属性。我们再次把思维焦点放过来。如何求 \(a^{p-1}\) 除 \(p\) 的余数呢？既然和第二问没关系了，我们就及时抛弃目前没啥用的信息。唯一和 \(\alpha\) 这个东西有关的是啥来着？回去看看，唯一的信息就是＂\(a\) 很牛掰，\(a\) 的从 0 到 \(p-2\) 的全部幂次除 \(p\) 余数都不一样，恰好涵盖 \(1 \cdots p-1\)＂。而我要研究的是 \(a^{p-1}\) 的余数，正好是刚才的一轮幂次取值的天花板，往上挪了一个单位。而我要证明它余数是 1 ，也像是画了个轮回一样，回到了 \(a^{0}\) 。而它不偏不倚，恰好在余数遍历了一圈 \(1 \cdots p-1\)之后再次回到 \(a^{0}\) 的起点。这就很有感觉了。我们要证明它转了一圈必须回到起点，可以用反证法，假如它没回到起点，跑到别的地方去了，比如跑到了 2 ，还是因为元问题：指数对数相互控制，它就会落到 \(\log (p)_{a} 2\) 身上。用数学语言严格来写的话：反设 \(a^{p-1, \otimes} \neq 1\) ，记它为 \(t\) ，再记 \(s=\log (p)_{a} t \leq p-2\) ，于是 \(a^{s, \otimes}=t \neq 1\) 。那如何尝试导出矛盾？自然想到尝试从 \(s\) 次幂前进到 \(p-1\) 次幂。再记 \(a^{p-1-s, \otimes}=t^{\prime}\) ，还是之前乘法余数的故技重施，我们知道 \(t \otimes t^{\prime}=t\) ，于是 \(t t^{\prime}-t=t\left(t^{\prime}-1\right)\) 就是 \(p\) 的倍数。但这显然不可能，因为这时从 \(s\) 次幂前进到 \(p-1\) 次幂并非原地踏步，也就是 \(p-1-s>0 \Rightarrow t^{\prime} \neq 1\) ，但是这样的话 \(t^{\prime}-1\) 和 \(t\) 都没有 \(p\) 这个因数，乘一起也没有，这就导出矛盾了。于是第二问得证。

回顾第二问，我们主要采用的战术为分析法：顺着所证结论倒推问题。在做完基本的要素分析，把问题归结为研究 \(a^{\mathrm{x} \oplus \mathrm{y}}\) 的余数后，我们通过遍历操纵杆的取值，发现要点之一是计算 \(a^{\mathrm{x}+\mathrm{y}}\) 的余数。计算后我们惊喜地发现问题直接解决一半。在另一半的分析中，仍然是遍历法，我们认识到：要点是证明＂任取 \(k \in X, a^{k+(p-1)}\) 与 \(a^{k}\) 与除 \(p\) 余数一样＂，并利用 \(a^{k+(p-1)}\) 形式上的提示与刚才算乘法的经验认识到问题归结为证明 \(a^{p-1, \otimes}=1\) 。然后隔离问题，放好思维焦点，结合抛弃思想抓取有用信息，最后成功解决全部矛盾，拿下问题至于书写步骤的问题，刚才一直在倒推，而步骤需要按照综合法的方式书写，因此我们需要先证明引理 \(a^{p-1, \otimes}=1\) ，再证明 \(a^{k+(p-1)}\) 与 \(a^{k}\) 与除 \(p\) 余数一样，然后计算乘法，再分类讨论。这样一梳理，我们还需要把乘法余数的性质也单独证明一下，这样的话不管是 \(a^{k+(p-1)}\) 的余数，还是 \(a^{\mathrm{x}+\mathrm{y}}\) 的余数，都好说清楚了。这就是步骤的书写方法，留给大家作为练习。

也有同学可能会问，上面证明 \(a^{p-1, \otimes}=1\) 的想法有些巧妙，我没想到怎么办？事实上，大家可以发现，即便你没想出如何证明 \(a^{p-1, \otimes}=1\) ，你也可以通过上面的思维过程认识到这道题证明这个结论是一个必然，而且这个结论也一定正确。因此，如果你不会证的话，把这个结论摆上，然后把其他过程填上，其他过程也能拿分。最重要的，这个结论的正确性，直接关系到第三问的解决。

我们继续看第三问。仍然沿用我们之前的思维方式。现在是 \(a^{n . \otimes}=b\) 给定了，然后又新加了 \(x\) 和 \(k\) 作为新的操纵杆。这下出来四个操纵杆 \(a, n, x, k\) 了，然后它们又控制了 \(y_{1}, y_{2}\) 。现在求证的是 \(x=y_{2} \otimes y_{1}^{n(p-2), ~ \otimes}\) ，那自然，我们沿着要素关系网去表达清楚 \(y_{2} \otimes y_{1}^{n(p-2), \otimes}\) 即可。那该咋算咋算呗。先是 \(y_{1}^{n(p-2), \otimes}\) 。 \(y_{1}\) 是 \(a^{k}\) 的余数，现在要算 \(y_{1}^{n(p-2)}\) 的余数，我们就得先把 \(y_{1}\) 写成 \(a, k\) 控制的式子，再用 \(y_{1}\) 去控制 \(y_{1}^{n(p-2)}\) ，从而沿着关系网完成表示。仍然是刚才用的被除数，除数那一套（毕竟这就是控制余数的元问题），记 \(a^{k}\) 除 \(p\) 之商为 \(r\) ，则 \(a^{k}=p r+y_{1}\) ，于是 \(y_{1}^{n(p-2)}=\left(a^{k}-p r\right)^{n(p-2)}\) ，这玩意如何算除 \(p\) 的余数？二项式展开不就完了。展开第一项是 \(\left(a^{k}\right)^{n(p-2)}\) ，其余各项都有 \(p\) 的因子，一除都没了，剩下 \(a^{k n(p-2), \otimes}\) 然后 \(y_{2}=x \otimes b^{k, \otimes}\) 。那得先用上 \(b\) 的控制 \(a^{n, \otimes=b}\) ，然后算 \(b^{k, \otimes}\) ，这不

跟刚才的路子一样吗？刚才揭鼓 \(y_{1}\) 的时候是算的 \(a^{k, \otimes}\) 再一个幂次取余数，最后是两个幂次乘起来，现在也是同理可得 \(b^{k, \otimes}=a^{k n, \otimes}\) ，于是 \(y_{2}\) 就是 \(x \cdot a^{k n, \otimes}\) 的余数。下一步控制：\(y_{2} \otimes y_{1}^{n(p-2), \otimes}\) 的两个东西都处理好了，得再用一次被除数表示：设 \(x \cdot a^{k n, \otimes}\) 除 \(p\) 之商为 \(g\) ，故 \(y_{2}=x \cdot a^{k n, \otimes}-g p\) 乘起来就是 \(\left(x \cdot a^{k n, \otimes}-g p\right) \cdot a^{k n(p-2), \otimes}\) ，它的除 \(p\) 余数就是最终答案。后面 \(g p \cdot a^{k n(p-2), \otimes}\) 除没了，剩下 \(x \cdot a^{k n, \otimes} \cdot a^{k n(p-2), \otimes}\) ，根据第二问的指数计算方法，不难证明它与 \(x \cdot a^{k n(p-1), ~} \otimes\) 除 \(p\) 的余数相同。

算出来之后就看所求了，最后所要证明的就是 \(x \cdot a^{k n(p-1), ~} \otimes\) 除 \(p\) 余数是 \(x\) 。这个东西被操纵杆 \(k, x, n\) 控制，但最后只和 \(x\) 有关，相信大家经过前面的过程，到这里都能一眼看出来了，这就等价于证明 \(a^{k n(p-1), ~ \otimes}=1\) 。里面又出现了 \(p-1\) ，正好就是第二问里面的那个 \(a^{p-1, \otimes}=1\) 的延伸，就是给 \(a^{p}-1\) 再套个任意的 \(k n\) 次方，证明它的余数还是1。刚才的过程实际上已经蕴含了它的证明（就是刚才处理 \(y_{1}\) 的方法），或者大家自行补充完整即可。

这道题对于没有竞赛数论基础的同学而言，难度绝对是称得上压轴题难度的。运算量不大，但是思维量很大，非常综合地运用到了要素关系网，动静关系，遍历法与充要替换隔离问题与思维焦点等等在up主＂高中数学经验分享系列＂中重点讲解的多种思维方式，并且使用上有一定的思维技巧性，例如颠倒要素关系网的处理思路。第二，三问由于是证明题，利用分析法进行逆推也是十分有效的思维方式，并且配合遍历的思维可以帮助我们抓出下一步应该关注的主要矛盾，让我们的思维焦点放对位置。

这道题虽然难度和创新度都很不错，但作为高考题目而言却并不是一道好题。原因既不是创新创死人，也没有出什么超纲知识，而是严重丧失了公平性。对于稍微接触过一点竞赛数论的同学而言，对带余除法，同余的基本性质稍加熟悉，就能很轻易地想清楚问题的本质，而从未接触过的同学则需要深刻掌握科学的思维方式，再配合一点灵感（主要灵感需求点来自 \(a^{p-1, \otimes}=1\) 的证明）才能完整解出题目。由于学习竞赛或者强基计划的学生比重并不小，因此有这样的知识优势的学生比重也就不小，这就造成了两个大比重学生群体面对此题的体感难度的不均衡。Up 主个人认为，想要创新，引人大学知识作为背景，是好事，但是不能违背两个更加基本的原则：一是要保证题目求解过程不必须使用超纲的知识；二是要保证超前学过这个知识的学生，相比于不知道这个知识的学生，不能获得太大的优势。第一点考验的是出题人对考纲的把握程度，而第二点中提到的优势无论如何都客观存在，因此只能通过人题知识选取，题干知识讲解，题干与各问号间的提示设置等方面尽量缩小二者的差距，这对出题人是一种考验。单看第二点来说，这道题的命制无疑是失败的。相比而言，去年四省联考的椭圆曲线压轴题就明显比此题更高明。Up 主对本题进行了少量的修改，改后的题目放在本文档的最后一页，感兴趣的同学可以参考。

\section*{（up 主修改后的题目）}
离散对数在密码学中有重要的应用。设 \(p \geq 2\) 是正整数，定义集合 \(X_{p}=\{0,1,2 \cdots p-1\}\) ，集合 \(X_{p}^{+}=\) \(\{1,2 \cdots p-1\}\) 。对 \(m \in N\) ，记 \([m]_{p} \in X_{p}\) 为 \(m\) 除以 \(p\) 的余数。\\
（1）给出下面两个命题：\\
\(A\) ：对 \(\forall m, n \in N\) ，有 \([m]_{p}[n]_{p}=[m n]_{p}\)\\
\(B:\) 对 \(\forall m, n \in N\) ，有 \(\left[{ }_{[m]_{p}}[n]_{p}\right]_{p}=[m n]_{p}\)\\
试判断命题 \(A\) 与 \(B\) 的真假性，若为真命题请给予证明，若为假命题请给出反例\\
（2）下设 \(p\) 是素数。 \(a \in X_{p}^{+}\)，满足 \(1,[a]_{p},\left[a^{2}\right]_{p} \cdots\left[a^{p-2}\right]_{p} \in X_{p}^{+}\)且两两不同。若 \(b \in X_{p}^{+}\)，且满足 \(\left[a^{n}\right]_{p}=b\)（其中 \(n \in\{0,1 \cdots p-2\}\) ），则定义 \(n\) 是 \(b\) 的以 \(a\) 为底的离散对数，并记为 \(n=\log (p)_{a} b\) 。\\
（i）证明：对 \(b, c \in X\) ，有 \(\log (p)_{a}\left([b c]_{p}\right)=\left[\log (p)_{a} b+\log (p)_{a} c\right]_{p-1}\)\\
（ii）已知 \(n=\log (p)_{a} b\) ，对 \(x \in X_{p}^{+}, k \in\{1,2 \cdots p-2\}\) ，令 \(y_{1}=\left[a^{k}\right]_{p}, y_{2}=\left[x \cdot\left[b^{k}\right]_{p}\right]_{p}\) ，证明：\(x=\left[y_{2} \cdot\left[y_{1}^{n(p-2)}\right]_{p}\right]_{p}\)

\section*{6．4．2 冻结理论＋思想实验}
多重最值方向何在？冻结理念配合思想实验寻找思维突破口（视频日期：2024年3月10日）题目：

定义 \(\max (x, y, z)\) 为 \(x, y, z\) 中最大的数。若 \(a, b, c>0\) ，则 \(\max \left(\frac{1}{\mathrm{ac}}+\mathrm{b}, \frac{1}{\mathrm{a}}+\mathrm{bc}, \frac{\mathrm{a}}{\mathrm{b}}+\mathrm{c}\right)\) 的最小值为 \(\qquad\) －

这道题目模仿今年九省联考的填空压轴题，其式子的形式更为复杂，严谨处理得出答案的过程也更加困难。如果想严谨地论证本题的答案，其难度将会很高。但是，这种多自由度问题的核心思维方式，是始终不变的，那就是冻结法（最早讲解于 2021 年 9 月 16 日视频）。所谓冻结法，就是对于实际自由度大于 1 的问题，选择先冻结某些量，降低问题的自由度数量，然后在该状态之下求出此时的最值（由于冻结的量此时作为背景，最值一定是受到它们控制的），这样就让自由度减少了 1。重复循环下去，就能消去全部自由度，得到真正的最值。不过，在较为困难的题目上面，冻结法的使用也会更加灵活。

这道题直接明摆着给了你三个主自由度 \(a, b, c>0\) ，因此我们肯定需要先冻结两个变量，并调试另一个变量使原式达到最小值，然后如此循环。那么，先冻结谁呢？这个我们暂时看不出来。但是，up 主说过：两三个自由度的时候，反正冻结的选项也不多，大不了挨种方案试试，总能做出来。本题的初始冻结手段显然有三种：冻 \(\mathrm{a}, \mathrm{b}\) ，冻 b，c，冻 a，c。我们可以选一种来尝试一下。

比如，来试试冻结 \(\mathrm{a}, \mathrm{b}\) 的手段。假如现在 \(\mathrm{a}, \mathrm{b}\) 均静止了，那我们就要调试（遍历） c ，使得所求

\[
\max \left(\frac{1}{\mathrm{ac}}+\mathrm{b}, \frac{1}{\mathrm{a}}+\mathrm{bc}, \frac{\mathrm{a}}{\mathrm{~b}}+\mathrm{c}\right)
\]

达到最小。先使用遍历法，遍历一下 c 找找感觉：如果 C 很小（比如 \(\frac{1}{10000}\) ），那么显然 \(\frac{1}{\mathrm{ac}}+\mathrm{b}\) 会很大，再取了最大的那个，所求也会变得爆炸性的大。如果 c 很大，那么同理，后两个

\[
\frac{1}{\mathrm{a}}+\mathrm{bc}, \frac{\mathrm{a}}{\mathrm{~b}}+\mathrm{c}
\]

也会一起变得很大，更不用说最后到所求了。因此直观上两头高，中间肯定有个最小的＂凹槽＂。他在哪呢？随着 c从很小增长到很大，这三个数中，一个从正无穷下降，两个增长到正无穷，然后我取这三个柱子的最大值。那么什么时候这个最大值最小呢？同样遍历一下，可以看到，肯定是下降的柱子与某个长高的柱子刚好平齐时，所谓的＂最大值＂达到最小。因此这个时候必是 \(\frac{1}{\mathrm{ac}}+\mathrm{b}\) 与另外俩的某一个一样。我们来对照一下。如果前两个相等，那么 \(\frac{1}{\mathrm{ac}}+\mathrm{b}=\frac{1}{\mathrm{a}}+\mathrm{bc}\) ，但是一旦我们的思维焦点随之过来，我们就不难注意到

\[
\frac{1}{\mathrm{ac}}+\mathrm{b}=\frac{1}{\mathrm{c}}\left(\frac{1}{\mathrm{a}}+\mathrm{bc}\right),
\]

因此前两个相等就等价于 \(\mathrm{c}=1\) ！！这也就是说，如果 \(0<\mathrm{c} \leq 1\) ，那么 \(\frac{1}{\mathrm{ac}}+\mathrm{b} \geq \frac{1}{\mathrm{a}}+\mathrm{bc}\) ，所求其实就可以等价为

\[
\max \left(\frac{1}{\mathrm{ac}}+\mathrm{b}, \frac{\mathrm{a}}{\mathrm{~b}}+\mathrm{c}\right)
\]

反之如果 \(\mathrm{c} \geq 1\) ，那所求就等价于

\[
\max \left(\frac{1}{\mathrm{a}}+\mathrm{bc}, \frac{\mathrm{a}}{\mathrm{~b}}+\mathrm{c}\right)
\]

这时，如果 \(\mathrm{c} \geq 1\) ，显然当 c 增大时，这俩数都在涨，因此要使最终结果最小，只能 \(\mathrm{C}=1\) ，这就可以归并到第一种情

况里去了。因此，我们只需要考虑 \(0<\mathrm{c} \leq 1\) 的 \(\max \left(\frac{1}{\mathrm{ac}}+\mathrm{b}, \frac{\mathrm{a}}{\mathrm{b}}+\mathrm{c}\right)\) ．

接下来很自然地就想要去考虑到底取什么 C 才能让上式最小了。类似地，仍然遍历 c ，我们会发现，仍然是两个柱子一个增长一个下降，取最值时应当刚好相等才对。但是，也有可能下降的还没来得及降到上升柱子的高度， c 就到 1 了，这时候最值应当取 \(\frac{1}{\mathrm{ac}}+\mathrm{b}\) 。但是，这时候都有一个问题。比如第一种情况，我知道应该让 \(\frac{1}{\mathrm{ac}}+\mathrm{b}=\frac{\mathrm{a}}{\mathrm{b}}+\mathrm{c}\)以实现 \(\mathrm{a}, \mathrm{b}\) 控制 C 的效果，但是显然这个控制式需要求解出 C 来，表达式会十分复杂，难以处理。更何况现在还需要分类讨论，事情会变得越来越棘手。

那怎么办呢？up 主说过，遇到抓不住方向的时候，首先可以做一做思想试验（最早讲解于 2023 年 3 月 5 日视频）。我们可以取点具体数字来看看情况。比如 \(a=1, b=1\) ，这时就是要看

\[
\max \left(\frac{1}{\mathrm{C}}+1,1+\mathrm{C}\right)
\]

但这太明显了，很显然 \(\mathrm{C}=1\) 时二者涨落到齐平，正好是 2 ，这个例子太简单。那如果 \(\mathrm{a}=2, \mathrm{~b}=1\) 呢？求

\[
\max \left(\frac{1}{2 \mathrm{C}}+1,2+\mathrm{c}\right) .
\]

虽然不知道是多少，但好像不可能刚才的 2 小了。咋回事？仔细一看，是 \(\mathrm{a}, \mathrm{b}\) 取的不好导致的。那我换一个， \(\mathrm{a}=2\) ， \(\mathrm{b}=\frac{3}{2}\) 呢？那就是

\[
\max \left(\frac{1}{2 \mathrm{C}}+\frac{3}{2}, \frac{4}{3}+\mathrm{c}\right)
\]

这时候就不太想计算，但是，最值一样可以去做思想试验啊。比如能不能小过刚才的 2 呢？那就是让上面两个都小于 2 ，看看有没有符合的 c 呗！于是尝试求解不等式组

\[
\frac{1}{2 \mathrm{C}}+\frac{3}{2} \leq 2, \frac{4}{3}+\mathrm{C} \leq 2
\]

解出来 \(1 \leq \mathrm{c} \leq \frac{2}{3}\) ，这怎么可能？\\
这么搞一搞思想试验，我就开始怀疑，是不是 2 就是最小的数呢？能不能证明，或者构造出更小的呢？如果要构造更小的话，那就得

\[
\frac{1}{\mathrm{ac}}+\mathrm{b} \leq 2, \frac{\mathrm{a}}{\mathrm{~b}}+\mathrm{c} \leq 2
\]

刚才做的操作就是在冻结 \(\mathrm{a}, \mathrm{b}\) 的前提下，去找 \(\mathrm{c} \in(0,1]\) 满足上述不等式。那我如法炮制，反正 \(\mathrm{a}, \mathrm{b}\) 视为静止，那我求解一下上述不等式，得到

\[
\frac{1}{a(2-b)} \leq c \leq 2-\frac{a}{b} \text { 且 } b<2 \text {. }
\]

这个不等式需要有解，且与 \((0,1]\) 有交集才行。因此，首先需要满足 \(\frac{1}{a(2-b)} \leq 2-\frac{a}{b}\) ，自然要做一下充要替换，得

\[
\frac{\mathrm{a}}{\mathrm{~b}}+\frac{1}{\mathrm{a}(2-\mathrm{b})} \leq 2
\]

什么样的 \(\mathrm{a}, \mathrm{b}\) 满足上述条件呢？仍然采用冻结法的思维方式，我先冻结住 \(b\) ，这时只需要关于 \(a\) 的函数 \(f(a)=\) \(\frac{a}{b}+\frac{1}{a(2-b)}\) 最小值不超过 2 ，我就能构造出对应的 \(a\) 来，而 \(f(a)\) 的最小值是 \(\frac{2}{\sqrt{b(2-b)}}\) ，因此再想找 \(b\) 就只需要

\[
\frac{2}{\sqrt{b(2-b)}} \leq 2
\]

整理得到 \((b-1)^{2} \leq 0\) ，显然只有 \(b=1\) 满足条件，那自然逆推回去，也只有 \(\mathrm{a}=1\) 了，这时就回到刚才最简单的情形，最小值是 2 ，等号可以成立，当且仅当 \(C=1\) 。所以使得所求值小于等于 2 的 \(a, b, c\) 只有这一组。那自然地，最

后的最小值就只能是 2 了，且取等条件就是最基本的 \(\mathrm{a}=\mathrm{b}=\mathrm{c}=1\) ，问题就这么被解决了。

回顾刚才的解题过程。首先，冻结法是一个基本的处理手段，从头贯穿到尾。但是在中间过程中我们遇到了挫折。遇到挫折怎么办？常见的打开新局面的手段就是我们之前讲到的思想试验，通过一些特殊例子寻找题目的感觉，从多个特殊例子中寻找到一般规律，随后大胆假设，小心求证，将刚才的直观感觉转化为严格的论证语言，从而打开新局面，突破难题。这道题单纯猜答案的话，取 \(\mathrm{a}=\mathrm{b}=\mathrm{c}=1\) 不难猜测 2 是正确答案，但是严格论证的话，通过思想试验的思考方式，也能沉稳找到突破口，一步步走向最后的结果。思想试验是带有灵感色彩的思考方式，面对一些思维上的难题常常有奇效，希望有志于突破 140 乃至 150 的同学能够结合这些视频，多多体会这种思考方式，让它配合基本的动静与控制的思维方式，协助自己取得更高的突破。

\section*{6．4．3 一道怪题的纯思维详解}
题目：\\
法国数学家韦达发现了一元二次方程根与系数的关系，并将其推广到高次方程。人们把这样的关系称为韦达定理。对实系数一元 \(n\) 次方程 \(a_{n} x^{n}+a_{n-1} x^{n-1} \cdots+a_{1} x+a_{0}=0\)（其中 \(a_{n} \neq 0\) ），若 \(x_{1} \cdots x_{n}\) 是它在复数集 \(\mathbb{C}\)内的 \(n\) 个根，则

\[
\left\{\begin{array}{l}
x_{1}+x_{2}+\cdots+x_{n}=-\frac{a_{n-1}}{a_{n}} \\
x_{1} x_{2}+x_{1} x_{3}+\cdots+x_{n-1} x_{n}=\frac{a_{n-2}}{a_{n}} \\
\cdots \\
x_{1} x_{2} \cdots x_{n}=(-1)^{n} \frac{a_{0}}{a_{n}}
\end{array} .\right.
\]

请用韦达定理解决下题：设 \(a, b \in \mathbb{R}\) ，关于 \(x\) 的方程 \(x^{3}+(2-a) x^{2}+b x-a=0(a>0)\) 有三个实数根，且至少有一个实数根落在区间 \((0, a)\) 内，则 \(2 a-b\) 的最大值为 \(\qquad\)。

这道题目比较有趣，难度也比较大。题目给了我们高次方程的韦达定理，并让我们拿去解决新问题。面对一个新知识，按照up 主在高中数学经验分享第五期的理念，高中阶段的知识点都要去刻画某种要素结构。在元问题的层面上尝试理解知识点，才能在解题中发挥出稳定性。这里出现的知识我们实际上很熟悉：我们学习过二次方程的韦达定理，而本题是把它推广到更高的次数。我们回忆一下二次方程韦达定理对应的元问题：一元二次方程 \(a_{2} x^{2}+a_{1} x+a_{0}=0\) 的两根 \(x_{1}, x_{2}\) 满足

\[
x_{1}+x_{2}=-\frac{a_{1}}{a_{2}}, x_{1} x_{2}=\frac{a_{0}}{a_{2}} .
\]

忽略大小无关后，控制一个二次方程需要两个自由度，也就是说两个自由度控制出一个确定的方程。进而可以求解方程得到确定的（复数）解 \(x_{1}, x_{2}\) 。因此是两个自由度控制 \(x_{1}, x_{2}\) 。而这种控制除了直接用方程实现之外，还可以用韦达定理的方式实现。给定两个大小无关下的自由度 \(\frac{a_{1}}{a_{2}}, \frac{a_{0}}{a_{2}}\) 之后，韦达定理可以直接作为方程组控制式来控制 \(x_{1}, x_{2}\) 。因此，韦达定理可以作为另一种控制模式。另外，up 主反复强调：颠倒要素网是非常重要的。在分析元问题对应的要素结构时，势必要关心它能怎样参与构建要素网。都是两个自由度，\(x_{1}, x_{2}\) 也可以作为主自由度，这样的话，它们就能直接控制方程本体。在2023年7月26日的视频中，我们直接了当地讲了韦达定理的这种元问题理解。

现在，我们来看高次方程的韦达定理。可以看到，与二次方程一样，一个 \(n\) 次方程去掉大小无关，一共有 \(n\) 个自由度。现在正好 \(x_{1} \cdots x_{n}\) 是 \(n\) 个操纵杆（主自由度），它们直接把方程里的系数全部合成出来了。这和我们熟

悉的二次方程情形一模一样，甚至题目的式子如此直截了当地让左边全是操纵杆在控制，本身也是对你直接地提示，提示你韦达定理这个元问题就是一种充要替换，一种要素网改写。

题目点明了让你使用韦达定理解题，就是直接告诉你：这道题用韦达定理来改写要素关系网，准没错！下面我们来读题：显然本题是两个自由度的问题。题目自带的操纵杆是 \(a, b\) ，并且剩下的全部是不等关系的限制，不会再减少自由度。直接思考确实感觉摸不着头脑，但是已知条件全部都是围绕三个根 \(x_{1}, x_{2}, x_{3}\) 的，再加上题目要求你用韦达，以及刚才对韦达的元问题分析，这路标已经用在你脸上了。

先不废话，把韦达式子写出来：

\[
x_{1}+x_{2}+x_{3}=a-2, x_{1} x_{2}+x_{2} x_{3}+x_{3} x_{1}=b, x_{1} x_{2} x_{3}=a
\]

如果我们尝试就以 \(x_{1}, x_{2}, x_{3}\) 为主自由度，那么已知要求至少一个根在 \((0, a)\) 内。由对称性，不妨就设 \(0<x_{1}<a\)。这样的话这一条已知的信息量就完全表达了。三个根都是实数没啥好说的。这样的话，三个操纵杆，整体两个自由度，我们还缺一个牵制。显然可以消去 \(a\) 得到牵制 \(x_{1}+x_{2}+x_{3}=x_{1} x_{2} x_{3}-2\) ，定义域也就等价于 \(0<x_{1}<\) \(x_{1} x_{2} x_{3} \Leftrightarrow x_{1}>0, x_{2} x_{3}>1\) 。这样，已知条件的信息量就完全完整了。

直接这样往下做当然是可行的（课后练习会引导大家继续沿着上面选 \(x_{1}, x_{2}, x_{3}\) 为主自由度的方法去做）。这里up主为了强调重构要素网的重要性，采用另一种方法为大家讲解。上面的这一大长串式子看着还是有点发毛。但是定义域 \(0<x_{1}<a\) 看上去却很简洁。本题自由度是二，能不能就直接选 \(x_{1}, a\) 作为主自由度呢？这样的话，我们就不需要烦人的牵制了。此时重构本题的要素网：给定 \(0<x_{1}<a\) 的 \(x_{1}, a\) ，首先利用 \(x_{2}+x_{3}=a-2-x_{1}\) ，以及 \(x_{2} x_{3}=\frac{a}{x_{1}}\) 控制出 \(x_{2}, x_{3}\) ，再直接控制 \(b\) 。既然如此，那就要保证 \(\left\{\begin{array}{l}x_{2}+x_{3}=a-2-x_{1} \\ x_{2} x_{3}=\frac{a}{x_{1}}\end{array}\right.\) 有实数解。up 主讲过，这就等价于

\[
\left(a-2-x_{1}\right)^{2} \geq \frac{4 a}{x_{1}}
\]

这个式子需要添加进定义域限制中。进而，我们用 \(x_{1}, a\) 来控制所求 \(2 a-b=2 a-\left(x_{1} x_{2}+x_{2} x_{3}+x_{3} x_{1}\right)\) 。注意 \(x_{2}, x_{3}\)都可以被 \(x_{1}, a\) 控制，而且因为都是对称式，正好可以直接带换：

\[
x_{1} x_{2}+x_{2} x_{3}+x_{3} x_{1}=x_{1}\left(x_{2}+x_{3}\right)+x_{2} x_{3}
\]

代人前面的式子得到 \(x_{1}\left(a-2-x_{1}\right)+\frac{a}{x_{1}}\) 。因此 \(2 a-b=2 a-x_{1}\left(a-2-x_{1}\right)-\frac{a}{x_{1}}\) 。\\
控制过程走完了，我们就该手握操纵杆去求最值了。首先得把刚才多出来的定义域不等式 \(\left(a-2-x_{1}\right)^{2} \geq \frac{4 a}{x_{1}}\)处理一下。光看 \(x_{1}\) 的话感觉式子结构很不舒服，但是这显然只是关于 \(a\) 的二次不等式。而且，注意到 \(a=x_{1}\) 时，这个不等式刚好等号成立！这个观察具有一定难度，主要是由于左边刚好减，右边刚好除，两边一起提示代人 \(a=x_{1}\)试试。不过，如果你具备冻结的思考方式，并且主动观察意识到先看 \(a\) 比较好了，那你就自然会关注不等式成立的二次方程的解的位置，也就会被刚才的提示引导着代入 \(a=x_{1}\) 尝试。这个观察表面上是灵感占主导，实则还是思维占主导。但不管如何，既然看出二次方程有一个解了，那就可以直接把不等式求解出来了。展开得到

\[
a^{2}-2\left(2+x_{1}\right) a-\frac{4 a}{x_{1}}+\left(2+x_{1}\right)^{2} \geq 0
\]

再由韦达定理得到它的另一个根是

\[
\frac{\left(2+x_{1}\right)^{2}}{x_{1}}=x_{1}+\frac{4}{x_{1}}+4
\]

因此，上述不等式等价于：

\[
\left.a \leq x_{1} \text { 或 } a \geq x_{1}+\frac{4}{x_{1}}+4 \text { (注意 } 0<x_{1}<a \text { 的限制 }\right)
\]

舍去一侧，得到定义域完整限制为 \(x_{1} \geq 0\) 且 \(a \geq x_{1}+\frac{4}{x_{1}}+4\)\\
定义域有了，接下来看所求 \(2 a-x_{1}\left(a-2-x_{1}\right)-\frac{a}{x_{1}}\) 。作为两个自由度的问题，使用冻结法不必多说（已经被 up 主讲烂的路子）。先冻结谁？所求关于 \(x_{1}\) 不明不白，但它显然是关于 \(a\) 的一次函数

\[
\left(2-x_{1}-\frac{1}{x_{1}}\right) a+x_{1}\left(2+x_{1}\right) .
\]

一次项系数 \(2-x_{1}-\frac{1}{x_{1}} \leq 0\) 显然，因此 \(a\) 越小，所求越大（或者持平）。因此冻结 \(x_{1}>0\) ，直接把 \(a\) 拉到最小值 \(x_{1}+\frac{4}{x_{1}}+4\) ，带回去得

\[
\left(2-x_{1}-\frac{1}{x_{1}}\right)\left(x_{1}+\frac{4}{x_{1}}+4\right)+x_{1}\left(2+x_{1}\right)=3+\frac{4}{x_{1}}-\frac{4}{x_{1}^{2}}
\]

接下来只需对 \(x_{1}>0\) 求 \(3+\frac{4}{x_{1}}-\frac{4}{x_{1}^{2}}\) 的最大值，已经很简单了，最后求出来就是 \(x_{1}=2\) 时取得最大值 4 ，这时有 \(a=8\) 。

最后做即时检查。我们的原则是带回原题验证，以破坏原有的计算路径。刚才绕了这么多得到的最大值 4 ，只要把 \(x_{1}=2\) 和 \(a=8\) 代人原题背景验证此时 \(2 a-b=4\) 吻合，我们就可以相信刚才绕的这一大圈没有错误。此时方程 \(x^{3}-6 x^{2}+b x-8=0\) 有 \(x_{1}=2\) 这个根。代人根得到 \(b=12\) ，因此此时就有 \(2 a-b=16-12=4\) 。检查成立。

做个总结。这道题整体难度较大，但是思路非常传统。始终坚守 up 主讲的路径：分析宏观的动静关系，预判与建构要素关系网，调用元问题解题。整道题目一直就在这样的路上走，没有一刻偏离出去。唯一的区别就是这道题是明晃晃给了你一个元问题让你使用，然后计算量大了些，达到了难题的标准。只要大家把基本的思维方式吃透并加以正确应用，这道题就难不倒你。

\section*{课后练习：}
（1）仿照视频中方法的大致过程，使用选 \(X_{1}, x_{2}, x_{3}\) 作为主自由度的路径重新求解本题。提示：你可以仿照 2021年10月3日的视频，试着把 \(x_{2} X_{3}\) 整体作为一个新主自由度。\\
（2）韦达定理的补充：我们需要一个前置知识，即复分析中著名的代数基本定理：对于任何不小于 1 次的复系数多项式 \(f(x)\) ，它在复数范围内都至少有一个零点。\\
（i）对于任意一个复系数 \(n \geq 1\) 次多项式 \(f(x)\) ，以及任意复数 \(a\) ，请你利用因式分解的思想，直接证明 \(f(x)-f(a)\)有因式 \(x-a\) 。据此说明：若 \(a\) 是 \(f(x)\) 的一个零点，则 \(f(x)\) 可以因式分解为 \((x-a) \cdot g(x)\) ，其中 \(g(x)\) 是一个 \(n-1\)次多项式。\\
（ii）用第（i）问证明：存在复数 \(x_{1} \cdots x_{n}\) ，使多项式 \(f(x)=a_{n} x^{n}+a_{n-1} x^{n-1} \cdots+a_{1} x+a_{0}\) 可以因式分解成

\[
a_{n}\left(x-x_{1}\right)\left(x-x_{2}\right) \cdots\left(x-x_{n}\right) .
\]

（iii）你能试着利用第（ii）问的结果，直接证明韦达定理吗？

\section*{6．4．4 XXXX}
\section*{6.5 圆雉曲线}
\section*{6．5．1 \(100 \%\) 稳收武汉二调 16 题}
\section*{题目（武汉二调填空压轴）：}
设 \(F\) 为双曲线 \(E: \frac{x^{2}}{a^{2}}-\frac{y^{2}}{b^{2}}=1(a>0, b>0)\) 的右焦点，\(A, B\) 分别为双曲线 \(E\) 的左右顶点点 \(P\) 为双曲线 \(E\)上异于 \(A, B\) 的动点，直线 \(l: x=t\) 使得过 \(F\) 作直线 \(A P\) 的垂线交直线 \(l\) 于点 \(Q\) 时总有 \(B, P, Q\) 三点共线，则 \(\frac{I}{a}\) 的最大值为 \(\qquad\) ．

边读题边画图，注意画图方式带来的要素关系网生成。双曲线本身还带着俩自由度。这没办法，只能先随手画一下。画完双曲线那么 \(F, A, B\) 都就跟着被控制了。然后新增一个自由度的动点 \(P\) 。继续读题，直线 \(l\) 似乎也是个带未知 \(t\) 自由度的东西，先随便画一条 \(l\) 试试，但后面还有一堆描述，先继续看：过 \(F\) 作垂线，这个被刚才的要素网也控制了，然后又和 \(l\) 生出 \(Q\) ，这时候牵制来了：三点共线！你如果刚才如果随便画一条 \(l\) 的话，这理会明显感觉到别扭：这看着不像三点共线啊！但我这个 \(Q\) ，从画图的角度来看，又是应该被 \(P\) 控制住了的。所以我想：这怎么能画的尽量准确一点呢？那非常容易发现，先把 \(Q\) 画出来不就完了？那我就先不画 \(l\) 了，先过 \(F\) 作 \(A P\) 的垂线，再和 \(B P\) 交于 \(Q\) ，你要搞三点共线那这个 \(l\) 就得是过 \(Q\) 且垂直于 \(x\) 轴的直线！这样一来，就画准确了。

但是，从句子的描述上来说，这么安逸地把 \(l\) 画出来似乎有点不太对劲。刚才说了，双曲线先自带俩自由度，这俩自由度的基础上，\(P\) 是个单自由度动点，它严格控制了后面一长串，后面出现的那些要素全都是在这个单自由度机械臂下控制运行的。但是题干原本的描述是：直线 \(l: x=t\) 使得不管 \(P\) 如何，都恒有三点共线，说明直线 \(l\)应该是这个双曲线的固有属性，不受 \(P\) 的控制，但是刚才画图法给我们的要素关系网显示：这个 \(l\) 是被 \(P\) 控制的啊！一分析到这里，马上就明白了，这正是我们在＂数学的干货＂的视频列表里重点强调的不变量问题，大家可以回去复习一下相关的视频。（先把双曲线看成静止背景）\(P\) 的自由度控制了 \(l\) ，或者说控制了 \(t, t\) 可以表示成 \(P\) 中的某个参量（主自由度）的函数。但这个函数碰巧是个常数函数，所以呈现出来的就是，\(t\) 是一个不变量，不被 \(P\) 的变化影响，而是受要索网史菲丽的因索（也就是双曲线的背景）的控制（当然，这个控制也有可能出来常数函数）。但是我们需要做的，和平时的问题没有什么差别，仍然是沿着要素网计算，只不过需要说明这确实是个常数函数才行。

那我们就沿着要素网去计算吧！首先设 \(P\) 的坐标是 \(\left(x_{0}, y_{0}\right)\) ，其满足牵制 \(\frac{x_{0}^{2}}{a^{2}}-\frac{y_{0}^{2}}{b^{2}}=1\) 这个式子承载了 \(P\) 的全部信息量，正好 \(x_{0}, y_{0}\) 两个参数加一个牵制，说明 \(P\) 的自由度是 1 。然后沿着要素网生成顺序去表示，\(B P\) 方程是 \(y=\frac{y_{0}}{x_{0}-a}(x-a), A P\) 斜率是 \(\frac{y_{0}}{x_{0}+a}\) ，因此垂线 \(F Q\) 方程为 \(y=-\frac{x_{0}+a}{y_{0}}(x-c)\) 。我们要求的 \(t\) 是交点的横坐标，所以直接联立消去 \(y\) ，即 \(\frac{y_{0}}{x_{0}-a}(x-a)=-\frac{x_{0}+a}{y_{0}}(x-c)\) ，解得 \(t=\frac{a y_{0}^{2}+c x_{0}^{2}-a^{2} c}{x_{0}^{2}+y_{0}^{2}-a^{2}}\) ．

刚才说过，全部的信息量都在 \(\frac{x_{0}^{2}}{a^{2}}-\frac{y_{0}^{2}}{b^{2}}=1\) 上，所以就是要用这个去证明 \(t\) 是定值。结合之前讲过的不变量问题统一思想，需要把表示出来的东西进行消元，一直消到式子中剩下的参数个数不超过实际自由度。这里 \(P\) 本身只有一个自由度，其余的都是双曲线自带的这些都先留着），所以只需要消到剩 \(x_{0}, y_{0}\) 中的一个参数即可，

比如把 \(y_{0}^{2}=\frac{b^{2} x_{0}^{2}}{a^{2}}-b^{2}\) 代人，整理得到

\[
t=\frac{a\left(\frac{b^{2} x_{0}^{2}}{a^{2}}-b^{2}\right)+c x_{0}^{2}-a^{2} c}{x_{0}^{2}-b^{2}+\frac{b^{2} x_{0}^{2}}{a^{2}}-a^{2}}=\frac{\left(b^{2}+a c\right)\left(\frac{x_{0}^{2}}{a}-a\right)}{\left(b^{2}+a^{2}\right)\left(\frac{x_{0}^{2}}{a^{2}}-1\right)}=\frac{a\left(b^{2}+a c\right)}{b^{2}+a^{2}}
\]

果然是定值。\\
现在，\(t\) 里已经全是双曲线背景相关的东西了，我们要求的就是遍历所有双曲线，看 \(\frac{t}{a}\) 何时最大。显然 \(\frac{t}{a}=\) \(\frac{b^{2}+a c}{b^{2}+a^{2}}\) ，这是个大小无关的东西，只与双曲线形状有关，被单自己的由度的形状参量控制。比如我们用 \(e\) 去表示它：

\[
\frac{t}{a}=\frac{c^{2}-a^{2}+a c}{c^{2}}=\frac{e^{2}+e-1}{e^{2}}=1+u-u^{2}
\]

这里有换元 \(u=\frac{1}{e} \in(0,1)\) ，那最大值显然就是 \(u=\frac{1}{2}\) 时取 \(\frac{5}{4}\) 。

做个总结。这道题里的关键步骤就是不带着 \(t\) 去算，而反过来去算 \(t\) ，若你只看答案解析，则不明白为什么要这样做。实际上，从我们最开始的分析来看，其实就是一个画图法寻找最舒服的要素关系网的过程。重构了要索网之后，很容易通过分析自由度得到这是一个不变量问题，不变量正是 \(t\) ，然后要素网本身结构很简单，计算过程很容易沿着要素网看清楚，直接上手计算就立刻把题目斩杀了。全程一直都是我们一直以来强调过的思维方式在起作用，思路非常自然，没有任何一步存在灵感要求，因此，对于平时学习我讲的思维方式较多的同学，这道题应该在考场上不失分。

\[
f(x) \sim \frac{a_{0}}{2}+\sum_{n=1}^{\infty}\left(a_{n} \cos \frac{n \pi}{L} x+b_{n} \sin \frac{n \pi}{L} x\right)
\]

考场时间规划：\\
读题＋画图法重构要素网： 1.5 min\\
预判计算路径＋选主自由度： 0.5 min\\
计算 \(t: 2 \mathrm{~min}\) 求最值： 1 min\\
合计：5min

\section*{课后思考：}
在我们计算出

\[
t=\frac{a\left(\frac{b^{2} x_{0}^{2}}{a^{2}}-b^{2}\right)+c x_{0}^{2}-a^{2} c}{x_{0}^{2}-b^{2}+\frac{b^{2} x_{0}^{2}}{a^{2}}-a^{2}}
\]

后，有一步因式分解。假设考场上你确信算出来的这个式子是对的，你是否能想出一种办法，不做因式分解，直接得到这个定值？\\
提示：先把 \(a, b, c\) 这些双曲线背景视为静止常量，分子和分母都是关于主自由度的什么函数？这种函数比如果要出现定值一定是怎样的？结合不变量的问题逻辑思考，这样得到的定值在逻辑上是否有保障？如果这样做了还有时间，该怎么通过即时检查来确保准确性？

\section*{6．5．2 方法预判助力圆梦高分}
\section*{题目：}
过坐标原点 \(O\) 作圆 \(C:(x+2)^{2}+y^{2}=3\) 的两条切线，设切点为 \(P, Q\) ，直线 \(P Q\) 恰为抛物线 \(E: y^{2}=2 p x(p>0)\)的准线\\
（1）求 \(E\) 的标准方程\\
（2）设 \(T\) 是圆 \(C\) 上的动点，扡物线上的四点 \(A, B, M, N\) 满足 \(\overrightarrow{T A}=2 \overrightarrow{T M}, \overrightarrow{T B}=2 \overrightarrow{T N}\) ，设线段 \(A B\) 中点为 \(D\)\\
（i）求直线 \(T D\) 的斜率\\
（ii）设 \(\triangle T A B\) 的面积为 \(S\) ，求 \(S\) 的最大值

题目容量比较大，第一问是个基本的纯静态问题，用于控制 \(P, Q\) 的元问题——直线与圆相切，也是基础问题。算出切线方程 \(y= \pm \sqrt{3} x\) ，以及切点横坐标 \(x=-\frac{1}{2}\)（这个式子也可以看成准线），就知道抛物线是 \(y^{2}=2 x\) 。值得提醒的是，第一问如果算错，后面一定就全完了，所以我第一问都会做简单的即时检查。比如这个，既可以联立直线与圆，也可以根据两条切线倾斜角 \(\frac{\pi}{3}, \frac{2 \pi}{3}\) 来直接几何法计算准线。方法与计算差异基本保证了不会出错。

关键来看第（2）问，整个第（2）问没出现过 \(P, Q\) ，所以直接从图里清除掉。 \(T\) 是静止背景下的一个动点，完全自由，肯定是作为控制关系的起始——主自由度。这时，一下子出现了四个点 \(A, B, M, N\) 。构建要素网一定不可能直接搞四个动点去瞎控制，需要根据已知的牵制关系去尝试在图中画出它们，看看它们有没有自由度，能不能被控制，若不能则考虑把主自由度给准。先握住控制杆 \(T\) ，尝试寻找 \(A, B, M, N\) 。它们的信息就是那两个已知的向量等式，并且这两个等式把 \(A, M\) 放一起，\(B, N\) 放一起，所以显然接下来的要素关系网需要分成两个平行枝，并且这两个枝还看上去是对称的。

我们先来看 \(A, M\) 。已知做一下充要替换简化，就是：\(M\) 是 \(T A\) 的中点，别的信息没了。这能够通过控制杆 \(T\) ，完全把 \(A, M\) 停下来吗？当然是可以的，理解方法也很多。比如：\(A, M\) 两个自由度，向量等式对应两个坐标相等，减两个自由度，静止；给 \(A\) 一个自由度，牵制为＂TA 中点在抛物线上＂，静止．总之，\(T\) 就是能控制住 \(A, M\) ，没问题。另一个枝 \(B, N\) 也是同理。注意控制并不代表只有一组唯一解，出现两组解也是自由度零，也是静止。这两个平行枝就对应＂\(T A\) 中点在抛物线上＂的两组解。再往后，\(D\) 就直接被 \(A B\) 控制，这个很简单。

所以，关系网就很清楚了：主自由度就是 \(T\) ，靠一个牵制＂TA 中点在抛物线上＂的两组解定死 \(A, B, M, N\) ，最后再控制一个 \(D\) 。第（i）问求 \(T D\) 斜率，自然就得沿着这条网走：先通过 \(T\) 把 \(A, B\) 的信息量表达出来，然后用这些信息量把 \(D\) 表达出来即可。但是，选好关系网，并不代表真正计算就完成了。而且，你的要素网建立的好不好，是不是这条路走下去的计算量就能接受，也还值得商榷。要素网的选择对于计算体验影响很大，这种事情在难题里经常出现。所以，面对难题，养成预判计算的习惯是很重要的。提前让大脑沿着要素网走一下，想想自己每一步都要算什么，这一步计算是不是能完成，计算得到的控制关系是不是保证了信息量，确保能继续算下去的。要素关系网的作用就在这里，你看着要素网，就知道下一步该研究谁与谁了，并且你的知识都是以元问题的姿态储存在大脑里的，随时调用各个控制的计算方法也很方便。反之，不做好预判的话，就有可能算了一大顿，到最后把自己绕糊涂了。

我们来预判一下（i）的计算。 \(T\) 是主自由度，肯定要设其坐标 \(\left(x_{0}, y_{0}\right)\) 作为控制杆的（当然还要加个约束\\
\(\left.\left(x_{0}+2\right)^{2}+y_{0}^{2}=3\right)\) 我们需要得到 \(D\) 的坐标信息，为此，需要先得到两点 \(A, B\) 的坐标信息，而且，最好是， \(x_{A}+x_{B}, y_{A}+y_{B}\) 能通过这些信息直接算出来（使用＂中点＂这一元问题进行预判）。现在思维焦点放到 \(A, B\) 上。刚才分析要素网就说了，两个向量等式的理解方法很多，而为了直接研究 \(A, B\) 的坐标，我们希望能把自由度就直接分配给这二者的坐标，然后利用这个坐标，把牵制变成方程。这样的话，\(T\) 是控制杆，\(A\) 是控制杆，就得让它们控制出 \(M\) ，然后另一侧的 \(B, N\) 计算则完全一样。所以，这就相当于研究一侧，得到方程，这个方程的两组解就分别对应 \(A\) 和 \(B\) 。这样的话，接下来的目标就很明确：先研究 \(A, M\) 这一组的控制与牵制，把它变成方程，然后直接解方程就出来了 \(A\) 和 \(B\) 的信息。从这个角度来讲，刚才设想 \(x_{A}+x_{B}, y_{A}+y_{B}\) 计算似乎立刻就成为了现实：这不就是韦达定理吗？我们一下子就看清了前面的道路，信心满满，确信接下来就需要把 \(A, M\) 这组牵制变成方程即可。

那么，怎么控制呢？\(M\) 可以用＂中点＂来控制，也可以用＂与抛物线的交点＂来控制。这二者的牵制就分别是＂ \(M\) 在抛物线上＂与＂\(M\) 恰为 \(T A\) 中点＂。于是，我们得到了两个列式方法。选哪个呢？就要做计算预判了：前者的话，利用 \(A\) 和 \(T\) 的坐标，表示 \(M\) 坐标很方便，＂\(M\) 在抛物线上＂列式也很方便；后者的话，需要先联立抛物线算交点，不过所幸 \(A\) 作为控制杆，在列式时默认是已知的，而＂已知直线与抛物线一交点，求另一交点＂正是我们讲过的解点模型。所以也能算。不过感觉上还是第一种控制更舒服一点，我们就选第一种控制吧。当然第二种的计算大家也可以去试试，肯定也能做。

开算：设 \(A\left(x_{A}, y_{A}\right)\) ，约束是 \(y_{A}^{2}=2 x_{A}\) ，那 \(T A\) 中点就是 \(\left(\frac{x_{0}+x_{A}}{2}, \frac{y_{0}+y_{A}}{2}\right)\) ，它在抛物线上，牵制就是 \(\left(y_{0}+y_{A}\right)^{2}=\) \(4\left(x_{0}+x_{A}\right)\) 。注意我们的目标是用 \(T\) 刻画 \(A, B\) 的信息，所以这个式子要把 \(x_{0}, y_{0}\) 先当已知，把 \(A, B\) 坐标当成未知数。这时，\(A, B\) 坐标的方程就列好了：

\[
\left\{\begin{array}{l}
\left(y_{0}+y\right)^{2}=4\left(x_{0}+x\right) \\
y^{2}=2 x
\end{array}\right.
\]

解这个方程得到的就是 \(A, B\) 的坐标，这是刚才的要素网告诉我们的。接下来肯定是消去 \(x\) 变成二次方程：\(\left(y_{0}+y\right)^{2}=\) \(4\left(x_{0}+\frac{y^{2}}{2}\right)\) ，展开整理一下，得到 \(y^{2}-2 y_{0} y+4 x_{0}-y_{0}^{2}=0\) ，所以用韦达定理得到 \(\left\{\begin{array}{l}y_{A}+y_{B}=2 y_{0} \\ y_{A} y_{B}=4 x_{0}-y_{0}^{2}\end{array}\right.\) ，这下就彻底用 \(T\) 把 \(A, B\) 的纵坐标信息刻画完整了，再配合 \(y^{2}=2 x\) 还能得到 \(x_{A}+x_{B}=3 y_{0}^{2}-4 x_{0}\) 。

信息量完整了，我们只需要关注上面的韦达定理式即可。（i）要算 \(T D\) 斜率，那就算呗！我现在控制式子在手，指哪打哪。

\[
D\left(\frac{3 y_{0}^{2}-4 x_{0}}{2}, y_{0}\right)
\]

䴲子都看得出来 \(k_{T D}=0\)（你甚至都不用算 \(D\) 的横坐标就能看出来）。\\
（ii）要算 \(\triangle T A B\) 的面积不是？\(T D\) 都水平了，算三角形面积这个元问题早就已经烂透了。圆雉曲线大题算面积，经典 \(S=\frac{1}{2}|T D| \cdot\left|y_{A}-y_{B}\right|\) 。这个式子现在什么东西我不知道？

\[
\left|y_{A}-y_{B}\right|=\sqrt{\left(y_{A}+y_{B}\right)^{2}-4 y_{A} y_{B}}=\sqrt{4 y_{0}^{2}-4\left(4 x_{0}-y_{0}^{2}\right)}=2 \sqrt{2} \sqrt{y_{0}^{2}-2 x_{0}}
\]

这不就一韦达控制式的附庸品？\(|T D|=\left|x_{0}-x_{D}\right|=\frac{3}{2}\left|y_{0}^{2}-2 x_{0}\right|\) ，也很简单。这下挺好，这俩玩意结构还差不多，最后就剩个 \(y_{0}^{2}-2 x_{0}\) 需要研究了。换个元，设 \(u=y_{0}^{2}-2 x_{0}\) ，那么最终面积就是 \(S=\frac{3}{2} \sqrt{2}|u| \sqrt{u}\) ，而 \(u\) 自然是直属于 \(T\) ，把 \(x_{0}, y_{0}\) 的原始约束 \(\left(x_{0}+2\right)^{2}+y_{0}^{2}=3\) 拿来，得

\[
u=3-\left(x_{0}+2\right)^{2}-2 x_{0}=-x_{0}^{2}-6 x_{0}-1=-\left(x_{0}+3\right)^{2}+8
\]

其中 \(x_{0}\) 遍历 \([-2-\sqrt{3},-2+\sqrt{3}]\) 。不难看出恒有 \(u>0\) ，且 \(u\) 最大就是 8 ，等号成立在 \(x_{0}=-3, S_{\max }=\frac{3}{2} \sqrt{2} \times 8 \times 2 \sqrt{2}=\) 48，题目结束。

\section*{总结：}
这道题的重点是想教会大家如何进行预判。对于复杂的题目而言，借助要素网提前进行预判是非常重要的。一道复杂题目可能有多种建立要素网的方式，我们需要从中选出那些必定能做出题来的，稳定性高的要素网，进行最终的计算。因此，提前在脑内预演一遍，沿着要素网调用一下元问题试试，想一想计算出来会是什么形式，信息量（充要替换）是否能保证？有没有无端的根号？这一步的计算结果大概支不支持下一步的元问题调用？如果不行，就试着换一种画图方法去思考。其实和大家平时做题没什么差别，无非就是：重视要素网，沿着要素网去预判会非常清晰；更改方法也要对本身的要素网结构下手，结合自己掌握的元问题，以及画图法的思维方式，去进行思考。

如何练习预判？其实也很简单，拿过一道题，先不要计算，就简单画画图，想想元问题，模仿我今天这道题前面的思维过程，去做做预判练习。做完之后再下笔计算，看看计算结果和答案方法和你之前预判的设想是否吻合。如果不吻合，那一定是预判的某些地方存在问题，把预判过程翻出来查出问题所在，是思维没转过来，还是元问题不熟练？如果答案有不同的方法，就想一想当时预判是否设想到了这种做法？和自己的做法比起来如何？之类的。一开始练不好很正常，但是如果你有志向冲击 140 分的话，预判是一种需要具备的能力。

注意！！预判是在大概有思路的前提下事先评估思路的方式，目的是更加清晰的看清问题，少走弯路。如果是那种难到乍一看没什么思路的问题，就要采取别的策略，比如试验，设问，遍历找感觉等等，参见之前其他的视频。

思考题：\\
在用 \(T\) 刻画 \(A, B\) 信息时，我们发现还有一种解点的思路，请沿着这条思路完成 \(T\) 对 \(A, B\) 的控制，这个结果和原题是否相同？你如何看待这个现象？这是否是一种逻辑上的必然？你能试着利用控制与信息量（充要替换）的原理给出解释吗？

\section*{6．5．3 长题干与巧解法都能一眼到底}
视频日期：2023年11月25日

题目：\\
2022年12月4日20点10分，神舟十四号返回舱顺利着陆，人们清楚全面地看到了神舟十四号返回舱成功着陆的直播盛况。根据搜救和直播的需要，在预设着陆场的某个坚直平面内设置了两个固定拍摄机位 \(A, B\) 和一个移动拍摄机位 \(C\) 。根据当时的气候与地理特征，点 \(C\) 在抛物线 \(\Gamma: y=\frac{1}{36} x^{2}\) 上（直线 \(y=0\) 与地平线重合，\(y\)轴垂直于水平面，单位：十米，下同），且 \(C\) 的横坐标 \(x_{C}>6 \sqrt{2}\) 。已知 \(A\) 的坐标为 \((-36,2)\) ，设点 \(D(0,-2)\) ，线段 \(A C, D C\) 分别交 \(\Gamma\) 于点 \(M, N, B\) 在线段 \(M N\) 上。则两固定机位 \(A, B\) 间的距离为（）\\
A， 360 m\\
B， 340 m\\
C， 320 m\\
D， 270 m

这道题目的题干属于现实情景类，而且比较长。但是即便是这样，动静视角总览全局，以及分析要素结构的思维方式还是不变的，无非是以前直接看已知构建要素关系网，与现在理解题意构建要素关系网的区别。情景是

安放拍摄机位，翻译成要素关系就是研究机位的位置（点要素）。首先是两个静止机位，也就是静态点 \(A, B\) ，然后一个单自由度动点 \(C\) 在 \(\Gamma: y=\frac{1}{36} x^{2}\) 的一部分上运动。然后要素关系网自然地往前走，\(A C, D C\) 分别交 \(\Gamma\) 于点 \(M, N\) ，这一句全都是控制关系。所以目前关系网就是：\(A, B\) 静态，\(A\) 位置已知，\(B\) 具体位置不明（但确定），然后 \(C\)带有一个自由度。因为题目并没有告诉我们 \(B\) 在哪，再加上刚才的要素关系网它也没出现，所以暂时放在一边。于是题目的状况就是一个单自由度的要素关系网再加一个不明位置的定点 \(B\) 。

但是这时，唐突出现了一个条件：\(B\) 在线段 \(M N\) 上。这是什么情况呢？首先，目前题目并没有告诉我们 \(B\)的具体位置，所以还在迷雾之中。而且画了要素关系网就知道，关系网到目前为止 \(B\) 都没有参与进来。所以这个已知条件是要告诉我们 \(B\) 到底在哪的。注意 \(C\) 作为移动机位，这个自由度是肯定有的，不可能被解出（不然就是题目出错了）。那么，线段 \(M N\) 就一定是一个单自由度系统控制的要素。所以题目的含义就很明确了：这个单自由度控制的 \(M N\) 需要过定点，而这个定点就是我们要找的 \(B\) ！这就告诉我们：这道题其实正是我们重点讲过的不变量问题，题目会给你一个自由度非 0 的系统，然后开上帝视角的天眼告诉你这个系统可以输出一个不变量，让你验证并计算这个不变量。因此，这个现实情景的长题干就被我们利用动静视角与要素关系网的思维给转化成一个纯数学的问题了。

转化完的问题如下：单自由度动点 \(C\) 在 \(\Gamma: y=\frac{1}{36} x^{2}(x>6 \sqrt{2})\) 上运动，点 \(A(-36,2)\) ，点 \(D(0,-2)\) 。线段 \(A C, D C\) 分别交 \(\Gamma\) 于点 \(M, N\) ，证明线段 \(M N\) 过定点，并求出该定点。下面自然就是依据要素关系网，选择合适的元问题展开计算了。要研究 \(M N\) ，自然就是要研究 \(M, N\) 两点，而这两点又是抛物线上动点分别与两个不同的定点连线，再与抛物线交出来的点。这一个控制关系能对应什么元问题呢？太明显了，不就是解点法吗？我在＂数学的干货＂视频列表里讲解点的时候，说的元问题和这个几乎一字不差了，但凡认真看过视频的同学，都不可能想不到。做两次解点把 \(M, N\) 控制出来，然后看定点即可。预判一下计算量大吗？抛物线解点本来就比椭圆双曲线简单，对好开口方向写直线列方程（上下开口列 \(y=k x+b\) ，左右开口列 \(x=t y+m\) ）计算即可，解两次点也可以接受。然后写个直线方程也轻而易举，判定可以算！

那就开算吹！先把主自由度给 \(C\left(t, \frac{1}{36} t^{2}\right)\) ，注意定义域 \(t>6 \sqrt{2}\)（设主自由度看定义域也强调多次了）。于是写出 \(C D\) 方程

\[
y=\frac{t^{2}+72}{36 t} x-2
\]

不想直接代入这么长串的斜率，可以用一下多级计算：设 \(k_{1}=\frac{t^{2}+72}{36 t}\) ，于是与抛物线联立得到 \(\frac{1}{36} x^{2}-k_{1} x+2=0\) ，于是用韦达定理得 \(t \cdot x_{N}=72\) ，这就控制了 \(x_{N}\) ，进而得

\[
N\left(\frac{72}{t}, \frac{144}{t^{2}}\right)
\]

同理，可以算出 \(A C\) 的斜率为

\[
k_{2}=\frac{t^{2}-72}{36(t+36)}
\]

于是其方程为 \(y-2=k_{2}(x+36)\)（当然用 \(C\) 写点斜式也无所谓，多级计算最后还原回去一定被 \(t\) 控制），代入抛物线并整理得

\[
\frac{1}{36} x^{2}-k_{2} x-36 k_{2}-2=0
\]

于是

\[
t+x_{M}=36 k_{2}=\frac{t^{2}-72}{t+36}
\]

进而得

\[
M\left(\frac{-36(t+2)}{t+36}, \frac{36(t+2)^{2}}{(t+36)^{2}}\right)
\]

我们已经实现了：用单自由度来控制 \(M, N\) 。接下来有两种方法可选，对应于不同的要素网理解。单自由度问题如何处理？一种处理方式是控制走到底：用 \(t\) 表达线段 \(M N\) 的方程。因此计算

\[
k_{M N}=\frac{\frac{(t+2)^{2}}{(t+36)^{2}}-\frac{4}{t^{2}}}{\frac{-(t+2)}{t+36}-\frac{2}{t}}=\frac{2}{t}-\frac{t+2}{t+36},
\]

于是直线 \(M N\) 的方程就可以直接写出来：

\[
y-\frac{144}{t^{2}}=\left(\frac{2}{t}-\frac{t+2}{t+36}\right)\left(x-\frac{72}{t}\right) \Leftrightarrow y=\frac{72-t^{2}}{t(t+36)} x+\frac{72(t+2)}{t(t+36)} .
\]

注意线段交点一定是直线交点，所以我们先看直线交点。这个方程能过某个定点吗？那就要注意观察了，我们希望 \(x\) 取某个数值的时候，分子的 \(-x t^{2}+72 t+72(2+x)\) 和 \(t(t+36)\) 可以彻底约分，那一个直接感受就是得让前者的常数项 \(72(2+x)=0 \Leftrightarrow x=-2\) ，这时分子变为 \(2 t(t+36)\) ，约掉分母变成常数 2 ，所以我们就看出来了它的确过定点 \((-2,2)\) 。

本题还有另一种处理思维，这是另一种对要素网的理解。 \(M, N\) 的坐标我们已经用单自由度刻画清楚了，而我们要研究直线 \(M N\) 这一对象，我们刻画的是它与抛物线的交点坐标。直线，与抛物线交点坐标，这一控制的元问题是什么？当然是联立韦达定理了！注意一般直线联立的时候，比如我们设 \(M N: y=k x+b\) ，因为直线方程含有两个主自由度，因此需要韦达定理的两个式子才能把 \(k, b\) 的信息量全部刻画出来。比如这里联立 \(M N\) 与抛物线得 \(\frac{1}{36} x^{2}-k x-b=0\) ，可知 \(x_{M}+x_{N}=36 k, x_{M} x_{N}=-36 b\) 。这两个式子都是必要的，丢失任何一个信息量都会流失，而不像解点那样任取一个韦达定理式子就能完成全部控制（因为解点法是已知了一个交点的信息了，信息量更多）。这样的话，我们希望用 \(M, N\) 的坐标之间的关系来等价刻画 \(k, b\) 之间的关系。因为 \(t\) 是一个自由度，所以如果想要充要替换成两个参数之间的式子，就应当是一个式子。获取方法就是消元法：利用 \(x_{N}=\frac{72}{t}\) 与 \(x_{M}=\frac{-36(t+2)}{t+36}\)消去 \(t\) 得

\[
x_{M}=\frac{-36\left(\frac{72}{x_{N}}+2\right)}{\frac{72}{x_{N}}+36} \Leftrightarrow 2\left(x_{M}+x_{N}\right)+x_{M} x_{N}=-72
\]

后面这个式子牵制了 \(x_{M}, x_{N}\) ，也就是牵制了 \(k, b\)（两个参数）之间的关系，于是自由度是 1 ，与原题总自由度守恒，因此信息量与原题保持一致。现在我们取了 \(k, b\) 作为主自由度来刻画直线 \(M N\) 的信息了，用韦达定理可以实现 \(k, b\) 对 \(x_{M}, x_{N}\) 的控制，接下来只需要把牵制关系转回 \(k, b\) 之间的式子，即可实现单自由度对直线 \(M N\) 的控制。代入韦达定理式得到 \(-2 k+b=2\) 。直线本身是 \(y=k x+b\) ，对照一下即可看出它过定点 \((-2,2)\) 。

不管用哪种方法，定点都算是找出来了，当然保险起见，大家还需要验证一下它确实是在线段 \(M N\) 上，这个留作练习。但定点肯定只有一个，只要题目没出错，那么 \(B\) 就一定是 \((-2,2)\) 。最后算出 \(A, B\) 之间距离为 \(|-36-(-2)|=\) 34 ，单位为十米，于是就是 340 m 。

做个总结。本题仍然是使用 up 主反复讲过的思维方式，按部就班做题，碰到该用什么思维了就直接用啥思维，都不带拐弯的。先把问题的要素关系网建构清楚，然后把要素关系网对应的元问题找出来，加以预判可计算之后直接开算。比较无脑的解法就是第一条路，直接用 \(t\) 一路控制到底，利用分子分母的多项式可以互相约分的特性把定点找出来。第二条路则是对于当前要素关系：＂坐标 \(X_{M}, x_{N}\) 控制直线＂的另一种理解，直线也可以反过来控制坐标，这是一个双向控制。遇到可以双向控制的要素的时候，颠倒要素网常常是一种可选的办法，要素网可以颠倒这件事情我们也早在经验分享第四期就给大家提到了。不管选用什么样的要素关系网，只要控制减牵制得到的总自由度不变，那么系统的运动幅度也就不变，信息量也就维持了稳定。而有时根据当前要素网的内容 （元问题）的认识，进行适当的要素网改写，常常能极大地改变计算的形式。如果预判某一步的计算量过于不稳

定，像本题第二条路这样，重新选取主自由度改写一下要素关系网，也是一条可选的，并且时不时就能带来惊喜的路径。可以看到，本题第二条路径的计算量明显低于第一条路径，就说明了这一点。

时间规划：\\
读题，建立要素关系网： 2 min\\
梳理要素关系网，选取并预判解题方法： 1 min\\
利用解点法得到 \(M, N\) 的坐标：3．5min\\
找出定点，得出答案： 2.5 min\\
总时间： 9 min

\section*{6．5．4 一道易错题的全思维解析}
题目：\\
已知点 \(A, B\) 位于抛物线 \(y^{2}=2 p x(p>0)\) 上，\(|A B|=20\) ，点 \(M\) 为线段 \(A B\) 的中点，记 \(M\) 到 \(y\) 轴的距离为 \(d\) 。若 \(d\) 的最小值为 \(\frac{64-2 p}{7}\) ，则当 \(d\) 取得该最小值时，线段 \(O M\) 的长度为 \(\qquad\) －

这道题目作为 2024 年高考前最后一题，目的就是再次强调一个 up 主反复强调多次了的要点内容。首先还是我们的惯例，读题分析要素结构。首先开局一个自由度 \(p\) ，作为背景把抛物线本体控制住。在此基础之上，抛物线上两个动点（两个主自由度）的 \(A, B\) ，被一个牵制 \(|A B|=20\) 锁在一起。如果把 \(p\) 暂时冻结，那这就是个单自由度系统。正好紧接着，题目给出了一个 \(d\) 。这个 \(d\) 明显是针对 \(A, B\) 这个单自由度系统的，因为后面那上就叙述到了 \(d\) 的最小值。如果把这个 \(d\) 看作是 \(p\) 和 \(A, B\) 组合共同控制的（被双自由度控制的）参数，那么它在两个参数意义下的最小值显然就是 0 （把抛物线口径无限撑大即可），这个题也就没啥意思了。题目的真实意义是：冻结 \(p\) ，考虑 \(d\) 被单自由度系统（主自由度信息在 \(A, B\) 组合中）控制时的最小值，这个最小值就成了只被 \(p\) 控制的参量。这个参量要求等于 \(\frac{64-2 p}{7}\) ，牵制关系一搭，\(p\) 就求解出来了，进而全题就静止了，该求啥求啥即可。

于是，我们第一层思维焦点就自然落在了：＂如何用 \(p\) 表示这个 \(d\) 的最小值＂上。思路自然也是老套路：选明确的参数作为主自由度，控制出 \(d\) ，然后处理函数求最值。选谁作为主自由度呢？\(A, B\) 是一个整体自由度为 1 的系统，但作为两个点配一个牵制的类型，想直接用一个主自由度控制到底不太现实。稳妥的考量还是选两个主自由度，然后把牵制表达成这两个主自由度之间的式子。而牵制是一个长度形式，调取和长度相关的元问题，发现这是在圆锥曲线问题里极为常见的类型：弦长问题。弦长问题选主自由度通常是设直线方程，而本题是抛物线，抛物线问题也可以选择直接设点的坐标。因此我们就有了两条选主自由度的方式。前者是设 \(A B\) 的方程。因为是横着的抛物线，应该选 \(x=t y+m\) 的形式，主自由度就是 \(t\) 和 \(m\) 。后者就很简单粗暴：设

\[
A\left(\frac{y_{1}^{2}}{2 p}, y_{1}\right), B\left(\frac{y_{2}^{2}}{2 p}, y_{2}\right)
\]

选好主自由度之后需要做简单的预判。前者的设法就是联立用弦长公式得到 \(t, m\) 之间的表达式；而所求的 \(d\) 就是弦中点的横坐标，这种东西显然韦达直接顺畅控制。整体预判感觉不错。后者就直接用距离公式表达，因为横坐标都是平方，所以作差后可以平方差分解，整体有公因式，是一个很大的简化（＂简化即胜利＂原则）；至于所求，仍然是中点，表示起来也很轻松。因此两种设法预判起来都没有明显的障碍阻止我们这样去做。

既然如此，那选啥都行了。比如我选第一种设直线。这里就出现 up 主另一个强调多次的要点——定义域了。主自由度必须必须必须确定好它们的定义域。首先，\(x=t y+m\) 能不能遍历所有的题目叙述的情形？要求是和抛

物线有两个交点，而这种设法只是不能表示水平直线，但水平直线肯定不会和抛物线有两个交点。因此不需要额外讨论。再一个就是两个交点等价于 \(\Delta>0\)（注意范围类判断必须充要替换）。定义域明确了，那就开始计算：联立得 \(y^{2}-2 p t y-2 p m=0\) ，于是 \(\Delta>0 \Leftrightarrow p t^{2}+2 m>0\) 。然后算弦长。由韦达定理得 \(y_{1}+y_{2}=2 p t, y_{1} y_{2}=-2 p m\) ，于是

\[
\left|y_{1}-y_{2}\right|=\sqrt{(2 p t)^{2}-4(-2 p m)}=2 \sqrt{p^{2} t^{2}+2 p m}
\]

于是弦长的平方为

\[
4\left(1+t^{2}\right)\left(p^{2} t^{2}+2 p m\right)=400
\]

因此

\[
\left(1+t^{2}\right)\left(p^{2} t^{2}+2 p m\right)=100
\]

把它和定义域结合起来看，很容易发现这个牵制式里面自带一个 \(p t^{2}+2 m\) 。由于 \(1+t^{2}>0\) ，因此在牵制式成立的前提下，\(p t^{2}+2 m\) 就自动为正了。所以这个定义域限制已经包含在牵制里面了，我们可以抛弃 \(\Delta>0\) 这个限制了。（但这并不代表刚才的判别式计算是不必要的！！定义域判断不能省略！！）

牵制搞完了，接下来让我们把思维焦点放到所求上。所求刚才也预判过了：韦达定理搞出

\[
\frac{x_{1}+x_{2}}{2}=\frac{t\left(y_{1}+y_{2}\right)+2 m}{2}=p t^{2}+m
\]

现在式子已经齐全，我们把思维焦点全部集中到刚才算出来的牵制和所求上来。注意现在 \(p\) 我们视为静止，只关注 \(t\) 和 \(m\) 。怎么通过已知的牵制搞出所求的最小值呢？双自由度加牵制的类型，如果没有什么特别好的办法的话，最自然的尝试就是消元，消到未知量数目等于总自由度数目 1 为止。换元产生多级计算结构或者转嫁自由度也是常见的技巧（这些结构层面的元问题理解一定要扎实）。这里显然，只看 \(m\) 的话，它只是个一次方程，因此用 \(t\) 表示 \(m\) 可谓轻而易举：

\[
m=\frac{50}{p\left(1+t^{2}\right)}-\frac{p t^{2}}{2}
\]

这个过程显然是充要替换，因此可以认为 \(t\) 定义域是整个 \(\mathbb{R}\) ，然后由 \(t\) 控制 \(m\) 。这样主自由度的定义域彻底搞清楚了，神清气爽。然后带人到所求里面就是

\[
p t^{2}+\frac{50}{p\left(1+t^{2}\right)}-\frac{p t^{2}}{2}
\]

因此所求变成了主自由度 \(t\) 控制的

\[
\frac{50}{p\left(1+t^{2}\right)}+\frac{p t^{2}}{2}
\]

还是别忘了 \(p\) 静止。这个式子的结构看上去是不是很亲切？是不是很有对号函数的感觉？配一下常数，改写成

\[
\frac{50}{p\left(1+t^{2}\right)}+\frac{p\left(t^{2}+1\right)}{2}-\frac{p}{2}
\]

再做个换元：令 \(u=p\left(1+t^{2}\right)\) ，定义域必须随时跟进：\(t\) 遍历 \(\mathbb{R}\) 时 \(u\) 遍历 \([p,+\infty)\) 。所求变成

\[
\frac{50}{u}+\frac{u}{2}-\frac{p}{2}
\]

接下来自然就是求最小值啦。这个函数在 \(u \in(0,10]\) 时递减，在 \(u \in[10,+\infty)\) 时递增。但现在我们并不知道 \(p\) 的具体值！！所以要根据 \(p\) 的范围情况讨论！！！！！如果 \(p \leq 10\) ，则对号函数的最小值点包含在定义域内，其最

小值就是 \(10-\frac{p}{2}\) ；如果 \(p>10\) ，则对号函数在 \([p,+\infty)\) 上一直增，其最小值应该是

\[
\frac{50}{p}+\frac{p}{2}-\frac{p}{2}=\frac{50}{p} .
\]

对于第一种情况，求解最后的牵制

\[
10-\frac{p}{2}=\frac{64-2 p}{7}
\]

得到的是 \(p=4\) ，在这个定义域范围内，保留；如果是第二种情况，则最后的牵制应该是 \(\frac{50}{p}=\frac{64-2 p}{7}\) ，解出来 \(p=25\) 或 7 。卡一下定义域范围，只保留 \(p=25\) 这个答案。

一切静止化完成后，来做最后的答案计算。 \(p=4\) 时，等号成立条件是 \(u=10\) 。立刻算出 \(t= \pm \frac{\sqrt{6}}{2}\) ，对应 \(m=2\)。此时 \(M\) 坐标很容易算出是 \((8, \pm 2 \sqrt{6})\) ，因此 \(|O M|=2 \sqrt{22} ; p=25\) 时等号成立条件是 \(u=p\) ，也就是 \(t=0\) ，即 \(A B \perp x\) 轴的情形。此时易算出 \(m=2\) ，以及 \(M(2,0)\) ，因此 \(|O M|=2\) 。最后的答案是 \(2 \sqrt{22}\) 或 2 ，你答对了吗？

做个总结。这道题其实就是我们的思维方式的完全按部就班的实战演练。其中的常见的要点和易错点也都涉及到了。基本的要素关系结构的思维方式，基本的预判加上元问题调用，定义域的判断，充要替换。这道题最大的易错点就是定义域的持续关注与充要替换。定义域在up主的视频里一直是超高频强调的，重点中的重点。遍历法的核心要点之一就是定义域，充要替换的重要步骤也是关注定义域。只要是要素网加控制就离不开定义域。希望这道题能让大家再把定义域的重视程度拔高一点，不要在高考中因为定义域出现失误，把会做的，该拿的分数都拿到，取得符合自己实力的成绩。

课后练习：\\
用我们视频里预判时提到的第二种设点的方法，把这道题再做一遍。

\section*{6．5．5 深度解析 2020 浙江圆曲大题}
\section*{题目：}
如图，已知椭圆 \(C_{1}: \frac{x^{2}}{2}+y^{2}=1\) ，抛物线 \(C_{2}: y^{2}=2 p x(p>0)\) ，点 \(A\) 是 \(C_{1}\) 与 \(C_{2}\) 的交点，过点 \(A\) 的直线 \(l\) 交 \(C_{1}\) 于点 \(B\) ，交 \(C_{2}\) 于 \(M(B, M\) 均异于 \(A)\) 若存在不过原点的直线 \(I\) 使 \(M\) 为 \(A B\) 的中点，求 \(p\) 的最大值\\
\includegraphics[max width=\textwidth, center]{2025_02_16_7c9aaf92bc99dc07f048g-255}

读完题，我们发现这是一个典型的单自由度封装命题，真自由度就是 \(p\) ，它能直接控制 \(A\)（怎么控制的？小小元问题），并且是一个双向控制。

然后是一个被封装的自由度：\(l\) 。当 \(A\) 确定下来之后，\(l\) 被一个自由度控制，但是题目要求＂存在 \(l \ldots\) ．．＂，这就是番外篇里讲到的封装了。＂存在 \(l\)＂的说法让 \(l\) 的自由度变成了约束变元，矛盾指向的还是 \(A\) 。因此，真自由度只有 \(A\) 一个。\\
（以上看不懂的建议先退出回去复习一下番外篇）

因此，接下来要做的，就是封装命题的一般思维：消去约束变元，把问题转化为一个真正的描述真自由度的要求的问题。思维方法番外篇也讲过了：先遍历加冻结，随便点一个真自由度（这里要点 \(A\) 的位置，有了 \(A\) 就可以返回去得到 \(p\) ，但是先点一个 \(p\) 还要求解 \(A\) ，我们的宗旨是逆着要素网进行冻结，这样才能一层层扒掉约束变元的伪装）

于是随便画一个抛物线（比如前面的那个图），现在该干什么呢？当然是继续思维：设问！什么样的 \(A\) 能找到这样的直线？一旦自己这么自问了，马上就想：既然我现在当成 \(A\) 静止了，那就尝试计算直线 \(l\) ，能算出来那这个 \(A\) 就好，算不出来就不好呗！（封装命题的基本逻辑就是这样，按照它封装的要求去尝试计算转化，＂任意＂就玩恒成立，＂存在＂就尝试计算寻找，都要在转化过程中考虑充要替换的方向）

先设一个 \(A\left(x_{1}, y_{1}\right)\)（暂时默认 \(x_{1}, y_{1}\) 静止，并在心里记住它可以被 \(p\) 直接控制）。那接下来不就该设直线了吗？直线设成什么样？（思维焦点稍微放过来）显然斜率存在且不为 0 ，那设 \(y=k x+b\) 就行（当然 \(x=t y+m\)也行，随意，这个不产生本质差别）

那就要尝试解出直线了（注意这里目标早已明确，\(A\) 静止，加上中点的牵制，直线的自由度应是 0 了，肯定要尝试解方程了）。确定 \(k, b\) 的第一个方程（牵制：\(A\) ）是 \(y_{1}=k x_{1}+b\)

第二个方程呢？第二个方程就是＂截粗圆的中点在抛物线上＂，我们当前的思维焦点是：利用这个条件，再列出一个可以方便我们研究，讨论的方程（这也是本题的难点所在，毕竟是浙江高考的压轴大题之一，一点波澜没有岂不是很没面子）。

我们描述的主语是＂截粗圆的中点＂，谓语是＂在抛物线上＂，所以研究主体自然就是截椭圆的中点。问题的主语决定了我们操作的着眼点，于是自然的把思维焦点放上来：如何研究截椭圆的中点？把问题隔离出来，瞅着＂截 \(F\) 圆的中点＂这么一想，稍微一思考自己学过什么和椭圆中点有关的东西（基本元问题的扎实程度），立刻就想到了点差法。

于是锁定点差法，迅速先完成这一套默写步骤：设直线 \(y=k x+b\) ，再设 \(B\left(x_{2}, y_{2}\right)\) ，把 \(A\) 和 \(B\) 代入栯圆方程，作差，分解，整理一套流程（大家自己练习），得 \(k_{O M} \cdot k=-\frac{1}{2}\)（注意这里 \(A B\) 不过原点，所以这些斜率都是存在的）

对点差法的基本理解（对大家元问题的扎实程度的考察，相信挑战这道题的同学完全不在话下）告诉我们，上述这个斜率式子就完全刻画了椭圆，把它直接和 \(A B\) 联合就完全确定 \(M\) ，于是问题充要替换为：存在 \(l\) ，使得 \(k_{O M} \cdot k=-\frac{1}{2}\) 的直线交过来正好交在抛物线上，这符合我们之前讲过的消点思维一一简化要素网。现在这个问题已经和椭圆没关系了（除了确定 \(A\) 点时），我们只需要考虑，是否存在过 \(A\) 的直线 \(l\) ，使得上述一条 \(O M\) 射过来正好在抛物线上。因此，继续充要替换列方程。

这里就有一个选择，我们从要素网的理念去思考这个问题。现在的思维焦点涉及到三个线（ \(l, O M\) 和抛物线），并且是要搞一个共点问题，因此在控制和牵制的选择上就会有分歧。我可以任选两条线交一下出来一个交点，然后把牵制选取为＂第三条线过这个点＂。我是让 \(O M\) 射到 \(l\) 上，交点代入抛物线列式？还是 \(O M\) 射到抛物线上，代人直线 \(l\) 列式？还是考虑 \(l\) 和抛物线的交点，让它经过 \(y=-\frac{1}{2 k} x\) ？

这三个选择在逻辑上是等价的，必然都能列出式子，那为什么还会出现这样的选择分歧？这是因为，当逻辑

保证了你的充要替换之后，剩下的就只有算式处理，而算式处理不可避免会存在代数的充要替换（算式等价变形）的功夫，有的甚至会有比较强的技巧性，这个无法通过单纯的逻辑思维来解决。这是本题的一大难点所在。不知道咋办了？那就先试试，反正这个地方只是个小思维焦点，试一下错也无妨。第一种，让 \(O M\) 射到 \(l\) 上，联立两个直线解出

\[
\left(x_{M}, y_{M}\right)=\left(-\frac{2 k b}{1+2 k^{2}}, \frac{b}{1+2 k^{2}}\right),
\]

代入抛物线化简，得到 \(b=-4 p k\left(1+2 k^{2}\right) \cdots \cdots\)

好丑啊！这怎么和我们之前列的第一个方程 \(y_{1}=k x_{1}+b\) 了．\\
联立啊！都出来三次项了。注意！当你列出了一个奇丑无比的式子，并且自由度也是吻合的时候，很可能是你的代数变形没有做好。这个时候剩下的工作理论上只是处理算式而已，但是不同方法导致的计算过程的差别显然也是有繁有简的。不过好在，这里我们拥有多个选项，我们可以再做尝试。

第二条路，\(O M\) 射到抛物线上，代人直线 \(l\) 列式。 \(O M\) 和抛物线联立，解得 \(x_{M}=8 p k^{2}\) ，进而解出 \(y_{M}=-4 p k\) 。代入直线，得到 \(b=-4 p k\left(1+2 k^{2}\right) \cdots \cdots\) 怎么绕回来了？？？看来，这两条路都不太行

第三条路，考虑 \(l\) 和抛物线的交点，让它经过 \(y=-\frac{1}{2 k} x\) 联立 \(l\) 和抛物线（抛物线是横着的，所以重设 \(x=t y+m\) ），整理得 \(y^{2}-2 p t y-2 p m=0\) ，因为我们把 \(A\) 看作静止，所以自然想到之前讲过的解点思维一一用韦达定理来完成 \(A\) 对另一个交点（也就是 \(M\) ）的控制。

但又面临选择，我韦达定理用加还是用乘？\(y_{1}+y_{M}=2 p t\) 和 \(y_{1} y_{M}=-2 p m\) 选谁？它们分别是

\[
y_{M}=2 p t-y_{1} \text { 和 } y_{M}=-\frac{2 p m}{y_{1}} \text {, }
\]

我先保留一下，大家继续看。\\
横坐标表示完了，该纵坐标了。如果我选第一条路，求横坐标 \(X_{M}\) 我的结果都不好看（形式都比较复杂），所以估计后续再让它经过 \(y=-\frac{1}{2 k} x\)（现在应该是 \(y=-\frac{t}{2} x\) 了）也不会太好看（大家可以花两分钟试一试）。但我选第二个的话，如果再代人抛物线求解 \(x_{M}\) ，那么得到的是

\[
x_{M}=\frac{2 p m^{2}}{y_{1}^{2}}
\]

这个时候代入 \(y=-\frac{t}{2} x\) 得到的就是纯粹的乘法式子，和原本的三次多项式发生了本质改变（纯乘法谁不喜欢啊），这条路有门！

代 \(M\) 横纵坐标入 \(y=-\frac{t}{2} x\) 化简，得到 \(t m=2 y_{1}\) ，多么漂亮的结果！这下可以联动原来的 \(x_{1}=t y_{1}+m\) 尝试求解了。问题已经被我们全部转化掉（自由度也完全对得上），所以问题充要替换为：以真自由度为参数的，关于 \(t, m\) 的方程组

\[
\left\{\begin{array}{l}
t m=2 y_{1} \\
x_{1}=t y_{1}+m
\end{array}\right.
\]

有非零解（不过原点）。

这个已经是非常基础的问题了，显然可以假定 \(x_{1}, y_{1}\) 均为正，消元出来个 \(\frac{x_{1}}{y_{1}}=\frac{2}{t}+t\) ，右边最小值是 \(2 \sqrt{2}\) ，因此想要存在 \(t\) 满足上式，等价于 \(\frac{y_{1}}{x_{1}} \leq \frac{\sqrt{2}}{4}\) ，也就是 \(k_{O A} \leq \frac{\sqrt{2}}{4}\) ，剩下的就是问：满足 \(A\) 在椭圆上，\(k_{O A} \leq \frac{\sqrt{2}}{4}\)

的抛物线是什么情况？已是最基础的题目了，大家自行练习，答案是 \(p \leq \frac{\sqrt{10}}{40}\)

\section*{用时分析（面向有志于考场本题满分的同学）：}
1．读题 + 封装命题分析： 2 min\\
2．点差法： 2 min\\
3．想清楚有三条路： 1.5 min\\
4．前两条路试错：各 3 min\\
5．失败的思考： 1 min\\
6．第三条路的分歧选择与探索： 3 min\\
7．得到 \(\mathrm{tm}=2 y_{1}: 1.5 \mathrm{~min}\)\\
8．完成余下部分： 3 min\\
9．合计： 20 min

\section*{思考题：}
为什么前两条路失败了？为什么第三条路找出了 \(y_{1}, y_{M}\) 之间的关系，但是前两条路没有找到，你能否尝试结合 \(p, x_{1}, y_{1}\) 之间的固有关系，把前两条路得到的丑式子充要替换成像 \(t m=2 y_{1}\) 这种简洁明快的式子？

提示：这里必定需要先搞一个类似于 \(y_{1} y_{M}=-2 p m\) 这样的式子

\section*{6．5．6 结构思维快速看穿全题结构}
题目：\\
几何学中，把满足某些特定条件的曲线组成的集合叫做曲线族。点 \(Q\) 是椭圆族 \(T\) 上的任意一点，椭圆族 \(T\)的元素满足以下条件：（1）长轴长为 4 ；（2）一个焦点为原点 \(O\) ；（3）过定点 \(P(0,3)\) ，则 \(|Q P|+|Q O|\) 的最大值是\\
\(\qquad\) －

老规矩，拿到一道题先分析它的系统结构，尝试构建要素们的关系网。显然原点 \(O\) 和点 \(P\) 都是静止要素，但是系统中另一个最重要的要素 \(Q\)（显然它是运动的）则是在一族看上去很混乱的歪斜粗圆上乱动，看不清它怎么去处理。尝试去数 \(Q\) 的主自由度，显然如果给我们一个确定的椭圆的话，它的自由度就是1。因此我们需要搞清楚这个椭圆族的自由度是多少，然后在此基础之上添加 \(Q\) 的运动。但一般位置的椭圆的自由度我们并不熟悉，更不用说处理已知中的条件了。因此，直接计算自由度的方法在这里并不是很有效。那么，我们应该怎么办呢？

遇到不熟悉的情境怎么处理？这就要再次强调这个强调过多次的观点了：面对要素关系网的梳理时，画图法与工程师意识是最为行之有效的方法。我们要尝试画出题中的椭圆，在画图的时候，用以画图的依据，遇到的限制，确定与不确定的事物，就能如实给我们一个构建体系要素框架，确定主自由度（操纵杆）与国制条件的方式。因此，我们尝试来画一画题干要求的椭圆，看看上述内容给我们怎样的反馈。

既然要画一个位置不标准（焦点未必在坐标轴上）的椭圆，我们就先来想想，给我们什么信息，我们能画出（控制出）一个椭圆？随便画一个斜的椭圆，我们就容易知道，例如，如果告诉我们中心坐标，长短轴的长度，以及它们的方向，就能画出来。再比如，更加几何一点的，给出两个焦点外加 \(a\) 的值，就能画出来。这些都是控制

比较顺畅的方式，也是最容易想到的方式。那么，题目现在的信息有 \(a\) ，有一部分焦点信息，还有个过定点，这三个已知。从确定椭圆的信息而言，前两个无疑是从画图法角度而言感觉最舒服的。因此我们先尝试用前两个来作为生成椭圆的要素网中的控制关系，用这两个条件先尝试作为控制粗圆的操纵杆。一个焦点是原点，那我只需再知道另一个焦点 \(F\) 的位置，配合 \(a=2\) 的已知，就能把粗圆完全确定下来。因此，\(F\) 完全控制了粗圆，很顺畅。

那么，还有一个条件＂过定点 \(P\)＂呢？要素关系网由控制与牵制两部分构成，我们如果选 \(F\)（平面上随机一个点）作为控制椭圆的主自由度的话，剩下的条件自然就作为牵制。牵制是干什么的？是给系统添加额外的限制，以牵制主自由度们的。那么，我随便点一个 \(F\) ，怎么把＂过定点＂的牵制变成对 \(F\) 自身的约束呢？隔离问题：椭圆确定时，＂过定点＂如何充要替换？很自然地就能想到，这就等价于 \(|P F|+|P O|=2 a\) ，它显然又可以充要替换为 \(|P F|=1\) 。这个条件太漂亮了，它告诉我们 \(F\) 就是一个圆上的动点，然后从系统静态背景与动点 \(F\) 出发就能控制粗圆。 \(F\) 显然占据的主自由度数目为 1 。然后有了椭圆再取一个粗圆上的动点 \(Q\) 。这就是整个问题的要素关系网，总自由度数目为 2 ，关系网一路控制到底。要素关系网构建完成。

既然是两个自由度的问题，冻结法自然要登场了。冻结方式也很显然：先冻结第一层主自由度 \(F\) ，让第二层主自由度 \(Q\) 跑，把所求的最大值搞出来（只被 \(F\) 控制），再解冻 \(F\) 把最终的最值求出来即可。冻结 \(F\) 之后，除了 \(Q\)以外的全静止。隔离出问题就是：定椭圆，\(O\) 是一焦点，\(P\) 是椭圆上一定点，\(Q\) 是椭圆上一动点，求 \(|Q P|+|Q O|\) 的最大值。隔离出来就很明显了：这是很老套的小题，求法就是 \(|Q P|+|Q O|=|Q P|+2 a-|Q F| \leq 2 a+|P F|\) ，且等号可以成立（当且仅当 \(Q, P, F\) 共线）。再强调一下冻结法一定要保证每一步的最值都可以取到。因此接下来解冻 \(F\) ，我们只需要求 \(|P F|\) 的最大值，但这不是个定值吗？求都不需要求了，就是 1 。再加上已知的 \(a=2\) ，最后答案就是 5 。

做个总结。这道题的外壳看上去比较怪异，但是仍然万变不离其宗。梳理要素关系，选好操纵杆（主自由度），根据关系网进行隔离与元问题对应，用结构思维来解决问题，这始终适用于绝大部分题目的科学的思维方式。通过画图的方法思考与尝试，我们构建出控制顺畅，令人舒服的要素关系网。随后的解决问题在这个要素关系网之下就变得易如反掌。这道题的重点就是要素关系的梳理与构建，不把要素关系理清，那么做本题时就会抓不住方向，全靠瞎碰运气。未做出本题的同学，需要反思自己在构建要素关系网时，所做的思考与视频讲的心路历程的差距在哪里，并结合差距进一步完善自己的思维方式。

时间规划：

\section*{读题，思考画图方式： 1.5 min}
抓住另一个焦点作为主自由度的要点，建立要素关系网： 1.5 min\\
利用冻结法解决题目：2 min\\
总时间： 5 min

\section*{6．5．7 XXXX}
\section*{6.6 函数与导数}
\section*{6．6．1 遍历法思维练习}
视频日期：2024年5月5日

题目：\\
已知函数 \(f(x)=a x^{3}+b x^{2}+c x+d\) 存在两个极值点 \(x_{1}, x_{2}\left(x_{1}<x_{2}\right)\) ，且有 \(f\left(x_{1}\right)=-x_{1}, f\left(x_{2}\right)=x_{2}\) 。设 \(f(x)\) 的零点个数为 \(m\) ，方程 \(3 a(f(x))^{2}+2 b f(x)+c=0\) 的实根个数为 \(n\) ，则 \(m+n\) 的取值集合为 \(\qquad\) －

这道题目条件较为繁杂。按照惯例，我们先简单分析一下自由度，看一下问题的结构。 \(f(x)\) 里有四个参数，也就是四个自由度。然后已知的两个等式减去两个自由度，因此系统总自由度是 2 。同时，系统的已知条件全部都围绕着两个极值点 \(x_{1}, x_{2}\) 做文章，既有不等关系，给的牵制也全是关于它俩的。这自然就提示我们选 \(x_{1}, x_{2}\) 作为主自由度，作为操纵杆控制整个系统。从刚才的自由度分析来看，只要知道了 \(x_{1}, x_{2}\) ，就能定下 \(f(x)\) 。而且既然有两个极值点，那 \(f(x)\) 必定是三次函数了，\(a \neq 0\) 。

怎么定呢？这个元问题是要去刻画 \(f(x)\) 的形态，显然直接关系到问题的解决，因此可以先做一个思考。运用画图的方法，把极值点尝试标在坐标系里，我们立刻就会发现，极值点的位置分布对 \(f(x)\) 的形态有巨大的影响。因此，我们不得不使用遍历法来进行分析，以保证所有情况不遗漏。这是一个双自由度的遍历过程，因此冻结法的使用也是必然的。我们可以先做点思想试验以大致找找感觉。遍历通常是以极端情况起手，比如主自由度很小或者很大。比如我们先假设 \(x_{1}=-100\) ，这样可以再取 \(x_{2}=-50\) 。在坐标系上画出两点 \((-100,100),(-50,-50)\) ，很容易看出 \(f(x)\) 是增减增的模式，且有三个零点。这样啊看来想分析清楚也不是很难啊。这种情况无非就是标标点，画画图，已经可以预见其他情况也无非就是这样了。于是，我们对这个元问题的分析大致有了信心了。

在真正起手遍历之前，我们还得把所求做一个充要替换，以方便我们遍历时直接观测。第一个 \(m\) 是零点个数，可以直接读图像读出来。第二个有点论异。这显然是两层方程：首先求解 \(3 a t^{2}+2 b t+c=0\) ，解出全部 \(t\) ，然后再把所有 \(f(x)=t\) 的解取并集（注意是取并集，这考察大家基本的元问题掌握）。既然是两层，我们自然需要先把思维焦点放到二次方程 \(3 a t^{2}+2 b t+c=0\) 上。这是什么？隔离出来观察，不难看出这就是 \(f^{\prime}(t)=0\) ，因此从题目描述来看，它正好对应 \(x_{1}, x_{2}\) 这两个解，所以原方程就等价于 \(f(x)=x_{1}, x_{2}\) 这两个方程解的并集。

基本的充要替换做完了，现在我们可以开始冻结与遍历了。首先，冻结谁呢？两个极值点地位也没啥区别，所以先冻结谁都行。我们就先冻结 \(x_{1}\) 吧。然后遍历 \(x_{2} \in\left(x_{1},+\infty\right)\) 。那冻结的 \(x_{1}\) 的位置也对题目有影响，所以也得先对 \(x_{1}\) 做点讨论。

首先，如果 \(x_{1}\) 非常小，比如是－10000，－1000 这样的。点一下 \(x_{1}\) ，然后遍历 \(x_{2}\) 。 \(x_{2}\) 起手自然是 \(x_{1}\) 附近，这时不难发现，和之前我们的例子是类似的，\(f(x)\) 的零点数 \(m=3\) 。看 \(n\) 的话需要再画出 \(y=x_{1}, y=x_{2}\) 两条横截线。显然它们与 \(f(x)\) 交点数目分别为 1,2 ，且零点彼此不重复，于是 \(n=3\) ，此时 \(m+n=6\) ，得到一个答案。

继续遍历，增大 \(x_{2}\) ，不难发现只要 \(x_{2}<0\) ，刚才的讨论都仍然适用，此时 \(m+n=6\) 继续保持。当 \(x_{2}=0\) 时，问题出现了转机。这时 \(f(x)\) 只有两个零点了，\(m=2\) ，但是 \(n\) 仍然是 3 ，因此 \(m+n=5\) ，又一个答案。

继续增大 \(x_{2}\) 。如果 \(x_{2}\) 只大了一点点，比如只是 0.01 ，这时 \(f(x)\) 还是增减增模式，此时 \(m=1\) ，但仍然 \(n=3\) ，因此第三个答案：\(m+n=4\) 。并且显然：如果 \(x_{2}\) 仍然不大的话这个模式会继续保持

继续增大 \(x_{2}\) 。如果 \(x_{2}=-x_{1}\) ，这时两个极值平齐，这对三次函数显然不可能，淘汰掉再增大 \(x_{2}\) 的话，\(f(x)\)就要翻转成减增减模式了。此时不难看出，仍然是 \(m=1, n=3\) ，因此没有新增的答案。继续让 \(x_{2}\) 跑到正无穷，也仍然是这种模式。

这种 \(x_{1}\) 冻结情况下，\(x_{2}\) 的遍历完成了。因此我们可以考虑其他的 \(x_{1}\) 了。不难看出：不管取 \(x_{1}\) 是－10000，－100还是－1。它的 \(x_{2}\) 遍历模式都一样。因此下一步我们直接上 \(x_{1}=0\) ，这时必须 \(x_{2}>0\) 。遍历 \(x_{2}\) ，显然此时 \(f(x)\) 只能

是减增减模式，\(m=2\) 。再考虑两条横截线 \(y=0, y=x_{2}\) ，不难看出它俩各自与 \(f(x)\) 有两个交点，且这四个零点彼此不重复，因此 \(n=4\) ，此时 \(m+n=6\) ，仍然是之前有过的答案。不管怎么遍历 \(x_{2}>0\) ，此模式都不变。

下一步只能继续增大 \(x_{1}\) ，然后让 \(x_{2}>x_{1}\) 。此时 \(f(x)\) 还是减增减，但是有三个零点了，\(m=3\) 。再画两条横截线 \(y=x_{1}, y=x_{2}\) ，显然它们与 \(f(x)\) 分别横截于 3,2 个点，且彼此不重复，因此 \(n=5\) ，此时 \(m+n=8\) ，有了一个新答案。不难看出 \(x_{2}>x_{1}>0\) 全是这种模式。

综合一下，最终答案就是 \(\{4,5,6,8\}\) 。顺便说一下这道题即时检查的方式。由于这道题在画横线的时候可能会出现错误，而且除了画图之外没有什么其他解法，也不方便大家重算，因此建议带入特殊值来进行检查。具体来说，就是上面的那些情况，你可以每种惜况取一个 \(x_{1}, x_{2}\) 的特殊值进行静态问题分析，具体地把那两条横截线的数值写出来。另外，这道题的检查可以往后拖一拖，制造一定的时间差，从而降低刚才解题时思维惯性的影响。

这道题虽然可以排到高考真题压轴填空的难度，但想做出来并不困难。考察的就是基本的思维方式：冻结法和遍历法的运用。近期有不少同学在私信提问的时候都暴露出了遍历法意识不强的问题。遍历法是经验分享第二期里重点出现的，十分重要的思维方式，需要引起大家的重视。它是用来不重复不遗漏地刻画自由度的行为，保持了系统的信息量：在相当多的问题中都有广泛的应用。大家应当重视这样的思维，通过使用这种思维来强化自己充要替换的意识，提升自己思维的严谨性。

\section*{6．6．2 一道导数难题深度思路讲解}
题目\\
已知函数 \(f(x)=a x \ln x-x^{2}+(3-a) x+1(a \in R)\) ，若 \(f(x)\) 存在两个极值点 \(x_{1}, x_{2}\left(x_{1}<x_{2}\right)\) ，则当 \(\frac{x_{2}}{x_{1}}\) 取得最小值时，实数 \(a\) 的值为 \(\qquad\)。

读完题，我们很容易发现，这是一个单自由度的导数问题，极值点条件相当于一个不等条件（有几个零点，极值点都是限定 \(a\) 的范围的），然后自然的，\(a\) 严格控制了 \(x_{1}, x_{2}\) 虽然要素关系网如此简单，但是所研究的对象 \(\frac{x_{2}}{x_{1}}\)却让人感觉有点摸不着头脑。我知道 \(a\) 肯定严格控制它

\[
\text { 要索网 } a \rightarrow\left\{\begin{array}{l}
x_{1} \\
x_{2}
\end{array} \rightarrow \frac{x_{2}}{x_{1}}\right.
\]

也知道题目就是要看这个控制（也就相当于个函数）的最值，但是乍一看，很不好写表达式，好像很不好处理。

遇到乍一看无从下手的问题的时候，我们在经验分享第五期讲过了，要主动去思考＂我能干什么＂，不然你永远也前进不了。既然是研究极值点，那求导讨论肯定是少不了的了。我们先把必要的默写步骤做完。\\
求导：\(f^{\prime}(x)=a \ln x-2 x+3\) ，这玩意存在两个变号零点（极值点的基础充要替换），看上去并不难，只是一个基础题水准的玩意，稍加讨论即可。

显然 \(a \neq 0\) ，然后遍历一下 \(a\) 看看（这里一个小技巧是让参数到直线上去，因为直线的参数几何意义明确），因此讨论 \(f^{\prime}(x)\) 的变号零点转化为讨论 \(\ln x-\frac{2 x-3}{a}\) 的，也就是对数函数和直线 \(y=\frac{2 x-3}{a}\) 的交点情况。遍历一下 \(a\)（即转动直线，基础元问题），不难看出\\
（1）如果 \(a<0\) ，那永远不行\\
（2）如果 \(a>0\) ，那一定有两个变号零点（图形上很容易看出）\\
注意第（2）种情况如果想要严格证明的话，需要利用零点存在性定理作为基本逻辑，找点进行论证。我们之前有详细讲过 \(100 \%\) 稳收找点问题的思路方法，大家可以在＂数学的干货＂中找到，可以复习一下。（由于是选择题，这里我们就跳过）

因此，\(a>0\) 时候都有两个这样的 \(x_{1}<x_{2}\)（直线与对数曲线的交点），并且很容易在图像上观察到 \(x_{1}<1, x_{2}>\) \(\frac{3}{2}\) ，当然这只是一些很浅层次的观察，但是，一些好的直观住往有可能帮助我们提高信心，辅助解题

我们还可以转转 \(a\) 来看看 \(\frac{x_{2}}{x_{1}}\) 大致的变化趋势。显然当 \(a\) 很小（斜率 \(\frac{2}{a}\) 很大）的时候，直线几乎坚直，这时候 \(x_{2} \approx \frac{3}{2}\)（在 \(\left(\frac{3}{2}, 0\right)\) 的头顶），但是 \(x_{1} \approx 0\)（很小），因此 \(\frac{x_{2}}{x_{1}}\) 会非常大但是当 \(a\) 很大（斜率 \(\frac{2}{a}\) 很小）的时候，直线很平，几乎贴着 \(x\) 轴了，因此 \(x_{1} \approx 1\) ，但是 \(x_{2}\) 非常大，因此 \(\frac{x_{2}}{x_{1}}\) 还是非常大。这说明 \(a\) 处在两个极端都会导致 \(\frac{x_{2}}{x_{1}}\) 非常大，因此从直观上，它也得有个最小值，所以这个题目问的还是有意义的（最小值到底在哪呢）。这样的一些简单的遍历观察虽然对解题没有实质性的帮助，但是会给我们对于题目要素网的控制关系提供一个大致的直观，这样的直观对于增加我们做题的信心，以及帮助即时检查（利用常识判断计算是否有很离谱的错误）都有益处。至少我在做这道题的时候，会花 1 分钟左右做一个这样的思考与感受。

大致确定了 \(x_{1}, x_{2}\) 的几何直观后，最根本的问题来了——怎么算？我们需要的，是利用一个主自由度去控制 \(\frac{x_{2}}{x_{1}}\) 。我们首先梳理我们有的已知（可以列式的关系），也就是两个极值点的控制式子：\(a \ln x_{1}-2 x_{1}+3=0\) 和 \(a \ln x_{2}-2 x_{2}+3=0\) 。这俩式子配合上我们刚才的几何直观（相当于范围的约束），就是我们已有的全部信息的充要替换，我们需要从这两个式子去控制 \(\frac{x_{2}}{x_{1}}\) 。

这里就涉及到一个做难题非常普遍的基本的逻辑技巧：如何选取主自由度？本题思维上最舒适的主自由度就是 \(a\) ，但是想直接通过上面的两个式子，就把 \(\frac{x_{2}}{x_{1}}\) 用 \(a\) 控制出来，看上去不太现实。所以，选取 \(a\) 作为计算用的自由度，并不合适。我们需要尝试把 \(\frac{x_{2}}{x_{1}}\) 用某个东西控制出来，这个＂某个东西＂（也就是计算用的主自由度）只要方便考虑其范围，就可以。下面就要尝试了。既然直接顺者做不好想，我们就要有一定的探索精神，也就是第五期讲到的试错。我们需要尝试去凑 \(\frac{x_{2}}{x_{1}}\) 这个东西才行。那怎么凑呢？观察上面那两个式子，似乎到的试错。我们需要尝试去凑 \(\frac{x_{2}}{x_{1}}\) 这个东西才行。那怎么凑呢？观察上面那两个式子，似乎相减看上去好点，减一下试试： \(a \ln \frac{x_{2}}{x_{1}}=2\left(x_{2}-x_{1}\right)\) ，虽说出来了需要的好东西，但又冒出来别的东西：\(x_{2}-x_{1}\) ，又需要用别的东西去控制它，想了想，似乎也没啥实质性的简化，工作量没有实质减少．．．．．．

肯定是信息没用全，没用好！事实上，我们刚才的变形确实流失了营养——实际自由度 1 个，但我减出来的式子是 3－1＝2，原本的两个式子里的营养到了下面这个，出现了信息流失！所以如果想继续往下走的话，肯定还要继续利用原始的两个式子，无论如何要保证信息不流失才行。

想到这个，接下来要试的路就很消楚了： 2 个式子 3 个参数， 1 个自由度，想保证信息不流失，那就需要不仅能建立 \(x_{1}\) 和 \(x_{2}\) 的联系，还要保证我得到的式子实际自由度是1。那怎么做能实现刚才的目标呢？（思维焦点放过来）显然是消元！把 2 个式子， 3 个参数，变成 1 个式子， 2 个参数，消去 \(a\) ！

实际操作一下，就是 \(a=\frac{2 x_{1}-3}{\ln x_{1}}=\frac{2 x_{2}-3}{\ln x_{2}}\) ，然后取后面一个等号（后文中大家将看到，这是可以走上正道的变形）。

战术分析：为什么会想到这样尝试\\
因为我们知道肯定要从上述两个式子出发去完成我们的控制，但是刚才的尝试不太成功，反思发现是信息流失，自由度没控好。所以改善的方向就是：在凑 \(\frac{x_{2}}{x_{1}}\) 的时候，保持实际自由度是 1 。要达到这个目标，需要先优先后者，因为凑 \(\frac{x_{2}}{x_{1}}\) 肯定是个麻烦的技术活，不太像能一眼看出来的样子，所以战术上要在控好自由度，保持信息不流失（充要替换）的基础上，去尝试配凑。实际上，既然你都信息不流失了，那逻辑上讲，你变形后的式子必定能够达到目标，所以为何不大胆一点去这么做呢？

另外，这一步把原本我们看着舒服的主自由度 \(a\) 给丢掉了，乍一看似乎不太舒服。但是，除了刚才自由度和充要替换的逻辑分析保障以外，希望有意冲击超高分的同学记住，动静视角和要素网是一个理念，不是一个照本宣科的东西。我刚才的转化虽然把 \(a\) 暂时丢掉了，但是我始终清楚，\(a\) 在控制 \(x_{1}, x_{2}\) ，并且控制方式（小元问题）就是我们原始的两个式子。只要这个控制一直在，那我完全可以在用某一个＂主自由度＂控制清楚 \(\frac{x_{2}}{x_{1}}\) 后，及时撤回来，完成 \(a\) 的研究。这样的理念加持下，我们得以不断向前开拓，又不忘初心，随时能拿来控制关系来处理问题。

回到解题中，刚才得到了 \(\frac{2 x_{1}-3}{\ln x_{1}}=\frac{2 x_{2}-3}{\ln x_{2}}\) ，这个式子左右形式高度一致。能从这玩意出发变形，实现我们的目标吗？我相信能，因为一个小经验：这种具有高度对称，一致的结果的构，一般都能通过变形得到一些好东西出来。

尝试一下，肯定是需要先通分的：\(\left(2 x_{1}-3\right) \ln x_{2}=\left(2 x_{2}-3\right) \ln x_{1}\) ，我一眼就看到两侧有一组东西可以直接合并： \(3 \ln x_{2}-3 \ln x_{1}=3 \ln \frac{x_{2}}{x_{1}}\) ，这是好东西，留着。

剩下一个 \(2 x_{1} \ln x_{2}-2 x_{2} \ln x_{1}\) ，这个乍一看没刚才那么直白（如果是 \(2 x_{1} \ln x_{2}-2 x_{2} \ln x_{2}\) 这种，就又可以直接往外提了），但是固有的一致结构还在，应该还有操作余地。因此，我们试着强行拆出一组 \(\frac{x_{2}}{x_{1}}\) 来，也就是如刚才括号中的，直接配一个出来试试 \(2 x_{1} \ln x_{2}-2 x_{2} \ln x_{1}=2 x_{1} \ln x_{2}-2 x_{2} \ln x_{2}+2 x_{2} \ln x_{2}-2 x_{2} \ln x_{1}\) 。然后右边分组合并一下，就变成了 \(2\left(x_{1}-x_{2}\right) \ln x_{2}+2 x_{2} \ln \frac{x_{2}}{x_{1}}\) ，因此得到 \(2\left(x_{1}-x_{2}\right) \ln x_{2}+2 x_{2} \ln \frac{x_{2}}{x_{1}}=3 \ln \frac{x_{2}}{x_{1}}\) 。把我们喜欢的，和 \(\frac{x_{2}}{x_{1}}\) 直接挂钩的东西放在一起，得到 \(\left(2 x_{2}-3\right) \ln \frac{x_{2}}{x_{1}}=2\left(x_{2}-x_{1}\right) \ln x_{2}\)

接下来是本题的一大难点：这个式子再怎么走？不要忘了我们做题的理念，通过一个主自由度去尝试控制 \(\frac{x_{2}}{x_{1}}\) 。我们刚才一直在围着 \(\frac{x_{2}}{x_{1}}\) 配凑，现在来试者观察一下，这个式子里面有没有那种，有希望做主自由度的变量？也就是某个东西，让整个式子里面只有它和 \(\frac{x_{2}}{x_{1}}\) ？这样想的话，我们立刻发现，这个式子里面，\(x_{2}\) 高强度出现！！只有一个 \(x_{1}\) 孤零零在边上，除此之外被 \(x_{2}\) 填满。难道，可以选 \(x_{2}\) 作为主自由度去尝试控制 \(\frac{x_{2}}{x_{1}}\) ？

还有个 \(x_{1}\) 碍事！怎么办？（思维焦点放过来）能不能把 \(x_{1}\) 用主自由度和被控制项表示出来，让整个式子里只有 \(x_{2}\) 和 \(\frac{x_{2}}{x_{1}}\) ？那显然，把这两个东西除一下，就是 \(x_{1}\) 。因此，我们继续变形（尝试把 \(x_{2}\) 放到一边，\(\frac{x_{2}}{x}\) 放到一边，让整个式子只有这俩玩意）

\[
\left(2 x_{2}-3\right) \ln \frac{x_{2}}{x_{1}}=2\left(x_{2}-x_{1}\right) \ln x_{2}=2 x_{2}\left(1-\frac{x_{1}}{x_{2}}\right) \ln x_{2}, \text { 因此 } \frac{2 x_{2}-3}{2 x_{2} \ln x_{2}}=\frac{1-\frac{x_{1}}{x_{2}}}{\ln \frac{x_{2}}{x_{1}}}
\]

这个式子就基本实现了我们的目标：首先 \(x_{2}\) 作为主自由度，控制了式子左侧，然后利用这个等号得到了关于 \(\frac{x_{2}}{x_{1}}\)的方程，通过这个方程就控制了 \(\frac{x_{2}}{x_{1}}\)（当然，说是＂基本实现＂，是因为方程可能存在多解或无解，可能需要取舍，但不管怎么说，我们已经前进了一大步）

既然控制方式找到了，那我们就换个元，方便叙述。设 \(t=\frac{x_{2}}{x_{1}}\) ，然后把上面的控制用函数严格写出来：设

\[
g(x)=\frac{2 x-3}{2 x \ln x}, h(x)=\frac{1-\frac{1}{x}}{\ln x}=\frac{x-1}{x \ln x},
\]

于是控制就变成了 \(g\left(x_{2}\right)=h(t)\) 我们最终要研究的是 \(t\) 的最小值。如果控制关系是 \(g\left(x_{2}\right)=t\) ，那舒服了，直接看 \(g\) 在定义域上的最小值就行了，但是事实是 \(t\) 也套了一层函数，所以我们得先看看 \(h\) 是个什么情况求导： \(h^{\prime}(x)=\frac{\ln x-x+1}{(x \ln x)^{2}}, 130\) 分以上的同学应该熟知 \(\ln x-x+1 \leq 0\) ，所以 \(h\) 在定义如内始终递减。想起之前画图直接看到的 \(x_{1}<1, x_{2}>\frac{3}{2}\) ，因此无论如何都有 \(t>\frac{3}{2}\) ，所以想要 \(t\) 最小，只需要看 \(h(t)\) 哈时候最大就行了。

那就是看 \(g(x)\) 哈时候最大了呗！求导：

\[
g^{\prime}(x)=\frac{3 \ln x-2 x+3}{2(x \ln x)^{2}}
\]

令 \(\varphi(x)=3 \ln x-2 x+3\) ，于是 \(\varphi^{\prime}(x)=\frac{3}{x}-2\) 。之前说过 \(x_{2}>\frac{3}{2}\) ，所以显然 \(\varphi(x)\) 单减，而 \(\varphi\left(\frac{3}{2}\right)=3 \ln \frac{3}{2}>0\) ，而且从量级上看（之前讲找点时候的＂望气＂），\(x_{2}\) 很大的时候肯定 \(\varphi\left(x_{2}\right)<0\) ，所以中间有个零点，记为 \(x^{*}\)（隐零点，静止）。所以很显然，\(g(x)\) 在 \(\left(\frac{3}{2}, x^{*}\right)\) 上增，在 \(\left(x^{*},+\infty\right)\) 上减，最大值就是在 \(x^{*}\) 处取的。（实际上，这个 \(x^{*}\) 确实能决定一个 \(t>\frac{3}{2}\) ，留作练习）。

现在我们得到了最大值的刻画： \(3 \ln x_{2}-2 x_{2}+3=0\) ，它可以完全定死 \(x_{2}>\frac{3}{2}\) 。控制已经完全了，目前的思维焦点只需要放在通过这个式子求 \(a\) 了。我们现在返回原始的控制关系：\(a=\frac{2 x_{2}-3}{\ln x_{2}}=3\) ，于是游戏结束！

用时分析（面向有志于考场拿下本题的同学）：

\begin{itemize}
  \item 读题＋要素网分析：1min
  \item 求导看出 \(a, x_{1}, x_{2}\) 的大致范围： 1 min
  \item 第一次尝试变形 + 失败总结 + 战术分析： 2 min
  \item 对 \(\frac{2 x_{1}-3}{\ln x}=\frac{2 x_{2}-3}{\ln x}\) 变形得到 \(g\left(x_{2}\right)=h(t): 2 \mathrm{~min}\)
  \item 余下部分： 2 min
  \item 合计： 8 min\\
（注：因为不需要写步骤，所以时间规划比较短）
\end{itemize}

\section*{总结：}
本题难度大，前期找出控制关系的每一步操作都不是非常明显的，但是几乎每一步的目的都是逻辑思维指引的结果。不仅要具备这样的思维方式（要素网，充要替换），还需要在这样思维方式的笼罩之下，去主动观察，思考，其中一个典型便是看出选取 \(x_{2}\) 作为计算用的主自由度是最简单的。当然，如果没有这种主自由度和控制的意识的话，即便列出了那个式子，也有可能与正道擦肩而过。

徐此之外，在逻辑告诉我们 \(\frac{2 x_{1}-3}{\ln x_{1}}=\frac{2 x_{2}-3}{\ln x_{2}}\) 最有希望能走上正道之后，还需要我们顽强的将这个目标推进下去。本题体现的最重要的代数变形技巧就是利用这种两侧一致的对称结构去配凑需要的东西，需要大家坚定的信心以及基本的代数变形功底。而且，本题的操作容量也不小，从最开始的讨论 \(a\) 的范围，形成各个参量

的几何直观，到代数变形，到后期两个函数的形态讨论，无不需要我们仔细操作，踏实完成（不过至少，当我们列出了控制关系式 \(g\left(x_{2}\right)=h(t)\) 之后，就应该能看出来本题已经迈过了难点，剩下的就是基本的元问题解决了，这样的逻辑认识能带给我们信心）。

综上所述，本题无论是逻辑思维，还是代数变形，还是操作容量，以及遇到困难，面对挫折的意志，都对大家提出了不小的要求，是一道不折不扣的难题。解题之中，需要大家始终沿着我们经验分享 1－5 期的思维方式前行，逢山开路，遇水架桥，才能最终解出。

本题是一位粉丝提问的某模拟题上的单选压轴（我把选项省去了），个人认为难度过大，高考不太可能考这么难，但是相信 130 分以上的同学都可以从中吸收到不少有价值的东西。

\section*{课后练习：}
1．将第一页中（2）的找点过程补全（当作大题练习）\\
2．第六页中提到：＂实际上，这个 \(x^{*}\) 确实能决定一个 \(t>\frac{3}{2}\) ，请严格说清楚这一点，确保这个 \(x^{*}\) 诱导的方程 \(h(t)=g\left(x^{*}\right)\) 确实有解，而不会出现 \(g\left(x^{*}\right)\) 跳出 \(h\) 值域的情况

3．另一个思考维度：还是刚才的题目 2 ，从这个 \(x^{*}\) 反向确定的 \(a\) 沿着要素网所控制的 \(t\) 是否必须是上述方程的解？你认为逻辑上是否足以说明上述方程有解？因此，可不可以说，逻辑已经保证了 \(g\left(x^{*}\right)\) 绝对不会跳出 \(h\) 的值域？提示：要素网已经存在，这是既定的客观事实\\
4．第四页中，我们通过 \(2 x_{1} \ln x_{2}-2 x_{2} \ln x_{1}=2 x_{1} \ln x_{2}-2 x_{2} \ln x_{2}+2 x_{2} \ln x_{2}-2 x_{2} \ln x_{1}\) 来进行配奏，其中配了一个 \(2 x_{2} \ln x_{2}\)（也就是用下角标 2 来配凑），如果我不选择 2 ，而选择 1 （也就是使用 \(2 x_{1} \ln x_{1}\) 进行配凑），会怎么样呢？请完成相关讨论，并在讨论完之后，把这种做法和视频中的方法进行类比

\section*{6．6．3 赋值推演：抽象函数的机构思维}
\section*{题目（不定项选择）：}
\(f(x)\) 是定义在 \(R\) 上的函数，对 \(\forall x, y \in R\) 有 \(f(x+y)-f(x-y)=f(x+2) f(y+2)\) 且 \(f(0) \neq 0\) ，则下列说法正确的有（）

A．\(f(x)-f(-x)=0\)

B．\(f(x+4)-f(x)=0\)

C．\(f(0)+f(2)+f(4)+\cdots+f(2024)=-2\)

D．\(f^{2}(1)+f^{2}(2)+f^{2}(3)+\cdots f^{2}(2024)=4048\)

这道题目是典型的抽象函数问题。给出一个抽象函数的表达式，让我们推算一些函数性质与函数值的属性。对于抽象函数问题，基本的操作手段就是赋值。凡是涉及到考察抽象函数的性质与函数值的问题，都必然离不开赋值法。但是，如何赋值，赋什么值，赋完值之后如何利用这些结果去推算，都是这类题目的难点。凡是操作就离不开行为理念，换而言之，就是思维方式。我们在数列经验分享中为大家分析过函数与数列问题的自由度。它们的自由度是无限的，但是利用控制的理念就能达到由少及多，从一小部分出发控制全局的效果。

在数列经验分享下篇中，我们为大家着重讲解了抽象函数的解题理念。对于研究函数性质的选项（奇偶性，周期性等函数形态），我们的赋值需要尝试实现相应的函数形态控制（如奇偶性是一半的定义域控制另一半，周期性是线段上的函数值递推到全体）。而对于具体函数值的问题，要么就是单纯求一个函数值，要么就是求很多个函数值的组合，后者常常都以数列的形式给出。此时的赋值过程要遵循的原则就是看你赋值出来的式子具有什么控制效果，能不能利用你的赋值式加上递推来求出所有想要的东西。以这道题为例，我们再一次给大家强化一下这类题目的解题策略。

首先看 A 选项，显然这等价于偶函数，因此我们赋值就希望让函数值里出现一正一负。对照题干，函数 \(f\) 下的括号里有 \(x+y, x-y, x+2, y+2\) 四部分。其实结合我们刚才的理念，大家已经可以看出来了，后两个就是拿来做我们上述的第二种递推的。而前两个混搭着 \(x\) 与 \(y\) ，还又加又减，可操作的余地很强。怎么赋值能搞出一正一负？显然取个 \(x=0\) 就造出来 A 选项需要的式子 \(f(y)-f(-y)=f(2) f(y+2)\) 了。这里要注意：研究一个问题需要的赋值未必能一步到位，我们需要结合赋值不断进行充要替换与抛弃来简化问题。这需要我们不断的切换思维焦点，寻找下一步能干的事情。A 选项显然等价于对任意 \(y\) ，都有 \(f(2) f(y+2)=0\) ，这个又等价于要么 \(f\) 恒为 0 ，要么 \(f(2)=0\) 。但是已知 \(f(0) \neq 0\) 就直接给第一种情况判了死刑，因此原命题就等价于 \(f(2)=0\) 。尝试证明或否定它显然就是下一步的目标（思维焦点转移过来）。既然是求一个具体函数值，那么就来试试直接把 2 搞出来吧。正好后两个都是有 +2 的，于是直接令 \(x=y=0\) ，得到 \(f(0)-f(0)=[f(2)]^{2}\) ，显然这就证明了 \(f(2)=0\) ，于是 A 正确。

对于 B 我们也是同理。现在需要看相差 4 的函数值之间的递推了。怎么赋值能让两个括号内的东西差出 4来？显然是 \(y=2\) ，得到 \(f(x+2)-f(x-2)=f(x+2) f(4)\) 。同理，B 选项就等价于 \(f(4)=0\) 。那么，怎么研究 \(f(4)\)呢？还是跟刚才一样，想要造出 4 那就让 \(x+2=y+2=4\) 即可，代入得到 \(f(4)-f(0)=[f(4)]^{2}\) 。这么看的话，正好 \(f(4)\) 与已知中的 \(f(0)\) 之间建立了严格的控制关系。既然 \(f(0) \neq 0\) ，因此自然可得到 \(f(4) \neq 0,1\) 。因此不仅 \(f(4) \neq 0\) ，而且 \(f\) 也不可能恒为 0 了，这就排除了 B 选项。

提问这道题的同学说，当时想着非要把周期给搞出来，其实这是完全没必要的，因为 C 和 D 研究的是数列求和问题。这类问题的常见思维就是找出所求各项之问的递推关系，再配合一下赋值得到的初始值来用数列技术求解。我们目前得到 \(f(2)=0\) 以及 \(f(4) \neq 0,1\) ，而 C 选项里，\(x\) 的值是公差为 2 的，因此我们最好能找到一个 \(f(x)\)控制 \(f(x+2)\) 的方式，或者至少得是个类似的控制。而刚才我们得到过一个 \(f(x+2)-f(x-2)=f(x+2) f(4)\) ，改写一下就是 \((1-f(4)) f(x+2)=f(x-2)\) 。而 \(1-f(4)\) 是个非零常数，于是这个式子本质上是一个用长度 4 的线段控制整个 \(R\) 上函数值的式子。这其实距离我们刚才的目标已经迈进一大步了。我们要研究 \(2,4,6 \cdots\) 的函数值，现在跳 4 步的控制实现了， 2 处的函数值也知道，如果我再知道一个 4 处的函数值，这个控制不就成立了吗？大家注意，有的同学可能会问，那我非要找出一个跳 2 步的控制行不行呢？我们在逻辑上分析一下这两条路对应的策略。即便你真的找到一个跳 2 步的控制，你也可以用它来由 2 算出 4 处的值。这说明这条路一样可以得到 \(f(4)\)的值，于是也就是说 \(f(4)\) 的值肯定是可以推出来的。既然这样的话，是强行再找一个递推简单呢？还是再直接算一个具体函数值简单呢？当然是后者直观上更加容易。因此，这里选择计算 \(f(4)\) 的策略，是最容易的。想出这种策略的思想依据，正是我们前面说到的控制。那么，现在我们有一个跳 4 步的递推，这时候 2 没法跳到 4 来，于是还得另外考虑 \(f(2)=0\) 以外的信息。大家注意：凡是你再赋值过程中得到的结果，最好都不要轻易忽视。刚才我们除了算出过 \(f(2)=0\) 之外，还得到过什么结果呢？是不是还有 A 选项的偶函数啊？结合当前目标考量一下，这不就是说 \(f(4)=f(-4)\) 吗？正好，我们刚才得到了一个跳 4 步的递推关系式（1－f（4））\(f(x+2)=f(x-2)\) ，所谓的＂公比＂也是被 \(f(4)\) 控制的。我从 \(f(4)\) 出发，往 \(x\) 轴负方向反跳两步就得到了 \(f(-4)\) ，正好用上这二者的关系！

因此，由 \((1-f(4)) f(4)=f(0),(1-f(4)) f(0)=f(-4)\) 以及 \(f(4)=f(-4)\) 可得 \(f(4)=(1-f(4))^{2} f(4)\) 。解这个方程，并舍去 0 的根，得到 \(f(4)=2\) 。于是递推关系也就彻底地定下来了：\(f(x+2)=-f(x-2)\) 。配合 \(f(2), f(4)\)的值，我们很容易从 \(f(0)=-2\) 一直递推到 C 选项中的 \(f(2024)=-2\) 。自变量从 0 开始，不断跳 2 ，函数值将经历 253 组 \(-2,0,2,0\) 循环后，到了末项 2024 重新变成 -2 ，因此求和也是 \(-2, \mathrm{C}\) 正确。

最后来看 D ，它说的是一堆平方。自变量公差变为 1 了，更加难以处理，因此我们考虑从平方入手。原始式子能不能搞出个平方来？当然可以，让 \(x, y\) 相等即可，于是我们赋值得到 \(f(2 x)-f(0)=[f(x+2)]^{2}\) ，重新整理得到 \([f(x)]^{2}=f(2 x-4)+2=2-f(2 x)\) 。故所求即 \((2-f(2))+(2-f(4))+\cdots+(2-f(4048))=4048-f(2)-f(4)-\cdots f(4048)\)。它的求法就和 C 项一样了，不难得出 \(f(2)+f(4)+\cdots f(4048)=506(-2+0+2+0)=0\) ，因此 D 正确。于是最后答案就是 ACD 。

做个总结。这道题在高考题中肯定属于多选题压轴的水平，但是可以看到，其解题理念完全没有脱离我们数列经验分享中的控制思想。抽象函数问题在研究单个函数值或者函数性质时经常是单纯依靠凑自变量的形式来解决。但是涉及到以周期性为代表的递推性质，或者是与递推性质紧密相关的数列结论时，根据我们经验分享里说到的，递推的控制结构才是核心。比如我有了一个 \(f(x)\) 与 \(f(x+4)\) 之间的等式，那么这时候重要的是认识到这个函数只需要长度 4 的区间就能控制整个 \(R\) 上的函数值。周期性啥的，数列啥的都是这个递推结构的自然衍生产物。递推结构这个事实是最重要最基本的。面对所求，只要用好数列的基本功与递推关系，就能稳稳拿下整道题目。不管做什么题目，up 主都在强调，控制与结构才是核心。搞清楚了题目的结构，针对结构选取恰当的知识与方法（元问题），问题自然就迎刃而解了。

\section*{时间规划参考}
观察 A 选项，决定赋值 \(x=0\) ，看出关键在于 \(f(2): 1 \mathrm{~min}\)\\
摆正思维焦点，赋值计算 \(f(2)\) ，判断 A 选项： 1 min\\
观察 B 选项，决定赋值 \(x=2\) ，看出关键在于 \(f(4): 0.5 \mathrm{~min}\)（有 A 的经验，速度加快）\\
摆正思维焦点，赋值判断 \(f(4)\) ，排除 B 选项： 1 min\\
认识到 C 选项本质是公差 2 的递推，以及 B 的讨论蕴含公差 4 的递推： 1 min\\
利用 B 中的递推与 A 的提示，求解 \(f(4): 1 \mathrm{~min}\)\\
计算判断 C 的正确性： 1.5 min\\
观察 D，并利用赋值把 D 转化为一次式： 1.5 min\\
计算判断 D 的正确性： 1 min 总时间： \(9-10 \mathrm{~min}\)

\section*{6．6．4 过于固定的思维方式}
一道高一难题稳定破解思路

题目：\\
设函数 \(f(x)=\left|x+\frac{4}{x}-a x-b\right|\) ，若对任意实数 \(a, b\) ，总存在 \(x_{0} \in[1,3]\) 使得 \(f\left(x_{0}\right) \geq m\) 成立，则实数 \(m\) 的最大值为 \(\qquad\) －

这是一道典型的封装命题类问题，不过其约束变元多，充要替换有一定难度，导致本题的整体难度比较高，

放在高一年级绝对是一道顶级难度的题目。不过，如果能较好地掌握我讲过的封装命题处理思维，解决这道题还是不会太依赖灵感和智商的。

所谓＂封装命题＂，就是某些自由度被量词（任意，存在）约束，导致该自由度所描述的主要矛盾不再是自身。例如，＂任意实数 \(x\) 的平方都是非负的＂，这里的 \(x\) 作为一个自由度遍历整个实数，但是因为前面加了＂任意＂这个量词，导致这个命题的主要矛盾不再是这个自由度 \(x\) ，而是整个实数轴这个静态对象，于是这个命题就成为了一个静态命题，而不是被动态的 \(x\) 控制的命题了。

关于封装命题类问题的处理方法，我们在＂高中数学经验分享番外篇＂中有着详细的讲解，建议在观看过那一期视频后再来学习本题。我们在番外篇中给出了封装命题的统一思维方式，就是冻结，遍历，设问的大致思维路线。当出现（多个）被封装的约束变元时，应当按照它们在要素关系网中出现的顺序，先冻结靠前的，研究靠后出现的约束变元，通过遍历，设问加充要替换的思维方式去尝试消去约束变元，在重复上述步骤直至消去所有约束变元。

我们来看本题。题目一共出现了三个约束变元 \(a, b, x_{0}\) ，其中 \(a, b\) 是同时出现的，都加的＂任意＂约束，所以看成 \(\forall a, \forall b\) 和 \(\forall b, \forall a\) 都可以。但是 \(x_{0}\) 肯定是最后出现的，于是我们先冻结（随便固定一组）\(a, b\) ，然后来看这时的命题：存在 \(x_{0} \in[1,3]\) 使得 \(f\left(x_{0}\right) \geq m\) 。这个命题的充要替换是非常基础的，都讲烂了，连遍历和设问都不值得去使用了。其等价于＂\(f(x)_{\max } \geq m "\) ，当然，要限制定义域为［1，3］。于是最后一层封装 \(x_{0}\) 就被我们消去了。

我们把思维焦点从 \(x_{0}\) 那里离开，继续看剩下的两个约束变元 \(a, b\) 。刚才说到，把它们看成 \(\forall a, \forall b\) 和 \(\forall b, \forall a\) 都可以，因此下一步的冻结既可以冻结 \(a\) ，也可以冻结 \(b\) ，有两条道路可以选了！因此，我们需要进行预判与选择。如果我们冻结 \(b\) ，去看 \(a\) ，那么问题就是：对任意 \(a\) 都有 \(\left|x+\frac{4}{x}-a x-b\right|_{\max } \geq m\) 。这是个什么问题呢？很明显，我们分出两个函数 \(g(x)=x+\frac{4}{x}\) 和 \(h(x)=a x+b\) ，其中 \(g(x)\) 是静止的一段对号函数，\(h(x)\)（在冻结 \(b\) 活动 \(a\) 时）是一个过 \(y\) 轴定点 \((0, b)\) 的旋转直线。我们要研究的是过这个定点的旋转和一段对号函数的最大差距，这个元问题好研究吗？\(|g(x)-h(x)|\) 本身作为两个函数的坚直高度差，这个东西要用 \(a\) 带来的旋转去控制，这本身感觉上就不搭边，预判的体感肯定不好研究。如果我们陈结 \(a\) ，去看 \(b\) 呢？刚才说到 \(|g(x)-h(x)|\) 是坚直高度差，正好 \(b\) 带来的控制就是坚直移动！这个感觉是对的，用操纵坚直移动的主自由度去研究坚直高度差，比刚才用 \(a\) 控制的预判体感要好很多。所以就决定接下来尝试先冻结 \(a\) ，去研究 \(b\) ，试着把 \(b\) 也消去。

现在的问题是：任意 \(b\) ，都有

\[
\left|x+\frac{4}{x}-a x-b\right|_{\max } \geq m
\]

这个问题虽然预判上的方向是正确的，但是似乎仍然有些模糊。大致的充要替换方向是有的：冻结 \(a\) ，遍历整个 \(b \in \mathbb{R}\) ，那么上面的 \(\left|x+\frac{4}{x}-a x-b\right|_{\max }\) 就只受 \(b\) 控制，相当于一个关于 \(b\) 的函数，记为 \(\varphi(b)\) 。我们希望这个 \(\varphi(b) \geq m\) 恒成立，也就是希望 \(\varphi(b)_{\min } \geq m\) ，所以关键就是搞清怎么研究 \(\varphi(b)_{\min }\) 。

思维仍然模糊的时候，我们之前说过，可以尝试思想试验（参见＂数学的干货＂视频列表）。我们既然冻结了 \(a\) ，那就做个试验，比如我取个 \(a=1\)（你取别的都行），先尝试研究一个

\[
\varphi(b)=\left|\frac{4}{x}-b\right|_{\max }
\]

这就是一段递减函数的上下平移，我们尝试动一动 \(b\) 看看。如果 \(b=100000\) ，那么整个反比例函数特别低，绝对值翻上去之后很高，\(\varphi(b)\) 很大；而如果 \(b=-100000\) ，同理，反比例函数特别高，绝对值直接去掉也很高，\(\varphi(b)\) 很大。所以 \(b\) 很大或者很小的时候肯定取不到 \(\varphi(b)_{\min }\) 。继续移动 \(b\) 的滑块，如果移动到 \(b=\frac{4}{3}\) ，减函数段刚好触碰到\\
\(x\) 轴地面，这时候 \(\varphi(b)\) 刚好就是这个减函数段本身的高低差 \(\left(4-\frac{4}{3}=\frac{8}{3}\right)\) 。继续变大 \(b\) ，减函数段有一部分下沉到 \(x\) 轴以下，这一部分通过绝对值翻到 \(x\) 轴上方。你会发现，此时 \(\varphi\)（ \(b\) ）继续变小（比减函数段本身的高低差 \(\frac{8}{3}\) 还要小了）。如果只是下沉下去一点，翻上来一小段，那么 \(\varphi(b)\) 也就比 \(\frac{8}{3}\) 小一点；如果下沉到正好 \(x\) 轴上方下方各一半，折上去之后 \(\varphi(b)\) 刚好就是高低差 \(\frac{8}{3}\) 的一半，也就是 \(\frac{4}{3}\) ；继续下沉的话，\(x\) 轴下方的部分就很多了，翻上去也会很高，会比 \(\frac{4}{3}\) 还高；再下沉下去的话，整个减函数段就全到 \(X\) 轴下方了，翻上去之后最大值就不会再小于高低差 \(\frac{8}{3}\) 了。

刚才的思想试验能带给我们什么呢？我们之前说过，做思想试验的过程本身不是目的，真正的目的是通过特殊值的试验明白背后的一般规律。刚才 \(\varphi(b)\) 取最小值是怎么过去的？首先得有一个高低差，然后需要适度上下平移来把高和低折半，这样才是 \(\varphi(b)\) 最小值。对于一般的函数，高低差是什么？当然是最大值和最小值的差啊！所以我们就明白了一个事情：\(\varphi(b)_{\min }\) 其实就是把 \(y=\frac{4}{x}+x-a x\) 在［1，3］上的最大值和最小值拿出来，作差（就相当于刚才的 \(\frac{8}{3}\) ），再除以 2 。 \(b\) 带来的上下平移不会影响高低差，所以这玩意和 \(b\) 本身无关。这样，我们就得到了一个原命题的充要替换：\(y=\frac{4}{x}+x-a x\) 的高低差 \(\geq 2 m\) 。预判一下，高低差就是求最大值和最小值，对于含参数的对号函数而言并不困难，因此预判可行。于是现在 \(b\) 也被我们消去了。

现在已经消去了两个约束变元，问题变成：令 \(\lambda(x)=\frac{4}{x}+x-a x, x \in[1,3]\) ，对任意 \(a\) 都有 \(\lambda(x)_{\max }-\lambda(x)_{\text {min }} \geq 2 m\)。注意这里 \(\lambda(x)=\frac{4}{x}+(1-a) x, a\) 遍历 \(\mathbb{R}\) 就相当于 \(1-a\) 遍历 \(\mathbb{R}\) ，所以我们不如直接就当 \(\lambda(x)=\frac{4}{x}+a x\) ，记号还简单。下面只需要求 \(\lambda(x)\) 在 \([1,3]\) 上的两个最值即可。

当 \(a \leq 0\) 时 \(\lambda(x)\) 递减；\(a>0\) 时 \(\lambda(x)\) 在 \(\left(0, \frac{2}{\sqrt{a}}\right)\) 上减，\(\left(\frac{2}{\sqrt{a}},+\infty\right)\) 上增，因此 \(\frac{2}{\sqrt{a}} \geq 3\) ，即 \(0<a \leq \frac{4}{9}\) 时 \(\lambda(x)\) 也在［1，3］上减，这时

\[
\lambda(x)_{\max }=\lambda(1)=a+4, \lambda(x)_{\min }=\lambda(3)=3 a+\frac{4}{3}
\]

于是高低差为 \(\frac{8}{3}-2 a\) 。当 \(\frac{2}{\sqrt{a}} \leq 1\) 即 \(a \geq 4\) 时高低差反过来，是 \(2 a-\frac{8}{3} 。 \frac{4}{9}<a<4\) 时，显然

\[
\lambda(x)_{\min }=\lambda\left(\frac{2}{\sqrt{a}}\right)=4 \sqrt{a}
\]

但是最大值要从 \(\lambda(1)\) 和 \(\lambda(3)\) 中选，不难看出 \(\frac{4}{9}<a \leq \frac{4}{3}\) 时 \(\lambda(1)\) 更大，此时高低差为 \(a+4-4 \sqrt{a}=(\sqrt{a}-2)^{2} ; \frac{4}{3} \leq\) \(a<4\) 时 \(\lambda\)（3）更大，此时高低差为

\[
3 a+\frac{4}{3}-4 \sqrt{a}=3\left(\sqrt{a}-\frac{2}{3}\right)^{2}
\]

现在高低差已经被表示成了 \(a\) 的函数（上面的四段），这个函数恒 \(\geq 2 m\) ，所以只需要求这个静态函数的最小值，这已经是轻而易举。直接就能判断出它在 \(a \leq \frac{4}{3}\) 时都递减，在 \(a \geq \frac{4}{3}\) 时都增，于是高低差函数的最小值在 \(a=\frac{4}{3}\) 时取，代人得到最小值是 \(\frac{16-8 \sqrt{3}}{3}\) 。因此最后 \(2 m \leq \frac{16-8 \sqrt{3}}{3}\) ，得到最后 \(m\) 的最大值是 \(\frac{8-4 \sqrt{3}}{3}\) 。

总结：\\
本题综合运用了封装命题的处理思想，通过不断地依靠冻结来消去靠后的约束变元，最后回归到问题本身的主要矛盾上去。因为题目的约束变元比较多，因此消去过程也比较多，从而增加了题目本身的思维容量。另外，在消去过程中，预判，思想试验，遍历等思想均得到了综合使用，最后的计算过程也很考察计算基本功。整

体而言，本题的面向群体是 130 分及以上的同学。本题或许也有更加简易的方法，但视频中呈现的方法则更容易想到，更加按部就班，且稳定性很高，熟悉本人之前讲过的这些思想，并且有较好的基本功就能基本稳定解出。希望大家通过本题，能够加深自己对＂数学的干货＂视频列表中补充思想的理解，提高自己的思维能力，体会冻结，遍历，设问，思想试验等思想的科学性和稳定性。

\section*{整体思路回顾时间规划：}
读题，识别出封装命题结构： 1 min\\
消去约束变元 \(x_{0}\) ，得到等价命题＂对任意 \(a, b\) 都有 \(\left|x+\frac{4}{x}-a x-b\right|_{\text {max }} \geq m ": 0.5 \mathrm{~min}\)预判下一步消谁，选择消 \(b: 1.5 \mathrm{~min}\)\\
通过思想试验，明白 \(\varphi(b)_{\min }\) 就是高低差除以 2 ，从而消去 \(b: 2.5 \mathrm{~min}\)\\
讨论 \(\lambda(x)\) 的高低差，并判断其关于 \(a\) 的单调性： 3.5 min\\
求出最后的答案： 1 min 合计： 10 min

\section*{课后练习：}
我们在视频中说到，\(\varphi(b)_{\min }\) 就是函数的高低差（最大值和最小值的差除以 2 ）。这一事实比较容易在直观上接受，但是我们并没有严格写出逻辑证明。请大家尝试给出该命题的严格描述，并逻辑严谨地给出证明。

\section*{6．6．5 两小题让结构深入你心}
结构的力量：把两道小题稳做，让结构深入你心

\section*{题目一：}
已知 \(a, b \in R\) ，且 \(a b \neq 0\) 。若对 \(\forall x>0\) ，均有 \((\ln x-a)(x-b)(x-a-b) \geq 0\) ，则下列四个选项中，可能成立的选项有 \(\qquad\)。（请选出所有可能成立的选项）\\
（1）\(a<0, b<0\)\\
（2）\(a<0, b>0\)\\
（3）\(a>0, b>0\)\\
（4）\(a>0, b<0\)

题目ニ：\\
已知 \(a, b \in R\) ，设函数 \(f(x)=\cos 2 x, g(x)=a-b \cos x\) ，若当 \(f(x) \leq g(x)\) 对任意的 \(x \in[m, n](m<n)\) 恒成立时，\(n-m\) 的最大值为 \(\frac{3 \pi}{2}\) ，则下列四个选项中，可能成立的选项有 \(\qquad\)。（请选出所有可能成立的选项）\\
（1）\(a \geq \sqrt{2}-1\)\\
（2）\(a \leq 1-\sqrt{2}\)\\
（3）\(b \geq 2-\sqrt{2}\)\\
（4）\(b \leq \sqrt{2}-2\)

这两道题有一个关键的共同点，就是它们都是在研究一个较为复杂的系统的行为。之所以说是＂较复杂＂，是因为这两道题的系统都是由两个主自由度所控制，在运动的自由程度上就＂高人一等＂。至于选项，则都是些关于主自由度的范围的提问。很明显，在做题之前，需要搞清楚题目的已知条件（更准确的说法应该是要求，约束这类词汇）对主自由度的限制是怎样的。但是，本身主自由度 \(a, b\) 对系统的控制情况就比较复杂，比如第一题相

当于考察被 \(a, b\) 控制的函数 \(f(x)=(\ln x-a)(x-b)(x-a-b)\) 的信息，考察它的正负，这个控制本身就看上去偏复杂；第二题也是相当于先去控制 \(h(x)=f(x)-g(x)\) 这个函数，研究的重心也是它。

我们之前讲过很长时间要素关系网的思维方式，这本质是一种控制。这是极为普适的理念，上面这两道题也不例外。控制过程（关系网本身的运作机理）常常不是出题考察的对象本身，出题目总是添加上一些额外的牵制，或者不等式关系等的约束条件，让我们去计算，去推导。但是，我所一直强调的做题理念就是：在开始最终解题之前，要先搞清楚系统本身长什么样，它的运作机理（关系网，内部的控制）是怎样的。这与传统的，以问题，解题为本的思维方式是截然不同的：传统思维往往是直接上手从已知设法推结论，算所求，检索知识点，不重视问题本身的结构。大家可以反思一下平时我讲题目的过程，很少上来就直接去思考怎么搞掉最终所求，大多都是先领大家好好分析清楚问题本身的结构。我所讲的一切思维方式，归根结底就是四个字：认清结构。本期视频这两道题，我们依然要沿着就直接去分析它的结构。

先来看第一道题。已知相当于要求 \(a, b\) 均非 0 。然后系统本身就是一个函数（被 \(a, b\) 完全控制）：\(f(x)=\) \((\ln x-a)(x-b)(x-a-b)\) ，讨论的不等式就是 \(f(x) \geq 0\) ，本质就是看 \(f(x)\) 在 \((0,+\infty)\) 上的正负情况。涉及到 3 个因式，要讨论正负，显然就要看那三个因式各自的零点，分别是 \(e^{a}, b, a+b\) 。移动约束变元 \(x\) ，简单的思想试验不难发现 \(f(x)\) 只有可能在经过上面这三个零点的时候都要变号，如果三个零点依次排列在正半轴上，那么就要变三次号；如果有哪个零点不是正的，那就直接无视掉；如果有零点重合的情况，比如两个零点重合，那么在 \(x\) 跨过这个地方的时候 \(f(x)\) 就不变号了。以上其实已经给出这个系统的结构了。有了 \(a, b\) 之后，把 \(e^{a}, b, a+b\) 里面的正的那些排列到正半轴上，这样直接就能读出 \(f(x)\) 的符号。

系统的行为明晰了之后呢，接下来就可以开始解题了。题目要求 \(f(x) \geq 0\) 恒成立，有的同学可能会直接开始思考＂如果要恒成立，那么零点的排列应该怎样呢＂，然后从问题出发去推理出 \(a, b\) 应满足的条件。因为现在系统清晰程度不错了，所以这样当然也可以。但是，除此之外还有一种方式：仍然从系统本身出发，直接对 \(a, b\) 本身分情况讨论，然后看各个情况之下的结论是否匹配题目要求。后者更加倾向于我们之前说的＂结构为上＂的思维方式：仍然是重点看系统本身的控制情况。这两种做法在思维方式上是完全不同的。

我们主要按照后者去讲解（前者的围绕问题的思考大家可自行去尝试）。根据前面对结构的分析，首先要看 \(e^{a}, b, a+b\) 里面哪些是正的。 \(e^{a}\) 一定正，后面两个不好说，但是显然 \(b\) 作为直接性的主自由度，其正负更好讨论。所以先分出 \(b<0\) 与 \(b>0\) 两个情况，然后在这个 \(b\) 的基础上（冻结思想）进一步讨论 \(a\) 。如果是 \(b<0\) ，那么直接扔掉 \(b\) ，这时 \(a+b\) 还不好说在哪。如果它不正，那么连它也扔了的话就显然没法不等式恒成立，于是 \(a+b\)得正。稍微移动一下就发现它必须和 \(e^{a}\) 重合，也就是 \(e^{a}-a=b\) ，直接把位置卡死了！顺着 \(b\) 去找 \(a\) 发现不等式无解。所以这种情况不可能。

只能 \(b>0\) 了！这时候两个正根了（相对位置暂不明）。思考发现急着讨论他俩相对位置也没啥意义，得先分析 \(a+b\) 的位置。如果 \(a+b\) 也正，那么三个正根排一起，最左边正根再往左，\(f(x)\) 就负了，不行。所以必须 \(a+b \leq 0\) ，然后扔掉它。为了 \(f(x) \geq 0\) 恒成立，像刚才那样，得两正根相等：\(e^{a}=b\) 。合起来看就是 \(a+b \leq 0\) 且 \(e^{a}=b\) 。注意刚才做的都是充要替换，所以这就是原命题的充要条件。 \(e^{a}=b\) 让 \(a, b\) 可以彼此控制，然后 \(a+b \leq 0\)就等价于 \(a+e^{a} \leq 0\) ，只要是这样的 \(a\)（再控制出 \(b\) ）就行了。显然这时一定有 \(a<0\) ，所以可能成立的选项只有（2）。

这道题在方法上其实就是穿根法的变种，如果基础知识足够扎实的话，直接用类似穿根法的思路分析也没问题，更加快捷。视频其实就是把穿根法用我们讲过的遍历法的思维方式再领大家认识一遍，帮助大家更好的理解知识背后的本质，用这种动态的视角去看问题。大家在平时学习知识的时候，要多用动静视角以及要素关系网的思维方式看问题，提炼其中的动静关系，以元问题的形式认清一个知识点的核心本质。

我们再来看第二道题。令 \(h(x)=\cos 2 x+b \cos x-a\) ，这就是要研究的函数。同样的，我们先看问题的含义。显然，问题的意思是：在 \(h(x) \leq 0\) 的解集里找一个最长的连续区间［ \(m, n\) ］，其长度 \(n-m=\frac{3 \pi}{2}\) 。因此，从结构的角度出发，我们先来关注 \(h(x) \leq 0\) 这个解集本身是什么情况。于是把思维焦点先抽出来：怎么求解 \(h(x) \leq 0\) ？

一旦你这么想，你就开始试着在脑内模拟求解 \(h(x) \leq 0 a, b\) 都是参数，就是要解上述关于 \(x\) 的不等式。这个就非常基础了：把 \(\cos 2 x\) 展开：\(h(x)=2 \cos ^{2} x+b \cos x-a-1\) 。因此令 \(\varphi(t)=2 t^{2}+b t-(a+1)\) 。先求解不等式 \(\varphi(t) \leq 0\) 得到一个 \(t\) 的取值范围（这个范围显然也是直接由 \(a, b\) 控制）。然后让 \(\cos x\) 在这个范围内，就可以定出 \(x\) 的最终范围。于是这个控制关系可以被如上拆解成两层。要看最终的 \(x\) 范围，就要根据这个控制关系去研究。

首先是 \(\varphi(t) \leq 0\) 。这时解出的 \(t\) 范围是什么？根据基本的二次不等式的常识，抛物线开口向上，有两个参数在控制它，因此要么无解或者仅一个点成立（ \(\Delta \leq 0\) ，这显然不满足已知，不考虑），要么就是一个闭区间，记为 \(\left[t_{1}, t_{2}\right]\)。当 \(\Delta>0\) 时，\(t_{1}, t_{2}\) 直接被 \(a, b\) 彻底控制。然后就要让 \(\cos x \in\left[t_{1}, t_{2}\right]\) ，去确定最终的 \(x\) 范围，而这个控制就太简单了，直接画个 \(y=\cos x\) 图像，两条横线 \(y=t_{1}, y=t_{2}\) 一划一夹，马上看清范围。

控制关系搞清楚了，我们就开始做题。题目要求最终 \(x\) 范围的最长的连续区间长度是 \(\frac{3 \pi}{2}\) ，那就画好 \(y=\cos x\) ，然后去尝试遍历地移动一下 \(t_{1}, t_{2}\) 的位置，看看什么样的 \(t_{1}, t_{2}\) 才能满足条件。既然要遍历，就先瞅瞅 \(t_{1}, t_{2}\) 能遍历的定义域。系统是两个自由度，所以不难看出 \(t_{1}\) 和 \(t_{2}\) 之间没有等量关系牵制，就是一个任取的闭区间。所以遍历的范围要遍布整个实数区域。多自由度想去遍历，仍然采用冻结，如先冻结 \(t_{1}\) ，然后把 \(t_{2}\) 放到它上面去滑动。

起始的 \(t_{1}\) 位置选取也是一个遍历过程，先考虑 \(t_{1}\) 很小（比如－100000），然后 \(t_{2}\) 一路往上滑，并同步观测，不难看出当且仅当 \(t_{2}=\frac{\sqrt{2}}{2}\) 时满足题意。然后慢慢把 \(t_{1}\) 往上推，易见这时候 \(t_{2}\) 的遍历过程都一样（都得 \(t_{2}=\frac{\sqrt{2}}{2}\) ），这种情况一直持续到 \(t_{1}=-1\) 。可见 \(t_{1} \leq-1\) 时必须且仅需 \(t_{2}=\frac{\sqrt{2}}{2}\) 。继续往上推 \(t_{1}\)（这时 \(-1<t_{1}<1\) ），然后固定 \(t_{1}\) 滑动 \(t_{2}\) ，显然如果 \(t_{2}<1\) ，那么解集是各个单调区间上的一段一段并起来，长度不可能到 \(\frac{3 \pi}{2}\) ，但是如果推到 \(t_{2} \geq 1\) ，这时候有的单调区间会拼起来，变得很长。这时候不难看出当且仅当 \(t_{1}=-\frac{\sqrt{2}}{2}\) 时满足题意，并且只要 \(t_{1}=-\frac{\sqrt{2}}{2}\) ，那么随便让 \(t_{2}\) 遍历 \([1,+\infty)\) 都行。这看上去就是把刚才的情形对称翻上去了。下一步，如果 \(t_{1} \geq 1\) ，那么显然无论怎么滑动 \(t_{2}\) 都不满足条件。于是我们得到的原题干的充要替换：要么 \(t_{1} \leq-1, t_{2}=\frac{\sqrt{2}}{2}\) ；要么 \(t_{2} \geq 1, t_{1}=-\frac{\sqrt{2}}{2}\) 。

现在扡弃原已知，就用上面的这个关于 \(t_{1}, t_{2}\) 的限制去分析 \(a, b\) 的信息。 \(\left[t_{1}, t_{2}\right]\) 是什么？它是 \(2 t^{2}+b t-(a+1) \leq 0\)的解集，因此很容易看出 \(t_{1}, t_{2}\) 就是 \(\varphi(t)\) 的零点，所以把元问题的控制内核提炼出来：题目就是要我们通过双自由度的二次函数的零点信息去反映主自由度的信息。一旦提炼出这个内核，你马上就会自然的考虑：怎么用零点的信息反映二次函数本体的信息？这对应什么知识点？那就显而易见了：二次函数的零点式，或者韦达定理，都可以起到这个作用。我们可以把 \(\varphi(t)\) 重新写成 \(\varphi(t)=2\left(t-t_{1}\right)\left(t-t_{2}\right)\) ，展开并与原系数对比得到 \(b=-2\left(t_{1}+t_{2}\right), a=-1-2 t_{1} t_{2}\)（这个你用韦达定理也可以直接得到）。韦达定理在元问题层面的理解就是：二次函数的本体和零点可以互相控制，零点对二次函数的控制在齐次意义（大小无关）下非常彻底。这样，一下子就把问题本质抓出来了。我在经验分享第五期强调的元问题的意义得到了充分体现：平时学知识的时候都花时间想清楚知识点背后的控制关系，然后做题时利用要素关系网的思维方式把题目背后的控制关系抓出来，然后匹配上对应的元问题（知识点）。只有这样，做题才能用对知识，用好知识。

到这里，问题已经毫无难度了，根据之前对 \(t_{1}, t_{2}\) 的限制讨论即可。若 \(t_{1} \leq-1, t_{2}=\frac{\sqrt{2}}{2}\) ，则

\[
b=-2\left(t_{1}+\frac{\sqrt{2}}{2}\right) \geq 2-\sqrt{2}, a=-1-\sqrt{2} t_{1} \geq \sqrt{2}-1,
\]

范围也有了；而若 \(t_{2} \geq 1, t_{1}=-\frac{\sqrt{2}}{2}\) ，则

\[
b=-2\left(t_{2}-\frac{\sqrt{2}}{2}\right) \leq \sqrt{2}-2, a=-1+\sqrt{2} t_{2} \geq \sqrt{2}-1 .
\]

对应到选项上就是（1）（3）（4）。

\section*{总结：}
这两道题都很好地反映了通过结构解题的理念。我们讲解自由度，要素关系网，元问题，归根结底都可以用结构二字来概括。在学知识的时候把知识点的控制关系想透，在做题的时候把题目的结构搞清，才能找到稳定正确的解题方向，才能恰当地使用知识。当拿到一道陌生的问题的时候，与其急于去想题目怎么做，去想该用什么知识点，不如先分析清楚题目的结构，搞清楚系统的要素关系，然后基于系统的要素关系本身去匹配知识点，题型，甚至二级结论等，这些你学过的所有东西。这需要大家把思维方式转变过来，也需要大家按照经验分享第五期的理念去把知识点本身形成自己结构性的理解。这是我们一直强调的理念，这也是我们不断实践的理念。希望大家能够更好地发挥结构思维的力量。

\section*{6．6．6 逆向思考提升数学思维}
题目：\\
已知函数 \(f(x)=a x^{2}+(a-2) x-\ln x(a \in \mathbb{R})\)\\
（1）讨论 \(f(x)\) 的单调性\\
（2）若 \(f(x)\) 有两个零点 \(x_{1}, x_{2}\) ，求 \(a\) 的取值范围，并证明：\(x_{1}+x_{2}>\frac{2}{a}\)

第一问是个单纯的单自由度讨论，

\[
f^{\prime}(x)=2 a x+(a-2) x-\frac{1}{x}=\frac{(a x-1)(2 x+1)}{x},
\]

定义域又是 \((0,+\infty)\) ，所以关键因子就在 \(a x-1\) 上，\(a \leq 0\) 时，\(f(x)\) 减；而 \(a>0\) 时，\(f(x)\) 在 \(\left(0, \frac{1}{a}\right)\) 上减，在 \(\left(\frac{1}{a},+\infty\right)\)上增

第二问的前半部分也不难，非常经典的找零点问题。很明显，必须 \(a>0\)（减函数出不来俩零点），而且低谷 \(f\left(\frac{1}{a}\right)<0\) ，还需要保证在 \(x \rightarrow 0\) 和 \(x \rightarrow+\infty\) 时能回到正值。先解不等式

\[
f\left(\frac{1}{a}\right)<0 \Leftrightarrow \ln a-\frac{1}{a}+1<0 \Leftrightarrow 0<a<1,
\]

然后还需要保证两个极限。这时候望气，\(x \rightarrow 0\) 时多项式部分变成常数，\(-\ln x\) 跑到正无穷去，可以！再看 \(x \rightarrow+\infty\) ，由于 \(a>0\) ，二次函数撑起门面，尽管对数拖后腿，但是之前讲过，多项式比对数快很多，直接碾压，跑到正无穷，没问题！所以 \(0<a<1\) 就是正确答案。

想要拿全分的话还是最好用找点的办法。之前我们也是单独做过找点的程序性完成方法（见＂数学的干货＂）。这里复习一下，做法就是＂扬长避短＂。首先看 \(x \rightarrow 0\) ，我们希望找到很小的 \(x\) 让 \(f(x)>0\) ，由于这里是两个有界量，一个无穷大量，于是先把有界量随便扔掉。把 \(a\) 看作静止，直接开扔：\(a x^{2}+(a-2) x>-2 x\) ，所以只需要让 \(-2 x-\ln x>0\) ，这个取 \(x=e^{-2}\) 即可满足（注意 \(\frac{1}{a}>1\) ）。跑到无穷的话得先把对数放大一点，还不能大过二次函数，那就用一下不等式 \(\ln x<x\)（这样放缩就够了），于是 \(f(x)>a x^{2}+(a-3) x\) ，于是会发现只要 \(x \geq \frac{3-a}{a}\) 即可，随便取一个就行（比如就 \(x=\frac{3-a}{a}\) ，它显然比 \(\frac{1}{a}\) 大）。所以我们就取了 \(x=e^{-2}\) 和 \(x=\frac{3-a}{a}\) ，它们的函数值都为正，再结合零点存在性定理和单调性就得知 \(f(x)\) 有两个零点。

前面的过程都是程序性的，套路化的。然后我们来看最难的 \(x_{1}+x_{2}>\frac{2}{a}\) 证明部分。从要素网的结构出发，就是一个主自由度 \(a \in(0,1)\) 先决定了 \(a x_{1}^{2}+(a-2) x_{1}=\ln x_{1}\) 的 \(x_{1}<\frac{1}{a}\) 以及 \(a x_{2}^{2}+(a-2) x_{2}=\ln x_{2}\) 的 \(x_{2}>\frac{1}{a}\) ，这里注意 \(x_{1}\) 和 \(a\) 和 \(x_{2}\) 是可以互相决定的（之前讲导数题也提到过这一点，也就是存在主自由度选取的任意性）。两个零点正好分居 \(\frac{1}{a}\) 两侧，所证就是它们的平均值相比 \(\frac{1}{a}\) 向右偏移，非常标准的极值点偏移型问题。不过我们今天想要讲另外一种解法，也就是在不会极值点偏移的情况下直接证明。之所以要讲这样的方法，是因为其涉及到另一种思维方式——分析法，这样的思维方式肯定是更广泛的，有助于我们处理更多问题的。

所谓分析法，就是逆推法，通过分析＂什么样的条件能推出结论＂来对证明过程进行逆推。不断的逆推最后会达到一个简洁易证的结果，完成这个结果的证明就完成了整个题目的证明。这便是分析法的思维方式，和我们平时做题使用的综合法有所区别。分析法并不是沿着要素网顺向推理，因此在计算型问题中比较少见，最常见于推理分析型问题，也就是证明题之中，而且对于偏难的证明题常常有重要的帮助。

我们的已知（也就是牵制关系）是两个式子，信息量全在这里头。因此从逻辑上来讲，肯定是对这两个式子进行处理就能做出题来。所证式子具有高度对称性，但是涉及到三个参量，相比于实际自由度 1 而言太多了，肯定没办法直接证明，需要把参数弄的少一点，才方便我们证明。那如何把参数弄的少一点（如何消元）？正向不知道该如何思考，我们先看所证。选 \(a\) 为主自由度去证它应当是不合适的，因为 \(x_{1}+x_{2}\) 没法写成一个关于 \(a\) 的显式表达式。所以必须得选 \(x_{1}\) 和 \(x_{2}\) 作为主自由度，附带一条牵制（小技巧：\(x_{1}\) 和 \(x_{2}\) 地位对等，在初步思考解法的时候我们一般把这两个都先留着）。这样的话我们就开始尝试先消去 \(a\) ，自然就是要利用已知的两个式子。为了保持 \(x_{1}\) 和 \(x_{2}\) 的地位对等，我们需要把两个式子进行整体操作，也就是：相加或者相减。对数的问题一般采用减法，因为容易出齐次式。

\section*{已知两式相减得}
\[
a\left(x_{1}-x_{2}\right)\left(x_{1}+x_{2}\right)+(a-2)\left(x_{1}-x_{2}\right)=\ln \frac{x_{1}}{x_{2}}
\]

多余的因式除到右边变成

\[
a\left(x_{1}+x_{2}+1\right)-2=\frac{\ln \frac{x_{1}}{x_{2}}}{x_{1}-x_{2}}
\]

这是我们直接变形得到的结果。注意所证式 \(x_{1}+x_{2}>\frac{2}{a}\) ，我们现在可以对它充要替换（逆向变形）了。

\[
a>\frac{2}{x_{1}+x_{2}} \Leftrightarrow \frac{\ln \frac{x_{1}}{x_{2}}}{x_{1}-x_{2}}>\frac{2\left(x_{1}+x_{2}+1\right)}{x_{1}+x_{2}}-2
\]

再等价一步就成了

\[
\frac{\ln \frac{x_{1}}{x_{2}}}{x_{1}-x_{2}}>\frac{2}{x_{1}+x_{2}}
\]

请大家注意这里的逻辑：

\[
a\left(x_{1}+x_{2}+1\right)-2=\frac{\ln \frac{x_{1}}{x_{2}}}{x_{1}-x_{2}}
\]

是亘古不变的真理，而我刚才对 \(x_{1}+x_{2}>\frac{2}{a}\) 做的变形都是等价变形，也就是：要证这个，即证那个，所以最后那个

\[
\frac{\ln \frac{x_{1}}{x_{2}}}{x_{1}-x_{2}}>\frac{2}{x_{1}+x_{2}}
\]

只要证出来了，反推回去原命题就得证了。而这已经是一个齐次式了，必然等价于 \(\frac{x_{1}}{x_{2}}\) 的范围问题，这么一充要替换，问题就变得很简洁了。我们最开始不妨设 \(x_{1}<x_{2}\) ，于是上式进一步等价变形为

\[
4 \ln \frac{x_{1}}{x_{2}}<\frac{2\left(x_{1}-x_{2}\right)}{x_{1}+x_{2}},
\]

令 \(t=\frac{x_{1}}{x_{2}}\) ，则上式进一步等价为

\[
\frac{2(t-1)}{t+1}-\ln t>0 \Leftrightarrow 0<t<1
\]

（基本的构造函数解不等式，不展开讲了）。一路逆向充要替换下来，竟成了原初假设 \(x_{1}<x_{2}\) ，不证自明，于是问题就解决了。

分析法写步骤则需要倒回去，先不妨设 \(x_{1}<x_{2}\) ，然后构造函数证明

\[
\ln \frac{x_{1}}{x_{2}}<\frac{2\left(x_{1}-x_{2}\right)}{x_{1}+x_{2}}
\]

然后和刚才的逆推反向，先推出

\[
\frac{\ln \frac{x_{1}}{x_{2}}}{x_{1}-x_{2}}>\frac{2}{x_{1}+x_{2}},
\]

再变成

\[
\frac{\ln \frac{x_{1}}{x_{2}}}{x_{1}-x_{2}}>\frac{2\left(x_{1}+x_{2}+1\right)}{x_{1}+x_{2}}-,
\]

然后左边换成 \(a\left(x_{1}+x_{2}+1\right)-2\) ，最后得到 \(a>\frac{2}{x_{1}+x_{2}}\) 从而得证。\\
由此可以看到，分析法的做题思维就是从所证结论入手，结合要素网带给我们的信息去尝试逆推，从而简化命题，让我们更加看清问题所在。需要注意的是，这里的逆推不需要非得保持充要替换，即使是充分替换，只要最后逆推回去的东西确实是对的，那也可以。不过，充分替换虽然符合数学的逻辑规则，但是有可能存在＂放缩过度＂的情况。像这个题一直是利用充要替换来进行回溯的，所以不可能出现上述情况，一定可以做出来。换句话说，当我看到回溯成一个齐次式 \(\frac{\ln \frac{x_{1}}{x_{2}}}{x_{1}-x_{2}}>\frac{2}{x_{1}+x_{2}}\) 的时候，我在心里就已经明白，这道题做完了。

\section*{思考题}
1.

请使用标准的极值点偏移的方法重新证明一下 \(x_{1}+x_{2}>\frac{2}{a}\) ，也就是：不妨设 \(x_{1}<x_{2}\) ，上式等价为 \(x_{2}>\frac{2}{a}-x_{1}\) ，然后论证这等价于 \(f\left(x_{1}\right)>f\left(\frac{2}{a}-x_{1}\right)\) ，令 \(g(x)=f(x)-f\left(\frac{2}{a}-x\right)\) ，论证清楚 \(g(x)\) 的单调性，然后得证原命题\\
2.\\
（1）根据提示来完成下面一道题的证明：\\
题目：函数 \(f(x)=\frac{\ln x}{x}\) ，若 \(f(x)=a\) 有两个解 \(x_{1}<x_{2}\) ，证明：\(x_{1} x_{2}>e^{2}\)\\
提示：先尝试证明

\[
\frac{\ln \frac{x_{1}}{x_{2}}}{x_{1}-x_{2}}=\frac{\ln x_{1} x_{2}}{x_{1}+x_{2}},
\]

然后利用这个式子构造函数直接证明原命\\
（2）在第（1）问的提示中给出的式子（在加上 \(x_{1}<e<x_{2}\) 的限制后）是否包含了题目已知的全部信息量？请结合问题的自由度进行简要说明\\
（3）你认为第（1）问的那个提示式是如何想到的？（提示：想想初中学分式的时候学过什么基本性质？）\\
（4）如果把第（1）问中的所证式改成\\
\(x_{1}+x_{2}>2 e\) ，请你说明，刚才在原第（1）问中证明的结论可以直接推出这个式子。如果考试时探索发现了

\[
\frac{\ln \frac{x_{1}}{x_{2}}}{x_{1}-x_{2}}=\frac{\ln x_{1} x_{2}}{x_{1}+x_{2}}
\]

你是否有信心先大胆猜 \(x_{1} x_{2}>e^{2}\) ，从而利用这个结论直接证明 \(x_{1}+x_{2}>2 e\) ？（提示：对于新构造的函数，请你观察它在 \(t=\frac{x_{1}}{x_{2}} \rightarrow 0\) 时的极限行为如何？这样的观察是否能给你带来信心？）\\
3.

对比视频中的题目，思考题 2 的第（1）问中的题目，以及 \(u p\) 主在2022年11月6日投稿的视频中的题目，这三道题目的处理方法有什么共性？结合刚才提到的 11 月 6 日视频中对思维方式的讲解，你对于控制理念是否有更深刻的感悟？

\section*{6．6．7 XXXX}
\section*{6.7 立体几何与概率}
\section*{6．7．1 高考算难题，思维下纯送}
视频日期：2023年11月11日

\section*{题目：}
已知在四棱雉 \(P-A B C D\) 中，底面 \(A B C D\) 是直角梯形，\(A B / / C D, \angle A B C=90^{\circ}\) ，且 \(A B=2, B C=2 \sqrt{3}\) 。若 \(P A=P D, P C=P B\) ，且三棱雉 \(P-A B C\) 外接球的表面积为 \(20 \pi\) 。则当四棱雉 \(P-A B C D\) 的体积取得最大值时， \(C D\) 的长为 \(\qquad\)。

这道题目的题干里涉及到的要素比较多，因此按照一个舒服的画图顺序来搭建题目的要素关系网就很重要。首先是＂四棱雉＂，碰到四棱雉，大部分情况都是先搭建底座，再架上锥点，此题也不例外。先叙述了底座 \(A B C D\) 是直角梯形，然后有两条边的长度已知。简单的画图法就可以知道底座只剩下一个自由度（并且这个主自由度分配给 \(C D\) 非常自然）。再后面给了两个条件 \(P A=P D, P C=P B\) 来刻画 \(P\) 的位置。但是在固定底座的前提下，需要三个自由度才能定下 \(P\) ，因此 \(P\) 还有一个自由度的活动空间。然后是一个外接球的已知，这显然就不是让我们直接用于搭建要素网用的控制条件，而是限制整个系统行为的牵制条件，它再减一个自由度。所以题目的总自由度为

1 。题目最后要我们考察 \(P-A B C D\) 体积（自然就是被一个自由度控制的单参数函数）的最值，并求出该最值取得的位置（说是问 \(C D\) 的长，其实说白了就是问最值取等的位置）。不过既然我们刚才说过 \(C D\) 作为底座的主自由度很舒服，现在又要去求它的值，不把它定为主自由度来做题都对不起出题人如此好心的提醒了。

底座很好画。接下来的已知两个刻画 \(P\) ，一个是外接球的牵制。我们讲过多次，牵制在研究好关系网，选好主自由度之前可以不着急，更何况外接球接的是三棱雉 \(P-A B C\) ，连 \(P\) 在什么位置还不清楚，如何研究外接球？因此下一步肯定就是去分析 \(P\) 的位置如何刻画，以及选取什么主自由度来控制 \(P\) 。明确好了这个思维焦点之后，我们来尝试分析 \(P\) 。先设一个 \(C D=t\) ，暂时把底座控制住。之前多次强调，设未知数就要看定义域。显然作为梯形，必有 \(t>0\) 且 \(t \neq 2\) 。

那么，\(P\) 的位置如何再用一个主自由度刻画呢？这就有不同的方法了。我们讲两种方法。第一种肯定是尝试利用已知信息推导，做充要替换。我们知道 \(P\) 自由度是 1 ，也就是说，它在某条线上运动，所以最好的结果就是把这条线找出来。再看已知：\(P A=P D\) 这个条件，和动态的 \(D\) 有关，但是 \(A, B, C\) 都可视为静止，所以先看一下条件 \(P C=P B\) 。什么样的点到 \(B, C\) 距离相等呢？想象一下桌子上摆两个点，到这两个点的距离相等，那不就是一把刀顺着中间国下去，这个辟面上的点吗？严格表述的话，取 \(B C\) 中点 \(M\) ，过 \(M\) 作平面 \(\alpha \perp B C\) 。显然 \(P M \perp B C\) ，因此 \(P M \subset \alpha\) 。所以我们只需要在这个平面 \(\alpha\) 里找能够使得 \(P A=P D\) 的点 \(P\) 了。但是这时候我们有一个观察，就是 \(\alpha\) 这一刀辟下来，在底座上划出的线正好是梯形 \(A B C D\) 的中位线。如果我们再取 \(A D\) 的中点 \(N\) ，那么 \(M N\) 就平行于梯形 \(A B C D\) 的上下底，也就垂直于 \(B C\) ，因此 \(M N \subset \alpha\) 。而刚才看不太清晰的 \(P A=P D\) 其实就是 \(P N \perp A D\) ，那这感觉上自然就得把 \(P N\) 在 \(\alpha\) 里面转直了。我们尝试把这种感觉严格说清楚：首先 \(P N\) 也在 \(\alpha\)里，这等价于 \(P N \perp B C\) 。然后现在又有 \(P N \perp A D\) ，注意梯形的限制要求 \(B C, A D\) 不平行，因此 \(P N\) 与底座垂直。这个＂与底座垂直＂的限制直接把 \(P\) 约束在了一条线上，实现了自由度为 1 的约束。这个约束条件是否完全刻画了 \(P\) 的信息（是否实现了充要替换）呢？自由度分析告诉我们：是。实际上大家也不难验证，只要 \(P\) 垂直于底面，那么就可以推导出 \(P A=P D, P C=P B\) 。这也体现了我们平时讲到过的动静视角用于分析信息量（充要替换）的思想：如果在动静关系上，一组条件和另一组条件对系统的约束效果一样（自由度减少量一致，都能定死系统等），那么这两组条件的信息量（对系统的约束效力）就是一样的。这种融合宏观信息与微观推理的能力希望大家学习。

有的同学可能会担心，自己没看出中位线咋办？其实，大家从上面的分析过程中可以看出，在看准了 \(\alpha\) 的前提下，顺着 \(\alpha\) 找 \(P\) ，视野自然就会沿着 \(\alpha\) 去关注 \(\alpha\) 和底座的交线。更何况两个对 \(P\) 的限制条件如此相像，第一个取了中点之后第二个如法炮制是很自然的考量，然后把这些思路放在一起，自然就能看到中位线了。如果还是觉得难想的话，up 主给大家还提供了第二种研究 \(P\) 的思路：建系计算。实际上，在很多需要定量研究的几何场景中，建系常常是非常好的元问题处理办法之一。不过，和很多其他元问题处理一样，在建系计算之前，都建议预判一下计算量。这里我们需要把 \(P A=P D, P C=P B\) 列成两个方程，进而搞清 \(P\) 的位置。而这种距离相等，会让平方项中的 \(x^{2}, y^{2}, z^{2}\) 项全都抵消掉，因此我们得到的就是两个一次方程，这显然是很好处理的。因此我们开始建系。为了把静止的要素 \(A, B, C\) 的优势发挥出来，我们选 \(B\) 为原点， \(\overrightarrow{B A}\) 方向为 \(x\) 轴， \(\overrightarrow{B C}\) 方向为 \(y\) 轴建系，坐标都很简单：\(A(2,0,0), C(0,2 \sqrt{3}, 0), D(t, 2 \sqrt{3}, 0)\) 。然后设 \(P(x, y, z)\) ，把 \(P A=P D\) 变成方程

\[
(x-2)^{2}+y^{2}+z^{2}=(x-t)^{2}+(y-2 \sqrt{3})^{2}+z^{2} \Leftrightarrow(2 t-4) x+4 \sqrt{3} y=t^{2}+8
\]

再把另一个 \(P C=P B\) 变成方程 \(x^{2}+(y-2 \sqrt{3})^{2}+z^{2}=x^{2}+y^{2}+z^{2} \Leftrightarrow y=\sqrt{3}\) ，代回第一个式子整理得 \(x=\frac{t+2}{2}\)。这两个式子就是 \(P\) 的全部约束，显然描述的是 \(P\) 的地面投影位置。主动寻找这个位置尝试标记，显然就能发现 \(P\) 的地面投影在 \(A D\) 中点。不过即使发现不了也没关系，有了定量的 \(P\) 点位置描述，何愁大事不成？

上述两种方法都告诉我们，\(P\) 就是在一条坚直旗杆上运动罢了。所以第二个主自由度 自然分配给 \(P\) 的距地面高度，设为 \(h(h>0)\) 。现在我们是两个主自由度 \(t, h\) ，接下来自然是把剩下的已知（牵制关系 \(S=4 \pi R^{2}=\) \(20 \pi \Leftrightarrow R=\sqrt{5}\) ）用上，来得到 \(t\) 之间的牵制关系，完成整个要素网的刻画。外接圆问题我们之前都放在＂数学的干货＂视频列表里（日期为 2023 年 1 月 15 日），这个元问题的处理方式很单一，就是穿刺法。研究对象：三棱雉 \(P-A B C\) 底面纯静止，外接圆圆心就是 \(A C\) 中点，记为 \(O^{\prime}\) 。再设 \(P-A B C\) 的外接球球心为 \(O\) ，于是在 \(O^{\prime}\) 处坚直穿刺一根轴，\(O\) 就在这个轴上。现在底面静止，半径还作为牵制给我们了，所以简单计算得到 \(O O^{\prime}=1\) 。穿刺法给出的式子是 \(O P=O A\) ，于是牵制就相当于 \(O P=\sqrt{5}\) ，因此我们只需要用 \(t, h\) 去表示 \(O P\) 即可。 \(O\) 和 \(P\) 的坚直高度差为 \(|h-1|\) 。至于水平差，那自然就是二者的地面投影点：\(O^{\prime}\) 与 \(N\) 的距离了。用几何法做到现在的同学，自然会关注这两个点是如何生成的，然后发现它刚好是 \(C D\) 的中位线，于是 \(O^{\prime} N=\frac{t}{2}\) ，因此 \(|O P|^{2}=\frac{t^{2}}{4}+(h-1)^{2}=5\)。而用建系法做到现在的同学，则会用更加简单粗暴的方法：坐标直接算。显然 \(O(1, \sqrt{3}, 1), P\left(\frac{t+1}{2}, \sqrt{3}, h\right)\) ，因此直接算距离也能得到 \(\frac{t^{2}}{4}+(h-1)^{2}=5\) 。

主自由度，定义域，牵制式全齐活了，剩下的自然就是完成最后一步求最值了。四棱雉体积为

\[
\frac{h}{3} S_{A B C D}=\frac{h}{3} \cdot \frac{1}{2} \cdot 2 \sqrt{3}(2+t) .
\]

由于我们只需关注最值取到的位置，因此只需要看 \(h(2+t)\) 。把思维焦点放过来：如何通过 \(\frac{t^{2}}{4}+(h-1)^{2}=5\) 来算出 \(h(2+t)\) 的最值呢？这就是对于最值问题的基本功考察了。一个基本处理是三角换元：已知形如平方和是定值，因此可设

\[
\frac{t}{2}=\sqrt{5} \cos \theta, h-1=\sqrt{5} \sin \theta\left(\text { 定义域限制 } t>0, t \neq 2 \Rightarrow \theta \in\left(-\frac{\pi}{2}, \frac{\pi}{2}\right), \cos \theta \neq \frac{1}{\sqrt{5}} \cdot\right)
\]

代入得

\[
h(2+t)=2(\sqrt{5} \sin \theta+1)(\sqrt{5} \cos \theta+1) .
\]

同样我们可以忽略 2 ，并展开式子可得到

\[
1+\sqrt{5}(\sin \theta+\cos \theta)+5 \sin \theta \cos \theta
\]

这就是经典问题了，再设 \(u=\sin \theta+\cos \theta\) ，由于恒等式 \((\sin \theta+\cos \theta)^{2}=1+2 \sin \theta \cos \theta\) ，上式又变成

\[
1+\sqrt{5} u+5 \cdot \frac{u^{2}-1}{2}
\]

由于 \(|u| \leq \sqrt{2}\) ，因此不难看出该二次函数在 \(u=\sqrt{2}\) 时最大，此时 \(\theta=\frac{\pi}{4}\) ，对应于 \(t=2 \sqrt{5} \cos \theta=\sqrt{10}\) ，这就是最后的答案。当然，这么简单的函数，求导计算然后找导函数零点解方程也可以。

做个小总结。这道题放在高考真题难度排名里的话是可以排进难题的，但是可以看到，我们踏踏实实用动静视角分析结构，选取主自由度，搭建要素关系网，确立好各步的思维焦点，隔离小问题，处理清楚元问题，步步为营，稳健前行，这道题可以做到不需要任何奇思妙想或者灵光乍现。哪怕是唯一一点的中位线观察，也可以用建系法彻底取代掉。通过这道题，再向大家强调两点。一是建系法，在处理很多需要定量的元问题时都是一种可行的处理办法，虽然不一定是最简方法，但稳定性大多较好。建系不一定是处理完整道题，可能中途的一个小元问题的处理，临时建一下坐标系算算，都能发挥作用，希望大家不要忽视。二是做题要踏实，关系网在脑子里想好，建立清楚再行动；计算前稍微预判一下，确信可算再动笔。相信自己掌握的思维方式的科学性，相信高考考场上题目的

稳定性，稳扎稳打，才能取得满意的成绩。

时间规划：\\
读题，分析好要素网，确立主自由度： 1.5 min\\
（1）几何法：在动静视角的指引下分析出 \(P\) 点的约束： 2 min\\
（2）建系法：直接爆算出 \(P\) 点的约束： 2 min\\
处理外接球元问题，得到主自由度之间的牵制： 2.5 min\\
最后求最值，解答案： 2 min\\
总时间（两种方法均为）： 8 min

\section*{6．7．2 XXXX}
\section*{6.8 不等式问题}
\section*{6．8．1 看似需要经验，实则思维秒解}
视频日期：2024年2月11日

\section*{题目：}
已知 \(x, y, z>0\) ，则 \(\frac{\left(x^{2}+2 y^{2}+z^{2}\right)^{2}+1}{x y+3 y z}\) 的最小值为 \(\qquad\) －

这道题目是一位普通的高中生同学私信提问的问题。在动态发出后，评论区的很多数学成绩优秀，见多识广的同学都给出了自己的解法。像强基计划中的待定系数法，大学数学中的偏导法等等，都是评论区的热词。如果提前学过这些解法，那么这道题就会变成一道毫无难度的二级结论题。关于待定系数法的解法，在视频的末尾会以习题的形式提供给大家，作为一个解法的补充。

下面，我们站在啥也没额外学过的普通高中生的角度，使用 up 主此前讲过的思维方式，来破解这道题。题面很简单，三个独立的操纵杆（主自由度）\(x, y, z>0\) ，让我们计算一个式子的最值。三个自由度的问题，多自由度问题，我们已经强调过无数遍，冻结法是最为普适的思维方式。三个主自由度很直白的告诉了你，就 \(x, y, z\) ，所以最直接的想法自然是挨个试试。所谓冻结，一定是先冻到只剩下一个变量，先研究单自由度问题，取到此时的最值，然后再逐层解冻，重复循环。比如我有三个独立的主自由度 \(x, y, z\) ，我可以考虑先冻结 \(x, y\) ，然后移动 \(z\) ，使得函数取得最值。这时使函数取最值的 \(z\) 一定受到背景的控制，也就是 \(z\) 被 \(x, y\) 控制。然后解冻 \(x, y\) ，此时函数中的 \(z\) 被 \(x, y\) 控制了，实际自由度减了 1 ，成为了一个关于 \(x, y\) 的式子。我们故技重施，对 \(x, y\) 的系统继续实施冻结法，再消掉一个，变成单自由度问题，再研究它，取到最值即可。

遇到多自由度问题，该选谁作为冻结的对象呢？通常都是挨个试一下，因为即便是三个自由度的问题，最开始也只有三种冻结方法，数目很少，完全可以挨个预判一下。但是，这道题从预判开始就出现了问题。无论先冻结 \(x, y, z\) 中的谁，式子似乎都不好研究。

这可怎么办呢？实际上，早在两年之前，在 2021 年 10 月 3 日的视频中，我们就已经可以得到答案了。在那

期视频里，考查了关于 \(a, b\) 的式子

\[
\frac{\sqrt{1+\left(a^{2}+b^{2}-\frac{10}{9} a b\right)^{2}}}{a b}
\]

而当时，我们既没冻 \(a\) ，也没有冻 \(b\) ，而是选择冻 \(a \cdot b\) 这个整体。究其本质，这相当于做了个换元：令 \(a b=\) \(x, a^{2}+b^{2}=y\) ，这样就一下子把式子结构看清楚了。所谓换元，其实就可以理解成代数问题中的重建要素网。原本是 \(a, b\) 作为主自由度，现在则是 \(x, y\) 做主自由度，数目对等，自由度守恒。通过换元，式子的结构（要素网控制形式，或者说元问题）就会发生改变。这不正好和我们在几何问题中讲到的＂要素关系网重构，选取最舒服的要素网＂的理念不谋而合吗？只不过不同的是，代数与几何问题的重构要素网的方式有所区别。几何问题通常是靠不同的画图方法来重建；而代数的换元更加简单粗暴，依据就是对式子结构的把握。此外，代数问题更加强调遍历法的思维，这一点我们后文会进一步解释。

我们来看这道题。既然原有的 \(x, y, z\) 系统不便于题中式子的研究，我们可以考虑进行换元。那么怎么换呢？仍然是参考两年前的那个视频。我们的换元实现了一个效果：靠冻结可以把分母定住，只需要考虑一个分子，这样就把分式结构给解消了。我们故技重施，通过换元，能不能配合冻结法，把分母或者分子之一全部㑈结住呢？题中式子的分子的核心是 \(x^{2}+2 y^{2}+z^{2}\) ，分母的核心是 \(x y+3 y z\) 。我们换元之后肯定也是三个参数，冻结也肯定是先冻两个。这样的话，如果我设 \(x y=a, y z=b\) ，再单独整一个别的参数。这样的话，我先冻结 \(a, b\) ，不就把分母冻住了？那我除了 \(a, b\) ，再加一个谁作为主自由度呢？观察刚才的换元，我再把 \(y\) 添回来不就行了？这就相当于原本是 \(x, y, z\) 系统，现在我改成 \(a, b, y\) 系统了。这时候，如何用 \(a, b, y\) 来表示原式子呢？那很简单啊，显然 \(x=\frac{a}{y}, z=\frac{b}{y}\) ，带回去：分母是 \(a+3 b\) ，分子的核心是 \(\frac{a^{2}+b^{2}}{y^{2}}+2 y^{2}\) 。现在我们抛弃 \(x, y, z\) 系统，只看当前的 \(a, b, y\) 系统。显然，目前只有 \(y\) 好研究，而且是非常好研究。冻结 \(a, b\) ，于是我们当前只需要考虑 \(\frac{a^{2}+b^{2}}{y^{2}}+2 y^{2}\) 的最小值，显然由基本不等式，其最小值为 \(2 \sqrt{2\left(a^{2}+b^{2}\right)}\) ，且等号可以成立。这样就消灭了自由度 \(y\) 。

此时式子变成了

\[
\frac{8\left(a^{2}+b^{2}\right)+1}{a+3 b}
\]

当前我们只需要关注它的最小值即可。有了新参数，定义域必须得跟上。注意到当原本的 \(x, y, z\) 遍历整个大于 0的区域时，\(a, b, y\) 也遍历整个大于 0 的区域，因此 \(a, b\) 的取值范围就是 \(a, b>0\) ，信息量没有流失。我们彻底抛弃原题，把思维焦点放到当前位置上来。这个式子的最小值如何求呢？同样的，我们再来一次换元，把分母的影响给彻底消掉。令 \(u=a+3 b\) ，那么 \(u, b\) 系统就可以替代 \(a, b\) 系统。这时，式子的分母是 \(u\) ，分子的核心是 \(b^{2}+(u-3 b)^{2}\) 。同样的，我们冻结 \(u\) ，这时我们只需要关注分子核心 \(b^{2}+(u-3 b)^{2}=10 b^{2}-6 u b+u^{2}\) 的最小值，也就是 \(10 b^{2}-6 u b\)的最小值。那自然我得看 \(b\) 的取值范围啊。因为 \(a, b>0\) ，所以我必须 \(u>3 b>0\) 。反之，只要 \(u>3 b>0\) ，那就可以反解出 \(a=u-3 b>0\) 。这说明所有满足 \(b \in\left(0, \frac{u}{3}\right)\) 的 \(b\) 都在定义域范围内。因此当前的思维焦点就是在 \(b \in\left(0, \frac{u}{3}\right)\) 的情形下求 \(10 b^{2}-6 u b\) 的最小值。这可太简单了，就是个普通的二次函数，对称轴的位置是 \(\frac{3 u}{10}\) ，正好落在 \(\left(0, \frac{u}{3}\right)\) 内，可以取到！

此时二次函数最小值为 \(-\frac{9}{10} u^{2}\) ，带回原式中，原式变成

\[
\frac{8\left(-\frac{9}{10} u^{2}+u^{2}\right)+1}{u}=\frac{4}{5} u+\frac{1}{u},
\]

从而消去了 \(b\) 。最后只需对 \(u>0\) 求这个最小值即可，易如反掌，最小值为 \(2 \sqrt{\frac{4}{5}}=\frac{4}{5} \sqrt{5}\) 。若问最开始的 \(x, y, z\) 取何值时等号成立？那就一步步倒回去，把取等条件逐级写出求解即可。所有的取等条件为

\[
\frac{4}{5} u=\frac{1}{u}, b=\frac{3 u}{10}, u=a+3 b, \frac{a^{2}+b^{2}}{y^{2}}=2 y^{2}, x y=a, y z=b,
\]

逐层求解得到

\[
u=\frac{\sqrt{5}}{2}, b=\frac{3 \sqrt{5}}{20}, a=\frac{\sqrt{5}}{20}, y=\frac{1}{2}, x=\frac{\sqrt{5}}{10}, z=\frac{3 \sqrt{5}}{10} .
\]

题目就彻底做完了。

做个总结。这道题知道待定系数法的话做起来会很快，但是啥也不知道，只是对up 主讲的冻结思想理解深刻的话，也是完全能做的，并且解题速度完全不慢。虽然思维讲解比较耗时间，但是实际的计算量并不大，三次求最值的过程都很简单。因此只要思维方式足够科学稳健，这道题仍然能够迅速解决。这道题与两年前的那道题的共同点就是通过换元以实现式子结构的简化。所谓换元法，就是代数中的重构要素网。相比较而言，几何问题有画图法的帮助，代数问题则更加凸显对式子结构的把握。二者在元问题与预判的层面，思维是一致的。只要掌握好科学的思维方式，让思维作为灯塔引领我们找寻方向，我们就能最大限度发挥好扎实的知识基础，破解难题，摘取高分。

\section*{课后练习}
1．把我们算出的 \(x, y, z\) 代回题千原式，验证其确实等于我们算出的最小值。你认为这个步骤是否符合我们之前讲过的科学即时检查原理？

2．观察题目如下的解法：

\[
x^{2}+2 y^{2}+z^{2}=\left(x^{2}+\frac{y^{2}}{5}\right)+\left(\frac{9 y^{2}}{5}+z^{2}\right) \geq \frac{2}{\sqrt{5}} x y+\frac{6}{\sqrt{5}} y z,
\]

于是原式

\[
\geq \frac{\frac{4}{5}(x y+3 y z)^{2}+1}{x y+3 y z}=\frac{4}{5}(x y+3 y z)+\frac{1}{x y+3 y z} \geq \frac{4 \sqrt{5}}{5},
\]

且不难求解出等号成立条件。\\
上述解法的第一步拆式子的系数是如何想到的？假如我想把 \(2 y^{2}\) 拆成 \(\lambda y^{2}\) 和 \((2-\lambda) y^{2}\) 两项（ \(\lambda\) 待定），实现上述求解，你能否根据分式中的系数，计算出 \(\lambda\) 的值？

\section*{6．8．2 思想试验辅佐看穿本质}
题目（江苏 2020 数学夏令营）：\\
已知正实数 \(a, b, c\) 满足 \(a b+b c+c a=16(a \geq 3)\) ，则 \(2 a+b+c\) 的最小值为 \(\qquad\) －

题干很短，就是一个简单粗暴的不等式小题，有高一的不等式知识储备就可以做。在讲这道题之前，我先展示一个本题的表演性做法，供大家欣赏：

\[
2 a+b+c=(a+b)+(a+c) \geq 2 \sqrt{(a+b)(a+c)}=2 \sqrt{a^{2}+a b+b c+c a}=2 \sqrt{a^{2}+16} \geq 10
\]

等号成立当且仅当 \(a+b=a+c\) 与 \(a=3\) 同时成立，解出 \(b=c=2\) ，题目做完。\\
是不是感觉很简单，两行就完事了。但是，我自认为脑袋愚笨，考场上绝对想不到这么巧妙的，所以此方法

只给大家展示一下，供头脑聪明的同学参考，觉得自己完全能在考场想到此方法的同学可以离开了。下面，我们正式来讲稳定求解此题的思维方式。

不等式求最值问题也要分析自由度，因为这关系到最后所求式会是单参数的函数，还是多变量的东西，后者则需要使用冻结思想。这道题有三个自由度 \(a, b, c\) ，配了一个牵制等式 \(a b+b c+c a=16\) ，外加一个 \(a \geq 3\) 的范围约束，实际自由度是 2 ，这已经说明必须采用冻结变量的思维方式了，不熟悉的同学可以去＂数学的干货＂里翻一下讲多变量问题的视频。冻结变量专门用于多自由度的问题，将多个自由度中的某一个先冻结住，视为静止，这样需要观测，研究的自由度就暂时少一个，更容易研究；这些东西研究明白（让它们先达到最值状态，从而消去这些自由度）以后，再解冻原来的冻结变量，从而简化问题思考。

这道题既然是两个自由度，那就得先冻结一个，冻结 \(a, b, c\) 中的谁呢？太明显了，就只有 \(a\) 单独配了个范围，不冻它冻谁。我先假定 \(a\) 是不小于 3 的一个常数，视为静止，这样的话，已知就相当于 \(a(b+c)+b c=16\) ，现在得先求 \(b+c\) 的最小值（所求之中 \(a\) 是常数，先不理它）。把元问题隔离出来，发现这就是一个已知 \(b+c\) 和 \(b c\) 的关系，求 \(b+c\) 最小值的问题。于是思维焦点先放上来：怎么处理这种问题呢？

碰到不太清晰的问题，设问是比较好用的开头方式：我自问一下，什么样 \(b=c\) 可以？什么叫＂可以＂？这是个很重要的问题，含义是：这样的 \(b+c\) 可以让 \(a(b+c)+b c=16\) 有合适的 \(b, c\) 的解（注意这里的思维过程，设问＂可以吗＂，然后很自然就主动想，＂可以＂一词的含义，这样主动引导自己思考是非常重要的）。有同学可能会问，引导了也想不到怎么办？这时候，就要试着去做一点思想试验，比如：随便想一个，\(b+c=5\) 行不行？这样的话，就相当于名面上思考一组方程组 \(\left\{\begin{array}{l}b+c=5 \\ b c=16-5 a\end{array}\right.\) 有没有正实数解了，这肯定都能看出来了。所以就看出来是要研究方程组问题了。思想试验的方法常常和设问配合，能够更好的引导你的思考。

既然是这样，那我就把刚才的思想试验一般化：设 \(b+c=k\) ，就要讨论什么样的 \(k\) ，使得方程组 \(\left\{\begin{array}{l}b+c=k \\ b c=16-a k\end{array}\right.\)有正实数解。那显然得 \(0<k<\frac{16}{a}\)（保证这俩都大于 0 ），不仅如此，还得满足基本不等式

\[
4 b c \leq(b+c)^{2} \Leftrightarrow 4 a k+k^{2} \geq 64
\]

注意：在都是正数的前提下，基本不等式成立就可以保证方程组一定有解了，所以这是充要替换，信息量都是不流失的。这样的话，如果把 \(a\) 看作静止，这已经把 \(k\) 描述的足够清楚了。为了更清楚一些，我们干脆直接把不等式解出来：

\[
(k+2 a)^{2} \geq 64+4 a^{2} \Leftrightarrow k \geq 2\left(\sqrt{a^{2}+16}-a\right)
\]

所以 \(k\) 的最小值已经得出了。还需注意的是，若这个最小值能取到，则必须保证

\[
\frac{8}{a}>\sqrt{a^{2}+16}-a
\]

这并不难说明，因为右边等于 \(\frac{16}{\sqrt{a^{2}+16}+a}\) ，所以显然成立。\\
这样的话，我们得到了 \(k\) 的最小值 \(2\left(\sqrt{a^{2}+16}-a\right)\) ，现在可以解冻 \(a\) 了：我任取一个不小于 3 的 \(a\) ，那么在这个 \(a\) 之下，\(b+c=k\) 的最小值是 \(2\left(\sqrt{a^{2}+16}-a\right)\) ，所以求整个式子 \(2 a+(b+c)\) 的最小值，就肯定得先让 \(b+c\) 直接站在 \(2\left(\sqrt{a^{2}+16}-a\right)\) 那里，这时候，所求式变成了 \(2 a+2\left(\sqrt{a^{2}+16}-a\right)=2 \sqrt{a^{2}+16}\) ，注意 \(a\) 可以随便遍历 \(a \geq 3\)的范围，那这玩意的最小值自然就是 \(2 \sqrt{3^{2}+16}=10\) 喽！于是题目做完。

做个简单的总结。这道题的思维方式很直白，都是讲过的东西。最主要的是冻结，这个处理多变量问题也是我们老生常谈的了。而本题还体现了一种思维方式就是：思想试验和设问的配合。设问是用来引导大家主动思考的，而如果设问了也不知道该干啥的时候，就可以试着做做思想试验，做法就是先随便猜一个具体的答案，再设问一下这个行不行？这时候，可以主动思考一下这个具体问题你会怎么处理，其本质是什么？具体问题永远都是比带着更多自由度的抽象问题简单的，从具体问题人手，通过具体例子来进行思考有助于更好的理解原本的抽象题意，更好的彰显＂主动思考＂这四个字的含义，从而帮助我们更好，更稳定的思考清楚接下来要干什么。

\section*{课后练习：}
此题的原题干所有已知都不变，只把所求的 \(2 a+b+c\) 改成 \(a+b+c\) ，请仿照视频的做法，完成这道改编题。你认为这道改编题能否使用与视频最开头的表演性做法类似的方法完成？

\section*{6．8．3 XXXX}
\section*{6.9 综合题目的思维练习}
\section*{6．9．1 如何识别隐藏的最值条件？}
\section*{题目1：}
已知在 \(\triangle A B C\) 中，角 \(A, B, C\) 所对的边分别是 \(a, b, c\) ，记 \(\triangle A B C\) 的面积为 \(S\) 。若满足 \(4 \sqrt{3} S=a^{2}+b^{2}+c^{2}\) ，则 \(\frac{2 a}{3 b+c}\) 的值为 \(\qquad\)。

题目2：\\
实数 \(a, b, c\) 满足 \(e^{a-b+c}+e^{a+b-c}=2 e^{2}(a-1)\) ，则 \(\frac{a b c}{a^{4}+b^{4}+c^{4}}\) 的最大值为 \(\qquad\)。

这两道题目都反映了一类新题型。我们先从第一题开始分析，之后上升到一般情况，为大家进行这类题型的讲解。

第一题。显然整个问题都是大小无关的。已知条件是一个等式，按照我们以往的经验，这一个等式应当减少一个自由度，因此实际自由度似乎应当是 1 。但是题目却要求我们输出一个确定的数值，因此回顾过去的经验的话，只有不变量问题符合这种情况。但是我们看已知：

\[
4 \sqrt{3} S=a^{2}+b^{2}+c^{2}
\]

这个条件是非常对称，如果我们轮换三边 \(a, b, c\) ，那么这个条件不会变。例如，假如 \(a, b, c=3,4,5\)（国写的）符合条件，那么 \(a, b, c=4,3,5\) 也得符合条件。这样一来的话，似乎所求值 \(\frac{2 a}{3 b+c}\) 是会变化的！！因此，从直观上，这个问题就不太像是个不变量问题，倒像是个真正的静止问题。

我们先沿用过去的要素结构的理念去尝试做一做这个题。题目只涉及到一个三角形，而三角形的主自由度选法常见的就两种，SAS 和 AAS（ASA）。已知涉及到三边和面积，边比较多，而面积又刚好可以用 SAS 很方便的控制出来（这是预判过程），因此我们先尝试选 SAS 作为主自由度。因为已知是对称的，选哪个 SAS 并不影响，我

们不妨选 \(a, b, C\) 作为主自由度。\\
既然如此，那 \(c, S\) 就都是次级要素了，我们要把它们用主自由度表示出来。非常简单：

\[
S=\frac{1}{2} a b \sin C
\]

以及

\[
c^{2}=a^{2}+b^{2}-2 a b \cos C .
\]

代回去得到

\[
a b(\sqrt{3} \sin C+\cos C)=a^{2}+b^{2}
\]

这里要提醒大家，如果你发现一个牵制式可以等价变形，使得等号左边和右边没有重叠的主自由度，那么一定要试着这样变形一下，这样做可以把牵制式的结构看得更清晰，更好地看清楚牵制式究竟牵制了什么。比如本题，三个主自由度 \(a, b, C\) ，如果我们把等号左边的 \(a b\) 除到右边，那么

\[
\sqrt{3} \sin C+\cos C=\frac{a^{2}+b^{2}}{a b}
\]

左边只有 \(C\) ，右边只有 \(a, b\) ，相当于原本的三个主自由度被泾渭分明地分成了两派，这样的话牵制条件就像是泾河与渭河之间架了一座桥，而不是彼此纠缠在一起，结构自然更清晰。

清晰了又如何呢？这样做思想试验更方便了。比如现在不是要求 \(\frac{2 a}{3 b+c}\) 吗？这个式子太洈异了，不好研究，那我可以先来个思想试验。大小无关我就当 \(a=1\) ，然后牵制式右边就是 \(b+\frac{1}{b}\) ，我倒要看看我取个 \(b\) 会导致怎样的 \(c\) ，以至于 \(3 b+c\) 还能是个定值？取 \(b=2\) ，上来就发现

\[
\sqrt{3} \sin C+\cos C=\frac{5}{2}
\]

居然没有解！！其他 \(b\) 呢？有了刚才的思想试验的信息提示，我们很快就发现右边 \(b+\frac{1}{b} \geq 2\) ，但是左边

\[
\sqrt{3} \sin C+\cos C \leq 2
\]

等号想取到只能这俩都取等，也就是 \(b=1, C=\frac{\pi}{3}!!\) 好家伙，原来这个等式带来的效果居然是减少两个自由度！不管怎么说，这一个式子带来的效果还真是静止掉了两个自由度，三角形只能是等边三角形，那题目最后的答案也就很简单了，就是 \(\frac{1}{2}\) 。

从这道题我们可以看出，有的时候一个式子也能静止掉多个自由度。这个道理其实很好理解，举个更简单的例子：考虑有两个主自由度 \(x, y\) 的系统，它们控制了一个次级要素 \(z=x^{2}+y^{2}\) 。那么牵制 \(z=1\) 自然就减少一个自由度，而且显然，加了这个牵制的单自由度系统就等价于一个圆。但是，如果牵制改成 \(z=0\) 呢？显然圆就退化成了一个点，这时候系统等价于 \(x=y=0\) 。究其原因，这个取值刚好使 \(x^{2}+y^{2}\) 取到最小值，单自由度的线退化成了零自由度的点。这是一种自由度出现问题的＂怪胎＂，可能使得＂一个等式减少一个自由度＂的规律失效。Up主将这类问题称为＂临界问题＂。关于这种导致退化的＂怪胎＂，本期视频最后会给大家作进一步的一点科普，供感兴趣的同学了解。

这种临界问题的存在可能会让大家感到担心，以往判断自由度的方法还能放心大胆地使用吗？事实上，数学研究表明（见视频末尾的科普），这种＂怪胎＂的存在是极为稀少的，几乎所有的情形下，原本判断自由度的方法都是生效的。而在高考考场上，专门考察这类＂怪胎＂情况的题目也非常少见，绝大部分高考真题套卷中完全没有临界问题，只在少部分套卷中可能偶尔出现一道题左右。不过，在考场上遇到的概率绝对不是 0 ，因此我们还是有必要说一下面对这类题目的甄别与处理办法的。

首先，如何甄别这类题目？既然是自由度出现问题，那么如果我们用以往的一般规律判断自由度的时候，如

果直接判断是静止，那显然就不是临界问题。在此之上，临界问题通常在题干中就会给予我们一些提示。第一，临界问题的题干中给出的牵制条件或者所求结果，通常长的会比较奇怪，乍一看似乎很难通过已知来获取什么漂亮的信息，或者很好的控制出所求。比如第一题的所求，不管怎么看都像是个静态题目问出来的问题。再比如第二题，已知长得就很难看，三个自由度还带着指数，所求的连四次方都出来了，已知和所求看上去风马牛不相及。第二，既然需要通过最值取等条件来多卡掉自由度，那就得有明显的可以求最值的结构。比如第一题，三角函数和齐次的单自由度式子，一个是辅助角，一个是基本不等式，明显都是容易求最值的结构。遇到一个式子里掺杂了这样的结构，就要提高警 F ，会不会是临界问题。

当然，有时候我们也不见得能立刻甄别出来。比如我把第一题改成单选：

\[
\frac{2 a}{3 b+c}
\]

可以取什么值？然后配四个选项，那你就很难用上述方式直接甄别了。在这种情况下，我们仍然要利用好要素结构的思维方式。从刚才第一题的过程可以看出，在不知道题目究竟是什么类型的前提下，我们仍然使用控制的理念，选取主自由度去充要替换已知条件。如同不变量问题最后一定会消去全部主自由度一样，临界问题只要我们通过主自由度去控制清楚了，那么也一定可以通过求最值判断恰好取到来解出。因此，此时的难点就是甄别出可以求最值的结构来。这时候，我们在讲第一题时提到的，尽量使得式子两边被不同的主自由度控制，就是一种很好的习惯。如果主自由度们相互纠缠，那么结构就很不清晰，这时候如果可以做到前面说的＂泾渭分明＂，让等式左右都被不同的主自由度控制，那么式子的结构就会更加清楚。这是让问题结构清晰的很好的结构性思维的习惯，不仅适用于甄别临界问题，十分建议养成。

再来看第二个问题，给了三个主自由度和一个牵制条件，牵制式同样很奇怪，所求也让人摸不着头脑。但是如果我们像刚才那样有这种使之泾渭分明的意识，注意左边的两个指数 \(a+b-c\) 和 \(a-b+c, a\) 是一样的，\(b, c\) 符号相反，而右边又只有 \(a\) ，因此按照刚才的泾渭分明的原则，我们可以把 \(a\) 全部放到等号一侧：

\[
e^{b-c}+e^{c-b}=\frac{2 e^{2}(a-1)}{e^{a}}
\]

这时，左边显然可以把 \(b-c\) 看成一体，先让 \(b, c\) 控制 \(x=e^{b-c}>0\) ，再让

\[
x+\frac{1}{x}=\frac{2 e^{2}(a-1)}{e^{a}}
\]

结构变得更加清晰。这时，右边的函数很好研究，求导易知其在 \(a=2\) 时取得最大值 2 ，这正好是左边的最小值，因此必须 \(a=2, b=c\) ，这也就是已知的充要替换，直接消掉两个自由度。因此所求的那个丑陃的式子也就变成了

\[
\frac{2 b^{2}}{16+2 b^{4}}=\frac{1}{b^{2}+\frac{8}{b^{2}}}
\]

\((b=0\) 时就是 0\()\) ，那最大值自然就是 \(\frac{\sqrt{2}}{8}\) 。\\
做个总结。我们这期视频讲到了一个等式可能消掉多个自由度的情况，up 主称其为临界问题。这类问题在高考中出现的频率非常低，但是如果真出现了，大家也要能够甄别出来。临界问题的题干通常都会有一些暗示，如丑陋的已知，丑国的所求，乍一看完全搭不上边的已知和所求等。另外，临界问题的牵制式通常都有明显的可以求最值的结构。这些信息可以让我们甄别出临界问题。

即便是没有通过读题感觉出来，只要大家把握好要素网的结构思维，踏踏实实选好主自由度去表达清楚牵制式，那么上面的那些暗示就会以更加清楚的姿态展现在我们面前。一个重要的甄别方式就是让式子左右分别被

不同的主自由度控制。这是一个一般的习惯，目的是让式子结构更加清晰，而不仅仅是为了甄别临界问题。有了这种习惯，我们就能更加轻松地识别出临界问题，破解它里面的陷阱。

补充科普：系统的临界点（与高考毫无关系！）\\
我们考虑一个一般性的系统，它的自由度为 \(n\) 。我们为它选取一组主自由度为 \(X_{1} \cdots X_{n}\) 。\\
刚才我们讨论的牵制式其实就是一个关于 \(x_{1} \cdots x_{n}\) 的等式，而等式就可以把含 \(x_{1} \cdots x_{n}\) 的东西挪到同一边，让另一边只剩下常数。如第一题，我可以等价为

\[
\frac{a^{2}+b^{2}+c^{2}}{S}=4 \sqrt{3}
\]

因此，牵制我们可以理解为：\(x_{1} \cdots x_{n}\) 控制某个次级要素 \(z\) ，牵制就是存在某个常数 \(c\) ，并考察方程 \(z=c\) 。所谓控制就可以理解成 \(z\) 是 \(X_{1} \cdots x_{n}\)（也就是系统的某一个具体状态）的函数（这是多个变量的函数，称为多元函数），因此我们可以记为 \(z=f\left(x_{1} \cdots x_{n}\right)\) 。而当我们加上一个牵制时，我们就相当于要研究这个大系统的一个子系统。比如考虑 \(n=2\) 的情形（主自由度记为 \(x, y\) ），那么大系统就是整个平面。当 \(f(x, y)=x^{2}+y^{2}\) 时，加上 \(z=1\) 的牵制就是考虑平面的子集

\[
\left\{(x, y) \mid x^{2}+y^{2}=1\right\}
\]

也就是把这个单位圆作为子系统。这个子系统是曲线，也就是说这个子系统的自由度为 1 。而 \(z=0\) 的牵制对应的子系统就只是原点，点的自由度就是 0 。我们通常的规则是：对于自由度为 \(n\) 的系统，一个方程 \(z=c\) 对应的子系统的自由度是 \(n-1\) 。而临界问题就是打破这个规则的情形。

自然的问题是：对于一个给定的控制式 \(z=f\left(x_{1} \cdots x_{n}\right)\) ，取什么样的 \(c\) 会打破上述规则？这样的 \(c\) 多不多？显然，如果 \(c\) 是极大值或者极小值的话，那么很有可能会导致上述情况。我们在学导数的时候，知道取极大值和极小值时导数一定为 0 ，但反之则不然。这里我们把条件进一步放宽，考虑更大的范围：也就是导数为 0 的点。不过现在 \(Z\) 是多个变量的函数，它的求导方法我们还没学，因此我们仅仅直观理解一下。如果在 \(X_{1} \cdots x_{n}\) 处＂导数为 0 ＂的话，那就像个＂小山包＂鼓了起来，鼓出了一个局部小山峰。当我们站在山顶的时候，沿着各个方向下山，那么起点的导数都为 0 。这里的＂各个方向＂就相当于经过峰顶的一条条直线。因此我们可以如下定义＂导数为 \(0 "\) ：考虑所有经过 \(\left(x_{1} \cdots x_{n}\right)\) 的自由度为 1 的子系统，若给这个子系统选取一个主自由度 \(t\) ，则恒有

\[
\left.\frac{d}{d t}\right|_{t=0} f\left(x_{1}(t) \cdots x_{n}(t)\right)=0
\]

满足上述条件的 \(\left(x_{1} \cdots x_{n}\right)\) 称之为系统的临界点（critical point）。显然所有的峰头或者谷底都是临界点。而我们关心的是那些临界点处的取值。如果 \(\left(x_{1} \cdots x_{n}\right)\) 是临界点（有高风险是极值点的地方），则 \(c=f\left(x_{1} \cdots x_{n}\right)\) 就称为临界值（critical value）。我们前面说的那些导致临界问题的 \(c\) 正是这些临界值。例如

\[
\frac{a^{2}+b^{2}+c^{2}}{S}=4 \sqrt{3}
\]

那么 \(4 \sqrt{3}\) 就是一个临界值。\\
不难看出，所谓临界值，就是存在临界点与之对应的实数值，用这样的实数值造出来的牵制条件就让我们很不放心。与之对应的，如果一个实数值不是临界值，也就是说它不是任何临界点的函数值（当然也有可能不在值域内，比如 \(f(x, y)=x^{2}+y^{2}\) ，我取 -1 就不在值域内），则称其为正则值（regular value）。我们希望的是，如果 \(c\) 是正则值，则 \(z=c\) 对应的子系统的自由度能是 \(n-1\) ，也就是恰好减一个自由度。事实上，确实如此：

\section*{定理 6.1 （正则值引理，简化版）}
记号如上。若 \(c\) 是 \(f\) 的正则值，且 \(c\) 在 \(f\) 的值域内，则 \(z=c\) 对应的子系统的自由度一定是 \(n-1\) 。

这个定理就告诉我们，去掉那些导数为 0 的讨厌的临界值（比如那些极大值或者极小值），剩下的牵制条件都不会出现自由度上的问题。那么，这种讨厌的临界值多不多呢？事实上，著名的 Sard 定理告诉我们：临界值很少，少到几乎可以忽略不计：

\section*{定理 6.2 （Sard 定理，简化版））}
若控制函数 \(f\) 足够好（无穷次可导），则几乎所有的实数都是它的正则值。

上述的两个定理都是微分几何中关于子流形的基本定理，它们的严格版本叙述如下：

\section*{定理 6.3 （正则值引理）}
设 \(M, N\) 都是光滑流形，其维数分别为 \(m, n, f: M \rightarrow N\) 是光滑映射，\(p \in N\) 满足 \(f^{-1}(p) \neq \varnothing\) 。若对 \(\forall x \in f^{-1}(p), f\) 在 \(x\) 处切映射为满射，则 \(f^{-1}(p)\) 是 \(M\) 的正则子流形，维数为 \(m-n\) 。

注这里的维数相当于系统的自由度数目；＂切映射为满射＂是正则值的一般定义，相当于 \(n\) 个牵制式有效且彼此独立；最后的＂维数＂就是子系统的实际自由度数目，即为 \(M\) 的自由度数目减去独立牵制式的数目

\section*{定理 6.4 （Sard 定理）}
设 \(M, N\) 都是光滑流形，\(f: M \rightarrow N\) 是光滑映射，则 \(f\) 的临界值之集为 \(N\) 的零测集。

注零测集直观上可以理解为少到在概率意义下可以忽略不计的集合，比如 \([0,1]\) 的子集 \(\left\{0, \frac{1}{2}, 1\right\}\) 相比于整个 \([0,1]\) 而言太少了，少到在 \([0,1]\) 内随便取一个数，它正好是 \(\left\{0, \frac{1}{2}, 1\right\}\) 中的元素的概率为 0 。而＂临界值少到可以忽略不计＂就相当于正则值多到几乎全是

\section*{6．9．2 要素关系建构与元问题预判练习}
题目一：\\
已知平面向量 \(\vec{a}, \vec{b}, \vec{c}\) 满足 \(|\vec{a}|=2,|\vec{b}|=3,|\vec{c}|=1\) ，且 \((\vec{a}-\vec{c}) \cdot(\vec{b}-\vec{c})=-4\) ，则 \(|\vec{a}-\vec{b}|\) 的取值范围是 \(\qquad\)。

题目ニ：\\
已知椭圆 \(C: \frac{x^{2}}{a^{2}}+\frac{y^{2}}{b^{2}}=1(a>b>0)\) 的左右焦点分别为 \(F_{1}, F_{2}\) ，离心率为 \(e\) 。点 \(P\) 在椭圆 \(C\) 上，连接 \(P F_{1}\) 并延长交 \(C\) 与点 \(Q\) ，连接 \(Q F_{2}\) 。若存在 \(P\) 点使得 \(|P Q|=\left|Q F_{2}\right|\) 成立，则 \(e^{2}\) 的取值范围是 \(\qquad\)。

这两道题目都需要构建一个顺畅的要素关系网。我们之间讲要素关系的时候就说过，构建出好的要素关系网是顺利做题的基础，因为这直接决定我们后续要算什么，调用什么元问题，计算体感如何。而构建好的要素网，最好的办法就是利用画图法，带着工程师的意识去操作。画图法就是考虑将题干描述的几何图像画出来，思考如何画能更加顺畅（多控制，少牵制），更容易画准确。这样的画图方法给出的要素生成关系就是易于列式的。后续再根据要素网里具体的控制关系，调用相应的元问题去计算即可。

要素关系网是一个抽象的观念，它只描述这个系统里谁＂能够＂控制谁，而不关心具体怎么控制。具体到怎么算则是需要把元问题展开，把具体已知的数据代进去计算的。因此，建立要素网的过程只是描述这个系统本身的结构罢了。这也传达出一个要点：建立要素网应该关注的是题目中的要素结构，理清题干中的系统的结构是做题的第一步。我经常说＂控制是更根本的＂，控制这个理念本身才是最核心的，大家做题需要先去关注题干中的系统的结构，不然的话做题只是在到处乱碰，没有目标性，没有踏实感。

既然控制的理念才是最根本的，我们在做更难的题目的时候，尤其是想冲击 130 以上的高分的时候，就需要让思维更加灵活，在动笔算之前要重点关注题目中的控制关系，去理顺它的结构。要素关系的语言就是说，系统的起始的主自由度手握操纵杆，去控制其他要素，效果是控制出题目中的重要事物，例如所求，例如牵制中的重要参量，例如多级计算中的中间量。这个控制关系本身是做题的基本，我们需要先从题目中准确识别这个事物。

我们先看题目一。我们先不必关心已知给的数字怎么参与计算，而先关注基础的要素关系。显然基础要素就是 \(\vec{a}, \vec{b}, \vec{c}\) ，牵制涉及到 \((\vec{a}-\vec{c}) \cdot(\vec{b}-\vec{c})\) ，是由三个向量的信息就能直接控制的参数；但是所求涉及到 \(|\vec{a}-\vec{b}|\) ，只是由 \(\vec{a}, \vec{b}\) 的信息控制。从这里就能感觉到，这个题要研究的系统的核心肯定是 \(\vec{a}, \vec{b}\) 构成的，\(\vec{c}\) 至多是陪因，所以如果要选取主自由度的话，它一定需要反应 \(\vec{a}, \vec{b}\) 的信息。再来看基础要素给出的信息，三个模长都是确定的，所以 \(\vec{a}, \vec{b}\) 系统只剩一个自由度了，所求 \(|\vec{a}-\vec{b}|\) 显然具备彻底控制死 \(\vec{a}, \vec{b}\) 的效果，因此题目的研究对象就明朗了：\(\vec{a}, \vec{b}\) 这个单自由度系统的主自由度的取值情况是什么？我强调过很多次，所求具体是谁并不重要，关键是要素关系，动静与控制的观念下问题本质是什么。

现在已经明确题目要研究的就是 \(\vec{a}, \vec{b}\) 单自由度系统，那我们首先把一个控制杆交给它。那这个 \(\vec{c}\) 是干嘛的？我们来看题：首先 \(\vec{c}\) 的模确定，所以至少 \(\vec{c}\) 目前只有一个自由度控制着。而另一个牵制：\((\vec{a}-\vec{c}) \cdot(\vec{b}-\vec{c})=-4\) ，显然，如果 \(\vec{a}, \vec{b}\) 的控制杆停下，那么这时候这个牵制就能把 \(\vec{c}\) 彻底静止化。所以我们把要素关系抽象出来：\(\vec{a}, \vec{b}\) 手握一个控制杆，\(\vec{c}\) 是一个单自由度的东西，有一个和此控制杆相关的牵制能够定死 \(\vec{c}\) ，求控制杆的活动范围。

上面这个关系网是一个纯抽象的东西，不含任何计算，这就是我们需要的东西。解题第一步完成，我们就可以根据这个关系网去调用元问题，开展计算。首先，我们来依据要素结构来识别问题类型。一个控制杆，一个单自由度参数，一个控制杆之下的牵制定死单自由度参数，问控制杆的范围，这是什么类型的问题？从脑海里调用元问题，识别出这是方程有解类问题，就和什么 \(x^{2}=a\) 有解，求 \(a\) 范围这种问题一样。再进一步调用元问题的处理方法：那就是先假定控制杆静止，列出方程，充要替换出有解的条件，再解冻控制杆，求解有解条件的不等式（大家看好了，元问题就是这么用的）。于是开始按照上面的蓝图展开操作。先让控制杆静止，也就是让 \(\vec{a}, \vec{b}\) 系统作为一个静止系统，相当于张开一个确定的三角形，主自由度因为方程还没列，所以先不急于选，看看方程怎么列方便再考虑选取。方程是 \((\vec{a}-\vec{c}) \cdot(\vec{b}-\vec{c})=-4\) ，对应元问题：向量等式变方程，自然是建系。怎么建系？开始预判，\(\vec{c}\) 肯定是

设 \((x, y)\) ，附带约束 \(|\vec{c}|=1\) ，于是选起点为原点就可以。而牵制方程需要 \(\vec{a}, \vec{b}\) 坐标，而且有了坐标，直接点乘，不含根号，式子形式肯定不坏，所以选好坐标系即可。选好原点的话，肯定靠位置无关把 \(\vec{a}\) 或者 \(\vec{b}\) 先静止，那另一个就是绕原点转了，主自由度给一个夹角即可，比如设 \(\vec{b}=(3,0)\) ，再设 \(\vec{a}=(2 \cos \theta, 2 \sin \theta)\) ，则 \(\theta\) 就是操纵杆，或者不会参数方程的话设 \(\vec{a}=\left(x_{0}, y_{0}\right)\) ，附带约束 \(x_{0}^{2}+y_{0}^{2}=4\) ，则 \(\left(x_{0}, y_{0}\right)\) 就是操纵杆。我们就选大家容易想到的第二种设法。

先假定 \(\left(x_{0}, y_{0}\right)\) 静止，于是算出

\[
\vec{a}-\vec{c}=(3-x,-y), \vec{b}-\vec{c}=\left(x_{0}-x, y_{0}-y\right),
\]

再点乘得

\[
(3-x)\left(x_{0}-x\right)-y\left(y_{0}-y\right)=-4,
\]

展开得

\[
x^{2}+y^{2}-\left(x_{0}+3\right) x-y_{0} y+3 x_{0}+4=0
\]

而 \((x, y)\) 自带约束 \(x^{2}+y^{2}=1\) ，所以方程全都列完了，问题充要替换成这个方程组有解。\\
其实大家应该已经看出来了，这就是两个圆，要它们相交。不过把两个方程隔离出来，单独看的话，很明显要先把第一个方程变成

\[
-\left(x_{0}+3\right) x-y_{0} y+3 x_{0}+5=0,
\]

这样和第二个方程组合起来，得到的方程组与原方程组等价，而且是直线与 \(x^{2}+y^{2}=1\) 的联立，非常简单。这时候直线和圆有交点这种元问题就非常基础：

\[
\frac{\left|3 x_{0}+5\right|}{\sqrt{\left(x_{0}+3\right)^{2}+y_{0}^{2}}} \leq 1
\]

隔离的思维习惯非常普遍，大家一定要养成。\\
充要替换已经结束了，这个式子只与原始的控制杆 \(\left(x_{0}, y_{0}\right)\) 有关了，配上 \(x_{0}^{2}+y_{0}^{2}=4\) 的约束。所以直接完成最后的求解即可。把现在的信息全放在一起（一等式一不等式一所求），不难看出把 \(y_{0}\) 消去即可：

\[
\frac{\left|3 x_{0}+5\right|}{\sqrt{\left(x_{0}+3\right)^{2}+4-x_{0}^{2}}} \leq 1
\]

解这个不等式得到

\[
-2 \leq x_{0} \leq-\frac{2}{3} .
\]

最后

\[
|\vec{a}-\vec{b}|^{2}=\left(x_{0}-3\right)^{2}+y_{0}^{2}=13-6 x_{0} \in[17,25]
\]

那答案自然就是 \([\sqrt{17}, 5]\) 。

可以看到，整个题目的思考过程就是先建立抽象的要素关系，根据要素关系识别问题类型，调用相应的元问题处理方法，然后顺着关系网，对每一步的计算做预判，确认计算可行性，然后按部就班完成计算即可。在要素关系建立，元问题预判都做好的前提下，完成整个题目就是砍瓜切菜。

我们再来看题目二。首先是背景的椭圆 \(C\) ，然后上面一个单自由度动点 \(P\)（不过它在后文来看是被封装了）， \(P\) 严格控制 \(Q\) 。这时题目给了牵制 \(|P Q|=\left|Q F_{2}\right|\) ，一个方程肯定可以求解，所以 \(P\) 是可以被求出来的。背景粗圆

在大小无关下只有 \(e\) 这一个自由度。而问的就是＂存在 \(P\)＂，即方程有解。看上去和上一题的单自由度系统的一个方程有解是一回事。

下面尝试去调用元问题，并进行预判。需要转化的方程是 \(|P Q|=\left|Q F_{2}\right|\) ，调用元问题：等腰如何充要替换成式子？显然有两条路，一个是直接算边长，一个是取底边 \(P F_{2}\) 中点 \(N\) ，式子变成 \(Q N \perp P F_{2}\) 。先假设粗圆静止，预判第一条路：如果把主自由度分给 \(P\) ，那么算 \(|P Q|\) 需要有 \(P Q\) 的斜率（这个容易），也需要 \(Q\) 坐标（这个元问题是已知一交点求另一交点，需要用解点法）；还需要 \(\left|Q F_{2}\right|\) ，这个得用 \(Q\) 坐标。不过如果把主自由度给 \(Q,\left|Q F_{2}\right|\) 的计算就容易一些。如果是第二条路，把主自由度给 \(P\) ，那么只需要通过解点算 \(Q\) 的坐标即可，比较便利。从预判来看，第二条路更加舒适，不过，既然涉及到解点，而且还是带着参数的解点，计算量肯定不小。平时练习处理过含参数解点这一元问题的同学应该会对自己的处理时间有一个了解，这个计算量已经达到了偏难大题的水平。

如果考场上预判到这里，而且想不到别的方法的话，已经可以跳过了。我们重新思考，这个方法不好，能不能换一个？换方法等价于换要素网，所以重新审视要素关系网。这个要素网显然不符合画图法的标准，等腰三角形根本画不准确，只能嗐点一个 \(P\) 来作为示意图，那有没有办法画准确呢？尝试画一下就明白了。我先画满足 \(|P Q|=\left|Q F_{2}\right|\) 的等腰三角形 \(P Q F_{2}\) ，就可以了。那椭圆怎么办呢？如果先画了 \(P Q F_{2}\) ，那么画椭圆就是画 \(F_{1}\) ，这只需要在 \(P Q\) 上找满足 \(\left|P F_{2}\right|+\left|P F_{1}\right|=\left|Q F_{2}\right|+\left|Q F_{1}\right|\) 的 \(F_{1}\) 即可，这个量一下长度就行了，很简单。这样的话整个图像都画出来了，而且控制很顺畅，是一个好的关系网！

这时候我们梳理关系网，起始是等腰三角形 \(P Q F_{2}\) ，大小无关的话只有一个自由度，然后它控制 \(F_{1}\) ，进而控制粗圆，控制离心率。那么这个问题就被重构成了已知等腰三角形，求离心率的取值范围，和一开始的方程有解问题差别很大。有点怪，那么这个时候就思考一下，这个重构是充要替换吗？从逻辑上理解一下，原题是考虑所有的椭圆，使得存在等腰能按题干放进去，求对应的哪些离心率符合条件；新题是考虑所有的等腰三角形，问控制出来的椭圆的离心率能遍历到哪些值。对于原题的范围，那样的椭圆肯定能放进合适的等腰三角形 \(P Q F_{2}\) ，所以原题的范围对新题都合适；反之，新题中给一个等腰三角形，控制出来的椭圆显然也符合原题。所以这两个问题确实是等价的。（大家注意，刚才的这个分析在考场上需要去做，这样才确保你的思维严谨）

刚才做的这个题目重构，实际上就是以要素关系网为灵魂的。这里的做题逻辑是：我先通过画图法认识到了一个新的构建要素关系网的办法，然后对于这个新要素网，我不知道它是不是和原题等价，所以我做了一点逻辑分析，发现确实是等价的。动静视角和要素关系网让我们具备了重构题目的能力，而充要替换的逻辑让我们能够相信自己的判断，踏实地向前走。

现在可以开始算了。首先大小无关，不妨设 \(P Q=1\) 。下面分配一个主自由度给 \(P Q F_{2}\) 。分配谁呢？可以分配给 \(\angle P Q F_{2}\) ，也可以分配给 \(\left|P F_{2}\right|\) 。我们来做一下计算预判。求的是离心率，也就是椭圆的 \(a\) 和 \(c\) 。而

\[
a=\frac{\left|P F_{2}\right|+|P Q|+\left|Q F_{2}\right|}{4}=\frac{2+\left|P F_{2}\right|}{4}, c=\frac{\left|F_{1} F_{2}\right|}{2}
\]

所以不管主自由度分配给谁，\(a\) 都好控制，分配给角的话稍微麻烦点。但是 \(c\) 则是需要控制 \(\left|F_{1} F_{2}\right|\) ，这个相当于等腰三角形腰上取一个定点 \(F_{1}\) 再算距离，隔离出这个元问题，根据要素关系思考一下，就能看出处理方法肯定就是先控制 \(\left|Q F_{1}\right|\) ，然后放在 \(\triangle Q F_{1} F_{2}\) 里用余弦定理。所以如果主自由度分配给 \(\angle P Q F_{2}\) 的话，则需要先表示 \(\left|Q F_{1}\right|\) ，这肯定需要 \(\left|P F_{2}\right|\) 参与，但 \(\left|P F_{2}\right|\) 也需要 \(\angle P Q F_{2}\) 来控制，这个控制肯定带根号，于是恐怕会出现根号相加的式子；但是如果主自由度直接分配给 \(\left|P F_{2}\right|\) 的话，则 \(\left|Q F_{1}\right|\) 很好控制，接下来只需要控制好 \(\cos \angle P Q F_{2}\) ，而这个放在等腰三角形里用余弦定理即可，不带根号，更好一些。所以主自由度给 \(\left|P F_{2}\right|\) 更优，把它设为 \(x \in(0,2)\) 。

此时 \(a=\frac{2+x}{4}\) ，于是

\[
\left|Q F_{1}\right|=2 a-\left|Q F_{2}\right|=\frac{x}{2} .
\]

又有

\[
\cos \angle P Q F_{2}=\frac{1+1-x^{2}}{2 \cdot 1 \cdot 1}=1-\frac{x^{2}}{2},
\]

因此

\[
\left|F_{1} F_{2}\right|^{2}=1+\frac{x^{2}}{4}-2 \cdot 1 \cdot \frac{x}{2} \cdot\left(1-\frac{x^{2}}{2}\right)=\frac{x^{3}}{2}+\frac{x^{2}}{4}-x+1=\frac{(x+2)\left(2 x^{2}-3 x+2\right)}{4}
\]

因此离心率

\[
e^{2}=\frac{c^{2}}{a^{2}}=\frac{\left|F_{1} F_{2}\right|^{2}}{4 a^{2}}=\frac{2 x^{2}-3 x+2}{x+2},
\]

定义域 \(x \in(0,2)\) 。最后求范围就不多说了，非常基础的函数求范围问题，注意一个基本技巧是把 \(x+2\) 换成 \(y\) ，这样求解更简单，大家自己去求一下即可，最后答案是 \([8 \sqrt{2}-11,1\) ）。注意，如果没看出可以因式分解约分的话，也没关系，大不了求导，时间会多花一些。不过换完 \(y\) 之后分子一定能够直接提出一个 \(y\) 来，也能够约分约掉。

总结：\\
本次习题课讲解的两道题都体现了要素网，元问题，预判，这三者结合的重要性。要素网体现了题目的结构，是做题的第一步，元问题的理解（基础知识的熟练以及与要素网的匹配）是正确预判与正确解题的基本要求。这也是我们平时一直以来讲解的重点。希望大家多多体会，多多实践，只有养成思维习惯才能发挥出这些思维方式的威力。多去思考问题背后的要素关系，多去用结构的眼光看问题，学知识和复习知识的时候都用结构，控制的眼光去看待，以元问题的形式学习，下笔之前多去预判，这样才能慢慢养成思维习惯，提高自己做题的稳定性与踏实感。

时间规划：\\
题目一：\\
识别出问题的要素结构，认识到是方程有解问题： 1 min\\
选取参数的设法，并大致预判计算可行性：1min\\
充要替换出

\[
\frac{\left|3 x_{0}+5\right|}{\sqrt{\left(x_{0}+3\right)^{2}+y_{0}^{2}}} \leq 1: 2 \mathrm{~min}
\]

完成最后的计算： 2 min\\
总时间： 6 min

\section*{题目二：}
识别出问题的要素结构，认识到是方程有解问题： 1 min\\
选取参数的设法，并预判计算可行性（发现计算量过大）：1min\\
更改方法，用画图法构建新的要素网： 2 min\\
选取参数的设法，并预判计算可行性： 1.5 min\\
完成最后的计算： 3.5 min

总时间： 9 min

\section*{课后练习：}
1．用参数方程的设法 \(\vec{a}=(2 \cos \theta, 2 \sin \theta)\) 重新完成题目一的全部计算。\\
2．尝试用原始的要素关系网，利用解点的思路完成题目二的计算，并与视频中给出的解法进行计算量上的对比。（一共四条可能的道路，试着分别预判一下，算完后和你的预判对照一下，看看符不符合）

\section*{6．9．3 要素网的转化与元问题的应用}
\section*{题目1：}
点 \(F_{1}, F_{2}\) 为椭圆 \(C: \frac{x^{2}}{a^{2}}+\frac{y^{2}}{b^{2}}=1(a>b>0)\) 的左右焦点，\(A, B\) 是粗圆 \(C\) 上关于国圆对称的两点，若椭圆 \(C\) 上存在点 \(P\) ，使得直线 \(A P, B P\) 斜率的平方和为 \(\frac{7}{6}\) ，且 \(\overrightarrow{P F_{1}} \cdot \overrightarrow{P F_{2}}=c^{2}\) ，则椭圆离心率的取值范围是 \(\qquad\)。

\section*{6．9．3．1 题目2：}
已知点 \(P\) 在以 \(F_{1}, F_{2}\) 为左右焦点的椭圆 \(C: \frac{x^{2}}{a^{2}}+\frac{y^{2}}{b^{2}}=1(a>b>0)\) 上，椭圆内（不含椭圆自身）一点 \(Q\) 在 \(P F_{2}\) 延长线（不含 \(F_{2}\) ）上，且满足 \(Q F_{1} \perp Q P\) ，若 \(\sin \angle F_{1} P Q=\frac{5}{13}\) ，则该椭圆离心率的取值范围是 \(\qquad\)。

这两道题目都是符合单选第 8 题难度的小题（已改为填空题），为不同粉丝的投稿，且质量都不错，而且还能反映出一些封装自由度的思想。因此本期合起来，做一个小的习题课练习。

先来看第一题。给出一个粗圆后，题目相继引入了一组对称 \(A, B\)（占 1 个主自由度），以及点 \(P\)（又一个主自由度），然后紧接着给出两个牵制，因此这应当把这两个主自由度全消掉了，因此 \(A, B\) 与 \(P\) 理论是可精确求解的。而题干中描述这些关系用到了＂存在＂这个量词，相当于做了一个封装。在语义上进行一个梳理，其实就是说，用上它给的牵制，确实能求解出 \(A, B\) 与 \(P\) 来（存在 \(P\) 满足一些方程式，和存在 \(x\) 满足 \(x^{2}-2 x=0\) 本质上都是同一种逻辑结构）。所以，我们在语义上重新理解一下这个问题，本质就是一个双变量的＂方程组有解问题＂。

容易引起混淆的一个问题是，这里 \(A, B\) 是一组自由的主自由度吗？\(P\) 是被封装的，如果 \(A, B\) 能自由的话，那么问题的主要矛盾应该是粗圆和 \(A, B\) 共同构成的系统。但是现在给了两个牵制，\(A, B\) 的位置也应当是可解的。换而言之，不是随便给一组 \(A, B\) ，都能存在 \(P\) 的。那么，设问一下，什么样的椭圆才是符合本题的呢？从遍历法的角度来看，给定一个椭圆，遍历所有 \(A, B, P\) ，只要存在一组合法的 \(A, B, P\) 满足题干中的两个牵制，那么这个粗圆就是符合题目的。因此，\(A, B\) 前面也同样可以加上＂存在＂这个量词，语义是不变的。这正是前面提到的＂双变量方程有解问题＂。这里要提醒大家一下，封装命题的识别，谁是约束变元，这是一个语文问题。像本题，虽然 \(A, B\) 前面没有显性加上＂存在＂一词，但是从句子语义的角度来讲，它们前面就是可以用这样一个量词封装。题干本身的叙述是自然语言，其准确语义需要大家通过语文理解进行体会，而写成封装命题则是十分严谨精确的数学语言。有时题干并不会写的特别精确严谨，这时候就需要大家通过设问，遍历法等帮助理解的方式来进行准确的数学语言翻译。这是一个很重要的事情，即便是高考大家也有可能遇到。如果大家在做题时发现有些问题的矛头指向不太清晰的话，就用上述的一套思考方式来分析一下题目中有没有隐藏的约束变元。

回到题目，既然现在 \(A, B\) 与 \(P\) 都是约束变元，而且根据前不久新人经验分享番外篇的约束变元顺序问题，这里的两个约束变元都被＂存在＂封装，因此可以有两种顺序来研究。我们应该按照怎样的顺序呢，那就要分析一

下要素结构了。首先，很容易注意到，两个已知的牵制里面，第二个 \(\overrightarrow{P F_{1}} \cdot \overrightarrow{P F_{2}}=c^{2}\) 只与 \(P\) 自己有关。 \(P\) 只是个单自由度的变元，这个条件足以定死它。因此，我们可以先说：存在 \(P\) 使得 \(\overrightarrow{P F_{1}} \cdot \overrightarrow{P F_{2}}=c^{2}\) 。然后在此基础之上，椭圆会受到一个限制，同时满足此限制的椭圆上就可以有确定的 \(P\) 的位置。然后，我们再＂存在 \(A, B\) 满足斜率平方和 balabala．．．．．．＂，进而给粗圆再加一个限制。显然这种要素结构优于另一种（因为先定下 \(P\) 怎么看也是一种简化）。这就是这个题的要素结构。

既然如此，我们先来分析一下第一个条件。是否存在这样的 \(P\) 呢？显然是很基础的元问题，直接计算分析即可。设 \(P(x, y)\) ，计算不难得出

\[
\overrightarrow{P F_{1}} \cdot \overrightarrow{P F_{2}}=c^{2} \Leftrightarrow x^{2}+y^{2}=2 c^{2}
\]

而我们知道粗圆上一点 \(P(x, y)\) 满足 \(x^{2}+y^{2} \in\left[b^{2}, a^{2}\right]\) ，两个等号成立的位置就是四个顶点。因此第一个对椭圆的限制就是

\[
b^{2} \leq 2 c^{2} \leq a^{2} \Leftrightarrow \frac{\sqrt{3}}{3} \leq e \leq \frac{\sqrt{2}}{2}
\]

联立两个方程就得到 \(P\) 的位置。\\
对于一个控制，我们可以不着急先把它的元问题处理掉。先对接下来的事情进行一个预判再说。有了 \(P\) 之后，我们要对运动的 \(A, B\) 来看 \(k_{A P}^{2}+k_{B P}^{2}\) 。显然这是分别先控制这两个斜率，然后再取平方和。因此我们先关注这两个斜率：关于原点对称的 \(A, B\) ，一个椭圆上的定点 \(P\) ，研究 \(K_{A P}\) 和 \(K_{B P}\) 有什么元问题可以对应吗？立刻想到经典的一个优美二级结论 论：

\[
k_{A P} \cdot k_{B P}=-\frac{b^{2}}{a^{2}}
\]

这个二级结论非常匹配我们在干货列表二级结论视频中提出的判断原则，简洁优美，不可替代，因此很值得学习。这个结论显然很契合当前局势。

Up 主在经验分享第五期中强调，学知识点（包括二级结论）要学习背后的要素结构，作为元问题来学习。大家来说说，这个二级结论如何从结构层面理解？那就先分析它自身的结构，是两个主自由度的系统输出一个不变量，也就是一个恒等式。等式本身按理说应该减一个自由度，但是现在作为不变量恒等式，无法减自由度，因此从要素结构层面分析，它一定不会产生新的信息量，而是包含在原本两个主自由度之中。因此它一定可以作为一个等式充要替换掉原本系统里的一个自由度的信息。比如，粗圆上一点 \(P\) ，我过 \(P\) 作两条斜率乘积 \(-\frac{b^{2}}{a^{2}}\) 的直线，它们与椭圆的交点必定关于原点对称；粗圆上两点 \(A, B\) 关于原点对称，过它们分别作两条直线斜率乘积 \(-\frac{b^{2}}{a^{2}}\) ，则两直线和椭圆三线共点（三线共点又可以出来三种要素网结构）。元问题学习就是要实现一个知识点与其背后结构的双向匹配：给出知识点，在要素结构层面深刻理解；做题遇到这种结构能反向匹配回这个知识点。大家反思一下，看过第五期后进行实践的时候，在学习一个知识点的时候，是否实践好了这样的学习方式呢？

回到这道题。我们希望存在 \(A, B\) 使得

\[
k_{A P}^{2}+k_{B P}^{2}=\frac{7}{6}
\]

刚才也分析到了，希望先控制清楚这两个斜率。现在 \(P\) 是确定的，我们希望去寻找这样的 \(A, B\) 。既然如此，根据控制越多越好的原则，我们直接把刚才的元问题拿出来，选一种要素重构方式，实现对 \(A, B\) 的控制那该多好！既然如此，我们何不用

\[
k_{A P} \cdot k_{B P}=-\frac{b^{2}}{a^{2}}
\]

射出两条直线，这样就相当于选一个主自由度 \(k\) ，过 \(P\) 射两条斜率分别为 \(k,-\frac{b^{2}}{a^{2} k}\) 的直线，让它们交粗圆为 \(A, B\)\\
，这样不就实现了对 \(A, B\) 的控制吗？！这样的话，我们只需要

\[
k^{2}+\frac{b^{4}}{a^{4} k^{2}}=\frac{7}{6}
\]

有解即可。但是需要注意，\(k\) 的遍历范围需要抠掉

\[
k_{O P} \text { 与 }-\frac{b^{2}}{a^{2} k_{O P}}
\]

不然的话这两个交点不就成了 \(P\) 和它的对径点了吗？肯定是不行的。当然，如果 \(P\) 是四个顶点之一，那么就不需要抠了。

考虑函数

\[
f(k)=k^{2}+\frac{b^{4}}{a^{4} k^{2}}
\]

不难知 \(k= \pm \frac{b}{a}\) 时 \(f(k)\) 取最小值 \(\frac{2 b^{2}}{a^{2}}\) 。如果 \(\frac{2 b^{2}}{a^{2}}<\frac{7}{6}\) ，即 \(e>\frac{\sqrt{15}}{6}\) 时，\(f(k)=\frac{7}{6}\) 有四个解，就算挖掉两个也还有俩，于是符合。如果 \(\frac{2 b^{2}}{a^{2}}=\frac{7}{6}\) ，这时候不难算出

\[
a: b: c=6: \sqrt{21}: \sqrt{15}
\]

此时必须过 \(P\) 的两条直线斜率是 \(\pm \frac{\sqrt{21}}{6}\) 。这两个斜率是不是需要被抠掉的呢？那就要计算一下 \(P\) 的位置了。此时大小无关，不妨就假设

\[
a=6, b=\sqrt{21}, c=\sqrt{15}
\]

于是联立 \(\frac{x^{2}}{36}+\frac{y^{2}}{21}=1\) 与 \(x^{2}+y^{2}=30\) ，解出 \(x^{2}=\frac{108}{5}, y^{2}=\frac{42}{5}\) ，因此被抠掉的斜率是 \(\frac{\sqrt{14}}{6},-\frac{\sqrt{14}}{4}\)（或者反过来），不是 \(\pm \frac{\sqrt{21}}{6}\) 。因此 \(e=\frac{\sqrt{15}}{6}\) 时也符合题意。综合来看，\(e\) 的取值范围就是 \(\left[\frac{\sqrt{15}}{6}, \frac{\sqrt{2}}{2}\right]\) 。

再来看第二道题。这道题同样是大小无关的。先画图，画出 \(P\) ，再延长到 \(Q\) ，我们很容易发现 \(\triangle P F_{1} Q\) 的形状是完全确定的。因此，虽然这道题和上一道题一样，相当于同时要求＂存在 \(P, Q\) 满足以下条件＂。但是 \(P\) 的张角就已经定死 \(P\) 了，而且这个张角本身并不好计算，因此定死 \(P\) 的计算也会很困难，更不用说加下来还要去定死 \(Q\) 了。类似第一题，先画 \(Q\) 也是不切实际的。与第一题相比，这道题从一开始的预判就让人摸不着头脑。

那怎么办呢？这时候不如回想一下，up 主讲过的，关于离心率问题，有什么招数是还没考虑过的？（干货列表里的视频不过十几个，哪怕把所有的思想都数一遍看看呢）不难想到当时有一个视频说过，离心率小题可以考虑淡化圆雉曲线本体的地位，利用纯几何的方式把要素网构建出来，然后通过圆雉曲线的几何定义来控制出曲线，让圆雉曲线本身成为附庸。这是一种很彻底的要素网重构方式，特化针对于圆锥曲线小题。

这道题因为 \(\triangle P F_{1} Q\) 形状确定，又是大小无关问题，因此如果视为静止的话，静态要素一下子就多了一大堆，看上去非常爽。如果 \(\triangle P F_{1} Q\) 真的静止了的话，那就只剩一个 \(F_{2}\) 的位置了，而且，如果确定了 \(F_{2}\) 的位置，就能够以 \(F_{1}, F_{2}\) 为焦点，经过 \(P_{1}\) 画出椭圆，原题的要素一样不少。然后对系统的限制就只剩下 \(Q\) 在椭圆内。看上去起到了很明显的简化效果。

这样重构的话，原本的两个约束变元全都消失了，看上去改写的过于猛了。那么，这种改写是充要替换吗？同样用一下设问与遍历。 \(\triangle P F_{1} Q\) 静止，是不是所有满足上述限制的 \(F_{2}\) 都符合题意？只要这个 \(F_{2}\) 能让 \(Q\) 在椭圆内，那么这个图形自然也符合原题干；反之，原题干合法的一个图形，我把 \(\triangle P F_{1} Q\) 和 \(F_{2}\) 构成的系统抽出来，自然也符合改写后的题目要求。因此，这个要素网重构是合法的。有的同学也许会奇怪，那我的约束变元跑哪去了？其实大家回想一下，我们在番外篇里讲到的，约束变元的矛盾点不在它自身，它指向的是系统本体，它自己只是个傀儡。因此，它的存在就等价于为整个系统提供一个限制，且大多是不等关系，不减少自由度。这道题椭圆去

掉大小自由度只剩一个自由度，后面的＂存在 \(P, Q\)＂都是对椭圆本体的限制，因此实际自由度就是 1 ；而我们重构之后的系统的总自由度也是 1 ，自由度也守恒，没有任何问题。

接下来我们需要考虑＂\(Q\) 在椭圆内＂这个限制如何充要替换。如果是原题的要素网，那么语境仍然放在解析几何下，那我思考这个元问题自然就是要用解析几何的思维方式去思考：算坐标，也就是设 \(Q(x, y)\) ，满足

\[
\frac{x^{2}}{a^{2}}+\frac{y^{2}}{b^{2}}<1
\]

但是现在我们是用几何法来构建要素网，原本的这种处理方法显然行不通了。那么，几何法中，圆雉曲线的构建方法就是侑圆的几何定义，到 \(F_{1}, F_{2}\) 距离和为定值，定值为 \(\left|P F_{1}\right|+\left|P F_{2}\right|\) 。因此，我们考虑 \(Q\) 的充要替换也要放在几何定义的语境下去考虑。＂在椭圆上＂是＂等于 \(\left|P F_{1}\right|+\left|P F_{2}\right| "\) ，那＂在粗圆内＂是什么？是不是就应该是＂小于 \(\left|P F_{1}\right|+\left|P F_{2}\right|\)＂吗？那么，＂\(Q\) 在椭圆内＂是不是就等价于 \(\left|Q F_{1}\right|+\left|Q F_{2}\right|<\left|P F_{1}\right|+\left|P F_{2}\right|\) 呢？这种元问题大家可能比较少见，平时没有积累过，但是只要你把思路往＂几何定义＂上去思考，几何定义是距离和，那能不能用距离和去刻画＂椭圆内＂呢？只要你这么想了，你就不难想到这样来充要替换了。这种充要替换对不对呢？直观上肯定是对的，所以大家考场上就这么猜就行了。严格也是可以证明的：现在 \(P\) 是确定的，那么设

\[
\left|P F_{1}\right|+\left|P F_{2}\right|=d,
\]

如果

\[
\left|Q F_{1}\right|+\left|Q F_{2}\right|=d^{\prime},
\]

那么 \(Q\) 就在另一个椭圆上运动。这个新椭圆和老椭圆焦点相同，但是长轴更短。两个椭圆显然不可能有交点（不然交点到 \(F_{1}, F_{2}\) 的距离和就同时等于 \(d\) 和 \(d^{\prime}\) ），同时新椭圆长轴更短，至少长轴端点包含在老组圆内，所以新椭圆就只能包含在老椭圆内了。

不管怎么说，只要大家能够基于当前的要素结构出发思考，就能想到 \(Q\) 的充要替换。现在，假设 \(\triangle P F_{1} Q\) 是一个 \(5: 12: 13\) 的三角形。线段 \(P Q\) 上的一动点（主自由度）\(F_{2}\) 满足

\[
\left|Q F_{1}\right|+\left|Q F_{2}\right|<\left|P F_{1}\right|+\left|P F_{2}\right|,
\]

这个不等式可太简单了。设 \(\left|Q F_{2}\right|=t\) ，则 \(5+t<13+12-t\) ，即 \(0<t<10\) 。离心率也很简单，就是

\[
\frac{\left|F_{1} F_{2}\right|}{\left|P F_{1}\right|+\left|P F_{2}\right|}=\frac{\sqrt{25+t^{2}}}{25-t}
\]

显然它是关于 \(t\) 的增函数，因此计算端点值就得到离心率范围是 \(\left(\frac{1}{5}, \frac{\sqrt{5}}{3}\right)\) 。\\
做个总结。这两道题都涉及到了要素关系网的重构与对元问题的把控。重构的过程都基于要素网的一般建立过程。先尝试建立不同的要素网，建立时对计算进行预判，预判就是基于要素结构调用元问题分析一下大致要计算什么，这些计算的可行性大概如何。如果预判发现计算很不舒服，就可以尝试换一种要素网构建方式，在多种要素网之间进行比较，选一个最稳妥的。使用元问题需要对它本身具有深度的理解，从要素结构的层面理解一个知识点刻画的动静关系，控制牵制，以及充要替换。这样才能配合对题目结构的把控打出组合拳。这两道题都深刻体现了题目结构把握与元问题结合的解题基本思想，思想很基本，但是用起来威力无穷。大家一定要注意，元问题的学习与积累在平时，结构思维的锻炼在平时，只有平时科学努力，把力气用在刀刃上，才能在实战中取得稳定的成绩。

\section*{课后练习}
1．对于第一题，如果我想把 \(\frac{7}{6}\) 改成另一个数，其他条件与所求不变，使得这道题的答案变成一个左开右闭区间，我应该改成什么？这时这道题的答案又会是多少？\\
提示：看看本题解题过程的哪一步可能会导致去掉一个解？\\
2．把本视频的第二题和2023年5月1日视频中的第二题合起来，体会它们之间方法与思想的共性。

6．9．4 XXXX

\section*{第7章 套卷相关的综合讲解}
\subsection*{7.1 2023年全国甲卷}
\section*{7．1．1 检查的重要性和基本方法}
很多同学在考后分析试卷，最痛苦的是犯了很多＂俊错＂。看错正负号，审错题，把6看成9等等。这些错误看似不起眼，但是可能会导致巨大的失分。相信绝大多数的读者也完全知道纠正它们的重要性，但却一直苦于没有一个系统化的解决方案而苦恼良多。希望这简单的一章可以为大家建立起＂检查基本法＂，筑起一个最基础的检查框架。

首先，我们要承认一点：没有人不会犯错误。大家在我的视频里也经常能看到，我在做题时也经常犯错。不等式符号者反，算错数的情况数不胜数。即使是计算机，也无法保证程序永远不出错，更不用说人脑了。所以，检查基本法的第一步，也是最重要的一步，就是调整思路：不要再去＂避免错误＂，而是＂找出错误＂。什么意思？错误是不可避免的，所以自己出现任何错误都是可接受的；同时，这些错误也都是可挽回的：只要我在考试结束之前，发现它们，改正它们，分数仍旧属于我。在改正它们之前，分数只是暂时不属于我。如果能够把检查这个像＂寻宝＂一样的游戏玩好，那么一定可以拿到一个理想的成绩。

这个＂寻宝＂游戏最难的两点，就是＂隐蔽性＂和＂时间限制＂。思维定势，惯性远比我们想象的要强大，所以只是用眼睛＂扫描＂，机械复算，是不可能找到错误的。你的大脑已经习惯了刚刚犯下的错误，会认为它是合理的。这个错误就像是藏在监控死角里的老鼠，如果不调整监控的视角，是不可能发现的。所以，在检查时，需要我们采取和第一遍相对不同的方法去重新看问题，这也是＂正向检查＂的奥义之一。

时间限制，是另一个困扰大家的问题。很多同学抱怨，我考试时的题都＂会做＂，但是就是做不完，算不完，所以导致我分数一直上不去。在这里，我也希望大家尽可能调整好一个心态：考试时间之内，做不完的题，就是你＂不会＂的题。什么意思？虽然可能每次考试你不会的题不同，是不一样的知识点，但是它们都是暂时在你能力范围之外的东西。它们之所以在能力范围之外，有非常多的原因：思路找得慢，计算慢等等。而考试时，最关键的，就是把所有你能力范围之内的分数全部拿到。检查就是帮助你拿到所有分数最关键的手段。

这里也涉及到一个心态的调整：＂快错＂不如＂慢对＂。很多同学学得还行，但是做题速度慢，所以，在考试时就会着急忙慌赶紧把所有题目做完。结果看上去所有题目都写了，写得也挺快，卷子满满当当一大片，但是得分率低，错了很多。甚至有可能，因为着急，大题上来就错，后面写的东西完全没有价值。这就是＂快错＂。这样的话，所有花在这道题上的时间就全都浪费了，不如把时间节约下来，留给前面＂可能出错＂，但是停下来检查就一定能对的题目。这就是＂慢对＂。很多同学之所以会出现＂快错＂，就是因为觉得自己可以＂快对＂：要知道，这是非常难，非常难的事情！既要快，又要对，不光基础要扎实，而且解题速度也要快，这对大家的要求很高。所以，循序渐进，先追求慢对，再到快对，才可能让成绩有一个本质的提升。

还是以我自己为例，如果我不检查，就会像视频中一样，出一些＂小错＂。这些小错可能不影响观看，但是考试时就是致命的错误。尤其是，选填只看结果，开闭弄反，少算一个 0 这种，就是一分没有，和不会没有区别。这种损失都是不可接受的。能力范围之外的题目不会是正常的，但是这种能力范围之内，完全可以避免的闪失，一定要尽可能想办法避免。

具体而言，有哪些检查的方法呢？大体上分为三类：\\
（1）正向检查：顾名思义，就是按着正常的思路再算一遍。由于思维定势的存在，简单的重复第一遍的过程是很

难找到错误的。因此，正面检查的关键，在于使用不同于第一遍的方法重新计算。正向检查最大的坏处是效率相对较低，但是最大的好处是实施起来相对简单。\\
（2）反向检查：这是我们小学时就一直在实践的检查方法。加法题用减法验算，减法题用加法演算等等。在高中阶段，题目更为复杂，我们不再是以简单的四则运算的逆运算的形式去验算。而是：将结果代人题目，重新检验它是否满足条件。无论是题目中的步骤，简单题，中档题还是难题，都可以用这个方法检查。它最大的优势在于：既不会掉入思维定势，实施起来也相对容易。所以，强烈推荐多使用反向检查的方法。\\
（3）结构检查：我们前面讲了这么多代数杜杆，它们的用处不仅仅是解题，还可以帮你查错。要知道，如果好的代数结构一直被保留，那么你总可以通过找到它，来检查你的结果。如果算着算着，对称式不对称，齐次式不齐次，肯定哪里出了问题。不光是代数结构，所有问题本身具有的属性，都可以被用来检查。这个方法最大的好处，就是逻辑链条短。不需要像（1），（2）一样，几乎要再完整走一遍整个过程。

\section*{7．1．2 正向检查—不同方法另算}
所谓正向检查简单来说就是再算一遍，但是方法不同。

\section*{7．1．2．1 不要＂一题多解＂，要＂一题多理解＂}
我们前面提到，一道题，用同样的方法算两遍，第二遍收益非常有限。这是因为，我们很难跳出自己的思维定势，看到自己前面犯过的错误。因此，如果想要正面检查，一定要采取和第一次方法不完全相同的做法。这些新的做法并非来源于异想天开，而是来自于我们对于这个问题不同角度的洞见和观察。

例如，我们在讲解代数杜杆时，遇到过很多具有不止一个结构的多元函数。比如：

\[
z=\frac{x^{2}+x y+y^{2}}{x^{2}+y^{2}}
\]

求 z 的取值范围。它具有这些好的代数结构：\\
（1）对称性：分子，分母都是关于 \(x, y\) 的对称式；\\
（2）齐次：分子，分母都是关于 \(x, y\) 的齐次式，整个分式是 0 次齐次；\\
（3）二次：分子，分母都是关于 \(x, y\) 的二次式。\\
这些好的代数结构分别指向不同的解题方法：\\
（1）对称性：令 \(x^{2}+y^{2}=m, x y=n\) ，那么根据均值不等式，

\[
m \geq 2 n, m \geq-2 n
\]

故 \(\frac{n}{m} \in\left[-\frac{1}{2}, \frac{1}{2}\right]\) ．此时

\[
z=\frac{m+n}{m}=1+\frac{n}{m} \in\left[\frac{1}{2}, \frac{3}{2}\right] .
\]

（2）齐次：令 \(t=\frac{x}{y}\) ，那么

\[
z=\frac{t^{2}+t+1}{t^{2}+1}=1+\frac{1}{t+\frac{1}{t}}
\]

注意到对勾函数 \(f(t)=t+\frac{1}{t}\) 的范围是 \(f(t) \geq 2\) 或 \(f(t) \leq-2\) ，所以

\[
z \leq 1+\frac{1}{2}=\frac{3}{2}, z \geq 1-\frac{1}{2}=\frac{1}{2}
\]

故 \(z \in\left[\frac{1}{2}, \frac{3}{2}\right]\) ．\\
（3）二次：注意到分子，分母都是二次，且 \(x, y\) 可以取任意实数（不同时为 0 ），所以可以考虑使用判别式：

\[
x^{2}+x y+y^{2}=z x^{2}+z y^{2}
\]

选 \(x\) 为主元，将它整理成关于 \(x\) 的二次方程：

\[
(1-z) x^{2}+y \cdot x+(1-z) y^{2}=0
\]

注意到二次项系数为 \(1-\mathrm{z}\) 可能为 0 ，需要分类讨论：\\
（1）若 \(1-z=0\) ，则 \(z=1\) ，此时 \(x y=0\) ．故 \(z=1\) 可以取到；\\
（2）若 \(1-z \neq 0\) ，则

\[
\Delta=y^{2}-4(1-z)^{2} y^{2}=y^{2}\left[1-4(1-z)^{2}\right] \geq 0
\]

由于 \(y^{2} \geq 0\) ，所以 \(1-4(1-z)^{2} \geq 0\) ．解得 \(\frac{1}{2} \leq z \leq \frac{3}{2}\) ．因此 \(z \in\left[\frac{1}{2}, \frac{3}{2}\right]\) ．

如果我们在学习代数结构的时候，可以掌握所有的这三个杠杆，那么处理这类问题的时候是不可能出错的。只要用不同的方法检查，得到同样的答案，那么结果大概率是正确的。

在这些不同的解法中，关键的不是解法本身，而是我们对于问题的观察。这些观察，构成了我们发展出这些解法的动机。换言之，也是要理清楚＂为什么这么想＂。我们换元，除以 \(y^{2}\) ，用判别式，不是因为这是一个＂固定的方法＂，我们碰运气一样去砸它，而是因为有了这些结构，它们才能够一一奏效。这些结构就像＂信号＂，给我们充分的提示去选择正确的方法。这也是为什么我们在这里强调＂一题多理解＂，而非＂一题多解＂：多解只是表象，藏在多解下面的原理和底层逻辑是我们对于结构的理解。总结这些理解，才能够发展出足够强大的解题逻辑，以应对所有可能出现的问题。

使用完全不同的方法检查，是一个理想情况。但是，在现实中，可行性较低：由于时间限制，我们不可能每道题做两遍，这几乎相当于做了两张卷子。所以，在这种情况下，我们只能用＂半个新方法＂去处理问题。什么叫＂半个新方法＂？意思就是：我们不完全另起炉灶，而是在过程中，找到另外的角度来理解。注意，这里的重点不再是给出全新的解法，而是要对题目的结构和答案有一个新的理解。通过这个理解，就可以检验我们答案的正确性。这也是我们说的＂结构检查＂。

\section*{7．1．3 结构检查一代数与合理性}
我们通过观察问题本身的结构，可以非常快地发现错误。

\section*{7．1．3．1 代数结构：}
当我们展开具有某个特定结构的式子时，如果结构没有被破坏，那么在化简过程中，我们始终都能观察到它。例如，展开表达式：

\[
\left(x^{2}+x y+y^{2}\right)\left(2 x^{2}+5 x y+2 y^{2}\right)
\]

如果我们有如下四个结果：\\
A， \(2 x^{6}+7 x^{4} y^{2}+9 x^{3} y^{3}+7 x^{2} y^{4}+2 y^{6}\)\\
B， \(2 x^{4}+7 x^{3} y+9 x^{2} y^{2}+7 x y^{3}+2 y^{4}\)\\
C， \(2 x^{4}+7 x^{2} y+9 x^{2} y^{2}+7 x y^{2}+2 y^{4}\)\\
D， \(2 x^{4}+7 x^{3} y+9 x^{2} y^{2}+6 x y^{3}+2 y^{4}\)

我们如何通过结构，直接找到正确的答案呢？注意到，原始的表达式有两个非常重要的结构：\\
（1）对称性：\(x^{2}+x y+y^{2}, 2 x^{2}+5 x y+2 y^{2}\) 是关于 \(x, y\) 的对称式。因此，最后展开的结果，也一定是对称的；\\
（2）齐次性：\(x^{2}+x y+y^{2}, ~ 2 x^{2}+5 x y+2 y^{2}\) 都是关于 \(x, y\) 的二次齐次式，因此，最后展开的结果，也一定是齐次式，且是 \(2+2=4\) 次齐次式。根据这两个理解，我们如何马上得到正确答案呢？我们来分别考察这几个选项：

选项 A：它是对称式，同时也是齐次式，但是它不是 4 次齐次，而是 6 次齐次，错误；\\
选项 B：它是对称式，同时也是齐次式，同时也是 4 次齐次式，可能正确；\\
选项 C：它是对称式，但不是齐次式。这里， \(2 x^{4}+9 x^{2} y^{2}+2 y^{4}\) 是一个 4 次齐次式，然而 \(7 x^{2} y+7 x y^{2}\) 则是一个 3次齐次式。所以，它不是齐次式，错误；\\
选项 D：它是齐次式，但不是对称式：\(x^{3} y, x y^{3}\) 两项系数不同，错误。\\
因此，这四个选项中只有 B 选项符合题意。在我们进行更复杂的计算（如圆锥曲线）等，也可以利用所有代数杠杆检查。

\section*{7．1．3．2 合理性：}
我们最终得到的结果和答案，一定要符合知识本身的结构，要是一个合理的答案。例如，我们在计算概率时，得到的答案一定是［ 0,1 ］之内的实数，这是由概率的定义决定的。如果我们得到了 \([0,1]\) 之外的数，那么答案一定是错误的；再比如，一个多元函数 \(z=\frac{x^{2}+y^{2}}{x^{2}+x y+y^{2}}\) ，它的分子，分母都是非负的，所以它的取值范围端点也一定是非负数，取值范围不可能是 \(\left[-\frac{1}{2}, \frac{1}{2}\right]\) 。所有这些判断，都源自于我们对于＂合理性＂的把握：这个答案得是什么样的，才是合理的。那么，有哪些常见的合理性的判断呢？\\
（1）量级：答案是几位数？大概是多少？要有一个基本的估计。例如，某人扔三次正方体骰子，那么第三次比第一次大的有多少种？如果计算出答案是 800 多，肯定是错的。因为整个样本空间中，只有 \(6^{3}=216\) 个样本点，因此，答案一定小于 216 ；\\
（2）符号：答案是正的，还是负的？例如刚才的概率，多元函数的例子；\\
（3）数论：在数的计算上，数论知识也是非常常用的验证合理性的手段。例如，如果将 4 不同的书，分给甲，乙两人，如果答案是 13 ，那么一定是错的。这是因为：答案一定是偶数。因为：甲，乙两人交换自己拿到的书，一定也是一种分法。因此，所有分法都可以和它相反的分法配成一对，答案必然是偶数。又比如，我们计算时可能会出现 \(16 \times 5=90\) 这样的错误。这也是不合理的：左侧 16,5 都不是 9 的倍数，但是结果却是一个 9 的倍数。\\
（4）代数结构：事实上，前面讲的所有代数结构，都可以作为合理性的判断依据。这是计算基本法的主要内容，所以我们在前面将它单独拿出，特别说明。但在结构检查法中，本质上也是对合理性的一种体会。

\section*{7．1．4 反向（即时性）检查}
反向检查法，顾名思义，就是从另一个方向检查（和第一章讲的即时性类似，都能让计算过程发生本质性变化）。这个方法在小学时就已经充分熟知了：加法用减法检查，乘法用除法检查等。在高中阶段，我们完全可以采用同样的思路：\\
（1）反向恒等变形：例如，当我们需要对代数式进行因式分解时，可以将结果展开，来验证是否成立。假设我们得到以下分解因式结果：

\[
x^{3}+2 x-3=(x-1)\left(x^{2}+x+3\right)
\]

那么我们将右侧的结果，利用杠杆 1 中的知识重新展开：\\
（1）\(x^{3}: x \cdot x^{2}=x^{3}\) ；\\
（2）\(x^{2}:-x^{2}+x \cdot x=0\) ；\\
（3）\(x: 3 x-x=2 x\) ；\\
（4） \(1:-1 \cdot 3=-3\) ．\\
我们发现，展开的结果就是 \(x^{3}+2 x-3\) ．因此，我们分解因式的结果正确。事实上，所有操作，都可以用它的＂逆运算＂进行检查：加法 \textbackslash ＆减法，乘法 \textbackslash ＆除法，展开 \textbackslash ＆分解因式，求导 \textbackslash ＆积分等等。＂逆运算＂往往藏在我们意想不到的地方，例如，我们学习过辅助角公式（也叫＂提斜＂）：

\[
a \sin x+b \cos x=\sqrt{a^{2}+b^{2}} \sin (x+\varphi)
\]

其中 \(\varphi\) 满足：

\[
\cos \varphi=\frac{a}{\sqrt{a^{2}+b^{2}}}, \sin \varphi=\frac{b}{\sqrt{a^{2}+b^{2}}}
\]

很多同学在使用这个公式的时候心惊胆战，生怕犯错。事实上，当我们算完后，从等式的右边重新推出等式的左边就可以了：根据和角公式，有

\[
\begin{aligned}
\sqrt{a^{2}+b^{2}} \sin (x+\varphi) & =\sqrt{a^{2}+b^{2}}(\sin x \cos \varphi+\sin \varphi \cos x) \\
& =\cos \varphi \cdot \sqrt{a^{2}+b^{2}} \cdot \sin x+\sin \varphi \cdot \sqrt{a^{2}+b^{2}} \cdot \cos x \\
& =\frac{a}{\sqrt{a^{2}+b^{2}}} \cdot \sqrt{a^{2}+b^{2}} \cdot \sin x+\frac{b}{\sqrt{a^{2}+b^{2}}} \cdot \sqrt{a^{2}+b^{2}} \cdot \cos x \\
& =a \sin x+b \cos x
\end{aligned}
\]

恰好等于左式。因此，结论成立，我们使用＂辅助角公式＂和使用＂和角公式＂，就可以被看作一对＂逆运算＂。具体而言，给出如下几个三角函数恒等变形，我们如何判断哪个是正确的呢？

A，\(\frac{1}{2} \sin x-\frac{\sqrt{3}}{2} \cos x=\sin \left(x-\frac{\pi}{6}\right)\) ；\\
B ，\(\frac{1}{2} \sin x-\frac{\sqrt{3}}{2} \cos x=\sin \left(x-\frac{\pi}{3}\right)\) ；\\
C，\(-\frac{\sqrt{3}}{2} \sin x+\frac{1}{1} \cos x=\sin \left(x+\frac{5 \pi}{6}\right)\) ；\\
D， \(3 \sqrt{3} \cos x+3 \sin x=6 \sin \left(x+\frac{\pi}{3}.\right)\)\\
（2）反向代入：我们解方程后，可以将得到的根代回原来的方程，确认它是否是根；分解因式后，可以代入特殊的 \(x\) 值来确认分解的结果是否正确。还是以刚才的例子为例，如果想要检查

\[
x^{3}+2 x-3=(x-1)\left(x^{2}+x+3\right)
\]

是否正确，可以代人特别的 \(x\) 去检查左右两侧是否相等。\\
（1）当 \(x=0\) 时，左式为 -3 ，右式为 \(-1 \cdot 3=-3\) ，成立；\\
（2）当 \(x=1\) 时，左式为 0 ，右式为 \(0 \cdot 5=0\) ，成立；\\
（3）当 \(x=-1\) 时，左式为 -6 ，右式为 \(-2 \cdot 3=-6\) ，成立；\\
（4）当 \(x=2\) 时，左式为 \(8+4-3=9\) ，右式为 \((2-1)\left(2^{2}+2+3\right)=1 \cdot 9=9\) ，成 \(\underline{V}_{\mathrm{o}}\) 对于三次函数的分解因式，我们只要代人四个不同的值，如果左右两侧都相等，就可以保证成立。这是因为：一个三次函数本身就有四个自由度，

\[
f(x)=a_{3} x^{3}+a_{2} x^{2}+a_{1} x+a_{0} .
\]

当我们有四个限制条件后，这个三次函数就确定下来了。左右两侧的表达式都是满足这四个条件的三次函数，那么必然相等。在实际应用中，如果时间有限，可以只代入 1 到 2 个，这样可以有效增加正确率和错误检出率。注

意，如果只代入一个值，尽量不要选择 \(x=0\) ．这是因为，它只能检查 \(a_{0}\) ：

\[
f(0)=a_{0} .
\]

它和别的系数无关，如果别的系数出错，无法检查出来。所以，除了 \(x=0\) ，最好再代入除了 \(x=0\) 以外的数进行检查。

一般地，\(n\) 次函数有 \((n+1)\) 个自由度，所以需要代入 \((n+1)\) 个点才能完全检查。实际应用中，根据情况代人少量点即可。

除了多项式的展开，其他计算也可以用代人检查。我们还是考虑前面三角函数的例子：对于

\[
\frac{1}{2} \sin x-\frac{\sqrt{3}}{2} \cos x=\sin \left(x-\frac{\pi}{6}\right)
\]

我们可以代人特殊的 \(x\) 值检查。例如，当 \(x=\frac{\pi}{3}\) 时，左式为

\[
\frac{1}{2} \sin \frac{\pi}{3}-\frac{\sqrt{3}}{2} \cos \frac{\pi}{3}=0
\]

而右式为

\[
\sin \left(\frac{\pi}{3}-\frac{\pi}{6}\right)=\sin \frac{\pi}{6}=\frac{1}{2}
\]

不成立。所以，等式不成立。而对于

\[
\frac{1}{2} \sin x-\frac{\sqrt{3}}{2} \cos x=\sin \left(x-\frac{\pi}{3}\right)
\]

我们可以代入：\\
（1）当 \(x=\frac{\pi}{3}\) 时，左式为 0 ，右式为 \(\sin \left(\frac{\pi}{3}-\frac{\pi}{3}\right)=0\) ，成立；\\
（2）当 \(x=\frac{\pi}{2}\) 时，左式为 \(\frac{1}{2} \cdot 1-\frac{\sqrt{3}}{2} \cdot 0=\frac{1}{2}\) ，右式为 \(\sin \left(\frac{\pi}{2}-\frac{\pi}{3}\right)=\sin \left(\frac{\pi}{6}\right)=\frac{1}{2}\) ，成立．\\
事实上，检验辅助角公式，只需要两个条件。这是因为，表达式 \(a \sin x+b \cos x\) 也只有两个自由度 \(a, b\) 。只要代入的角度的终边不重合，就可以保证答案正确。

更一般地，我们可以把得到的答案或结论代回到原来的题目，验证是否满足已知条件。例如，对于三角函数大题，当我们求出解析式后，可以反过来验证它是否满足已知条件，如周期，最值等，也可以代人特殊的 \(x\) 值检验三角函数值；对于数列大题，如果已知 \(a_{n}\) 求 \(S_{n}\) ，我们可以利用 \(S_{n}-S_{n-1}=a_{n}\) 验证结果。我们将在练习中使用这些方法。

\section*{7．1．5 先对后快—快错与慢对}
所谓先对后快即＂快错＂永远不如＂慢对＂。我们在最开始提到，很多的同学担心做不完，导致做题时毛手毛脚，题是做完了，但是正确率很低，导致成绩也受到影响。要知道，我们最终要完成的目标是＂又快又好＂，也就是＂快对＂。我们从一开始的＂又慢又错＂，发展到＂又快又对＂，是不可能一蹴而就的，一定有中间步骤。这里的中间步骤，只可能是＂慢但是对＂，而不可能是＂快但是错＂。

原因很简单：＂慢对＂可以发展到＂快对＂，但是＂快错＂无法发展到＂快对＂。＂慢但是对＂，意味着我们已经可以做到彻底理清题目，知识的底层逻辑，欠缺的只是熟练度。那么加以训练，一定可以慢慢向＂快＂的方向进化；而＂快但是错＂，是无法通过增加熟练度来向＂快对＂进化的，因为它没有把题目正确的底层逻辑完全理清楚。反复训练，只能训练错误的路径，那么成绩就永远有一个天花板。想要逃离＂快错＂的恶性循环，只能打破＂错＂，走向＂对＂，学习全新的正确思路。全新思路的使用必然是慢的，所以只能先向＂慢对＂转化。因此，即使现在是＂快错＂，也必然先经历＂慢对＂，再到＂快对＂。

这和我们在各种运动的训练中是一样的：一上来，都是练习＂正确的＂的基本动作，这是＂慢对＂。它和最终熟练掌握（快对）相比，有很大的距离，但是是它必须要经过的阶段。我们不可能在一开始就追求＂快＂，但是练习错误动作，也就是＂快错＂。这样一来，我们的错误动作反而会定型，后续会付出更大的代价去纠正错误，得不偿失。

曾经选择＂快错＂，也完全不需要自责。这只是考试机制带来的影响。题量大，计算复杂，会很自然地引导部分同学选择＂快错＂：看到这么多题，我得赶紧做完才行，不然肯定考不好。如果你有这样的焦虑，一定要宽慰自己：＂慢＂就是＂快＂。只有一开始的＂慢对＂，才有可能发展到最后的＂快对＂。在我们水平有限的时候，大胆地选择这条看上去更慢，但是更正确的那条道。它一定比那条看上去快的捷径能走更远，有更高的上限。

我们以 2023 年甲卷为例，看看计算杜杆里的检查思想是如何在其中发挥作用的。

注因本节内容的主要目的是应用这些检查方法，所以无需提前做题，故没有放出纯享版的试卷，如有需要请自行查找。

\section*{7．1．6 选择题 01－12}
1．设全集 \(U=\mathbb{Z}, M=\{x \mid x=3 k+1, k \in \mathbb{Z}\}, N=\{x \mid x=3 k+2, k \in \mathbb{Z}\}\) ，则 \(C_{U}(M \cup N)=\) \(\qquad\)。\\
A，\(\{x \mid x=3 k, k \in \mathbb{Z}\}\)\\
B ，\(\{x \mid x=3 k-1, k \in \mathbb{Z}\}\)\\
C，\(\{x \mid x=3 k-2, k \in \mathbb{Z}\}\)\\
\(\emptyset\)

直接处理即可。集合 \(M\) 代表所有除以 3 余 1 的整数，集合 \(N\) 代表所有除以 3 余 2 的整数，那么 \(M, N\) 并集的补集就是所有 3 的倍数，故选 A．在这里，如果不好处理或想要检查，可以利用代入法检查。\\
（1）对于集合 \(M\) ，我们代入 \(k=0,1,2,3, \ldots\) ，可以得到 \(M\) 含有如下几个元素：1，4，7，10， \(13, \ldots\) 显然，它们是除以 3余 1 的。将 \(M\) 中的数排成数列，通项公式 \(a_{k}=3 k+1\) 也是公差为 3 的等差数列，所以每一项都比前一项大 3 ，它就可以遍历所有除以 3 余 1 的数；\\
（2）对于集合 \(N\) ，我们代入 \(k=0,1,2,3, \ldots\) ，可以得到 \(N\) 含有如下几个元素：2，5， \(8,11,14 \ldots\) 显然，它们是除以 3 余 2 的。将 \(M\) 中的数排成数列，通项公式 \(a_{k}=3 k+2\) 也是公差为 3 的等差数列，所以每一项都比前一项大 3 ，它就可以遍历所有除以 3 余 1 的数。\\
将它们从整数序列中剔除，剩下的就是 \(0, ~ 3, ~ 6,9, \ldots\) ，故选 A．注意，这里 \(M, N\) 和选项 A 也都包含了负整数。

2．设 \(a \in \mathbb{R}\) ，且 \((a+i)(1-a i)=2\) ，则 \(a=\) \(\qquad\) －\\
A，-2\\
B ，-1\\
C，1\\
D， 2

直接展开括号，得

\[
(a+i)(1-a i)=a-a^{2} i+i+a=2 a+\left(1-a^{2}\right) i
\]

在这里展开的时候，我们可以不使用任何额外的技巧，也可以选 \(i\) 为主元。这是因为：我们得到的是一个虚数，它必然可以写成 \(x+y i\) 的形式，其中 \(x, y \in \mathbb{R}\) 。这样一来，

\[
(i+a)(-a i+1)=-a i^{2}+\left(1-a^{2}\right) i+a=2 a+\left(1-a^{2}\right) i
\]

根据题目条件，由于它是 2 ，且 \(a\) 为实数，所以

\[
2 a=2,1-a^{2}=0
\]

\section*{解得 \(a=1\) ．}
检查的时候，可以反向代人：当 \(a=1\) 时，根据平方差公式，

\[
(1+i)(1-i)=1-i^{2}=2
\]

3．执行下面的的程序框图，则输出的 \(B=\) \(\qquad\) －\\
A， 21\\
B， 34\\
C， 55\\
D， 89

我们可以将 \(A, B\) 的值都记录下来。\\
（1）\(k=1\) 时，进循环。 \(A=1+2=3, B=3+2=5\) ．循环结束时 \(A=3, B=5\) ；\\
（2）\(k=2\) 时，进循环。 \(A=3+5=8, B=5+8=13\) ．循环结束时 \(A=8, B=13\) ；\\
（3）\(k=3\) 时，进循环。 \(A=8+13=21, B=13+21=34\) ．循环结束时 \(A=21, B=34\) ．\\
（4）当 \(k=4\) 时，出循环。此时 \(B=34\) ．\\
4．已知向量 \(\vec{a}, \vec{b}, \vec{c}\) 满足 \(|\vec{a}|=|\vec{b}|=1,|\vec{c}|=\sqrt{2}\) ，且 \(\vec{a}+\vec{b}+\vec{c}=\overrightarrow{0}\) ．则 \(\cos (\vec{a}-\vec{c}, \vec{b}-\vec{c})\)\\
\(\vec{c})=\) \(\qquad\) －\\
A，\(-\frac{4}{5}\)\\
B，\(-\frac{2}{5}\)\\
C，\(\frac{2}{5}\)\\
D，\(\frac{4}{5}\)

当我们看到三个向量的长度分别为 \(1, ~ 1, \sqrt{2}\) 时，我们马上想到腰长为 1 的等腰直角三角形。由于 \(\vec{a}+\vec{b}+\vec{c}=\overrightarrow{0}\) ，所以它们可以构成首尾相接的三角形，这个三角形就是 \(45^{\circ}, 45^{\circ}, 90^{\circ}\) 的三角形。因此，\(\vec{a} \perp \vec{b}\) ．另外，也可以将它们顶点重合，设为 \(O \cdot \vec{a}, \vec{b}, \vec{c}\) 的终点分别记作 \(A, B, C\) ．\\
这个题有多种切入角度，非常适合使用正向检查法。\\
【方法一】余弦定理。 \(\vec{a}-\vec{c}=\overrightarrow{C A}, \vec{b}-\vec{c}=\overrightarrow{C B}\) ，故所求就是 \(\cos \angle A C B\) ．我们可以使用余弦定理。过 \(C\) 作 \(A\) 垂线，垂足为 \(D\) ，显然 \(A D=1+1=2, C D=1\) ，所以 \(|A C|=\sqrt{5}\) ．同理，\(|B C|=\sqrt{5}\) ．根据余弦定理，

\[
\cos \angle A C B=\frac{|C A|^{2}+|C B|^{2}-|A B|^{2}}{2|C A| \cdot|C B|}=\frac{10-2}{2 \cdot \sqrt{5} \cdot \sqrt{5}}=\frac{4}{5} .
\]

故所求为 \(\frac{4}{5}\) ．\\
【方法二】建系。以 \(\overrightarrow{O B}, \overrightarrow{O A}\) 方向分别为 \(x, y\) 轴正方向建立平面直角坐标系，那么可以写出所有向量的坐标表示：

\[
\vec{a}=(0,1), \vec{b}=(1,0), \vec{c}=(-1,-1) .
\]

所以

\[
\vec{a}-\vec{c}=(1,2), \vec{b}-\vec{c}=(2,1) .
\]

所以

\[
\cos \langle\vec{a}-\vec{c}, \vec{b}-\vec{c}\rangle=\frac{(\vec{a}-\vec{c}) \cdot(\vec{b}-\vec{c})}{|\vec{a}-\vec{c}| \cdot|\vec{b}-\vec{c}|}=\frac{4}{\sqrt{5} \cdot \sqrt{5}}=\frac{4}{5} .
\]

【方法三】直接计算。注意到 \(\vec{a}, \vec{c}\) 的夹角为 \(135^{\circ}, \vec{b}, \vec{c}\) 的夹角也是 \(135^{\circ}\) ，所以可以直接计算：

\[
\begin{gathered}
|\vec{a}-\vec{c}|=\sqrt{(\vec{a}-\vec{c}) \cdot(\vec{a}-\vec{c})}=\sqrt{|\vec{a}|^{2}+|\vec{c}|^{2}-2|\vec{a}| \cdot|\vec{c}| \cdot \cos 135^{\circ}} \\
=\sqrt{1+2-2 \cdot \sqrt{2} \cdot\left(-\frac{\sqrt{2}}{2}\right)}=\sqrt{5}
\end{gathered}
\]

\(|\vec{b}-\vec{c}|=\sqrt{5}\),

\[
(\vec{a}-\vec{c}) \cdot(\vec{b}-\vec{c})=\vec{a} \cdot \vec{b}-\vec{a} \cdot \vec{c}-\vec{b} \cdot \vec{c}+|\vec{c}|^{2}=0-2 \cdot 1 \cdot \sqrt{2} \cdot\left(-\frac{\sqrt{2}}{2}\right)+2=4
\]

所以

\[
\cos \langle\vec{a}-\vec{c}, \vec{b}-\vec{c}\rangle=\frac{(\vec{a}-\vec{c}) \cdot(\vec{b}-\vec{c})}{|\vec{a}-\vec{c}| \cdot|\vec{b}-\vec{c}|}=\frac{4}{5} .
\]

【注】本题的关键是观察到特殊三角形 \(45^{\circ}, 45^{\circ}, 90^{\circ}\) ．

5．设等比数列 \(\left\{a_{n}\right\}\) 的各项均为正数，前 \(n\) 项和为 \(S_{n}\) ，若 \(a_{1}=1, S_{5}=5 S_{3}-4\) ，则 \(S_{4}\) 的值为 \(\qquad\) ＿。\\
A，\(\frac{15}{8}\)\\
B，\(\frac{65}{8}\)\\
C，15\\
D， 40

【方法一】注意到题目给的是首项为 1 的等比数列，设公比为 \(q\) ，则 \(a_{n}=q^{n-1}\) ，

\[
S_{n}=\frac{q^{n}-1}{q-1} .
\]

所以直接根据题目条件，可以得到

\[
\begin{aligned}
S_{5}= & 5 S_{3}-4 \Longrightarrow \frac{q^{5}-1}{q-1}=5 \cdot \frac{q^{3}-1}{q-1}-4 \\
& q^{5}-1=5\left(q^{3}-1\right)-4(q-1) \\
& q^{5}-5 q^{3}+4 q=0 \\
& q\left(q^{4}-5 q^{2}+4\right)=0 \\
& q\left(q^{2}-1\right)\left(q^{2}-4\right) .
\end{aligned}
\]

所以 \(q=0, \pm 1, \pm 2\) ．由于等比数列各项都是正数，所以 \(q=1\) 或 \(q=2\) ．\\
（1）当 \(q=1\) 时，\(a_{n}=1, S_{5}=5, S_{3}=3, S_{5} \neq 5 S_{3}-4\) ，不成立；\\
（2））当 \(q=2\) 时，\(a_{n}=2^{n}-1, S_{4}=1+2+4+8=15\) ，故选 C．\\
这里，我们【反向代入】进行检查：此时，\(S_{5}=1+2+4+8+16=31, S_{3}=1+2+4=7\) ，所以

\[
5 S_{3}-4=31=S_{5}
\]

故题目条件成立！\\
这个方法的＂问题＂是没有考虑 \(q=1\) 的情况，在一开始就直接使用了 \(q \neq 1\) 时的求和公式，相对不够严谨。更严谨的做法是：\\
【方法二】直接写成求和的形式，即

\[
\begin{aligned}
& S_{5}=1+q+q^{2}+q^{3}+q^{4} \\
& S_{3}=1+q+q^{2} .
\end{aligned}
\]

所以

\[
\begin{aligned}
& S_{5}=5 S_{3}-4 \Longrightarrow 1+q+q^{2}+q^{3}+q^{4}=5\left(1+q+q^{2}\right)-4 \\
& q^{4}+q^{3}-4 q^{2}-4 q=0 \\
& q\left(q^{3}+q^{2}-4 q-4\right)=0 \\
& q(q+1)\left(q^{2}-4\right)=0 .
\end{aligned}
\]

以上若干步的分解因式在前面的代数杜杆中都有所涉及。由于 \(q>0\) ，所以 \(q=2\) ．也可以得到正确答案。

【方法三】利用齐次结构。所给的已知条件，并不是齐次式：4 是零次，而 \(S_{5}, S_{3}\) 都是关于项数的一次式。因此，我们可以用 4 S ，来替代 4 ：

\[
S_{5}=5 S_{3}-4 S_{1} .
\]

故

\[
S_{5}-S_{3}=4\left(S_{3}-S_{1}\right)
\]

而 \(S_{5}-S_{3}=q^{2}\left(S_{3}-S_{1}\right)\) ，所以 \(q^{2}=4\) ，又由于各项均为正数，所以 \(q=2\) 。关于 \(S_{5}-S_{3}\) 刚好是 \(S_{3}-S_{1}\) 的 \(q^{2}\) 倍，这在我们的讲解数列时提到过：

\[
\begin{aligned}
& S_{3}-S_{1}=a_{2}+a_{3}, \\
& S_{5}-S_{3}=a_{4}+a_{5}=q^{2}\left(a_{2}+a_{3}\right), \\
& \frac{S_{5}-S_{3}}{S_{3}-S_{1}}=\frac{q^{2}\left(a_{2}+a_{3}\right)}{a_{2}+a_{3}}=q^{2} .
\end{aligned}
\]

在这里，我们不光使用了正向检查，还是使用了反向检查。\\
6．某地的中学生中有 \(60 \%\) 的同学爱好滑冰， \(50 \%\) 的同学爱好滑雪， \(70 \%\) 的同学爱好滑冰或爱好滑雪。在该地的中学生中随机调查一位同学，若该同学爱好滑雪，则该同学也爱好滑冰的概率为 \(\qquad\) －\\
A， 0.8\\
B， 0.6\\
C， 0.5\\
D， 0.4

可以用韦恩图来表示有这两个爱好的中学生的比例：根据图示，既喜欢滑冰，又喜欢滑雪的同学所占的比例为\\
\includegraphics[max width=\textwidth, center]{2025_02_16_7c9aaf92bc99dc07f048g-306}\\
\(50 \%+60 x \%-70 \%=40 \%\) ．设事件 \(A\) 为该同学爱好滑冰，事件 \(B\) 为该同学爱好滑雪，那么

\[
\mathbb{P}(A)=0.6, \mathbb{P}(B)=0.5, \mathbb{P}(A \cup B)=0.7, \mathbb{P}(A \cap B)=0.4
\]

根据题意，所求为条件概率：

\[
\mathbb{P}(A \mid B)=\frac{\mathbb{P}(A \cap B)}{\mathbb{P}(B)}=\frac{0.4}{0.5}=0.8
\]

这道题中涉及计算的部分较少，主要是对于条件概率这个概念本身的理解，所以不需要过分检查。\\
7．设甲： \(\sin ^{2} \alpha+\sin ^{2} \beta=1, Z: \sin \alpha+\cos \beta=0\) ，则：\\
A，甲是乙的充分但不是必要条件\\
B，甲是乙的必要但不是充分条件\\
C，甲是乙的充要条件\\
D，甲既不是乙的充分条件也不是乙的必要条件

根据勾股定理，甲命题等价于 \(\sin ^{2} \alpha=1-\sin ^{2} \beta=\cos ^{2} \beta\) ，故

\[
(\sin \alpha-\cos \beta)(\sin \alpha+\cos \beta)=0
\]

\begin{center}
\includegraphics[max width=\textwidth]{2025_02_16_7c9aaf92bc99dc07f048g-307}
\end{center}

第 3 题图\\
\includegraphics[max width=\textwidth, center]{2025_02_16_7c9aaf92bc99dc07f048g-307(1)}

第 18 题图

而乙命题是＂ \(\sin \alpha+\cos \beta=0\)＂，所以乙成立时，甲成立；但是甲成立时，乙不一定成立。故甲是乙的必要不充分条件。

8．已知双曲线 \(C: \frac{x^{2}}{a^{2}}-\frac{y^{2}}{b^{2}}=1(a>0, b>0)\) 的离心率为 \(\sqrt{5}, c\) 的一条渐近线与圆 \((x-2)^{2}+(y-3)^{2}=1\) 交于 \(A, B\)两点，则 \(|A B|=\) \(\qquad\)\\
A，\(\frac{\sqrt{5}}{5}\)\\
B，\(\frac{2 \sqrt{5}}{5}\)\\
C，\(\frac{3 \sqrt{5}}{5}\)\\
D，\(\frac{4 \sqrt{5}}{5}\)

根据已知条件，\(c=\sqrt{5} a\) ．故

\[
c^{2}=a^{2}+b^{2}=5 a^{2} .
\]

故 \(b^{2}=4 a^{2}\) ，即 \(b=2 a\) ．那么双曲线的渐近线方程为 \(y=2 x\) 和 \(y=-2 x\) ．显然，\(y=2 x\) 和圆 \((x-2)^{2}+(y-3)^{2}=1\)有焦点。这里可以使用正向检查法，因为我们同样有多于一种方法。\\
【方法一】利用圆的性质。圆心为 \(D(2,3)\) ，设 \(D\) 到直线 \(y=2 x\) 的距离为 \(d\) ，那么根据垂径定理，

\[
\left(\frac{|A B|}{2}\right)^{2}+d^{2}=r^{2}
\]

所以 \(|A B|=2 \sqrt{r^{2}-d^{2}}=2 \sqrt{1-d^{2}}\) ．只需求 \(d\) 即可。根据点到直线距离公式，有

\[
d=\frac{|3-4|}{\sqrt{5}}=\frac{\sqrt{5}}{5} .
\]

故

\[
|A B|=2 \cdot \sqrt{1-\left(\frac{1}{5}\right)}=\frac{4 \sqrt{5}}{5} .
\]

【方法二】联立。虽然直接求出 \(A, B\) 两点坐标比较复杂，但是可以通过联立的方法得到它们之间的关系。将 \(y=2 x\) 和圆方程联立，得到

\[
\begin{aligned}
& (x-2)^{2}+(2 x-3)^{2}=1 \\
& 5 x^{2}-16 x+12=0 .
\end{aligned}
\]

为了求 \(|A B|\) ，可以利用两点间距离公式：

\[
|A B|=\sqrt{k^{2}+1} \cdot\left|x_{A}-x_{B}\right|=\sqrt{5} \cdot\left|x_{A}-x_{B}\right| .
\]

同时，

\[
\left|x_{A}-x_{B}\right|=\frac{\sqrt{\Delta}}{5}=\frac{\sqrt{16^{2}-4 \cdot 5 \cdot 12}}{5}=\frac{4}{5} .
\]

故

\[
|A B|=\sqrt{5}\left|x_{A}-x_{B}\right|=\frac{4 \sqrt{5}}{5} .
\]

【注】这里最好不要直接求 \(A, B\) 两点坐标，计算量相对较大。

9．现有 5 名志愿者报名参加公益活动，在某一星期的星期六，星期日两天，每天从这 5 人中安排 2 人参加公益活动，则恰有 1 人在这两天都参加的不同安排方式共有 \(\qquad\) －\\
A， 120 种\\
B， 50 种\\
C， 30 种\\
D， 20 种

我们可以把所要做的事情拆解为：\\
（1）首先选择一人，参加星期六和星期天两天的活动，共 5 种选择；\\
（2）再选择一人，参加星期六的公益活动，（1）中已选走 1 人，共 4 种选择；\\
（3）再选择一人，参加星期日的活动，（1），2 已选走 2 人，共 3 种选择。这三步构成了整个事情，根据分步乘法原理：总方法数为 \(5 \times 4 \times 3=60\) 。

我们可以使用结构检查法：答案必然是 10 的倍数，因为我们可以整体将这三人替换成另外的三人。而从五人中选择三人，共有 \(\mathrm{C}_{5}^{3}=10\) 种选法。\\
同样，可以正向检查：我们发现，相当于这三个人扮演了三个不同的岗位：两天都去，星期六去和星期天去。我们只需要选择三人，让三人分别担任这些岗位。共种数为 \(\mathrm{A}_{5}^{3}=60\) ，或理解为 \(\mathrm{C}_{5}^{3} \cdot \mathrm{~A}_{3}^{3}\) 也可以。

10．函数 \(y=f(x)\) 的图像由函数 \(y=\cos \left(2 x+\frac{\pi}{6}\right)\) 的图像向左平移 \(\frac{\pi}{6}\) 个单位长度得到，则 \(y=f(x)\) 的图像与直线 \(y=\frac{1}{2} x-\frac{1}{2}\) 的交点个数为 \(\qquad\) －\\
A， 1\\
B， 2\\
C， 3\\
D， 4

我们可以直接利用平移的性质，将图像向左平移，相当于将 \(x\) 换为 \(x+\frac{\pi}{2}\) ：

\[
f(x)=\cos \left(2\left(x+\frac{\pi}{6}\right)+\frac{\pi}{6}\right)=\cos \left(2 x+\frac{\pi}{2}\right)=-\sin 2 x .
\]

由于有相位，所以可以考虑用诱导公式去掉相位。再画出 \(y=\frac{1}{2} x-\frac{1}{2}\) 的图像，发现它们有三个交点。故选 C．这里可以用代入法检查。如果不确定平移后函数的解析式，可以令 \(x=\frac{\pi}{6}\) 代入计算。此时 \(\cos \left(2 x+\frac{\pi}{6}\right)=0\) ，所以平移前，函数图像经过 \(\left(\frac{\pi}{6}, 0\right)\) ．平移后，这个点坐标为 \((0,0)\) ，确实在 \(y=-\sin 2 x\) 的图像上。所以平移无误。

11．已知四棱雉 \(P-A B C D\) 的底面是边长为 4 的正方形，\(P C=P D=3, \angle P C A=45^{\circ}\) ，则 \(\triangle P B C\) 的面积为 \(\qquad\) －\\
A， \(2 \sqrt{2}\)\\
B， \(3 \sqrt{2}\)\\
C， \(4 \sqrt{2}\)\\
D， \(5 \sqrt{2}\)

我们可以从自由度的角度来理解。四棱雉 \(P-A B C D\) 想要确定通过，只需要确定 \(P\) 点的位置。题目条件有两个，两条直线可以确定一个点，那么 \(P\) 点是如何确定的呢？\\
（1）\(P C=P D=3\) ，那么三角形 \(P C D\) 是一个边长为 \(3,3,4\) 的等腰三角形。设 \(C D\) 中点为 \(E\) ，那么 \(P E\) 的长度就是定值。我们可以想象，整个三角形 \(P C D\) 绕着 \(C D\) 边转动，\(P\) 的轨迹是一个圆：它以 \(E\) 为圆心，且这个平面垂直平分底面的正方形。\\
（2）\(\angle P C A=45^{\circ}\) ，由于 \(P C\) 长度确定了，那么三角形 \(P C A\) 也确定了：它是一个某角为 \(45^{\circ}\) ，两条边长为 \(3,4 \sqrt{2}\) 。因此，也可以想象整个三角形绕着 AC 边转动， P 的轨迹也是一个圆。\\
由（1），（2）中的两个圆相交，就确定了 \(P\) 点的位置。显然，这两个圆有两个交点，但是它们关于底面 \(A B C D\) 对称，所以只考虑一种情况即可。

为了确定三角形 PBC 的面积，由于 \(P C, B C\) 已知，根据自由度的理论，只需要再确定 \(P B\) 即可。根据（1）中的分析，\(P\) 所在平面垂直平分 \(A B C D\) ，因此 \(P A=P B\) 。那么，可以根据（2）中的结论求解 \(P A, P B\) 。根据余弦定理，

\[
P A=\sqrt{|P C|^{2}+|A C|^{2}-2|P C| \cdot|A C| \cdot \cos \angle P C A}=\sqrt{9+32-\sqrt{2} \cdot 3 \cdot 4 \sqrt{2}}=\sqrt{17} .
\]

因此，\(P B=P A=\sqrt{17}\) ．三角形 \(P C B\) 的三边长为 \(3,4, \sqrt{17}\) ，可以用勾股定理求它的面积。过 P 作 \(B C\) 垂线，垂足为 \(F\) ，设 \(C F=x\) ，那么

\[
\begin{aligned}
& h^{2}=P C^{2}-C D^{2}=P B^{2}-B D^{2} \\
& 17-(4-x)^{2}=3-x^{2}
\end{aligned}
\]

解得 \(x=1\) ．故 \(h=2 \sqrt{2}\) ，所以

\[
S_{\triangle P B C}=\frac{|B C| \cdot h}{2}=4 \sqrt{2}
\]

12．设 \(O\) 为坐标原点，\(F_{1}, F_{2}\) 为椭圆 \(C: \frac{x^{2}}{9}+\frac{y^{2}}{6}=1\) 的两个焦点，\(P\) 在 \(C\) 上， \(\cos \angle F_{1} P F_{2}=\frac{3}{5}\) ，则 \(|P O|=\) \(\qquad\) \(-\)\\
A，\(\frac{13}{5}\)\\
B，\(\frac{\sqrt{30}}{2}\)\\
C，\(\frac{14}{5}\)\\
D，\(\frac{\sqrt{35}}{2}\)

这是我们在章节＜代数结构的实战应用＞8 中学习过的内容。根据自由度的理论，三角形 \(P F_{1} F_{2}\) 已经完全确定了：因为有三个关于它的已知条件：（1）\(F_{1} F_{2}=2 c=2 \sqrt{3}\) ；（2）\(P F_{1}+\lim _{n \rightarrow \infty} \frac{1}{n^{2}} P F_{2}=6\) ；（3） \(\cos \angle F_{1} P F_{2}=\frac{3}{5}\) 。因此，我们可以利用这几个条件确定 \(|P O|\) ．\\
在这里，我们有两种选择：求 \(P\) 点坐标，或者求 \(P F_{1}, P F_{2}\) 的长度。在这里，直接求解边长更方便，因为（1）， （2）两个条件都是关于边长的，（3）余弦值也可以很快用余弦定理翻译。设 \(P F_{1}=x, P F_{2}=y\) ，那么

\[
\left\{\begin{array}{l}
x+y=6 \\
\left|F_{1} F_{2}\right|^{2}=x^{2}+y^{2}-\frac{6}{5} x y=12
\end{array}\right.
\]

这样就得到了一个关于 \(x, y\) 对称的方程组。注意到所求是 \(|O P|\) ，根据中线长公式，

\[
|P O|^{2}+\left|F_{1} O\right|^{2}=\frac{1}{2} \cdot\left(\left|P F_{1}\right|^{2}+\left|P F_{2}\right|^{2}\right)=\frac{1}{2}\left(x^{2}+y^{2}\right)
\]

故只需要求 \(x^{2}+y^{2}\) 即可。根据 \(x+y=6\) ，可得

\[
x^{2}+y^{2}+2 x y=(x+y)^{2}=36
\]

故我们可以将它看成关于 \(m=x^{2}+y^{2}, n=x y\) 的方程，即

\[
\left\{\begin{array}{l}
m+2 n=36 \\
m-\frac{6}{5} n=12
\end{array}\right.
\]

解得 \(n=\frac{15}{2}, m=21\) ．故

\[
|P O|^{2}=\frac{1}{2} m-3=\frac{15}{2},|P O|=\sqrt{\frac{15}{2}}=\frac{\sqrt{30}}{2} .
\]

可以简单地【代人检查】，当 \(m=21, n=\frac{15}{2}\) 时，

\[
m+2 n=21+15=36, m-\frac{6}{5} n=21-\frac{6}{5} \cdot \frac{15}{2}=12 .
\]

\section*{7．1．7 填空题 13－16}
13．若 \(f(x)=(x-1)^{2}+a x+\sin \left(x+\frac{\pi}{2}\right)\) 为偶函数，则 \(a=\) \(\qquad\) －\\
可以根据定义求解。首先注意到 \(\sin \left(x+\frac{\pi}{2}\right)=\cos x\) ，那么

\[
f(-x)=(-x-1)^{2}-a x+\cos (-x) .
\]

由于 \(f(x)\) 是偶函数，所以

\[
\begin{aligned}
& f(x)=f(-x) \\
& (x-1)^{2}+a x+\cos x=(-x-1)^{2}-a x+\cos (-x) \\
& x^{2}-2 x+1+a x+\cos x=x^{2}+2 x+1-a x+\cos x \\
& (2 a-4) x=0 .
\end{aligned}
\]

注意到，这对于任意 \(x \in \mathbb{R}\) 都成立，所以 \(2 a-4=0\) ，即 \(a=2\) ．对于奇函数，偶函数性质更为熟悉的同学，可以通过【结构检查】法来检查：我们可以将 \(f(x)\) 写成奇函数，偶函数的和：

\[
f(x)=x^{2}-2 x+1+a x+\cos x=\left(x^{2}+\cos x+1\right)+(2-a) x .
\]

由于 \(y=x^{2}, y=\cos x, y=1\) 都是偶函数，所以 \(y=x^{2}+\cos x+1\) 也是偶函数。为了让 \(f(x)\) 也是偶函数，必须要求奇函数部分 \((2-a) x\) 是 0 函数。所以 \(a=2\) ．直接【代人检查】也可以验证答案，当 \(a=2\) 时，

\[
f(x)=x^{2}+1+\cos x
\]

确实是偶函数。\\
14．若 \(x, y\) 满足约束条件 \(\left\{\begin{array}{l}3 x-2 y \leq 3 \\ -2 x+3 y \leq 3 \text { ，设 } z=3 x+2 y \text { 的最大值为 } \\ x+y \geq 1\end{array}\right.\) \(\qquad\) ．

这是线性规划问题。我们需要将直线都在直角坐标系中画出。我们关心的三条支线为：

\[
l_{1}: 3 x-2 y=3, l_{2}:-2 x+3 y=3, l_{3}: x+y=1 .
\]

在绘图过程中，尽量保证图形准确，我们需要求解三条支线的两两交点。\\
（1）\(l_{1}, l_{2}\) 的交点满足

\[
\left\{\begin{array}{l}
3 x-2 y=3 \\
-2 x+3 y=3
\end{array}\right.
\]

解得 \(x=3, y=3\) ．所以交点为 \((3,3)\) ．\\
（2）\(l_{1}, l_{3}\) 的交点满足

\[
\left\{\begin{array}{c}
3 x-2 y=3 \\
x+y=1
\end{array}\right.
\]

解得 \(x=1, y=0\) ．所以交点为 \((1,0)\) ．\\
（3）\(l_{2}, l_{3}\) 的交点满足

\[
\left\{\begin{array}{c}
-2 x+3 y=3 \\
x+y=1
\end{array}\right.
\]

解得 \(x=0, y=1\) ．所以交点为 \((0,1)\) ．\\
由这三个不等式构成的区域是哪部分呢？\\
（1） \(3 x-2 y \leq 3\) ，所以 \(y \geq \frac{3}{x}-32\) ，因此它表示直线 \(l_{1}\) 上方的位置；\\
（2）\(-2 x+3 y \leq 3\) ，所以 \(y \leq \frac{2}{x}+33\) ，因此它表示直线 \(l_{2}\) 下方的位置；\\
（3）\(x+y \geq 1\) ，所以 \(y \geq-x+1\) ，因此它表示直线 \(l_{3}\) 上方的位置。\\
这样一来，不等式组代表的区域是这三条直线围成的三角形区域。这里，我们可以【代人检查】来验证我们的判断。我们看 \(P(1,1)\) 是否满足题意：当 \(x=1, y=1\) 时，

\[
3 x-2 y=1 \leq 3,-2 x+3 y=1 \leq 3, x+y=2 \geq 1
\]

成立。所以 \(P(1,1)\) 满足题意，而且它确实在这个三角形区域内，所以我们的结论正确。\\
为了求 \(3 x+2 y\) 的最大值，可以令 \(z=3 x+2 y\) 然后看截距，但是这里我们注意到：在这个区域内，\(x \leq 3, y \leq 3\) ，而且 \(x=3, y=3\) 可以同时取等，所以

\[
3 x+2 y \leq 3 \cdot 3+2 \cdot 3=15
\]

当 \(x=y=3\) 时取等。

15．在正方体 \(A B C D-A_{1} B_{1} C_{1} D_{1}\) 中，\(E, F\) 分别为 \(A B, C_{1} D_{1}\) 的中点，以 \(E F\) 为直径的球的球面与该正方体的棱共有 \(\qquad\)个公共点．

设正方体边长为 2 。显然，这个球面的球心恰为正方体的中心 \(o\) ，且半径为 \(\frac{E F}{2}=\sqrt{2}\) ．所以我们可以在每条棱上找到距离正方体中心 \(\sqrt{2}\) 的点。我们发现，由于 \(O A=O B\) ，所以 \(O E \perp A B\) ，所以 \(O E\) 恰好就是 \(O\) 到 \(A B\) 的距离。那么，根据对称性，\(O\) 到每条棱的距离都是 \(\sqrt{2}\) ，垂足都是中点。正方体共有 12 条棱，所以共有 12 个交点。\\
如果对结果有顾虑，可以用【正向检查】，建系处理。以 \(D\) 为坐标原点， \(\overrightarrow{D A}, \overrightarrow{D C}, \overrightarrow{D D_{1}}\) 方向为 \(x, y, z\) 轴正方向建立空间直角坐标系。那么正方体中心为 \(O(1,1,1)\) ，我们可以检查若干个中点到 \(o\) 的距离：由于 \(E(2,1,0), F(0,1,2)\) ，所以

\[
|E F|=\sqrt{2^{2}+0^{2}+2^{2}}=2 \sqrt{2} .
\]

故半径为 \(\sqrt{2}\) ，正确。棱 \(C C_{1}\) 的中点 \(M\) 为 \((0,2,1)\) ，那么

\[
|O M|=\sqrt{1^{2}+1^{2}+0^{2}}=\sqrt{2}
\]

正确。

16．在 \(\triangle A B C\) 中，\(\angle B A C=60^{\circ}, A B=2, B C=\sqrt{6}, \angle B A C\) 的角平分线交 \(B C\) 于 \(D\) ，则 \(A D=\) \(\qquad\) ．

首先从自由度的角度理解这个题目中的三角形。共有三个已知条件，所以三角形＂确定＂了（有不超过 2 种可能）。然而，注意到这三个条件是 SSA，并不符合全等三角形的判定定理，所以要小心是否需要分类讨论。根据章节＜代数结构的实战应用＞中的讨论，由于 \(B C=\sqrt{6}>A B=2\) ，所以不需要分类讨论，这个图形唯一确定。

处理此题的核心就是处理角平分线。首先，我们求出 \(A C\) 边的长度：过 \(B\) 作 \(A C\) 垂线，垂足为 \(E\) ．那么

\[
A E=A B \cdot \cos 60^{\circ}=1, C E=\sqrt{B C^{2}-B E^{2}}=\sqrt{3}
\]

故 \(A C=A E+C E=\sqrt{3}+1\) ．最自然的方法，就是设 \(A D=x\) ，利用余弦定理。\\
【方法一】根据角平分线定理，

\[
\frac{B D}{C D}=\frac{A B}{A C}=\frac{2}{\sqrt{3}+1}=\sqrt{3}-1
\]

所以

\[
\frac{C D}{B C}=\frac{C D}{C D+B D}=\frac{1}{\sqrt{3}} .
\]

所以 \(C D=\sqrt{2}\) ．根据余弦定理，

\[
C D^{2}=A D^{2}+A C^{2}-2 A C \cdot A D \cdot \cos \angle C A D .
\]

所以

\[
\begin{gathered}
x^{2}+(\sqrt{3}+1)^{2}-2 x \cdot(\sqrt{3}+1) \cdot \frac{\sqrt{3}}{2}=2 \\
x^{2}-(3+\sqrt{3}) x+2+2 \sqrt{3}=0
\end{gathered}
\]

（1）求根公式：\(\Delta=(3+\sqrt{3})^{2}-4(2+2 \sqrt{3})=4-2 \sqrt{3}=(\sqrt{3}-1)^{2}\) ．所以

\[
x_{1,2}=\frac{3+\sqrt{3} \pm(\sqrt{3}-1)}{2}
\]

即 \(x_{1}=\sqrt{3}+1, x_{2}=2\) ．\\
（2）分解因式：我们可以直接观察到 \(2+2 \sqrt{3}=2(\sqrt{3}+1)\) ，所以直接十字相乘，

\[
x^{2}-(3+\sqrt{3}) x+2+2 \sqrt{3}=(x-2)(x-\sqrt{3}-1) .
\]

即 \(x_{1}=\sqrt{3}+1, x_{2}=2\) ．也可以选 \(\sqrt{3}\) 为主元：

\[
\begin{aligned}
x^{2}-(3+\sqrt{3}) x+2+2 \sqrt{3} & =(2-x) \sqrt{3}+x^{2}-3 x+2 \\
& =(2-x) \sqrt{3}+(x-1)(x-2) \\
& =(x-2)(x-1-\sqrt{3}) .
\end{aligned}
\]

也可以得到正确答案。\\
总而言之，这里有两个根。但是，根据前面的分析，我们知道，答案必然是唯一确定的，所以要舍掉较大的根 \(x_{1}\) ，保留 \(x_{2}=2\) ．这是因为：我们可以从图中看到，直线 \(A D\) 上确实有两个点，可以使得 \(C D=\sqrt{2}\) ．而由于 \(\angle A D C\)是钝角，所以我们应该取靠近 \(A\) 的，即较小根。所以 \(A D=2\) ．\\
这道题我们可以【正向检查】或者【代人检查】。如果代人检查，我们就需要把这个三角形中的各个角度等求解出来去验证，这个方法可以做到。但是，我们也可以使用另一个方法：\\
【方法二】正弦定理。注意到 \(\sqrt{6}\) 所对边的角度 \(60^{\circ}\) 已经确定，又知道了 \(A B=2\) ，所以可以直接确定 \(\angle A C B\) 的大小。记 \(A B=c, B C=a, A C=b\) ．那么，根据正弦定理，

\[
\frac{c}{\sin C}=\frac{a}{\sin \angle B A C} .
\]

所以

\[
\sin C=\frac{A B}{B C} \cdot \sin \angle B A C=\frac{2}{\sqrt{6}} \cdot \frac{\sqrt{3}}{2}=\frac{\sqrt{2}}{2} .
\]

所以 \(C=45^{\circ}\) 或 \(135^{\circ}\) ．而 \(\angle B A C=60^{\circ}\) ，所以 \(C=45^{\circ}\) 。因此，这个三角形是一个特殊度数的三角形，我们可以求出题目中的所有角度：

\[
\begin{gathered}
\angle B=180^{\circ}-\angle C-\angle B A C=75^{\circ} \\
\angle A D B=\angle C+\frac{1}{2} \angle B A C=75^{\circ} .
\end{gathered}
\]

我们惊喜地发现，三角形 ADB 是一个等腰三角形。所以

\[
A D=A B=2
\]

直接得到了答案。\\
【方法三】角平分线长可以计算

\[
A D^{2}=A B \cdot A C-B D \cdot C D
\]

根据方法一中的计算，\(A C=\sqrt{3}=1, C D=\sqrt{2}, B D=\sqrt{6}-\sqrt{2}\) ，所以

\[
A D^{2}=2(\sqrt{3}+1)-\sqrt{2}(\sqrt{6}-\sqrt{2})=4
\]

所以 \(A D=2\) ．\\
【注】这道题平面几何色彩非常重，需要对解三角形，自由度有一个非常好的理解才可以快速完成。

\section*{7．1．8 解答题 17－21}
17．设 \(S_{n}\) 为数列 \(\left\{a_{n}\right\}\) 的前 \(n\) 项和，已知 \(a_{2}=1,2 S_{n}=n a_{n}\) ．\\
（1）求 \(\left\{a_{n}\right\}\) 的通项公式；\\
（2）求数列 \(\left\{\frac{a}{n+1} 2^{n}\right\}\) 的前 \(n\) 项和 \(T_{n}\) ．\\
（I）最直接的方法，就是看到 \(S_{n}, a_{n}\) ，就利用 \(S_{n}-S_{n}-1=a_{n}\) 即可。将 \(n\) 替换成 \(n-1\) 代人，知

\[
2 S_{n-1}=(n-1) a_{n-1}
\]

用 \(2 S_{n}=n a_{n}\) 和它相减，得到

\[
\begin{aligned}
& 2 a_{n}=n a_{n}-(n-1) a_{n-1} \\
& (n-2) a_{n}=(n-1) a_{n-1}
\end{aligned}
\]

这和＂等比数列＂的递推式接近，可以考虑累乘；更快的方法是，注意到 \(n-1, n-2\) 两个角标也差 1 ，所以，当 \(n \neq 2\) 时，

\[
\frac{a_{n}}{n-1}=\frac{a_{n-1}}{n-2}
\]

所以，\(b_{n}=\frac{a_{n}}{n-1}\) 是一个常数列。而 \(b_{2}=\frac{a_{2}}{1}=1\) ，所以

\[
\frac{a_{n}}{n-1}=b_{n}=b_{2}=1
\]

所以 \(a_{n}=n-1\) ．注意，由于我们带入了 \(S_{n-1}\) ，所以这只对于 \(n-1 \geq 1\) 即 \(n \geq 2\) 成立。我们还需要单独检验 \(n=1\) ．将 \(n=1\) 代人已知条件，

\[
2 S_{1}=a_{1}
\]

所以 \(a_{1}=0=1-1\) ．所以，\(a_{n}=n-1\) 对 \(n \geq 1\) 都成立，故 \(a_{n}=n-1\) ．

我们可以用【代入检查】法进行检查：当 \(a_{n}=n-1\) 时，数列 \(\left\{a_{n}\right\}\) 是一个等差数列，所以

\[
S_{n}=\frac{a_{1}+a_{n}}{2} \cdot n=\frac{(n-1) n}{2}
\]

所以

\[
2 S_{n}=n(n-1)=n a_{n} .
\]

也可以【结构检查】来理解。我们观察到 \(2 S_{n}=n a_{n}\) 的左侧是 \(S_{n}\) ，右侧有 \(n\) ，所以可以变形成成

\[
\frac{S_{n}}{n}=\frac{a_{n}}{2}
\]

要知道，＂如果 \(\left\{a_{n}\right\}\) 是等差数列＂，那么 \(\frac{s}{n} n\) 代表了前 \(n\) 项的平均数（章节＜代数结构的实战应用＞中讨论过），那么

\[
\frac{S_{n}}{n}=\frac{a_{n}+a_{1}}{2}
\]

而题目条件是 \(\frac{S_{n}}{n}=\frac{a_{n}}{2}\) ，所以 \(a_{1}=0\) 。所以，＂如果 \(\left\{a_{n}\right\}\) 是等差数列＂，必然 \(a_{1}=0\) 。又 \(a_{2}=1\) ，这个等差数列必为 \(a_{n}=n-1\) ，这是符合题意的。注意，这个不是证明：因为前提条件是＂\(\left\{a_{n}\right\}\) 是等差数列＂，但是我们做题时不知道这一点。所以，它只能用来帮我们验证和理解问题本身的结构。\\
（II）根据（I）的结果，

\[
\frac{a_{n+1}}{2^{n}}=\frac{n}{2^{n}}
\]

错位相减即可。

\[
\begin{gathered}
T_{n}=\frac{1}{2}+\frac{2}{2^{2}}+\frac{3}{2^{3}}+\cdots+\frac{n-1}{2^{n-1}}+\frac{n}{2^{n}} \\
2 T_{n}=1+\frac{2}{2}+\frac{3}{2^{2}}+\frac{4}{2^{3}}+\cdots+\frac{n}{2^{n-1}}
\end{gathered}
\]

下式减上式，得

\[
T_{n}=1+\frac{1}{2}+\frac{1}{2^{2}}+\frac{1}{2^{3}}+\cdots+\frac{1}{2^{n-1}}-\frac{n}{2^{n}}=2-\frac{1}{2^{n-1}}-\frac{n}{2^{n}}=2-\frac{n+2}{2^{n}}
\]

可以【代入检查】：（1）当 \(n=1\) 时，

\[
T_{1}=\frac{a_{2}}{2}=\frac{1}{2}
\]

根据上面的结果，\(T_{1}=2-\frac{3}{2}=\frac{1}{2}\) ，成立；\\
（2）当 \(n=3\) 时，

\[
T_{3}=\frac{1}{2}+\frac{2}{2^{2}}+\frac{3}{2^{3}}=\frac{1}{2}+\frac{1}{2}+\frac{3}{8}=\frac{11}{8},
\]

根据上面的结果，\(T_{3}=2-\frac{5}{8}=\frac{11}{8}\) ，成立。\\
也可以【反向检查】，若

\[
T_{n}=2-\frac{n+2}{2^{n}}
\]

那么

\[
T_{n}-T_{n-1}=\left(2-\frac{n+2}{2^{n}}\right)-\left(2-\frac{n+1}{2^{n-1}}\right)=\frac{2 n+2}{2^{n}}-\frac{n+2}{2^{n}}=\frac{n}{2^{n}}=\frac{a_{n+1}}{2^{n}} .
\]

符合题意。

18．如图，在三棱柱 \(A B C-A_{1} B_{1} C_{1}\) 中，\(A_{1} C \perp\) 底面 \(A B C, \angle A C B=90^{\circ}, A A_{1}=2, A_{1}\) 到平面 \(B C C_{1} B_{1}\) 的距离为 1 ．\\
（1）证明：\(A_{1} C=A C\) ；\\
（2）已知 \(A A_{1}\) 与 \(B B_{1}\) 的距离为 2 ，求 \(A B_{1}\) 与平面 \(B C C_{1} B_{1}\) 所成角的正弦值．\\
（I）题目给了 4 个条件，其中最难利用的就是 \(A_{1}\) 到平面 \(B C C_{1} B_{1}\) 的距离。我们可以聚焦这个条件考虑。最自然的思路，就是把垂线作出来。但是，直接做点到平面的垂线，我们很难确定垂足的位置，所以我们可以先考察图形的性质。

我们发现，前面两个条件都是和垂直有关：线面垂，线线垂。由于 \(A_{1} C 1\) 平面 \(A B C\) ，所以 \(B C \perp C A_{1}, A C \perp C A_{1}\) ；又因为 \(\angle A C B=90^{\circ}\) ，所以 \(B C, A C, C A_{1}\) 三条直线两两垂直，我们可以依托它建立空间直角坐标系，这是一种可能的方法。\\
同时，我们发现：\(B C\) 垂直于平面 \(A_{1} A C C_{1}\) 。根据面面垂的判定定理，平面 \(A_{1} C C_{1}\) 和平面 \(B C C_{1} B_{1}\) 垂直。这样的话，我们只需要过 \(A_{1}\) 作 \(C C_{1}\) 的垂线，这个垂线就是 \(A_{1}\) 到平面 \(B C C_{1} B_{1}\) 的垂线。这样的话，垂足就可以确定位置了。再利用 \(A A_{1}=2\) 和距离为 1 这个条件，问题就解决了。这部分的证明大部分方法都会使用，我们先展示出来：\\
【证明】 \(\because A_{1} C \perp\) 平面 \(A B C, \therefore B C \perp C A_{1}, A C \perp C A_{1}\) ．\\
\(\because \angle A C B=90^{\circ}\) ，所以 \(B C \perp A C\) ．\\
\(\therefore B C \perp\) 平面 \(A A_{1} C_{1} C\) ．\\
\(\because\) 平面 \(B B_{1} C_{1} C\) 过直线 \(B C\) ，根据面面垂的判定定理，平面 \(B B_{1} C_{1} C\) 和平面 \(A A_{1} C_{1} C\) 垂直。\\
【方法一】直接作垂线。过 \(A_{1}\) 作 \(C C_{1}\) 垂线，垂足为 \(O_{1}\) ，根据面面垂性质定理，\(A_{1} O \perp\) 平面 \(B C C_{1} B_{1}\) ．所以 \(A_{1} O=1\) ．又因为 \(A_{1} A=2\) ，所以 \(C C_{1}=A A_{1}=2\) ．考察三角形 \(A_{1} C_{1} C_{1} C\) ，由于 \(A_{1} C_{1} / / A C\) ，所以 \(\angle C A_{1} C_{1}=90^{\circ}\) ．故三角形 \(A_{1} C_{1} C_{1} C\) 是直角三角形。其斜边为 \(C C_{1}=2\) ，斜边上的高为 \(A_{1} O=1\) 。那么根据三角形面积公式：

\[
A_{1} C_{1} \cdot A_{1} C=A_{1} O \cdot \bar{C} C_{1}=2
\]

又根据勾股定理：

\[
\left|A_{1} C_{1}\right|^{2}+\left|A_{1} C\right|^{2}=\left|C C_{1}\right|^{2}=4
\]

解得 \(A_{1} C_{1}=A_{1} C=\sqrt{2}\) ．所以 \(A C=A_{1} C_{1}=A_{1} C=\sqrt{2}\) ．\\
【方法二】建系。以 \(C\) 为原点， \(\overrightarrow{C A}, \overrightarrow{C B}, \overrightarrow{C A_{1}}\) 为 \(x, y, z\) 轴正方向建立空间直角坐标系。\\
设 \(A C=x, C A_{1}=z\) ，那么根据勾股定理，\(x^{2}+z^{2}=A A_{1}^{2}=4\) ．

\[
A(x, 0,0), A_{1}(0,0, z), C_{1}(-x, 0, z)
\]

设平面 \(B B_{1} C_{1} C\) 的一个法向量为 \(\vec{n}=(a, b, c)\) ，那么

\[
\left\{\begin{array}{c}
\vec{n} \cdot \overline{C B}=0 \Longrightarrow b=0 \\
\vec{n} \cdot \overline{C C_{1}}=0 \Longrightarrow-a x+c z=0
\end{array}\right.
\]

令 \(a=z\) ，那么 \(c=x\) ，所以 \(\vec{n}=(z, 0, x)\) 。由于 \(A_{1}\) 到平面 \(B B_{1} C_{1} C\) 的距离为 1 ，所以 \(C A_{1}\) 在 \(\vec{n}\) 方向上的投影为 1 ，即

\[
\frac{\left|\overrightarrow{C A_{1}} \cdot \vec{n}\right|}{|\vec{n}|}=1
\]

即

\[
\frac{|x z|}{\sqrt{x^{2}+z^{2}}}=1
\]

由于 \(x^{2}+z^{2}=4\) ，所以 \(|x z|=2\) ．解得 \(x=z=\sqrt{2}\) ．所以

\[
A_{1} C=A C=\sqrt{2} .
\]

【方法三】利用体积。由于 \(A_{1}\) 到平面 \(B C C_{1} B_{1}\) 的距离为 1 这个条件不好翻译，所以我们想到了利用体积翻译。我们用两种方法计算体积。首先，设 \(A C=x, B C=y, A_{1} C=z\) 。\\
（1）以 \(B B_{1} C_{1} C\) 为底面。由于 \(B C \mathrm{~L}\) 平面 \(A C C_{1}\) ，所以 \(B C \perp C C_{1}\) ．所以底面面积为

\[
S_{1}=B C \cdot C C_{1}=2 y .
\]

由于 \(A_{1}\) 到底面距离为 \(h_{1}=1\) ，且它的体积是平行六面体的一半，所以三棱柱的体积为

\[
V=\frac{1}{2} S_{1} \cdot h_{1}=y .
\]

（2）以 \(A B C\) 为底面。那么底面面积为

\[
S_{2}=S_{\triangle A B C}=\frac{x y}{2} .
\]

高为 z ，所以三棱柱的体积为

\[
V=S_{2} \cdot z=\frac{x y z}{2} .
\]

故 \(y=\frac{x y z}{2}\) ，即 \(x z=2\) ．又根据勾股定理，

\[
x^{2}+z^{2}=|A C|^{2}+\left|A_{1} C\right|^{2}=\left|A A_{1}\right|^{2}=4
\]

解得 \(x=z=2\) ．所以 \(A_{1} C=A C=\sqrt{2}\) ．\\
【方法四】由于三角形 \(A_{1} C_{1} C\) 是直角三角形，取 \(C C_{1}\) 中点 \(D\) ，由于直角三角形斜边中线等于斜边一半，所以 \(A_{1} D=\frac{1}{2} C C_{1}=1\) ，恰好就是 \(A_{1}\) 到底面 \(B B_{1} C_{1} C\) 的距离。由于垂线段最短，所以 \(A_{1} D \perp B C C_{1} B_{1}\) 。所以 \(A_{1} D \perp\) \(C C_{1}\) 。根据等腰三角形三线合－，\(A_{1} C=A_{1} C_{1}=A C\) 。\\
【注】以上所有的方法都可以用作【正向检查】，【代人检查】。\\
（II）平行线间的距离，就是点到另一条平行线的距离。我们可以沿用（I）中建系的方法。\\
【方法一】设 \(B C=y\) ，此时 \(B(0, y, 0)\) ．由于

\[
\begin{array}{r}
|A B|^{2}=|A C|^{2}+|B C|^{2} \\
\left|A_{1} B\right|^{2}=\left|A_{1} C\right|^{2}+|B C|^{2} .
\end{array}
\]

且 \(A C=A_{1} C\) ，所以 \(A B=A_{1} B\) ．根据等腰三角形三线合一，取 \(A A_{1}\) 中点 \(D\) ，必然有 \(B D \perp A A_{1}\) ．且 \(B D=2\) ．根据勾股定理，

\[
|A B|^{2}=|B D|^{2}+|A D|^{2}=4+1=5 .
\]

所以

\[
|B C|^{2}=|A B|^{2}-|A C|^{2}=5-2=3 .
\]

所以 \(|B C|=\sqrt{3}\) ．故 \(B(0, \sqrt{3}, 0)\) ．根据（I）中的计算，平面 \(B C C C_{1} B_{1}\) 的一个法向量为 \(\vec{n}=(\sqrt{2}, 0, \sqrt{2})\) ．由于 \(A(\sqrt{2}, 0,0), B_{1}(-\sqrt{2}, \sqrt{3}, \sqrt{2})\) 。所以

\[
\overrightarrow{A B_{1}}=(-2 \sqrt{2}, \sqrt{3}, \sqrt{2})
\]

那么 \(A B_{1}\) 与平面 \(B C C_{1} B_{1}\) 所称的正弦值为

\[
\frac{\left|\overrightarrow{A B_{1}} \cdot \vec{n}\right|}{|\vec{n}| \cdot\left|\overrightarrow{A B_{1}}\right|}=\frac{2}{2 \cdot \sqrt{13}}=\frac{\sqrt{13}}{13}
\]

我们可以【正向检查】，利用几何性质处理。由于已经知道 \(A_{1}\) 到平面 \(B C C_{1} B_{1}\) 的距离为 1 ，那么 \(A\) 到平面 \(B C C_{1} B_{1}\) 的距离也是 1 ．那么 \(A B_{1}\) 和平面 \(B C C_{1} B_{1}\) 的夹角正弦值为

\[
\sin \theta=\frac{1}{\left|A B_{1}\right|}
\]

根据勾股定理，

\[
\left|A B_{1}\right|^{2}=\left|B_{1} C_{1}\right|^{2}+\left|A C_{1}\right|^{2}=13
\]

所以 \(\left|A B_{1}\right|=\sqrt{13}\) ，故 \(\sin \theta=\frac{1}{\sqrt{13}}=\frac{\sqrt{13}}{13}\) ．\\
【注】这样看来，建系的方法是＂整体上看＂最方便的做法。

19．一项试验旨在研究臭氧效应。实验方案如下：选 40 只小白鼠，随机地将其中 20 只分配到实验组，另外 20 只分配到对照组，实验组的小白鼠饲养在高浓度臭氧环境，对照组的小白鼠饲养在正常环境，一段时间后统计每只小白鼠体重的增加量（单位：g）。\\
（1）设 \(X\) 表示指定的两只小白鼠中分配到对照组的只数，求 \(X\) 的分布列和数学期望；\\
（2）实验结果如下：\\
对照组的小白鼠体重的增加量从小到大排序为：\\
15.218 .820 .221 .322 .523 .225 .826 .527 .530 .1\\
32.634 .334 .835 .635 .635 .836 .237 .340 .543 .2

对照组的小白鼠体重的增加量从小到大排序为：\\
7.89 .211 .412 .413 .215 .516 .518 .018 .819 .2\\
19.820 .221 .622 .823 .623 .925 .128 .232 .336 .5\\
（i）求 40 只小鼠体重的增加量的中位数 \(m\) ，再分别统计两样本中小于 \(m\) 与不小于的数据的个数，完成如下列联表

\begin{center}
\begin{tabular}{|l|l|l|}
\hline
 & \(<m\) & \(\geq m\) \\
\hline
对照组 &  &  \\
\hline
实验组 &  &  \\
\hline
\end{tabular}
\end{center}

（ii）根据（i）中的列联表，能否有 \(95 \%\) 的把握认为小白鼠在高浓度臭氧环境中与正常环境中体重的增加量有差异．附：\(K^{2}=\frac{n(a d-b c)^{2}}{(a+b)(c+d)(a+c)(b+d)}\) ，

\begin{center}
\begin{tabular}{|c|l|l|c|}
\hline
\(k_{0}\) & 0.100 & 0.050 & 0.010 \\
\hline
\(P\left(k^{2} \geq k_{0}\right)\) & 2.706 & 3.841 & 6.635 \\
\hline
\end{tabular}
\end{center}

（I）由于两只小白鼠不同，且小白鼠数量不大，\(X\) 应服从超几何分布。这样思考可以，但是非常容易乱，因为我们很多时候对＂放回＂和＂不放回＂的理解是不准确的。事实上，我们可以更朴素地来思考这个问题。\\
【方法一】假设指定的两只小白鼠分别为 1 号， 2 号。设事件 \(A\) 为 1 号小白鼠在试验组，事件 \(B\) 为 2 号小白鼠在试验组。那么：\\
（1） \(\mathbb{P}(X=2)=\mathbb{P}(A \cap B)\) ．由于 \(A, B\) 并不独立（这是重点），所以我们可以一步一步想这个问题。显然， \(\mathbb{P}(A)=\)\\
\(\frac{20}{40}=\frac{1}{2}\) ，当 1 号已经在试验组时，试验组只有 19 个位置了。此时，还剩 39 只小鼠，那么 2 号小鼠在试验组的概率只能是 \(\mathbb{P}(B \mid A)=\frac{19}{39}\) 。注意，这里是条件概率：因为我们求的是 1 号小鼠已经在试验组时， 2 号小鼠在试验组的概率。因此：

\[
\mathbb{P}(A \cap B)=\mathbb{P}(A) \cdot \mathbb{P}(B \mid A)=\frac{1}{2} \cdot \frac{19}{39}=\frac{19}{78}
\]

（2）\(\backslash \mathbb{P}(X=1)=\mathbb{P}\left(A^{c} \cap B\right)+\mathbb{P}\left(A \cap B^{c}\right)\) ．显然， \(\mathbb{P}\left(A^{c} \cap B\right)=\mathbb{P}\left(A \cap B^{c}\right)\) ，求一个即可。还是和 同样的思路，当 1 号小鼠在试验组时，对照组有 20 个位置，一共还有 39 只小鼠，那么 2 号小鼠不在试验组（在对照组）的概率就是

\[
\mathbb{P}\left(B^{c} \mid A\right)=\frac{20}{39} .
\]

所以 \(\mathbb{P}\left(A \cap B^{c}\right)=\mathbb{P}(A) \cdot \mathbb{P}\left(B^{c} \mid A\right)=\frac{1}{2} \cdot \frac{20}{39}=\frac{10}{39}\) ．故 \(\mathbb{P}(X=1)=2 \mathbb{P}\left(A \cap B^{c}\right)=\frac{20}{39}\) ．\\
（3） \(\mathbb{P}(X=0)=\mathbb{P}\left(A^{c} \cap B^{c}\right)\) 。当 1 号小鼠不在试验组时，对照组有 19 个位置共还有 39 只小鼠，那么 2 号小鼠不在试验组的概率就是

\[
\mathbb{P}\left(B^{c} \mid A^{c}\right)=\frac{19}{39} .
\]

所以 \(\mathbb{P}(X=0)=\mathbb{P}\left(A^{c} \cap B^{c}\right)=\mathbb{P}\left(A^{c}\right) \cdot \mathbb{P}\left(B^{c} \mid A^{c}\right)=\frac{1}{2} \cdot \frac{19}{39}=\frac{19}{78}\) ．\\
【方法二】由于白鼠随机分配，那么我们考察指定的白鼠，就相当于将白鼠先分

\[
\begin{gathered}
\mathbb{P}(X=2)=\frac{\mathrm{C}_{20}^{2}}{\mathrm{C}_{40}^{2}}=\frac{20 \cdot 19}{40 \cdot 39}=\frac{19}{78}, \\
\mathbb{P}(X=1)=\frac{\mathrm{C}_{20}^{1} \mathrm{C}_{20}^{1}}{\mathrm{C}_{40}^{2}}=\frac{20 \cdot 20}{40 \cdot 39}=\frac{20}{39}, \\
\mathbb{P}(X=0)=\frac{\mathrm{C}_{20}^{2}}{\mathrm{C}_{40}^{2}}=\frac{20 \cdot 10}{40 \cdot 39}=\frac{19}{78} .
\end{gathered}
\]

也能得到同样的答案。\\
【方法三】古典概型。我们直接去考虑所有可能的分组方法，找到符合条件的并计算概率。总方法数为 \(C_{4} 0^{20}\) ．\\
（1）两只小鼠都在试验组，相当于从剩下 38 只中选 18 只进入试验组，其余进人对照组。所以

\[
\mathbb{P}(X=2)=\frac{\mathrm{C}_{38}^{18}}{\mathrm{C}_{40}^{20}}=\frac{38 \cdot 37 \cdot \ldots \cdot 21}{18!} \cdot \frac{20!}{40 \cdot 39 \cdot \ldots \cdot 21}=\frac{20 \cdot 19}{40 \cdot 39}=\frac{19}{78} .
\]

（2）一只小鼠在试验组，另一只再对照组，相当于从剩下 38 只中选 19 只进入试验组，其余进入对照组。所以

\[
\mathbb{P}(X=1)=\frac{\mathrm{C}_{38}^{19}}{\mathrm{C}_{40}^{20}}=\frac{38 \cdot 37 \cdot \ldots \cdot 20}{19!} \cdot \frac{20!}{40 \cdot 39 \cdot \ldots \cdot 21}=\frac{20 \cdot 20}{40 \cdot 39}=\frac{20}{39} .
\]

（3）两只都在对照组，和两只都在试验组完全一样，因为两组只数相同。所以

\[
\mathbb{P}(X=0)=\mathbb{P}(X=2)=\frac{19}{78} .
\]

因此，\(X\) 的分布列为

\begin{center}
\begin{tabular}{|c|c|c|c|}
\hline
\(x\) & 0 & 1 & 2 \\
\hline
\(\mathbb{P}\) & \(\frac{19}{78}\) & \(\frac{20}{39}\) & \(\frac{19}{78}\) \\
\hline
\end{tabular}
\end{center}

在这里，我们可以使用【结构检查法】。由于全空间概率为 1 ，所以可以验证分布列所有概率之和是否为 1 ：

\[
\frac{19}{78}+\frac{20}{39}+\frac{19}{78}=\frac{38}{78}+\frac{40}{78}=1 .
\]

成立。数学期望为：

\[
\mathbb{E} X=\frac{19}{78} \cdot 2+\frac{20}{39} \cdot 1+\frac{19}{78} \cdot 0=1 .
\]

这里我们也可以使用【结构检查法】。由于每只小鼠各有一半的概率进入试验组，所以平均应该有 \(2 \cdot 0.5=1\) 只小鼠在试验组，这和 \(\mathrm{E} X=1\) 相符。这也和超几何分布的期望相符合：在 \(M=40\) 个样品中，有 \(N=20\) 个次品，从中抽取 \(n=2\) 个，那么次品个数的数学期望为

\[
n \cdot \frac{N}{M}=2 \cdot \frac{1}{2}=1 .
\]

（II）（i）需要对所有数据进行排序。显然第二组的数据较小，如果从小到大排序，大部分数据都会从第二组出现。这 40 个数据从小到达排序为：7．89．211．412．413．215．215．516．518．018．818．819．219．820．220．221．321．622．522．823．223．623．925．125．82从小到大第 20 个数据是 23.2 ，第 21 个数据是 23.6 ，所以中位数 \(m=23.4\) ．检查的时候，可以从大到小再检查一下。显然，列联表为：

\begin{center}
\begin{tabular}{|c|c|c|}
\hline
 & \(<m\) & \(\geq m\) \\
\hline
对照组 & 6 & 14 \\
\hline
试验组 & 14 & 6 \\
\hline
\end{tabular}
\end{center}

根据 \(K^{2}\) 的计算公式，

\[
K^{2}=40 \cdot \frac{\left(6^{2}-14^{2}\right)^{2}}{20 \cdot 20 \cdot 20 \cdot 20}=2 \cdot \frac{160^{2}}{20 \cdot 20 \cdot 20}=\frac{2 \cdot 8 \cdot 8}{20}=6.4>3.841 .
\]

因此，有 \(95 \%\) 的把握认为小白鼠在高浓度臭氧环境中与在正常环境中体重的增加量有差异。\\
20．已知直线 \(x-2 y+1=0\) 与抛物线 \(C: y^{2}=2 p x(p>0)\) 交于 \(A, B\) 两点，且 \(|A B|=4 \sqrt{15}\) ．\\
（1）求 \(P\) ；\\
（2）设 \(F\) 为 \(C\) 的焦点，\(M, N\) 为 \(C\) 上两点， \(\overrightarrow{F M} \cdot \overrightarrow{F N}=0\) ，求 \(\triangle M F N\) 面积的最小值．\\
（I）\(p=2 ;\)（2） \(12-8 \sqrt{2}\) ．\\
我们在章节＜代数结构的实战应用 \(>8\) 中已经讲过这道题。在这里，我们只展现其中一种方法，重点放在如何检查。\\
（1）将直线方程和扡物线方程联立，得

\[
y^{2}=2 p(2 y-1) y^{2}-4 p y+2 p=0 .
\]

根据韦达定理，

\[
\left\{\begin{array}{c}
y_{1}+y_{2}=4 p \\
y_{1} y_{2}=2 p
\end{array}\right.
\]

根据两点间距离公式，有：

\[
\begin{aligned}
|A B|=\sqrt{1+2^{2}} \cdot\left|y_{1}-y_{2}\right| & =\sqrt{5} \cdot \sqrt{\left(y_{1}+y_{2}\right)^{2}-4 y_{1} y_{2}}=\sqrt{5} \cdot \sqrt{16 p^{2}-8 p} \\
& =2 \sqrt{5} \cdot \sqrt{4 p^{2}-2 p}=4 \sqrt{15} .
\end{aligned}
\]

故

\[
\begin{aligned}
& \sqrt{4 p^{2}-2 p}=2 \sqrt{3} \\
& 4 p^{2}-2 p-12=0 \\
& 2 p^{2}-p-6=0 \\
& (p-2)(2 p+3)=0 .
\end{aligned}
\]

由于 \(p>0\) ，所以 \(p=2\) ．\\
得到这个结果后，可以【代入检查 \(\mathbf{1}\) 。此时，拖物线为 \(y^{2}=4 x\) ．那么拖物线和直线两个交点的纵坐标满足：

\[
\begin{array}{r}
\frac{y^{2}}{4}-2 y+1=0 \\
y^{2}-8 y+4=0 \\
(y-4)^{2}=12 .
\end{array}
\]

所以 \(y_{1}=4+2 \sqrt{3}, y_{2}=4-2 \sqrt{3}\) ．那么

\[
\begin{aligned}
& x_{1}=2 y_{1}+1=9+4 \sqrt{3}, \\
& x_{2}=2 y_{2}+1=9-4 \sqrt{3} .
\end{aligned}
\]

所以

\[
|A B|=\sqrt{\left(x_{1}-x_{2}\right)^{2}+\left(y_{1}-y_{2}\right)^{2}}=\sqrt{(8 \sqrt{3})^{2}+(4 \sqrt{3})^{2}}=4 \sqrt{3} \cdot \sqrt{5}=4 \sqrt{15} .
\]

在这里，也有计算的小技巧需要注意：很多同学直接去计算 \(4 \sqrt{3}, 8 \sqrt{3}\) 的平方，但是可以利用提公因数巧妙规避计算：

\[
\sqrt{(8 \sqrt{3})^{2}+(4 \sqrt{3})^{2}}=\sqrt{(4 \sqrt{3})^{2} \cdot\left(2^{2}+1\right)}=4 \sqrt{3} \cdot \sqrt{5}=4 \sqrt{15} .
\]

经检验，满足题目条件。\\
这里的验算是非常有必要的。大题的第一问和后面有着很强的关联，如果第一间出错，后面也不可能正确。所以，一定要在做完第一问后马上停下来检查。\\
（II）由于 \(p=2\) ，所以焦点坐标为 \(F(1,0)\) ．设直线 \(M N\) 的方程为 \(x=m y+t\) ．和抛物线方程联立：

\[
\begin{aligned}
& y^{2}=4(m y+t) \\
& y^{2}-4 m y-4 t=0 .
\end{aligned}
\]

根据韦达定理，

\[
\left\{\begin{array}{c}
y_{1}+y_{2}=4 m \\
y_{1} y_{2}=-4 t
\end{array}\right.
\]

那么

\[
\overrightarrow{F M} \cdot \overrightarrow{F N}=y_{1}^{2} y_{2}^{2}-4\left(y_{1}^{2}+y_{2}^{2}\right)+16+16 y_{1} y_{2}=0 .
\]

根据韦达定理，

\[
y_{1}^{2}+y_{2}^{2}=\left(y_{1}+y_{2}\right)^{2}-2 y_{1} y_{2}=4 m^{2}+8 t .
\]

故

\[
\begin{aligned}
& 16 t^{2}-4\left(4 m^{2}+8 t\right)+16-64 t=0 \\
& 16 t^{2}-16 m^{2}-96 t+16=0 \\
& m^{2}-t^{2}-6 t+1
\end{aligned}
\]

根据方法一中的运算，

\[
S_{\triangle F M N}=\frac{1}{16}\left(y_{1} y_{2}+4\right)^{2}=\frac{1}{16}(-4 t+4)^{2}=(t-1)^{2} .
\]

故求 \(t\) 的取值范围即可。由于

\[
t^{2}-6 t+1=m^{2} \geq 0
\]

所以

\[
\begin{aligned}
& t^{2}-6 t+1 \geq 0 \\
& (t-3)^{2} \geq 8
\end{aligned}
\]

故 \(t \geq 3+2 \sqrt{2}\) 或 \(t \leq 3-2 \sqrt{2}\) ．\\
（i）若 \(t \geq 3+2 \sqrt{2}\) ，则 \(t-1 \geq 2+2 \sqrt{2}>0\) ，则

\[
S_{\triangle F M N}=(t-1)^{2} \geq(2+2 \sqrt{2})^{2}=12+8 \sqrt{2} \text {. }
\]

（ii）若 \(t \leq 3-2 \sqrt{2}\) ，则 \(t-1 \leq 2-2 \sqrt{2}<0\) ，则

\[
S_{\triangle F M N}=(t-1)^{2} \geq(2-2 \sqrt{2})^{2}=12-8 \sqrt{2} .
\]

综合（i），（ii），\(S_{\triangle F M N}\) 面积的最小值为 \(12-8 \sqrt{2}\) ．当 \(t=3-2 \sqrt{2}, m=0\) 时取等。依旧可以【代人检查】。当 \(t=3-2 \sqrt{2}, m=0\) 时，直线方程为 \(x=3-2 \sqrt{2}\) ．由于抛物线方程为 \(y^{2}=4 x\) ，所以它和扡物线的两个交点的纵坐标满足：

\[
y^{2}=4(3-2 \sqrt{2}) .
\]

注意到 \(3-2 \sqrt{2}=(\sqrt{2}-1)^{2}\) ，所以

\[
\begin{gathered}
y^{2}=2^{2}(\sqrt{2}-1)^{2} \\
y= \pm(2 \sqrt{2}-2) .
\end{gathered}
\]

所以 \(M(3-2 \sqrt{2}, 2 \sqrt{2}-2), N(3-2 \sqrt{2}, 2-2 \sqrt{2})\) ．我们验证 \(F M, F N\) 是否垂直：事实上，

\[
\begin{aligned}
& k_{F M}=\frac{(2 \sqrt{2}-2)-0}{3-2 \sqrt{2}-1}=\frac{2 \sqrt{2}-2}{2-2 \sqrt{2}}=-1, \\
& k_{F N}=\frac{(2-2 \sqrt{2})-0}{3-2 \sqrt{2}-1}=\frac{2-2 \sqrt{2}}{2-2 \sqrt{2}}=1 .
\end{aligned}
\]

所以 \(k_{F M}, k_{F N}\) 和 \(x\) 轴的夹角都是 \(45^{\circ}\) ，显然垂直，此时刚好是一个等腰直角三角形。

\[
|F M|=|F N|=x+1=4-2 \sqrt{2} \text {. }
\]

根据扡物线的定义，故面积为故面积为

\[
S=\frac{1}{2}|F M| \cdot|F N|=\frac{1}{2} \cdot(4-2 \sqrt{2})^{2}=\frac{24-16 \sqrt{2}}{2}=12-8 \sqrt{2} .
\]

符合题意。\\
【注】圆锥曲线大题，尽量都检查取等条件。虽然计算量可能稍显复杂，但是可以显著提高正确率。

21．已知函数 \(f(x)=a x-\frac{\sin x}{\cos ^{3} x}, x \in\left(0, \frac{\pi}{2}\right)\)\\
（1）当 \(a=8\) 时，讨论 \(f(x)\) 的单调性；\\
（2）若 \(f(x)<\sin 2 x\) 恒成立，求 \(a\) 的取值范围．\\
（I）讨论单调性，只需要直接求导即可。当 \(a=8\) 时，

\[
f(x)=8 x-\frac{\sin x}{\cos ^{3} x}
\]

注意到后面的三角函数是一个分式，求导可能非常复杂，分母会出现 \(\cos ^{6} x\) ，所以可以考虑在章节＜代数结构的实战应用＞中讲过的用乘法法则计算：

【方法一】直接除法法则。

\[
\begin{aligned}
& \left(\frac{\sin x}{\cos ^{3} x}\right)^{\prime}=\frac{\cos x \cdot \cos ^{3} x-\sin x \cdot\left(\cos ^{3} x\right)^{\prime}}{\cos ^{6} x}=\frac{\cos ^{4} x-\sin x\left(-\sin x \cdot 3 \cos ^{2} x\right)}{\cos ^{6} x} \\
& =\frac{\cos ^{4} x+3 \sin ^{2} x \cos ^{2} x}{\cos ^{6} x}=\frac{\cos ^{2} x+3 \sin ^{2} x}{\cos ^{4} x} .
\end{aligned}
\]

【方法二】利用乘法法则。

\[
\begin{aligned}
\left(\frac{\sin x}{\cos ^{3} x}\right)^{\prime} & =(\sin x)^{\prime} \cdot \frac{1}{\cos ^{3} x}+\sin x \cdot\left(\frac{1}{\cos ^{3} x}\right)^{\prime} \\
& =\frac{\cos x}{\cos ^{3} x}+\sin x \cdot(-\sin x) \cdot\left(-\frac{3}{\cos ^{4} x}\right) \\
& =\frac{1}{\cos ^{2} x}+\frac{3 \sin ^{2} x}{\cos ^{4} x}
\end{aligned}
\]

所以

\[
f^{\prime}(x)=8-\left(\frac{\sin x}{\cos ^{3} x}\right)^{\prime}=8-\frac{1}{\cos ^{2} x}-\frac{3 \sin ^{2} x}{\cos ^{4} x}
\]

注意到这里又有 \(\sin ^{2} x\) 又有 \(\cos ^{2} x\) ，所以可以利用勾股定理消元。将 \(\sin ^{2} x=1-\cos ^{2} x\) 代入，得

\[
f^{\prime}(x)=8-\frac{1}{\cos ^{2} x}-\frac{3-3 \cos ^{2} x}{\cos ^{4} x}=8+\frac{2}{\cos ^{2} x}-\frac{3}{\cos ^{4} x}
\]

在这里，我们既然可以把它看作关于 \(t=\frac{1}{\cos ^{2} x}\) 的二次函数，也可以直接通分。通分的做法相对技巧性更低，更直观：

\[
f^{\prime}(x)=\frac{8 \cos ^{4} x+2 \cos ^{2} x-3}{\cos ^{4} x}
\]

令 \(m=\cos ^{2} x\) ，那么十字相乘得：

\[
f^{\prime}(x)=\frac{8 m^{2}+2 m-3}{m^{4}}=\frac{(2 m-1)(4 m+3)}{m^{2}}
\]

因此，令 \(f^{\prime}(x)=0\) ．有 \(m=\frac{1}{2}\) ．\\
（1）当 \(m<\frac{1}{2}\) 时，即 \(\cos x<\frac{\sqrt{2}}{2}\) ，即 \(x \in\left(\frac{\pi}{4}, \frac{\pi}{2}\right)\) 时，\(f(x)\) 单调递减；\\
（2）当 \(m>\frac{1}{2}\) 时，即 \(\cos x>\frac{\sqrt{2}}{2}\) ，即 \(x \in\left(0, \frac{\pi}{4}\right)\) 时，\(f(x)\) 单调递增。\\
所以，\(f(x)\) 在 \(\left(0, \frac{\pi}{4}\right)\) 上单调递增，在 \(\left(\frac{\pi}{4}, \frac{\pi}{2}\right)\) 上单调递减．\\
可以使用【代入检查法】。根据计算，当 \(x=\frac{\pi}{4}\) 时，应有 \(f^{\prime}(x)=0\) 。我们可以代入到充分化简前的 \(f^{\prime}(x)\) 的表达式中检验：

\[
f^{\prime}\left(\frac{\pi}{4}\right)=8-\frac{1}{\cos ^{2} \frac{\pi}{4}}-\frac{3 \sin ^{2} \frac{\pi}{4}}{\cos ^{4} \frac{\pi}{4}}=8-\frac{1}{\left(\frac{1}{2}\right)}-\frac{\frac{3}{2}}{\left(\frac{1}{4}\right)}=8-2-6=0 .
\]

\section*{符合题意。}
（II）令 \(g(x)=f(x)-\sin 2 x\) ．那么 \(g(x)<0\) 对 \(x \in\left(0, \frac{\pi}{2}\right)\) 成立。此时

\[
\begin{aligned}
g^{\prime}(x) & =f^{\prime}(x)-(\sin 2 x)^{\prime}=a+\left(\frac{2}{\cos ^{2} x}-\frac{3}{\cos ^{4} x}\right)-2 \cos 2 x \\
& =a+\frac{2}{\cos ^{2} x}-\frac{3}{\cos ^{4} x}-2\left(2 \cos ^{2} x-1\right) \\
& =\frac{2}{\cos ^{2} x}-4 \cos ^{2} x-\frac{3}{\cos ^{4} x}+a+2 .
\end{aligned}
\]

由于 \(g(0)=f(0)=0\) ，所以其实 \(g(x)<0\) 等价于 \(g(x)<g(0)\) ．这意味着：\(g(x)\) 在 \(x=0\) 处的导数小于等于 0 ．部分教辅称之为＂端点效应＂，但是本质上就是一个单调性的判断。我们先承认这一点，在后面通过分类讨论来深人理解它。由于 \(g^{\prime}(0) \leq 0\) ，所以

\[
g^{\prime}(0)=2-4-3+a+2=a-3 \leq 0 .
\]

所以 \(a \leq 3\) ．\\
当 \(a \leq 3\) 时，我们尝试着证明 \(g^{\prime}(x) \leq 0\) ．我们之所以得到这个判断，是因为看到了 \(\cos ^{2} x\) 的次数结构。注意到

\[
\frac{2}{\cos ^{2} x}-4 \cos ^{2} x-\frac{3}{\cos ^{4} x}
\]

的存在，决定了它＂有上界＂。这是因为：\\
（1）当 \(x\) 接近 0 时， \(\cos ^{2} x\) 会接近 1 ，那么这个数是个负数，接近 -5 ；\\
（2）当 \(x\) 接近 \(\frac{\pi}{2}\) 时， \(\cos ^{2} x\) 会接近 0 ，但是由于 \(\frac{1}{\cos ^{4} x}\) 远比 \(\frac{1}{\cos ^{2} x}\) 要大，所以这个数也会是个负数；\\
（3）当 \(x\) 取其他角度时， \(\cos ^{2} x\) 也会是一个小于 1 的数，和（2）相同，\(\frac{1}{\cos ^{4} x}\) 都会比 \(\frac{1}{\cos ^{2} x}\) 大，所以它不会特别大。这些＂微判断＂共同帮助我们做出了一个大胆的猜测，就是：当 \(a \leq 3\) 时，\(g^{\prime}(x) \leq 0\) 。我们接下来尝试证明这一点。此时，

\[
g^{\prime}(x)=\frac{2}{\cos ^{2} x}-4 \cos ^{2} x-\frac{3}{\cos ^{4} x}+a+2 \leq \frac{2}{\cos ^{2} x}-4 \cos ^{2} x-\frac{3}{\cos ^{4} x}+5 .
\]

故只需要证明

\[
\frac{2}{\cos ^{2} x}-4 \cos ^{2} x-\frac{3}{\cos ^{4} x}+5 \leq 0 .
\]

令 \(\cos ^{2} x=t\) ，通分得

\[
\frac{5 t^{2}-4 t^{3}+2 t-3}{t^{2}} \leq 0
\]

令 \(h(t)=-4 t^{3}+5 t^{2}+2 t-3\) ，这是一个三次函数，且 \(h(1)=0\) 。根据杜杆 1 ，杜杆 6 中的结论，\(h(t)\) 必然有 \((t-1)\)这个因式。利用待定系数法，口算得：

\[
-4 t^{3}+5 t^{2}+2 t-3=-\left(4 t^{3}-5 t^{2}-2 t+3\right)=-(t-1)\left(4 t^{2}-t-3\right) .
\]

还可以继续分解因式：

\[
-(t-1)\left(4 t^{2}-t-3\right)=-(t-1)(t-1)(4 t+3)=-(t-1)^{2}(4 t+3) .
\]

由于 \(t=\cos ^{2} x \geq 0\) ，所以 \(h(t) \leq 0\) ．证毕！因此，当 \(a \leq 3\) 时，\(g^{\prime}(x) \leq 0\) 对 \(x \in\left(0, \frac{\pi}{2}\right)\) 恒成立。最后，我们来解释端点效应，或者说：为什么一定要求 \(g^{\prime}(0) \leq 0\) ，才有可能 \(g(x)<0\) 恒成立？\\
（1）当 \(a>3\) 时，\(g^{\prime}(0)>0\) ．那么，\(g(x)\) 必然在 0 ＂附近＂单调递增。因此，存在 \(\delta>0\) ，是的 \(g^{\prime}(x)>0\) 对 \(x \in(0, \delta)\)成立。那么 \(g(x)>g(0)=0\) ，与 \(g(x)<0\) 矛盾。不成立；\\
（2）当 \(a \leq 3\) 时，根据前面的计算，\(g^{\prime}(x) \leq 0\) ．所以 \(g(x)\) 在 \(\left(0, \frac{\pi}{2}\right)\) 上单调递减，所以

\[
g(x)<g(0)=0 .
\]

成立。\\
综合（1），（2），有 \(a \leq 3\) 。\\
依然可以使用【代入检查法】。为了验证 \(a \leq 3\) ，我们可以代入一些特殊的情况检查。\\
（1）当 \(a=0\) 时，\(f(x)=-\frac{\sin }{x} \cos ^{3} x\) ．此时，当 \(x \in\left(0, \frac{\pi}{2}\right)\) 时，\(f(x)<0\) ．而 \(\sin 2 x>0\) ，成立！\\
（2）当 \(a=5\) 时，\(f(x)=5 x-\frac{\sin }{x} \cos ^{3} x . \stackrel{\Delta}{\Leftrightarrow} x=\frac{\pi}{4}\) ，那么

\[
f\left(\frac{\pi}{4}\right)=5 \cdot \frac{\pi}{4}-2=\frac{5 \pi-8}{4}>\frac{12-8}{4}=1 .
\]

此时 \(\sin 2 x=\sin \frac{\pi}{2}=1\) ．所以 \(f(x)>\sin 2 x\) ，不符合题意。\\
这两个特殊情况都和 \(a \leq 3\) 这个范围没有矛盾。我们还可以进一步验证：\\
【注】如果想要把 中的附近说清楚，需要找点的技巧，较为繁琐。事实上，当 \(a>3\) 时，令 \(t=\cos ^{2} x\) ，那么

\[
g(x)=\frac{-4 t^{3}+(a+2) t^{2}+2 t-3}{t^{2}}
\]

考察分子 \(p(t)=-4 t^{3}+(a+2) t^{2}+2 t-3\) ．注意到 \(p(1)=a-3>0, p(0)=-3<0\) ，所以根据零点存在定理，存在 \(t_{0} \in(0,1)\) 使得 \(p\left(t_{0}\right)=0\) ．记最靠近 1 的零点为 \(t_{0}\) ，那么当 \(t \in\left(t_{0}, 1\right)\) 时，\(p(t)>0\) ，即 \(g(x)>0\) ．设 \(x_{0}\) 满足 \(\cos x_{0}=\sqrt{t_{0}}\) ．那么，当 \(x \in\left(0, x_{0}\right)\) 时，\(t=\cos ^{2} x \in\left(t_{0}, 1\right)\) ，此时 \(g(x)>0\) 。所以 \(g(x)\) 在 \(\left(0, x_{0}\right)\) 上单调递增，所以 \(g\left(x_{0}\right)>0\) ，不成立。

【注】还可以从导数的定义理解。由于我们要求了 \(g(x)<0\) ，所以

\[
g^{\prime}(0)=\lim _{x \rightarrow 0} \frac{g(x)-g(0)}{x} .
\]

由于 \(g(0)=0, g(x)<0\) ，所以 \(g(x)-g(0)<0\) 。而我们只关心 \(x \in\left(0, \frac{\pi}{2}\right)\) 的情况，所以 \(x-0>0\) 。所以

\[
\frac{g(x)-g(0)}{x}<0 .
\]

对一个负数取极限，结果只能是非负的。所以 \(g^{\prime}(0) \leq 0\) ．由于高中学习过导数的定义，所以这个结果也可以接受。

\subsection*{7.2 2023 年新高考 I 卷}
\section*{XXX}
\section*{开场白}
XXX

\section*{7．2．1 稳定全对不失误 01－06}
XXX

\section*{7．2．2 基本功与认识自己 07－08}
\section*{XXX}
7．2．3 分析法思考问题 09－11 XXX

7．2．4 考场利益最大化 11－12\\
XXX

7．2．5 画图法解压轴 13－16\\
XXX

7．2．6 考场稳定全对 17－19 XXX

7．2．7 预判与大小无关 20 数列\\
XXX

7．2．8 数列要素关系的核心 21 概率\\
XXX

7．2．9 讲烪了的冻结法 22 解几\\
XXX

总结\\
XXX

\subsection*{7.3 2024年新高考 I 卷}
XXX

7．3．1 无任何失误的可能 01－08\\
XXX

7．3．2 强调烂了的画图 09－11\\
XXX

\section*{7．3．3 树状图魅力时刻 12－14}
XXX

7．3．4 小小元问题，彰显基本功 15－16\\
XXX

\section*{7．3．5 预判配合即时检查 17 立几}
XXX

\section*{7．3．6 普通的找零点 18 导数}
\section*{XXX}
\section*{7．3．7 思想实验大放送 19 数列}
思想试验：没有感觉时先做做简单情形或者静态情况的尝试，寻找突破

设 \(m\) 为正整数，数列 \(a_{1}, a_{2}, \cdots, a_{4 m+2}\) 是公差不为 0 的等差数列，若从中删去两项 \(a_{i}\) 和 \(a_{j}(i<j)\) 后制余的 \(4 m\) 项可被平均分为 \(m\) 组，且每组的 4 个数都能构成等差数列，则称数列 \(a_{1}, a_{2}, \cdots, a_{4 m+2}\) 是 \((i, j)\) 一可分数列。\\
（1）写出所有的 \((i, j), 1 \leqslant i<j \leqslant 6\) ，使得数列 \(a_{1}, a_{2}, \cdots, a_{6}\) 足 \((i, j)\) 一可分数列；\\
（2）当 \(m \geqslant 3\) 时，证明：数列 \(a_{1}, a_{2}, \cdots, a_{4 m+2}\) 是 \((2,13)\)－可分数列；\\
（3）从 \(1,2, \cdots, 4 m+2\) 中一次任取两个数 \(i\) 和 \(j(i<j)\) ，记数列 \(a_{1}, a_{2}, \cdots, a_{4 m+2}\) 是 \((i, j)\) 一可分数列的概率为 \(P_{m}\) ，证明：\(P_{m}>\frac{1}{8}\) ．

因为大小无关，所以能只看下标，比如 \(a_{1}, ~ a_{3}, ~ a_{5}, ~ a_{7}\) 。

对答案复盘\\
XXX

总结

\section*{XXX}
\subsection*{7.42023 年四省联考}
2023年3月11日视频

开场白：稳定性与排名才是王道\\
XXX

7．4．1 单选的临场发挥 01－07 Xxx

7．4．2 多选的临场发挥 09－11 XXX

\section*{7．4．3 填空的临场发挥 13－15}
 XXX7．4．4＂恬不知耻＂的骗分 08 XXX

7．4．5 快速寻找可能出路 12 XXX

\section*{7．4．6 耐心数数 16}
 XXX7．4．7 敏捷思维与手速 17－19 Xxx

7．4．8 思维带来破题人口 20－21 Xxx

7．4．9 抢分圣战 22 椭圆 XxX

\section*{7．5 阳哥高质量原创题}
2023年1月6日标题只写了新高考题

7．5．1 简单题稳扎稳打 01－03＋13 xxx

7．5．2 有趣的情景题 04\\
XXX

7．5．3 基本的比大小问题 05 XXX

\section*{7．5．4 反思与批判 14}
XXX

7．5．5 线线角也需要思维 06 XXX

7．5．6 冻结思想的深刻性 07 xxx

7．5．7 简单又新颖的小题 15\\
7．5．8 充分必要也能考出花 08\\
xxx

7．5．9 不依赖二级结论一样做 16 Xxx

7．5．10 动态眼光＋充要替换 15 XxX

7．5．11 重视要素关系网 12 XXX\\
（旧高考独有）新定义数列问题 12 xxx

7．5．12 基本的描述性数列 17 XXX

7．5．13 第一问的新考法 18 Xxx

7．5．14 分析好自由度 19\\
XxX

7．5．15 纸老虎的假自由度 20\\
XXX

7．5．16 源于游戏的概率难题 21 XxX

7．5．17 短小精悍的导数大题 22\\
XXX

旧高考选考题\\
XXX

总结\\
XXX

\subsection*{7.6 2022年新高考 I 卷}
开场白\\
本节讲解主要目的是让平时（较稳定）考 130 分以上的同学在这套试卷上接近 130 。

7．6．1 速度与正确率 01－03 XXX

7．6．2 知识障碍也要抢分 04－05 Xxx

7．6．3 稳分析与抢时间 06－08 Xxx

7．6．4 动静视角快速收割 09－11 XXX

7．6．5 选择题总结 +12 多选 XXX

7．6．6 熟悉的问题迅速转化 13－15 XXX

7．6．7 选填回顾总结＋ 16 填空 XXX

7．6．8 迅速抓主要矛盾并解决 17 数列 XXX

7．6．9 思维范式胜过一切 18 三角 XXX

7．6．10 思维打击变成送分 19 立几 XXX

7．6．11 基础充要替换随手拿下 20 概率 XXX

7．6．12 辅助思想大放异彩 21 圆曲 XXX

\section*{7．6．13 生死时速抢分大战 22 导数}
 XXX总结 XXX

\section*{第8章 大学数学科普系列}
本章内容为 up 主 ninchijou 的涂山苏苏＂北大数院up的大学数学科普＂系列视频对应文稿的整合。\\
【第一期一数学的基础：逻辑】https：／／b23．tv／umg6c0I\\
【第二期—重要的概念：集合】https：／／b23．tv／1QuOceg\\
【第三期—联系万物：映射】https：／／b23．tv／wGoI8Ee\\
【第四期—优美的仿射几何（上）】https：／／b23．tv／IwTFTdQ\\
【第五期—优美的仿射几何（下）】https：／／b23．tv／NWDkAz9\\
【第六期—数学规则篇（上）：＂数学存在＂与实数理论】https：／／b23．tv／GCnwgoy

\section*{8.1 数学的基础}
事前声明：本次内容主要为大学数学的数理逻辑部分，为了方便高中生理解，内容中的部分定义和表述进行了直观化的简化处理，因此可能部分地方不够严谨，感兴趣的同学可以自行查阅大学的数理逻辑内容

\begin{displayquote}
我们从小学到现在, 一直在学数学, 大家心目中的数学是什么呢?\\
是枯燥无味的计算? 难度巨大的题目? 还是别的?
\end{displayquote}

常见观点1：学数学的话，就要整天做好多繁杂的计算，做好多需要高智商的难题\\
常见观点 2：数学好的人，一定非常聪明，反应速度很快\\
常见观点 3：数学好的人，计算能力一定很强\\
常见观点 4：数学家都是相当聪明的人，都是通过各种奇思妙想的技巧来解决世界难题\\
实际上，数学这门学科真正在干的事情，和这些都有一定的区别

\section*{8．1．1 数学的逻辑}
大家都知道，数学是讲逻辑的，每一步推理都要求是有根有据，那大家能简单阐释一下，什么叫＂有根有据＂呢？数学只承认一种推理，即

\section*{定义8．1（演绎推理）。）}
从一系列基本前提（也就是公理）出发，通过某些确定的推理规则，推出结论的过程叫做演绎推理。

这一点相信大家在初中多少有所接触，所谓公理，就是推理的起点，是不证自明的事实\\
数学上承认的推理，不能依靠直觉，只能依靠严格的演绎比如，那户人家，父母都是小偷，则孩子也是小偷又比如，那个人面色䵢黑，手上茧子很多，如果从农民和白领里面选，我们一定会选农民

这些推理都不能称之为数学推理，最多只能称之为归纳推理而且，这些推理还有可能出现错误那么，我想问大家两个问题：

1．公理从何而来？\\
2．数学中承认的推理规则是什么？\\
我可以先回答一下大家：\\
1．公理是人为规定的，尽管是根据经验而给出的，但真正的数学公理都是人为规定的！也就是说：公理是可以修改的

2．数学中承认的最基本的推理规则称为分离规则，也就是

\[
\text { 若 } A, A \rightarrow B \text { 则 } B
\]

其他的推理规则都可以被严格的演绎推理而证明\\
所谓定理，就是由公理出发经过演绎推理得到的命题

\section*{定理 8.1}
对于命题 \(A\) ，若存在一个命题序列 \(A_{1}, A_{2} \cdots A_{n}\) ，使得对任意 \(\mathrm{i}=1,2 \cdots n\) ，要么 \(\mathrm{A}_{1}\) 是公理集中的元素，要么它是 \(\mathrm{A}_{1} \cdots \mathrm{~A}_{\mathrm{i}-1}\) 中的某些元素经过推理规则直接得到，则称 A 是该系统的定理，一个系统所有定理构成的集合称为定理集

因此，一个数学的系统应当配备一个公理集，该系统具有一个定理集。下面，我们围绕公理和证明这两个角度，带大家看一看数学究竟在做什么

\section*{8．1．2 一个著名的例子一一欧式几何与非欧几何}
大家从初中开始学习平面几何，都听说过欧几里得和他的《几何原本》，欧几里得把平面几何中的推理形式化成为了演绎推理，给出了一些对象，例如点，直线，平面，角的定义，提出了重要的 5 条公设（也属于公理）\\
1．由任意一点到另外任意一点可以作直线\\
2．一条有限直线可以继续延长\\
3．以任意点为心，以及任意距离可以画圆\\
4．凡直角都彼此相等\\
5．同平面内一条直线和另外两条直线相交，若在某一侧两个内角的和小于二直角的和，则这二直线经过无限延长之后在这一侧相交\\
以及一些基本的公理，例如相等可以传递，等量加等量还是等量，整体大于部分等等\\
欧几里得从这些公理公设出发，通过严格的演绎推理证明了大量的平面几何的命题。这时候，人们便有了一个问题：

5 条公设中，其他 4 条都是十分自然且基本的，但是唯独第 5 条比较特别可以证明：第五公设等价于：过直线外一点有且仅有一条直线与已知直线平行它比较长，叙述繁琐，而且其表达的事实看上去与我们的经验一致，但是又没那么显然：凭什么说一直延伸就没有交点？

它在逻辑上等价于：两直线平行，则同位角相等\\
事实上，通过前 4 条公设（不使用第五公设），欧几里得证明了：同位角相等，两直线平行。人们自然而然想问：第五公设可不可以去掉？也就是：通过前四条公设以及其他的定义，公理，能否在逻辑上证明第五公设？注 这个问题是有价值的，这牵扯到一个数学系统的根基——公理集的独立性与简洁性

大家觉得可以吗？从直觉上来讲，第五公设叙述的事实是显然的。但是，在逻辑上（也就是我们前面讲到的

逻辑推理）能否独立的证明出来呢？

\section*{焦点：直觉上的真理 VS 逻辑上的真理}
事实上，这个问题曾困扰了数学家们 2000 年在过去，数学家们大多相信是可以的，于是投人了不少精力在论证上历史上，出现了很多个证明，但是全部被检查出了错误其中还包括大数学家勒让德，他曾经给出过一个精巧的证明，但是后来他自己发现，证明过程隐含使用了和第五公设等价的命题

这个问题一直没有进展，于是一些数学家便怀疑这个结论是否是错的，也就是：第五公设不可能由其他公理证明得到

其中有位数学家叫罗巴切夫斯基，他把第五公设改成了：过直线外一点有无数条直线和已知直线无交点，并以此演绎推理了一系列与我们日常生活经验大相径庭的定理。如：\\
例 8.1 三角形内角和小于 \(\pi\) ，所有三角形的面积不超过 \(\pi\) ，三个内角相同的三角形全等，圆周长与半径不成正比．．．．．．

这些事情看上去十分难以接受，以至于连罗巴切夫斯基自己都称这个几何为＂伪几何＂因此，问题的关键变成了：能否找到一个模型，符合罗巴切夫斯基几何的公理？

这里有一个关键的概念，就是模型：\\
欧几里得叙述平面几何的概念和公理时，并没有强调生活中的具体平面，而是抽象的谈论点，线，角等等，因此，他的概念可以放在抽象的世界里谈论——就是集合

我们可以取一个集合，把它称为平面，里面的元素叫做点，里面的一些子集叫做直线，每两个点组成的点对赋予一个正数，叫做距离

只要它们的选取完全符合欧几里得的概念和公理，就可以称之为一个平面几何！\\
事实上，我们也可以理解，欧几里得的证明是逻辑上的证明，也就是演绎推理，也就是说，即使我们换了一个集合，只要推理的过程是严密的（符合我们之前的定义），那么换了集合（整了个新集合当平面），那些定理仍然应该是成立的！

\section*{定义 8.2 （（解䉽与模型）．）}
对于一个数学系统，我们可以把它的所有概念搬到一个具体的集合上去，称之为这个系统的解释；如果这个解释满足这个数学系统全部的概念和公理，那么称这个解释是该数学系统的模型

因此，如果能够构造一个罗巴切夫斯基几何的模型，那么就能够说明，第五公设不可能被其他公理推出（否则，任何的去掉第五公设的几何模型都应该满足第五公设才对）

结果，还真给构造出来了！\\
最早的模型是贝尔特拉米给出的，他在一种特殊的曲面（也就是伪球面）上定义了点和直线，以及距离，满足欧几里得的其他公理以及罗巴切夫斯基的＂第五公设＂，不过这个模型只是一个局部模型，不是一个整体模型

后来，克莱因和庞加莱分别在单位圆盘 \(\mathrm{x}^{2}+\mathrm{y}^{2}<1\) 上定义了直线，以及相应的距离（只不过克莱因定义的直线就是我们一般认知的线段，而庞加莱定义的直线是我们一般认知的圆弧，当然，尽管克莱因定义的＂线＂和我们的一般认知一样，上面的距离却截然不同，具体如何定义的目前大家理解起来可能有点麻烦），使得满足我们之前的全部要求，这宣告了第五公设相对于欧几里得其他公理的独立性

\section*{8．1．3 形式系统与公理化}
从刚才这个故事，我们可以看出很深刻的道理——我们直觉上的显然与推理，当真正碰上严格的数学推理的时候，是有可能失效的！真正研究数学，必须保证证明过程可以用演绎推理的形式写出来！

另一方面，我们在刚才的故事中引人了解释和模型的概念，实际上，在数学的基础——逻辑中，研究的更多是形式上的，抽象的系统，当我们赋予了具体的解释后，也就是用形式上的系统来解释数学中有价值的系统，用具有数学意义的模型来给予形式系统发挥逻辑推理作用的平台。\\
例 8.2 欧几里得的平面几何公理系统，后来，希尔伯特在 《几何基础》中提出的公理系统弥补了欧几里得的系统不是特别严谨的缺点\\
例 8.3 朴素算术（自然数的加法与乘法）的形式系统——皮亚诺公理系统\\
一个形式系统要承担它演绎推理，探求真理的价值，就要赋予一个公理集，公理集中的元素是一些指派好的命题\\
例8．4 欧几里得的形式系统里的公理集就是 5 条公设和一些常识性的与数量相关的公理\\
还要赋予一个规则集，表示演绎推理所允许的规则\\
例 8.5 最基本的分离规则，一阶逻辑中的概括规则\\
例 8.6 以自然数为基础的归纳规则（数学归纳法）\\
公理化方法在现在的数学中普遍使用，是建立数学理论的重要方法\\
例 8.7 集合论的形式系统（形式集论），公理集为 ZF 体系的 8 条公理（下一期会讲一下）\\
例 8.8 线性空间的形式系统，公理集为有关加法和数乘的 8 条公理\\
例 8.9 拓扑学的形式系统，公理集为关于开集的 3 条公理

\section*{8．1．4 一致性问题}
\begin{itemize}
  \item 致性是逻辑中十分重要的概念
  \item 个形式系统，它未必就是好的，比如考虑一个形式系统，它的公理集包含这样两条公理：
\end{itemize}

\section*{A 对}
A 错\\
这不是扯淡吗！！这样的形式系统本身就自相矛盾，而自相矛盾的东西，我们数学不欢迎！数学是探求真理的，既然要探求真理，那就要保证你的理论不能有矛盾！

\section*{定义8．3（（一致性）．）}
形式系统称为一致的，若不存在命题 A ，使得 A 和 \(\neg \mathrm{A}\) 都是该系统的定理

一致性也是看上去极为显然的，我们平时使用的算术，集合，代数这些东西，似乎都是一致的，不可能自相矛盾。事实上真的是这样吗？

以大家高中学习的集合论为例\\
高中的集合论中，把一些具有某一类性质的对象放在一起，称之为集合（这句话似乎就可以拿来当一个公理），似乎非常的直观

集合要求了确定性，互异性等特点，规定了并，交等运算。看上去非常正常，不就考虑一些对象吗，这肯定是一致吧，这不是废话吗？但是！请看这样一个集合：\(R=\{x\) 是集合 \(\mid \mathrm{x} \notin \mathrm{x}\}\) 按照我们的直观，这个 \(R\) 符合集合

的要求，当然算是个集合。但是问题来了：命题 \(R \in R\) 为真还是为假？如果 \(R \in R\) 为真，那么 \(R\) 本身要满足 \(R\)中元素的要求，那么当然需要 \(R \notin R\) ；反过来，如果 \(R \notin R\) ，那么 \(R\) 不在 \(R\) 中，自然不满足 \(R\) 中元素的性质，那就有 \(\mathrm{R} \in \mathrm{R}\) ，这么一来．．．．．．好像出现矛盾了吧

事实上，根据基本的演绎推理，我们是可以证明 \(R \in R\) ，也可以证明 \(R \notin R\) ．因此，就出现了这样的悖论：

\section*{命题 8.1 （（罗素悖论）．）}
在满足一致性的集合论中，集合 \(R=\{x\) 是集合 \(\mid x \notin x\}\) 不存在

证明 反设存在，于是我们可以证明：\\
\(R \in R\) 与 \(R \notin R\) 都是该集合论的定理这矛盾于一致性\\
因此，朴素的集合论就不满足一致性，这样的东西当然不能存在于数学研究中，因此，真正的公理化集合论是对公理增加了一些限制的

从这里大家可以看出，保持体系的一致性对于数学研究是非常重要的。事实上，真正的数学研究经常需要对现有的理论体系进行推广，一个大家在高中最熟悉的例子便是实数到复数的扩充

问题在于，这样的扩充能保持一致性吗？虚数这种＂虚无飘渺＂的东西，看着就不像好人，造出这么丑陋的东西来，数学家良心不会痛吗？

实际上，复数这个东西，虽然违反直觉，但是在数学上是不违反什么东西的，实际上，复数好比在平面坐标系 \(\mathrm{xOy}:\{(\mathrm{x}, \mathrm{y}) \mid \mathrm{x}, \mathrm{y} \in \mathrm{R}\}\) 上定义了加减乘除，加减同向量加减，乘除则是 \((\mathrm{a}, \mathrm{b}) \cdot(\mathrm{c}, \mathrm{d})=(\mathrm{ac}-\mathrm{bd}, \mathrm{ad}+\mathrm{bc})\)（相当于 \((a+b i)(c+d i)=(a c-b d)+(a d+b c) i)\) 除法类似

如果单纯看这样的定义，就是在一个集合 xOy 上定义了新的四则运算，完全没毛病啊！而且，这样定义的复数还是很有价值的，它是多项式 \(x^{2}+1\) 的分裂域，具有良好的拓扑结构，其上定义的微积分还具有相当优越的性质．．．．．．

但是，胡乱创造一套理论，并且闭门造车自己瞎推理，却很容易翻车，不仅很可能出现演绎推理上的漏洞，还有可能，自己造的理论本身连＂符合直觉的一致性＂的基本要求都显然无法通过，甚至直接就可以证明它不一致

\section*{8．1．5 有关＂数学＂}
刚才我们谈了这么多，回到最原始的问题——有关数学\\
真正的数学，就是根据真理探索的需求，建立了一个个形式系统，规定了一系列公理，在一个个模型上建立演绎推理的大厦，得到一系列解构自然的真理，每一种不同的形式系统就对应不同的数学分支

不同的形式系统可能会有大小包含的关系，例如比较基本的自然数算术系统，在此基础上引进一批新的符号（负数），把加减法和乘法的概念推广到上面去，就得到了整数的算术

在整数的基础上，引进一批新的符号（分数），定义新的运算（除法）等，就可以建立有理数算术．．．．．．\\
例8．10 加法\\
我们可以在很多地方看到加法，比如数的加法，式的加法，向量加法．．．．．．\\
解 我们定义一个形式系统，考虑一个集合 \(A\) ，对于任意一组 \(a, b \in A\) ，我们规定一种运算 \(a+b\) ，使得存在唯一的 \(c \in A c=a+b\)（集合元素的互异性保证我们可以定义等号）\\
我们给出如下下列公理：\\
1）加法交换律： \(\mathrm{a}+\mathrm{b}=\mathrm{b}+\mathrm{a}, \forall \mathrm{a}, \mathrm{b} \in \mathrm{A}\)

2）加法结合律：\((a+b)+C=a+(b+c), \forall a, b, c \in A\)\\
3）存在零元：存在一个 \(0 \in A\) ，使得 \(\forall a \in A a+0=0+a=a\)\\
4）存在相反数：\(\forall \mathrm{a} \in \mathrm{A}, \exists \mathrm{b} \in \mathrm{A}, \mathrm{a}+\mathrm{b}=\mathrm{b}+\mathrm{a}=0\)\\
如果满足上述公理1）2）3），则称 A 上定义了加法运算（加法半群），如果还满足 4），则称 A 上定义了加法群

\section*{定理 8.2}
A 上若定义了加法运算，则 A 中的零元唯一

证明 考虑 A 中的两个零元（未必互异） \(0_{1}, 0_{2}\) ，由公理3）

\[
0_{1}+0_{2}=0_{1}, 0_{1}+0_{2}=0_{2},
\]

于是 \(0_{1}=0_{2}\)\\
从而我们可以看出，公理3）就能证明零元的唯一性（逻辑的强大），而刚才的这个内容，其实就是数学的一个极为简单的例子

下面，回答一下最开始的几个问题\\
1．学数学的话，就要整天做好多繁杂的计算，做好多需要高智商的难题\\
答：不是这样，可能学分析领域，搞偏微分方程的计算会多一些，但是像代数领域，拓扑领域这些，与其说计算，思考一些结构性的问题，发挥想象力会很多，反而计算很少

2．数学好的人，一定非常聪明，反应速度很快\\
答：聪明的人学数学有优势，但是即使没那么聪明，只要对一些结构有着深刻的理解，热爱思考，见多识广（指阅读当今成果），同样可以做下去

\section*{3．数学好的人，计算能力一定很强}
答：与其说学数学计算厉害，不如说学工科计算厉害（

\section*{4．数学家都是相当聪明的人，都是通过各种奇思妙想的技巧来解决世界难题}
答：数学的发展也是需要经历长时间的铺垫，不排除有天才靠奇思妙想的能力推动数学发展，但是更多的还是一步一个脚印，无数数学研究者不断铺垫得到的发展，新的形式系统的建立，发现某个领域的内容与另一个领域的联系，通过不断试错，得到一个新的定理的证明．．．．．这些都需要时间的沉淀以及无数数学工作者的努力

\section*{8．1．6 不完全性定理}
\section*{疑问：添加了一些限制之后，一致性就解决了吗？}
\section*{疑问：既然现在都用形式系统以及演绎推理的方法研究数学，那么这个东西本身是否强大？}
也就是说（以自然数的算术为例），我们建立皮亚诺公理系统来研究算术，那么是不是所有的朴素算术的真理都可以被形式上的逻辑推理证明出来呢？

遗憾的是，哥德尔在上个世纪证明了不完全性定理

\section*{命题 8.2 （（ \(\omega\) 一致性）））}
若一个形式系统中不存在与自然数 \(n\) 有关的命题 \(A(n)\) ，使得 \(\forall n \in N, A(n)\) 都是该系统的定理，且 \(\neg(\forall \mathrm{n}) \mathrm{A}(\mathrm{n})\) 也是该系统的定理

\section*{定理 8.3 （（不完金性定理）．）}
若皮亚诺公理系统是 \(\omega\) —致的，则存在一个朴素算术命题为真，但是它不是皮亚诺公理系统的定理\\
注 \(\omega\) —致性是必要的，为了保证其合适公理集递归，这个太难，就不展开了\\
另外，任何满足 \(\omega\) 一致性的系统都有不完全性定理\\
而且关于一致性，事实上，哥德尔第二不完全性定理告诉我们：

\section*{定现 8.4 （（第二不完金性定理）．）}
包含算术系统的形式系统自身不能证明它本身的一致性

目前为止，一致性获得证明的形式系统，有命题逻辑，有一阶逻辑，但是绝大部分，甚至包括算术的形式系统（皮亚诺公理系统），以及 ZF 公理化集合论，人类都还没能严格证明其一致性，更不用说在此基础上衍生出的实数理论等一系列数学理论了

感觉上的一致性不能取代严格的证明！这是个极为艰难的问题，至今为止人类还没有找到特别好的方法来解决

那怎么办？没有一致性的保障，数学难道就停滞不前？数学家们也做出了妥协，先继续研究，哪天出来新的漏洞（罗素悖论二号）再接着添加新公理，排除悖论（事实上，这是以布尔巴基学派为代表的结构主义数学的观点，现在作为数学界普遍接受这种观点作为主流）

在人类等到逻辑发展到可以解决这些问题之前，让数学继续发展吧！（暂时默认这些都是一致的，慢慢接着研究）

\section*{8．1．7 总结：数学的基础}
逻辑是数学的基础，数学通过建立各种形式系统来研究各个数学分支，各种数学对象在演绎推理的基础上向人类展示丰富的定理（客观真理）

一个形式系统有一系列公理，这些公理决定了一个数学分支的根基，在符合研究需求的前提下，修改或扩大公理集，将会提升一个领域探求真理的能力

但是，由于不完全性定理，我们通过演绎推理很难穷尽所有真理，就比如集合论的选择公理，添加了这个之后，集合论便具有了更加强大的威力，可以构造出相当多的丰富的对象，但是像连续统假设，目前的集合论还没有办法证明或否定这个命题（集合论的这些内容，下一次视频会给大家简单科普一下）

但是，在我心目中，尽管演绎推理无法穷尽客观真理，但是它告诉我们，数学研究的道路上还有无限风光，不断追寻真理，是一件无比浪漫的事情

最后给大家留下两个思考题（高考一定不会考！！），感兴趣的同学可以留作思考\\
1．视频中提到的加法群（是满足那 4 条公理的），请尝试证明：若 A 上定义了加法群，则对于 \(\forall \mathrm{a} \in \mathrm{A}, \mathrm{a}\) 的相反数是唯一的\\
2．我们考虑集合 \(R^{+}=\{x \in R \mid x>0\}\) ，我们定义加法 \(x \oplus y=x y\)（指＂乘＂为＂加＂），请验证 \(\mathrm{R}^{+}\)关于这个＂加法＂构成加法群，这个加法群的＂零元＂是谁？某一个数的＂相反数＂又是谁？

\section*{8.2 重要的概念}
事前声明：本次内容主要为集合论的基础知识，属于大学数学的入门衔接知识，与高考没有关系。本讲内容面向对数学有兴趣的观众。为了方便高中生理解，内容中的部分定义和表述进行了直观化的简化处理，因此可能部分地方不够严谨，感兴趣的同学可以自行查阅大学的集合论相关内容。

凡是刚上高一的同学，都要学习集合，但是，大家是否会有一种＂多此一举＂的感觉？也就是：为什么要引入集合的概念？我在初中写 \(0<x<1\) 写的挺舒适的，上了高中整出那么多集合来，写成 \(x \in(0,1)\) ，何必呢？数学家何必搞出这么个东西来呢？

今天，为大家带来有关集合论的入门级科普，旨在让大家理解集合这一概念的必要性，以及了解集合论公理化的必要性。本次视频主要面向对数学有兴趣的观众。

\section*{8．2．1 数学基础的巩固}
\section*{8．2．2 第三次数学危机与集合论公理化}
\section*{8．2．3 一些集合论的基础知识}
\section*{8．2．4 总结：重要的概念}
\section*{8.3 联系万物}
事前声明：本次内容主要为映射的基础知识，属于大学数学的入门衔接知识，与高考没有任何关系。本讲内容面向对数学有兴趣的观众，前置知识为集合论，看过上一期集合的科普视频的话，本期内容会更好理解，不过即使只有高中的集合知识储备，同样可以观看。为了方便高中生理解，内容中的部分定义和表述进行了直观化的简化处理，因此可能部分地方不够严谨，感兴趣的同学可以自行查阅大学数学中映射的相关内容。

另外，本次内容相对难度会稍微大一点，因为涉及到的逻辑推理比较多，相信也能够提高部分想要数学拔尖的同学的逻辑分析能力和处理信息的能力。本讲最后设置了一些比较有挑战性的思考题，偏向新信息处理和数学探究，提供给对数学有兴趣的同志们。

集合论是当今数学的基本，有了集合论，我们才能在各式各样的空间中建立丰富多彩的数学分支。集合的基本思想就是把具有相同的某种性质的对象放在一起来研究，同时，用集合的语言来描述空间也是十分自然，合理的。当我们在数学中谈到＂空间＂一词时，无一例外，都是具备了某些结构或性质的集合。例如二维欧式空间 \(R^{2}=\{(x, y) \mid x, y \in R\}\) ，就是我们通常认知的平面世界的数学语言刻画，在上面定义距离，直线等结构，引入平行，垂直，夹角等概念，就是我们熟悉的平面几何。

不过，数学可不止于此。光是在各个集合上面＂闭门造车＂，各研究各的，可不会让数学家们满意。数学的美感很大程度上来源于不同事物之间千丝万缕的联系。比如，圆和椭圆看上去形状差不多，都是个圈，但是直线和圆就不一样，因此，从结构上来看，圆和椭圆可以归为一类，直线可以归为一类，这两种＂空间＂的结构就完全不同。

因此，研究不同集合之间的关系，是数学中十分重要且基本的内容。而不同集合要建立关系，搭上钩，映射就是一个十分自然的概念了。

8．3．1 从对应关系到咉射\\
8．3．2 单射，满射与双射\\
8．3．3 映射的性质\\
8．3．4 形形色色的映射\\
8．3．5 总结：联系万物

\section*{8.4 优美的仿射几何}
\section*{8.5 实数理论}

\end{document}